\documentclass{article}

\usepackage[utf8]{inputenc}
\usepackage[T1]{fontenc}
\usepackage{helvet}
\usepackage{geometry}
\usepackage{xcolor}
\usepackage{eso-pic}
\usepackage{fancyhdr}
\usepackage{parskip}
\usepackage{microtype}
\usepackage[english, ukrainian]{babel}
\usepackage{url}
\usepackage{amsmath}
\usepackage{amssymb}
\usepackage{amsthm}
\usepackage{longtable}
\usepackage{booktabs}
\usepackage{hyperref}
\usepackage{titlesec}
\usepackage{tocloft}

\renewcommand{\cftsecfont}{\small\bfseries}
\renewcommand{\cftsubsecfont}{\footnotesize}
\renewcommand{\cftsubsubsecfont}{\scriptsize}
\renewcommand{\cftsecpagefont}{\small\bfseries}
\renewcommand{\cftsubsecpagefont}{\footnotesize}
\renewcommand{\cftsubsubsecpagefont}{\scriptsize}

\definecolor{nato}{HTML}{293757}
\definecolor{footergray}{gray}{0.65}

\geometry{a4paper,left=3cm,right=3cm,top=3cm,bottom=3cm}
\renewcommand{\familydefault}{\sfdefault}
\setcounter{secnumdepth}{3}

\titleformat{\section}
  {\color{nato}\Huge\bfseries\raggedright}{\thesection.}{1em}{}
\titlespacing*{\section}{0pt}{30pt}{12pt}

\titleformat{\subsection}
  {\color{nato}\LARGE\bfseries\raggedright}{\thesubsection.}{1em}{}
\titlespacing*{\subsection}{0pt}{20pt}{8pt}

\titleformat{\subsubsection}
  {\color{nato}\Large\bfseries\raggedright}{\thesubsubsection.}{1em}{}
\titlespacing*{\subsubsection}{0pt}{10pt}{6pt}

\fancyhf{}
\fancyhead[R]{\small\color{footergray} 2 червня 2026}
\fancyfoot[L]{\small\color{footergray} Комплексна система захисту інформації. Повний каталог заходів із захисту інформації (1003).}
\fancyfoot[R]{\small\color{footergray}\thepage}

\newcommand\HUGE[1]{\fontsize{40}{40}\selectfont #1}

\renewcommand{\headrulewidth}{0pt}
\renewcommand{\footrulewidth}{0.4pt}
\renewcommand{\footrule}{\color{footergray}\hrule width \textwidth height 0.4pt}

\begin{document}

\pagestyle{fancy}

\begin{titlepage}
  \AddToShipoutPictureBG*{\AtPageLowerLeft{\color{nato}\rule{\paperwidth}{\paperheight}}}
  \color{white}\thispagestyle{empty}\vspace*{2cm}\centering
  {\Huge  \textbf{Комплексна система захисту інформації} \par} \vspace{1cm}
  {\HUGE  \textbf{Повний каталог заходів із захисту інформації (1003)} \par} \vspace{1cm}
  \vfill
  {\large 2026 \copyright\ Критпографічні Телесистеми \par}
\end{titlepage}

\newpage

\begin{center}
  {\Huge\color{nato}\bfseries Зміст\par}
\end{center}
\vspace{40pt}
\tableofcontents

\begin{abstract}
Повний каталог (Bible) — містить абсолютно всі заходи, посилення та субконтролі згідно з повним скоупом НД ТЗІ КСЗІ.
\end{abstract}

\section{AC}
Клас заходів захисту AC --- УПРАВЛІННЯ ДОСТУПОМ

Цей клас зосереджується на обмеженні доступу до інформаційних систем, ресурсів та функцій лише для авторизованих користувачів, програм і пристроїв.

\textbf{Перелік заходів захисту:} Політика та процедури управління доступом (AC-1); Управління обліковими записами (AC-2); Автоматизоване управління обліковими записами системи (AC-2(1)); Видалення тимчасових та екстрених облікових записів (AC-2(2)); Деактивація облікових записів (AC-2(3)); Дії при автоматизованому аудиті (AC-2(4)); Вихід із системи за відсутності активності (AC-2(5)); Динамічне управління привілеями (AC-2(6)); Схеми, засновані на ролях (AC-2(7)); Динамічне управління обліковими записами (AC-2(8)); Обмеження на використання спільних та групових облікових записів (AC-2(9)); Зміна даних спільних і групових облікових записів (AC-2(10)) [Вилучено]; Умови використання (AC-2(11)); Моніторинг нетипового використання облікових записів (AC-2(12)); Деактивація облікових записів осіб з високим рівнем ризику (AC-2(13)); Забезпечення доступу (AC-3); Обмежений доступ до привілейованих функцій (AC-3(1)) [Вилучено]; Подвійна авторизація (AC-3(2)); Мандатне управління доступом (AC-3(3)); Дискреційне управління доступом (AC-3(4)); Інформація щодо безпеки (AC-3(5)); Захист інформації користувача та системи (AC-3(6)) [Вилучено]; Управління доступом на основі ролей (AC-3(7)); Анулювання прав доступу (AC-3(8)); КЕРОВАНА ПЕРЕДАЧА (ПУБЛІКАЦІЯ) ІНФОРМАЦІЇ (AC-3(9)); Перегляд аудитом механізмів контролю доступу (AC-3(10)); Обмеження доступу до спеціальної інформації (AC-3(11)); Встановлення та забезпечення доступу до застосунків (AC-3(12)); Управління доступом на основі атрибутів (AC-3(13)); Індивідуальний доступ (AC-3(14)); Дискреційний та обов'язковий доступ (AC-3(15)); Управління інформаційними потоками (AC-4); Атрибути безпеки об'єкту (AC-4(1)); Домени обробробки даних (AC-4(2)); Динамічне управління інформаційним потоком (AC-4(3)); Управління потоком зашифрованої інформації (AC-4(4)); Вбудовування типів даних (AC-4(5)); Метадані (AC-4(6)); Механізми одностороннього потоку (AC-4(7)); Фільтри політики безпеки (AC-4(8)); Перевірки, що проводить персонал (AC-4(9)); Активація та деактивація фільтрів політики безпеки (AC-4(10)); Конфігурація фільтрів політики безпеки (AC-4(11)); Ідентифікатори типу даних (AC-4(12)); Декомпозиція на відповідні політиці субкомпоненти (AC-4(13)); Обмеження фільтра політики безпеки (AC-4(14)); Виявлення несанкціонованої інформації (AC-4(15)); Передача інформації про взаємопов'язані системи (AC-4(16)) [Вилучено]; Автентифікація домену (AC-4(17)); Прив'язка атрибуту безпеки (AC-4(18)) [Вилучено]; Перевірка метаданих (AC-4(19)); Затверджені рішення (AC-4(20)); Фізичне та логічне відділення інформаційних потоків (AC-4(21)); Єдиний доступ (AC-4(22)); Модифікована інформація, яка не підлягає оприлюдненню (AC-4(23)); Внутрішній нормалізований формат (AC-4(24)); Очищення даних (AC-4(25)); Дії з фільтрації аудиту (AC-4(26)); Надлишкові/незалежні фільтруючі механізми (AC-4(27)); Лінійні фільтрувальні канали (AC-4(28)); Фільтр механізмів оркестровки (AC-4(29)); Механізми фільтрації з використанням кількох процесів (AC-4(30)); Запобігання спробам передачі вмісту, який не пройшов перевірку фільтрації (AC-4(31)); Вимоги до процесу передачі інформації (AC-4(32)); Розмежування обов'язків (AC-5); Мінімізація повноважень (AC-6); Авторизований доступ до функцій безпеки (AC-6(1)); Непривілейований доступ до незахищених функцій (AC-6(2)); Мережевий доступ до привілейованих команд (AC-6(3)); Роздільні домени обробки (AC-6(4)); Привілейовані облікові записи (AC-6(5)); Привілейований доступ користувачами, що не належать до організації (AC-6(6)); Перегляд повноважень користувача (AC-6(7)); Рівні привілеїв для виконання коду (AC-6(8)); Аудит використання привілейованих функцій (AC-6(9)); Заборона непривілейованим користувачам виконувати привілейовані функції (AC-6(10)); Невдалі спроби входу в систему (AC-7); Автоматичне блокування облікового запису (AC-7(1)) [Вилучено]; Очищення або стирання мобільного пристрою (AC-7(2)); Обмеження на спроби біометричного входу (AC-7(3)); Використання альтернативного фактора (AC-7(4)); Попередження про використання системи (AC-8); СПОВІЩЕННЯ ПРО ПОПЕРЕДНІЙ ВХІД (ДОСТУП) (AC-9); Невдалі спроби входу до системи (AC-9(1)); Успішні та невдалі спроби входу до системи (AC-9(2)); Повідомлення про зміни в обліковому записі (AC-9(3)); СПОВІЩЕННЯ ПРО ПОПЕРЕДНІЙ ВХІД (ДОСТУП) -- ДОДАТКОВА ІНФОРМАЦІЯ ПРО ВХІД (AC-9(4)); Управління паралельною сесією (AC-10); Блокування пристрою (AC-11); Приховані дисплеї (AC-11(1)); Припинення сеансу (AC-12); Ініційоване користувачем блокування (AC-12(1)); Повідомлення про припинення сеансу (AC-12(2)); Застережне повідомлення про те, що час сесії добігає кінця (AC-12(3)); Управління доступом (AC-13) [Вилучено]; Дозволені дії без ідентифікації або автентифікації (AC-14); Дозволені дії без ідентифікації необхідне використання (AC-14(1)) [Вилучено]; Автоматизоване маркування (AC-15) [Вилучено]; Атрибути безпеки та приватності (AC-16); Динамічне пов'язання атрибутів (AC-16(1)); Зміна значень атрибутів авторизованими особами (AC-16(2)); Підтримка системою пов'язання атрибутів (AC-16(3)); Пов'язання атрибутів авторизованими особами (AC-16(4)); Відображення атрибутів на пристроях виведення (AC-16(5)); Підтримка пов'язання атрибутів організацією (AC-16(6)); Послідовна інтерпретація атрибутів (AC-16(7)); Техніки та технології пов'язання атрибутів (AC-16(8)); Перепризначення атрибутів (AC-16(9)); Конфігурація атрибутів уповноваженими особами (AC-16(10)); Віддалений доступ (AC-17); Автоматизований моніторинг та управління (AC-17(1)); Захист конфіденційності та цілісності за допомогою шифрування (AC-17(2)); Керовані точки контролю доступу (AC-17(3)); Привілейовані команди та доступ (AC-17(4)); Моніторинг для неавторизованих підключень (AC-17(5)) [Вилучено]; Захист інформації (AC-17(6)); Додатковий захист для доступу до функцій безпеки (AC-17(7)) [Вилучено]; Деактивація незахищених протоколів мережі (AC-17(8)) [Вилучено]; Відключення або деактивація доступу (AC-17(9)); (10) автентифікація віддалених команд (AC-17(10)); Бездротовий доступ (AC-18); Автентифікація та шифрування (AC-18(1)); Бездротовий підключень (AC-18(2)) [Вилучено]; Відключення бездротової мережі (AC-18(3)); Бездротовий доступ стувачами (AC-18(4)); Антени та рівень потужності передачі (AC-18(5)); Контроль доступу для мобільних пристроїв (AC-19); Використання письмових та портативний пристроїв для зберігання даних (AC-19(1)) [Вилучено]; Використання персональних портативних пристроїв зберігання даних (AC-19(2)) [Вилучено]; Використання портативних пристроїв зберігання даних з неідентифікованим власником (AC-19(3)); Обмеження для засекреченої інформації (AC-19(4)); Повне шифрування пристроїв та сховищ інформації (AC-19(5)); Використання зовнішніх систем (AC-20); Обмеження на авторизоване використання (AC-20(1)); Переносні пристрої зберігання даних (AC-20(2)); Системи та компоненти, що не знаходяться у власності організації (AC-20(3)); Пристрої для зберігання даних, які можуть мати доступ до мережі (AC-20(4)); Портативні пристрої для зберігання даних -- заборона використання (AC-20(5)); Розповсюдження інформації (AC-21); Автоматична підтримка ухвалення рішень (AC-21(1)); Розповсюдження інформації (AC-21(2)); Публічно доступний контент (AC-22); Захист від несанкціонованого інтелектуального аналізу даних (AC-23); Рішення щодо управління доступом (AC-24); Інформація про передачу авторизованого доступу (AC-24(1)); Відсутність ідентифікації користувача або процесу, що діє від імені користувача (AC-24(2)); Диспетчер доступу (AC-25).

\subsection{ПОЛІТИКА ТА ПРОЦЕДУРИ УПРАВЛІННЯ ДОСТУПОМ (AC-1)}
a. Розробити, задокументувати та поширити [Призначення: серед визначеного організацією персоналу або ролей]:\\ 
1. 2. [Вибір (один або декілька): Рівень організації; Рівень місії/бізнес-процесу; рівень системи] політики контролю доступу, яка:\\ 
(a) містить мету, сферу застосування, ролі, відповідальність, зобов'язання керівництва, координацію між організаційними підрозділами та систему контролю відповідності (compliances);\\ 
(b) відповідає чинному законодавству, нормативним документам, директивам, нормам, політикам, стандартам і керівним документам. Процедури, що сприяють реалізації політики управління доступом і відповідних заходів управління доступом.\\ 
b. Призначити на посаду [Призначення: визначену організацією посадову особу] для управління, документування і розповсюдження політики та процедур контролю доступом.\\ 
c. Переглянути та оновити:\\ 
1. поточну політику управління доступом [Призначення: з визначеною організацією частотою] та [Призначення: події, визначені організацією];\\ 
2. поточні процедури управління доступом [Призначення: з визначеною організацією частотою] та [Завдання: події, визначені організацією].\\ 


\vspace{10pt}
{\footnotesize\bfseries
No: 1\\
Name: ac\_\allowbreak{}1\_\allowbreak{}odp\_\allowbreak{}01\\
Type: list\\
Default: [\textquotedbl{}admin\textquotedbl{}, \textquotedbl{}security\_\allowbreak{}officer\textquotedbl{}]
}

{\footnotesize Визначено персонал або ролі, на які поширюється політика контролю доступу}


{\footnotesize\bfseries
No: 2\\
Name: ac\_\allowbreak{}1\_\allowbreak{}odp\_\allowbreak{}02\\
Type: list\\
Default: [\textquotedbl{}admin\textquotedbl{}, \textquotedbl{}security\_\allowbreak{}officer\textquotedbl{}]
}

{\footnotesize Визначено персонал або ролі, на які поширюються процедури контролю доступу}


{\footnotesize\bfseries
No: 3\\
Name: ac\_\allowbreak{}1\_\allowbreak{}odp\_\allowbreak{}03\\
Type: string\\
Default: nil
}

{\footnotesize Вибрано одне або декілька з наступних ЗНАЧЕНЬ ПАРАМЕТРІВ: \{рівень організації; рівень місії/бізнес-процесу; рівень системи\}}


{\footnotesize\bfseries
No: 4\\
Name: ac\_\allowbreak{}1\_\allowbreak{}odp\_\allowbreak{}04\\
Type: list\\
Default: [\textquotedbl{}admin\textquotedbl{}, \textquotedbl{}security\_\allowbreak{}officer\textquotedbl{}]
}

{\footnotesize Визначено посадову особу, яка керуватиме політикою та процедурами контролю доступу}


{\footnotesize\bfseries
No: 5\\
Name: ac\_\allowbreak{}1\_\allowbreak{}odp\_\allowbreak{}05\\
Type: list\\
Default: [\textquotedbl{}default\_\allowbreak{}deny\_\allowbreak{}rule\textquotedbl{}, \textquotedbl{}abac\_\allowbreak{}rule\_\allowbreak{}1\textquotedbl{}]
}

{\footnotesize Визначено частоту, з якою переглядається та оновлюється поточна політика контролю доступу}


{\footnotesize\bfseries
No: 6\\
Name: ac\_\allowbreak{}1\_\allowbreak{}odp\_\allowbreak{}06\\
Type: list\\
Default: [\textquotedbl{}default\_\allowbreak{}deny\_\allowbreak{}rule\textquotedbl{}, \textquotedbl{}abac\_\allowbreak{}rule\_\allowbreak{}1\textquotedbl{}]
}

{\footnotesize Визначено події, які вимагають перегляду та оновлення поточної політики контролю доступу}


{\footnotesize\bfseries
No: 7\\
Name: ac\_\allowbreak{}1\_\allowbreak{}odp\_\allowbreak{}07\\
Type: integer\\
Default: 30
}

{\footnotesize Визначено частоту, з якою переглядаються та оновлюються поточні процедури контролю доступу}


{\footnotesize\bfseries
No: 8\\
Name: ac\_\allowbreak{}1\_\allowbreak{}odp\_\allowbreak{}08\\
Type: integer\\
Default: 30
}

{\footnotesize Визначено події, які потребують перегляду та оновлення процедур}


{\footnotesize\bfseries
No: 9\\
Name: ac\_\allowbreak{}1\_\allowbreak{}a\_\allowbreak{}01\\
Type: integer\\
Default: 30
}

{\footnotesize Розроблено та задокументовано політику контролю доступу}


{\footnotesize\bfseries
No: 10\\
Name: ac\_\allowbreak{}1\_\allowbreak{}a\_\allowbreak{}02\\
Type: integer\\
Default: 30
}

{\footnotesize Політика контролю доступу поширюється на <AC-01\_\allowbreak{}ODP[01] персонал або ролі>;}


{\footnotesize\bfseries
No: 11\\
Name: ac\_\allowbreak{}1\_\allowbreak{}a\_\allowbreak{}03\\
Type: integer\\
Default: 30
}

{\footnotesize Розроблені та задокументовані процедури контролю доступу для полегшення впровадження політики контролю доступу та пов'язаних з нею заходів захисту}


{\footnotesize\bfseries
No: 12\\
Name: ac\_\allowbreak{}1\_\allowbreak{}a\_\allowbreak{}04\\
Type: integer\\
Default: 30
}

{\footnotesize Процедури контролю доступу поширюються на <AC-01\_\allowbreak{}ODP[02] персонал або ролі>}


{\footnotesize\bfseries
No: 13\\
Name: ac\_\allowbreak{}1\_\allowbreak{}a\_\allowbreak{}1\_\allowbreak{}a\_\allowbreak{}01\\
Type: integer\\
Default: 30
}

{\footnotesize Політика контролю доступу <AC-01\_\allowbreak{}ODP[03] ВИБІРКОВЕ ЗНАЧЕННЯ ПАРАМЕТРА(ів)> містить мету}


{\footnotesize\bfseries
No: 14\\
Name: ac\_\allowbreak{}1\_\allowbreak{}a\_\allowbreak{}1\_\allowbreak{}a\_\allowbreak{}02\\
Type: integer\\
Default: 30
}

{\footnotesize Політика контролю доступу <AC-01\_\allowbreak{}ODP[03] ВИБІРКОВЕ ЗНАЧЕННЯ ПАРАМЕТРА(ів)> містить сферу застосування}


{\footnotesize\bfseries
No: 15\\
Name: ac\_\allowbreak{}1\_\allowbreak{}a\_\allowbreak{}1\_\allowbreak{}a\_\allowbreak{}04\\
Type: integer\\
Default: 30
}

{\footnotesize Політика контролю доступу <AC-01\_\allowbreak{}ODP[03] ВИБІРКОВЕ ЗНАЧЕННЯ ПАРАМЕТРА(ів)> містить сферу застосування}


{\footnotesize\bfseries
No: 16\\
Name: ac\_\allowbreak{}1\_\allowbreak{}a\_\allowbreak{}1\_\allowbreak{}a\_\allowbreak{}05\\
Type: integer\\
Default: 30
}

{\footnotesize Політика контролю доступу <AC-01\_\allowbreak{}ODP[03] ВИБІРКОВЕ ЗНАЧЕННЯ ПАРАМЕТРА(ів)> містить сферу застосування}


{\footnotesize\bfseries
No: 17\\
Name: ac\_\allowbreak{}1\_\allowbreak{}a\_\allowbreak{}1\_\allowbreak{}a\_\allowbreak{}06\\
Type: integer\\
Default: 30
}

{\footnotesize Політика контролю доступу <AC-01\_\allowbreak{}ODP[03] ВИБІРКОВЕ ЗНАЧЕННЯ ПАРАМЕТРА(ів)> містить сферу застосування}


{\footnotesize\bfseries
No: 18\\
Name: ac\_\allowbreak{}1\_\allowbreak{}a\_\allowbreak{}1\_\allowbreak{}a\_\allowbreak{}07\\
Type: integer\\
Default: 30
}

{\footnotesize Політика контролю доступу <AC-01\_\allowbreak{}ODP[03] ВИБІРКОВЕ ЗНАЧЕННЯ ПАРАМЕТРА(ів)> містить сферу застосування}


{\footnotesize\bfseries
No: 19\\
Name: ac\_\allowbreak{}1\_\allowbreak{}a\_\allowbreak{}1\_\allowbreak{}b\\
Type: integer\\
Default: 30
}

{\footnotesize Політика контролю доступу <AC-01\_\allowbreak{}ODP[03] ВИБРАНЕ ЗНАЧЕННЯ ПАРАМЕТРА(ів)> відповідає чинним законам, виконавчим наказам, директивам, положенням, політикам, стандартам та настановам}


{\footnotesize\bfseries
No: 20\\
Name: ac\_\allowbreak{}1\_\allowbreak{}b\\
Type: integer\\
Default: 30
}

{\footnotesize <AC-01\_\allowbreak{}ODP[04] посадова особа> призначається для управління розробкою, документуванням та розповсюдженням політики та процедур контролю доступу}


{\footnotesize\bfseries
No: 21\\
Name: ac\_\allowbreak{}1\_\allowbreak{}c\_\allowbreak{}01\_\allowbreak{}1\\
Type: integer\\
Default: 30
}

{\footnotesize Переглядається та оновлюється поточна політика контролю доступу <AC-01\_\allowbreak{}ODP[05] частота>}


{\footnotesize\bfseries
No: 22\\
Name: ac\_\allowbreak{}1\_\allowbreak{}c\_\allowbreak{}01\_\allowbreak{}2\\
Type: integer\\
Default: 30
}

{\footnotesize Поточну політику контролю доступу переглянуто та оновлено після <AC-01\_\allowbreak{}ODP[06] подій>}


{\footnotesize\bfseries
No: 23\\
Name: ac\_\allowbreak{}1\_\allowbreak{}c\_\allowbreak{}02\_\allowbreak{}1\\
Type: integer\\
Default: 30
}

{\footnotesize Переглядаються та оновлюються поточні процедури контролю доступу <AC-01\_\allowbreak{}ODP[07] частота>}


{\footnotesize\bfseries
No: 24\\
Name: ac\_\allowbreak{}1\_\allowbreak{}c\_\allowbreak{}02\_\allowbreak{}2\\
Type: integer\\
Default: 30
}

{\footnotesize Поточні процедури контролю доступу переглядаються та оновлюються після <AC-01\_\allowbreak{}ODP[08] подій>}




\subsection{УПРАВЛІННЯ ОБЛІКОВИМИ ЗАПИСАМИ (AC-2)}
a. Визначити та задокументувати типи облікових записів системи, дозволених для використання в ІС для підтримки цілей, завдань, функцій і процесів організації.\\ 
b. Призначити менеджерів облікових записів для управління системними обліковими записами.\\ 
c. Створити умови для групового та рольового членства.\\ 
d. Визначити авторизованих користувачів інформаційної системи, членство в групі та ролі, а також дозволи доступу (наприклад, привілеї) та інші атрибути (за потреби) для кожного облікового запису.\\ 
e. Вимагати схвалення [Призначення: визначеною організацією відповідальною особою або роллю] запитів на створення облікових записів системи.\\ 
f. Створювати, активувати, змінювати, деактивувати та видаляти системні облікові записи відповідно до [Призначення: визначених організацією політики, процедур та умов].\\ 
g. Впровадити моніторинг використання облікових записів системи.\\ 
h. Повідомляти адміністраторів облікових записів у межах [Призначення: визначеного організацією часового періоду для кожної ситуації]:\\ 
1. коли облікові записи більше не потрібні;\\ 
2. коли користувачі звільнені чи переведені;\\ 
3. коли використовуються індивідуальні системи або наявні зміни, які потребують нових знань.\\ 
i. Авторизувати доступ до системи на основі:\\ 
1. Дійсної авторизації доступу.\\ 
2. Передбачуваного використання системи.\\ 
3. Інших атрибутів, що вимагаються організацією.\\ 
j. Проводити перегляд облікових записів на відповідність вимогам управління обліковими записами з [Призначення: визначеною організацією частотою].\\ 
k. Впровадити процес повторного випуску облікових даних спільного/групового облікового запису (якщо він буде розгорнутий), коли особи виходять з групи.\\ 
l. Узгодити процеси управління обліковими записами з процесами звільнення та переводу (передачі повноважень) персоналу.

\vspace{10pt}
{\footnotesize\bfseries
No: 1\\
Name: ac\_\allowbreak{}2\_\allowbreak{}odp\_\allowbreak{}01\\
Type: string\\
Default: nil
}

{\footnotesize визначено передумови та критерії членства в групах і ролях}


{\footnotesize\bfseries
No: 2\\
Name: ac\_\allowbreak{}2\_\allowbreak{}odp\_\allowbreak{}02\\
Type: string\\
Default: nil
}

{\footnotesize визначено атрибути (за необхідності) для кожного облікового запису}


{\footnotesize\bfseries
No: 3\\
Name: ac\_\allowbreak{}2\_\allowbreak{}odp\_\allowbreak{}03\\
Type: list\\
Default: [\textquotedbl{}admin\textquotedbl{}, \textquotedbl{}security\_\allowbreak{}officer\textquotedbl{}]
}

{\footnotesize визначено персонал або ролі, необхідні для затвердження запитів на створення облікових записів}


{\footnotesize\bfseries
No: 4\\
Name: ac\_\allowbreak{}2\_\allowbreak{}odp\_\allowbreak{}04\\
Type: string\\
Default: nil
}

{\footnotesize визначено політику, процедури, передумови та критерії створення, активації, зміни, деактивації та видалення облікових записів}


{\footnotesize\bfseries
No: 5\\
Name: ac\_\allowbreak{}2\_\allowbreak{}odp\_\allowbreak{}05\\
Type: list\\
Default: [\textquotedbl{}admin\textquotedbl{}, \textquotedbl{}security\_\allowbreak{}officer\textquotedbl{}]
}

{\footnotesize визначено персонал або ролі, які мають бути повідомлені}


{\footnotesize\bfseries
No: 6\\
Name: ac\_\allowbreak{}2\_\allowbreak{}odp\_\allowbreak{}06\\
Type: integer\\
Default: 30
}

{\footnotesize визначено період часу, протягом якого адміністратори облікових записів повинні бути повідомлені про те, що облікові записи більше не потрібні}


{\footnotesize\bfseries
No: 7\\
Name: ac\_\allowbreak{}2\_\allowbreak{}odp\_\allowbreak{}07\\
Type: integer\\
Default: 30
}

{\footnotesize визначено термін, протягом якого необхідно повідомляти адміністраторів облікових записів про звільнення або переведення користувачів}


{\footnotesize\bfseries
No: 8\\
Name: ac\_\allowbreak{}2\_\allowbreak{}odp\_\allowbreak{}08\\
Type: integer\\
Default: 30
}

{\footnotesize визначено період часу, протягом якого необхідно повідомляти адміністраторів облікових записів про зміни у використанні системи або необхідність знати про зміни для окремої особи}


{\footnotesize\bfseries
No: 9\\
Name: ac\_\allowbreak{}2\_\allowbreak{}odp\_\allowbreak{}09\\
Type: string\\
Default: nil
}

{\footnotesize визначено атрибути, необхідні для авторизації доступу до системи (за потреби)}


{\footnotesize\bfseries
No: 10\\
Name: ac\_\allowbreak{}2\_\allowbreak{}odp\_\allowbreak{}10\\
Type: integer\\
Default: 30
}

{\footnotesize визначено періодичність перегляду облікових записів}


{\footnotesize\bfseries
No: 11\\
Name: ac\_\allowbreak{}2\_\allowbreak{}a\_\allowbreak{}01\\
Type: string\\
Default: nil
}

{\footnotesize визначено та задокументовано типи облікових записів, дозволених для використання в системі}


{\footnotesize\bfseries
No: 12\\
Name: ac\_\allowbreak{}2\_\allowbreak{}a\_\allowbreak{}02\\
Type: string\\
Default: nil
}

{\footnotesize визначено та задокументовано типи облікових записів, які заборонено використовувати в системі}


{\footnotesize\bfseries
No: 13\\
Name: ac\_\allowbreak{}2\_\allowbreak{}b\\
Type: string\\
Default: nil
}

{\footnotesize призначені менеджери облікових записів}


{\footnotesize\bfseries
No: 14\\
Name: ac\_\allowbreak{}2\_\allowbreak{}c\\
Type: string\\
Default: nil
}

{\footnotesize необхідні <AC-02\_\allowbreak{}ODP[01] умови та критерії> для членства в групах та ролях}


{\footnotesize\bfseries
No: 15\\
Name: ac\_\allowbreak{}2\_\allowbreak{}d\_\allowbreak{}01\\
Type: string\\
Default: nil
}

{\footnotesize визначено авторизованих користувачів системи}


{\footnotesize\bfseries
No: 16\\
Name: ac\_\allowbreak{}2\_\allowbreak{}d\_\allowbreak{}02\\
Type: string\\
Default: nil
}

{\footnotesize вказано приналежність до групи або ролі}


{\footnotesize\bfseries
No: 17\\
Name: ac\_\allowbreak{}2\_\allowbreak{}d\_\allowbreak{}03\_\allowbreak{}01\\
Type: string\\
Default: nil
}

{\footnotesize для кожного облікового запису вказуються повноваження доступу (тобто привілеї)}


{\footnotesize\bfseries
No: 18\\
Name: ac\_\allowbreak{}2\_\allowbreak{}d\_\allowbreak{}03\_\allowbreak{}02\\
Type: string\\
Default: nil
}

{\footnotesize <AC-02\_\allowbreak{}ODP[02] атрибути (за необхідності)> вказуються для кожного облікового запису}


{\footnotesize\bfseries
No: 19\\
Name: ac\_\allowbreak{}2\_\allowbreak{}e\\
Type: string\\
Default: nil
}

{\footnotesize для запитів на створення облікових записів потрібні схвалення від <AC-02\_\allowbreak{}ODP[03] персоналу або ролей>}


{\footnotesize\bfseries
No: 20\\
Name: ac\_\allowbreak{}2\_\allowbreak{}f\_\allowbreak{}01\\
Type: string\\
Default: nil
}

{\footnotesize облікові записи створюються відповідно до <AC-02\_\allowbreak{}ODP[04] політики, процедур, передумов та критеріїв>}


{\footnotesize\bfseries
No: 21\\
Name: ac\_\allowbreak{}2\_\allowbreak{}f\_\allowbreak{}02\\
Type: string\\
Default: nil
}

{\footnotesize облікові записи активуються відповідно до <AC-02\_\allowbreak{}ODP[04] політики, процедур, передумов та критеріїв>}


{\footnotesize\bfseries
No: 22\\
Name: ac\_\allowbreak{}2\_\allowbreak{}f\_\allowbreak{}03\\
Type: string\\
Default: nil
}

{\footnotesize облікові записи змінюються відповідно до <AC-02\_\allowbreak{}ODP[04] політики, процедур, передумов та критеріїв>}


{\footnotesize\bfseries
No: 23\\
Name: ac\_\allowbreak{}2\_\allowbreak{}f\_\allowbreak{}04\\
Type: string\\
Default: nil
}

{\footnotesize облікові записи деактивуються відповідно до <AC-02\_\allowbreak{}ODP[04] політики, процедур, передумов та критеріїв>}


{\footnotesize\bfseries
No: 24\\
Name: ac\_\allowbreak{}2\_\allowbreak{}f\_\allowbreak{}05\\
Type: string\\
Default: nil
}

{\footnotesize облікові записи видаляються відповідно до <AC-02\_\allowbreak{}ODP[04] політики, процедур, передумов та критеріїв>}


{\footnotesize\bfseries
No: 25\\
Name: ac\_\allowbreak{}2\_\allowbreak{}g\\
Type: string\\
Default: nil
}

{\footnotesize контролюється використання облікових записів}


{\footnotesize\bfseries
No: 26\\
Name: ac\_\allowbreak{}2\_\allowbreak{}h\_\allowbreak{}01\\
Type: string\\
Default: nil
}

{\footnotesize адміністратори облікових записів та <AC-02\_\allowbreak{}ODP[05] персонал або ролі> отримують повідомлення протягом <AC-02\_\allowbreak{}ODP[06] періоду часу>, коли облікові записи більше не потрібні}


{\footnotesize\bfseries
No: 27\\
Name: ac\_\allowbreak{}2\_\allowbreak{}h\_\allowbreak{}02\\
Type: string\\
Default: nil
}

{\footnotesize адміністратори облікових записів та <AC-02\_\allowbreak{}ODP[05] персонал або ролі> отримують повідомлення протягом <AC-02\_\allowbreak{}ODP[07] періоду часу>, коли користувачі звільнені чи переведені}


{\footnotesize\bfseries
No: 28\\
Name: ac\_\allowbreak{}2\_\allowbreak{}h\_\allowbreak{}03\\
Type: string\\
Default: nil
}

{\footnotesize адміністратори облікових записів та <AC-02\_\allowbreak{}ODP[05] персонал або ролі> отримують повідомлення протягом <AC-02\_\allowbreak{}ODP[08] періоду часу>, коли використовуються індивідуальні системи або наявні зміни, які потребують нових знань}


{\footnotesize\bfseries
No: 29\\
Name: ac\_\allowbreak{}2\_\allowbreak{}i\_\allowbreak{}01\\
Type: string\\
Default: nil
}

{\footnotesize доступ до системи здійснюється на підставі дійсної авторизації доступу}


{\footnotesize\bfseries
No: 30\\
Name: ac\_\allowbreak{}2\_\allowbreak{}i\_\allowbreak{}02\\
Type: string\\
Default: nil
}

{\footnotesize доступ до системи авторизується на основі передбачуваного використання системи}


{\footnotesize\bfseries
No: 31\\
Name: ac\_\allowbreak{}2\_\allowbreak{}i\_\allowbreak{}03\\
Type: string\\
Default: nil
}

{\footnotesize доступ до системи авторизовано на основі атрибутів <AC-02\_\allowbreak{}ODP[09] (за необхідності)>}


{\footnotesize\bfseries
No: 32\\
Name: ac\_\allowbreak{}2\_\allowbreak{}j\\
Type: string\\
Default: nil
}

{\footnotesize облікові записи переглядаються на відповідність вимогам управління обліковими записами <AC-02\_\allowbreak{}ODP[10] частота>}


{\footnotesize\bfseries
No: 33\\
Name: ac\_\allowbreak{}2\_\allowbreak{}k\_\allowbreak{}01\\
Type: string\\
Default: nil
}

{\footnotesize створено процес повторного випуску облікових даних спільного доступу або групових облікових записів (якщо вони розгорнуті), коли користувачів вилучено з групи}


{\footnotesize\bfseries
No: 34\\
Name: ac\_\allowbreak{}2\_\allowbreak{}k\_\allowbreak{}02\\
Type: string\\
Default: nil
}

{\footnotesize впроваджено процес повторного випуску облікових даних спільного доступу або групових облікових записів (якщо вони розгорнуті), коли користувачів вилучено з групи}


{\footnotesize\bfseries
No: 35\\
Name: ac\_\allowbreak{}2\_\allowbreak{}l\_\allowbreak{}01\\
Type: string\\
Default: nil
}

{\footnotesize процеси управління обліковими записами узгоджуються з процесами звільнення персоналу}


{\footnotesize\bfseries
No: 36\\
Name: ac\_\allowbreak{}2\_\allowbreak{}l\_\allowbreak{}02\\
Type: string\\
Default: nil
}

{\footnotesize процеси управління обліковими записами узгоджуються з процесами переводу персоналу}




\subsubsection{АВТОМАТИЗОВАНЕ УПРАВЛІННЯ ОБЛІКОВИМИ ЗАПИСАМИ СИСТЕМИ (AC-2(1))}
Використовувати автоматизовані системними обліковими записами. механізми для підтримки управління

\vspace{10pt}
{\footnotesize\bfseries
No: 1\\
Name: ac\_\allowbreak{}2\_\allowbreak{}1\_\allowbreak{}01\\
Type: string\\
Default: \textquotedbl{}автоматизований засіб моніторингу\textquotedbl{}
}

{\footnotesize Управління обліковими записами системи підтримується за допомогою автоматизовані механізми}


{\footnotesize\bfseries
No: 2\\
Name: ac\_\allowbreak{}2\_\allowbreak{}1\_\allowbreak{}odp\\
Type: string\\
Default: \textquotedbl{}автоматизований засіб моніторингу\textquotedbl{}
}

{\footnotesize Визначено автоматизовані механізми, що використовуються для підтримки управління обліковими записами системи}




\subsubsection{ВИДАЛЕННЯ ТИМЧАСОВИХ ТА ЕКСТРЕНИХ ОБЛІКОВИХ ЗАПИСІВ (AC-2(2))}
Автоматично видаляти або деактивувати тимчасові та екстрені облікові записи через [Призначення: визначений організацією період часу].

\vspace{10pt}
{\footnotesize\bfseries
No: 1\\
Name: ac\_\allowbreak{}2\_\allowbreak{}2\_\allowbreak{}odp\_\allowbreak{}01\\
Type: string\\
Default: \textquotedbl{}деактивувати\textquotedbl{}
}

{\footnotesize вибрано одне з наступних ЗНАЧЕНЬ ПАРАМЕТРІВ: \{видаляти; деактивувати\}}


{\footnotesize\bfseries
No: 2\\
Name: ac\_\allowbreak{}2\_\allowbreak{}2\_\allowbreak{}odp\_\allowbreak{}02\\
Type: integer\\
Default: 30
}

{\footnotesize визначено період часу, після якого автоматично видаляються або деактивуються тимчасові або екстрені облікові записи}


{\footnotesize\bfseries
No: 3\\
Name: ac\_\allowbreak{}2\_\allowbreak{}2\_\allowbreak{}01\\
Type: string\\
Default: nil
}

{\footnotesize тимчасові та екстрені облікові записи автоматично <AC-02(02) \_\allowbreak{}ODP[01] ВИБРАНЕ ЗНАЧЕННЯ ПАРАМЕТРА> через <AC-02(02) \_\allowbreak{}ODP[02] період часу>}




\subsubsection{ДЕАКТИВАЦІЯ ОБЛІКОВИХ ЗАПИСІВ (AC-2(3))}
Автоматично деактивувати облікові записи коли: a) їх строк дії минув; b) вони більше не пов'язані з користувачем; c) вони порушують організаційну політику; d) вони були неактивними впродовж [Призначення: визначеного організацією періоду часу].

\vspace{10pt}
{\footnotesize\bfseries
No: 1\\
Name: ac\_\allowbreak{}2\_\allowbreak{}3\_\allowbreak{}odp\_\allowbreak{}01\\
Type: integer\\
Default: 30
}

{\footnotesize Визначено період часу, протягом якого необхідно деактивувати облікові записи}


{\footnotesize\bfseries
No: 2\\
Name: ac\_\allowbreak{}2\_\allowbreak{}3\_\allowbreak{}odp\_\allowbreak{}02\\
Type: integer\\
Default: 30
}

{\footnotesize Визначено період часу неактивності, після закінчення якого облікові записи будуть деактивовані}


{\footnotesize\bfseries
No: 3\\
Name: ac\_\allowbreak{}2\_\allowbreak{}3\_\allowbreak{}a\\
Type: string\\
Default: nil
}

{\footnotesize облікові записи деактивуються протягом часового періоду, коли термін дії облікових записів минув}


{\footnotesize\bfseries
No: 4\\
Name: ac\_\allowbreak{}2\_\allowbreak{}3\_\allowbreak{}b\\
Type: string\\
Default: nil
}

{\footnotesize Облікові записи деактивуються протягом часового періоду, коли облікові записи більше не пов'язані з користувачем або фізичною особою}


{\footnotesize\bfseries
No: 5\\
Name: ac\_\allowbreak{}2\_\allowbreak{}3\_\allowbreak{}c\\
Type: string\\
Default: nil
}

{\footnotesize Облікові записи відключено протягом часового періоду, коли облікові записи порушують політику організації}


{\footnotesize\bfseries
No: 6\\
Name: ac\_\allowbreak{}2\_\allowbreak{}3\_\allowbreak{}d\\
Type: string\\
Default: nil
}

{\footnotesize облікові записи деактивуються протягом часового періоду, якщо вони були неактивними протягом визначеного періоду часу}




\subsubsection{ДІЇ ПРИ АВТОМАТИЗОВАНОМУ АУДИТІ (AC-2(4))}
Проводити автоматизований аудит створення, модифікації, деактивації та видалення облікових записів і сповіщення про дії. ПРИ активації,

\vspace{10pt}
{\footnotesize\bfseries
No: 1\\
Name: ac\_\allowbreak{}2\_\allowbreak{}4\_\allowbreak{}01\\
Type: string\\
Default: nil
}

{\footnotesize Створення облікового запису автоматично аудитується}


{\footnotesize\bfseries
No: 2\\
Name: ac\_\allowbreak{}2\_\allowbreak{}4\_\allowbreak{}02\\
Type: string\\
Default: nil
}

{\footnotesize Модифікація облікового запису автоматично аудитується}


{\footnotesize\bfseries
No: 3\\
Name: ac\_\allowbreak{}2\_\allowbreak{}4\_\allowbreak{}03\\
Type: string\\
Default: nil
}

{\footnotesize Активація облікового запису автоматично аудитується}


{\footnotesize\bfseries
No: 4\\
Name: ac\_\allowbreak{}2\_\allowbreak{}4\_\allowbreak{}04\\
Type: string\\
Default: nil
}

{\footnotesize Деактивація облікового запису автоматично аудитується}


{\footnotesize\bfseries
No: 5\\
Name: ac\_\allowbreak{}2\_\allowbreak{}4\_\allowbreak{}05\\
Type: string\\
Default: nil
}

{\footnotesize Видалення облікового запису автоматично аудитується}




\subsubsection{ВИХІД ІЗ СИСТЕМИ ЗА ВІДСУТНОСТІ АКТИВНОСТІ (AC-2(5))}
Вимагати від користувачів виходити із системи, коли [Призначення: вичерпано визначений організацією періоду часу очікування або опис того, коли необхідно вийти із системи].

\vspace{10pt}
{\footnotesize\bfseries
No: 1\\
Name: ac\_\allowbreak{}2\_\allowbreak{}5\_\allowbreak{}01\\
Type: integer\\
Default: 30
}

{\footnotesize Користувачі повинні виходити з системи, коли період очікуваної бездіяльності або опис часу, коли потрібно вийти з системи}


{\footnotesize\bfseries
No: 2\\
Name: ac\_\allowbreak{}2\_\allowbreak{}5\_\allowbreak{}odp\\
Type: integer\\
Default: 30
}

{\footnotesize Визначено часовий період очікуваної бездіяльності або опис, коли потрібно вийти з системи}




\subsubsection{ДИНАМІЧНЕ УПРАВЛІННЯ ПРИВІЛЕЯМИ (AC-2(6))}
Реалізувати такі можливості динамічного управління привілеями: [Призначення: визначений організацією перелік можливостей динамічного управління привілеями].

\vspace{10pt}
{\footnotesize\bfseries
No: 1\\
Name: ac\_\allowbreak{}2\_\allowbreak{}6\_\allowbreak{}01\\
Type: string\\
Default: nil
}

{\footnotesize Реалізовано можливості динамічного управління привілеями}


{\footnotesize\bfseries
No: 2\\
Name: ac\_\allowbreak{}2\_\allowbreak{}6\_\allowbreak{}odp\\
Type: string\\
Default: nil
}

{\footnotesize Визначено можливості динамічного управління привілеями}




\subsubsection{СХЕМИ, ЗАСНОВАНІ НА РОЛЯХ (AC-2(7))}
a) Створювати й адмініструвати привілейовані облікові записи користувачів відповідно до схеми доступу на основі ролей (role-based), яка реалізує дозволений доступ до системи та призначення привілеїв для ролей. b) Проводити моніторинг призначення привілейованих ролей. c) Відстежувати зміни ролей або атрибутів. d) Скасовувати доступ, коли призначені привілейовані ролі більше не потрібні.

\vspace{10pt}
{\footnotesize\bfseries
No: 1\\
Name: ac\_\allowbreak{}2\_\allowbreak{}7\_\allowbreak{}b\\
Type: string\\
Default: nil
}

{\footnotesize Проводиться моніторинг призначення привілейованих ролей або атрибутів}


{\footnotesize\bfseries
No: 2\\
Name: ac\_\allowbreak{}2\_\allowbreak{}7\_\allowbreak{}c\\
Type: string\\
Default: nil
}

{\footnotesize Відстежуються зміни ролей або атрибутів}


{\footnotesize\bfseries
No: 3\\
Name: ac\_\allowbreak{}2\_\allowbreak{}7\_\allowbreak{}d\\
Type: list\\
Default: [\textquotedbl{}admin\textquotedbl{}, \textquotedbl{}security\_\allowbreak{}officer\textquotedbl{}]
}

{\footnotesize Доступ скасовується, коли призначені привілейовані ролі більше не потрібні.0}


{\footnotesize\bfseries
No: 4\\
Name: ac\_\allowbreak{}2\_\allowbreak{}7\_\allowbreak{}odp\\
Type: string\\
Default: nil
}

{\footnotesize Вибрано одне з наступних ЗНАЧЕНЬ ПАРАМЕТРІВ: \{схема доступу на основі ролей; схема доступу на основі атрибутів\}}




\subsubsection{ДИНАМІЧНЕ УПРАВЛІННЯ ОБЛІКОВИМИ ЗАПИСАМИ (AC-2(8))}
Створювати, активувати, управляти та деактивувати [Призначення: системні облікові записи, визначені організацією] динамічно.

\vspace{10pt}
{\footnotesize\bfseries
No: 1\\
Name: ac\_\allowbreak{}2\_\allowbreak{}8\_\allowbreak{}01\\
Type: string\\
Default: nil
}

{\footnotesize УПРАВЛІННЯ ОБЛІКОВИМИ ЗАПИСАМИ - ДИНАМІЧНЕ УПРАВЛІННЯ ОБЛІКОВИМИ ЗАПИСАМИ МЕТА ОЦІНКИ: Визначити, чи:}


{\footnotesize\bfseries
No: 2\\
Name: ac\_\allowbreak{}2\_\allowbreak{}8\_\allowbreak{}02\\
Type: string\\
Default: nil
}

{\footnotesize Облікові записи системи активуються динамічно}


{\footnotesize\bfseries
No: 3\\
Name: ac\_\allowbreak{}2\_\allowbreak{}8\_\allowbreak{}03\\
Type: string\\
Default: nil
}

{\footnotesize Облікові записи системи активуються динамічно}


{\footnotesize\bfseries
No: 4\\
Name: ac\_\allowbreak{}2\_\allowbreak{}8\_\allowbreak{}04\\
Type: string\\
Default: nil
}

{\footnotesize Облікові записи системи деактивуються динамічно}


{\footnotesize\bfseries
No: 5\\
Name: ac\_\allowbreak{}2\_\allowbreak{}8\_\allowbreak{}odp\\
Type: string\\
Default: nil
}

{\footnotesize Визначено облікові записи системи, які динамічно створюються, активуються, управляються та деактивуються}




\subsubsection{ОБМЕЖЕННЯ НА ВИКОРИСТАННЯ СПІЛЬНИХ ТА ГРУПОВИХ ОБЛІКОВИХ ЗАПИСІВ (AC-2(9))}
Використовувати лише ті спільні та групові облікові записи, які відповідають [Призначення: визначеним організацією умовам для створення спільних та групових облікових записів].

\vspace{10pt}
{\footnotesize\bfseries
No: 1\\
Name: ac\_\allowbreak{}2\_\allowbreak{}9\_\allowbreak{}01\\
Type: list\\
Default: []
}

{\footnotesize Використання спільних та групових облікових записів дозволено лише за умови дотримання умов}


{\footnotesize\bfseries
No: 2\\
Name: ac\_\allowbreak{}2\_\allowbreak{}9\_\allowbreak{}odp\\
Type: list\\
Default: []
}

{\footnotesize Визначено умови створення спільних та групових облікових записів}




\subsubsection{ЗМІНА ДАНИХ СПІЛЬНИХ І ГРУПОВИХ ОБЛІКОВИХ ЗАПИСІВ (AC-2(10)) [Вилучено]}
[Вилучено: включено до AC-02(k)]

\vspace{10pt}
\textit{Немає параметрів для цього контролю.}
\vspace{10pt}


\subsubsection{УМОВИ ВИКОРИСТАННЯ (AC-2(11))}
Забезпечити дотримання [Призначення: обставин та/або умов використання, визначених організацією] для [Призначення: визначених організацією облікових записів системи].

\vspace{10pt}
{\footnotesize\bfseries
No: 1\\
Name: ac\_\allowbreak{}2\_\allowbreak{}11\_\allowbreak{}01\\
Type: list\\
Default: []
}

{\footnotesize Обставини та/або умови використання для облікових записів системи застосовуються}


{\footnotesize\bfseries
No: 2\\
Name: ac\_\allowbreak{}2\_\allowbreak{}11\_\allowbreak{}odp\_\allowbreak{}01\\
Type: list\\
Default: []
}

{\footnotesize Визначено обставини та/або умови використання визначених облікових записів системи}


{\footnotesize\bfseries
No: 3\\
Name: ac\_\allowbreak{}2\_\allowbreak{}11\_\allowbreak{}odp\_\allowbreak{}02\\
Type: string\\
Default: nil
}

{\footnotesize Визначені облікові записи системи, що підлягають виконанню обставин та/або умов використання}




\subsubsection{МОНІТОРИНГ НЕТИПОВОГО ВИКОРИСТАННЯ ОБЛІКОВИХ ЗАПИСІВ (AC-2(12))}
a) Проводити моніторинг облікових записів системи на [Призначення: визначене організацією нетипове використання]. b) Повідомляти про нетипове використання облікових записів системи [Призначення: визначеного організацією персоналу або ролей].

\vspace{10pt}
{\footnotesize\bfseries
No: 1\\
Name: ac\_\allowbreak{}2\_\allowbreak{}12\_\allowbreak{}a\\
Type: list\\
Default: []
}

{\footnotesize Облікові записи системи відстежуються на предмет обставини та/або умови використання}


{\footnotesize\bfseries
No: 2\\
Name: ac\_\allowbreak{}2\_\allowbreak{}12\_\allowbreak{}b\\
Type: list\\
Default: [\textquotedbl{}admin\textquotedbl{}, \textquotedbl{}security\_\allowbreak{}officer\textquotedbl{}]
}

{\footnotesize Про визначені обставини та/або умови використання системних облікових записів повідомляється персонал або ролі}


{\footnotesize\bfseries
No: 3\\
Name: ac\_\allowbreak{}2\_\allowbreak{}12\_\allowbreak{}odp\_\allowbreak{}01\\
Type: list\\
Default: []
}

{\footnotesize Визначено обставини та/або умови використання, для яких необхідно здійснювати моніторинг обліко- вих записів системи}


{\footnotesize\bfseries
No: 4\\
Name: ac\_\allowbreak{}2\_\allowbreak{}12\_\allowbreak{}odp\_\allowbreak{}02\\
Type: list\\
Default: [\textquotedbl{}admin\textquotedbl{}, \textquotedbl{}security\_\allowbreak{}officer\textquotedbl{}]
}

{\footnotesize Визначено персонал або ролі, яким належить повідомляти про визначені обставини та/або умови використання}




\subsubsection{ДЕАКТИВАЦІЯ ОБЛІКОВИХ ЗАПИСІВ ОСІБ З ВИСОКИМ РІВНЕМ РИЗИКУ (AC-2(13))}
Деактивувати облікові записи користувачів, які становлять значний ризик, у межах [Призначення: визначеного організацією періоду часу] після виявлення ризику.

\vspace{10pt}
{\footnotesize\bfseries
No: 1\\
Name: ac\_\allowbreak{}2\_\allowbreak{}13\_\allowbreak{}01\\
Type: integer\\
Default: 30
}

{\footnotesize Облікові записи користувачів деактивуються протягом періоду часу з моменту виявлення значних ризиків}


{\footnotesize\bfseries
No: 2\\
Name: ac\_\allowbreak{}2\_\allowbreak{}13\_\allowbreak{}odp\_\allowbreak{}01\\
Type: integer\\
Default: 30
}

{\footnotesize Визначено період часу, протягом якого необхідно деактивувати облікові записи фізичних осіб, які становлять значний ризик}


{\footnotesize\bfseries
No: 3\\
Name: ac\_\allowbreak{}2\_\allowbreak{}13\_\allowbreak{}odp\_\allowbreak{}02\\
Type: string\\
Default: nil
}

{\footnotesize Визначено значні ризики, що призводять до деактивації облікових записів}




\subsection{ЗАБЕЗПЕЧЕННЯ ДОСТУПУ (AC-3)}
Застосовувати затверджені повноваження для логічного доступу до інформації та ресурсів системи відповідно до чинної політики (правил) управління доступом.

\vspace{10pt}
{\footnotesize\bfseries
No: 1\\
Name: ac\_\allowbreak{}3\_\allowbreak{}01\\
Type: list\\
Default: [\textquotedbl{}default\_\allowbreak{}deny\_\allowbreak{}rule\textquotedbl{}, \textquotedbl{}abac\_\allowbreak{}rule\_\allowbreak{}1\textquotedbl{}]
}

{\footnotesize Затверджені повноваження на логічний доступ до інформації та ресурсів системи виконуються відповідно до чинних політик(правил) управління доступом}




\subsubsection{ОБМЕЖЕНИЙ ДОСТУП ДО ПРИВІЛЕЙОВАНИХ ФУНКЦІЙ (AC-3(1)) [Вилучено]}
[Вилучено: включено до складу AC-06]

\vspace{10pt}
\textit{Немає параметрів для цього контролю.}
\vspace{10pt}


\subsubsection{ПОДВІЙНА АВТОРИЗАЦІЯ (AC-3(2))}
Забезпечити подвійну авторизацію для [Призначення: визначених організацією привілейованих команд та/або інших дій, визначених організацією].

\vspace{10pt}
{\footnotesize\bfseries
No: 1\\
Name: ac\_\allowbreak{}3\_\allowbreak{}2\_\allowbreak{}01\\
Type: string\\
Default: nil
}

{\footnotesize Подвійна авторизація застосовується для привілейованих команд та/або інших дій}


{\footnotesize\bfseries
No: 2\\
Name: ac\_\allowbreak{}3\_\allowbreak{}2\_\allowbreak{}odp\\
Type: list\\
Default: [\textquotedbl{}login\textquotedbl{}, \textquotedbl{}logout\textquotedbl{}, \textquotedbl{}failed\_\allowbreak{}attempt\textquotedbl{}]
}

{\footnotesize Визначено привілейовані команди та/або інші дії, що потребують подвійної авторизації}




\subsubsection{МАНДАТНЕ УПРАВЛІННЯ ДОСТУПОМ (AC-3(3))}
Застосовувати [Призначення: визначену організацією мандатну (mandatory) політику управління доступом] щодо всіх суб'єктів і об'єктів доступу, у яких політика:\\ 
(a) одноманітно застосовується для всіх суб'єктів і об'єктів у межах системи;\\ 
(b) вказує, що суб'єкт, якому було надано доступ до інформації, обмежений у виконанні будь-якої з таких дій:\\ 
(c) (1) передача інформації неавторизованим суб'єктам або об'єктам; (2) надання іншим суб'єктам привілеїв; (3) зміна одного чи декількох атрибутів безпеки суб'єкта, об'єкта, системи або компонентів системи; (4) вибір атрибутів безпеки та значень атрибутів, які повинні бути пов'язані з новоствореними або зміненими об'єктами; (5) зміна правил, що регулюють управління доступом; має бути вказано, що [Призначення: визначеним організацією суб'єктам] можуть бути явно надані [Призначення: визначені організацією привілеї], так що вони не обмежуються будь-яким з перелічених вище обмежень.

\vspace{10pt}
{\footnotesize\bfseries
No: 1\\
Name: ac\_\allowbreak{}3\_\allowbreak{}3\_\allowbreak{}odp\_\allowbreak{}01\\
Type: string\\
Default: nil
}

{\footnotesize Визначено мандатну політику контролю доступу, що застосовується до набору охоплених суб'єктів}


{\footnotesize\bfseries
No: 2\\
Name: ac\_\allowbreak{}3\_\allowbreak{}3\_\allowbreak{}odp\_\allowbreak{}02\\
Type: string\\
Default: nil
}

{\footnotesize Визначено мандатну політику контролю доступу, що застосовується до набору охоплених об'єктів}


{\footnotesize\bfseries
No: 3\\
Name: ac\_\allowbreak{}3\_\allowbreak{}3\_\allowbreak{}odp\_\allowbreak{}03\\
Type: string\\
Default: nil
}

{\footnotesize Визначені суб'єкти, яким явно надаються привілеї}


{\footnotesize\bfseries
No: 4\\
Name: ac\_\allowbreak{}3\_\allowbreak{}3\_\allowbreak{}odp\_\allowbreak{}04\\
Type: string\\
Default: nil
}

{\footnotesize Визначено привілеї, які мають бути прямо надані суб'єктам}


{\footnotesize\bfseries
No: 5\\
Name: ac\_\allowbreak{}3\_\allowbreak{}3\_\allowbreak{}01\\
Type: string\\
Default: nil
}

{\footnotesize <AC-03(03)\_\allowbreak{}ODP[01] мандатна політика контролю доступу> застосовується до набору охоплених суб'єктів, зазначених у політиці}


{\footnotesize\bfseries
No: 6\\
Name: ac\_\allowbreak{}3\_\allowbreak{}3\_\allowbreak{}02\\
Type: string\\
Default: nil
}

{\footnotesize <AC-03(03)\_\allowbreak{}ODP[02] мандатна політика контролю доступу> застосовується до набору охоплених об'єктів, зазначених у політиці}


{\footnotesize\bfseries
No: 7\\
Name: ac\_\allowbreak{}3\_\allowbreak{}3\_\allowbreak{}a\_\allowbreak{}01\\
Type: string\\
Default: nil
}

{\footnotesize <AC-03(03)\_\allowbreak{}ODP[01] мандатна політика контролю доступу> застосовується одноманітно до всіх суб'єктів системи}


{\footnotesize\bfseries
No: 8\\
Name: ac\_\allowbreak{}3\_\allowbreak{}3\_\allowbreak{}a\_\allowbreak{}02\\
Type: string\\
Default: nil
}

{\footnotesize <AC-03(03)\_\allowbreak{}ODP[02] мандатна політика контролю доступу> застосовується одноманітно до всіх об'єктів системи}


{\footnotesize\bfseries
No: 9\\
Name: ac\_\allowbreak{}3\_\allowbreak{}3\_\allowbreak{}b\_\allowbreak{}01\\
Type: string\\
Default: nil
}

{\footnotesize <AC-03(03)\_\allowbreak{}ODP[01] мандатна політика контролю доступу> та <AC-03(03)\_\allowbreak{}ODP[02] мандатна політика контролю доступу> визначають, що суб'єкт, якому надано доступ до інформації, зобов'язаний не передавати інформацію неавторизованим суб'єктам або об'єктам}


{\footnotesize\bfseries
No: 10\\
Name: ac\_\allowbreak{}3\_\allowbreak{}3\_\allowbreak{}b\_\allowbreak{}02\\
Type: string\\
Default: nil
}

{\footnotesize <AC-03(03)\_\allowbreak{}ODP[01] мандатна політика контролю доступу> та <AC-03(03)\_\allowbreak{}ODP[02] мандатна політика контролю доступу> визначають, що суб'єкт, якому надано доступ до інформації, обмежений у наданні своїх привілеїв іншим суб'єктам}


{\footnotesize\bfseries
No: 11\\
Name: ac\_\allowbreak{}3\_\allowbreak{}3\_\allowbreak{}b\_\allowbreak{}03\\
Type: string\\
Default: nil
}

{\footnotesize <AC-03(03)\_\allowbreak{}ODP[01] мандатна політика контролю доступу> та <AC-03(03)\_\allowbreak{}ODP[02] мандатна політика контролю доступу> визначають, що суб'єкт, якому надано доступ до інформації, не може змінювати один або декілька атрибутів безпеки (визначених політикою) суб'єктів, об'єктів, системи або компонентів системи}


{\footnotesize\bfseries
No: 12\\
Name: ac\_\allowbreak{}3\_\allowbreak{}3\_\allowbreak{}b\_\allowbreak{}04\\
Type: string\\
Default: nil
}

{\footnotesize <AC-03(03)\_\allowbreak{}ODP[01] мандатна політика контролю доступу> та <AC-03(03)\_\allowbreak{}ODP[02] мандатна політика контролю доступу> визначають, що суб'єкт, якому надано доступ до інформації, обмежений у виборі атрибутів безпеки та значень атрибутів (визначених політикою), що повинні бути пов'язані з новостворюваними або зміненими об'єктами}


{\footnotesize\bfseries
No: 13\\
Name: ac\_\allowbreak{}3\_\allowbreak{}3\_\allowbreak{}b\_\allowbreak{}05\\
Type: string\\
Default: nil
}

{\footnotesize <AC-03(03)\_\allowbreak{}ODP[01] мандатна політика контролю доступу> та <AC-03(03)\_\allowbreak{}ODP[02] мандатна політика контролю доступу> визначають, що суб'єкт, якому надано доступ до інформації, не має права змінювати правила, що регулюють управління доступом}


{\footnotesize\bfseries
No: 14\\
Name: ac\_\allowbreak{}3\_\allowbreak{}3\_\allowbreak{}c\\
Type: string\\
Default: nil
}

{\footnotesize <AC-03(03)\_\allowbreak{}ODP[01] мандатна політика контролю доступу> та <AC-03(03)\_\allowbreak{}ODP[02] мандатна політика контролю доступу> визначають, що <AC-03(03)\_\allowbreak{}ODP[03] суб'єктам> можуть бути явно надані <AC-03(03)\_\allowbreak{}ODP[04] привілеї> таким чином, щоб вони не були обмежені будь-якою визначеною підмножиною (або всіма) з наведених вище обмежень}




\subsubsection{ДИСКРЕЦІЙНЕ УПРАВЛІННЯ ДОСТУПОМ (AC-3(4))}
Застосовувати [Призначення: визначену організацією дискреційну політику управління доступом] щодо визначених суб'єктів і об'єктів доступу, для яких політика визначає, що суб'єкт, якому було надано доступ до інформації, може виконати одну чи більше з таких дій:\\ 
(a) передача інформацію будь-яким іншим суб'єктам чи об'єктам;\\ 
(b) призначення своїх привілей іншим суб'єктам;\\ 
(c) зміна атрибутів безпеки суб'єктів, об'єктів, систем або компонентів системи;\\ 
(d) вибір атрибутів безпеки, які будуть пов'язані з новоствореними або переглянутими об'єктами;\\ 
(e) зміна правил, що регулюють управління доступом.

\vspace{10pt}
{\footnotesize\bfseries
No: 1\\
Name: ac\_\allowbreak{}3\_\allowbreak{}4\_\allowbreak{}odp\_\allowbreak{}01\\
Type: string\\
Default: nil
}

{\footnotesize Визначено дискреційну політику управління доступом, яка застосовується до набору охоплених суб'єктів}


{\footnotesize\bfseries
No: 2\\
Name: ac\_\allowbreak{}3\_\allowbreak{}4\_\allowbreak{}odp\_\allowbreak{}02\\
Type: string\\
Default: nil
}

{\footnotesize Визначено дискреційну політику управління доступом, яка застосовується до набору охоплених об'єктів}


{\footnotesize\bfseries
No: 3\\
Name: ac\_\allowbreak{}3\_\allowbreak{}4\_\allowbreak{}01\\
Type: string\\
Default: nil
}

{\footnotesize <AC-03(04)\_\allowbreak{}ODP[01] дискреційна політика управління доступом> застосовується до набору охоплених суб'єктів, зазначених у політиці}


{\footnotesize\bfseries
No: 4\\
Name: ac\_\allowbreak{}3\_\allowbreak{}4\_\allowbreak{}02\\
Type: string\\
Default: nil
}

{\footnotesize <AC-03(04)\_\allowbreak{}ODP[02] дискреційна політика управління доступом> застосовується до набору охоплених об'єктів, зазначених у політиці}


{\footnotesize\bfseries
No: 5\\
Name: ac\_\allowbreak{}3\_\allowbreak{}4\_\allowbreak{}a\\
Type: string\\
Default: nil
}

{\footnotesize <AC-03(04)\_\allowbreak{}ODP[01] дискреційна політика управління доступом> та <AC-03(04)\_\allowbreak{}ODP[02] дискреційна політика управління доступом> визначають, що суб'єкт, якому надано доступ до інформації, може передавати інформацію будь-яким іншим суб'єктам або об'єктам}


{\footnotesize\bfseries
No: 6\\
Name: ac\_\allowbreak{}3\_\allowbreak{}4\_\allowbreak{}b\\
Type: string\\
Default: nil
}

{\footnotesize <AC-03(04)\_\allowbreak{}ODP[01] дискреційна політика управління доступом> та <AC-03(04)\_\allowbreak{}ODP[02] дискреційна політика управління доступом> визначають, що суб'єкт, якому надано доступ до інформації, може надавати свої привілеї іншим суб'єктам}


{\footnotesize\bfseries
No: 7\\
Name: ac\_\allowbreak{}3\_\allowbreak{}4\_\allowbreak{}c\\
Type: string\\
Default: nil
}

{\footnotesize <AC-03(04)\_\allowbreak{}ODP[01] дискреційна політика управління доступом> та <AC-03(04)\_\allowbreak{}ODP[02] дискреційна політика управління доступом> визначають, що суб'єкт, якому надано доступ до інформації, може змінювати атрибути безпеки суб'єктів, об'єктів, системи або компонентів системи}


{\footnotesize\bfseries
No: 8\\
Name: ac\_\allowbreak{}3\_\allowbreak{}4\_\allowbreak{}d\\
Type: string\\
Default: nil
}

{\footnotesize <AC-03(04)\_\allowbreak{}ODP[01] дискреційна політика управління доступом> та <AC-03(04)\_\allowbreak{}ODP[02] дискреційна політика управління доступом> визначають, що суб'єкт, якому надано доступ до інформації, може вибирати атрибути безпеки, які будуть пов'язані з новоствореними або переглянутими об'єктами}


{\footnotesize\bfseries
No: 9\\
Name: ac\_\allowbreak{}3\_\allowbreak{}4\_\allowbreak{}e\\
Type: string\\
Default: nil
}

{\footnotesize <AC-03(04)\_\allowbreak{}ODP[01] дискреційна політика управління доступом> та <AC-03(04)\_\allowbreak{}ODP[02] дискреційна політика управління доступом> визначають, що суб'єкт, якому надано доступ до інформації, може змінювати правила управління доступом}




\subsubsection{ІНФОРМАЦІЯ ЩОДО БЕЗПЕКИ (AC-3(5))}
Запобігати доступу до [Призначення: інформації щодо безпеки, яка визначена організацією], за винятком випадків, коли наявні безпечні неробочі стани системи.

\vspace{10pt}
{\footnotesize\bfseries
No: 1\\
Name: ac\_\allowbreak{}3\_\allowbreak{}5\_\allowbreak{}01\\
Type: string\\
Default: nil
}

{\footnotesize Доступ до інформація щодо безпеки заборонено, за винятком випадків, коли наявні безпечні неробочі стани системи}


{\footnotesize\bfseries
No: 2\\
Name: ac\_\allowbreak{}3\_\allowbreak{}5\_\allowbreak{}odp\\
Type: string\\
Default: nil
}

{\footnotesize Визначено інформацію щодо безпеки, доступ до якої заборонено, за винятком випадків, коли наявні безпечні неробочі стани системи}




\subsubsection{ЗАХИСТ ІНФОРМАЦІЇ КОРИСТУВАЧА ТА СИСТЕМИ (AC-3(6)) [Вилучено]}
[Вилучено: Включено в MP-04 та SC-28]

\vspace{10pt}
\textit{Немає параметрів для цього контролю.}
\vspace{10pt}


\subsubsection{УПРАВЛІННЯ ДОСТУПОМ НА ОСНОВІ РОЛЕЙ (AC-3(7))}
Застосовувати політику управління доступом на основі ролей щодо визначених суб'єктів і об'єктів та управління доступом на основі [Призначення: визначених організацією ролей та користувачів, уповноважених приймати такі ролі].

\vspace{10pt}
{\footnotesize\bfseries
No: 1\\
Name: ac\_\allowbreak{}3\_\allowbreak{}7\_\allowbreak{}01\\
Type: list\\
Default: [\textquotedbl{}default\_\allowbreak{}deny\_\allowbreak{}rule\textquotedbl{}, \textquotedbl{}abac\_\allowbreak{}rule\_\allowbreak{}1\textquotedbl{}]
}

{\footnotesize Політика управління доступом на основі ролей застосовується до визначених суб'єктів}


{\footnotesize\bfseries
No: 2\\
Name: ac\_\allowbreak{}3\_\allowbreak{}7\_\allowbreak{}02\\
Type: list\\
Default: [\textquotedbl{}default\_\allowbreak{}deny\_\allowbreak{}rule\textquotedbl{}, \textquotedbl{}abac\_\allowbreak{}rule\_\allowbreak{}1\textquotedbl{}]
}

{\footnotesize Політика управління доступом на основі ролей застосовується до визначених об'єктів}


{\footnotesize\bfseries
No: 3\\
Name: ac\_\allowbreak{}3\_\allowbreak{}7\_\allowbreak{}03\\
Type: list\\
Default: [\textquotedbl{}admin\textquotedbl{}, \textquotedbl{}security\_\allowbreak{}officer\textquotedbl{}]
}

{\footnotesize Доступ контролюється на основі ролей та користувачів, яким дозволено приймати такі ролі}


{\footnotesize\bfseries
No: 4\\
Name: ac\_\allowbreak{}3\_\allowbreak{}7\_\allowbreak{}odp\_\allowbreak{}01\\
Type: list\\
Default: [\textquotedbl{}admin\textquotedbl{}, \textquotedbl{}security\_\allowbreak{}officer\textquotedbl{}]
}

{\footnotesize Визначено ролі, на яких базується управління доступом}


{\footnotesize\bfseries
No: 5\\
Name: ac\_\allowbreak{}3\_\allowbreak{}7\_\allowbreak{}odp\_\allowbreak{}02\\
Type: string\\
Default: nil
}

{\footnotesize Визначено користувачів, уповноважених на прийняття ролей (визначених у AC-03(07)\_\allowbreak{}ODP[01])}




\subsubsection{АНУЛЮВАННЯ ПРАВ ДОСТУПУ (AC-3(8))}
Здійснювати анулювання прав доступу в результаті змін атрибутів безпеки суб'єктів і об'єктів на основі [Призначення: визначених організацією правил, що регулюють терміни скасування прав доступу].

\vspace{10pt}
{\footnotesize\bfseries
No: 1\\
Name: ac\_\allowbreak{}3\_\allowbreak{}8\_\allowbreak{}01\\
Type: string\\
Default: nil
}

{\footnotesize Здійснюється анулювання прав доступу в результаті зміни атрибутів безпеки суб'єктів на основі правил}


{\footnotesize\bfseries
No: 2\\
Name: ac\_\allowbreak{}3\_\allowbreak{}8\_\allowbreak{}02\\
Type: string\\
Default: nil
}

{\footnotesize Здійснюється анулювання прав доступу в результаті зміни атрибутів безпеки об'єктів на основі пра- вил}


{\footnotesize\bfseries
No: 3\\
Name: ac\_\allowbreak{}3\_\allowbreak{}8\_\allowbreak{}odp\\
Type: list\\
Default: [\textquotedbl{}default\_\allowbreak{}deny\_\allowbreak{}rule\textquotedbl{}, \textquotedbl{}abac\_\allowbreak{}rule\_\allowbreak{}1\textquotedbl{}]
}

{\footnotesize Визначено правила, що регулюють терміни скасування дозволів на доступ}




\subsubsection{КЕРОВАНА ПЕРЕДАЧА (ПУБЛІКАЦІЯ) ІНФОРМАЦІЇ (AC-3(9))}
Передавати (публікувати) інформацію за межами встановленої межі системи можливо, якщо: a) Приймальна [Призначення: визначена організацією система або компонент системи] забезпечує [Призначення: визначені організацією заходи безпеки]; b) [Призначення: визначені організацією заходи безпеки] використовуються для підтвердження відповідності інформації, призначеної для керованих передач (публікації).

\vspace{10pt}
{\footnotesize\bfseries
No: 1\\
Name: ac\_\allowbreak{}3\_\allowbreak{}9\_\allowbreak{}a\\
Type: string\\
Default: nil
}

{\footnotesize Інформація випускається за межі системи, тільки якщо отримуюча система або компонент системи забезпечує заходи захисту}


{\footnotesize\bfseries
No: 2\\
Name: ac\_\allowbreak{}3\_\allowbreak{}9\_\allowbreak{}b\\
Type: string\\
Default: nil
}

{\footnotesize Інформація публікується за межами системи, тільки якщо заходи захисту використовуються для перевірки відповідності інформації, призначеної для публікації}




\subsubsection{ПЕРЕГЛЯД АУДИТОМ МЕХАНІЗМІВ КОНТРОЛЮ ДОСТУПУ (AC-3(10))}
Застосувати перегляд аудитом механізмів автоматизованого управління доступу при [Призначення: визначених організацією умовах] [Призначення: визначеними організацією ролями].

\vspace{10pt}
{\footnotesize\bfseries
No: 1\\
Name: ac\_\allowbreak{}3\_\allowbreak{}10\_\allowbreak{}01\\
Type: string\\
Default: nil
}

{\footnotesize За умов застосовується перегляд аудитом механізмів автоматизованого контролю доступу за допомогою ролей}


{\footnotesize\bfseries
No: 2\\
Name: ac\_\allowbreak{}3\_\allowbreak{}10\_\allowbreak{}odp\_\allowbreak{}01\\
Type: list\\
Default: []
}

{\footnotesize Визначено умови, за яких можна застосовувати перегляд аудитом механізмів автоматизованого управління доступом}


{\footnotesize\bfseries
No: 3\\
Name: ac\_\allowbreak{}3\_\allowbreak{}10\_\allowbreak{}odp\_\allowbreak{}02\\
Type: list\\
Default: [\textquotedbl{}admin\textquotedbl{}, \textquotedbl{}security\_\allowbreak{}officer\textquotedbl{}]
}

{\footnotesize Визначено ролі, яким дозволено використовувати перегляд аутом механізмів автоматизованого управління доступом}




\subsubsection{ОБМЕЖЕННЯ ДОСТУПУ ДО СПЕЦІАЛЬНОЇ ІНФОРМАЦІЇ (AC-3(11))}
Обмежити прямий доступ до сховищ даних, що містять [Призначення: визначені організацією типи інформації].

\vspace{10pt}
{\footnotesize\bfseries
No: 1\\
Name: ac\_\allowbreak{}3\_\allowbreak{}11\_\allowbreak{}01\\
Type: string\\
Default: nil
}

{\footnotesize Обмежено доступ до сховищ даних, що містять типи інформації}




\subsubsection{ВСТАНОВЛЕННЯ ТА ЗАБЕЗПЕЧЕННЯ ДОСТУПУ ДО ЗАСТОСУНКІВ (AC-3(12))}
a) Вимагати від застосунків встановити в процесі інсталяції доступ до таких застосунків системи і функцій: [Призначення: визначених організацією програм та функції системи]; b) Впровадити механізм примусового застосування, щоб запобігти доступу, відмінному від заявленого. c) Схвалити зміни доступу після початкового встановлення застосунків.

\vspace{10pt}
{\footnotesize\bfseries
No: 1\\
Name: ac\_\allowbreak{}3\_\allowbreak{}12\_\allowbreak{}a\\
Type: string\\
Default: nil
}

{\footnotesize У процесі інсталяції програми повинні встановити доступ до таких системних застосунків і функцій системи: програм та функції системи}


{\footnotesize\bfseries
No: 2\\
Name: ac\_\allowbreak{}3\_\allowbreak{}12\_\allowbreak{}b\\
Type: string\\
Default: nil
}

{\footnotesize Передбачено механізм примусового застосування запобігання несанкціонованому доступу}


{\footnotesize\bfseries
No: 3\\
Name: ac\_\allowbreak{}3\_\allowbreak{}12\_\allowbreak{}c\\
Type: string\\
Default: nil
}

{\footnotesize Зміни доступу після первинної інсталяції програми схвалено}


{\footnotesize\bfseries
No: 4\\
Name: ac\_\allowbreak{}3\_\allowbreak{}12\_\allowbreak{}odp\\
Type: string\\
Default: nil
}

{\footnotesize Визначено програми та функції системи, яким необхідно встановити права доступу}




\subsubsection{УПРАВЛІННЯ ДОСТУПОМ НА ОСНОВІ АТРИБУТІВ (AC-3(13))}
Здійснювати політику управління доступу на основі атрибутів (atribute-based) для визначених суб'єктів і об'єктів доступу й управляти доступом на основі [Призначення: визначених організацією атрибутів для ухвалення рішень про доступ].

\vspace{10pt}
{\footnotesize\bfseries
No: 1\\
Name: ac\_\allowbreak{}3\_\allowbreak{}13\_\allowbreak{}1\\
Type: list\\
Default: [\textquotedbl{}default\_\allowbreak{}deny\_\allowbreak{}rule\textquotedbl{}, \textquotedbl{}abac\_\allowbreak{}rule\_\allowbreak{}1\textquotedbl{}]
}

{\footnotesize Політика управління доступом на здійснюється до визначених суб'єктів; основі атрибутів AC-03(13)[2] політика управління доступом на здійснюється до визначених об'єктів; основі атрибутів}


{\footnotesize\bfseries
No: 2\\
Name: ac\_\allowbreak{}3\_\allowbreak{}13\_\allowbreak{}3\\
Type: string\\
Default: nil
}

{\footnotesize Доступ контролюється на основі атрибутів}


{\footnotesize\bfseries
No: 3\\
Name: ac\_\allowbreak{}3\_\allowbreak{}13\_\allowbreak{}odp\\
Type: list\\
Default: [\textquotedbl{}default\_\allowbreak{}deny\_\allowbreak{}rule\textquotedbl{}, \textquotedbl{}abac\_\allowbreak{}rule\_\allowbreak{}1\textquotedbl{}]
}

{\footnotesize Визначено атрибути для визначення прав доступу}




\subsubsection{ІНДИВІДУАЛЬНИЙ ДОСТУП (AC-3(14))}
Надайте [Призначення: механізми, визначені організацією], щоб дозволити особам мати доступ до певних елементів їх особистої інформації: [Призначення: елементи, визначені організацією]

\vspace{10pt}
{\footnotesize\bfseries
No: 1\\
Name: ac\_\allowbreak{}3\_\allowbreak{}14\_\allowbreak{}01\\
Type: list\\
Default: [\textquotedbl{}admin\textquotedbl{}, \textquotedbl{}security\_\allowbreak{}officer\textquotedbl{}]
}

{\footnotesize Механізми надаються для того, щоб дозволити особам мати доступ до елементів їхньої персональної інформації}


{\footnotesize\bfseries
No: 2\\
Name: ac\_\allowbreak{}3\_\allowbreak{}14\_\allowbreak{}odp\_\allowbreak{}01\\
Type: list\\
Default: [\textquotedbl{}admin\textquotedbl{}, \textquotedbl{}security\_\allowbreak{}officer\textquotedbl{}]
}

{\footnotesize Визначено механізми, що дозволяють фізичним особам мати доступ до елементів їхньої персональної інформації}


{\footnotesize\bfseries
No: 3\\
Name: ac\_\allowbreak{}3\_\allowbreak{}14\_\allowbreak{}odp\_\allowbreak{}02\\
Type: list\\
Default: [\textquotedbl{}admin\textquotedbl{}, \textquotedbl{}security\_\allowbreak{}officer\textquotedbl{}]
}

{\footnotesize Визначено елементи інформації, що ідентифікує особу, до якої мають доступ фізичні особи}




\subsubsection{ДИСКРЕЦІЙНИЙ ТА ОБОВ'ЯЗКОВИЙ ДОСТУП (AC-3(15))}
(a) Застосовувати [Призначення: визначену організацією політику обов'язкового контролю доступу] до набору охоплених суб'єктів і об'єктів, указаних у політиці;\\ 
(b) Застосування [Призначення: визначена організацією дискреційна політика контролю доступу] до набору охоплених суб'єктів і об'єктів, указаних у політиці.

\vspace{10pt}
{\footnotesize\bfseries
No: 1\\
Name: ac\_\allowbreak{}3\_\allowbreak{}15\_\allowbreak{}a\_\allowbreak{}01\\
Type: list\\
Default: [\textquotedbl{}default\_\allowbreak{}deny\_\allowbreak{}rule\textquotedbl{}, \textquotedbl{}abac\_\allowbreak{}rule\_\allowbreak{}1\textquotedbl{}]
}

{\footnotesize Політика обов'язкового контролю доступу застосовується до набору охоплених суб'єктів, зазначених у політиці}


{\footnotesize\bfseries
No: 2\\
Name: ac\_\allowbreak{}3\_\allowbreak{}15\_\allowbreak{}a\_\allowbreak{}02\\
Type: list\\
Default: [\textquotedbl{}default\_\allowbreak{}deny\_\allowbreak{}rule\textquotedbl{}, \textquotedbl{}abac\_\allowbreak{}rule\_\allowbreak{}1\textquotedbl{}]
}

{\footnotesize Політика обов'язкового контролю доступу застосовується до набору охоплених об'єктів, зазначених у політиці}


{\footnotesize\bfseries
No: 3\\
Name: ac\_\allowbreak{}3\_\allowbreak{}15\_\allowbreak{}b\_\allowbreak{}01\\
Type: list\\
Default: [\textquotedbl{}default\_\allowbreak{}deny\_\allowbreak{}rule\textquotedbl{}, \textquotedbl{}abac\_\allowbreak{}rule\_\allowbreak{}1\textquotedbl{}]
}

{\footnotesize Дискреційна політика контролю доступу застосовується до набору суб'єктів, зазначених у політиці}


{\footnotesize\bfseries
No: 4\\
Name: ac\_\allowbreak{}3\_\allowbreak{}15\_\allowbreak{}b\_\allowbreak{}02\\
Type: list\\
Default: [\textquotedbl{}default\_\allowbreak{}deny\_\allowbreak{}rule\textquotedbl{}, \textquotedbl{}abac\_\allowbreak{}rule\_\allowbreak{}1\textquotedbl{}]
}

{\footnotesize Дискреційна політика контролю доступу застосовується до набору об'єктів, зазначених у політиці}


{\footnotesize\bfseries
No: 5\\
Name: ac\_\allowbreak{}3\_\allowbreak{}15\_\allowbreak{}odp\_\allowbreak{}01\\
Type: list\\
Default: [\textquotedbl{}default\_\allowbreak{}deny\_\allowbreak{}rule\textquotedbl{}, \textquotedbl{}abac\_\allowbreak{}rule\_\allowbreak{}1\textquotedbl{}]
}

{\footnotesize Визначено обов'язкову політику контролю доступу, яка застосовується до набору суб'єктів, зазначених у політиці}


{\footnotesize\bfseries
No: 6\\
Name: ac\_\allowbreak{}3\_\allowbreak{}15\_\allowbreak{}odp\_\allowbreak{}02\\
Type: list\\
Default: [\textquotedbl{}default\_\allowbreak{}deny\_\allowbreak{}rule\textquotedbl{}, \textquotedbl{}abac\_\allowbreak{}rule\_\allowbreak{}1\textquotedbl{}]
}

{\footnotesize Визначено обов'язкову політику контролю доступу, яка застосовується до набору об'єктів, зазначених у політиці}


{\footnotesize\bfseries
No: 7\\
Name: ac\_\allowbreak{}3\_\allowbreak{}15\_\allowbreak{}odp\_\allowbreak{}03\\
Type: list\\
Default: [\textquotedbl{}default\_\allowbreak{}deny\_\allowbreak{}rule\textquotedbl{}, \textquotedbl{}abac\_\allowbreak{}rule\_\allowbreak{}1\textquotedbl{}]
}

{\footnotesize Визначено дискреційну політику контролю доступу, яка застосовується до набору суб'єктів, зазначених у політиці}


{\footnotesize\bfseries
No: 8\\
Name: ac\_\allowbreak{}3\_\allowbreak{}15\_\allowbreak{}odp\_\allowbreak{}04\\
Type: list\\
Default: [\textquotedbl{}default\_\allowbreak{}deny\_\allowbreak{}rule\textquotedbl{}, \textquotedbl{}abac\_\allowbreak{}rule\_\allowbreak{}1\textquotedbl{}]
}

{\footnotesize Визначено дискреційну політику контролю доступу, яка застосовується до набору об'єктів, зазначених у політиці}




\subsection{УПРАВЛІННЯ ІНФОРМАЦІЙНИМИ ПОТОКАМИ (AC-4)}
Застосувати затверджені повноваження для управління потоком інформації всередині системи та між пов'язаними системами на основі [Призначення: визначеними організацією політиками управління інформаційним потоком].

\vspace{10pt}
{\footnotesize\bfseries
No: 1\\
Name: ac\_\allowbreak{}4\_\allowbreak{}01\\
Type: list\\
Default: [\textquotedbl{}default\_\allowbreak{}deny\_\allowbreak{}rule\textquotedbl{}, \textquotedbl{}abac\_\allowbreak{}rule\_\allowbreak{}1\textquotedbl{}]
}

{\footnotesize Затверджені повноваження застосовуються для контролю потоку інформації всередині системи та між підключеними системами на основі політики управління інформаційними потоками}


{\footnotesize\bfseries
No: 2\\
Name: ac\_\allowbreak{}4\_\allowbreak{}odp\\
Type: list\\
Default: [\textquotedbl{}default\_\allowbreak{}deny\_\allowbreak{}rule\textquotedbl{}, \textquotedbl{}abac\_\allowbreak{}rule\_\allowbreak{}1\textquotedbl{}]
}

{\footnotesize Визначено політики управління інформаційними потоками всередині системи та між підключеними системами}




\subsubsection{АТРИБУТИ БЕЗПЕКИ ОБ'ЄКТУ (AC-4(1))}
Використовувати [Призначення: визначені організацією атрибути безпеки], пов'язані з [Призначення: визначеними організацією інформацією, джерелами та об'єктами призначення], щоб запровадити [Призначення: визначену організацією політику управління потоками інформації] як основу для ухвалення рішень щодо управління потоками.

\vspace{10pt}
{\footnotesize\bfseries
No: 1\\
Name: ac\_\allowbreak{}4\_\allowbreak{}1\_\allowbreak{}01\\
Type: list\\
Default: [\textquotedbl{}default\_\allowbreak{}deny\_\allowbreak{}rule\textquotedbl{}, \textquotedbl{}abac\_\allowbreak{}rule\_\allowbreak{}1\textquotedbl{}]
}

{\footnotesize Атрибути безпеки, пов'язані з об'єктами інформації, джерела об'єктів та об'єктами призначення, використовуються для забезпечення виконання політик управління інформаційними потоками як основи для прийняття рішень щодо управління потоками}


{\footnotesize\bfseries
No: 2\\
Name: ac\_\allowbreak{}4\_\allowbreak{}1\_\allowbreak{}02\\
Type: list\\
Default: [\textquotedbl{}default\_\allowbreak{}deny\_\allowbreak{}rule\textquotedbl{}, \textquotedbl{}abac\_\allowbreak{}rule\_\allowbreak{}1\textquotedbl{}]
}

{\footnotesize Атрибути конфіденційності, пов'язані з об'єктами інформації, джерела об'єктів та об'єктами призначення, використовуються для забезпечення виконання політик управління інформаційними потоками як основи для прийняття рішень щодо управління потоками}


{\footnotesize\bfseries
No: 3\\
Name: ac\_\allowbreak{}4\_\allowbreak{}1\_\allowbreak{}odp\_\allowbreak{}01\\
Type: list\\
Default: [\textquotedbl{}default\_\allowbreak{}deny\_\allowbreak{}rule\textquotedbl{}, \textquotedbl{}abac\_\allowbreak{}rule\_\allowbreak{}1\textquotedbl{}]
}

{\footnotesize Визначено атрибути безпеки, які будуть пов'язані з інформацією, джерелом та об'єктами призначення}


{\footnotesize\bfseries
No: 4\\
Name: ac\_\allowbreak{}4\_\allowbreak{}1\_\allowbreak{}odp\_\allowbreak{}02\\
Type: list\\
Default: [\textquotedbl{}default\_\allowbreak{}deny\_\allowbreak{}rule\textquotedbl{}, \textquotedbl{}abac\_\allowbreak{}rule\_\allowbreak{}1\textquotedbl{}]
}

{\footnotesize Визначено атрибути конфіденційності, які будуть пов'язані з інформацією, джерелом та об'єктами призначення}


{\footnotesize\bfseries
No: 5\\
Name: ac\_\allowbreak{}4\_\allowbreak{}1\_\allowbreak{}odp\_\allowbreak{}03\\
Type: string\\
Default: nil
}

{\footnotesize Визначено об'єкти інформації, які будуть пов'язані з атрибутами безпеки}


{\footnotesize\bfseries
No: 6\\
Name: ac\_\allowbreak{}4\_\allowbreak{}1\_\allowbreak{}odp\_\allowbreak{}04\\
Type: string\\
Default: nil
}

{\footnotesize Визначено об'єкти інформації, які будуть пов'язані з атрибутами конфіденційності}


{\footnotesize\bfseries
No: 7\\
Name: ac\_\allowbreak{}4\_\allowbreak{}1\_\allowbreak{}odp\_\allowbreak{}05\\
Type: string\\
Default: nil
}

{\footnotesize Визначено джерела об'єктів, які будуть пов'язані з атрибутами безпеки}


{\footnotesize\bfseries
No: 8\\
Name: ac\_\allowbreak{}4\_\allowbreak{}1\_\allowbreak{}odp\_\allowbreak{}06\\
Type: string\\
Default: nil
}

{\footnotesize Визначено джерела об'єктів, які будуть пов'язані з атрибутами конфіденційності}


{\footnotesize\bfseries
No: 9\\
Name: ac\_\allowbreak{}4\_\allowbreak{}1\_\allowbreak{}odp\_\allowbreak{}07\\
Type: string\\
Default: nil
}

{\footnotesize Визначено об'єкти призначення, які будуть пов'язані з атрибутами безпеки}


{\footnotesize\bfseries
No: 10\\
Name: ac\_\allowbreak{}4\_\allowbreak{}1\_\allowbreak{}odp\_\allowbreak{}08\\
Type: string\\
Default: nil
}

{\footnotesize Визначено об'єкти призначення, які будуть пов'язані з атрибутами конфіденційності}


{\footnotesize\bfseries
No: 11\\
Name: ac\_\allowbreak{}4\_\allowbreak{}1\_\allowbreak{}odp\_\allowbreak{}09\\
Type: list\\
Default: [\textquotedbl{}default\_\allowbreak{}deny\_\allowbreak{}rule\textquotedbl{}, \textquotedbl{}abac\_\allowbreak{}rule\_\allowbreak{}1\textquotedbl{}]
}

{\footnotesize Визначено політику управління інформаційними потоками як основу для ухвалення рішень щодо управління потоками}




\subsubsection{ДОМЕНИ ОБРОБРОБКИ ДАНИХ (AC-4(2))}
Використовувати захищені домени обробки даних для забезпечення [Призначення: визначеної організацією політики управління потоками інформації] як основу для ухвалення рішень щодо управління потоками.

\vspace{10pt}
{\footnotesize\bfseries
No: 1\\
Name: ac\_\allowbreak{}4\_\allowbreak{}2\_\allowbreak{}01\\
Type: list\\
Default: [\textquotedbl{}default\_\allowbreak{}deny\_\allowbreak{}rule\textquotedbl{}, \textquotedbl{}abac\_\allowbreak{}rule\_\allowbreak{}1\textquotedbl{}]
}

{\footnotesize Захищені домени обробки використовуються для забезпечення дотримання політики управління інформаційними потоками як основи для ухвалення рішень щодо управління потоками}


{\footnotesize\bfseries
No: 2\\
Name: ac\_\allowbreak{}4\_\allowbreak{}2\_\allowbreak{}odp\\
Type: list\\
Default: [\textquotedbl{}default\_\allowbreak{}deny\_\allowbreak{}rule\textquotedbl{}, \textquotedbl{}abac\_\allowbreak{}rule\_\allowbreak{}1\textquotedbl{}]
}

{\footnotesize Визначено політики управління інформаційними потоками, які будуть застосовуватися з використанням захищених доменів обробки}




\subsubsection{ДИНАМІЧНЕ УПРАВЛІННЯ ІНФОРМАЦІЙНИМ ПОТОКОМ (AC-4(3))}
Здійснювати динамічне управління потоком інформації на основі [Призначення: визначених організацією політик (правил)].

\vspace{10pt}
{\footnotesize\bfseries
No: 1\\
Name: ac\_\allowbreak{}4\_\allowbreak{}3\_\allowbreak{}01\\
Type: list\\
Default: [\textquotedbl{}default\_\allowbreak{}deny\_\allowbreak{}rule\textquotedbl{}, \textquotedbl{}abac\_\allowbreak{}rule\_\allowbreak{}1\textquotedbl{}]
}

{\footnotesize Здійснюється динамічне управління потоком інформації на основі політики управління інформаційними потоками}


{\footnotesize\bfseries
No: 2\\
Name: ac\_\allowbreak{}4\_\allowbreak{}3\_\allowbreak{}odp\\
Type: list\\
Default: [\textquotedbl{}default\_\allowbreak{}deny\_\allowbreak{}rule\textquotedbl{}, \textquotedbl{}abac\_\allowbreak{}rule\_\allowbreak{}1\textquotedbl{}]
}

{\footnotesize Визначені політики контролю інформаційних потоків, які необхідно впроваджувати}




\subsubsection{УПРАВЛІННЯ ПОТОКОМ ЗАШИФРОВАНОЇ ІНФОРМАЦІЇ (AC-4(4))}
Запобігати обходу [Призначення: механізмів управління потоками, визначених організацією] зашифрованої інформації шляхом [Вибір (один або декілька): дешифрування інформації; блокування потоку зашифрованої інформації; завершення сеансів зв'язку, що намагаються передавати зашифровану інформацію; [Призначення: визначеними організацією процедурою або методом]].

\vspace{10pt}
{\footnotesize\bfseries
No: 1\\
Name: ac\_\allowbreak{}4\_\allowbreak{}4\_\allowbreak{}odp\_\allowbreak{}01\\
Type: string\\
Default: \textquotedbl{}автоматизований засіб моніторингу\textquotedbl{}
}

{\footnotesize Визначено механізми контролю інформаційних потоків, які унеможливлюють обхід зашифрованої інформації}


{\footnotesize\bfseries
No: 2\\
Name: ac\_\allowbreak{}4\_\allowbreak{}4\_\allowbreak{}odp\_\allowbreak{}02\\
Type: string\\
Default: \textquotedbl{}AES-256-GCM\textquotedbl{}
}

{\footnotesize Вибрано одне або декілька з наступних ЗНАЧЕНЬ ПАРАМЕТРІВ: \{дешифрування інформації; блокування потоку зашифрованої інформації; завершення сеансів зв'язку що намагаються передавати зашифровану інформацію; визначена організацією процедура або метод\}}


{\footnotesize\bfseries
No: 3\\
Name: ac\_\allowbreak{}4\_\allowbreak{}4\_\allowbreak{}odp\_\allowbreak{}03\\
Type: string\\
Default: \textquotedbl{}автоматизований засіб моніторингу\textquotedbl{}
}

{\footnotesize Визначено організацією процедуру або метод, що використовується для запобігання обходу зашифрованої інформації через механізми контролю інформаційних потоків (якщо вибрано)}




\subsubsection{ВБУДОВУВАННЯ ТИПІВ ДАНИХ (AC-4(5))}
Впровадити [Призначення: визначені організацією вбудовування типів даних в інші типи даних. обмеження] для

\vspace{10pt}
{\footnotesize\bfseries
No: 1\\
Name: ac\_\allowbreak{}4\_\allowbreak{}5\_\allowbreak{}odp\\
Type: string\\
Default: nil
}

{\footnotesize Визначеного обмеження, які слід застосовувати щодо вбудовування типів даних в інші типи даних; обмеження накладаються на вбудовування типів даних у інші типи даних}




\subsubsection{МЕТАДАНІ (AC-4(6))}
Здійснювати управління інформаційним потоком на основі [Призначення: визначених організацією метаданих].

\vspace{10pt}
{\footnotesize\bfseries
No: 1\\
Name: ac\_\allowbreak{}4\_\allowbreak{}6\_\allowbreak{}01\\
Type: string\\
Default: nil
}

{\footnotesize Інформаційна система здійснює управління інформаційним потоком на основі метадані}


{\footnotesize\bfseries
No: 2\\
Name: ac\_\allowbreak{}4\_\allowbreak{}6\_\allowbreak{}odp\\
Type: string\\
Default: nil
}

{\footnotesize Визначено метадані, які слід використовувати як засіб управління інформаційним потоком}




\subsubsection{МЕХАНІЗМИ ОДНОСТОРОННЬОГО ПОТОКУ (AC-4(7))}
Впровадити [Призначення: визначені організацією односторонні інформаційні потоки] за допомогою апаратних механізмів.

\vspace{10pt}
{\footnotesize\bfseries
No: 1\\
Name: ac\_\allowbreak{}4\_\allowbreak{}7\_\allowbreak{}01\\
Type: string\\
Default: nil
}

{\footnotesize Односторонні інформаційні потоки забезпечуються за допомогою апаратних механізмів управління потоками. ПОТЕНЦІЙНІ МЕТОДИ ТА ОБ'ЄКТИ ОЦІНКИ:}




\subsubsection{ФІЛЬТРИ ПОЛІТИКИ БЕЗПЕКИ (AC-4(8))}
a) Забезпечити контроль над потоком інформації, використовуючи [Призначення: визначені організацією фільтри безпеки або політики конфіденційності] як основу для рішень щодо керування потоком для [Призначення: визначені організацією потоки інформації]; b) [Вибір (один або кілька): Блокування; Зміна; Карантин] даних після помилки обробки фільтра відповідно до [Призначення: політика безпеки або конфіденційності, визначена організацією].

\vspace{10pt}
{\footnotesize\bfseries
No: 1\\
Name: ac\_\allowbreak{}4\_\allowbreak{}8\_\allowbreak{}01\\
Type: list\\
Default: [\textquotedbl{}default\_\allowbreak{}deny\_\allowbreak{}rule\textquotedbl{}, \textquotedbl{}abac\_\allowbreak{}rule\_\allowbreak{}1\textquotedbl{}]
}

{\footnotesize ODP[02] визначено фільтри політики конфіденційності, які будуть використовуватися як основа для забезпечення керування інформаційними потоками}


{\footnotesize\bfseries
No: 2\\
Name: ac\_\allowbreak{}4\_\allowbreak{}8\_\allowbreak{}a\_\allowbreak{}01\\
Type: list\\
Default: [\textquotedbl{}default\_\allowbreak{}deny\_\allowbreak{}rule\textquotedbl{}, \textquotedbl{}abac\_\allowbreak{}rule\_\allowbreak{}1\textquotedbl{}]
}

{\footnotesize Управління інформаційними потоками здійснюється за допомогою <AC-04(08) \_\allowbreak{}ODP[01] фільтр політики безпеки> як основи для прийняття рішень щодо управління потоками для <AC-04(08) \_\allowbreak{}ODP[03] інформаційних потоків>; AC-04(08)(a)[01] контроль інформаційних потоків здійснюється за допомогою <AC-04(08) \_\allowbreak{}ODP[02] фільтр політики конфіденційності> як основи для прийняття рішень щодо контролю потоків для <AC-04(08) \_\allowbreak{}ODP[04] інформаційних потоків>}




\subsubsection{ПЕРЕВІРКИ, ЩО ПРОВОДИТЬ ПЕРСОНАЛ (AC-4(9))}
Примусово використовувати перевірку персоналом [Призначення: потоки інформації, визначені організацією] за таких умов: [Призначення: умови, визначені організацією].

\vspace{10pt}
{\footnotesize\bfseries
No: 1\\
Name: ac\_\allowbreak{}4\_\allowbreak{}9\_\allowbreak{}01\\
Type: list\\
Default: [\textquotedbl{}admin\textquotedbl{}, \textquotedbl{}security\_\allowbreak{}officer\textquotedbl{}]
}

{\footnotesize Перевірка персоналом використовуються для інформаційних потоків за умов}


{\footnotesize\bfseries
No: 2\\
Name: ac\_\allowbreak{}4\_\allowbreak{}9\_\allowbreak{}odp\_\allowbreak{}01\\
Type: list\\
Default: [\textquotedbl{}admin\textquotedbl{}, \textquotedbl{}security\_\allowbreak{}officer\textquotedbl{}]
}

{\footnotesize Визначено інформаційні потоки, які потребують використання перевірку персоналом}


{\footnotesize\bfseries
No: 3\\
Name: ac\_\allowbreak{}4\_\allowbreak{}9\_\allowbreak{}odp\_\allowbreak{}02\\
Type: list\\
Default: [\textquotedbl{}admin\textquotedbl{}, \textquotedbl{}security\_\allowbreak{}officer\textquotedbl{}]
}

{\footnotesize Визначено умови, за яких використання перевірки персоналом на інформаційні потоки має бути обов'язковим}




\subsubsection{АКТИВАЦІЯ ТА ДЕАКТИВАЦІЯ ФІЛЬТРІВ ПОЛІТИКИ БЕЗПЕКИ (AC-4(10))}
Впровадити можливість для привілейованих адміністраторів активувати та деактивувати [Призначення: фільтри політики безпеки, що визначаються організацією] за таких умов: [Призначення: визначені організацією умови].

\vspace{10pt}
{\footnotesize\bfseries
No: 1\\
Name: ac\_\allowbreak{}4\_\allowbreak{}10\_\allowbreak{}01\\
Type: string\\
Default: nil
}

{\footnotesize Привілейованим адміністраторам надано можливість активувати та деактивувати <AC-04(10) \_\allowbreak{}ODP[01] фільтри безпеки> за <AC-04(10) \_\allowbreak{}ODP[03] умов>}


{\footnotesize\bfseries
No: 2\\
Name: ac\_\allowbreak{}4\_\allowbreak{}10\_\allowbreak{}02\\
Type: string\\
Default: nil
}

{\footnotesize Привілейованим адміністраторам надано можливість активувати та деактивувати <AC-04(10) \_\allowbreak{}ODP[02] фільтри конфіденційності> за <AC-04(10) \_\allowbreak{}ODP[04] умов>}




\subsubsection{КОНФІГУРАЦІЯ ФІЛЬТРІВ ПОЛІТИКИ БЕЗПЕКИ (AC-4(11))}
Впровадити можливість для привілейованих адміністраторів налаштувати [Призначення: визначені організацією фільтри політики безпеки] для підтримки різних політик безпеки.

\vspace{10pt}
{\footnotesize\bfseries
No: 1\\
Name: ac\_\allowbreak{}4\_\allowbreak{}11\_\allowbreak{}01\\
Type: list\\
Default: [\textquotedbl{}default\_\allowbreak{}deny\_\allowbreak{}rule\textquotedbl{}, \textquotedbl{}abac\_\allowbreak{}rule\_\allowbreak{}1\textquotedbl{}]
}

{\footnotesize Привілейованим адміністраторам надано можливість налаштовувати фільтри політики безпеки для підтримки різних політик безпеки або конфіденційності}


{\footnotesize\bfseries
No: 2\\
Name: ac\_\allowbreak{}4\_\allowbreak{}11\_\allowbreak{}02\\
Type: list\\
Default: [\textquotedbl{}default\_\allowbreak{}deny\_\allowbreak{}rule\textquotedbl{}, \textquotedbl{}abac\_\allowbreak{}rule\_\allowbreak{}1\textquotedbl{}]
}

{\footnotesize Привілейованим адміністраторам надано можливість налаштовувати фільтри політики конфіденційності для підтримки різних політик безпеки або конфіденційності}


{\footnotesize\bfseries
No: 3\\
Name: ac\_\allowbreak{}4\_\allowbreak{}11\_\allowbreak{}odp\_\allowbreak{}01\\
Type: list\\
Default: [\textquotedbl{}default\_\allowbreak{}deny\_\allowbreak{}rule\textquotedbl{}, \textquotedbl{}abac\_\allowbreak{}rule\_\allowbreak{}1\textquotedbl{}]
}

{\footnotesize Визначено фільтри політики безпеки, які привілейовані адміністратори можуть налаштовувати для підтримки різних політик безпеки та конфіденційності}


{\footnotesize\bfseries
No: 4\\
Name: ac\_\allowbreak{}4\_\allowbreak{}11\_\allowbreak{}odp\_\allowbreak{}02\\
Type: list\\
Default: [\textquotedbl{}default\_\allowbreak{}deny\_\allowbreak{}rule\textquotedbl{}, \textquotedbl{}abac\_\allowbreak{}rule\_\allowbreak{}1\textquotedbl{}]
}

{\footnotesize Визначено фільтри політики конфіденційності, які привілейовані адміністратори можуть налаштовувати для підтримки різних політик безпеки та конфіденційності}




\subsubsection{ІДЕНТИФІКАТОРИ ТИПУ ДАНИХ (AC-4(12))}
При передачі інформації між різними захищеними доменами використовувати [Призначення: визначені організацією ідентифікатори типів даних] для перевірки даних, необхідних для ухвалення рішень щодо інформаційного потоку.

\vspace{10pt}
{\footnotesize\bfseries
No: 1\\
Name: ac\_\allowbreak{}4\_\allowbreak{}12\_\allowbreak{}odp\\
Type: string\\
Default: nil
}

{\footnotesize Визначено ідентифікатори типів даних, які будуть використовуватися для перевірки даних, необхідних для ухвалення рішень щодо інформаційних потоків; AC-04(12) при передачі інформації між різними доменами безпеки, ідентифікатори типів даних використовуються для перевірки даних, необхідних для ухвалення рішень щодо інформаційних потоків}




\subsubsection{ДЕКОМПОЗИЦІЯ НА ВІДПОВІДНІ ПОЛІТИЦІ СУБКОМПОНЕНТИ (AC-4(13))}
При передачі інформації між різними захищеними доменами здійснювати декомпозицію інформації на [Призначення: визначені організацією субкомпоненти, що відповідають політиці] для представлення в механізмах реалізації політики.

\vspace{10pt}
{\footnotesize\bfseries
No: 1\\
Name: ac\_\allowbreak{}4\_\allowbreak{}13\_\allowbreak{}01\\
Type: list\\
Default: [\textquotedbl{}default\_\allowbreak{}deny\_\allowbreak{}rule\textquotedbl{}, \textquotedbl{}abac\_\allowbreak{}rule\_\allowbreak{}1\textquotedbl{}]
}

{\footnotesize При передачі інформації між різними доменами безпеки інформація розкладається на субкомпоненти політики для подання механізмам забезпечення дотримання політики}


{\footnotesize\bfseries
No: 2\\
Name: ac\_\allowbreak{}4\_\allowbreak{}13\_\allowbreak{}odp\\
Type: list\\
Default: [\textquotedbl{}default\_\allowbreak{}deny\_\allowbreak{}rule\textquotedbl{}, \textquotedbl{}abac\_\allowbreak{}rule\_\allowbreak{}1\textquotedbl{}]
}

{\footnotesize Визначено субкомпоненти політики, на які слід розкласти інформацію для подання до механізмів реалізації політики}




\subsubsection{ОБМЕЖЕННЯ ФІЛЬТРА ПОЛІТИКИ БЕЗПЕКИ (AC-4(14))}
При передачі інформації між різними захищеними доменами реалізувати [Призначення: визначені організацією фільтри політики безпеки], що вимагають повного переліку форматів, які обмежують структуру та зміст даних.

\vspace{10pt}
{\footnotesize\bfseries
No: 1\\
Name: ac\_\allowbreak{}4\_\allowbreak{}14\_\allowbreak{}01\\
Type: list\\
Default: [\textquotedbl{}default\_\allowbreak{}deny\_\allowbreak{}rule\textquotedbl{}, \textquotedbl{}abac\_\allowbreak{}rule\_\allowbreak{}1\textquotedbl{}]
}

{\footnotesize При передачі інформації між різними захищеними доменами, реалізовані фільтри політики безпеки вимагають повністю перелічених форматів, які обмежують структуру та вміст даних}


{\footnotesize\bfseries
No: 2\\
Name: ac\_\allowbreak{}4\_\allowbreak{}14\_\allowbreak{}02\\
Type: list\\
Default: [\textquotedbl{}default\_\allowbreak{}deny\_\allowbreak{}rule\textquotedbl{}, \textquotedbl{}abac\_\allowbreak{}rule\_\allowbreak{}1\textquotedbl{}]
}

{\footnotesize При передачі інформації між різними захищеними доменами, реалізовані фільтри політики конфіденційності вимагають повністю перелічених форматів, які обмежують структуру та вміст даних}


{\footnotesize\bfseries
No: 3\\
Name: ac\_\allowbreak{}4\_\allowbreak{}14\_\allowbreak{}odp\_\allowbreak{}02\\
Type: list\\
Default: [\textquotedbl{}default\_\allowbreak{}deny\_\allowbreak{}rule\textquotedbl{}, \textquotedbl{}abac\_\allowbreak{}rule\_\allowbreak{}1\textquotedbl{}]
}

{\footnotesize Визначено фільтри політики конфіденційності, які вимагають повного переліку форматів, що обмежують структуру та зміст даних}




\subsubsection{ВИЯВЛЕННЯ НЕСАНКЦІОНОВАНОЇ ІНФОРМАЦІЇ (AC-4(15))}
При передачі інформації між різними захищеними доменами перевіряти інформацію на наявність [Призначення: визначеної організацією несанкціонованої інформації] та забороняти передачу такої інформації відповідно до [Призначення: визначеної організацією політики безпеки].

\vspace{10pt}
{\footnotesize\bfseries
No: 1\\
Name: ac\_\allowbreak{}4\_\allowbreak{}15\_\allowbreak{}01\\
Type: string\\
Default: nil
}

{\footnotesize При передачі інформації між різними доменами безпеки інформація перевіряється на наявність <AC- 04(15)\_\allowbreak{}ODP[01] несанкціонованої інформації>}


{\footnotesize\bfseries
No: 2\\
Name: ac\_\allowbreak{}4\_\allowbreak{}15\_\allowbreak{}02\\
Type: list\\
Default: [\textquotedbl{}default\_\allowbreak{}deny\_\allowbreak{}rule\textquotedbl{}, \textquotedbl{}abac\_\allowbreak{}rule\_\allowbreak{}1\textquotedbl{}]
}

{\footnotesize При передачі інформації між різними доменами безпеки забороняється передача несанкціонованої інформації відповідно до політики безпеки}


{\footnotesize\bfseries
No: 3\\
Name: ac\_\allowbreak{}4\_\allowbreak{}15\_\allowbreak{}03\\
Type: list\\
Default: [\textquotedbl{}default\_\allowbreak{}deny\_\allowbreak{}rule\textquotedbl{}, \textquotedbl{}abac\_\allowbreak{}rule\_\allowbreak{}1\textquotedbl{}]
}

{\footnotesize При передачі інформації між різними доменами безпеки забороняється передача несанкціонованої інформації відповідно до політики конфіденційності}


{\footnotesize\bfseries
No: 4\\
Name: ac\_\allowbreak{}4\_\allowbreak{}15\_\allowbreak{}odp\_\allowbreak{}01\\
Type: string\\
Default: nil
}

{\footnotesize Визначено несанкціоновану інформацію, яку потрібно виявляти}


{\footnotesize\bfseries
No: 5\\
Name: ac\_\allowbreak{}4\_\allowbreak{}15\_\allowbreak{}odp\_\allowbreak{}02\\
Type: list\\
Default: [\textquotedbl{}default\_\allowbreak{}deny\_\allowbreak{}rule\textquotedbl{}, \textquotedbl{}abac\_\allowbreak{}rule\_\allowbreak{}1\textquotedbl{}]
}

{\footnotesize Визначено політику безпеки, яка вимагає заборонити передачу несанкціонованої інформації між різними доменами безпеки (якщо вибрано)}


{\footnotesize\bfseries
No: 6\\
Name: ac\_\allowbreak{}4\_\allowbreak{}15\_\allowbreak{}odp\_\allowbreak{}03\\
Type: list\\
Default: [\textquotedbl{}default\_\allowbreak{}deny\_\allowbreak{}rule\textquotedbl{}, \textquotedbl{}abac\_\allowbreak{}rule\_\allowbreak{}1\textquotedbl{}]
}

{\footnotesize Визначено політику конфіденційності, яка вимагає заборонити передачу визначеної організацією несанкціонованої інформації між різними доменами безпеки (якщо вибрано)}




\subsubsection{ПЕРЕДАЧА ІНФОРМАЦІЇ ПРО ВЗАЄМОПОВ'ЯЗАНІ СИСТЕМИ (AC-4(16)) [Вилучено]}
[Вилучено: Включено в АС-04]

\vspace{10pt}
\textit{Немає параметрів для цього контролю.}
\vspace{10pt}


\subsubsection{АВТЕНТИФІКАЦІЯ ДОМЕНУ (AC-4(17))}
Автентифікувати [Призначення: визначені організацією домени] до [Призначення: дії, визначеної організацією] з обміну інформацією.

\vspace{10pt}
{\footnotesize\bfseries
No: 1\\
Name: ac\_\allowbreak{}4\_\allowbreak{}17\_\allowbreak{}odp\\
Type: string\\
Default: nil
}

{\footnotesize Вибрано одне або декілька з наступних ЗНАЧЕНЬ ПАРАМЕТРІВ: \{організація, система, програми, служби, індивід\};}


{\footnotesize\bfseries
No: 2\\
Name: ac\_\allowbreak{}4\_\allowbreak{}17\_\allowbreak{}01\\
Type: string\\
Default: nil
}

{\footnotesize Для передачі інформації пункти відправлення та призначення унікально ідентифікуються та аутентифікуються за допомогою <AC-04(17)\_\allowbreak{}ODP ВИБРАНЕ ЗНАЧЕННЯ(Я) ПАРАМЕТРА(ів)>}




\subsubsection{ПРИВ'ЯЗКА АТРИБУТУ БЕЗПЕКИ (AC-4(18)) [Вилучено]}
[Вилучено: Включено в АС-16]

\vspace{10pt}
\textit{Немає параметрів для цього контролю.}
\vspace{10pt}


\subsubsection{ПЕРЕВІРКА МЕТАДАНИХ (AC-4(19))}
Під час передачі інформації між різними захищеними доменами застосовувати до метаданих ту ж політику безпеки фільтрації, що й для корисних даних.

\vspace{10pt}
{\footnotesize\bfseries
No: 1\\
Name: ac\_\allowbreak{}4\_\allowbreak{}19\_\allowbreak{}odp\_\allowbreak{}01\\
Type: string\\
Default: nil
}

{\footnotesize Визначено фільтри політики безпеки, які буде застосовано до метаданих}


{\footnotesize\bfseries
No: 2\\
Name: ac\_\allowbreak{}4\_\allowbreak{}19\_\allowbreak{}odp\_\allowbreak{}02\\
Type: string\\
Default: nil
}

{\footnotesize Визначено фільтри політики конфіденційності, які буде застосовано до метаданих}


{\footnotesize\bfseries
No: 3\\
Name: ac\_\allowbreak{}4\_\allowbreak{}19\_\allowbreak{}01\\
Type: string\\
Default: nil
}

{\footnotesize При передачі інформації між різними доменами безпеки, <AC-04(19)\_\allowbreak{}ODP[01] фільтри політики безпеки> реалізовано на метаданих}


{\footnotesize\bfseries
No: 4\\
Name: ac\_\allowbreak{}4\_\allowbreak{}19\_\allowbreak{}02\\
Type: string\\
Default: nil
}

{\footnotesize При передачі інформації між різними доменами безпеки, <AC-04(19)\_\allowbreak{}ODP[02] фільтри політики конфіденційності> реалізовано на метаданих}




\subsubsection{ЗАТВЕРДЖЕНІ РІШЕННЯ (AC-4(20))}
Впровадити [Призначення: визначені організацією рішення про схвалені конфігурації] для керування потоком [Призначення: інформації, визначеної організацією] через захищені домени.

\vspace{10pt}
{\footnotesize\bfseries
No: 1\\
Name: ac\_\allowbreak{}4\_\allowbreak{}20\_\allowbreak{}01\\
Type: string\\
Default: nil
}

{\footnotesize <AC-04(20)\_\allowbreak{}ODP[01] рішення> використовуються для контролю потоку <AC-04(20)\_\allowbreak{}ODP[02] інформації> між захищеними доменами}


{\footnotesize\bfseries
No: 2\\
Name: ac\_\allowbreak{}4\_\allowbreak{}20\_\allowbreak{}odp\_\allowbreak{}01\\
Type: string\\
Default: nil
}

{\footnotesize Визначені рішення про схвалені конфігурації для керування потоками інформації через захищені домени}


{\footnotesize\bfseries
No: 3\\
Name: ac\_\allowbreak{}4\_\allowbreak{}20\_\allowbreak{}odp\_\allowbreak{}02\\
Type: string\\
Default: nil
}

{\footnotesize Визначено інформацію, якою потрібно керувати, коли вона проходить через захищені домени}




\subsubsection{ФІЗИЧНЕ ТА ЛОГІЧНЕ ВІДДІЛЕННЯ ІНФОРМАЦІЙНИХ ПОТОКІВ (AC-4(21))}
Відокремлювати потоки інформації логічно або фізично, використовуючи [Призначення: визначені організацією механізми та/або методи] для досягнення [Призначення: визначеного організацією необхідного поділу за типами інформації].

\vspace{10pt}
{\footnotesize\bfseries
No: 1\\
Name: ac\_\allowbreak{}4\_\allowbreak{}21\_\allowbreak{}odp\_\allowbreak{}01\\
Type: string\\
Default: \textquotedbl{}автоматизований засіб моніторингу\textquotedbl{}
}

{\footnotesize Визначено механізми та/або методи, що використовуються для логічного розділення інформаційних потоків}


{\footnotesize\bfseries
No: 2\\
Name: ac\_\allowbreak{}4\_\allowbreak{}21\_\allowbreak{}odp\_\allowbreak{}02\\
Type: string\\
Default: \textquotedbl{}автоматизований засіб моніторингу\textquotedbl{}
}

{\footnotesize Визначено механізми та/або методи, що використовуються для фізичного розділення інформаційних потоків}


{\footnotesize\bfseries
No: 3\\
Name: ac\_\allowbreak{}4\_\allowbreak{}21\_\allowbreak{}odp\_\allowbreak{}03\\
Type: string\\
Default: nil
}

{\footnotesize Визначено необхідні поділи за типами інформації}


{\footnotesize\bfseries
No: 4\\
Name: ac\_\allowbreak{}4\_\allowbreak{}21\_\allowbreak{}01\\
Type: string\\
Default: nil
}

{\footnotesize Інформаційні потоки логічно розділені за допомогою <AC-04(21)\_\allowbreak{}ODP[01] механізмів та/або методів> для досягнення <AC-04(21)\_\allowbreak{}ODP[03] необхідного поділу за типами інформації>}


{\footnotesize\bfseries
No: 5\\
Name: ac\_\allowbreak{}4\_\allowbreak{}21\_\allowbreak{}02\\
Type: string\\
Default: nil
}

{\footnotesize Інформаційні потоки фізично розділені за допомогою <AC-04(21)\_\allowbreak{}ODP[02] механізмів та/або методів> для досягнення <AC-04(21)\_\allowbreak{}ODP[03] необхідного поділу за типами інформації>}




\subsubsection{ЄДИНИЙ ДОСТУП (AC-4(22))}
Забезпечити доступ з одного пристрою до обчислювальних платформ, застосунків або даних, що розташовуються в декількох різних захищених доменах, одночасно запобігаючи передачі будь-якого потоку інформації між різними захищеними доменами.

\vspace{10pt}
{\footnotesize\bfseries
No: 1\\
Name: ac\_\allowbreak{}4\_\allowbreak{}22\_\allowbreak{}01\\
Type: integer\\
Default: 30
}

{\footnotesize Доступ забезпечується з одного пристрою до обчислювальних платформ, застосунків або даних, що розташовуються в декількох різних захищених доменах, одночасно запобігаючи передачі інформації між різними захищеними доменами}




\subsubsection{МОДИФІКОВАНА ІНФОРМАЦІЯ, ЯКА НЕ ПІДЛЯГАЄ ОПРИЛЮДНЕННЮ (AC-4(23))}
Під час передачі інформації між різними доменами безпеки змінюйте інформацію, яка не підлягає оприлюдненню, реалізувавши [Призначення: визначена організацією дія модифікації]

\vspace{10pt}
{\footnotesize\bfseries
No: 1\\
Name: ac\_\allowbreak{}4\_\allowbreak{}23\_\allowbreak{}odp\\
Type: \\
Default: \textquotedbl{}\textquotedbl{}
}

{\footnotesize Визначено дію модифікації, що застосовується до інформації, яка не підлягає оприлюдненню;}


{\footnotesize\bfseries
No: 2\\
Name: ac\_\allowbreak{}4\_\allowbreak{}23\_\allowbreak{}01\\
Type: string\\
Default: nil
}

{\footnotesize AC-04(23) при передачі інформації між доменами безпеки інформація, що не підлягає оприлюдненню, модифікується шляхом реалізації дія модифікації}




\subsubsection{ВНУТРІШНІЙ НОРМАЛІЗОВАНИЙ ФОРМАТ (AC-4(24))}
Під час передачі інформації між різними доменами безпеки аналізуйте вхідні дані у внутрішньому нормалізованому форматі та повторно генеруйте дані, щоб вони відповідали призначеній специфікації.

\vspace{10pt}
{\footnotesize\bfseries
No: 1\\
Name: ac\_\allowbreak{}4\_\allowbreak{}24\_\allowbreak{}1\\
Type: string\\
Default: nil
}

{\footnotesize При передачі інформації між різними доменами безпеки вхідні дані розбираються у внутрішній, нормалізований формат}


{\footnotesize\bfseries
No: 2\\
Name: ac\_\allowbreak{}4\_\allowbreak{}24\_\allowbreak{}2\\
Type: string\\
Default: nil
}

{\footnotesize При передачі інформації між різними доменами безпеки дані регенеруються, щоб відповідати їхній специфікації}




\subsubsection{ОЧИЩЕННЯ ДАНИХ (AC-4(25))}
Під час передачі інформації між різними доменами безпеки очищуйте дані, щоб мінімізувати [Вибір (один або кілька): доставка зловмисного вмісту, керування та керування зловмисним кодом, доповнення зловмисного коду та стеганографічно закодовані дані; витік конфіденційної інформації] відповідно до [Призначення: політика, визначена організацією]].

\vspace{10pt}
{\footnotesize\bfseries
No: 1\\
Name: ac\_\allowbreak{}4\_\allowbreak{}25\_\allowbreak{}odp\_\allowbreak{}01\\
Type: string\\
Default: nil
}

{\footnotesize Вибрано одне або декілька з наступних ЗНАЧЕНЬ ПАРАМЕТРІВ: \{доставка шкідливого коду, керування та контроль шкідливого коду, доповнення шкідливого коду та даних, закодованих стеганографією; витік конфіденційної інформації\}}


{\footnotesize\bfseries
No: 2\\
Name: ac\_\allowbreak{}4\_\allowbreak{}25\_\allowbreak{}odp\_\allowbreak{}02\\
Type: string\\
Default: nil
}

{\footnotesize Визначено політику очищення даних}


{\footnotesize\bfseries
No: 3\\
Name: ac\_\allowbreak{}4\_\allowbreak{}25\_\allowbreak{}01\\
Type: string\\
Default: nil
}

{\footnotesize Під час передачі інформації між різними доменами безпеки дані очищуються, щоб мінімізувати <AC-04(25)\_\allowbreak{}ODP[01] вибрані ризики> відповідно до <AC-04(25)\_\allowbreak{}ODP[02] політики>}




\subsubsection{ДІЇ З ФІЛЬТРАЦІЇ АУДИТУ (AC-4(26))}
Під час передачі інформації між різними доменами безпеки записуйте та перевіряйте дії фільтрації вмісту та результати для інформації, що фільтрується.

\vspace{10pt}
{\footnotesize\bfseries
No: 1\\
Name: ac\_\allowbreak{}4\_\allowbreak{}26\_\allowbreak{}01\\
Type: string\\
Default: nil
}

{\footnotesize При передачі інформації між різними доменами безпеки дії з фільтрації вмісту фіксуються і перевіряються}


{\footnotesize\bfseries
No: 2\\
Name: ac\_\allowbreak{}4\_\allowbreak{}26\_\allowbreak{}02\\
Type: string\\
Default: nil
}

{\footnotesize При передачі інформації між різними доменами безпеки, результати для інформації, що фільтрується, записуються і перевіряються}




\subsubsection{НАДЛИШКОВІ/НЕЗАЛЕЖНІ ФІЛЬТРУЮЧІ МЕХАНІЗМИ (AC-4(27))}
Під час передачі інформації між різними доменами безпеки впроваджуйте рішення фільтрації вмісту, які забезпечують надлишкові та незалежні механізми фільтрації для кожного типу даних.

\vspace{10pt}
{\footnotesize\bfseries
No: 1\\
Name: ac\_\allowbreak{}4\_\allowbreak{}27\_\allowbreak{}01\\
Type: integer\\
Default: 30
}

{\footnotesize Під час передачі інформації між системами безпеки впроваджені рішення для фільтрації контенту забезпечують надлишкові та незалежні механізми фільтрації для кожного типу даних}




\subsubsection{ЛІНІЙНІ ФІЛЬТРУВАЛЬНІ КАНАЛИ (AC-4(28))}
Під час передачі інформації між різними доменами безпеки запровадьте конвеєр лінійного фільтрування вмісту, який забезпечується дискреційним і обов'язковим контролем доступу.

\vspace{10pt}
{\footnotesize\bfseries
No: 1\\
Name: ac\_\allowbreak{}4\_\allowbreak{}28\_\allowbreak{}01\\
Type: string\\
Default: nil
}

{\footnotesize При передачі інформації між доменами безпеки реалізовано лінійний конвеєр фільтрації контенту, який забезпечується дискретними та обов'язковими засобами контролю доступу}




\subsubsection{ФІЛЬТР МЕХАНІЗМІВ ОРКЕСТРОВКИ (AC-4(29))}
Під час передачі інформації між різними доменами безпеки використовуйте механізми оркестровки фільтрів вмісту, щоб забезпечити:\\ 
a. Механізми фільтрації вмісту успішно завершують виконання без помилок;\\ 
b. Дії фільтрації вмісту виконуються в правильному порядку та відповідають [Призначення: політика, визначена організацією]

\vspace{10pt}
{\footnotesize\bfseries
No: 1\\
Name: ac\_\allowbreak{}4\_\allowbreak{}29\_\allowbreak{}a\\
Type: string\\
Default: \textquotedbl{}автоматизований засіб моніторингу\textquotedbl{}
}

{\footnotesize При передачі інформації між доменами безпеки використовуються механізми оркестрування фільтрації контенту, які гарантують, що механізми фільтрації контенту успішно завершать виконання без помилок}


{\footnotesize\bfseries
No: 2\\
Name: ac\_\allowbreak{}4\_\allowbreak{}29\_\allowbreak{}b\_\allowbreak{}01\\
Type: list\\
Default: [\textquotedbl{}login\textquotedbl{}, \textquotedbl{}logout\textquotedbl{}, \textquotedbl{}failed\_\allowbreak{}attempt\textquotedbl{}]
}

{\footnotesize При передачі інформації між доменами безпеки використовуються механізми оркестрування фільтрації контенту, які гарантують, що дії з фільтрації контенту відбуваються в правильному порядку}


{\footnotesize\bfseries
No: 3\\
Name: ac\_\allowbreak{}4\_\allowbreak{}29\_\allowbreak{}b\_\allowbreak{}02\\
Type: list\\
Default: [\textquotedbl{}login\textquotedbl{}, \textquotedbl{}logout\textquotedbl{}, \textquotedbl{}failed\_\allowbreak{}attempt\textquotedbl{}]
}

{\footnotesize При передачі інформації між доменами безпеки використовуються механізми оркестрування фільтрації контенту, які гарантують, що дії з фільтрації контенту відповідають політиці. }


{\footnotesize\bfseries
No: 4\\
Name: ac\_\allowbreak{}4\_\allowbreak{}29\_\allowbreak{}odp\\
Type: list\\
Default: [\textquotedbl{}default\_\allowbreak{}deny\_\allowbreak{}rule\textquotedbl{}, \textquotedbl{}abac\_\allowbreak{}rule\_\allowbreak{}1\textquotedbl{}]
}

{\footnotesize Визначено політику щодо дій з фільтрації контенту}




\subsubsection{МЕХАНІЗМИ ФІЛЬТРАЦІЇ З ВИКОРИСТАННЯМ КІЛЬКОХ ПРОЦЕСІВ (AC-4(30))}
Під час передачі інформації між різними доменами безпеки реалізуйте механізми фільтрації вмісту за допомогою кількох процесів.

\vspace{10pt}
{\footnotesize\bfseries
No: 1\\
Name: ac\_\allowbreak{}4\_\allowbreak{}30\_\allowbreak{}01\\
Type: string\\
Default: \textquotedbl{}автоматизований засіб моніторингу\textquotedbl{}
}

{\footnotesize При передачі інформації між доменами безпеки реалізовані механізми контент-фільтрації з використанням декількох процесів. ПОТЕНЦІЙНІ МЕТОДИ ТА ОБ'ЄКТИ ОЦІНКИ:}




\subsubsection{ЗАПОБІГАННЯ СПРОБАМ ПЕРЕДАЧІ ВМІСТУ, ЯКИЙ НЕ ПРОЙШОВ ПЕРЕВІРКУ ФІЛЬТРАЦІЇ (AC-4(31))}
Під час передачі інформації між різними доменами безпеки запобігайте передачі вмісту, який не пройшов перевірку фільтрації до домену-одержувача.

\vspace{10pt}
{\footnotesize\bfseries
No: 1\\
Name: ac\_\allowbreak{}4\_\allowbreak{}31\_\allowbreak{}01\\
Type: string\\
Default: nil
}

{\footnotesize При передачі інформації між різними доменами безпеки запобігається передача вмісту який не пройшов фільрацію}




\subsubsection{ВИМОГИ ДО ПРОЦЕСУ ПЕРЕДАЧІ ІНФОРМАЦІЇ (AC-4(32))}
Під час передачі інформації між різними доменами безпеки, процес, який передає інформацію між конвеєрами фільтрації:\\ 
a. не фільтрує вміст повідомлення;\\ 
b. перевіряє метадані фільтрації;\\ 
c. забезпечує успішне завершення фільтрації вмісту, пов'язаного з метаданими фільтрації; і\\ 
d. передає вміст до цільового фільтруючого конвеєра.

\vspace{10pt}
{\footnotesize\bfseries
No: 1\\
Name: ac\_\allowbreak{}4\_\allowbreak{}32\_\allowbreak{}01\\
Type: string\\
Default: nil
}

{\footnotesize Під час передачі інформації між різними доменами безпеки, процес, який передає інформацію між конвеєрами фільтрації, не фільтрує вміст повідомлення}


{\footnotesize\bfseries
No: 2\\
Name: ac\_\allowbreak{}4\_\allowbreak{}32\_\allowbreak{}02\\
Type: string\\
Default: nil
}

{\footnotesize Під час передачі інформації між різними доменами безпеки, процес, який передає інформацію між конвеєрами фільтрації, перевіряє метадані фільтрації}


{\footnotesize\bfseries
No: 3\\
Name: ac\_\allowbreak{}4\_\allowbreak{}32\_\allowbreak{}03\\
Type: string\\
Default: nil
}

{\footnotesize Під час передачі інформації між різними доменами безпеки, процес, який передає інформацію між конвеєрами фільтрації, забезпечує успішне завершення фільтрації вмісту, пов'язаного з метаданими фільтрації}


{\footnotesize\bfseries
No: 4\\
Name: ac\_\allowbreak{}4\_\allowbreak{}32\_\allowbreak{}04\\
Type: string\\
Default: nil
}

{\footnotesize Під час передачі інформації між різними доменами безпеки, процес, який передає інформацію між конвеєрами фільтрації, передає вміст до цільового фільтруючого конвеєра}




\subsection{РОЗМЕЖУВАННЯ ОБОВ'ЯЗКІВ (AC-5)}
a. Розмежувати і документувати [Призначення: визначені організацією обов'язки окремих осіб].\\ 
b. Установити правила авторизації доступу для підтримки розмежування обов'язків.

\vspace{10pt}
{\footnotesize\bfseries
No: 1\\
Name: ac\_\allowbreak{}5\_\allowbreak{}odp\_\allowbreak{}01\\
Type: string\\
Default: nil
}

{\footnotesize Визначено обов'язки осіб, які потребують розмежування}


{\footnotesize\bfseries
No: 2\\
Name: ac\_\allowbreak{}5\_\allowbreak{}01\\
Type: string\\
Default: nil
}

{\footnotesize Обов'язки осіб <AC-05\_\allowbreak{}ODP[01] визначені організацією> розмежовані та задокументовані}


{\footnotesize\bfseries
No: 3\\
Name: ac\_\allowbreak{}5\_\allowbreak{}02\\
Type: string\\
Default: nil
}

{\footnotesize Встановлено правила авторизації доступу до системи для підтримки розмежування обов'язків}




\subsection{МІНІМІЗАЦІЯ ПОВНОВАЖЕНЬ (AC-6)}
Впровадити принцип мінімізації повноважень, який дозволяє користувачам (або процесам, що діють від імені користувачів) здійснювати лише такі авторизовані звернення, які необхідні для виконання визначених завдань відповідно до цілей (призначення, місії) організації та функцій.

\vspace{10pt}
{\footnotesize\bfseries
No: 1\\
Name: ac\_\allowbreak{}6\_\allowbreak{}01\\
Type: string\\
Default: nil
}

{\footnotesize Застосовується принцип мінімалізації повноважень, який дозволяє користувачам (або процесам, що діють від імені користувачів) здійснювати лише такі авторизовані звернення, які необхідні для виконання поставлених завдань організації}




\subsubsection{АВТОРИЗОВАНИЙ ДОСТУП ДО ФУНКЦІЙ БЕЗПЕКИ (AC-6(1))}
Авторизувати доступ до [Призначення: визначені організацією функції безпеки (розгорнуті в апаратному, програмному та мікропрограмному забезпеченні)] та [Призначення: визначена організацією інформація, що має відношення до безпеки] для [Призначення: визначені організацією особи або ролі].

\vspace{10pt}
{\footnotesize\bfseries
No: 1\\
Name: ac\_\allowbreak{}6\_\allowbreak{}1\_\allowbreak{}odp\_\allowbreak{}01\\
Type: string\\
Default: nil
}

{\footnotesize Визначені особи або ролі з авторизованим доступом до функцій безпеки та інформації, що має відношення до безпеки}


{\footnotesize\bfseries
No: 2\\
Name: ac\_\allowbreak{}6\_\allowbreak{}1\_\allowbreak{}odp\_\allowbreak{}02\\
Type: string\\
Default: nil
}

{\footnotesize Визначені функції безпеки (розгорнуті в апаратному забезпеченні) для авторизованого доступу}


{\footnotesize\bfseries
No: 3\\
Name: ac\_\allowbreak{}6\_\allowbreak{}1\_\allowbreak{}odp\_\allowbreak{}03\\
Type: string\\
Default: nil
}

{\footnotesize Визначені функції безпеки (розгорнуті в програмному забезпеченні) для авторизованого доступу}


{\footnotesize\bfseries
No: 4\\
Name: ac\_\allowbreak{}6\_\allowbreak{}1\_\allowbreak{}odp\_\allowbreak{}04\\
Type: string\\
Default: nil
}

{\footnotesize Визначені функції безпеки (розгорнуті в мікропрограмному забезпеченні) для авторизованого доступу}


{\footnotesize\bfseries
No: 5\\
Name: ac\_\allowbreak{}6\_\allowbreak{}1\_\allowbreak{}odp\_\allowbreak{}05\\
Type: string\\
Default: nil
}

{\footnotesize Визначено інформацію, важливу для забезпечення безпеки, для авторизованого доступу}


{\footnotesize\bfseries
No: 6\\
Name: ac\_\allowbreak{}6\_\allowbreak{}1\_\allowbreak{}a\_\allowbreak{}01\\
Type: string\\
Default: nil
}

{\footnotesize Авторизовано доступ для <AC-06(01)\_\allowbreak{}ODP[01] осіб та ролей> до <AC-06(01)\_\allowbreak{}ODP[02] функцій безпеки (розгорнутих на апаратному забезпеченні)>}


{\footnotesize\bfseries
No: 7\\
Name: ac\_\allowbreak{}6\_\allowbreak{}1\_\allowbreak{}a\_\allowbreak{}02\\
Type: string\\
Default: nil
}

{\footnotesize Авторизовано доступ для <AC-06(01)\_\allowbreak{}ODP[01] осіб та ролей> до <AC-06(01)\_\allowbreak{}ODP[03] функцій безпеки (розгорнутих на програмному забезпеченні)>}


{\footnotesize\bfseries
No: 8\\
Name: ac\_\allowbreak{}6\_\allowbreak{}1\_\allowbreak{}a\_\allowbreak{}03\\
Type: string\\
Default: nil
}

{\footnotesize Авторизовано доступ для <AC-06(01)\_\allowbreak{}ODP[01] осіб та ролей> до <AC-06(01)\_\allowbreak{}ODP[04] функцій безпеки (розгорнутих на мікропрограмному забезпеченні)>}


{\footnotesize\bfseries
No: 9\\
Name: ac\_\allowbreak{}6\_\allowbreak{}1\_\allowbreak{}b\\
Type: string\\
Default: nil
}

{\footnotesize Авторизовано доступ для <AC-06(01)\_\allowbreak{}ODP[01] осіб та ролей> до <AC-06(01)\_\allowbreak{}ODP[05] інформації, що має відношення до безпеки>}




\subsubsection{НЕПРИВІЛЕЙОВАНИЙ ДОСТУП ДО НЕЗАХИЩЕНИХ ФУНКЦІЙ (AC-6(2))}
Вимагати від користувачів облікових записів системи або ролей, які мають доступ до [Призначення: визначених організацією функцій безпеки або інформації, що стосується безпеки], використовувати непривілейовані облікові записи чи ролі під час доступу до незахищених функцій.

\vspace{10pt}
{\footnotesize\bfseries
No: 1\\
Name: ac\_\allowbreak{}6\_\allowbreak{}2\_\allowbreak{}01\\
Type: list\\
Default: [\textquotedbl{}admin\textquotedbl{}, \textquotedbl{}security\_\allowbreak{}officer\textquotedbl{}]
}

{\footnotesize Користувачі облікових записів (або ролей) системи з доступом до функцій безпеки або інформації, що стосується безпеки, повинні використовувати непривілейовані облікові записи або ролі під час доступу до незахищених функцій}


{\footnotesize\bfseries
No: 2\\
Name: ac\_\allowbreak{}6\_\allowbreak{}2\_\allowbreak{}odp\\
Type: string\\
Default: nil
}

{\footnotesize Визначені функції безпеки або інформація, що стосується безпеки}




\subsubsection{МЕРЕЖЕВИЙ ДОСТУП ДО ПРИВІЛЕЙОВАНИХ КОМАНД (AC-6(3))}
Авторизувати мережевий доступ до [Призначення: визначені організацією привілейовані команди] лише для [Призначення: визначені організацією потреби] та задокументувати обґрунтування такого доступу в плані безпеки системи.

\vspace{10pt}
{\footnotesize\bfseries
No: 1\\
Name: ac\_\allowbreak{}6\_\allowbreak{}3\_\allowbreak{}odp\_\allowbreak{}01\\
Type: string\\
Default: nil
}

{\footnotesize Визначено привілейовані команди для яких дозволяється мережевий доступ}


{\footnotesize\bfseries
No: 2\\
Name: ac\_\allowbreak{}6\_\allowbreak{}3\_\allowbreak{}odp\_\allowbreak{}02\\
Type: string\\
Default: nil
}

{\footnotesize Визначено потреби для яких авторизується мережевий доступ до привілейованих команд}


{\footnotesize\bfseries
No: 3\\
Name: ac\_\allowbreak{}6\_\allowbreak{}3\_\allowbreak{}01\\
Type: string\\
Default: nil
}

{\footnotesize Мережевий доступ до <AC-06(03)\_\allowbreak{}ODP[01] привілейованих команд> авторизовано лише для <AC-06(03)\_\allowbreak{}ODP[02] потреб>}


{\footnotesize\bfseries
No: 4\\
Name: ac\_\allowbreak{}6\_\allowbreak{}3\_\allowbreak{}02\\
Type: string\\
Default: nil
}

{\footnotesize Обґрунтування авторизованого мережевого доступу до привілейованих команд задокументовано в плані безпеки системи}




\subsubsection{РОЗДІЛЬНІ ДОМЕНИ ОБРОБКИ (AC-6(4))}
Надати окремі домени обробки даних для забезпечення більш точного розподілу повноважень користувача.

\vspace{10pt}
{\footnotesize\bfseries
No: 1\\
Name: ac\_\allowbreak{}6\_\allowbreak{}4\_\allowbreak{}01\\
Type: string\\
Default: nil
}

{\footnotesize Надаються окремі домени обробки для більш тонкого розподілу повноважень користувачів}




\subsubsection{ПРИВІЛЕЙОВАНІ ОБЛІКОВІ ЗАПИСИ (AC-6(5))}
Обмежити привілейовані облікові записи в системі згідно з [Призначення: визначеним організацією персоналом або ролями].

\vspace{10pt}
{\footnotesize\bfseries
No: 1\\
Name: ac\_\allowbreak{}6\_\allowbreak{}5\_\allowbreak{}01\\
Type: list\\
Default: [\textquotedbl{}admin\textquotedbl{}, \textquotedbl{}security\_\allowbreak{}officer\textquotedbl{}]
}

{\footnotesize Привілейовані облікові записи в системі обмежено персонал або ролі}


{\footnotesize\bfseries
No: 2\\
Name: ac\_\allowbreak{}6\_\allowbreak{}5\_\allowbreak{}odp\\
Type: list\\
Default: [\textquotedbl{}admin\textquotedbl{}, \textquotedbl{}security\_\allowbreak{}officer\textquotedbl{}]
}

{\footnotesize Визначено персонал або ролі, яким мають бути обмежені привілейовані облікові записи в системі}




\subsubsection{ПРИВІЛЕЙОВАНИЙ ДОСТУП КОРИСТУВАЧАМИ, ЩО НЕ НАЛЕЖАТЬ ДО ОРГАНІЗАЦІЇ (AC-6(6))}
Заборонити привілейований доступ до системи користувачам, які не належать до організації.

\vspace{10pt}
{\footnotesize\bfseries
No: 1\\
Name: ac\_\allowbreak{}6\_\allowbreak{}6\_\allowbreak{}01\\
Type: string\\
Default: nil
}

{\footnotesize Привілейований доступ до системи для користувачів, які не є членами організації, заборонено}




\subsubsection{ПЕРЕГЛЯД ПОВНОВАЖЕНЬ КОРИСТУВАЧА (AC-6(7))}
a) Переглядати [Призначення: з визначеною організацією частотою] повноваження призначених для [Призначення: визначених організацією посад або класів користувачів] для перевірки необхідності таких повноважень; b) За необхідності перепризначити або зняти повноваження, правильного відображення цілей (місії) організації та потреб організації. для

\vspace{10pt}
{\footnotesize\bfseries
No: 1\\
Name: ac\_\allowbreak{}6\_\allowbreak{}7\_\allowbreak{}a\\
Type: integer\\
Default: 30
}

{\footnotesize Повноваження, призначені ролям і класам, переглядаються з частотою для перевірки необхідності таких повноважень}


{\footnotesize\bfseries
No: 2\\
Name: ac\_\allowbreak{}6\_\allowbreak{}7\_\allowbreak{}b\\
Type: string\\
Default: nil
}

{\footnotesize Привілеї перепризначаються або знімаються, якщо це необхідно, для правильного відображення місії організації та потреб}


{\footnotesize\bfseries
No: 3\\
Name: ac\_\allowbreak{}6\_\allowbreak{}7\_\allowbreak{}odp\_\allowbreak{}01\\
Type: integer\\
Default: 30
}

{\footnotesize Визначено частоту перегляду повноважень, призначених ролям або класам користувачів}


{\footnotesize\bfseries
No: 4\\
Name: ac\_\allowbreak{}6\_\allowbreak{}7\_\allowbreak{}odp\_\allowbreak{}02\\
Type: list\\
Default: [\textquotedbl{}admin\textquotedbl{}, \textquotedbl{}security\_\allowbreak{}officer\textquotedbl{}]
}

{\footnotesize Визначено ролі або класи користувачів, яким призначено повноваження }




\subsubsection{РІВНІ ПРИВІЛЕЇВ ДЛЯ ВИКОНАННЯ КОДУ (AC-6(8))}
Запобігати виконанню програмного забезпечення на рівні привілеїв вищому, ніж доступний користувачеві, який використовує програмне забезпечення [Призначення: визначене організацією програмне забезпечення].

\vspace{10pt}
{\footnotesize\bfseries
No: 1\\
Name: ac\_\allowbreak{}6\_\allowbreak{}8\_\allowbreak{}odp\_\allowbreak{}01\\
Type: string\\
Default: nil
}

{\footnotesize Визначене програмне забезпечення для якого запобігається виконання на вищому рівні привілеїв}


{\footnotesize\bfseries
No: 2\\
Name: ac\_\allowbreak{}6\_\allowbreak{}8\_\allowbreak{}01\\
Type: string\\
Default: nil
}

{\footnotesize Запобігається виконання <AC-06(08)\_\allowbreak{}ODP[01] програмного забезпечення> на рівні привілеїв вищому, ніж доступний користувачеві, який його використовує}




\subsubsection{АУДИТ ВИКОРИСТАННЯ ПРИВІЛЕЙОВАНИХ ФУНКЦІЙ (AC-6(9))}
Реєструвати виконання привілейованих функцій.

\vspace{10pt}
{\footnotesize\bfseries
No: 1\\
Name: ac\_\allowbreak{}6\_\allowbreak{}9\_\allowbreak{}01\\
Type: string\\
Default: nil
}

{\footnotesize Проводиться аудит виконання привілейованих функцій}




\subsubsection{ЗАБОРОНА НЕПРИВІЛЕЙОВАНИМ КОРИСТУВАЧАМ ВИКОНУВАТИ ПРИВІЛЕЙОВАНІ ФУНКЦІЇ (AC-6(10))}
Вжити заходи для запобігання можливості виконувати привілейовані функції непривілейованими користувачами.

\vspace{10pt}
{\footnotesize\bfseries
No: 1\\
Name: ac\_\allowbreak{}6\_\allowbreak{}10\_\allowbreak{}01\\
Type: string\\
Default: nil
}

{\footnotesize Непривілейовані користувачі не можуть виконувати привілейовані функції}




\subsection{НЕВДАЛІ СПРОБИ ВХОДУ В СИСТЕМУ (AC-7)}
a. Встановити обмеження на [Призначення: визначену організацією кількість] послідовних неуспішних спроб входу користувача в систему впродовж [Призначення: визначеного організацією часового періоду].\\ 
b. Автоматично виконати [Вибір (один або декілька): блокування облікового запису/вузла на [Призначення: визначений організацією часовий період]; блокування облікового запису/вузла, доки він не буде розблокований адміністратором; затримання наступної команди входу в систему за [Надання: визначеним організацією алгоритмом затримки]; виконати [Призначення: визначені організацією дії]], коли перевищено максимальну кількість невдалих спроб входу в систему.

\vspace{10pt}
{\footnotesize\bfseries
No: 1\\
Name: ac\_\allowbreak{}7\_\allowbreak{}odp\_\allowbreak{}01\\
Type: integer\\
Default: 3
}

{\footnotesize Визначено кількість послідовних неуспішних спроб входу користувача, дозволених протягом певного періоду часу}


{\footnotesize\bfseries
No: 2\\
Name: ac\_\allowbreak{}7\_\allowbreak{}odp\_\allowbreak{}02\\
Type: string\\
Default: \textquotedbl{}15m\textquotedbl{}
}

{\footnotesize Визначено період часу, яким обмежується кількість послідовних неуспішних спроб входу користувача}


{\footnotesize\bfseries
No: 3\\
Name: ac\_\allowbreak{}7\_\allowbreak{}odp\_\allowbreak{}03\\
Type: string\\
Default: nil
}

{\footnotesize Вибрано одне або декілька з наступних ЗНАЧЕНЬ ПАРАМЕТРІВ: \{заблокувати обліковий запис або вузол на період часу; заблокувати обліковий запис або вузол до зняття адміністратором; затримати наступний запит на вхід за алгоритмом затримки; повідомити системного адміністратора; виконати іншу дію\}}


{\footnotesize\bfseries
No: 4\\
Name: ac\_\allowbreak{}7\_\allowbreak{}odp\_\allowbreak{}04\\
Type: string\\
Default: \textquotedbl{}15m\textquotedbl{}
}

{\footnotesize Період часу, на який буде заблоковано обліковий запис або вузол (якщо вибрано)}


{\footnotesize\bfseries
No: 5\\
Name: ac\_\allowbreak{}7\_\allowbreak{}odp\_\allowbreak{}05\\
Type: string\\
Default: nil
}

{\footnotesize Визначено алгоритм затримки наступного запиту на вхід (якщо вибрано)}


{\footnotesize\bfseries
No: 6\\
Name: ac\_\allowbreak{}7\_\allowbreak{}odp\_\allowbreak{}06\\
Type: string\\
Default: nil
}

{\footnotesize Інша дія, яка буде виконана після перевищення максимальної кількості невдалих спроб (якщо вибрано)}


{\footnotesize\bfseries
No: 7\\
Name: ac\_\allowbreak{}7\_\allowbreak{}01\\
Type: string\\
Default: nil
}

{\footnotesize Встановлено обмеження на <AC-07\_\allowbreak{}ODP[01] кількість> послідовних неуспішних спроб входу користувача в систему впродовж <AC-07\_\allowbreak{}ODP[02] часового періоду>}


{\footnotesize\bfseries
No: 8\\
Name: ac\_\allowbreak{}7\_\allowbreak{}02\\
Type: string\\
Default: nil
}

{\footnotesize Автоматично виконується <AC-07\_\allowbreak{}ODP[03] вибрана дія> (із врахуванням <AC-07\_\allowbreak{}ODP[04] періоду блокування>, <AC-07\_\allowbreak{}ODP[05] алгоритму затримки> або <AC-07\_\allowbreak{}ODP[06] інших дій>), коли перевищено максимальну кількість невдалих спроб входу в систему}




\subsubsection{АВТОМАТИЧНЕ БЛОКУВАННЯ ОБЛІКОВОГО ЗАПИСУ (AC-7(1)) [Вилучено]}
[Вилучено: Включено в АС-07]

\vspace{10pt}
\textit{Немає параметрів для цього контролю.}
\vspace{10pt}


\subsubsection{ОЧИЩЕННЯ АБО СТИРАННЯ МОБІЛЬНОГО ПРИСТРОЮ (AC-7(2))}
Очистити або стерти інформацію з [Призначення: визначених організацією мобільних пристроїв] на основі [Призначення: визначених організацією вимог та методик очищення чи стирання] після [Призначення: визначеної організацією кількості] послідовних невдалих спроб входу в систему з пристрою.

\vspace{10pt}
{\footnotesize\bfseries
No: 1\\
Name: ac\_\allowbreak{}7\_\allowbreak{}2\_\allowbreak{}01\\
Type: integer\\
Default: 3
}

{\footnotesize Інформація очищується або стирається з мобільних пристроїв на основі вимог або методів очищення або стирання після <AC-07(02)\_\allowbreak{}ODP[03 кількість> послідовних, невдалих спроб входу на пристрій}


{\footnotesize\bfseries
No: 2\\
Name: ac\_\allowbreak{}7\_\allowbreak{}2\_\allowbreak{}odp\_\allowbreak{}01\\
Type: string\\
Default: nil
}

{\footnotesize Визначено мобільні пристрої, які підлягають очищенню або стиранню інформації}


{\footnotesize\bfseries
No: 3\\
Name: ac\_\allowbreak{}7\_\allowbreak{}2\_\allowbreak{}odp\_\allowbreak{}02\\
Type: string\\
Default: nil
}

{\footnotesize Визначено вимоги та методи очищення чи стирання інформації з мобільних пристроїв}


{\footnotesize\bfseries
No: 4\\
Name: ac\_\allowbreak{}7\_\allowbreak{}2\_\allowbreak{}odp\_\allowbreak{}03\\
Type: integer\\
Default: 3
}

{\footnotesize Визначається кількість послідовних невдалих спроб входу в систему до того, як інформація буде очищена або стерта з мобільних пристроїв}




\subsubsection{ОБМЕЖЕННЯ НА СПРОБИ БІОМЕТРИЧНОГО ВХОДУ (AC-7(3))}
Обмежити кількість невдалих спроб входу за допомогою біометрики [Призначення: визначена організацією кількість].

\vspace{10pt}
{\footnotesize\bfseries
No: 1\\
Name: ac\_\allowbreak{}7\_\allowbreak{}3\_\allowbreak{}odp\_\allowbreak{}01\\
Type: integer\\
Default: 3
}

{\footnotesize Визначено кількість невдалих спроб входу за допомогою біометрики}


{\footnotesize\bfseries
No: 2\\
Name: ac\_\allowbreak{}7\_\allowbreak{}3\_\allowbreak{}01\\
Type: string\\
Default: nil
}

{\footnotesize Обмежено кількість невдалих спроб входу за допомогою біометрики до <AC-07(03)\_\allowbreak{}ODP[01] кількості>}




\subsubsection{ВИКОРИСТАННЯ АЛЬТЕРНАТИВНОГО ФАКТОРА (AC-7(4))}
a) Дозволити використання [Призначення: визначені організацією фактори автентифікації], які відрізняються від основних факторів автентифікації після перевищення визначеної організацією кількості послідовних невдалих спроб входу в систему; b) Обмежити [Призначення: визначена організацією кількість] послідовних невдалих спроб входу за допомогою використання альтернативних факторів користувачем протягом [Призначення: визначеного організацією періоду часу].

\vspace{10pt}
{\footnotesize\bfseries
No: 1\\
Name: ac\_\allowbreak{}7\_\allowbreak{}4\_\allowbreak{}odp\_\allowbreak{}01\\
Type: string\\
Default: nil
}

{\footnotesize Визначено фактори автентифікації, які дозволено використовувати, що відрізняються від основних факторів}


{\footnotesize\bfseries
No: 2\\
Name: ac\_\allowbreak{}7\_\allowbreak{}4\_\allowbreak{}odp\_\allowbreak{}02\\
Type: integer\\
Default: 3
}

{\footnotesize Визначено кількість послідовних невдалих спроб входу за допомогою альтернативних факторів}


{\footnotesize\bfseries
No: 3\\
Name: ac\_\allowbreak{}7\_\allowbreak{}4\_\allowbreak{}odp\_\allowbreak{}03\\
Type: string\\
Default: \textquotedbl{}15m\textquotedbl{}
}

{\footnotesize Визначено період часу, протягом якого обмежуються спроби через альтернативні фактори}


{\footnotesize\bfseries
No: 4\\
Name: ac\_\allowbreak{}7\_\allowbreak{}4\_\allowbreak{}a\\
Type: string\\
Default: nil
}

{\footnotesize Дозволяється використання <AC-07(04)\_\allowbreak{}ODP[01] факторів автентифікації> після перевищення визначеної організацією кількості послідовних невдалих спроб входу}


{\footnotesize\bfseries
No: 5\\
Name: ac\_\allowbreak{}7\_\allowbreak{}4\_\allowbreak{}b\\
Type: string\\
Default: nil
}

{\footnotesize Обмежено до <AC-07(04)\_\allowbreak{}ODP[02] кількості> послідовних невдалих спроб входу за допомогою використання альтернативних факторів протягом <AC-07(04)\_\allowbreak{}ODP[03] періоду часу>}




\subsection{ПОПЕРЕДЖЕННЯ ПРО ВИКОРИСТАННЯ СИСТЕМИ (AC-8)}
a.\\ 
b. Демонструвати користувачам [Призначення: визначене організацією сповіщення або банер про використання системи] перед тим, як надавати доступ до системи, що забезпечує безпеку та приватність відповідно до чинних законів, нормативних документів, наказів, директив, політик, правил, стандартів і керівних принципів, які зазначають, що:\\ 
1. користувачі здійснюють доступ до урядової системи;\\ 
2. використання системи може контролюватися, реєструватися та підлягати аудиту;\\ 
3. несанкціоноване використання системи забороняється та приводить до кримінальної та цивільної відповідальності;\\ 
4. використання системи означає згоду на моніторинг і запис дій користувача. Зберігати сповіщення або банер на екрані, доки користувачі не визнають умови використання та не приймуть явних дій для входу в систему або подальшого доступу до системи.\\ 
c. Для загальнодоступних систем:\\ 
1. демонструвати інформацію про умови використання системи [Призначення: визначені організацією умови], перш ніж надавати подальший доступ до загальнодоступної системи;\\ 
2. демонструвати посилання, якщо такі є, на моніторинг, запис або аудит, які узгоджуються з акомодацією приватності для таких систем, які зазвичай забороняють такі дії;\\ 
3. мати опис авторизованого використання системи.

\vspace{10pt}
{\footnotesize\bfseries
No: 1\\
Name: ac\_\allowbreak{}8\_\allowbreak{}odp\_\allowbreak{}01\\
Type: string\\
Default: nil
}

{\footnotesize Визначено сповіщення або банер про використання системи}


{\footnotesize\bfseries
No: 2\\
Name: ac\_\allowbreak{}8\_\allowbreak{}odp\_\allowbreak{}02\\
Type: string\\
Default: nil
}

{\footnotesize Визначені умови використання загальнодоступної системи}


{\footnotesize\bfseries
No: 3\\
Name: ac\_\allowbreak{}8\_\allowbreak{}01\\
Type: string\\
Default: nil
}

{\footnotesize Користувачам демонструється <AC-08\_\allowbreak{}ODP[01] сповіщення або банер> перед тим, як надавати доступ до системи}


{\footnotesize\bfseries
No: 4\\
Name: ac\_\allowbreak{}8\_\allowbreak{}02\\
Type: string\\
Default: nil
}

{\footnotesize Сповіщення зазначає: доступ до урядової системи, можливість контролю/аудиту, заборону несанкціонованого використання та згоду на моніторинг}


{\footnotesize\bfseries
No: 5\\
Name: ac\_\allowbreak{}8\_\allowbreak{}03\\
Type: string\\
Default: nil
}

{\footnotesize Сповіщення або банер залишається на екрані, доки користувачі не визнають умови використання та не приймуть явних дій для входу}


{\footnotesize\bfseries
No: 6\\
Name: ac\_\allowbreak{}8\_\allowbreak{}04\\
Type: string\\
Default: nil
}

{\footnotesize Для загальнодоступних систем демонструється інформація про <AC-08\_\allowbreak{}ODP[02] умови використання>, перш ніж надавати подальший доступ}


{\footnotesize\bfseries
No: 7\\
Name: ac\_\allowbreak{}8\_\allowbreak{}05\\
Type: string\\
Default: nil
}

{\footnotesize Для загальнодоступних систем демонструються посилання на моніторинг, запис або аудит}


{\footnotesize\bfseries
No: 8\\
Name: ac\_\allowbreak{}8\_\allowbreak{}06\\
Type: string\\
Default: nil
}

{\footnotesize Для загальнодоступних систем наведено опис авторизованого використання системи}




\subsection{СПОВІЩЕННЯ ПРО ПОПЕРЕДНІЙ ВХІД (ДОСТУП) (AC-9)}
Сповіщати користувача після успішного входу (доступу) до системи про дату та час останнього входу (доступу).

\vspace{10pt}
{\footnotesize\bfseries
No: 1\\
Name: ac\_\allowbreak{}9\_\allowbreak{}01\\
Type: integer\\
Default: 30
}

{\footnotesize Система повідомляє користувача при успішному вході (доступі) до системи про дату та час останнього входу (доступу)}




\subsubsection{НЕВДАЛІ СПРОБИ ВХОДУ ДО СИСТЕМИ (AC-9(1))}


\vspace{10pt}
{\footnotesize\bfseries
No: 1\\
Name: ac\_\allowbreak{}9\_\allowbreak{}1\_\allowbreak{}01\\
Type: integer\\
Default: 3
}

{\footnotesize Система сповіщає користувача після успішного входу / доступу про кількість невдалих спроб входу / доступу з моменту останнього успішного входу / доступу}




\subsubsection{УСПІШНІ ТА НЕВДАЛІ СПРОБИ ВХОДУ ДО СИСТЕМИ (AC-9(2))}
Сповіщати користувача, після успішного входу/доступу до системи про кількість [Вибір: успішних спроб доступу/входу; невдалих спроб входу/доступу; обидва варіанти] за [Призначення: визначений організацією період часу].

\vspace{10pt}
{\footnotesize\bfseries
No: 1\\
Name: ac\_\allowbreak{}9\_\allowbreak{}2\_\allowbreak{}01\\
Type: integer\\
Default: 30
}

{\footnotesize Після успішного входу в систему користувач отримує повідомлення про кількість ЗНАЧЕННЯ ВИБРАНОГО ПАРАМЕТРА протягом періоду часу}


{\footnotesize\bfseries
No: 2\\
Name: ac\_\allowbreak{}9\_\allowbreak{}2\_\allowbreak{}odp\_\allowbreak{}01\\
Type: integer\\
Default: 3
}

{\footnotesize Вибрано одне з наступних ЗНАЧЕНЬ ПАРАМЕТРІВ: \{успішних спроб доступу/входу; невдалих спроб входу/доступу; обидва варіанти\}}


{\footnotesize\bfseries
No: 3\\
Name: ac\_\allowbreak{}9\_\allowbreak{}2\_\allowbreak{}odp\_\allowbreak{}02\\
Type: integer\\
Default: 30
}

{\footnotesize Визначається період часу, протягом якого система повідомляє користувача про кількість успішних спроб входу в систему, невдалих спроб входу або про обидва випадки}




\subsubsection{ПОВІДОМЛЕННЯ ПРО ЗМІНИ В ОБЛІКОВОМУ ЗАПИСІ (AC-9(3))}
Сповіщати користувача, після успішного входу/доступу, про внесення змін до [Призначення: певних характеристик/параметрів облікового запису користувача, визначених організацією] протягом [Призначення: визначеного організацією періоду часу].

\vspace{10pt}
{\footnotesize\bfseries
No: 1\\
Name: ac\_\allowbreak{}9\_\allowbreak{}3\_\allowbreak{}odp\_\allowbreak{}01\\
Type: string\\
Default: nil
}

{\footnotesize Визначено характеристики або параметри облікового запису користувача, зміни яких потребують сповіщення}


{\footnotesize\bfseries
No: 2\\
Name: ac\_\allowbreak{}9\_\allowbreak{}3\_\allowbreak{}odp\_\allowbreak{}02\\
Type: string\\
Default: \textquotedbl{}15m\textquotedbl{}
}

{\footnotesize Визначено період часу, протягом якого вносились зміни}


{\footnotesize\bfseries
No: 3\\
Name: ac\_\allowbreak{}9\_\allowbreak{}3\_\allowbreak{}01\\
Type: string\\
Default: nil
}

{\footnotesize Користувач сповіщається після успішного входу про внесення змін до <AC-09(03)\_\allowbreak{}ODP[01] характеристик/параметрів облікового запису> протягом <AC-09(03)\_\allowbreak{}ODP[02] періоду часу>}




\subsubsection{СПОВІЩЕННЯ ПРО ПОПЕРЕДНІЙ ВХІД (ДОСТУП) -- ДОДАТКОВА ІНФОРМАЦІЯ ПРО ВХІД (AC-9(4))}
Повідомляти користувачеві, після успішного входу/доступу, наступну додаткову інформацію: [Призначення: інформація, визначена організацією, яка повинна бути включена на додаток до дати та часу останнього входу/доступу].

\vspace{10pt}
{\footnotesize\bfseries
No: 1\\
Name: ac\_\allowbreak{}9\_\allowbreak{}4\_\allowbreak{}01\\
Type: string\\
Default: nil
}

{\footnotesize Після успішного входу користувач отримує повідомлення додаткова інформація}


{\footnotesize\bfseries
No: 2\\
Name: ac\_\allowbreak{}9\_\allowbreak{}4\_\allowbreak{}odp\\
Type: string\\
Default: nil
}

{\footnotesize Визначено додаткову інформацію, про яку слід повідомити користувача}




\subsection{УПРАВЛІННЯ ПАРАЛЕЛЬНОЮ СЕСІЄЮ (AC-10)}
Обмежити кількість одночасних сеансів для кожного [Призначення: визначеного організацією облікового запису та/або типу облікового запису] до [Призначення: визначеної організацією кількості].

\vspace{10pt}
{\footnotesize\bfseries
No: 1\\
Name: ac\_\allowbreak{}10\_\allowbreak{}odp\_\allowbreak{}01\\
Type: string\\
Default: nil
}

{\footnotesize Визначено обліковий запис та/або тип облікового запису, для якого обмежуються одночасні сеанси}


{\footnotesize\bfseries
No: 2\\
Name: ac\_\allowbreak{}10\_\allowbreak{}odp\_\allowbreak{}02\\
Type: integer\\
Default: 3
}

{\footnotesize Визначено кількість одночасних сеансів, дозволених для кожного облікового запису та/або типу облікового запису}


{\footnotesize\bfseries
No: 3\\
Name: ac\_\allowbreak{}10\_\allowbreak{}01\\
Type: string\\
Default: nil
}

{\footnotesize Кількість одночасних сеансів для кожного <AC-10\_\allowbreak{}ODP[01] облікового запису та/або типу> обмежена до <AC-10\_\allowbreak{}ODP[02] кількості>}




\subsection{БЛОКУВАННЯ ПРИСТРОЮ (AC-11)}
a. Заборонити подальший доступ до системи шляхом ініціювання блокування пристрою після [Призначення: визначеного організацією періоду] бездіяльності або після отримання запиту від користувача.\\ 
b. Зберігати блокування пристрою, поки користувач не відновить доступ, використовуючи встановлені процедури ідентифікації та автентифікації.

\vspace{10pt}
{\footnotesize\bfseries
No: 1\\
Name: ac\_\allowbreak{}11\_\allowbreak{}odp\_\allowbreak{}01\\
Type: string\\
Default: nil
}

{\footnotesize Вибрано одне або декілька з наступних значень: \{ініціювання блокування пристрою після періоду неактивності; вимога до користувача ініціювати блокування пристрою перед тим, як залишити систему без нагляду\}}


{\footnotesize\bfseries
No: 2\\
Name: ac\_\allowbreak{}11\_\allowbreak{}odp\_\allowbreak{}02\\
Type: string\\
Default: \textquotedbl{}15m\textquotedbl{}
}

{\footnotesize Визначено часовий проміжок бездіяльності, після якого ініціюється блокування пристрою}


{\footnotesize\bfseries
No: 3\\
Name: ac\_\allowbreak{}11\_\allowbreak{}a\\
Type: string\\
Default: nil
}

{\footnotesize Заборонено подальший доступ до системи шляхом <AC-11\_\allowbreak{}ODP[01] вибору умов блокування>, зокрема після <AC-11\_\allowbreak{}ODP[02] періоду> бездіяльності}


{\footnotesize\bfseries
No: 4\\
Name: ac\_\allowbreak{}11\_\allowbreak{}b\\
Type: string\\
Default: nil
}

{\footnotesize Блокування пристрою зберігається, поки користувач не відновить доступ, використовуючи встановлені процедури ідентифікації та автентифікації}




\subsubsection{ПРИХОВАНІ ДИСПЛЕЇ (AC-11(1))}


\vspace{10pt}
{\footnotesize\bfseries
No: 1\\
Name: ac\_\allowbreak{}11\_\allowbreak{}1\_\allowbreak{}01\\
Type: string\\
Default: nil
}

{\footnotesize Система приховує (через блокування сеансу) інформацію, попередньо видиму на дисплеї, загальнодоступним зображенням}




\subsection{ПРИПИНЕННЯ СЕАНСУ (AC-12)}
Сеанс користувача має завершуватися автоматично після [Призначення: визначених організацією умов або тригерних подій, що вимагають припинення сеансу].

\vspace{10pt}
{\footnotesize\bfseries
No: 1\\
Name: ac\_\allowbreak{}12\_\allowbreak{}01\\
Type: string\\
Default: nil
}

{\footnotesize Сеанс користувача автоматично завершується після виконання умов або подій}


{\footnotesize\bfseries
No: 2\\
Name: ac\_\allowbreak{}12\_\allowbreak{}odp\\
Type: list\\
Default: [\textquotedbl{}login\textquotedbl{}, \textquotedbl{}logout\textquotedbl{}, \textquotedbl{}failed\_\allowbreak{}attempt\textquotedbl{}]
}

{\footnotesize Визначено умови або події, що вимагають припинення сеансу}




\subsubsection{ІНІЦІЙОВАНЕ КОРИСТУВАЧЕМ БЛОКУВАННЯ (AC-12(1))}
Забезпечити можливість припинення сеансів зв'язку з ініціативи користувача, коли автентифікація використовується для отримання доступу до [Призначення: визначених організацією інформаційних ресурсів].

\vspace{10pt}
{\footnotesize\bfseries
No: 1\\
Name: ac\_\allowbreak{}12\_\allowbreak{}1\_\allowbreak{}odp\_\allowbreak{}01\\
Type: string\\
Default: nil
}

{\footnotesize Визначено інформаційні ресурси, для доступу до яких використовується автентифікація}


{\footnotesize\bfseries
No: 2\\
Name: ac\_\allowbreak{}12\_\allowbreak{}1\_\allowbreak{}01\\
Type: string\\
Default: nil
}

{\footnotesize Забезпечується можливість припинення сеансів зв'язку з ініціативи користувача, коли автентифікація використовується для отримання доступу до <AC-12(01)\_\allowbreak{}ODP[01] інформаційних ресурсів>}




\subsubsection{ПОВІДОМЛЕННЯ ПРО ПРИПИНЕННЯ СЕАНСУ (AC-12(2))}
Відобразити виразне повідомлення для користувача, що вказує на достовірне припинення автентифікованих сеансів зв'язку.

\vspace{10pt}
{\footnotesize\bfseries
No: 1\\
Name: ac\_\allowbreak{}12\_\allowbreak{}2\_\allowbreak{}01\\
Type: string\\
Default: nil
}

{\footnotesize Користувачам буде показано явне повідомлення про завершення сеансу автентифікованого зв'язку}




\subsubsection{ЗАСТЕРЕЖНЕ ПОВІДОМЛЕННЯ ПРО ТЕ, ЩО ЧАС СЕСІЇ ДОБІГАЄ КІНЦЯ (AC-12(3))}
Відобразити виразне повідомлення користувачам, що вказує, що сесія добігає кінця [Завдання: визначений організацією час до кінця сесії].

\vspace{10pt}
{\footnotesize\bfseries
No: 1\\
Name: ac\_\allowbreak{}12\_\allowbreak{}3\_\allowbreak{}odp\_\allowbreak{}01\\
Type: string\\
Default: \textquotedbl{}5m\textquotedbl{}
}

{\footnotesize Визначений час до кінця сесії}


{\footnotesize\bfseries
No: 2\\
Name: ac\_\allowbreak{}12\_\allowbreak{}3\_\allowbreak{}01\\
Type: string\\
Default: nil
}

{\footnotesize Користувачам відображається виразне повідомлення, що вказує, що сесія добігає кінця за <AC-12(03)\_\allowbreak{}ODP[01] визначений час>}




\subsection{НАГЛЯД ТА ОГЛЯД - УПРАВЛІННЯ ДОСТУПОМ (AC-13) [Вилучено]}
[Вилучено: включено в АС-02 та АU-06]

\vspace{10pt}
\textit{Немає параметрів для цього контролю.}
\vspace{10pt}


\subsection{ДОЗВОЛЕНІ ДІЇ БЕЗ ІДЕНТИФІКАЦІЇ АБО АВТЕНТИФІКАЦІЇ (AC-14)}
a. Визначити [Призначення: дозволені організацією дії користувачів], які можуть виконуватися в системі без ідентифікації або автентифікації відповідно до завдань та функцій організації.\\ 
b. Документувати та визначити відповідне обґрунтування в плані безпеки системи дій користувача, які не потребують ідентифікації або автентифікації.

\vspace{10pt}
{\footnotesize\bfseries
No: 1\\
Name: ac\_\allowbreak{}14\_\allowbreak{}a\\
Type: list\\
Default: [\textquotedbl{}login\textquotedbl{}, \textquotedbl{}logout\textquotedbl{}, \textquotedbl{}failed\_\allowbreak{}attempt\textquotedbl{}]
}

{\footnotesize Визначено дії користувача, які можуть бути виконані в системі без ідентифікації або автентифікації, що відповідають місії та функціям організації}


{\footnotesize\bfseries
No: 2\\
Name: ac\_\allowbreak{}14\_\allowbreak{}a\_\allowbreak{}02\\
Type: string\\
Default: nil
}

{\footnotesize Обґрунтування дій користувачів, які не потребують ідентифікації або автентифікації, надається в плані захисту інформації}


{\footnotesize\bfseries
No: 3\\
Name: ac\_\allowbreak{}14\_\allowbreak{}b\_\allowbreak{}01\\
Type: list\\
Default: [\textquotedbl{}login\textquotedbl{}, \textquotedbl{}logout\textquotedbl{}, \textquotedbl{}failed\_\allowbreak{}attempt\textquotedbl{}]
}

{\footnotesize Дії користувачів, які не потребують ідентифікації або автентифікації, задокументовані в плані захисту інформації}


{\footnotesize\bfseries
No: 4\\
Name: ac\_\allowbreak{}14\_\allowbreak{}odp\\
Type: list\\
Default: [\textquotedbl{}login\textquotedbl{}, \textquotedbl{}logout\textquotedbl{}, \textquotedbl{}failed\_\allowbreak{}attempt\textquotedbl{}]
}

{\footnotesize Визначено дії користувача, які можуть бути виконані в системі без ідентифікації або автентифікації}




\subsubsection{ДОЗВОЛЕНІ ДІЇ БЕЗ ІДЕНТИФІКАЦІЇ НЕОБХІДНЕ ВИКОРИСТАННЯ (AC-14(1)) [Вилучено]}
[Вилучено: включено до AC-14]

\vspace{10pt}
\textit{Немає параметрів для цього контролю.}
\vspace{10pt}


\subsection{АВТОМАТИЗОВАНЕ МАРКУВАННЯ (AC-15) [Вилучено]}
[Вилучено: включено до MP-03]

\vspace{10pt}
\textit{Немає параметрів для цього контролю.}
\vspace{10pt}


\subsection{АТРИБУТИ БЕЗПЕКИ ТА ПРИВАТНОСТІ (AC-16)}
a. Визначити засоби для асоціювання (пов'язання) [Призначення: визначених організацією типів атрибутів безпеки та приватності], що приймають [Призначення: визначені організацією значення атрибутів безпеки та приватності] з інформацією, яка зберігається, обробляється та/або передається.\\ 
b. Пов'язані атрибути безпеки та приватності мають створюватися і зберігатися разом з інформацією.\\ 
c. Встановити дозволені [Призначення: визначені організацією атрибути безпеки та приватності] для [Призначення: систем, визначених організацією].\\ 
d. Визначити дозволені [Призначення: визначені організацією значення або діапазони] для кожного з встановлених атрибутів безпеки та приватності.\\ 
e. Проводити аудит змін атрибутів.\\ 
f. Переглядати атрибути безпеки та приватності на відповідність з [Призначення: визначеною організацією частотою].

\vspace{10pt}
{\footnotesize\bfseries
No: 1\\
Name: ac\_\allowbreak{}16\_\allowbreak{}odp\_\allowbreak{}01\\
Type: string\\
Default: nil
}

{\footnotesize Визначено типи атрибутів безпеки}


{\footnotesize\bfseries
No: 2\\
Name: ac\_\allowbreak{}16\_\allowbreak{}odp\_\allowbreak{}02\\
Type: string\\
Default: nil
}

{\footnotesize Визначено типи атрибутів приватності (конфіденційності)}


{\footnotesize\bfseries
No: 3\\
Name: ac\_\allowbreak{}16\_\allowbreak{}odp\_\allowbreak{}03\\
Type: string\\
Default: nil
}

{\footnotesize Визначено значення атрибутів безпеки}


{\footnotesize\bfseries
No: 4\\
Name: ac\_\allowbreak{}16\_\allowbreak{}odp\_\allowbreak{}04\\
Type: string\\
Default: nil
}

{\footnotesize Визначено значення атрибутів приватності (конфіденційності)}


{\footnotesize\bfseries
No: 5\\
Name: ac\_\allowbreak{}16\_\allowbreak{}odp\_\allowbreak{}05\\
Type: string\\
Default: nil
}

{\footnotesize Визначено системи, для яких мають бути встановлені дозволені атрибути безпеки}


{\footnotesize\bfseries
No: 6\\
Name: ac\_\allowbreak{}16\_\allowbreak{}odp\_\allowbreak{}06\\
Type: string\\
Default: nil
}

{\footnotesize Визначено системи, для яких мають бути встановлені дозволені атрибути приватності}


{\footnotesize\bfseries
No: 7\\
Name: ac\_\allowbreak{}16\_\allowbreak{}odp\_\allowbreak{}07\\
Type: string\\
Default: nil
}

{\footnotesize Визначено атрибути безпеки, які дозволені для систем}


{\footnotesize\bfseries
No: 8\\
Name: ac\_\allowbreak{}16\_\allowbreak{}odp\_\allowbreak{}08\\
Type: string\\
Default: nil
}

{\footnotesize Визначено атрибути приватності, які дозволені для систем}


{\footnotesize\bfseries
No: 9\\
Name: ac\_\allowbreak{}16\_\allowbreak{}odp\_\allowbreak{}09\\
Type: string\\
Default: nil
}

{\footnotesize Визначено значення атрибутів або діапазони для встановлених атрибутів}


{\footnotesize\bfseries
No: 10\\
Name: ac\_\allowbreak{}16\_\allowbreak{}odp\_\allowbreak{}10\\
Type: string\\
Default: \textquotedbl{}1 year\textquotedbl{}
}

{\footnotesize Визначено частоту, з якою слід переглядати атрибути безпеки на предмет відповідності}


{\footnotesize\bfseries
No: 11\\
Name: ac\_\allowbreak{}16\_\allowbreak{}odp\_\allowbreak{}11\\
Type: string\\
Default: \textquotedbl{}1 year\textquotedbl{}
}

{\footnotesize Визначено частоту, з якою слід переглядати атрибути приватності на предмет відповідності}


{\footnotesize\bfseries
No: 12\\
Name: ac\_\allowbreak{}16\_\allowbreak{}a\_\allowbreak{}01\\
Type: string\\
Default: nil
}

{\footnotesize Визначено засоби для асоціювання <AC-16\_\allowbreak{}ODP[01] типів атрибутів безпеки> з їх <AC-16\_\allowbreak{}ODP[03] значеннями> для інформації}


{\footnotesize\bfseries
No: 13\\
Name: ac\_\allowbreak{}16\_\allowbreak{}a\_\allowbreak{}02\\
Type: string\\
Default: nil
}

{\footnotesize Визначено засоби для асоціювання <AC-16\_\allowbreak{}ODP[02] типів атрибутів приватності> з їх <AC-16\_\allowbreak{}ODP[04] значеннями> для інформації}


{\footnotesize\bfseries
No: 14\\
Name: ac\_\allowbreak{}16\_\allowbreak{}b\_\allowbreak{}01\\
Type: string\\
Default: nil
}

{\footnotesize Пов'язані атрибути безпеки створюються і зберігаються разом з інформацією}


{\footnotesize\bfseries
No: 15\\
Name: ac\_\allowbreak{}16\_\allowbreak{}b\_\allowbreak{}02\\
Type: string\\
Default: nil
}

{\footnotesize Пов'язані атрибути приватності створюються і зберігаються разом з інформацією}


{\footnotesize\bfseries
No: 16\\
Name: ac\_\allowbreak{}16\_\allowbreak{}c\_\allowbreak{}01\\
Type: string\\
Default: nil
}

{\footnotesize Встановлюються дозволені <AC-16\_\allowbreak{}ODP[07] атрибути безпеки> для <AC-16\_\allowbreak{}ODP[05] систем>}


{\footnotesize\bfseries
No: 17\\
Name: ac\_\allowbreak{}16\_\allowbreak{}c\_\allowbreak{}02\\
Type: string\\
Default: nil
}

{\footnotesize Встановлюються дозволені <AC-16\_\allowbreak{}ODP[08] атрибути приватності> для <AC-16\_\allowbreak{}ODP[06] систем>}


{\footnotesize\bfseries
No: 18\\
Name: ac\_\allowbreak{}16\_\allowbreak{}d\_\allowbreak{}01\\
Type: string\\
Default: nil
}

{\footnotesize Визначено дозволені <AC-16\_\allowbreak{}ODP[09] значення або діапазони> для кожного з встановлених атрибутів безпеки та приватності}


{\footnotesize\bfseries
No: 19\\
Name: ac\_\allowbreak{}16\_\allowbreak{}e\_\allowbreak{}01\\
Type: string\\
Default: nil
}

{\footnotesize Проводиться аудит змін до атрибутів}


{\footnotesize\bfseries
No: 20\\
Name: ac\_\allowbreak{}16\_\allowbreak{}f\_\allowbreak{}01\\
Type: string\\
Default: nil
}

{\footnotesize Атрибути безпеки перевіряються на відповідність із <AC-16\_\allowbreak{}ODP[10] частотою>}


{\footnotesize\bfseries
No: 21\\
Name: ac\_\allowbreak{}16\_\allowbreak{}f\_\allowbreak{}02\\
Type: string\\
Default: nil
}

{\footnotesize Атрибути приватності перевіряються на відповідність із <AC-16\_\allowbreak{}ODP[11] частотою>}




\subsubsection{ДИНАМІЧНЕ ПОВ'ЯЗАННЯ АТРИБУТІВ (AC-16(1))}


\vspace{10pt}
{\footnotesize\bfseries
No: 1\\
Name: ac\_\allowbreak{}16\_\allowbreak{}1\_\allowbreak{}01\\
Type: list\\
Default: [\textquotedbl{}default\_\allowbreak{}deny\_\allowbreak{}rule\textquotedbl{}, \textquotedbl{}abac\_\allowbreak{}rule\_\allowbreak{}1\textquotedbl{}]
}

{\footnotesize Атрибути безпеки динамічно пов'язуються з суб'єктами відповідно до наведених нижче політик безпеки під час створення та комбінування інформації: політики безпеки}


{\footnotesize\bfseries
No: 2\\
Name: ac\_\allowbreak{}16\_\allowbreak{}1\_\allowbreak{}02\\
Type: list\\
Default: [\textquotedbl{}default\_\allowbreak{}deny\_\allowbreak{}rule\textquotedbl{}, \textquotedbl{}abac\_\allowbreak{}rule\_\allowbreak{}1\textquotedbl{}]
}

{\footnotesize Атрибути безпеки динамічно пов'язуються з об'єктами відповідно до наведених нижче політик безпеки під час створення та комбінування інформації: політики безпеки}


{\footnotesize\bfseries
No: 3\\
Name: ac\_\allowbreak{}16\_\allowbreak{}1\_\allowbreak{}03\\
Type: list\\
Default: [\textquotedbl{}default\_\allowbreak{}deny\_\allowbreak{}rule\textquotedbl{}, \textquotedbl{}abac\_\allowbreak{}rule\_\allowbreak{}1\textquotedbl{}]
}

{\footnotesize Атрибути конфіденційності динамічно пов'язуються з суб'єктами відповідно до наведених нижче політик конфіденційності під час створення та комбінування інформації: політики конфіденційності}


{\footnotesize\bfseries
No: 4\\
Name: ac\_\allowbreak{}16\_\allowbreak{}1\_\allowbreak{}04\\
Type: list\\
Default: [\textquotedbl{}default\_\allowbreak{}deny\_\allowbreak{}rule\textquotedbl{}, \textquotedbl{}abac\_\allowbreak{}rule\_\allowbreak{}1\textquotedbl{}]
}

{\footnotesize Атрибути конфіденційності динамічно пов'язуються з об'єктами відповідно до наведених нижче політик конфіденційності під час створення та комбінування інформації: політики конфіденційності}


{\footnotesize\bfseries
No: 5\\
Name: ac\_\allowbreak{}16\_\allowbreak{}1\_\allowbreak{}odp\_\allowbreak{}01\\
Type: list\\
Default: [\textquotedbl{}default\_\allowbreak{}deny\_\allowbreak{}rule\textquotedbl{}, \textquotedbl{}abac\_\allowbreak{}rule\_\allowbreak{}1\textquotedbl{}]
}

{\footnotesize Визначено суб'єкти, з якими атрибути безпеки повинні динамічно пов'язуватися при створенні та комбінуванні інформації}


{\footnotesize\bfseries
No: 6\\
Name: ac\_\allowbreak{}16\_\allowbreak{}1\_\allowbreak{}odp\_\allowbreak{}02\\
Type: list\\
Default: [\textquotedbl{}default\_\allowbreak{}deny\_\allowbreak{}rule\textquotedbl{}, \textquotedbl{}abac\_\allowbreak{}rule\_\allowbreak{}1\textquotedbl{}]
}

{\footnotesize Визначено об'єкти, з якими атрибути безпеки повинні динамічно пов'язуватися при створенні та комбінуванні інформації}


{\footnotesize\bfseries
No: 7\\
Name: ac\_\allowbreak{}16\_\allowbreak{}1\_\allowbreak{}odp\_\allowbreak{}03\\
Type: list\\
Default: [\textquotedbl{}default\_\allowbreak{}deny\_\allowbreak{}rule\textquotedbl{}, \textquotedbl{}abac\_\allowbreak{}rule\_\allowbreak{}1\textquotedbl{}]
}

{\footnotesize Визначені суб'єкти, з якими атрибути конфіденційності повинні динамічно пов'язуватися при створенні та комбінуванні інформації}


{\footnotesize\bfseries
No: 8\\
Name: ac\_\allowbreak{}16\_\allowbreak{}1\_\allowbreak{}odp\_\allowbreak{}04\\
Type: list\\
Default: [\textquotedbl{}default\_\allowbreak{}deny\_\allowbreak{}rule\textquotedbl{}, \textquotedbl{}abac\_\allowbreak{}rule\_\allowbreak{}1\textquotedbl{}]
}

{\footnotesize Визначені об'єкти, з якими атрибути конфіденційності повинні динамічно пов'язуватися при створенні та комбінуванні інформації}


{\footnotesize\bfseries
No: 9\\
Name: ac\_\allowbreak{}16\_\allowbreak{}1\_\allowbreak{}odp\_\allowbreak{}05\\
Type: list\\
Default: [\textquotedbl{}default\_\allowbreak{}deny\_\allowbreak{}rule\textquotedbl{}, \textquotedbl{}abac\_\allowbreak{}rule\_\allowbreak{}1\textquotedbl{}]
}

{\footnotesize Визначено політики безпеки, що вимагають динамічного пов'язування атрибутів безпеки з суб'єктами та об'єктами}


{\footnotesize\bfseries
No: 10\\
Name: ac\_\allowbreak{}16\_\allowbreak{}1\_\allowbreak{}odp\_\allowbreak{}06\\
Type: list\\
Default: [\textquotedbl{}default\_\allowbreak{}deny\_\allowbreak{}rule\textquotedbl{}, \textquotedbl{}abac\_\allowbreak{}rule\_\allowbreak{}1\textquotedbl{}]
}

{\footnotesize Визначено політики конфіденційності, що вимагають динамічного пов'язування атрибутів безпеки з суб'єктами та об'єктами; AC-16(01)[01] атрибути безпеки динамічно пов'язуються з суб'єктами відповідно до наведених нижче політик безпеки під час створення та комбінування інформації: політики безпеки}




\subsubsection{ЗМІНА ЗНАЧЕНЬ АТРИБУТІВ АВТОРИЗОВАНИМИ ОСОБАМИ (AC-16(2))}
Надати уповноваженим особам (або процесам, що діють від імені фізичних осіб) можливість визначати або змінювати значення відповідних атрибутів безпеки та приватності.

\vspace{10pt}
{\footnotesize\bfseries
No: 1\\
Name: ac\_\allowbreak{}16\_\allowbreak{}2\_\allowbreak{}01\\
Type: string\\
Default: nil
}

{\footnotesize Уповноважені особи (або процеси, що діють від імені осіб) мають можливість визначати або змінювати значення пов'язаних з ними атрибутів безпеки}


{\footnotesize\bfseries
No: 2\\
Name: ac\_\allowbreak{}16\_\allowbreak{}2\_\allowbreak{}02\\
Type: string\\
Default: nil
}

{\footnotesize Уповноважені особи (або процеси, що діють від імені осіб) мають можливість визначати або змінювати значення пов'язаних з ними атрибутів конфіденційності}




\subsubsection{ПІДТРИМКА СИСТЕМОЮ ПОВ'ЯЗАННЯ АТРИБУТІВ (AC-16(3))}
ПОВ'ЯЗАННЯ АТРИБУТІВ Підтримати пов'язання та цілісність [Призначення: визначених організацією атрибутів безпеки та приватності] з [Призначення: визначених організацією суб'єктів і об'єктів].

\vspace{10pt}
{\footnotesize\bfseries
No: 1\\
Name: ac\_\allowbreak{}16\_\allowbreak{}3\_\allowbreak{}01\\
Type: string\\
Default: nil
}

{\footnotesize Підтримується зв'язок та цілісність атрибутів безпеки з суб'єктами}


{\footnotesize\bfseries
No: 2\\
Name: ac\_\allowbreak{}16\_\allowbreak{}3\_\allowbreak{}02\\
Type: string\\
Default: nil
}

{\footnotesize Підтримується зв'язок та цілісність атрибутів безпеки з об'єктами}


{\footnotesize\bfseries
No: 3\\
Name: ac\_\allowbreak{}16\_\allowbreak{}3\_\allowbreak{}03\\
Type: string\\
Default: nil
}

{\footnotesize Підтримується зв'язок та цілісність атрибутів конфіденційності з суб'єктами}


{\footnotesize\bfseries
No: 4\\
Name: ac\_\allowbreak{}16\_\allowbreak{}3\_\allowbreak{}04\\
Type: string\\
Default: nil
}

{\footnotesize Підтримується зв'язок та цілісність атрибутів конфіденційності з об'єктами; ПОТЕНЦІЙНІ МЕТОДИ ТА ОБ'ЄКТИ ОЦІНКИ:}


{\footnotesize\bfseries
No: 5\\
Name: ac\_\allowbreak{}16\_\allowbreak{}3\_\allowbreak{}odp\_\allowbreak{}01\\
Type: list\\
Default: [\textquotedbl{}default\_\allowbreak{}deny\_\allowbreak{}rule\textquotedbl{}, \textquotedbl{}abac\_\allowbreak{}rule\_\allowbreak{}1\textquotedbl{}]
}

{\footnotesize Визначено атрибути безпеки, які потребують підтримки асоціацій та цілісності}


{\footnotesize\bfseries
No: 6\\
Name: ac\_\allowbreak{}16\_\allowbreak{}3\_\allowbreak{}odp\_\allowbreak{}02\\
Type: list\\
Default: [\textquotedbl{}default\_\allowbreak{}deny\_\allowbreak{}rule\textquotedbl{}, \textquotedbl{}abac\_\allowbreak{}rule\_\allowbreak{}1\textquotedbl{}]
}

{\footnotesize Визначено атрибути конфіденційності, які потребують підтримки асоціації та цілісності}


{\footnotesize\bfseries
No: 7\\
Name: ac\_\allowbreak{}16\_\allowbreak{}3\_\allowbreak{}odp\_\allowbreak{}03\\
Type: string\\
Default: nil
}

{\footnotesize Визначено суб'єкти, які потребують об'єднання та збереження цілісності атрибутів безпеки, що належать до таких суб'єктів}


{\footnotesize\bfseries
No: 8\\
Name: ac\_\allowbreak{}16\_\allowbreak{}3\_\allowbreak{}odp\_\allowbreak{}04\\
Type: string\\
Default: nil
}

{\footnotesize Визначено об'єкти, які потребують об'єднання та збереження цілісності атрибутів безпеки, що належать до таких об'єктів}


{\footnotesize\bfseries
No: 9\\
Name: ac\_\allowbreak{}16\_\allowbreak{}3\_\allowbreak{}odp\_\allowbreak{}05\\
Type: string\\
Default: nil
}

{\footnotesize Визначено суб'єктів, які потребують об'єднання та збереження цілісності атрибутів конфіденційності щодо таких суб'єктів}


{\footnotesize\bfseries
No: 10\\
Name: ac\_\allowbreak{}16\_\allowbreak{}3\_\allowbreak{}odp\_\allowbreak{}06\\
Type: string\\
Default: nil
}

{\footnotesize Визначено об'єктів, які потребують об'єднання та збереження цілісності атрибутів конфіденційності щодо таких об'єктів}




\subsubsection{ПОВ'ЯЗАННЯ АТРИБУТІВ АВТОРИЗОВАНИМИ ОСОБАМИ (AC-16(4))}
Впровадити можливість пов'язувати [Призначення: визначені організацією атрибути безпеки та приватності] з [Призначення: визначеними організацією суб'єктами та об'єктами] уповноваженими особами (або процесами, що діють від імені фізичних осіб).

\vspace{10pt}
{\footnotesize\bfseries
No: 1\\
Name: ac\_\allowbreak{}16\_\allowbreak{}4\_\allowbreak{}odp\_\allowbreak{}01\\
Type: string\\
Default: nil
}

{\footnotesize Визначено атрибути безпеки, які пов'язуються з суб'єктами}


{\footnotesize\bfseries
No: 2\\
Name: ac\_\allowbreak{}16\_\allowbreak{}4\_\allowbreak{}odp\_\allowbreak{}02\\
Type: string\\
Default: nil
}

{\footnotesize Визначено атрибути безпеки, які пов'язуються з об'єктами}


{\footnotesize\bfseries
No: 3\\
Name: ac\_\allowbreak{}16\_\allowbreak{}4\_\allowbreak{}odp\_\allowbreak{}03\\
Type: string\\
Default: nil
}

{\footnotesize Визначено атрибути приватності, які пов'язуються з суб'єктами}


{\footnotesize\bfseries
No: 4\\
Name: ac\_\allowbreak{}16\_\allowbreak{}4\_\allowbreak{}odp\_\allowbreak{}04\\
Type: string\\
Default: nil
}

{\footnotesize Визначено атрибути приватності, які пов'язуються з об'єктами}


{\footnotesize\bfseries
No: 5\\
Name: ac\_\allowbreak{}16\_\allowbreak{}4\_\allowbreak{}odp\_\allowbreak{}05\\
Type: string\\
Default: nil
}

{\footnotesize Визначено суб'єкти, з якими пов'язуються атрибути безпеки}


{\footnotesize\bfseries
No: 6\\
Name: ac\_\allowbreak{}16\_\allowbreak{}4\_\allowbreak{}odp\_\allowbreak{}06\\
Type: string\\
Default: nil
}

{\footnotesize Визначено об'єкти, з якими пов'язуються атрибути безпеки}


{\footnotesize\bfseries
No: 7\\
Name: ac\_\allowbreak{}16\_\allowbreak{}4\_\allowbreak{}odp\_\allowbreak{}07\\
Type: string\\
Default: nil
}

{\footnotesize Визначено суб'єкти, з якими пов'язуються атрибути приватності}


{\footnotesize\bfseries
No: 8\\
Name: ac\_\allowbreak{}16\_\allowbreak{}4\_\allowbreak{}odp\_\allowbreak{}08\\
Type: string\\
Default: nil
}

{\footnotesize Визначено об'єкти, з якими пов'язуються атрибути приватності}


{\footnotesize\bfseries
No: 9\\
Name: ac\_\allowbreak{}16\_\allowbreak{}4\_\allowbreak{}01\\
Type: string\\
Default: nil
}

{\footnotesize Впроваджено можливість уповноваженим особам (або процесам) пов'язувати <AC-16(04)\_\allowbreak{}ODP[01] атрибути безпеки> з <AC-16(04)\_\allowbreak{}ODP[05] суб'єктами>}


{\footnotesize\bfseries
No: 10\\
Name: ac\_\allowbreak{}16\_\allowbreak{}4\_\allowbreak{}02\\
Type: string\\
Default: nil
}

{\footnotesize Впроваджено можливість уповноваженим особам (або процесам) пов'язувати <AC-16(04)\_\allowbreak{}ODP[02] атрибути безпеки> з <AC-16(04)\_\allowbreak{}ODP[06] об'єктами>}


{\footnotesize\bfseries
No: 11\\
Name: ac\_\allowbreak{}16\_\allowbreak{}4\_\allowbreak{}03\\
Type: string\\
Default: nil
}

{\footnotesize Впроваджено можливість уповноваженим особам (або процесам) пов'язувати <AC-16(04)\_\allowbreak{}ODP[03] атрибути приватності> з <AC-16(04)\_\allowbreak{}ODP[07] суб'єктами>}


{\footnotesize\bfseries
No: 12\\
Name: ac\_\allowbreak{}16\_\allowbreak{}4\_\allowbreak{}04\\
Type: string\\
Default: nil
}

{\footnotesize Впроваджено можливість уповноваженим особам (або процесам) пов'язувати <AC-16(04)\_\allowbreak{}ODP[04] атрибути приватності> з <AC-16(04)\_\allowbreak{}ODP[08] об'єктами>}




\subsubsection{ВІДОБРАЖЕННЯ АТРИБУТІВ НА ПРИСТРОЯХ ВИВЕДЕННЯ (AC-16(5))}
Відображати атрибути безпеки та приватності в зручній для людини формі для кожного об'єкту, який система передає на пристрої виведення, щоб ідентифікувати [Призначення: визначені організацією спеціальні інструкції щодо поширення, обробки чи наступного розподілу інформації], використовуючи [Призначення: визначену організацією ідентифікацію, у зручній для людини формі про стандартні угоди про присвоєння імен].

\vspace{10pt}
{\footnotesize\bfseries
No: 1\\
Name: ac\_\allowbreak{}16\_\allowbreak{}5\_\allowbreak{}odp\_\allowbreak{}01\\
Type: string\\
Default: nil
}

{\footnotesize Визначено спеціальні інструкції щодо поширення, обробки чи наступного розподілу інформації}


{\footnotesize\bfseries
No: 2\\
Name: ac\_\allowbreak{}16\_\allowbreak{}5\_\allowbreak{}odp\_\allowbreak{}02\\
Type: string\\
Default: nil
}

{\footnotesize Визначено стандартні угоди про ідентифікацію (присвоєння імен) у зручній для людини формі}


{\footnotesize\bfseries
No: 3\\
Name: ac\_\allowbreak{}16\_\allowbreak{}5\_\allowbreak{}01\\
Type: string\\
Default: nil
}

{\footnotesize Атрибути безпеки відображаються у зручній для людини формі для кожного об'єкту, який система передає на пристрої виведення, щоб ідентифікувати <AC-16(05)\_\allowbreak{}ODP[01] спеціальні інструкції>, використовуючи <AC-16(05)\_\allowbreak{}ODP[02] стандартні угоди>}


{\footnotesize\bfseries
No: 4\\
Name: ac\_\allowbreak{}16\_\allowbreak{}5\_\allowbreak{}02\\
Type: string\\
Default: nil
}

{\footnotesize Атрибути приватності відображаються у зручній для людини формі для кожного об'єкту, який система передає на пристрої виведення, щоб ідентифікувати <AC-16(05)\_\allowbreak{}ODP[01] спеціальні інструкції>, використовуючи <AC-16(05)\_\allowbreak{}ODP[02] стандартні угоди>}




\subsubsection{ПІДТРИМКА ПОВ'ЯЗАННЯ АТРИБУТІВ ОРГАНІЗАЦІЄЮ (AC-16(6))}
Вимагати від персоналу пов'язувати та підтримувати асоціацію [Призначення: визначених організацією атрибутів безпеки та приватності] з [Призначенням: визначеними організацією суб'єктами та об'єктами] відповідно до [Призначення: визначеної організацією політики безпеки та приватності].

\vspace{10pt}
{\footnotesize\bfseries
No: 1\\
Name: ac\_\allowbreak{}16\_\allowbreak{}6\_\allowbreak{}01\\
Type: list\\
Default: [\textquotedbl{}admin\textquotedbl{}, \textquotedbl{}security\_\allowbreak{}officer\textquotedbl{}]
}

{\footnotesize Персонал зобов'язаний пов'язувати та підтримувати зв'язок <AC-16(06)\_\allowbreak{}ODP[01] атрибутів безпеки> з <AC-16(06)\_\allowbreak{}ODP[05] суб'єктами> відповідно до <AC-16(06)\_\allowbreak{}ODP[09] політик безпеки>}


{\footnotesize\bfseries
No: 2\\
Name: ac\_\allowbreak{}16\_\allowbreak{}6\_\allowbreak{}02\\
Type: list\\
Default: [\textquotedbl{}admin\textquotedbl{}, \textquotedbl{}security\_\allowbreak{}officer\textquotedbl{}]
}

{\footnotesize Персонал зобов'язаний пов'язувати та підтримувати зв'язок <AC-16(06)\_\allowbreak{}ODP[02] атрибутів безпеки> з <AC-16(06)\_\allowbreak{}ODP[06] об'єктами> відповідно до <AC-16(06)\_\allowbreak{}ODP[09] політик безпеки>}


{\footnotesize\bfseries
No: 3\\
Name: ac\_\allowbreak{}16\_\allowbreak{}6\_\allowbreak{}03\\
Type: list\\
Default: [\textquotedbl{}admin\textquotedbl{}, \textquotedbl{}security\_\allowbreak{}officer\textquotedbl{}]
}

{\footnotesize Персонал зобов'язаний пов'язувати та підтримувати зв'язок <AC-16(06)\_\allowbreak{}ODP[03] атрибутів конфіденційності> з <AC-16(06)\_\allowbreak{}ODP[07] суб'єктами> відповідно до <AC-16(06)\_\allowbreak{}ODP[10] політик безпеки>}


{\footnotesize\bfseries
No: 4\\
Name: ac\_\allowbreak{}16\_\allowbreak{}6\_\allowbreak{}04\\
Type: list\\
Default: [\textquotedbl{}admin\textquotedbl{}, \textquotedbl{}security\_\allowbreak{}officer\textquotedbl{}]
}

{\footnotesize Персонал зобов'язаний пов'язувати та підтримувати зв'язок <AC-16(06)\_\allowbreak{}ODP[04] атрибутів конфіденційності> з <AC-16(06)\_\allowbreak{}ODP[08] об'єктами> відповідно до <AC-16(06)\_\allowbreak{}ODP[10] політик безпеки>}


{\footnotesize\bfseries
No: 5\\
Name: ac\_\allowbreak{}16\_\allowbreak{}6\_\allowbreak{}odp\_\allowbreak{}01\\
Type: list\\
Default: [\textquotedbl{}default\_\allowbreak{}deny\_\allowbreak{}rule\textquotedbl{}, \textquotedbl{}abac\_\allowbreak{}rule\_\allowbreak{}1\textquotedbl{}]
}

{\footnotesize Визначено атрибути безпеки, які будуть пов'язані з суб'єктами}


{\footnotesize\bfseries
No: 6\\
Name: ac\_\allowbreak{}16\_\allowbreak{}6\_\allowbreak{}odp\_\allowbreak{}02\\
Type: list\\
Default: [\textquotedbl{}default\_\allowbreak{}deny\_\allowbreak{}rule\textquotedbl{}, \textquotedbl{}abac\_\allowbreak{}rule\_\allowbreak{}1\textquotedbl{}]
}

{\footnotesize Визначено атрибути безпеки, які будуть пов'язані з об'єктами}


{\footnotesize\bfseries
No: 7\\
Name: ac\_\allowbreak{}16\_\allowbreak{}6\_\allowbreak{}odp\_\allowbreak{}03\\
Type: list\\
Default: [\textquotedbl{}default\_\allowbreak{}deny\_\allowbreak{}rule\textquotedbl{}, \textquotedbl{}abac\_\allowbreak{}rule\_\allowbreak{}1\textquotedbl{}]
}

{\footnotesize Визначено атрибути конфіденційності, які будуть пов'язані з суб'єктами}


{\footnotesize\bfseries
No: 8\\
Name: ac\_\allowbreak{}16\_\allowbreak{}6\_\allowbreak{}odp\_\allowbreak{}04\\
Type: list\\
Default: [\textquotedbl{}default\_\allowbreak{}deny\_\allowbreak{}rule\textquotedbl{}, \textquotedbl{}abac\_\allowbreak{}rule\_\allowbreak{}1\textquotedbl{}]
}

{\footnotesize Визначено атрибути конфіденційності, які будуть пов'язані з об'єктами}


{\footnotesize\bfseries
No: 9\\
Name: ac\_\allowbreak{}16\_\allowbreak{}6\_\allowbreak{}odp\_\allowbreak{}05\\
Type: string\\
Default: nil
}

{\footnotesize Визначені суб'єкти, які будуть пов'язані з атрибутами безпеки}


{\footnotesize\bfseries
No: 10\\
Name: ac\_\allowbreak{}16\_\allowbreak{}6\_\allowbreak{}odp\_\allowbreak{}06\\
Type: string\\
Default: nil
}

{\footnotesize Визначені об'єкти, які будуть пов'язані з атрибутами безпеки}


{\footnotesize\bfseries
No: 11\\
Name: ac\_\allowbreak{}16\_\allowbreak{}6\_\allowbreak{}odp\_\allowbreak{}07\\
Type: string\\
Default: nil
}

{\footnotesize Визначені суб'єкти, які будуть пов'язані з атрибутами конфіденційності}


{\footnotesize\bfseries
No: 12\\
Name: ac\_\allowbreak{}16\_\allowbreak{}6\_\allowbreak{}odp\_\allowbreak{}08\\
Type: string\\
Default: nil
}

{\footnotesize Визначені об'єкти, які будуть пов'язані з атрибутами конфіденційності}


{\footnotesize\bfseries
No: 13\\
Name: ac\_\allowbreak{}16\_\allowbreak{}6\_\allowbreak{}odp\_\allowbreak{}09\\
Type: list\\
Default: [\textquotedbl{}admin\textquotedbl{}, \textquotedbl{}security\_\allowbreak{}officer\textquotedbl{}]
}

{\footnotesize Політики безпеки, які вимагають від персоналу пов'язувати та підтримувати зв'язок атрибутів безпеки та конфіденційності з суб'єктами та об'єктами}


{\footnotesize\bfseries
No: 14\\
Name: ac\_\allowbreak{}16\_\allowbreak{}6\_\allowbreak{}odp\_\allowbreak{}10\\
Type: list\\
Default: [\textquotedbl{}admin\textquotedbl{}, \textquotedbl{}security\_\allowbreak{}officer\textquotedbl{}]
}

{\footnotesize Політики конфіденційності, які вимагають від персоналу пов'язувати та підтримувати зв'язок атрибутів безпеки та конфіденційності з суб'єктами та об'єктами}




\subsubsection{ПОСЛІДОВНА ІНТЕРПРЕТАЦІЯ АТРИБУТІВ (AC-16(7))}
Забезпечити послідовну інтерпретацію атрибутів безпеки та приватності, що передаються між розподіленими компонентами системи.

\vspace{10pt}
{\footnotesize\bfseries
No: 1\\
Name: ac\_\allowbreak{}16\_\allowbreak{}7\_\allowbreak{}01\\
Type: string\\
Default: nil
}

{\footnotesize Забезпечується послідовна інтерпретація атрибутів безпеки, що передаються між розподіленими компонентами системи}


{\footnotesize\bfseries
No: 2\\
Name: ac\_\allowbreak{}16\_\allowbreak{}7\_\allowbreak{}02\\
Type: string\\
Default: nil
}

{\footnotesize Забезпечується послідовна інтерпретація атрибутів приватності (конфіденційності), що передаються між розподіленими компонентами системи}




\subsubsection{ТЕХНІКИ ТА ТЕХНОЛОГІЇ ПОВ'ЯЗАННЯ АТРИБУТІВ (AC-16(8))}
ПОВ'ЯЗАННЯ АТРИБУТІВ Реалізація [Призначення: методи та технології, визначені організацією] для пов'язування атрибутів безпеки та конфіденційності з інформацією.

\vspace{10pt}
{\footnotesize\bfseries
No: 1\\
Name: ac\_\allowbreak{}16\_\allowbreak{}8\_\allowbreak{}01\\
Type: string\\
Default: nil
}

{\footnotesize Методи та технології застосовуються для пов'язування атрибутів безпеки з інформацією}


{\footnotesize\bfseries
No: 2\\
Name: ac\_\allowbreak{}16\_\allowbreak{}8\_\allowbreak{}02\\
Type: string\\
Default: nil
}

{\footnotesize Методи та технології застосовуються для пов'язування атрибутів конфіденційності з інформацією}


{\footnotesize\bfseries
No: 3\\
Name: ac\_\allowbreak{}16\_\allowbreak{}8\_\allowbreak{}odp\_\allowbreak{}01\\
Type: string\\
Default: nil
}

{\footnotesize Визначено методи та технології, які необхідно застосувати для пов'язання інформації з атрибутами безпеки}


{\footnotesize\bfseries
No: 4\\
Name: ac\_\allowbreak{}16\_\allowbreak{}8\_\allowbreak{}odp\_\allowbreak{}02\\
Type: string\\
Default: nil
}

{\footnotesize Визначено методи та технології, які необхідно застосувати для пов'язання інформації з атрибутами конфіденційності}




\subsubsection{ПЕРЕПРИЗНАЧЕННЯ АТРИБУТІВ (AC-16(9))}
Перепризначення атрибутів безпеки та приватності, пов'язаних з інформацією, здійснювати лише за допомогою механізмів перегляду, перевірених з використанням [Призначення: визначених організацією технік або процедур].

\vspace{10pt}
{\footnotesize\bfseries
No: 1\\
Name: ac\_\allowbreak{}16\_\allowbreak{}9\_\allowbreak{}odp\_\allowbreak{}01\\
Type: string\\
Default: \textquotedbl{}автоматизований засіб моніторингу\textquotedbl{}
}

{\footnotesize Визначено техніки або процедури, що використовуються для перевірки механізмів перегляду при перепризначенні атрибутів безпеки}


{\footnotesize\bfseries
No: 2\\
Name: ac\_\allowbreak{}16\_\allowbreak{}9\_\allowbreak{}odp\_\allowbreak{}02\\
Type: string\\
Default: \textquotedbl{}автоматизований засіб моніторингу\textquotedbl{}
}

{\footnotesize Визначено техніки або процедури, що використовуються для перевірки механізмів перегляду при перепризначенні атрибутів приватності (конфіденційності)}


{\footnotesize\bfseries
No: 3\\
Name: ac\_\allowbreak{}16\_\allowbreak{}9\_\allowbreak{}01\\
Type: string\\
Default: nil
}

{\footnotesize Перепризначення атрибутів безпеки, пов'язаних з інформацією, здійснюється лише за допомогою механізмів перегляду, перевірених з використанням <AC-16(09)\_\allowbreak{}ODP[01] технік або процедур>}


{\footnotesize\bfseries
No: 4\\
Name: ac\_\allowbreak{}16\_\allowbreak{}9\_\allowbreak{}02\\
Type: string\\
Default: nil
}

{\footnotesize Перепризначення атрибутів приватності, пов'язаних з інформацією, здійснюється лише за допомогою механізмів перегляду, перевірених з використанням <AC-16(09)\_\allowbreak{}ODP[02] технік або процедур>}




\subsubsection{КОНФІГУРАЦІЯ АТРИБУТІВ УПОВНОВАЖЕНИМИ ОСОБАМИ (AC-16(10))}
Надати уповноваженим особам можливість визначати або змінювати тип і значення атрибутів безпеки та приватності, доступних для пов'язання із суб'єктами та об'єктами.

\vspace{10pt}
{\footnotesize\bfseries
No: 1\\
Name: ac\_\allowbreak{}16\_\allowbreak{}10\_\allowbreak{}01\\
Type: string\\
Default: nil
}

{\footnotesize Уповноваженим особам надається можливість визначати або змінювати тип і значення атрибутів безпеки, доступних для пов'язання з суб'єктами та об'єктами}


{\footnotesize\bfseries
No: 2\\
Name: ac\_\allowbreak{}16\_\allowbreak{}10\_\allowbreak{}02\\
Type: string\\
Default: nil
}

{\footnotesize Уповноваженим особам надається можливість визначати або змінювати тип і значення атрибутів приватності (конфіденційності), доступних для пов'язання з суб'єктами та об'єктами}




\subsection{ВІДДАЛЕНИЙ ДОСТУП (AC-17)}
a. Встановити та задокументувати обмеження на використання, вимоги до конфігурації/підключення та рекомендації щодо здійснення кожного типу віддаленого доступу.\\ 
b. Авторизувати віддалений доступ до системи, перш ніж будуть дозволені такі підключення.

\vspace{10pt}
{\footnotesize\bfseries
No: 1\\
Name: ac\_\allowbreak{}17\_\allowbreak{}a\_\allowbreak{}01\\
Type: string\\
Default: nil
}

{\footnotesize Для кожного типу дозволеного віддаленого доступу встановлені та задокументовані обмеження на використання}


{\footnotesize\bfseries
No: 2\\
Name: ac\_\allowbreak{}17\_\allowbreak{}a\_\allowbreak{}02\\
Type: string\\
Default: nil
}

{\footnotesize Для кожного типу дозволеного віддаленого доступу встановлені та задокументовані вимоги до конфігурації/підключення}


{\footnotesize\bfseries
No: 3\\
Name: ac\_\allowbreak{}17\_\allowbreak{}a\_\allowbreak{}03\\
Type: string\\
Default: nil
}

{\footnotesize Для кожного типу дозволеного віддаленого доступу встановлені та задокументовані рекомендації щодо здійснення}


{\footnotesize\bfseries
No: 4\\
Name: ac\_\allowbreak{}17\_\allowbreak{}b\_\allowbreak{}01\\
Type: string\\
Default: nil
}

{\footnotesize Кожен тип віддаленого доступу до системи авторизується перед тим, як дозволити такі підключення}




\subsubsection{АВТОМАТИЗОВАНИЙ МОНІТОРИНГ ТА УПРАВЛІННЯ (AC-17(1))}


\vspace{10pt}
{\footnotesize\bfseries
No: 1\\
Name: ac\_\allowbreak{}17\_\allowbreak{}1\_\allowbreak{}01\\
Type: string\\
Default: nil
}

{\footnotesize Проводиться моніторинг методами віддаленого доступу}


{\footnotesize\bfseries
No: 2\\
Name: ac\_\allowbreak{}17\_\allowbreak{}1\_\allowbreak{}02\\
Type: string\\
Default: nil
}

{\footnotesize Проводиться управління методами віддаленого доступу}




\subsubsection{ЗАХИСТ КОНФІДЕНЦІЙНОСТІ ТА ЦІЛІСНОСТІ ЗА ДОПОМОГОЮ ШИФРУВАННЯ (AC-17(2))}


\vspace{10pt}
{\footnotesize\bfseries
No: 1\\
Name: ac\_\allowbreak{}17\_\allowbreak{}2\_\allowbreak{}01\\
Type: string\\
Default: \textquotedbl{}AES-256-GCM\textquotedbl{}
}

{\footnotesize Запроваджено криптографічні механізми для захисту конфіденційності та цілісності сесій віддаленого доступу.}




\subsubsection{КЕРОВАНІ ТОЧКИ КОНТРОЛЮ ДОСТУПУ (AC-17(3))}


\vspace{10pt}
{\footnotesize\bfseries
No: 1\\
Name: ac\_\allowbreak{}17\_\allowbreak{}3\_\allowbreak{}01\\
Type: string\\
Default: nil
}

{\footnotesize Віддалений доступ маршрутизується через авторизовані та керовані точки контролю доступу до мережі}




\subsubsection{ПРИВІЛЕЙОВАНІ КОМАНДИ ТА ДОСТУП (AC-17(4))}


\vspace{10pt}
{\footnotesize\bfseries
No: 1\\
Name: ac\_\allowbreak{}17\_\allowbreak{}4\_\allowbreak{}odp\_\allowbreak{}01\\
Type: string\\
Default: nil
}

{\footnotesize Визначено потреби, що потребують виконання привілейованих команд за допомогою віддаленого доступу}


{\footnotesize\bfseries
No: 2\\
Name: ac\_\allowbreak{}17\_\allowbreak{}4\_\allowbreak{}odp\_\allowbreak{}02\\
Type: string\\
Default: nil
}

{\footnotesize Визначено потреби, що вимагають доступу до інформації, що стосується безпеки, за допомогою віддаленого доступу}


{\footnotesize\bfseries
No: 3\\
Name: ac\_\allowbreak{}17\_\allowbreak{}4\_\allowbreak{}a\_\allowbreak{}01\\
Type: string\\
Default: nil
}

{\footnotesize Виконання привілейованих команд за допомогою віддаленого доступу дозволено лише для наступних потреб: <AC-17(04)\_\allowbreak{}ODP[01] потреби >}


{\footnotesize\bfseries
No: 4\\
Name: ac\_\allowbreak{}17\_\allowbreak{}4\_\allowbreak{}a\_\allowbreak{}02\\
Type: string\\
Default: nil
}

{\footnotesize Доступ до інформації, важливої для безпеки, за допомогою віддаленого доступу дозволяється лише для наступних потреб: <AC-17(04)\_\allowbreak{}ODP[02] потреби >}


{\footnotesize\bfseries
No: 5\\
Name: ac\_\allowbreak{}17\_\allowbreak{}4\_\allowbreak{}b\\
Type: string\\
Default: nil
}

{\footnotesize Обґрунтування віддаленого доступу задокументовано в плані захисту інформації}




\subsubsection{МОНІТОРИНГ ДЛЯ НЕАВТОРИЗОВАНИХ ПІДКЛЮЧЕНЬ (AC-17(5)) [Вилучено]}
[Вилучено: Включено в СІ-04]

\vspace{10pt}
\textit{Немає параметрів для цього контролю.}
\vspace{10pt}


\subsubsection{ЗАХИСТ ІНФОРМАЦІЇ (AC-17(6))}


\vspace{10pt}
{\footnotesize\bfseries
No: 1\\
Name: ac\_\allowbreak{}17\_\allowbreak{}6\_\allowbreak{}01\\
Type: string\\
Default: \textquotedbl{}автоматизований засіб моніторингу\textquotedbl{}
}

{\footnotesize Інформація про механізми віддаленого доступу захищена від неавторизованого використання та розкриття}




\subsubsection{ДОДАТКОВИЙ ЗАХИСТ ДЛЯ ДОСТУПУ ДО ФУНКЦІЙ БЕЗПЕКИ (AC-17(7)) [Вилучено]}
[Вилучено: Включено в AC-03(10)]

\vspace{10pt}
\textit{Немає параметрів для цього контролю.}
\vspace{10pt}


\subsubsection{ДЕАКТИВАЦІЯ НЕЗАХИЩЕНИХ ПРОТОКОЛІВ МЕРЕЖІ (AC-17(8)) [Вилучено]}
[Вилучено: Включено в CM-07]

\vspace{10pt}
\textit{Немає параметрів для цього контролю.}
\vspace{10pt}


\subsubsection{ВІДКЛЮЧЕННЯ АБО ДЕАКТИВАЦІЯ ДОСТУПУ (AC-17(9))}


\vspace{10pt}
{\footnotesize\bfseries
No: 1\\
Name: ac\_\allowbreak{}17\_\allowbreak{}9\_\allowbreak{}01\\
Type: integer\\
Default: 30
}

{\footnotesize Передбачена можливість відключення або деактивації віддаленого доступу до системи протягом <AC-17(09)\_\allowbreak{}ODP> періоду часу>}


{\footnotesize\bfseries
No: 2\\
Name: ac\_\allowbreak{}17\_\allowbreak{}9\_\allowbreak{}odp\\
Type: integer\\
Default: 30
}

{\footnotesize Визначено період часу, протягом якого потрібно відключити або деактивувати віддалений доступ до системи}




\subsubsection{(10) АВТЕНТИФІКАЦІЯ ВІДДАЛЕНИХ КОМАНД (AC-17(10))}


\vspace{10pt}
{\footnotesize\bfseries
No: 1\\
Name: ac\_\allowbreak{}17\_\allowbreak{}10\_\allowbreak{}01\\
Type: string\\
Default: \textquotedbl{}автоматизований засіб моніторингу\textquotedbl{}
}

{\footnotesize Визначені <AC-17(10)\_\allowbreak{}ODP[01] механізми> атентифікують визначені <AC-17(10)\_\allowbreak{}ODP[02] віддалені команди>}


{\footnotesize\bfseries
No: 2\\
Name: ac\_\allowbreak{}17\_\allowbreak{}10\_\allowbreak{}odp\_\allowbreak{}01\\
Type: string\\
Default: \textquotedbl{}автоматизований засіб моніторингу\textquotedbl{}
}

{\footnotesize Визначено механізми, реалізовані для автентифікації віддалених команд}


{\footnotesize\bfseries
No: 3\\
Name: ac\_\allowbreak{}17\_\allowbreak{}10\_\allowbreak{}odp\_\allowbreak{}02\\
Type: string\\
Default: nil
}

{\footnotesize Визначено віддалені команди, які мають бути автентифіковані механізмами}




\subsection{БЕЗДРОТОВИЙ ДОСТУП (AC-18)}
a. Установити обмеження на використання, вимоги до конфігурації/підключення та рекомендації щодо здійснення бездротового доступу.\\ 
b. Авторизувати бездротовий доступ до системи, перш ніж будуть дозволені такі підключення.

\vspace{10pt}
{\footnotesize\bfseries
No: 1\\
Name: ac\_\allowbreak{}18\_\allowbreak{}a\_\allowbreak{}01\\
Type: string\\
Default: nil
}

{\footnotesize Встановлено обмеження на використання щодо здійснення бездротового доступу}


{\footnotesize\bfseries
No: 2\\
Name: ac\_\allowbreak{}18\_\allowbreak{}a\_\allowbreak{}02\\
Type: string\\
Default: nil
}

{\footnotesize Встановлено вимоги до конфігурації або підключення щодо здійснення бездротового доступу}


{\footnotesize\bfseries
No: 3\\
Name: ac\_\allowbreak{}18\_\allowbreak{}a\_\allowbreak{}03\\
Type: string\\
Default: nil
}

{\footnotesize Встановлено рекомендації щодо здійснення бездротового доступу}


{\footnotesize\bfseries
No: 4\\
Name: ac\_\allowbreak{}18\_\allowbreak{}b\\
Type: string\\
Default: nil
}

{\footnotesize Авторизується бездротовий доступ до системи перед тим, як дозволяти такі з'єднання.}




\subsubsection{АВТЕНТИФІКАЦІЯ ТА ШИФРУВАННЯ (AC-18(1))}


\vspace{10pt}
{\footnotesize\bfseries
No: 1\\
Name: ac\_\allowbreak{}18\_\allowbreak{}1\_\allowbreak{}odp\\
Type: list\\
Default: [\textquotedbl{}користувачі\textquotedbl{}, \textquotedbl{}пристрої\textquotedbl{}]
}

{\footnotesize Вибрано одне або декілька з наступних ЗНАЧЕНЬ ПАРАМЕТРІВ: \{користувачі; пристрої\}}


{\footnotesize\bfseries
No: 2\\
Name: ac\_\allowbreak{}18\_\allowbreak{}1\_\allowbreak{}01\\
Type: string\\
Default: nil
}

{\footnotesize Бездротовий доступ до системи захищено за допомогою автентифікації <AC-18(01)\_\allowbreak{}ODP ВИБРАНЕ ЗНАЧЕННЯ ПАРАМЕТРА(ів)>}


{\footnotesize\bfseries
No: 3\\
Name: ac\_\allowbreak{}18\_\allowbreak{}1\_\allowbreak{}02\\
Type: string\\
Default: \textquotedbl{}AES-256-GCM\textquotedbl{}
}

{\footnotesize Бездротовий доступ до системи захищений за допомогою шифрування}




\subsubsection{БЕЗДРОТОВИЙ ПІДКЛЮЧЕНЬ (AC-18(2)) [Вилучено]}
[Вилучено: Включено в SI-04]

\vspace{10pt}
\textit{Немає параметрів для цього контролю.}
\vspace{10pt}


\subsubsection{ВІДКЛЮЧЕННЯ БЕЗДРОТОВОЇ МЕРЕЖІ (AC-18(3))}


\vspace{10pt}
{\footnotesize\bfseries
No: 1\\
Name: ac\_\allowbreak{}18\_\allowbreak{}3\_\allowbreak{}01\\
Type: string\\
Default: nil
}

{\footnotesize Відключено, у разі відсутності необхідності у використанні, вбудовані в компоненти системи можливості бездротових мереж до їх виклику та розгортання}




\subsubsection{БЕЗДРОТОВИЙ ДОСТУП СТУВАЧАМИ (AC-18(4))}


\vspace{10pt}
{\footnotesize\bfseries
No: 1\\
Name: ac\_\allowbreak{}18\_\allowbreak{}4\_\allowbreak{}01\\
Type: string\\
Default: nil
}

{\footnotesize Встановлено користувачів, яким дозволено самостійно налаштовувати можливості бездротової мережі}


{\footnotesize\bfseries
No: 2\\
Name: ac\_\allowbreak{}18\_\allowbreak{}4\_\allowbreak{}02\\
Type: string\\
Default: nil
}

{\footnotesize Явно авторизуються визначені користувачі, яким дозволено самостійно налаштовувати можливості бездротової мережі}




\subsubsection{АНТЕНИ ТА РІВЕНЬ ПОТУЖНОСТІ ПЕРЕДАЧІ (AC-18(5))}


\vspace{10pt}
{\footnotesize\bfseries
No: 1\\
Name: ac\_\allowbreak{}18\_\allowbreak{}5\_\allowbreak{}01\\
Type: string\\
Default: nil
}

{\footnotesize Вибирано такі радіо антени, які зменшують ймовірність того, що корисні сигнали можуть прийматися за межами контрольованих організацією меж}


{\footnotesize\bfseries
No: 2\\
Name: ac\_\allowbreak{}18\_\allowbreak{}5\_\allowbreak{}02\\
Type: string\\
Default: nil
}

{\footnotesize Калібруються рівні потужності передачі, щоб зменшити ймовірність того, що корисні сигнали можуть прийматися за контрольованими організацією межами}




\subsection{КОНТРОЛЬ ДОСТУПУ ДЛЯ МОБІЛЬНИХ ПРИСТРОЇВ (AC-19)}
a. Встановити обмеження на використання, вимоги до конфігурації, вимоги до підключення і рекомендації щодо впровадження мобільних пристроїв, контрольованих організацією.\\ 
b. Авторизувати підключення мобільних пристроїв до систем, які експлуатуються організацією.

\vspace{10pt}
{\footnotesize\bfseries
No: 1\\
Name: ac\_\allowbreak{}19\_\allowbreak{}a\_\allowbreak{}01\\
Type: string\\
Default: nil
}

{\footnotesize Встановлено вимоги до конфігурації мобільних пристроїв, що контролюються організацією, в тому числі, коли такі пристрої перебувають за межами контрольованої території}


{\footnotesize\bfseries
No: 2\\
Name: ac\_\allowbreak{}19\_\allowbreak{}a\_\allowbreak{}02\\
Type: string\\
Default: nil
}

{\footnotesize Встановлюються вимоги до підключення для мобільних пристроїв, що контролюються організацією, в тому числі, коли такі пристрої знаходяться за межами контрольованої території}


{\footnotesize\bfseries
No: 3\\
Name: ac\_\allowbreak{}19\_\allowbreak{}a\_\allowbreak{}03\\
Type: string\\
Default: nil
}

{\footnotesize Розроблено рекомендації щодо впровадження для мобільних пристроїв, які контролюються організацією, в тому числі, коли такі пристрої перебувають за межами контрольованої території}


{\footnotesize\bfseries
No: 4\\
Name: ac\_\allowbreak{}19\_\allowbreak{}b\\
Type: string\\
Default: nil
}

{\footnotesize Авторизовано підключення мобільних пристроїв до систем організації}




\subsubsection{ВИКОРИСТАННЯ ПИСЬМОВИХ ТА ПОРТАТИВНИЙ ПРИСТРОЇВ ДЛЯ ЗБЕРІГАННЯ ДАНИХ (AC-19(1)) [Вилучено]}
[Вилучено: Включено в MP-07]

\vspace{10pt}
\textit{Немає параметрів для цього контролю.}
\vspace{10pt}


\subsubsection{ВИКОРИСТАННЯ ПЕРСОНАЛЬНИХ ПОРТАТИВНИХ ПРИСТРОЇВ ЗБЕРІГАННЯ ДАНИХ (AC-19(2)) [Вилучено]}
[Вилучено: Включено в MP-07]

\vspace{10pt}
\textit{Немає параметрів для цього контролю.}
\vspace{10pt}


\subsubsection{ВИКОРИСТАННЯ ПОРТАТИВНИХ ПРИСТРОЇВ ЗБЕРІГАННЯ ДАНИХ З НЕІДЕНТИФІКОВАНИМ ВЛАСНИКОМ (AC-19(3))}


\vspace{10pt}
\textit{Немає параметрів для цього контролю.}
\vspace{10pt}


\subsubsection{ОБМЕЖЕННЯ ДЛЯ ЗАСЕКРЕЧЕНОЇ ІНФОРМАЦІЇ (AC-19(4))}


\vspace{10pt}
{\footnotesize\bfseries
No: 1\\
Name: ac\_\allowbreak{}19\_\allowbreak{}4\_\allowbreak{}odp\_\allowbreak{}01\\
Type: list\\
Default: [\textquotedbl{}admin\textquotedbl{}, \textquotedbl{}security\_\allowbreak{}officer\textquotedbl{}]
}

{\footnotesize Визначено посадових осіб із захисту інформації, відповідальних за огляд та перевірку захищених мобільних пристроїв та інформації, що зберігається на цих пристроях}


{\footnotesize\bfseries
No: 2\\
Name: ac\_\allowbreak{}19\_\allowbreak{}4\_\allowbreak{}odp\_\allowbreak{}02\\
Type: list\\
Default: [\textquotedbl{}default\_\allowbreak{}deny\_\allowbreak{}rule\textquotedbl{}, \textquotedbl{}abac\_\allowbreak{}rule\_\allowbreak{}1\textquotedbl{}]
}

{\footnotesize Визначено політики безпеки, що обмежують підключення захищених мобільних пристроїв до засекречених систем}


{\footnotesize\bfseries
No: 3\\
Name: ac\_\allowbreak{}19\_\allowbreak{}4\_\allowbreak{}a\\
Type: string\\
Default: nil
}

{\footnotesize Використання незахищених мобільних пристроїв на об'єктах, що містять системи, які обробляють, зберігають або передають секретну інформацію, заборонено, за винятком випадків, коли на це є спеціальний дозвіл уповноваженої посадової особи}


{\footnotesize\bfseries
No: 4\\
Name: ac\_\allowbreak{}19\_\allowbreak{}4\_\allowbreak{}b\_\allowbreak{}01\\
Type: string\\
Default: nil
}

{\footnotesize Заборона підключення незахищених мобільних пристроїв до систем з обмеженим доступом застосовується до осіб, яким уповноважена посадова особа дозволила використовувати незахищені мобільні пристрої на об'єктах, що містять системи, які обробляють, зберігають або передають інформацію з обмеженим доступом}


{\footnotesize\bfseries
No: 5\\
Name: ac\_\allowbreak{}19\_\allowbreak{}4\_\allowbreak{}b\_\allowbreak{}02\\
Type: string\\
Default: nil
}

{\footnotesize Дозвіл уповноваженої посадової особи на підключення незахищених мобільних пристроїв до незахищених систем вимагається від осіб, яким дозволено використовувати незахищені мобільні пристрої на об'єктах, що містять системи, які обробляють, зберігають або передають інформацію з обмеженим доступом}


{\footnotesize\bfseries
No: 6\\
Name: ac\_\allowbreak{}19\_\allowbreak{}4\_\allowbreak{}b\_\allowbreak{}03\\
Type: string\\
Default: nil
}

{\footnotesize Заборона використання внутрішніх або зовнішніх модемів чи бездротових інтерфейсів у складі незахищених мобільних пристроїв поширюється на осіб, яким уповноваженою посадовою особою дозволено використовувати незахищені мобільні пристрої під час виконання службових обов'язків, а також на осіб, які не мають права на використання таких пристроїв}


{\footnotesize\bfseries
No: 7\\
Name: ac\_\allowbreak{}19\_\allowbreak{}4\_\allowbreak{}b\_\allowbreak{}04\_\allowbreak{}01\\
Type: string\\
Default: nil
}

{\footnotesize Вибірковий огляд та перевірка незахищених мобільних пристроїв та інформації, що зберігається на них, <AC-19(04)\_\allowbreak{}ODP[01] посадовими особами> є обов'язковими}


{\footnotesize\bfseries
No: 8\\
Name: ac\_\allowbreak{}19\_\allowbreak{}4\_\allowbreak{}b\_\allowbreak{}04\_\allowbreak{}02\\
Type: string\\
Default: nil
}

{\footnotesize Дотримання політики обробки інцидентів застосовується у разі виявлення секретної інформації під час випадкового огляду та перевірки незахищених мобільних пристроїв}


{\footnotesize\bfseries
No: 9\\
Name: ac\_\allowbreak{}19\_\allowbreak{}4\_\allowbreak{}c\\
Type: string\\
Default: nil
}

{\footnotesize Підключення захищенихмобільних пристроїв до засекречених систем обмежено відповідно до <AC-19(04)\_\allowbreak{}ODP[02] політики безпеки>}




\subsubsection{ПОВНЕ ШИФРУВАННЯ ПРИСТРОЇВ ТА СХОВИЩ ІНФОРМАЦІЇ (AC-19(5))}


\vspace{10pt}
{\footnotesize\bfseries
No: 1\\
Name: ac\_\allowbreak{}19\_\allowbreak{}5\_\allowbreak{}odp\_\allowbreak{}01\\
Type: list\\
Default: [\textquotedbl{}повне шифрування пристроїв\textquotedbl{}, \textquotedbl{}шифрування сховищ інформації\textquotedbl{}]
}

{\footnotesize Вибрано одне з наступних ЗНАЧЕНЬ ПАРАМЕТРІВ: \{повне шифрування пристроїв; шифрування сховищ інформації\}}


{\footnotesize\bfseries
No: 2\\
Name: ac\_\allowbreak{}19\_\allowbreak{}5\_\allowbreak{}odp\_\allowbreak{}02\\
Type: string\\
Default: nil
}

{\footnotesize Визначено мобільні пристрої, на яких слід використовувати шифрування}


{\footnotesize\bfseries
No: 3\\
Name: ac\_\allowbreak{}19\_\allowbreak{}5\_\allowbreak{}01\\
Type: string\\
Default: nil
}

{\footnotesize <AC-19(05)\_\allowbreak{}ODP[01] ВИБРАНЕ ЗНАЧЕННЯ ПАРАМЕТРІВ> використовується для захисту конфіденційності та цілісності інформації на <AC-19(05)\_\allowbreak{}ODP[02] мобільних пристроях>}




\subsection{ВИКОРИСТАННЯ ЗОВНІШНІХ СИСТЕМ (AC-20)}
a. [Вибір (один або кілька): Встановіть [Призначення: умови, визначені організацією]; Визначте [Призначення: визначені організацією засоби контролю, які, як стверджується, будуть реалізовані на зовнішніх системах]], узгоджені з довірчими відносинами, встановленими з іншими організаціями, які володіють, експлуатують та/або обслуговують зовнішні системи, дозволяючи уповноваженим особам:\\ 
1. доступ до системи із зовнішніх систем;\\ 
2. обробляти, зберігати або передавати керовану організацією інформацію за допомогою зовнішніх систем;\\ 
b. Заборонити використання [Призначення: організаційно-визначені типи зовнішніх систем].

\vspace{10pt}
{\footnotesize\bfseries
No: 1\\
Name: ac\_\allowbreak{}20\_\allowbreak{}odp\_\allowbreak{}01\\
Type: list\\
Default: [\textquotedbl{}умови та положення\textquotedbl{}, \textquotedbl{}заходи захисту\textquotedbl{}]
}

{\footnotesize Вибрано одне або більше з наступних ЗНАЧЕНЬ ПАРАМЕТРІВ: \{встановити <AC -20\_\allowbreak{}ODP[02] умови та положення>; визначити <AC-20\_\allowbreak{}ODP[03] заходи захисту>\}}


{\footnotesize\bfseries
No: 2\\
Name: ac\_\allowbreak{}20\_\allowbreak{}odp\_\allowbreak{}02\\
Type: list\\
Default: []
}

{\footnotesize Визначено умови та положення, що відповідають довірчим відносинам, встановленим з іншими організаціями, які володіють, експлуатують та/або обслуговують зовнішні системи (якщо вибрано)}


{\footnotesize\bfseries
No: 3\\
Name: ac\_\allowbreak{}20\_\allowbreak{}odp\_\allowbreak{}03\\
Type: string\\
Default: nil
}

{\footnotesize Визначено заходи захисту, які мають бути застосовані до зовнішніх систем відповідно до довірчих відносин, встановлених з іншими організаціями, що володіють, експлуатують та/або обслуговують зовнішні системи (якщо обрано)}


{\footnotesize\bfseries
No: 4\\
Name: ac\_\allowbreak{}20\_\allowbreak{}odp\_\allowbreak{}04\\
Type: string\\
Default: nil
}

{\footnotesize Визначено типи зовнішніх систем, заборонених до використання}


{\footnotesize\bfseries
No: 5\\
Name: ac\_\allowbreak{}20\_\allowbreak{}a\_\allowbreak{}01\\
Type: string\\
Default: nil
}

{\footnotesize <AC-20\_\allowbreak{}ODP[01] ВИБРАНЕ ЗНАЧЕННЯ ПАРАМЕТРА(ів)> узгоджується(ються) з довірчими відносинами, встановленими з іншими організаціями, які володіють, експлуатують та/або обслуговують зовнішні системи, що дозволяє уповноваженим особам отримувати доступ до системи із зовнішніх систем (якщо це застосовно)}


{\footnotesize\bfseries
No: 6\\
Name: ac\_\allowbreak{}20\_\allowbreak{}a\_\allowbreak{}02\\
Type: string\\
Default: nil
}

{\footnotesize <AC-20\_\allowbreak{}ODP[01] ВИБРАНЕ ЗНАЧЕННЯ ПАРАМЕТРА(ів)> відповідає(ють) довірчим відносинам, встановленим з іншими організаціями, які володіють, експлуатують та/або підтримують зовнішні системи, що дозволяє уповноваженим особам обробляти, зберігати або передавати інформацію, за допомогою зовнішніх систем (якщо це застосовно)}


{\footnotesize\bfseries
No: 7\\
Name: ac\_\allowbreak{}20\_\allowbreak{}b\\
Type: string\\
Default: nil
}

{\footnotesize Заборонено використання <AC-20\_\allowbreak{}ODP[04] заборонені типи зовнішніх систем> (якщо застосовно)}




\subsubsection{ОБМЕЖЕННЯ НА АВТОРИЗОВАНЕ ВИКОРИСТАННЯ (AC-20(1))}


\vspace{10pt}
{\footnotesize\bfseries
No: 1\\
Name: ac\_\allowbreak{}20\_\allowbreak{}1\_\allowbreak{}a\\
Type: list\\
Default: [\textquotedbl{}admin\textquotedbl{}, \textquotedbl{}security\_\allowbreak{}officer\textquotedbl{}]
}

{\footnotesize Авторизовані особи мають право використовувати зовнішню систему для доступу до системи або для обробки, зберігання чи передачі інформації, що контролюється організацією, лише після перевірки виконання заходів безпеки та конфіденційності, зазначених у політиці безпеки та конфіденційності організації, а також планах безпеки та конфіденційності (якщо такі застосовуються)}


{\footnotesize\bfseries
No: 2\\
Name: ac\_\allowbreak{}20\_\allowbreak{}1\_\allowbreak{}b\\
Type: list\\
Default: [\textquotedbl{}admin\textquotedbl{}, \textquotedbl{}security\_\allowbreak{}officer\textquotedbl{}]
}

{\footnotesize Авторизовані особи мають право використовувати зовнішню систему для доступу до системи або для обробки, зберігання чи передачі інформації, що контролюється організацією, лише після збереження погоджених угод про підключення або обробку системи з структурою організації, на якій розміщена зовнішня система}




\subsubsection{ПЕРЕНОСНІ ПРИСТРОЇ ЗБЕРІГАННЯ ДАНИХ (AC-20(2))}


\vspace{10pt}
{\footnotesize\bfseries
No: 1\\
Name: ac\_\allowbreak{}20\_\allowbreak{}2\_\allowbreak{}01\\
Type: list\\
Default: [\textquotedbl{}admin\textquotedbl{}, \textquotedbl{}security\_\allowbreak{}officer\textquotedbl{}]
}

{\footnotesize Використання портативних пристроїв носіїв інформації уповноваженими особами обмежено у зовнішніх системах за допомогою обмеження }


{\footnotesize\bfseries
No: 2\\
Name: ac\_\allowbreak{}20\_\allowbreak{}2\_\allowbreak{}odp\\
Type: list\\
Default: [\textquotedbl{}admin\textquotedbl{}, \textquotedbl{}security\_\allowbreak{}officer\textquotedbl{}]
}

{\footnotesize Визначено обмеження на використання авторизованими особами портативних носіїв інформації у зовнішніх системах}




\subsubsection{СИСТЕМИ ТА КОМПОНЕНТИ, ЩО НЕ ЗНАХОДЯТЬСЯ У ВЛАСНОСТІ ОРГАНІЗАЦІЇ (AC-20(3))}


\vspace{10pt}
{\footnotesize\bfseries
No: 1\\
Name: ac\_\allowbreak{}20\_\allowbreak{}3\_\allowbreak{}01\\
Type: string\\
Default: nil
}

{\footnotesize Використання систем або компонентів систем, що не належать організації, для обробки, зберігання або передачі інформації, що належить організації, обмежується за допомогою обмеження}


{\footnotesize\bfseries
No: 2\\
Name: ac\_\allowbreak{}20\_\allowbreak{}3\_\allowbreak{}odp\\
Type: string\\
Default: nil
}

{\footnotesize Визначено обмеження на використання систем або компонентів систем, що не належать організації, для обробки, зберігання або передачі інформації організації}




\subsubsection{ПРИСТРОЇ ДЛЯ ЗБЕРІГАННЯ ДАНИХ, ЯКІ МОЖУТЬ МАТИ ДОСТУП ДО МЕРЕЖІ (AC-20(4))}


\vspace{10pt}
{\footnotesize\bfseries
No: 1\\
Name: ac\_\allowbreak{}20\_\allowbreak{}4\_\allowbreak{}odp\\
Type: string\\
Default: nil
}

{\footnotesize Визначано мережеві носії інформації , які можуть мати доступ до мережі}


{\footnotesize\bfseries
No: 2\\
Name: ac\_\allowbreak{}20\_\allowbreak{}4\_\allowbreak{}01\\
Type: string\\
Default: nil
}

{\footnotesize Заборонено використовувати <AC-20(03)\_\allowbreak{}ODP мережеві носії інформації> у зовнішніх системах}




\subsubsection{ПОРТАТИВНІ ПРИСТРОЇ ДЛЯ ЗБЕРІГАННЯ ДАНИХ -- ЗАБОРОНА ВИКОРИСТАННЯ (AC-20(5))}


\vspace{10pt}
{\footnotesize\bfseries
No: 1\\
Name: ac\_\allowbreak{}20\_\allowbreak{}5\_\allowbreak{}01\\
Type: list\\
Default: [\textquotedbl{}admin\textquotedbl{}, \textquotedbl{}security\_\allowbreak{}officer\textquotedbl{}]
}

{\footnotesize Використання уповноваженими особами зовнішніх носіїв інформації, підконтрольних організації, на зовнішніх системах заборонено}




\subsection{РОЗПОВСЮДЖЕННЯ ІНФОРМАЦІЇ (AC-21)}
a. Спростити обмін інформацією, надаючи авторизованим користувачам змогу визначати, чи відповідають повноваження на доступ, що призначені партнерам для обміну, обмеженням доступу та повноваженням з приватності щодо інформації для [Призначення: визначених організацією обставин обміну інформацією, коли це необхідно користувачу].\\ 
b. Використовувати [Призначення: визначені організацією автоматизовані механізми або ручні процеси], щоб допомогти користувачам в ухваленні рішень щодо обміну інформацією та співпраці.

\vspace{10pt}
{\footnotesize\bfseries
No: 1\\
Name: ac\_\allowbreak{}21\_\allowbreak{}a\\
Type: string\\
Default: nil
}

{\footnotesize Авторизованим користувачам дозволено визначати, чи відповідають повноваження доступу, призначені партнеру з обміну, обмеженням доступу та використання інформації для обставин обміну інформацією}


{\footnotesize\bfseries
No: 2\\
Name: ac\_\allowbreak{}21\_\allowbreak{}b\\
Type: string\\
Default: \textquotedbl{}автоматизований засіб моніторингу\textquotedbl{}
}

{\footnotesize Автоматизовані механізми використовуються для допомоги користувачам у прийнятті рішень щодо обміну інформацією та співпраці}


{\footnotesize\bfseries
No: 3\\
Name: ac\_\allowbreak{}21\_\allowbreak{}odp\_\allowbreak{}01\\
Type: string\\
Default: nil
}

{\footnotesize Визначені обставини обміну інформацією, за яких користувач повинен на власний розсуд визначати, чи відповідають повноваження доступу, надані партнеру з обміну, обмеженням доступу та використання інформації}


{\footnotesize\bfseries
No: 4\\
Name: ac\_\allowbreak{}21\_\allowbreak{}odp\_\allowbreak{}02\\
Type: string\\
Default: \textquotedbl{}автоматизований засіб моніторингу\textquotedbl{}
}

{\footnotesize Визначено автоматизовані механізми або ручні процеси, які допомагають користувачам у ухваленні рішень щодо обміну інформацією та співпраці}




\subsubsection{АВТОМАТИЧНА ПІДТРИМКА УХВАЛЕННЯ РІШЕНЬ (AC-21(1))}


\vspace{10pt}
{\footnotesize\bfseries
No: 1\\
Name: ac\_\allowbreak{}21\_\allowbreak{}1\_\allowbreak{}01\\
Type: string\\
Default: \textquotedbl{}автоматизований засіб моніторингу\textquotedbl{}
}

{\footnotesize Автоматизовані механізми використовуються для забезпечення виконання рішень щодо обміну інформацією уповноваженими користувачами на основі дозволів доступу партнерів з обміну та обмежень доступу до інформації, що підлягає обміну}


{\footnotesize\bfseries
No: 2\\
Name: ac\_\allowbreak{}21\_\allowbreak{}1\_\allowbreak{}odp\\
Type: string\\
Default: \textquotedbl{}автоматизований засіб моніторингу\textquotedbl{}
}

{\footnotesize Визначено автоматизовані механізми, що застосовуються для забезпечення виконання рішень про спільний доступ до інформації авторизованими користувачами}




\subsubsection{РОЗПОВСЮДЖЕННЯ ІНФОРМАЦІЇ (AC-21(2))}


\vspace{10pt}
{\footnotesize\bfseries
No: 1\\
Name: ac\_\allowbreak{}21\_\allowbreak{}2\_\allowbreak{}01\\
Type: string\\
Default: nil
}

{\footnotesize Впроваджено сервіси пошуку та перевірки інформації, які застосовують e обмеження щодо обміну інформацією}


{\footnotesize\bfseries
No: 2\\
Name: ac\_\allowbreak{}21\_\allowbreak{}2\_\allowbreak{}odp\\
Type: string\\
Default: nil
}

{\footnotesize Визначено обмеження щодо обміну інформацією}




\subsection{ПУБЛІЧНО ДОСТУПНИЙ КОНТЕНТ (AC-22)}
a. Призначити осіб, що уповноважені на розміщення інформації в загальнодоступній системі.\\ 
b. Навчати уповноважених осіб тому, щоб загальнодоступна інформація не містила інформацію з обмеженим доступом.\\ 
c. Переглядати запропонований зміст інформації до публікації в загальнодоступній системі, щоб гарантувати, що там не міститься інформація з обмеженим доступом.\\ 
d. Переглядати вміст загальнодоступної системи на предмет наявності там інформації з обмеженим доступом з [Призначення: визначеною організацією частотою]; така інформація має бути видалена в разі її виявлення.

\vspace{10pt}
{\footnotesize\bfseries
No: 1\\
Name: ac\_\allowbreak{}22\_\allowbreak{}odp\_\allowbreak{}01\\
Type: string\\
Default: nil
}

{\footnotesize Визначено частоту, з якою слід переглядати вміст загальнодоступної системи на предмет наявності там інформації з обмеженим доступом}


{\footnotesize\bfseries
No: 2\\
Name: ac\_\allowbreak{}22\_\allowbreak{}a\\
Type: string\\
Default: nil
}

{\footnotesize Визначені особи уповноважені на розміщення інформації в загальнодоступній системі}


{\footnotesize\bfseries
No: 3\\
Name: ac\_\allowbreak{}22\_\allowbreak{}b\\
Type: string\\
Default: nil
}

{\footnotesize Уповноважені особи проходять навчання, щоб гарантувати, що загальнодоступна інформація не містить інформацію з обмеженим доступом}


{\footnotesize\bfseries
No: 4\\
Name: ac\_\allowbreak{}22\_\allowbreak{}c\\
Type: string\\
Default: nil
}

{\footnotesize Запропонований зміст інформації перевіряється до публікації в загальнодоступній системі}


{\footnotesize\bfseries
No: 5\\
Name: ac\_\allowbreak{}22\_\allowbreak{}d\_\allowbreak{}01\\
Type: string\\
Default: nil
}

{\footnotesize Зміст у загальнодоступній системі перевіряється на наявність інформації з обмеженим доступом з <AC-22\_\allowbreak{}ODP[01] частотою>}


{\footnotesize\bfseries
No: 6\\
Name: ac\_\allowbreak{}22\_\allowbreak{}d\_\allowbreak{}02\\
Type: string\\
Default: nil
}

{\footnotesize Інформація з обмеженим доступом видаляється з загальнодоступної системи, якщо її виявлено}




\subsection{ЗАХИСТ ВІД НЕСАНКЦІОНОВАНОГО ІНТЕЛЕКТУАЛЬНОГО АНАЛІЗУ ДАНИХ (AC-23)}
Використовувати [Призначення: визначені організацією техніки виявлення та попередження витоку даних] для [Призначення: визначених організацією об'єктів зберігання даних] для виявлення та захисту від несанкціонованого інтелектуального аналізу даних.

\vspace{10pt}
{\footnotesize\bfseries
No: 1\\
Name: ac\_\allowbreak{}23\_\allowbreak{}odp\_\allowbreak{}01\\
Type: string\\
Default: nil
}

{\footnotesize Визначено техніки виявлення та попередження витоку даних}


{\footnotesize\bfseries
No: 2\\
Name: ac\_\allowbreak{}23\_\allowbreak{}odp\_\allowbreak{}02\\
Type: string\\
Default: nil
}

{\footnotesize Визначено об'єкти зберігання даних}


{\footnotesize\bfseries
No: 3\\
Name: ac\_\allowbreak{}23\_\allowbreak{}01\\
Type: string\\
Default: nil
}

{\footnotesize <AC-23\_\allowbreak{}ODP[01] Техніки> використовуються для <AC-23\_\allowbreak{}ODP[02] об'єктів зберігання даних> для виявлення та захисту від несанкціонованого інтелектуального аналізу даних}




\subsection{РІШЕННЯ ЩОДО УПРАВЛІННЯ ДОСТУПОМ (AC-24)}
[Вибір: Встановити процедури; Запровадити механізми], щоб забезпечити застосування [Призначення: визначені організацією рішення щодо контролю доступу] до кожного запиту щодо доступу до виконання доступу.

\vspace{10pt}
{\footnotesize\bfseries
No: 1\\
Name: ac\_\allowbreak{}24\_\allowbreak{}odp\_\allowbreak{}01\\
Type: string\\
Default: nil
}

{\footnotesize Вибрано: встановити процедури; запровадити механізми}


{\footnotesize\bfseries
No: 2\\
Name: ac\_\allowbreak{}24\_\allowbreak{}odp\_\allowbreak{}02\\
Type: string\\
Default: nil
}

{\footnotesize Визначено рішення щодо контролю доступу}


{\footnotesize\bfseries
No: 3\\
Name: ac\_\allowbreak{}24\_\allowbreak{}01\\
Type: string\\
Default: nil
}

{\footnotesize Відповідно до <AC-24\_\allowbreak{}ODP[01] вибору>, забезпечується застосування <AC-24\_\allowbreak{}ODP[02] рішень щодо контролю доступу> до кожного запиту щодо доступу до виконання доступу}




\subsubsection{ІНФОРМАЦІЯ ПРО ПЕРЕДАЧУ АВТОРИЗОВАНОГО ДОСТУПУ (AC-24(1))}


\vspace{10pt}
{\footnotesize\bfseries
No: 1\\
Name: ac\_\allowbreak{}24\_\allowbreak{}1\_\allowbreak{}01\\
Type: string\\
Default: nil
}

{\footnotesize Інформація щодо авторизації доступу передається за допомогою заходів безпеки до систем, які забезпечують ухвалення рішень щодо управління доступом. ухвалення}


{\footnotesize\bfseries
No: 2\\
Name: ac\_\allowbreak{}24\_\allowbreak{}1\_\allowbreak{}odp\_\allowbreak{}01\\
Type: string\\
Default: nil
}

{\footnotesize Визначено інформацію щодо авторизації доступу, яка передається до систем, які забезпечують ухвалення рішень щодо управління доступом}


{\footnotesize\bfseries
No: 3\\
Name: ac\_\allowbreak{}24\_\allowbreak{}1\_\allowbreak{}odp\_\allowbreak{}02\\
Type: string\\
Default: nil
}

{\footnotesize Визначено заходи безпеки, які слід використовувати, коли інформація про авторизацію передається до систем, що забезпечують виконання рішень щодо управління доступом}


{\footnotesize\bfseries
No: 4\\
Name: ac\_\allowbreak{}24\_\allowbreak{}1\_\allowbreak{}odp\_\allowbreak{}03\\
Type: string\\
Default: nil
}

{\footnotesize Визначено системи, які забезпечують рішень щодо управління доступом}




\subsubsection{ВІДСУТНІСТЬ ІДЕНТИФІКАЦІЇ КОРИСТУВАЧА АБО ПРОЦЕСУ, ЩО ДІЄ ВІД ІМЕНІ КОРИСТУВАЧА (AC-24(2))}
Здійснювати ухвалення рішень щодо управління доступом, засновуючись на [Призначення: визначених організацією атрибутах безпеки], які не охоплюють ідентифікацію користувача або процесу, що діє від імені користувача.

\vspace{10pt}
{\footnotesize\bfseries
No: 1\\
Name: ac\_\allowbreak{}24\_\allowbreak{}2\_\allowbreak{}odp\_\allowbreak{}01\\
Type: string\\
Default: nil
}

{\footnotesize Визначено атрибути безпеки, які не охоплюють ідентифікацію користувача або процесу, що діє від імені користувача (якщо вибрано)}


{\footnotesize\bfseries
No: 2\\
Name: ac\_\allowbreak{}24\_\allowbreak{}2\_\allowbreak{}odp\_\allowbreak{}02\\
Type: string\\
Default: nil
}

{\footnotesize Визначено атрибути конфіденційності, які не охоплюють ідентифікацію користувача або процесу, що діє від імені користувача (якщо вибрано)}


{\footnotesize\bfseries
No: 3\\
Name: ac\_\allowbreak{}24\_\allowbreak{}2\_\allowbreak{}01\\
Type: string\\
Default: nil
}

{\footnotesize Рішення щодо управління доступом здійснюються на основі <AC-24(02)\_\allowbreak{}ODP[01] атрибутів безпеки> , які не охоплюють ідентифікацію користувача або процесу, що діє від імені користувача (якщо вибрано)}


{\footnotesize\bfseries
No: 4\\
Name: ac\_\allowbreak{}24\_\allowbreak{}2\_\allowbreak{}02\\
Type: string\\
Default: nil
}

{\footnotesize Рішення щодо управління доступом здійснюються на основі <AC-24(02)\_\allowbreak{}ODP[02] атрибутів конфіденційності>, які не охоплюють ідентифікацію користувача або процесу, що діє від імені користувача (якщо вибрано)}




\subsection{ДИСПЕТЧЕР ДОСТУПУ (AC-25)}
Впровадити диспетчер доступу для [Призначення: визначеної організацією політики контролю доступу], який захищений від несанкціонованого доступу, завжди був доступний для виклику та досить компактний, щоб бути підданим аналізу й тестуванню, надійність якого може бути гарантована.

\vspace{10pt}
{\footnotesize\bfseries
No: 1\\
Name: ac\_\allowbreak{}25\_\allowbreak{}01\\
Type: list\\
Default: [\textquotedbl{}default\_\allowbreak{}deny\_\allowbreak{}rule\textquotedbl{}, \textquotedbl{}abac\_\allowbreak{}rule\_\allowbreak{}1\textquotedbl{}]
}

{\footnotesize Реалізовано испетчер доступу для <AC-25\_\allowbreak{}ODP політики контролю доступу>, який захищений від несанкціонованого доступу, завжди доступний для виклику та досить компактний, щоб бути підданим аналізу й тестуванню, надійність якого може бути гарантована}


{\footnotesize\bfseries
No: 2\\
Name: ac\_\allowbreak{}25\_\allowbreak{}odp\\
Type: list\\
Default: [\textquotedbl{}default\_\allowbreak{}deny\_\allowbreak{}rule\textquotedbl{}, \textquotedbl{}abac\_\allowbreak{}rule\_\allowbreak{}1\textquotedbl{}]
}

{\footnotesize Визначено політики контролю доступу, для яких реалізовано диспетчер доступу}




\section{AT}
Клас заходів захисту AT --- ОБІЗНАНІСТЬ ТА НАВЧАННЯ

Цей клас спрямований на забезпечення належного рівня знань персоналу щодо загроз інформаційній безпеці та їхніх обов'язків у цій сфері.

\textbf{Перелік заходів захисту:} Політика та процедури підвищення обізнаності та навчання (AT-1); Навчання з підвищення обізнаності (AT-2); Практичні заняття (AT-2(1)); Внутрішні загрози (AT-2(2)); Соціальна інженерія та соціальний інтелектуальний аналіз даних (AT-2(3)); Підозрілі повідомлення та аномальна поведінка системи (AT-2(4)); Вдосконалена стійка загроза (AT-2(5)); Середовище кіберзагроз (AT-2(6)); Рольове навчання (AT-3); Заходи безпеки робочого середовища (AT-3(1)); Фізичні заходи безпеки (AT-3(2)); Практичні заняття (AT-3(3)); Підозрілі зв'язки та аномальна поведінка системи (AT-3(4)); Обробка персональних даних (AT-3(5)); Навчальні записи (AT-4); Контакти з групами безпеки та асоціаціями (AT-5) [Вилучено]; Відгуки про проведені навчання (AT-6).

\subsection{ПОЛІТИКА ТА ПРОЦЕДУРИ ПІДВИЩЕННЯ ОБІЗНАНОСТІ ТА НАВЧАННЯ (AT-1)}
a. Розробити, задокументувати та поширити [Призначення: серед визначеного організацією персоналу або ролей]:\\ 
1. [Вибір (один або декілька): Рівень організації; Рівень місії/бізнес-процесу; рівень системи] політики обізнаності та навчання у сфері забезпечення безпеки та приватності, яка:\\ 
(a) містить мету, сферу застосування, ролі, обов'язки, відповідальність керівництва, координацію між організаційними підрозділами та систему контролю відповідності (complaince);\\ 
(b) відповідає чинним законам, нормативним документам, директивам, нормам, політикам, стандартам та керівним документам.\\ 
2. Процедури, що сприяють реалізації політики підвищення обізнаності та професійної підготовки в галузі безпеки, приватності, а також пов'язаних з ними заходів захисту інформації та персональних даних.\\ 
b. Призначити [Призначення: визначену організацією посадову особу] для управління політикою та процедурами підвищення обізнаності та навчання у сфері забезпечення безпеки та приватності.\\ 
c. Переглядати та оновлювати:\\ 
1. Поточну політика [Призначення: частота, визначена організацією] і наступне [Призначення: події, визначені організацією];\\ 
2. Процедури [Призначення: частота, визначена організацією] та наступні [Призначення: події, визначені організацією].

\vspace{10pt}
{\footnotesize\bfseries
No: 1\\
Name: at\_\allowbreak{}1\_\allowbreak{}odp\_\allowbreak{}01\\
Type: string\\
Default: nil
}

{\footnotesize Визначено персонал або ролі, серед яких має бути поширена політика обізнаності та навчання}


{\footnotesize\bfseries
No: 2\\
Name: at\_\allowbreak{}1\_\allowbreak{}odp\_\allowbreak{}02\\
Type: string\\
Default: nil
}

{\footnotesize Визначено персонал або ролі, серед яких мають бути поширені процедури, що сприяють реалізації політики підвищення обізнаності}


{\footnotesize\bfseries
No: 3\\
Name: at\_\allowbreak{}1\_\allowbreak{}odp\_\allowbreak{}03\\
Type: list\\
Default: [\textquotedbl{}рівень організації\textquotedbl{}, \textquotedbl{}рівень місії/бізнес-процесу\textquotedbl{}, \textquotedbl{}рівень системи\textquotedbl{}]
}

{\footnotesize Вибрано одне або декілька з наступних ЗНАЧЕНЬ ПАРАМЕТРІВ: \{рівень організації; рівень місії/бізнес-процесу; рівень системи\}}


{\footnotesize\bfseries
No: 4\\
Name: at\_\allowbreak{}1\_\allowbreak{}odp\_\allowbreak{}04\\
Type: string\\
Default: nil
}

{\footnotesize Визначено посадову особу, яка керуватиме політикою та процедурами підвищення обізнаності та навчання}


{\footnotesize\bfseries
No: 5\\
Name: at\_\allowbreak{}1\_\allowbreak{}odp\_\allowbreak{}05\\
Type: string\\
Default: nil
}

{\footnotesize Визначено частоту, з якою переглядається та оновлюється поточна політика інформування}


{\footnotesize\bfseries
No: 6\\
Name: at\_\allowbreak{}1\_\allowbreak{}odp\_\allowbreak{}06\\
Type: string\\
Default: nil
}

{\footnotesize Визначено події, які потребують перегляду та оновлення поточної політики}


{\footnotesize\bfseries
No: 7\\
Name: at\_\allowbreak{}1\_\allowbreak{}odp\_\allowbreak{}07\\
Type: string\\
Default: nil
}

{\footnotesize Визначено частоту, з якою переглядаються та оновлюються поточні процедури}


{\footnotesize\bfseries
No: 8\\
Name: at\_\allowbreak{}1\_\allowbreak{}odp\_\allowbreak{}08\\
Type: string\\
Default: nil
}

{\footnotesize Визначено події, які потребують перегляду та оновлення поточних процедур}


{\footnotesize\bfseries
No: 9\\
Name: at\_\allowbreak{}1\_\allowbreak{}a\_\allowbreak{}01\\
Type: string\\
Default: nil
}

{\footnotesize Розроблено та задокументовано політику обізнаності та навчання}


{\footnotesize\bfseries
No: 10\\
Name: at\_\allowbreak{}1\_\allowbreak{}a\_\allowbreak{}02\\
Type: string\\
Default: nil
}

{\footnotesize Політика обізнаності та навчання поширюється на <AT-01\_\allowbreak{}ODP[01] персонал або ролі>}


{\footnotesize\bfseries
No: 11\\
Name: at\_\allowbreak{}1\_\allowbreak{}a\_\allowbreak{}03\\
Type: string\\
Default: nil
}

{\footnotesize Розроблені та задокументовані процедури, що сприяють впровадженню політики обізнаності та навчання і пов'язаних з нею засобів контролю доступу}


{\footnotesize\bfseries
No: 12\\
Name: at\_\allowbreak{}1\_\allowbreak{}a\_\allowbreak{}04\\
Type: string\\
Default: nil
}

{\footnotesize Процедури поширюються на <AT-01\_\allowbreak{}ODP[02] персонал або ролі>}


{\footnotesize\bfseries
No: 13\\
Name: at\_\allowbreak{}1\_\allowbreak{}a\_\allowbreak{}01\_\allowbreak{}a\_\allowbreak{}01\\
Type: string\\
Default: nil
}

{\footnotesize <AT-01\_\allowbreak{}ODP[03] ВИБРАНЕ ЗНАЧЕННЯ ПАРАМЕТРА(ІВ)> політика обізнаності та навчання містить мету}


{\footnotesize\bfseries
No: 14\\
Name: at\_\allowbreak{}1\_\allowbreak{}a\_\allowbreak{}01\_\allowbreak{}a\_\allowbreak{}02\\
Type: string\\
Default: nil
}

{\footnotesize <AT-01\_\allowbreak{}ODP[03] ВИБРАНЕ ЗНАЧЕННЯ ПАРАМЕТРА(ІВ)> політика обізнаності та навчання містить сферу застосування}


{\footnotesize\bfseries
No: 15\\
Name: at\_\allowbreak{}1\_\allowbreak{}a\_\allowbreak{}01\_\allowbreak{}a\_\allowbreak{}03\\
Type: string\\
Default: nil
}

{\footnotesize <AT-01\_\allowbreak{}ODP[03] ВИБРАНЕ ЗНАЧЕННЯ ПАРАМЕТРА(ІВ)> політика обізнаності та навчання містить ролі}


{\footnotesize\bfseries
No: 16\\
Name: at\_\allowbreak{}1\_\allowbreak{}a\_\allowbreak{}01\_\allowbreak{}a\_\allowbreak{}04\\
Type: string\\
Default: nil
}

{\footnotesize <AT-01\_\allowbreak{}ODP[03] ВИБРАНЕ ЗНАЧЕННЯ ПАРАМЕТРА(ІВ)> політика обізнаності та навчання містить обов'язки}


{\footnotesize\bfseries
No: 17\\
Name: at\_\allowbreak{}1\_\allowbreak{}a\_\allowbreak{}01\_\allowbreak{}a\_\allowbreak{}05\\
Type: string\\
Default: nil
}

{\footnotesize <AT-01\_\allowbreak{}ODP[03] ВИБРАНЕ ЗНАЧЕННЯ ПАРАМЕТРА(ІВ)> політика обізнаності та навчання містить відповідальність керівництва}


{\footnotesize\bfseries
No: 18\\
Name: at\_\allowbreak{}1\_\allowbreak{}a\_\allowbreak{}01\_\allowbreak{}a\_\allowbreak{}06\\
Type: string\\
Default: nil
}

{\footnotesize <AT-01\_\allowbreak{}ODP[03] ВИБРАНЕ ЗНАЧЕННЯ ПАРАМЕТРА(ІВ)> політика обізнаності та навчання містить координацію між підрозділами організації}


{\footnotesize\bfseries
No: 19\\
Name: at\_\allowbreak{}1\_\allowbreak{}a\_\allowbreak{}01\_\allowbreak{}a\_\allowbreak{}07\\
Type: string\\
Default: nil
}

{\footnotesize <AT-01\_\allowbreak{}ODP[03] ВИБРАНЕ ЗНАЧЕННЯ ПАРАМЕТРА(ІВ)> політика обізнаності та навчання містить систему контролю відповідності}


{\footnotesize\bfseries
No: 20\\
Name: at\_\allowbreak{}1\_\allowbreak{}a\_\allowbreak{}01\_\allowbreak{}b\\
Type: string\\
Default: nil
}

{\footnotesize <AT-01\_\allowbreak{}ODP[03] ВИБРАНЕ ЗНАЧЕННЯ ПАРАМЕТРА(ІВ)> політика обізнаності та навчання відповідає чинним законам, нормативним документам, директивам, нормам, політикам, стандартам та керівним документам}


{\footnotesize\bfseries
No: 21\\
Name: at\_\allowbreak{}1\_\allowbreak{}b\\
Type: string\\
Default: nil
}

{\footnotesize <AT-01\_\allowbreak{}ODP[04] посадова особа> призначається для управління політикою та процедурами підвищення обізнаності та навчання у сфері забезпечення безпеки та конфіденційності}


{\footnotesize\bfseries
No: 22\\
Name: at\_\allowbreak{}1\_\allowbreak{}c\_\allowbreak{}01\_\allowbreak{}01\\
Type: string\\
Default: nil
}

{\footnotesize Переглядається та оновлюється поточна політика обізнаності та навчання з <AT-01\_\allowbreak{}ODP[05] частототою>}


{\footnotesize\bfseries
No: 23\\
Name: at\_\allowbreak{}1\_\allowbreak{}c\_\allowbreak{}01\_\allowbreak{}02\\
Type: string\\
Default: nil
}

{\footnotesize Переглядається та оновлюється поточна політика обізнаності та навчання після <AT-01\_\allowbreak{}ODP[06] подій>}


{\footnotesize\bfseries
No: 24\\
Name: at\_\allowbreak{}1\_\allowbreak{}c\_\allowbreak{}02\_\allowbreak{}01\\
Type: string\\
Default: nil
}

{\footnotesize Переглядаються та оновлюються поточні процедури обізнаності та навчання з <AT-01\_\allowbreak{}ODP[07] частототою>}


{\footnotesize\bfseries
No: 25\\
Name: at\_\allowbreak{}1\_\allowbreak{}c\_\allowbreak{}02\_\allowbreak{}02\\
Type: string\\
Default: nil
}

{\footnotesize Переглядаються та оновлюються поточні процедури обізнаності та навчання після <AT-01\_\allowbreak{}ODP[08] подій>}




\subsection{НАВЧАННЯ З ПІДВИЩЕННЯ ОБІЗНАНОСТІ (AT-2)}
Впровадити базові тренінги з підвищення обізнаності у сфері безпеки та приватності для користувачів системи (включно з менеджерами, керівниками компаній і підрядниками):\\ 
a. Забезпечити навчання грамотності з питань безпеки та конфіденційності для користувачів системи (включаючи менеджерів, керівників вищої ланки та підрядників):\\ 
1. як частину початкового навчання для нових користувачів і [Призначення: частота, визначена організацією] після цього;\\ 
2. якщо цього потребують системні зміни або наступні [Призначення: події, визначені організацією].\\ 
b. Використовувати наведені нижче методи, щоб підвищити рівень безпеки та конфіденційності користувачів системи [Завдання: визначені організацією методи поінформованості];\\ 
c. Оновлювати навчання грамотності та зміст обізнаності [Завдання: частота, визначена організацією] і наступні [Завдання: події, визначені організацією];\\ 
d. Включити уроки, отримані з внутрішніх або зовнішніх інцидентів безпеки або порушень, у навчання грамотності та методи підвищення обізнаності.

\vspace{10pt}
{\footnotesize\bfseries
No: 1\\
Name: at\_\allowbreak{}2\_\allowbreak{}odp\_\allowbreak{}01\\
Type: string\\
Default: \textquotedbl{}щорічно\textquotedbl{}
}

{\footnotesize Визначено періодичність проведення навчання грамотності з питань безпеки для користувачів системи (в тому числі менеджерів, вищого керівництва та підрядників) після початкового тренінгу}


{\footnotesize\bfseries
No: 2\\
Name: at\_\allowbreak{}2\_\allowbreak{}odp\_\allowbreak{}02\\
Type: string\\
Default: \textquotedbl{}щорічно\textquotedbl{}
}

{\footnotesize Визначено періодичність проведення навчання грамотності з питань конфіденційності для користувачів системи (в тому числі менеджерів, вищого керівництва та підрядників) після початкового тренінгу}


{\footnotesize\bfseries
No: 3\\
Name: at\_\allowbreak{}2\_\allowbreak{}odp\_\allowbreak{}03\\
Type: string\\
Default: nil
}

{\footnotesize Визначено події, які потребують навчання користувачів системи грамотності з питань безпеки}


{\footnotesize\bfseries
No: 4\\
Name: at\_\allowbreak{}2\_\allowbreak{}odp\_\allowbreak{}04\\
Type: string\\
Default: nil
}

{\footnotesize Визначено події, які потребують навчання користувачів системи грамотності з питань конфіденційності}


{\footnotesize\bfseries
No: 5\\
Name: at\_\allowbreak{}2\_\allowbreak{}odp\_\allowbreak{}05\\
Type: string\\
Default: nil
}

{\footnotesize Визначено методи, які слід застосовувати для підвищення обізнаності користувачів системи щодо безпеки та конфіденційністі}


{\footnotesize\bfseries
No: 6\\
Name: at\_\allowbreak{}2\_\allowbreak{}odp\_\allowbreak{}06\\
Type: string\\
Default: nil
}

{\footnotesize Визначено частоту оновлення навчання грамотності та змісту обізнаності}


{\footnotesize\bfseries
No: 7\\
Name: at\_\allowbreak{}2\_\allowbreak{}odp\_\allowbreak{}07\\
Type: string\\
Default: nil
}

{\footnotesize Визначено події після яких необхідне оновлення навчання грамотності та змісту обізнаності}


{\footnotesize\bfseries
No: 8\\
Name: at\_\allowbreak{}2\_\allowbreak{}a\_\allowbreak{}01\_\allowbreak{}01\\
Type: string\\
Default: nil
}

{\footnotesize Навчання з грамотності з питань безпеки надається користувачам системи (включаючи менеджерів, керівників вищої ланки та підрядників) як частина початкового навчання для нових користувачів}


{\footnotesize\bfseries
No: 9\\
Name: at\_\allowbreak{}2\_\allowbreak{}a\_\allowbreak{}01\_\allowbreak{}02\\
Type: string\\
Default: nil
}

{\footnotesize Навчання з грамотності з питань конфіденційності надається користувачам системи (включаючи менеджерів, керівників вищої ланки та підрядників) як частина початкового навчання для нових користувачів}


{\footnotesize\bfseries
No: 10\\
Name: at\_\allowbreak{}2\_\allowbreak{}a\_\allowbreak{}01\_\allowbreak{}03\\
Type: string\\
Default: nil
}

{\footnotesize Для користувачів системи (включно з менеджерами, вищим керівництвом та підрядниками) проводиться навчання з безпеки з <AT-02\_\allowbreak{}ODP[01] періодичністю> після цього}


{\footnotesize\bfseries
No: 11\\
Name: at\_\allowbreak{}2\_\allowbreak{}a\_\allowbreak{}01\_\allowbreak{}04\\
Type: string\\
Default: nil
}

{\footnotesize Для користувачів системи (включно з менеджерами, вищим керівництвом та підрядниками) проводиться навчання з конфіденційності з <AT-02\_\allowbreak{}ODP[02] періодичністю> після цього}


{\footnotesize\bfseries
No: 12\\
Name: at\_\allowbreak{}2\_\allowbreak{}a\_\allowbreak{}02\_\allowbreak{}01\\
Type: string\\
Default: nil
}

{\footnotesize Тренінги з грамотності з питань безпеки проводяться для користувачів системи (включаючи менеджерів, керівників вищої ланки та підрядників), коли цього вимагають зміни в системі або після <AT-02\_\allowbreak{}ODP[03] подій>}


{\footnotesize\bfseries
No: 13\\
Name: at\_\allowbreak{}2\_\allowbreak{}a\_\allowbreak{}02\_\allowbreak{}02\\
Type: string\\
Default: nil
}

{\footnotesize Тренінги з грамотності з питань конфіденціності проводяться для користувачів системи (включаючи менеджерів, керівників вищої ланки та підрядників), коли цього вимагають зміни в системі або після <AT-02\_\allowbreak{}ODP[04] подій>}


{\footnotesize\bfseries
No: 14\\
Name: at\_\allowbreak{}2\_\allowbreak{}b\\
Type: string\\
Default: nil
}

{\footnotesize <AT-02\_\allowbreak{}ODP[05] методи> застосовуються для підвищення обізнаності користувачів системи щодо безпеки та конфіденційності}


{\footnotesize\bfseries
No: 15\\
Name: at\_\allowbreak{}2\_\allowbreak{}c\_\allowbreak{}01\\
Type: string\\
Default: nil
}

{\footnotesize Оновлюється зміст навчання грамотності та підвищення обізнаності з <AT-02\_\allowbreak{}ODP[06] частотою>}


{\footnotesize\bfseries
No: 16\\
Name: at\_\allowbreak{}2\_\allowbreak{}c\_\allowbreak{}02\\
Type: string\\
Default: nil
}

{\footnotesize Оновлюється зміст навчання грамотності та підвищення обізнаності після <AT-02\_\allowbreak{}ODP[07] подій>}


{\footnotesize\bfseries
No: 17\\
Name: at\_\allowbreak{}2\_\allowbreak{}d\\
Type: string\\
Default: nil
}

{\footnotesize Уроки, отримані в результаті внутрішніх або зовнішніх інцидентів або порушень безпеки, включені в методи навчання та підвищення обізнаності}




\subsubsection{ПРАКТИЧНІ ЗАНЯТТЯ (AT-2(1))}


\vspace{10pt}
{\footnotesize\bfseries
No: 1\\
Name: at\_\allowbreak{}2\_\allowbreak{}1\_\allowbreak{}01\\
Type: string\\
Default: nil
}

{\footnotesize Передбачені практичні вправи з тренування обізнаності, які імітують інциденти в області безпеки та конфіденційності}




\subsubsection{ВНУТРІШНІ ЗАГРОЗИ (AT-2(2))}
Ввести до програми навчання вправи з розпізнавання та виявлення потенційних індикаторів внутрішніх загроз.

\vspace{10pt}
{\footnotesize\bfseries
No: 1\\
Name: at\_\allowbreak{}2\_\allowbreak{}2\_\allowbreak{}01\\
Type: string\\
Default: nil
}

{\footnotesize Введено до програми навчання вправи з розпізнавання потенційних індикаторів внутрішніх загроз}


{\footnotesize\bfseries
No: 2\\
Name: at\_\allowbreak{}2\_\allowbreak{}2\_\allowbreak{}02\\
Type: string\\
Default: nil
}

{\footnotesize Введено до програми навчання вправи з виявлення потенційних індикаторів внутрішніх загроз}




\subsubsection{СОЦІАЛЬНА ІНЖЕНЕРІЯ ТА СОЦІАЛЬНИЙ ІНТЕЛЕКТУАЛЬНИЙ АНАЛІЗ ДАНИХ (AT-2(3))}
Ввести до програми навчання вправи з підвищення обізнаності щодо розпізнавання та повідомлення про потенційні та фактичні атаки, з використанням методів соціальної інженерії та інтелектуального аналізу соціальних даних.

\vspace{10pt}
{\footnotesize\bfseries
No: 1\\
Name: at\_\allowbreak{}2\_\allowbreak{}3\_\allowbreak{}01\\
Type: string\\
Default: nil
}

{\footnotesize До програми навчання введено вправи з розпізнавання потенційних та фактичних випадків соціального інжинірингу}


{\footnotesize\bfseries
No: 2\\
Name: at\_\allowbreak{}2\_\allowbreak{}3\_\allowbreak{}02\\
Type: string\\
Default: nil
}

{\footnotesize До програми навчання введено вправи з повідомлення про потенційні та фактичні випадки соціального інжинірингу}


{\footnotesize\bfseries
No: 3\\
Name: at\_\allowbreak{}2\_\allowbreak{}3\_\allowbreak{}03\\
Type: string\\
Default: nil
}

{\footnotesize До програми навчання введено вправи з розпізнавання потенційних та фактичних випадків інтелектуального аналізу соціальних даних}


{\footnotesize\bfseries
No: 4\\
Name: at\_\allowbreak{}2\_\allowbreak{}3\_\allowbreak{}04\\
Type: string\\
Default: nil
}

{\footnotesize До програми навчання введено вправи з повідомлення про потенційні та фактичні випадки інтелектуального аналізу соціальних даних}




\subsubsection{ПІДОЗРІЛІ ПОВІДОМЛЕННЯ ТА АНОМАЛЬНА ПОВЕДІНКА СИСТЕМИ (AT-2(4))}


\vspace{10pt}
{\footnotesize\bfseries
No: 1\\
Name: at\_\allowbreak{}2\_\allowbreak{}4\_\allowbreak{}odp\\
Type: string\\
Default: nil
}

{\footnotesize Визначено індикатори шкідливого коду}


{\footnotesize\bfseries
No: 2\\
Name: at\_\allowbreak{}2\_\allowbreak{}4\_\allowbreak{}01\\
Type: string\\
Default: nil
}

{\footnotesize Навчання грамотності щодо розпізнавання підозрілих повідомлень та аномальної поведінки у системах організації проводиться з використанням <AT-02(04)\_\allowbreak{}ODP індикаторів>}




\subsubsection{ВДОСКОНАЛЕНА СТІЙКА ЗАГРОЗА (AT-2(5))}


\vspace{10pt}
{\footnotesize\bfseries
No: 1\\
Name: at\_\allowbreak{}2\_\allowbreak{}5\_\allowbreak{}01\\
Type: string\\
Default: nil
}

{\footnotesize Забезпечено навчання грамотності щодо стійкої постійної загрози}




\subsubsection{СЕРЕДОВИЩЕ КІБЕРЗАГРОЗ (AT-2(6))}


\vspace{10pt}
{\footnotesize\bfseries
No: 1\\
Name: at\_\allowbreak{}2\_\allowbreak{}6\_\allowbreak{}a\\
Type: string\\
Default: nil
}

{\footnotesize Забезпечено навчання грамотності щодо середовища кіберзагроз}


{\footnotesize\bfseries
No: 2\\
Name: at\_\allowbreak{}2\_\allowbreak{}6\_\allowbreak{}b\\
Type: string\\
Default: nil
}

{\footnotesize Відображається поточна інформація про кіберзагрози в операціях системи}




\subsection{РОЛЬОВЕ НАВЧАННЯ (AT-3)}
a. Забезпечити проведення навчання з питань безпеки та приватності на основі ролей для працівників з ролями та обов'язками: [Призначення: визначені організацією ролі та обов'язки]:\\ 
1. перед авторизацією доступу до системи, інформації або виконанням призначених обов'язків і [Призначення: частота, визначена організацією] після цього;\\ 
2. коли цього потребують системні зміни.\\ 
b. Оновити навчальний контент на основі ролей [Призначення: частота, визначена організацією] і наступні [Призначення: події, визначені організацією];\\ 
c. Включіть у рольове навчання, інформацію, отриману з внутрішніх або зовнішніх інцидентів та порушень безпеки.

\vspace{10pt}
{\footnotesize\bfseries
No: 1\\
Name: at\_\allowbreak{}3\_\allowbreak{}odp\_\allowbreak{}01\\
Type: list\\
Default: [\textquotedbl{}admin\textquotedbl{}, \textquotedbl{}security\_\allowbreak{}officer\textquotedbl{}]
}

{\footnotesize Визначено ролі та обов'язки для тренінгів з безпеки на основі ролей}


{\footnotesize\bfseries
No: 2\\
Name: at\_\allowbreak{}3\_\allowbreak{}odp\_\allowbreak{}02\\
Type: list\\
Default: [\textquotedbl{}admin\textquotedbl{}, \textquotedbl{}security\_\allowbreak{}officer\textquotedbl{}]
}

{\footnotesize Визначено ролі та обов'язки для тренінгів з конфіденційності на основі ролей}


{\footnotesize\bfseries
No: 3\\
Name: at\_\allowbreak{}3\_\allowbreak{}odp\_\allowbreak{}03\\
Type: string\\
Default: nil
}

{\footnotesize Визначено частоту проведення тренінгів на основі ролей з безпеки та конфіденційності для призначеного персоналу після початкової підготовки}


{\footnotesize\bfseries
No: 4\\
Name: at\_\allowbreak{}3\_\allowbreak{}odp\_\allowbreak{}04\\
Type: integer\\
Default: 30
}

{\footnotesize Визначено частоту оновлення змісту навчання на основі ролей}


{\footnotesize\bfseries
No: 5\\
Name: at\_\allowbreak{}3\_\allowbreak{}odp\_\allowbreak{}05\\
Type: string\\
Default: nil
}

{\footnotesize Визначено події, які потребують оновлення змісту навчання на основі ролей}


{\footnotesize\bfseries
No: 6\\
Name: at\_\allowbreak{}3\_\allowbreak{}a\_\allowbreak{}01\_\allowbreak{}01\\
Type: string\\
Default: nil
}

{\footnotesize Навчання з безпеки на основі ролей проводиться для <AT-03\_\allowbreak{}ODP[01] ролей та обов'язків> перед авторизацією доступу до системи, інформації або виконанням призначених обов'язків}


{\footnotesize\bfseries
No: 7\\
Name: at\_\allowbreak{}3\_\allowbreak{}a\_\allowbreak{}01\_\allowbreak{}02\\
Type: string\\
Default: nil
}

{\footnotesize Навчання з конфіденційності на основі ролей проводиться для <AT-03\_\allowbreak{}ODP[02] ролей та обов'язків> перед авторизацією доступу до системи, інформації або виконанням призначених обов'язків}


{\footnotesize\bfseries
No: 8\\
Name: at\_\allowbreak{}3\_\allowbreak{}a\_\allowbreak{}01\_\allowbreak{}03\\
Type: string\\
Default: nil
}

{\footnotesize Навчання з безпеки на основі ролей проводиться для <AT-03\_\allowbreak{}ODP[01] ролей та обов'язків> з <AT-03\_\allowbreak{}ODP[03] частотою> після цього}


{\footnotesize\bfseries
No: 9\\
Name: at\_\allowbreak{}3\_\allowbreak{}a\_\allowbreak{}01\_\allowbreak{}04\\
Type: string\\
Default: nil
}

{\footnotesize Навчання з конфіденційності на основі ролей проводиться для <AT-03\_\allowbreak{}ODP[02] ролей та обов'язків> з <AT-03\_\allowbreak{}ODP[03] частотою> після цього}


{\footnotesize\bfseries
No: 10\\
Name: at\_\allowbreak{}3\_\allowbreak{}a\_\allowbreak{}02\_\allowbreak{}01\\
Type: string\\
Default: nil
}

{\footnotesize Навчання з питань безпеки на основі ролей проводиться для персоналу, який виконує певні ролі та обов'язки у сфері безпеки, коли цього вимагають зміни в системі}


{\footnotesize\bfseries
No: 11\\
Name: at\_\allowbreak{}3\_\allowbreak{}a\_\allowbreak{}02\_\allowbreak{}02\\
Type: string\\
Default: nil
}

{\footnotesize Навчання з питань конфіденційності на основі ролей проводиться для персоналу, який виконує певні ролі та обов'язки у сфері безпеки, коли цього вимагають зміни в системі}


{\footnotesize\bfseries
No: 12\\
Name: at\_\allowbreak{}3\_\allowbreak{}b\_\allowbreak{}01\\
Type: string\\
Default: nil
}

{\footnotesize Оновлюється вміст навчання на основі ролей з <AT-03\_\allowbreak{}ODP[04] частотою>}


{\footnotesize\bfseries
No: 13\\
Name: at\_\allowbreak{}3\_\allowbreak{}b\_\allowbreak{}02\\
Type: string\\
Default: nil
}

{\footnotesize Оновлюється вміст навчання на основі ролей після <AT-03\_\allowbreak{}ODP[05] подій>}


{\footnotesize\bfseries
No: 14\\
Name: at\_\allowbreak{}3\_\allowbreak{}c\\
Type: string\\
Default: nil
}

{\footnotesize Інформація отримана з внутрішніх чи зовнішніх інцидентів або порушень безпеки, включається в навчання на основі ролей}




\subsubsection{ЗАХОДИ БЕЗПЕКИ РОБОЧОГО СЕРЕДОВИЩА (AT-3(1))}


\vspace{10pt}
{\footnotesize\bfseries
No: 1\\
Name: at\_\allowbreak{}3\_\allowbreak{}1\_\allowbreak{}odp\_\allowbreak{}01\\
Type: string\\
Default: nil
}

{\footnotesize Визначено персонал або ролі, які мають бути забезпечені початковим навчанням та підвищенням кваліфікації з питань застосування та експлуатації заходів захисту робочого середовища}


{\footnotesize\bfseries
No: 2\\
Name: at\_\allowbreak{}3\_\allowbreak{}1\_\allowbreak{}odp\_\allowbreak{}02\\
Type: string\\
Default: nil
}

{\footnotesize Визначено частоту проведення підвищення кваліфікації в галузі операцій та функціонування заходів захисту робочого середовища}


{\footnotesize\bfseries
No: 3\\
Name: at\_\allowbreak{}3\_\allowbreak{}1\_\allowbreak{}01\\
Type: string\\
Default: nil
}

{\footnotesize <AT-03(01)\_\allowbreak{}ODP[01] персонал або ролі> забезпечується підвищенням кваліфікації в галузі зайнятості та функціонування заходів захисту робочого середовища з <AT-03(01)\_\allowbreak{}ODP[02] частотою>}




\subsubsection{ФІЗИЧНІ ЗАХОДИ БЕЗПЕКИ (AT-3(2))}


\vspace{10pt}
{\footnotesize\bfseries
No: 1\\
Name: at\_\allowbreak{}3\_\allowbreak{}2\_\allowbreak{}odp\_\allowbreak{}01\\
Type: string\\
Default: nil
}

{\footnotesize Визначаено персонал або ролі, які мають бути забезпечені підготовкою з питань застосування та експлуатації заходів фізичної безпеки}


{\footnotesize\bfseries
No: 2\\
Name: at\_\allowbreak{}3\_\allowbreak{}2\_\allowbreak{}odp\_\allowbreak{}02\\
Type: string\\
Default: nil
}

{\footnotesize Визначаеночастоту проведення підготовки з питань застосування та експлуатації заходів фізичної безпеки}


{\footnotesize\bfseries
No: 3\\
Name: at\_\allowbreak{}3\_\allowbreak{}2\_\allowbreak{}01\\
Type: string\\
Default: nil
}

{\footnotesize <AT-03(02)\_\allowbreak{}ODP[01] персонал або ролі> забезпечуються підготовкою з питань застосування та експлуатації заходів фізичної безпеки з <AT-03(02)\_\allowbreak{}ODP[02] частотою>}




\subsubsection{ПРАКТИЧНІ ЗАНЯТТЯ (AT-3(3))}


\vspace{10pt}
{\footnotesize\bfseries
No: 1\\
Name: at\_\allowbreak{}3\_\allowbreak{}3\_\allowbreak{}01\\
Type: string\\
Default: nil
}

{\footnotesize Програма навчання включає практичні заняття з безпеки, які мають підкріпити досягнення цілей навчання}


{\footnotesize\bfseries
No: 2\\
Name: at\_\allowbreak{}3\_\allowbreak{}3\_\allowbreak{}02\\
Type: string\\
Default: nil
}

{\footnotesize Програма навчання включає практичні заняття з конфіденційності, які мають підкріпити досягнення цілей навчання}




\subsubsection{ПІДОЗРІЛІ ЗВ'ЯЗКИ ТА АНОМАЛЬНА ПОВЕДІНКА СИСТЕМИ (AT-3(4))}
[Вилучено: включено до АТ-02(04)].

\vspace{10pt}
\textit{Немає параметрів для цього контролю.}
\vspace{10pt}


\subsubsection{ОБРОБКА ПЕРСОНАЛЬНИХ ДАНИХ (AT-3(5))}


\vspace{10pt}
{\footnotesize\bfseries
No: 1\\
Name: at\_\allowbreak{}3\_\allowbreak{}5\_\allowbreak{}odp\_\allowbreak{}01\\
Type: string\\
Default: nil
}

{\footnotesize Визначено персонал або посади, які мають пройти навчання з використання та управління обробкою персональних даних та контролю прозорості}


{\footnotesize\bfseries
No: 2\\
Name: at\_\allowbreak{}3\_\allowbreak{}5\_\allowbreak{}odp\_\allowbreak{}02\\
Type: string\\
Default: nil
}

{\footnotesize Визначено періодичність проведення навчання з використання та управління обробкою персональних даних та контролю прозорості}


{\footnotesize\bfseries
No: 3\\
Name: at\_\allowbreak{}3\_\allowbreak{}5\_\allowbreak{}01\\
Type: string\\
Default: nil
}

{\footnotesize <AT-03(05)\_\allowbreak{}ODP[01] персонал або ролі> навчання з <AT-03(05)\_\allowbreak{}ODP[02] частотою> з використання та управління обробкою персональних даних та контролю прозорості}




\subsection{НАВЧАЛЬНІ ЗАПИСИ (AT-4)}
a. Документувати та відстежувати індивідуальні навчальні заходи із забезпечення безпеки та приватності, включно з базовою підготовкою з питань безпеки та приватності, а також спеціальною підготовкою з питань безпеки та приватності визначених посадових осіб.\\ 
b. Зберігати індивідуальні записи про навчання впродовж [Призначення: визначеного організацією періоду часу].

\vspace{10pt}
{\footnotesize\bfseries
No: 1\\
Name: at\_\allowbreak{}4\_\allowbreak{}a\_\allowbreak{}01\\
Type: string\\
Default: nil
}

{\footnotesize Задокументовані індивідуальні навчальні заходи із забезпечення безпеки та конфіденційності інформації, включно з базовою підготовкою з питань безпеки та конфіденційності, а також спеціальною підготовкою з безпеки та конфіденційності}


{\footnotesize\bfseries
No: 2\\
Name: at\_\allowbreak{}4\_\allowbreak{}a\_\allowbreak{}02\\
Type: string\\
Default: nil
}

{\footnotesize Відстежуються індивідуальні навчальні заходи із забезпечення безпеки та конфіденційності інформації, включно з базовою підготовкою з питань безпеки та конфіденційності, а також спеціальною підготовкою з безпеки та конфіденційності}


{\footnotesize\bfseries
No: 3\\
Name: at\_\allowbreak{}4\_\allowbreak{}b\\
Type: integer\\
Default: 30
}

{\footnotesize Індивідуальні записи про навчання зберігаються протягом період часу}


{\footnotesize\bfseries
No: 4\\
Name: at\_\allowbreak{}4\_\allowbreak{}odp\\
Type: integer\\
Default: 30
}

{\footnotesize Визначено період зберігання індивідуальних записів про навчання}




\subsection{КОНТАКТИ З ГРУПАМИ БЕЗПЕКИ ТА АСОЦІАЦІЯМИ (AT-5) [Вилучено]}
[Вилучено: включено в PM-15]

\vspace{10pt}
\textit{Немає параметрів для цього контролю.}
\vspace{10pt}


\subsection{ВІДГУКИ ПРО ПРОВЕДЕНІ НАВЧАННЯ (AT-6)}
Надати відгук про результати організаційного навчання наступному персоналу [Призначення: з визначеною організацією частотою та визначеному організацією персоналу]

\vspace{10pt}
{\footnotesize\bfseries
No: 1\\
Name: at\_\allowbreak{}6\_\allowbreak{}01\\
Type: list\\
Default: [\textquotedbl{}admin\textquotedbl{}, \textquotedbl{}security\_\allowbreak{}officer\textquotedbl{}]
}

{\footnotesize Відгуки про результати навчання надаютсья з визначеною частотою до персоналу}


{\footnotesize\bfseries
No: 2\\
Name: at\_\allowbreak{}6\_\allowbreak{}odp\_\allowbreak{}01\\
Type: integer\\
Default: 30
}

{\footnotesize Визначено частоту надання відгуків щодо результатів навчання в організації}


{\footnotesize\bfseries
No: 3\\
Name: at\_\allowbreak{}6\_\allowbreak{}odp\_\allowbreak{}02\\
Type: list\\
Default: [\textquotedbl{}admin\textquotedbl{}, \textquotedbl{}security\_\allowbreak{}officer\textquotedbl{}]
}

{\footnotesize Призначено персонал, якому надаватиметься відгуки щодо результатів навчання в організації}




\section{AU}
Клас заходів захисту AU --- АУДИТ ТА ПІДЗВІТНІСТЬ

Цей клас забезпечує можливість відстеження дій у системі, генерування, захист та аналіз записів аудиту для виявлення порушень політики безпеки.

\textbf{Перелік заходів захисту:} Політика та процедури аудиту та підзвітності (AU-1); Події аудиту (AU-2); Узагальнення записів про аудит з декількох джерел (AU-2(1)) [Вилучено]; Вибір події аудиту за компонентами (AU-2(2)) [Вилучено]; Перегляд та оновлення (AU-2(3)) [Вилучено]; Привілейовані функції (AU-2(4)) [Вилучено]; Зміст записів аудиту (AU-3); Додаткова інформація про аудит (AU-3(1)); Централізоване управління планованим змістом записів аудиту (AU-3(2)) [Вилучено]; Обмеження елементів персональних даних (AU-3(3)); Місткість сховища записів аудиту (AU-4); Передача до альтернативного сховища (AU-4(1)); Реагування на відмови обробки даних аудиту (AU-5); Місткість сховища записів аудиту (AU-5(1)); Тривожне сповіщення в реальному часі (AU-5(2)); Налаштування порогового обсягу трафіку (AU-5(3)); Вимкнення у разі відмови (AU-5(4)); Можливість альтернативного журналювання аудиту (AU-5(5)); Огляд, аналіз і звітність аудиту (AU-6); Автоматизована інтеграція процесів (AU-6(1)); Огляд, аналіз і звітність аудиту сповіщення про порушення безпеку (AU-6(2)) [Вилучено]; Зіставляння сховищ аудиту (AU-6(3)); Централізований перегляд та аналіз (AU-6(4)); Інтегрований аналіз записів аудиту (AU-6(5)); Кореляція з фізичним моніторингом (AU-6(6)); Дозволені дії (AU-6(7)); Аналіз повного тексту привілейованих команд (AU-6(8)); Кореляція з інформацією з нетехнічних джерел (AU-6(9)); Регулювання рівня аудиту (AU-6(10)) [Вилучено]; Скорочення записів аудиту та формування звіту (AU-7); Автоматична обробка (AU-7(1)); Автоматичне сортування та пошук (AU-7(2)) [Вилучено]; Позначка часу (AU-8); Синхронізація з авторитетним джерелом часу (AU-8(1)) [Вилучено]; Вторинне авторитетне джерело часу (AU-8(2)) [Вилучено]; Захист інформації аудиту (AU-9); Апаратні носії інформації одноразового запису (AU-9(1)); Зберігання на окремих фізичних системах або компонентах (AU-9(2)); Криптографічний захист (AU-9(3)); Доступ, який надається через членство в підмножини привілейованих користувачів (AU-9(4)); Подвійна авторизація (AU-9(5)); Доступ тільки для читання (AU-9(6)); Зберігання на компоненті іншої операційної системи (AU-9(7)); Неспростовність (AU-10); Асоціація ідентичності (AU-10(1)); Ратифікація прив'язки інформації про ідентичність виробника (AU-10(2)); Ланцюжок збереження доказів (AU-10(3)); Валідація зв'язку ідентичності перегля- (AU-10(4)); Цифрові підписи (AU-10(5)) [Вилучено]; Збереження записів аудиту (AU-11); Довгострокова можливість отримання (AU-11(1)); Генерація даних аудиту (AU-12); Загальносистемний та синхронізований за часом журналу аудиту (AU-12(1)); Стандартизовані формати (AU-12(2)); Зміни, що вносять авторизовані особи (AU-12(3)); Аудит запитів персональних даних (AU-12(4)); Моніторинг розкриття інформації (AU-13); Використання автоматичних засобів (AU-13(1)); Огляд сайтів, що підлягають моніторингу (AU-13(2)); Авторизоване копіювання інформації (AU-13(3)); Аудит сесії (AU-14); Система запуску (AU-14(1)); Захоплення та запис інформації (AU-14(2)) [Вилучено]; Віддалений перегляд та прослуховування (AU-14(3)); Альтернативна можливість аудиту (AU-15) [Вилучено]; Міжорганізаційний аудит (AU-16); Збереження ідентичності (AU-16(1)); Обмін інформацією аудиту (AU-16(2)); Розмежування (AU-16(3)).

\subsection{ПОЛІТИКА ТА ПРОЦЕДУРИ АУДИТУ ТА ПІДЗВІТНОСТІ (AU-1)}
a. Розробити, задокументувати та поширити [Призначення: серед персоналу або ролей, що їх визначила організація]:\\ 
1. [Вибір (один або декілька): Рівень організації; Рівень місії/бізнес-процесу; рівень системи] політика аудиту та підзвітності, яка:\\ 
(a) містить мету, сферу застосування, ролі, обов'язки, відповідальність керівництва, координацію між організаційними підрозділами та систему контролю відповідності (complaince);\\ 
(b) відповідає чинним законам, нормативним документам, директивам, нормам, політикам, стандартам та керівним документам.\\ 
2. Процедури, що сприяють здійсненню політики аудиту та підзвітності, а також пов'язані з ними заходи аудиту та підзвітності.\\ 
b. Призначити [Призначення: визначену організацією старшу посадову особу] для управління політикою та процедурами аудиту та підзвітності.\\ 
c. Переглядати та оновлювати поточний аудит та підзвітність:\\ 
1. політику [Призначення: частота, визначена організацією] та наступне [Призначення: події, визначені організацією];\\ 
2. процедури аудиту [Призначення: визначеною організацією частотою] та [Завдання: події, визначені організацією].

\vspace{10pt}
{\footnotesize\bfseries
No: 1\\
Name: au\_\allowbreak{}1\_\allowbreak{}a\_\allowbreak{}01\\
Type: list\\
Default: [\textquotedbl{}default\_\allowbreak{}deny\_\allowbreak{}rule\textquotedbl{}, \textquotedbl{}abac\_\allowbreak{}rule\_\allowbreak{}1\textquotedbl{}]
}

{\footnotesize Розроблено та задокументовано політику аудиту та підзвітності}


{\footnotesize\bfseries
No: 2\\
Name: au\_\allowbreak{}1\_\allowbreak{}a\_\allowbreak{}02\\
Type: list\\
Default: [\textquotedbl{}admin\textquotedbl{}, \textquotedbl{}security\_\allowbreak{}officer\textquotedbl{}]
}

{\footnotesize Політика аудиту та підзвітності доведена до персонал або ролі}


{\footnotesize\bfseries
No: 3\\
Name: au\_\allowbreak{}1\_\allowbreak{}a\_\allowbreak{}03\\
Type: list\\
Default: [\textquotedbl{}default\_\allowbreak{}deny\_\allowbreak{}rule\textquotedbl{}, \textquotedbl{}abac\_\allowbreak{}rule\_\allowbreak{}1\textquotedbl{}]
}

{\footnotesize Розроблені та задокументовані процедури аудиту та підзвітності, що сприяють впровадженню політики аудиту та підзвітності, а також відповідні заходи контролю аудиту та підзвітності}


{\footnotesize\bfseries
No: 4\\
Name: au\_\allowbreak{}1\_\allowbreak{}b\\
Type: list\\
Default: [\textquotedbl{}admin\textquotedbl{}, \textquotedbl{}security\_\allowbreak{}officer\textquotedbl{}]
}

{\footnotesize Посадова особа призначається для управління політикою та процедурами аудиту та підзвітності AU-01(c)[01][01] переглядається та оновлюється поточна політика аудиту та підзвітності з частота; AU-01(c)[01][02] переглядається та оновлюється поточна політика аудиту та підзвітності після подій; AU-01(c)[02][01] переглядається та оновлюється поточні процедури аудиту та підзвітності з частота; AU-01(c)[02][02] переглядається та оновлюється поточні процедури аудиту та підзвітності після подій}


{\footnotesize\bfseries
No: 5\\
Name: au\_\allowbreak{}1\_\allowbreak{}odp\_\allowbreak{}01\\
Type: list\\
Default: [\textquotedbl{}admin\textquotedbl{}, \textquotedbl{}security\_\allowbreak{}officer\textquotedbl{}]
}

{\footnotesize Визначено персонал або ролі, до яких має бути доведена політика аудиту та підзвітності}


{\footnotesize\bfseries
No: 6\\
Name: au\_\allowbreak{}1\_\allowbreak{}odp\_\allowbreak{}02\\
Type: list\\
Default: [\textquotedbl{}admin\textquotedbl{}, \textquotedbl{}security\_\allowbreak{}officer\textquotedbl{}]
}

{\footnotesize Визначено персонал або ролі, на які поширюються процедури аудиту та підзвітності}


{\footnotesize\bfseries
No: 7\\
Name: au\_\allowbreak{}1\_\allowbreak{}odp\_\allowbreak{}03\\
Type: string\\
Default: nil
}

{\footnotesize Вибрано одне або декілька з наступних ЗНАЧЕНЬ ПАРАМЕТРІВ: \{рівень організації; рівень місії/бізнес-процесу; рівень системи\}}


{\footnotesize\bfseries
No: 8\\
Name: au\_\allowbreak{}1\_\allowbreak{}odp\_\allowbreak{}04\\
Type: list\\
Default: [\textquotedbl{}admin\textquotedbl{}, \textquotedbl{}security\_\allowbreak{}officer\textquotedbl{}]
}

{\footnotesize Визначено посадову особу, яка управлятиме політикою та процедурами аудиту та підзвітності}


{\footnotesize\bfseries
No: 9\\
Name: au\_\allowbreak{}1\_\allowbreak{}odp\_\allowbreak{}05\\
Type: list\\
Default: [\textquotedbl{}default\_\allowbreak{}deny\_\allowbreak{}rule\textquotedbl{}, \textquotedbl{}abac\_\allowbreak{}rule\_\allowbreak{}1\textquotedbl{}]
}

{\footnotesize Визначено частоту, з якою переглядається та оновлюється поточна політика аудиту та підзвітності}


{\footnotesize\bfseries
No: 10\\
Name: au\_\allowbreak{}1\_\allowbreak{}odp\_\allowbreak{}06\\
Type: list\\
Default: [\textquotedbl{}default\_\allowbreak{}deny\_\allowbreak{}rule\textquotedbl{}, \textquotedbl{}abac\_\allowbreak{}rule\_\allowbreak{}1\textquotedbl{}]
}

{\footnotesize Визначено події, які потребують перегляду та оновлення поточної політики аудиту та підзвітності}


{\footnotesize\bfseries
No: 11\\
Name: au\_\allowbreak{}1\_\allowbreak{}odp\_\allowbreak{}07\\
Type: integer\\
Default: 30
}

{\footnotesize Визначено частоту, з якою переглядаються та оновлюються поточні процедури аудиту та підзвітності}


{\footnotesize\bfseries
No: 12\\
Name: au\_\allowbreak{}1\_\allowbreak{}odp\_\allowbreak{}08\\
Type: list\\
Default: [\textquotedbl{}login\textquotedbl{}, \textquotedbl{}logout\textquotedbl{}, \textquotedbl{}failed\_\allowbreak{}attempt\textquotedbl{}]
}

{\footnotesize Визначено події, які потребують перегляду та оновлення поточної процедури аудиту та підзвітності}




\subsection{ПОДІЇ АУДИТУ (AU-2)}
a. Визначити типи подій, які система може реєструвати для підтримки функції аудиту: [Призначення: типи подій, визначені організацією, які система здатна реєструвати];\\ 
b. Координувати функції аудиту безпеки з іншими організаційними підрозділами, які вимагають інформації, пов'язаної з аудитом, для посилення взаємної підтримки та допомоги у виборі типів подій, що перевіряються;\\ 
c. Визначити, які типи подій підлягають аудиту: [Призначення: визначені організацією події, що підлягають аудиту (підмножина подій, що підлягають аудиту, визначених в AU-2 a.), а також частота (або ситуація, що вимагає) проведення аудиту для кожної ідентифікованої події]\\ 
d. Обґрунтувати, чому типи подій, що перевіряються, вважаються достатніми для підтримки розслідувань інцидентів (постфактум), пов'язаних з безпекою та приватністю;\\ 
e. Перегляньте й оновіть типи подій, вибрані для журналювання [Призначення: частота, визначена організацією].

\vspace{10pt}
{\footnotesize\bfseries
No: 1\\
Name: au\_\allowbreak{}2\_\allowbreak{}a\\
Type: string\\
Default: nil
}

{\footnotesize Типи подій, які система здатна реєструвати, визначено для підтримки функції аудиту}


{\footnotesize\bfseries
No: 2\\
Name: au\_\allowbreak{}2\_\allowbreak{}b\\
Type: string\\
Default: nil
}

{\footnotesize Функція аудиту безпеки координується з іншими підрозділами організації, які вимагають інформації, пов'язаної з аудитом, длдля посилення взаємної підтримки та допомоги у виборі типів подій, що перевіряються}


{\footnotesize\bfseries
No: 3\\
Name: au\_\allowbreak{}2\_\allowbreak{}c\_\allowbreak{}01\\
Type: string\\
Default: nil
}

{\footnotesize Типи подій (підмножина 02\_\allowbreak{}ODP[01]) визначаються для реєстрації у системі; AU-}


{\footnotesize\bfseries
No: 4\\
Name: au\_\allowbreak{}2\_\allowbreak{}c\_\allowbreak{}02\\
Type: string\\
Default: \textquotedbl{}щорічно\textquotedbl{}
}

{\footnotesize Зазначені типи подій реєструються 02\_\allowbreak{}ODP[03] частота або ситуація>; <AU-}


{\footnotesize\bfseries
No: 5\\
Name: au\_\allowbreak{}2\_\allowbreak{}d\\
Type: string\\
Default: nil
}

{\footnotesize Надається обґрунтування, чому типи подій, що перевіряються, вважаються достатніми для підтримки розслідувань інцидентів (постфактум), пов'язаних з безпекою та конфіденційністю}


{\footnotesize\bfseries
No: 6\\
Name: au\_\allowbreak{}2\_\allowbreak{}e\\
Type: string\\
Default: \textquotedbl{}щорічно\textquotedbl{}
}

{\footnotesize Переглядаються та оновлюються типи подій, вибрані для реєстрації, частота. у системі}


{\footnotesize\bfseries
No: 7\\
Name: au\_\allowbreak{}2\_\allowbreak{}odp\_\allowbreak{}01\\
Type: string\\
Default: nil
}

{\footnotesize Визначено типи подій, які система може реєструвати для підтримки функції аудиту}


{\footnotesize\bfseries
No: 8\\
Name: au\_\allowbreak{}2\_\allowbreak{}odp\_\allowbreak{}02\\
Type: string\\
Default: nil
}

{\footnotesize Визначено типи подій (підмножина AU-02\_\allowbreak{}ODP[01]) що підлягають аудиту у системі}


{\footnotesize\bfseries
No: 9\\
Name: au\_\allowbreak{}2\_\allowbreak{}odp\_\allowbreak{}03\\
Type: list\\
Default: [\textquotedbl{}login\textquotedbl{}, \textquotedbl{}logout\textquotedbl{}, \textquotedbl{}failed\_\allowbreak{}attempt\textquotedbl{}]
}

{\footnotesize Визначено частоту або ситуацію, що вимагає проведення аудиту для кожної ідентифікованої події}


{\footnotesize\bfseries
No: 10\\
Name: au\_\allowbreak{}2\_\allowbreak{}odp\_\allowbreak{}04\\
Type: string\\
Default: \textquotedbl{}щорічно\textquotedbl{}
}

{\footnotesize Частота перегляду та оновлення типів подій, обраних для журналювання}




\subsubsection{УЗАГАЛЬНЕННЯ ЗАПИСІВ ПРО АУДИТ З ДЕКІЛЬКОХ ДЖЕРЕЛ (AU-2(1)) [Вилучено]}
[Вилучено: Включено в АU-12]

\vspace{10pt}
\textit{Немає параметрів для цього контролю.}
\vspace{10pt}


\subsubsection{ВИБІР ПОДІЇ АУДИТУ ЗА КОМПОНЕНТАМИ (AU-2(2)) [Вилучено]}
[Вилучено: Включено в АU-12]

\vspace{10pt}
\textit{Немає параметрів для цього контролю.}
\vspace{10pt}


\subsubsection{ПЕРЕГЛЯД ТА ОНОВЛЕННЯ (AU-2(3)) [Вилучено]}
[Вилучено: Включено в АU-02]

\vspace{10pt}
\textit{Немає параметрів для цього контролю.}
\vspace{10pt}


\subsubsection{ПРИВІЛЕЙОВАНІ ФУНКЦІЇ (AU-2(4)) [Вилучено]}
[Вилучено: Включено в AC-06(09)]

\vspace{10pt}
\textit{Немає параметрів для цього контролю.}
\vspace{10pt}


\subsection{ЗМІСТ ЗАПИСІВ АУДИТУ (AU-3)}
Переконатися, що записи аудиту містять інформацію, яка встановлює наступне:\\ 
a. який тип події стався;\\ 
b. коли відбулася подія;\\ 
c. де відбулася подія;\\ 
d. джерело події;\\ 
e. наслідки події;\\ 
f. результат події та ідентифікатор будь-яких осіб або суб'єктів, пов'язаних з подією.

\vspace{10pt}
{\footnotesize\bfseries
No: 1\\
Name: au\_\allowbreak{}3\_\allowbreak{}a\\
Type: list\\
Default: [\textquotedbl{}login\textquotedbl{}, \textquotedbl{}logout\textquotedbl{}, \textquotedbl{}failed\_\allowbreak{}attempt\textquotedbl{}]
}

{\footnotesize Записи аудиту містять інформацію, яка встановлює який тип події стався}


{\footnotesize\bfseries
No: 2\\
Name: au\_\allowbreak{}3\_\allowbreak{}b\\
Type: string\\
Default: nil
}

{\footnotesize Записи аудиту містять інформацію, яка встановлює коли подія сталася}


{\footnotesize\bfseries
No: 3\\
Name: au\_\allowbreak{}3\_\allowbreak{}c\\
Type: string\\
Default: nil
}

{\footnotesize Записи аудиту містять інформацію, яка встановлює де відбулася подія}


{\footnotesize\bfseries
No: 4\\
Name: au\_\allowbreak{}3\_\allowbreak{}d\\
Type: list\\
Default: [\textquotedbl{}login\textquotedbl{}, \textquotedbl{}logout\textquotedbl{}, \textquotedbl{}failed\_\allowbreak{}attempt\textquotedbl{}]
}

{\footnotesize Записи аудиту містять інформацію, яка встановлює джерело події}


{\footnotesize\bfseries
No: 5\\
Name: au\_\allowbreak{}3\_\allowbreak{}e\\
Type: list\\
Default: [\textquotedbl{}login\textquotedbl{}, \textquotedbl{}logout\textquotedbl{}, \textquotedbl{}failed\_\allowbreak{}attempt\textquotedbl{}]
}

{\footnotesize Записи аудиту містять інформацію, яка встановлює наслідки події}


{\footnotesize\bfseries
No: 6\\
Name: au\_\allowbreak{}3\_\allowbreak{}f\\
Type: list\\
Default: [\textquotedbl{}login\textquotedbl{}, \textquotedbl{}logout\textquotedbl{}, \textquotedbl{}failed\_\allowbreak{}attempt\textquotedbl{}]
}

{\footnotesize Записи аудиту містять інформацію, яка встановлює результат події та ідентифікатор будь-яких осіб або суб'єктів, пов'язаних з подією}




\subsubsection{ДОДАТКОВА ІНФОРМАЦІЯ ПРО АУДИТ (AU-3(1))}


\vspace{10pt}
{\footnotesize\bfseries
No: 1\\
Name: au\_\allowbreak{}3\_\allowbreak{}1\_\allowbreak{}01\\
Type: string\\
Default: nil
}

{\footnotesize Сформовані записи аудиту містять наступну <AU-03(01)\_\allowbreak{}ODP додаткова інформація>}


{\footnotesize\bfseries
No: 2\\
Name: au\_\allowbreak{}3\_\allowbreak{}1\_\allowbreak{}odp\\
Type: string\\
Default: nil
}

{\footnotesize Визначено додаткову інформацію, яка має бути включена до записів аудиту}




\subsubsection{ЦЕНТРАЛІЗОВАНЕ УПРАВЛІННЯ ПЛАНОВАНИМ ЗМІСТОМ ЗАПИСІВ АУДИТУ (AU-3(2)) [Вилучено]}
[Вилучено: Включено до PL-09]

\vspace{10pt}
\textit{Немає параметрів для цього контролю.}
\vspace{10pt}


\subsubsection{ОБМЕЖЕННЯ ЕЛЕМЕНТІВ ПЕРСОНАЛЬНИХ ДАНИХ (AU-3(3))}
Обмежити персональні дані, що міститься в записах аудиту, до таких елементів, які визначені в оцінці ризику приватності: [Призначення: визначені організацією елементи].

\vspace{10pt}
{\footnotesize\bfseries
No: 1\\
Name: au\_\allowbreak{}3\_\allowbreak{}3\_\allowbreak{}01\\
Type: list\\
Default: [\textquotedbl{}admin\textquotedbl{}, \textquotedbl{}security\_\allowbreak{}officer\textquotedbl{}]
}

{\footnotesize Інформація, щщо ідентифікує особу, яка міститься в записах аудиту, обмежується <AU-03(03)\_\allowbreak{}ODP елементами>, визначеними в оцінці ризиків конфіденційності}


{\footnotesize\bfseries
No: 2\\
Name: au\_\allowbreak{}3\_\allowbreak{}3\_\allowbreak{}odp\\
Type: string\\
Default: nil
}

{\footnotesize Визначаються елементи, визначені конфіденційності; в оцінці ризику}




\subsection{МІСТКІСТЬ СХОВИЩА ЗАПИСІВ АУДИТУ (AU-4)}
Розподіляти місткість сховища записів аудиту у відповідності до [Призначення: визначених організацією вимог до зберігання записів аудиту].

\vspace{10pt}
{\footnotesize\bfseries
No: 1\\
Name: au\_\allowbreak{}4\_\allowbreak{}01\\
Type: string\\
Default: nil
}

{\footnotesize Розподілено ємність для зберігання записів аудиту відповідно до вимог}


{\footnotesize\bfseries
No: 2\\
Name: au\_\allowbreak{}4\_\allowbreak{}odp\\
Type: string\\
Default: nil
}

{\footnotesize Визначено вимоги до зберігання записів аудиту}




\subsubsection{ПЕРЕДАЧА ДО АЛЬТЕРНАТИВНОГО СХОВИЩА (AU-4(1))}


\vspace{10pt}
{\footnotesize\bfseries
No: 1\\
Name: au\_\allowbreak{}4\_\allowbreak{}1\_\allowbreak{}01\\
Type: integer\\
Default: 30
}

{\footnotesize Записи аудиту вивантажуються на іншу систему або носій інформації, з системи, що перевіряється, з частотою}


{\footnotesize\bfseries
No: 2\\
Name: au\_\allowbreak{}4\_\allowbreak{}1\_\allowbreak{}odp\\
Type: integer\\
Default: 30
}

{\footnotesize Визначено частоту завантаження записів аудиту на іншу систему чи носій інформації, з системи що перевіряється}




\subsection{РЕАГУВАННЯ НА ВІДМОВИ ОБРОБКИ ДАНИХ АУДИТУ (AU-5)}
a. Сповіщати [Призначення: визначені організацією персонал або посади] у разі збою обробки даних аудиту в [Призначення: визначений організацією період часу].\\ 
b. Виконати наступні додаткові дії: [Призначення: визначені організацією дії, які необхідно зробити].

\vspace{10pt}
{\footnotesize\bfseries
No: 1\\
Name: au\_\allowbreak{}5\_\allowbreak{}a\\
Type: list\\
Default: [\textquotedbl{}admin\textquotedbl{}, \textquotedbl{}security\_\allowbreak{}officer\textquotedbl{}]
}

{\footnotesize Персонал або ролі отримують сповіщення у разі збою процесу обробки даних аудиту періоду часу}


{\footnotesize\bfseries
No: 2\\
Name: au\_\allowbreak{}5\_\allowbreak{}b\\
Type: list\\
Default: [\textquotedbl{}login\textquotedbl{}, \textquotedbl{}logout\textquotedbl{}, \textquotedbl{}failed\_\allowbreak{}attempt\textquotedbl{}]
}

{\footnotesize Додаткові дії виконуються у разі збою процесу обробки даних аудиту}


{\footnotesize\bfseries
No: 3\\
Name: au\_\allowbreak{}5\_\allowbreak{}odp\_\allowbreak{}01\\
Type: list\\
Default: [\textquotedbl{}admin\textquotedbl{}, \textquotedbl{}security\_\allowbreak{}officer\textquotedbl{}]
}

{\footnotesize Визначено персонал або ролі, які отримують сповіщення про збої в процесі обробки даних аудиту }


{\footnotesize\bfseries
No: 4\\
Name: au\_\allowbreak{}5\_\allowbreak{}odp\_\allowbreak{}02\\
Type: list\\
Default: [\textquotedbl{}admin\textquotedbl{}, \textquotedbl{}security\_\allowbreak{}officer\textquotedbl{}]
}

{\footnotesize Визначено період часу, протягом якого персонал або ролі отримують сповіщення про збої в процесі обробки даних аудиту}


{\footnotesize\bfseries
No: 5\\
Name: au\_\allowbreak{}5\_\allowbreak{}odp\_\allowbreak{}03\\
Type: list\\
Default: [\textquotedbl{}login\textquotedbl{}, \textquotedbl{}logout\textquotedbl{}, \textquotedbl{}failed\_\allowbreak{}attempt\textquotedbl{}]
}

{\footnotesize Визначено додаткові дії, яких слід вжити у випадку збою в процесі обробки даних аудиту }




\subsubsection{МІСТКІСТЬ СХОВИЩА ЗАПИСІВ АУДИТУ (AU-5(1))}


\vspace{10pt}
{\footnotesize\bfseries
No: 1\\
Name: au\_\allowbreak{}5\_\allowbreak{}1\_\allowbreak{}odp\_\allowbreak{}01\\
Type: list\\
Default: [\textquotedbl{}admin\textquotedbl{}, \textquotedbl{}security\_\allowbreak{}officer\textquotedbl{}]
}

{\footnotesize Визначено персонал або ролі, які мають бути попереджені, коли обсяг записів аудиту, що зберігаються, досягає максимуму місткості сховища}


{\footnotesize\bfseries
No: 2\\
Name: au\_\allowbreak{}5\_\allowbreak{}1\_\allowbreak{}odp\_\allowbreak{}02\\
Type: string\\
Default: nil
}

{\footnotesize Визначено період часу, протягом якого визначений персонал або ролі будуть попереджені}


{\footnotesize\bfseries
No: 3\\
Name: au\_\allowbreak{}5\_\allowbreak{}1\_\allowbreak{}odp\_\allowbreak{}03\\
Type: string\\
Default: nil
}

{\footnotesize Визначено відсоток максимальної ємності сховища для зберігання журналів аудиту}


{\footnotesize\bfseries
No: 4\\
Name: au\_\allowbreak{}5\_\allowbreak{}1\_\allowbreak{}01\\
Type: string\\
Default: nil
}

{\footnotesize Попередження надається <AU-05(01)\_\allowbreak{}ODP[01] персоналу або ролям> протягом <AU-05(01)\_\allowbreak{}ODP[02] періоду часу>, коли виділений обсяг сховища журналів аудиту досягає <AU-05(01)\_\allowbreak{}ODP[03] відсотків> від максимального обсягу сховища журналів аудиту}




\subsubsection{ТРИВОЖНЕ СПОВІЩЕННЯ В РЕАЛЬНОМУ ЧАСІ (AU-5(2))}
Забезпечити сповіщення в [Призначення: визначений організацією період реального часу] [Призначення: визначених організацією персоналу, ролей та/або місць], коли відбуваються такі події збою аудиту: [Призначення: визначені організацією події, пов'язані зі збоями та помилками аудиту, які вимагають тривоги в реальному часі].

\vspace{10pt}
{\footnotesize\bfseries
No: 1\\
Name: au\_\allowbreak{}5\_\allowbreak{}2\_\allowbreak{}01\\
Type: list\\
Default: [\textquotedbl{}admin\textquotedbl{}, \textquotedbl{}security\_\allowbreak{}officer\textquotedbl{}]
}

{\footnotesize Протягом періоду реального часу надається сповіщення персоналу або ролям, коли виникають події}


{\footnotesize\bfseries
No: 2\\
Name: au\_\allowbreak{}5\_\allowbreak{}2\_\allowbreak{}odp\_\allowbreak{}01\\
Type: integer\\
Default: 30
}

{\footnotesize Визначено період реального часу, за який потрібно надсилати сповіщення при виникненні подій збою аудиту (визначених у AU-05(02)\_\allowbreak{}ODP[03])}


{\footnotesize\bfseries
No: 3\\
Name: au\_\allowbreak{}5\_\allowbreak{}2\_\allowbreak{}odp\_\allowbreak{}02\\
Type: list\\
Default: [\textquotedbl{}admin\textquotedbl{}, \textquotedbl{}security\_\allowbreak{}officer\textquotedbl{}]
}

{\footnotesize Визначено персонал або ролі, які мають бути сповіщені при виникненні подій збоїв в аудиті (визначених в AU05(02)\_\allowbreak{}ODP[03])}


{\footnotesize\bfseries
No: 4\\
Name: au\_\allowbreak{}5\_\allowbreak{}2\_\allowbreak{}odp\_\allowbreak{}03\\
Type: list\\
Default: [\textquotedbl{}login\textquotedbl{}, \textquotedbl{}logout\textquotedbl{}, \textquotedbl{}failed\_\allowbreak{}attempt\textquotedbl{}]
}

{\footnotesize Визначено події, пов'язані зі збоями та помилками аудиту}




\subsubsection{НАЛАШТУВАННЯ ПОРОГОВОГО ОБСЯГУ ТРАФІКУ (AU-5(3))}
Здійснювати налаштування порогових значень обсягу трафіку комунікаційних мереж, що відображають обмеження на можливості аудиту та [Вибір: відхиляти; затримувати] мережевий трафік, якщо він перевищує цей поріг.

\vspace{10pt}
{\footnotesize\bfseries
No: 1\\
Name: au\_\allowbreak{}5\_\allowbreak{}3\_\allowbreak{}odp\\
Type: string\\
Default: nil
}

{\footnotesize Вибрано одне або декілька з наступних ЗНАЧЕНЬ ПАРАМЕТРІВ: \{відхилити; затримати\}}


{\footnotesize\bfseries
No: 2\\
Name: au\_\allowbreak{}5\_\allowbreak{}3\_\allowbreak{}01\\
Type: string\\
Default: nil
}

{\footnotesize Застосовуються налаштовані порогові значення обсягу трафіку комунікаційних мереж, що відображають обмеження на можливості аудиту}


{\footnotesize\bfseries
No: 3\\
Name: au\_\allowbreak{}5\_\allowbreak{}3\_\allowbreak{}02\\
Type: string\\
Default: nil
}

{\footnotesize Мережевий трафік <AU-05(03)\_\allowbreak{}ODP ВИБРАНЕ ЗНАЧЕННЯ ПАРАМЕТРА(ів)>, якщо обсяг мережевого трафіку перевищує налаштовані порогові значення}




\subsubsection{ВИМКНЕННЯ У РАЗІ ВІДМОВИ (AU-5(4))}
Застосовувати [Вибір: повне вимикання системи; часткове вимикання системи; знижений режим роботи з обмеженням доступної/цільової функціональності] у разі [Призначення: визначених організацією збоїв аудиту], якщо немає альтернативної можливості аудиту.

\vspace{10pt}
{\footnotesize\bfseries
No: 1\\
Name: au\_\allowbreak{}5\_\allowbreak{}4\_\allowbreak{}odp\_\allowbreak{}01\\
Type: string\\
Default: nil
}

{\footnotesize Вибрано одне або декілька з наступних ЗНАЧЕНЬ ПАРАМЕТРІВ: \{повне вимикання системи; часткове вимикання системи; знижений режим роботи з обмеженням доступної/цільової функціональності\}}


{\footnotesize\bfseries
No: 2\\
Name: au\_\allowbreak{}5\_\allowbreak{}4\_\allowbreak{}odp\_\allowbreak{}02\\
Type: string\\
Default: nil
}

{\footnotesize Визначено збої аудиту, які спричиняють зміну режиму роботи}


{\footnotesize\bfseries
No: 3\\
Name: au\_\allowbreak{}5\_\allowbreak{}4\_\allowbreak{}01\\
Type: string\\
Default: nil
}

{\footnotesize <AU-05(04)\_\allowbreak{}ODP[01] ВИБРАНЕ ЗНАЧЕННЯ ПАРАМЕТРА(ів)> викликається/викликаються у випадку <AU-05(04)\_\allowbreak{}ODP[02] збоїв аудиту>, якщо не існує альтернативної можливості ведення аудиту}




\subsubsection{МОЖЛИВІСТЬ АЛЬТЕРНАТИВНОГО ЖУРНАЛЮВАННЯ АУДИТУ (AU-5(5))}
Надання альтернативної можливості журналювання аудиту в разі збою основної можливості журналювання аудиту, яка реалізується [Призначення: визначена організацією функція альтернативного журналювання аудиту]

\vspace{10pt}
{\footnotesize\bfseries
No: 1\\
Name: au\_\allowbreak{}5\_\allowbreak{}5\_\allowbreak{}odp\\
Type: string\\
Default: nil
}

{\footnotesize Визначено альтернативний функціонал ведення журналу аудиту на випадок збою в роботі основної функції ведення журналу аудиту}


{\footnotesize\bfseries
No: 2\\
Name: au\_\allowbreak{}5\_\allowbreak{}5\_\allowbreak{}01\\
Type: string\\
Default: nil
}

{\footnotesize Альтернативна можливість ведення журналу аудиту надається на випадок відмови основної можливості ведення журналу аудиту, який реалізує <AU-05(05)\_\allowbreak{}ODP визначений альтернативний функціонал ведення журналу аудиту>}




\subsection{ОГЛЯД, АНАЛІЗ І ЗВІТНІСТЬ АУДИТУ (AU-6)}
a. Переглядати та аналізувати записи системного аудиту [Призначення: з визначеною організацією частотою] для виявлення [Призначення: визначеної організацією неналежної або незвичайної діяльності].\\ 
b. Відправляти звіт про аудит [Призначення: визначеним організацією персоналу або посадам].\\ 
c. Налаштувати рівні огляду аудиту, аналізу та звітності в рамках системи, коли змінюється рівень ризику на основі інформації від правоохоронних органів, розвідувальної інформації або від інших достовірних джерел інформації.

\vspace{10pt}
{\footnotesize\bfseries
No: 1\\
Name: au\_\allowbreak{}6\_\allowbreak{}a\\
Type: string\\
Default: \textquotedbl{}щотижня\textquotedbl{}
}

{\footnotesize Записи аудиту системи переглядаються та аналізуються частота для виявлення ознак неналежної або незвичної діяльності та потенційного впливу неналежної або незвичної діяльності}


{\footnotesize\bfseries
No: 2\\
Name: au\_\allowbreak{}6\_\allowbreak{}b\\
Type: list\\
Default: [\textquotedbl{}admin\textquotedbl{}, \textquotedbl{}security\_\allowbreak{}officer\textquotedbl{}]
}

{\footnotesize Звіт аудиту відправляється персоналу або ролям}


{\footnotesize\bfseries
No: 3\\
Name: au\_\allowbreak{}6\_\allowbreak{}c\\
Type: string\\
Default: nil
}

{\footnotesize Рівень перевірки, аналізу та звітування записів аудиту в системі коригується у разі зміни ризиків на основі інформації правоохоронних органів, розвідувальної інформації або інших достовірних джерел інформації}


{\footnotesize\bfseries
No: 4\\
Name: au\_\allowbreak{}6\_\allowbreak{}odp\_\allowbreak{}01\\
Type: integer\\
Default: 30
}

{\footnotesize Визначено частоту, з якою переглядаються та аналізуються записи аудиту системи}


{\footnotesize\bfseries
No: 5\\
Name: au\_\allowbreak{}6\_\allowbreak{}odp\_\allowbreak{}02\\
Type: string\\
Default: nil
}

{\footnotesize Визначена неналежна або незвична діяльність}


{\footnotesize\bfseries
No: 6\\
Name: au\_\allowbreak{}6\_\allowbreak{}odp\_\allowbreak{}03\\
Type: list\\
Default: [\textquotedbl{}admin\textquotedbl{}, \textquotedbl{}security\_\allowbreak{}officer\textquotedbl{}]
}

{\footnotesize Визначено персонал або ролі які отримують результати оглядів та аналізів системних записів}




\subsubsection{АВТОМАТИЗОВАНА ІНТЕГРАЦІЯ ПРОЦЕСІВ (AU-6(1))}


\vspace{10pt}
{\footnotesize\bfseries
No: 1\\
Name: au\_\allowbreak{}6\_\allowbreak{}1\_\allowbreak{}01\\
Type: string\\
Default: nil
}

{\footnotesize Процеси перегляду, аналізу та звітності інтегровані з використанням автоматизованих механізмів}


{\footnotesize\bfseries
No: 2\\
Name: au\_\allowbreak{}6\_\allowbreak{}1\_\allowbreak{}odp\\
Type: string\\
Default: \textquotedbl{}автоматизований засіб моніторингу\textquotedbl{}
}

{\footnotesize Визначено автоматизовані механізми, що використовуються для інтеграції процесів перегляду, аналізу та звітності записів аудиту}




\subsubsection{ОГЛЯД, АНАЛІЗ І ЗВІТНІСТЬ АУДИТУ СПОВІЩЕННЯ ПРО ПОРУШЕННЯ БЕЗПЕКУ (AU-6(2)) [Вилучено]}
[Вилучено: Включено до SI-4]

\vspace{10pt}
\textit{Немає параметрів для цього контролю.}
\vspace{10pt}


\subsubsection{ЗІСТАВЛЯННЯ СХОВИЩ АУДИТУ (AU-6(3))}
Аналізувати та зіставляти записи аудиту в різних сховищах задля забезпечення ситуативної обізнаності в масштабах організації.

\vspace{10pt}
{\footnotesize\bfseries
No: 1\\
Name: au\_\allowbreak{}6\_\allowbreak{}3\_\allowbreak{}01\\
Type: string\\
Default: nil
}

{\footnotesize Аналізуєються та зіставляються записи аудиту в різних сховищах, задля забезпечення ситуативної обізнаності в масштабах організації}




\subsubsection{ЦЕНТРАЛІЗОВАНИЙ ПЕРЕГЛЯД ТА АНАЛІЗ (AU-6(4))}
Забезпечити та впровадити можливість централізованого перегляду та аналізу записів аудиту з декількох компонентів у системі.

\vspace{10pt}
{\footnotesize\bfseries
No: 1\\
Name: au\_\allowbreak{}6\_\allowbreak{}4\_\allowbreak{}01\\
Type: string\\
Default: nil
}

{\footnotesize Забезпечино можливість централізованого перегляду та аналізу записів аудиту з декількох компонентів у системі}


{\footnotesize\bfseries
No: 2\\
Name: au\_\allowbreak{}6\_\allowbreak{}4\_\allowbreak{}02\\
Type: string\\
Default: nil
}

{\footnotesize Впроваджено можливість централізованого перегляду та аналізу записів аудиту з декількох компонентів у системі}




\subsubsection{ІНТЕГРОВАНИЙ АНАЛІЗ ЗАПИСІВ АУДИТУ (AU-6(5))}
Інтегрувати аналіз записів аудиту з аналізом [Вибір (один або більше): інформації про сканування уразливостей; даних про продуктивність; інформації про моніторинг системи; [Призначення: визначених організацією даних/інформації, зібраних з інших джерел]] для подальшого підвищення здатності виявляти неприйнятну або незвичайну діяльність.

\vspace{10pt}
{\footnotesize\bfseries
No: 1\\
Name: au\_\allowbreak{}6\_\allowbreak{}5\_\allowbreak{}odp\_\allowbreak{}01\\
Type: string\\
Default: nil
}

{\footnotesize Вибрано одне або декілька з наступних ЗНАЧЕНЬ ПАРАМЕТРІВ: \{інформація про сканування вразливостей; дані про продуктивність; інформація про моніторинг системи; <AU-06(05)\_\allowbreak{}ODP[02] дані/інформація, зібрана з інших джерел>\}}


{\footnotesize\bfseries
No: 2\\
Name: au\_\allowbreak{}6\_\allowbreak{}5\_\allowbreak{}odp\_\allowbreak{}02\\
Type: string\\
Default: nil
}

{\footnotesize Визначено дані/інформацію, зібрані з інших джерел, що підлягають аналізу (якщо вони були обрані)}


{\footnotesize\bfseries
No: 3\\
Name: au\_\allowbreak{}6\_\allowbreak{}5\_\allowbreak{}01\\
Type: string\\
Default: nil
}

{\footnotesize Аналіз записів аудиту інтегровано з аналізом <AU-06(05)\_\allowbreak{}ODP[01] ВИБРАНЕ ЗНАЧЕННЯ ПАРАМЕТРА(ів)>, для подальшого підвищення здатності виявляти неприйнятну або незвичайну діяльність}




\subsubsection{КОРЕЛЯЦІЯ З ФІЗИЧНИМ МОНІТОРИНГОМ (AU-6(6))}
Зіставляти інформацію із записів аудиту з інформацією, отриманою від моніторингу фізичного доступу, для подальшого підвищення здатності ідентифікувати підозрілу, неприйнятну, незвичайну або зловмисну діяльність.

\vspace{10pt}
{\footnotesize\bfseries
No: 1\\
Name: au\_\allowbreak{}6\_\allowbreak{}6\_\allowbreak{}01\\
Type: string\\
Default: nil
}

{\footnotesize Інформація з записів аудиту співвідноситься з інформацією, отриманою в результаті моніторингу фізичного доступу, для подальшого посилення здатності виявляти підозрілу, невідповідну, незвичну або зловмисну діяльність}




\subsubsection{ДОЗВОЛЕНІ ДІЇ (AU-6(7))}
Визначити дозволені дії для кожного [Вибір (один або кілька): системного процесу; ролі; користувача], пов'язаного з переглядом, аналізом та поданням інформації про аудит.

\vspace{10pt}
{\footnotesize\bfseries
No: 1\\
Name: au\_\allowbreak{}6\_\allowbreak{}7\_\allowbreak{}odp\\
Type: list\\
Default: [\textquotedbl{}admin\textquotedbl{}]
}

{\footnotesize вибрано одне або декілька з наступних ЗНАЧЕНЬ ПАРА- МЕТРІВ: \{процес системи; роль; користувач\};}


{\footnotesize\bfseries
No: 2\\
Name: au\_\allowbreak{}6\_\allowbreak{}7\_\allowbreak{}01\\
Type: string\\
Default: nil
}

{\footnotesize визначено дозволені дії для кожного <AU-06(07)\_\allowbreak{}ODP ВИБРАНЕ ЗНАЧЕННЯ ПАРАМЕТРА(ів)> , пов'язані з переглядом, аналізом та поданням інформації про аудит.}




\subsubsection{АНАЛІЗ ПОВНОГО ТЕКСТУ ПРИВІЛЕЙОВАНИХ КОМАНД (AU-6(8))}
Виконувати повний аналіз тексту привілейованих команд аудиту у фізично окремому компоненті чи підсистемі або іншій системі, яка може виконувати такий аналіз.

\vspace{10pt}
{\footnotesize\bfseries
No: 1\\
Name: au\_\allowbreak{}6\_\allowbreak{}8\_\allowbreak{}01\\
Type: string\\
Default: nil
}

{\footnotesize виконується повний аналіз тексту привілейованих команд аудиту у фізично окремому компоненті чи підсистемі або іншій системі, яка може виконувати такий аналіз. 135}




\subsubsection{КОРЕЛЯЦІЯ З ІНФОРМАЦІЄЮ З НЕТЕХНІЧНИХ ДЖЕРЕЛ (AU-6(9))}
Зіставляти інформацію з нетехнічних джерел з інформацією аудиту з метою посилення організаційної обізнаності.

\vspace{10pt}
{\footnotesize\bfseries
No: 1\\
Name: au\_\allowbreak{}6\_\allowbreak{}9\_\allowbreak{}01\\
Type: string\\
Default: nil
}

{\footnotesize Зіставляється інформація з нетехнічних джерел з інформацією аудиту з метою посилення організаційної обізнаності.}




\subsubsection{РЕГУЛЮВАННЯ РІВНЯ АУДИТУ (AU-6(10)) [Вилучено]}
[Вилучено: Включено до AU-06]

\vspace{10pt}
\textit{Немає параметрів для цього контролю.}
\vspace{10pt}


\subsection{СКОРОЧЕННЯ ЗАПИСІВ АУДИТУ ТА ФОРМУВАННЯ ЗВІТУ (AU-7)}
Забезпечити та реалізувати можливості скорочення записів перевірок аудитом і звітів, до рівня, який:\\ 
a. підтримує перевірку, аналіз і звітність аудиту на вимогу та розслідування (постфактум) інцидентів безпеки;\\ 
b. не змінює оригінальний вміст або час упорядкування записів аудиту.

\vspace{10pt}
{\footnotesize\bfseries
No: 1\\
Name: au\_\allowbreak{}7\_\allowbreak{}a\_\allowbreak{}1\\
Type: string\\
Default: nil
}

{\footnotesize забезпечено можливість скорочення записів перевірок аудитом та звітів, до рінвня що підтримує перевірку, аналіз і звітність аудиту на вимогу та розслідування (постфактум) інцидентів безпеки;}


{\footnotesize\bfseries
No: 2\\
Name: au\_\allowbreak{}7\_\allowbreak{}a\_\allowbreak{}2\\
Type: string\\
Default: nil
}

{\footnotesize реалізовано можливість скорочення записів перевірок аудитом та звітів, до рінвня що підтримує перевірку, аналіз і звітність аудиту на вимогу та розслідування (постфактум) інцидентів безпеки;}


{\footnotesize\bfseries
No: 3\\
Name: au\_\allowbreak{}7\_\allowbreak{}b\_\allowbreak{}1\\
Type: string\\
Default: nil
}

{\footnotesize забезпечено можливість скорочення записів перевірок аудитом та звітів, які не змінюють оригінальний зміст або час упорядкування записів аудиту;}


{\footnotesize\bfseries
No: 4\\
Name: au\_\allowbreak{}7\_\allowbreak{}b\_\allowbreak{}2\\
Type: string\\
Default: nil
}

{\footnotesize реалізовано можливість скорочення записів перевірок аудитом та звітів, які не змінюють оригінальний зміст або час упорядкування записів аудиту;}




\subsubsection{АВТОМАТИЧНА ОБРОБКА (AU-7(1))}


\vspace{10pt}
{\footnotesize\bfseries
No: 1\\
Name: au\_\allowbreak{}7\_\allowbreak{}1\_\allowbreak{}odp\\
Type: string\\
Default: nil
}

{\footnotesize визначено поля в записах аудиту, які можна обробляти, сортувати або шукати;}


{\footnotesize\bfseries
No: 2\\
Name: au\_\allowbreak{}7\_\allowbreak{}1\_\allowbreak{}1\\
Type: string\\
Default: nil
}

{\footnotesize забезпечити можливість обробки записів аудиту для подій, що представляють інтерес, на основі <AU-07(01)\_\allowbreak{}ODP полей в записах аудиту>}


{\footnotesize\bfseries
No: 3\\
Name: au\_\allowbreak{}7\_\allowbreak{}1\_\allowbreak{}2\\
Type: string\\
Default: nil
}

{\footnotesize реалізувати можливість обробки записів аудиту для подій, що представляють інтерес, на основі <AU-07(01)\_\allowbreak{}ODP полей в записах аудиту>}




\subsubsection{АВТОМАТИЧНЕ СОРТУВАННЯ ТА ПОШУК (AU-7(2)) [Вилучено]}
[Вилучено: Включено до AU-07(01)]

\vspace{10pt}
\textit{Немає параметрів для цього контролю.}
\vspace{10pt}


\subsection{ПОЗНАЧКА ЧАСУ (AU-8)}
a. Використовувати внутрішньосистемний годинник для створення позначок часу для записів аудиту.\\ 
b. Застосовувати позначки часу, які відповідають [Призначення: деталізація вимірювання часу, визначена організацією] і використовують всесвітній координований час, мають фіксоване зміщення місцевого часу відносно всесвітнього координованого часу або включають зміщення місцевого часу як частину позначки часу.

\vspace{10pt}
{\footnotesize\bfseries
No: 1\\
Name: au\_\allowbreak{}8\_\allowbreak{}odp\\
Type: string\\
Default: nil
}

{\footnotesize визначено деталізацію вимірювання часу для часових позначок записів аудиту;}


{\footnotesize\bfseries
No: 2\\
Name: au\_\allowbreak{}8\_\allowbreak{}a\\
Type: string\\
Default: nil
}

{\footnotesize внутрішній системний годинник використовується для створення позначок часу для записів аудиту;}


{\footnotesize\bfseries
No: 3\\
Name: au\_\allowbreak{}8\_\allowbreak{}b\\
Type: string\\
Default: nil
}

{\footnotesize позначки часу застосовуються для записів аудиту, які відповідають <AU-08\_\allowbreak{}ODP деталізація вимірювання часу>  і які використовують всесвітній координований час, мають фіксоване місцеве зміщення місцевого часу від всесвітнього координованого часу або включають зміщення місцевого часу як частину позначки часу.}




\subsubsection{СИНХРОНІЗАЦІЯ З АВТОРИТЕТНИМ ДЖЕРЕЛОМ ЧАСУ (AU-8(1)) [Вилучено]}
[Вилучено: Включено до SC-45(01)]

\vspace{10pt}
\textit{Немає параметрів для цього контролю.}
\vspace{10pt}


\subsubsection{ВТОРИННЕ АВТОРИТЕТНЕ ДЖЕРЕЛО ЧАСУ (AU-8(2)) [Вилучено]}
[Вилучено: Включено до SC-45(02)]

\vspace{10pt}
\textit{Немає параметрів для цього контролю.}
\vspace{10pt}


\subsection{ЗАХИСТ ІНФОРМАЦІЇ АУДИТУ (AU-9)}
a. Захист інформації аудиту та інструментів несанкціонованого доступу, зміни та видалення; журналювання аудиту від\\ 
b. Сповіщення [Призначення: персонал або ролі, визначені організацією] у разі виявлення несанкціонованого доступу, зміни або видалення інформації аудиту.

\vspace{10pt}
{\footnotesize\bfseries
No: 1\\
Name: au\_\allowbreak{}9\_\allowbreak{}odp\\
Type: list\\
Default: [\textquotedbl{}admin\textquotedbl{}]
}

{\footnotesize визначено персонал або ролі, які мають бути сповіщені при виявленні несанкціонованого доступу, зміни або видалення інформації аудиту;}


{\footnotesize\bfseries
No: 2\\
Name: au\_\allowbreak{}9\_\allowbreak{}a\\
Type: string\\
Default: nil
}

{\footnotesize інформація про аудит та інструменти журналювання аудиту захищені від несанкціонованого доступу, зміни та видалення;}


{\footnotesize\bfseries
No: 3\\
Name: au\_\allowbreak{}9\_\allowbreak{}b\\
Type: list\\
Default: [\textquotedbl{}admin\textquotedbl{}]
}

{\footnotesize <AU-09\_\allowbreak{}ODP персонал або ролі>  отримують сповіщення при виявленні несанкціонованого доступу, зміни або видалення інформації аудиту.}




\subsubsection{АПАРАТНІ НОСІЇ ІНФОРМАЦІЇ ОДНОРАЗОВОГО ЗАПИСУ (AU-9(1))}


\vspace{10pt}
{\footnotesize\bfseries
No: 1\\
Name: au\_\allowbreak{}9\_\allowbreak{}1\_\allowbreak{}01\\
Type: string\\
Default: nil
}

{\footnotesize журнали аудиту записані на апаратні носії інформації з одноразовим записом.}




\subsubsection{ЗБЕРІГАННЯ НА ОКРЕМИХ ФІЗИЧНИХ СИСТЕМАХ АБО КОМПОНЕНТАХ (AU-9(2))}
Зберігати записи аудиту з [Призначення: визначеною організацією з частотою] у репозиторії, який є частиною фізично іншої системи або компонента системи, ніж система або компонент, який перевіряється.

\vspace{10pt}
{\footnotesize\bfseries
No: 1\\
Name: au\_\allowbreak{}9\_\allowbreak{}2\_\allowbreak{}odp\\
Type: integer\\
Default: 30
}

{\footnotesize визначено частоту з якою необхідно зберігати записи аудиту;}


{\footnotesize\bfseries
No: 2\\
Name: au\_\allowbreak{}9\_\allowbreak{}2\_\allowbreak{}01\\
Type: string\\
Default: nil
}

{\footnotesize зберігати записи аудиту з <AU-09(02)\_\allowbreak{}ODP частотою> у репозиторії, який є частиною іншої системи або компонента системи, не частиною системи або компонента системи, який перевіряється.}




\subsubsection{КРИПТОГРАФІЧНИЙ ЗАХИСТ (AU-9(3))}
Запровадити криптографічні механізми для захисту цілісності інформації аудиту та інструментів аудиту.

\vspace{10pt}
{\footnotesize\bfseries
No: 1\\
Name: au\_\allowbreak{}9\_\allowbreak{}3\_\allowbreak{}01\\
Type: string\\
Default: nil
}

{\footnotesize впроваджено криптографічні механізми для захисту цілісності інформації аудиту}




\subsubsection{ДОСТУП, ЯКИЙ НАДАЄТЬСЯ ЧЕРЕЗ ЧЛЕНСТВО В ПІДМНОЖИНИ ПРИВІЛЕЙОВАНИХ КОРИСТУВАЧІВ (AU-9(4))}
Авторизувати доступ до управління функціональністю аудиту тільки для [Призначення: визначеної організацією підмножини привілейованих користувачів].

\vspace{10pt}
{\footnotesize\bfseries
No: 1\\
Name: au\_\allowbreak{}9\_\allowbreak{}4\_\allowbreak{}odp\\
Type: string\\
Default: nil
}

{\footnotesize визначено підмножину привілейованих користувачів;}


{\footnotesize\bfseries
No: 2\\
Name: au\_\allowbreak{}9\_\allowbreak{}4\_\allowbreak{}01\\
Type: string\\
Default: nil
}

{\footnotesize авторизувати доступ до управління функціональністю аудиту тільки для <AU-09(04)\_\allowbreak{}ODP підмножини привілейованих користувачів>.}




\subsubsection{ПОДВІЙНА АВТОРИЗАЦІЯ (AU-9(5))}
Здійснювати подвійну авторизацію для [Вибір (один або кілька): переміщення; видалення] [Призначення: визначеної організацією інформації аудиту].

\vspace{10pt}
{\footnotesize\bfseries
No: 1\\
Name: au\_\allowbreak{}9\_\allowbreak{}5\_\allowbreak{}odp\_\allowbreak{}1\\
Type: integer\\
Default: 30
}

{\footnotesize вибрано одне або декілька з наступних ЗНАЧЕНЬ ПАРАМЕТРІВ: \{переміщення; видалення\};}


{\footnotesize\bfseries
No: 2\\
Name: au\_\allowbreak{}9\_\allowbreak{}5\_\allowbreak{}odp\_\allowbreak{}2\\
Type: string\\
Default: nil
}

{\footnotesize визначено інформацію аудиту, для якої має бути застосована подвійна авторизація;}


{\footnotesize\bfseries
No: 3\\
Name: au\_\allowbreak{}9\_\allowbreak{}5\_\allowbreak{}01\\
Type: string\\
Default: nil
}

{\footnotesize подвійна авторизація застосовується для <AU09(05)\_\allowbreak{}ODP[01] ВИБІРАНЕ ЗНАЧЕННЯ ПАРАМЕТРА(ів)> з <AU-09(05)\_\allowbreak{}ODP[02] інформації аудиту>.}




\subsubsection{ДОСТУП ТІЛЬКИ ДЛЯ ЧИТАННЯ (AU-9(6))}
Авторизувати доступ лише для читання інформації аудиту для [Призначення: визначеної організацією підмножини привілейованих користувачів].

\vspace{10pt}
{\footnotesize\bfseries
No: 1\\
Name: au\_\allowbreak{}9\_\allowbreak{}6\_\allowbreak{}odp\\
Type: string\\
Default: nil
}

{\footnotesize визначено підмножину привілейованих користувачів для яких доступ авторизовано тільки для читання.}


{\footnotesize\bfseries
No: 2\\
Name: au\_\allowbreak{}9\_\allowbreak{}6\_\allowbreak{}01\\
Type: string\\
Default: nil
}

{\footnotesize авторизувати доступ лише для читання інформації аудиту для <AU-09(06)\_\allowbreak{}ODP підмножини привілейованих користувачів>.}




\subsubsection{ЗБЕРІГАННЯ НА КОМПОНЕНТІ ІНШОЇ ОПЕРАЦІЙНОЇ СИСТЕМИ (AU-9(7))}
Зберігати інформацію про аудит на компоненті, що працює з іншою операційною системою, ніж система або компонент, який проходить аудит.

\vspace{10pt}
{\footnotesize\bfseries
No: 1\\
Name: au\_\allowbreak{}9\_\allowbreak{}7\_\allowbreak{}01\\
Type: string\\
Default: nil
}

{\footnotesize інформація про аудит зберігається на компоненті, що працює з іншою операційною системою, ніж система або компонент, який проходить аудит}




\subsection{НЕСПРОСТОВНІСТЬ (AU-10)}
Надавайте неспростовні докази того, що особа (або процес, який діє від імені особи) виконала [Призначення: дії, визначені організацією, на які поширюється принцип неспростовності].

\vspace{10pt}
{\footnotesize\bfseries
No: 1\\
Name: au\_\allowbreak{}10\_\allowbreak{}odp\\
Type: string\\
Default: nil
}

{\footnotesize визначено дії, на які поширюється принцип неспростовності;}


{\footnotesize\bfseries
No: 2\\
Name: au\_\allowbreak{}10\_\allowbreak{}01\\
Type: string\\
Default: nil
}

{\footnotesize надаються неспростовні докази того, що особа (або процес, що діє від імені особи) виконала <AU-10\_\allowbreak{}ODP дії>.}




\subsubsection{АСОЦІАЦІЯ ІДЕНТИЧНОСТІ (AU-10(1))}


\vspace{10pt}
{\footnotesize\bfseries
No: 1\\
Name: au\_\allowbreak{}10\_\allowbreak{}1\_\allowbreak{}odp\\
Type: string\\
Default: nil
}

{\footnotesize визначено міцність зв'язку між особистістю джерела інформації та інформацією;}


{\footnotesize\bfseries
No: 2\\
Name: au\_\allowbreak{}10\_\allowbreak{}1\_\allowbreak{}a\\
Type: string\\
Default: nil
}

{\footnotesize особистість джерела інформації зв'язується з інформацією з <AU-10(01)\_\allowbreak{}ODP сила зв'язування>;}


{\footnotesize\bfseries
No: 3\\
Name: au\_\allowbreak{}10\_\allowbreak{}1\_\allowbreak{}b\\
Type: string\\
Default: nil
}

{\footnotesize впроваджено засоби, якими уповноважені особи можуть визначити особу виробника інформації.}




\subsubsection{РАТИФІКАЦІЯ ПРИВ'ЯЗКИ ІНФОРМАЦІЇ ПРО ІДЕНТИЧНІСТЬ ВИРОБНИКА (AU-10(2))}
a) Підтвердити прив'язку інформації про ідентичність джерела до інформації з [Призначення: визначеною організацією частотою]. b) Виконати [Призначення: визначені організацією дії] у разі помилки перевірки.

\vspace{10pt}
{\footnotesize\bfseries
No: 1\\
Name: au\_\allowbreak{}10\_\allowbreak{}2\_\allowbreak{}01\\
Type: string\\
Default: nil
}

{\footnotesize НЕСПРОСТОВНІСТЬ - РАТИФІКАЦІЯ ПРИВ'ЯЗКИ ІНФОРМАЦІЇ ПРО ІДЕНТИЧНІСТЬ ВИРОБНИКА МЕТА ОЦІНКИ: Визначити, чи:}


{\footnotesize\bfseries
No: 2\\
Name: au\_\allowbreak{}10\_\allowbreak{}2\_\allowbreak{}a\\
Type: integer\\
Default: 30
}

{\footnotesize Прив'язка інформації про ідентичність джерела до інформації підтвержується з частотою}


{\footnotesize\bfseries
No: 3\\
Name: au\_\allowbreak{}10\_\allowbreak{}2\_\allowbreak{}b\\
Type: list\\
Default: [\textquotedbl{}login\textquotedbl{}, \textquotedbl{}logout\textquotedbl{}, \textquotedbl{}failed\_\allowbreak{}attempt\textquotedbl{}]
}

{\footnotesize Виконуються дії у разі помилки перевірки}


{\footnotesize\bfseries
No: 4\\
Name: au\_\allowbreak{}10\_\allowbreak{}2\_\allowbreak{}odp\_\allowbreak{}01\\
Type: integer\\
Default: 30
}

{\footnotesize Визначено частоту, з якою необхідно підтверджувати прив'язку інформації про ідентичність джерела до інформації}




\subsubsection{ЛАНЦЮЖОК ЗБЕРЕЖЕННЯ ДОКАЗІВ (AU-10(3))}
Підтримувати перегляд і випуск ідентичності та повноважень у межах встановленого ланцюжка збереження доказів для всієї переглянутої або оприлюдненої інформації.

\vspace{10pt}
{\footnotesize\bfseries
No: 1\\
Name: au\_\allowbreak{}10\_\allowbreak{}3\_\allowbreak{}01\\
Type: string\\
Default: nil
}

{\footnotesize Підтримується перегляд і випуск ідентичності та повноважень у межах встановленого ланцюжка збереження доказів для всієї переглянутої або оприлюдненої інформації}




\subsubsection{ВАЛІДАЦІЯ ЗВ'ЯЗКУ ІДЕНТИЧНОСТІ ПЕРЕГЛЯ- (AU-10(4))}
ЗВ'ЯЗКУ ІДЕНТИЧНОСТІ a) Підтвердити прив'язку особистості рецензента до інформації в точках передачі або видачі до її випуску або передачі між [Призначення: визначеними організацією домени безпеки]. b) Виконати [Призначення: визначені організацією дії] у разі помилки перевірки.

\vspace{10pt}
{\footnotesize\bfseries
No: 1\\
Name: au\_\allowbreak{}10\_\allowbreak{}4\_\allowbreak{}a\\
Type: list\\
Default: [\textquotedbl{}admin\textquotedbl{}, \textquotedbl{}security\_\allowbreak{}officer\textquotedbl{}]
}

{\footnotesize Прив'язка особистості рецензента інформації до інформації в точках передачі або видачі до її випуску або передачі між доменами безпеки підтверджується}


{\footnotesize\bfseries
No: 2\\
Name: au\_\allowbreak{}10\_\allowbreak{}4\_\allowbreak{}b\\
Type: list\\
Default: [\textquotedbl{}login\textquotedbl{}, \textquotedbl{}logout\textquotedbl{}, \textquotedbl{}failed\_\allowbreak{}attempt\textquotedbl{}]
}

{\footnotesize Дії виконуються у випадку помилки перевірки}


{\footnotesize\bfseries
No: 3\\
Name: au\_\allowbreak{}10\_\allowbreak{}4\_\allowbreak{}odp\_\allowbreak{}01\\
Type: list\\
Default: [\textquotedbl{}admin\textquotedbl{}, \textquotedbl{}security\_\allowbreak{}officer\textquotedbl{}]
}

{\footnotesize Визначено домени безпеки, для яких прив'язка особи рецензента інформації до інформації повинна бути підтверджена в точках передачі або видачі}


{\footnotesize\bfseries
No: 4\\
Name: au\_\allowbreak{}10\_\allowbreak{}4\_\allowbreak{}odp\_\allowbreak{}02\\
Type: list\\
Default: [\textquotedbl{}login\textquotedbl{}, \textquotedbl{}logout\textquotedbl{}, \textquotedbl{}failed\_\allowbreak{}attempt\textquotedbl{}]
}

{\footnotesize Визначено дії, які мають бути виконані у випадку помилки перевірки}




\subsubsection{ЦИФРОВІ ПІДПИСИ (AU-10(5)) [Вилучено]}
[Вилучено: Включено до SI-07]

\vspace{10pt}
\textit{Немає параметрів для цього контролю.}
\vspace{10pt}


\subsection{ЗБЕРЕЖЕННЯ ЗАПИСІВ АУДИТУ (AU-11)}
Зберігати записи аудиту впродовж [Призначення: визначеного організацією періоду часу, відповідно політиці зберігання записів], щоб забезпечити підтримку розслідувань (постфактум) інцидентів безпеки та приватності, а також для задоволення вимог нормативних і документів організації щодо збереження даних аудиту.

\vspace{10pt}
{\footnotesize\bfseries
No: 1\\
Name: au\_\allowbreak{}11\_\allowbreak{}odp\\
Type: string\\
Default: nil
}

{\footnotesize визначено період часу для зберігання записів аудиту, який узгоджується з політикою зберігання записів;}


{\footnotesize\bfseries
No: 2\\
Name: au\_\allowbreak{}11\_\allowbreak{}01\\
Type: string\\
Default: nil
}

{\footnotesize записи аудиту зберігаються впродовж <AU-11\_\allowbreak{}ODP період часу>, щоб забезпечити підтримку розслідування (постфактум) інцидентів безпеки та конфіденційності, а також відповідати нормативним та вимогам організації щодо збереження даних аудиту.}




\subsubsection{ДОВГОСТРОКОВА МОЖЛИВІСТЬ ОТРИМАННЯ (AU-11(1))}


\vspace{10pt}
{\footnotesize\bfseries
No: 1\\
Name: au\_\allowbreak{}11\_\allowbreak{}1\_\allowbreak{}odp\\
Type: string\\
Default: nil
}

{\footnotesize визначено заходи, необхідні для реалізації довгострокової можливості отримання записів аудиту}


{\footnotesize\bfseries
No: 2\\
Name: au\_\allowbreak{}11\_\allowbreak{}1\_\allowbreak{}01\\
Type: string\\
Default: nil
}

{\footnotesize впровадити <AU-11(01)\_\allowbreak{}ODP заходи>, щоб гарантувати, що довгострокові записи аудиту, можуть бути отримані.}




\subsection{ГЕНЕРАЦІЯ ДАНИХ АУДИТУ (AU-12)}
a. Забезпечити генерацію даних аудиту для типів подій, що перевіряються в AU-2а в [Призначення: визначених організацією компонентах системи].\\ 
b. Дозволити [Призначення: визначеному організацією персоналу або посадам] вибирати, які типи подій, що перевіряються, повинні перевірятися окремими компонентами системи;\\ 
c. Генерувати записи аудиту для типів подій, визначених в AU-2с. з вмістом згідно з AU-3.

\vspace{10pt}
{\footnotesize\bfseries
No: 1\\
Name: au\_\allowbreak{}12\_\allowbreak{}odp\_\allowbreak{}1\\
Type: string\\
Default: nil
}

{\footnotesize визначено компоненти системи, які забезпечують можливість}


{\footnotesize\bfseries
No: 2\\
Name: au\_\allowbreak{}12\_\allowbreak{}odp\_\allowbreak{}2\\
Type: list\\
Default: [\textquotedbl{}admin\textquotedbl{}]
}

{\footnotesize визначено персонал або ролі, яким дозволено обирати типи подій, що мають реєструватися певними компонентами системи;}


{\footnotesize\bfseries
No: 3\\
Name: au\_\allowbreak{}12\_\allowbreak{}a\\
Type: string\\
Default: nil
}

{\footnotesize можливість генерації записів аудиту для типів подій, які система}


{\footnotesize\bfseries
No: 4\\
Name: au\_\allowbreak{}12\_\allowbreak{}b\\
Type: list\\
Default: [\textquotedbl{}admin\textquotedbl{}]
}

{\footnotesize <AU-12\_\allowbreak{}ODP[02] персонал або ролі> може/можуть вибирати типи подій, які будуть реєструватися певними компонентами системи; 146}


{\footnotesize\bfseries
No: 5\\
Name: au\_\allowbreak{}12\_\allowbreak{}c\\
Type: string\\
Default: nil
}

{\footnotesize згенеровано записи аудиту для типів подій, визначених у AU02\_\allowbreak{}ODP[02], які включають вміст записів аудиту, визначений у AU03.}




\subsubsection{ЗАГАЛЬНОСИСТЕМНИЙ ТА СИНХРОНІЗОВАНИЙ ЗА ЧАСОМ ЖУРНАЛУ АУДИТУ (AU-12(1))}


\vspace{10pt}
{\footnotesize\bfseries
No: 1\\
Name: au\_\allowbreak{}12\_\allowbreak{}1\_\allowbreak{}odp\_\allowbreak{}1\\
Type: string\\
Default: nil
}

{\footnotesize визначено компоненти системи, з яких записи аудиту мають бути зібрані в загальносистемний (логічний або фізичний) журнал аудиту;}


{\footnotesize\bfseries
No: 2\\
Name: au\_\allowbreak{}12\_\allowbreak{}1\_\allowbreak{}odp\_\allowbreak{}1\\
Type: string\\
Default: nil
}

{\footnotesize визначено рівень взаємозв'язку між мітками часу окремих записів у журналах аудиту;}


{\footnotesize\bfseries
No: 3\\
Name: au\_\allowbreak{}12\_\allowbreak{}1\_\allowbreak{}01\\
Type: string\\
Default: nil
}

{\footnotesize записи аудиту з <AU-12(01)\_\allowbreak{}ODP[01] компонентів системи> збираються у загальносистемний (логічний або фізичний) журнал аудиту, який синхронізується у часі в межах}




\subsubsection{СТАНДАРТИЗОВАНІ ФОРМАТИ (AU-12(2))}
Створити загальносистемний (логічний або фізичний) журнал аудиту, що складається із записів аудиту в стандартизованому форматі.

\vspace{10pt}
{\footnotesize\bfseries
No: 1\\
Name: au\_\allowbreak{}12\_\allowbreak{}2\_\allowbreak{}01\\
Type: string\\
Default: nil
}

{\footnotesize Створюється загальносистемний (логічний) журнал аудиту, що складається з записів аудиту в стандартизованому форматі}




\subsubsection{ЗМІНИ, ЩО ВНОСЯТЬ АВТОРИЗОВАНІ ОСОБИ (AU-12(3))}
Забезпечити та реалізувати можливість для [Призначення: визначених організацією окремих осіб або ролей] змінити аудит, який виконуватиметься на [Призначення: визначених організацією компонентах системи] на основі [Призначення: визначених організацією критеріїв вибору подій] у межах [Призначення: визначених організацією часових порогів].

\vspace{10pt}
{\footnotesize\bfseries
No: 1\\
Name: au\_\allowbreak{}12\_\allowbreak{}3\_\allowbreak{}odp\_\allowbreak{}1\\
Type: list\\
Default: [\textquotedbl{}admin\textquotedbl{}]
}

{\footnotesize Визначено осіб або ролі, яким дозволено змінювати аудит компонентів системи;}


{\footnotesize\bfseries
No: 2\\
Name: au\_\allowbreak{}12\_\allowbreak{}3\_\allowbreak{}odp\_\allowbreak{}2\\
Type: string\\
Default: nil
}

{\footnotesize Визначено компоненти системи, на яких має виконуватися аудит;}


{\footnotesize\bfseries
No: 3\\
Name: au\_\allowbreak{}12\_\allowbreak{}3\_\allowbreak{}odp\_\allowbreak{}3\\
Type: string\\
Default: nil
}

{\footnotesize Визначено критерії вибору подій;}


{\footnotesize\bfseries
No: 4\\
Name: au\_\allowbreak{}12\_\allowbreak{}3\_\allowbreak{}odp\_\allowbreak{}4\\
Type: string\\
Default: nil
}

{\footnotesize Визначено часові пороги, в яких мають змінюватися аудит;}


{\footnotesize\bfseries
No: 5\\
Name: au\_\allowbreak{}12\_\allowbreak{}3\_\allowbreak{}1\\
Type: string\\
Default: nil
}

{\footnotesize Забезпечено можливість <AU-12(03)\_\allowbreak{}ODP[01] особам або часових порогів>;}


{\footnotesize\bfseries
No: 6\\
Name: au\_\allowbreak{}12\_\allowbreak{}3\_\allowbreak{}2\\
Type: string\\
Default: nil
}

{\footnotesize Реалізовано можливість <AU-12(03)\_\allowbreak{}ODP[01] особам або часових порогів>;}




\subsubsection{АУДИТ ЗАПИТІВ ПЕРСОНАЛЬНИХ ДАНИХ (AU-12(4))}
Забезпечити та реалізувати можливості аудиту параметрів подій запитів користувачів для наборів даних, що містять персональні дані.

\vspace{10pt}
{\footnotesize\bfseries
No: 1\\
Name: au\_\allowbreak{}12\_\allowbreak{}4\_\allowbreak{}1\\
Type: list\\
Default: [\textquotedbl{}admin\textquotedbl{}]
}

{\footnotesize забезпечена можливості аудиту параметрів подій запитів користувачів для наборів даних, що містять персональну ідентифікаційну інформацію.}


{\footnotesize\bfseries
No: 2\\
Name: au\_\allowbreak{}12\_\allowbreak{}4\_\allowbreak{}2\\
Type: list\\
Default: [\textquotedbl{}admin\textquotedbl{}]
}

{\footnotesize реалізувана можливості аудиту параметрів подій запитів користувачів для наборів даних, що містять персональну ідентифікаційну інформацію.}




\subsection{МОНІТОРИНГ РОЗКРИТТЯ ІНФОРМАЦІЇ (AU-13)}
a. Моніторинг [Завдання: визначена організацією інформація з відкритих джерел та/або інформаційних сайтів] [Завдання: частота, визначена організацією] на наявність доказів неавторизованого розголошення конфіденційної інформації;\\ 
b. Якщо виявлено розголошення інформації:\\ 
1. Повідомити [Призначення: персонал або ролі, визначені організацією];\\ 
2. Виконайте такі додаткові дії: [Призначення: додаткові дії, визначені організацією].

\vspace{10pt}
{\footnotesize\bfseries
No: 1\\
Name: au\_\allowbreak{}13\_\allowbreak{}odp\_\allowbreak{}1\\
Type: string\\
Default: nil
}

{\footnotesize визначено інформацію з відкритих джерел та/або інформаційні сайти, що підлягають моніторингу на наявність доказів неавторизованого розголошення конфіденційної інформації;}


{\footnotesize\bfseries
No: 2\\
Name: au\_\allowbreak{}13\_\allowbreak{}odp\_\allowbreak{}2\\
Type: integer\\
Default: 30
}

{\footnotesize визначено частоту моніторингу інформації з відкритих джерел та/або інформаційних сайтів на наявність доказів неавторизованого розголошення конфіденційної інформації;}


{\footnotesize\bfseries
No: 3\\
Name: au\_\allowbreak{}13\_\allowbreak{}odp\_\allowbreak{}3\\
Type: list\\
Default: [\textquotedbl{}admin\textquotedbl{}]
}

{\footnotesize визначено персонал або ролі, які мають бути повідомлені в 149 разі виявлення розголошення інформації;}


{\footnotesize\bfseries
No: 4\\
Name: au\_\allowbreak{}13\_\allowbreak{}odp\_\allowbreak{}4\\
Type: string\\
Default: nil
}

{\footnotesize визначено додаткові дії, як і необхідно вжити у разі виявлення розголошення інформації;}


{\footnotesize\bfseries
No: 5\\
Name: au\_\allowbreak{}13\_\allowbreak{}a\\
Type: string\\
Default: nil
}

{\footnotesize <AU-13\_\allowbreak{}ODP[01] інформація з відкритих джерел та/або інформаційні сайти>  відстежуються <AU-13\_\allowbreak{}ODP[02] частота>  на наявність доказів неавторизованого розголошення конфіденційної інформації;}


{\footnotesize\bfseries
No: 6\\
Name: au\_\allowbreak{}13\_\allowbreak{}b\_\allowbreak{}1\\
Type: list\\
Default: [\textquotedbl{}admin\textquotedbl{}]
}

{\footnotesize <AU-13\_\allowbreak{}ODP[03] персонал або ролі>  буде повідомлено, якщо буде виявлено розголошення інформації;}


{\footnotesize\bfseries
No: 7\\
Name: au\_\allowbreak{}13\_\allowbreak{}b\_\allowbreak{}2\\
Type: string\\
Default: nil
}

{\footnotesize <AU-13\_\allowbreak{}ODP[04] додаткові дії> вживаються, якщо виявлено розголошення інформації.}




\subsubsection{ВИКОРИСТАННЯ АВТОМАТИЧНИХ ЗАСОБІВ (AU-13(1))}


\vspace{10pt}
{\footnotesize\bfseries
No: 1\\
Name: au\_\allowbreak{}13\_\allowbreak{}1\_\allowbreak{}odp\\
Type: string\\
Default: nil
}

{\footnotesize визначено автоматизовані механізми моніторингу інформації з відкритих джерел та інформаційних сайтів;}


{\footnotesize\bfseries
No: 2\\
Name: au\_\allowbreak{}13\_\allowbreak{}1\_\allowbreak{}01\\
Type: string\\
Default: nil
}

{\footnotesize моніторинг інформації з відкритих джерел та інформаційних сайтів здійснюється за допомогою  <AU-13(01)\_\allowbreak{}ODP автоматизовані механізми>.}




\subsubsection{ОГЛЯД САЙТІВ, ЩО ПІДЛЯГАЮТЬ МОНІТОРИНГУ (AU-13(2))}
Проводити огляд відкритих інформаційних сайтів, що підлягають моніторингу [Призначення: з визначеною організацією частотою].

\vspace{10pt}
{\footnotesize\bfseries
No: 1\\
Name: au\_\allowbreak{}13\_\allowbreak{}2\_\allowbreak{}odp\\
Type: integer\\
Default: 30
}

{\footnotesize визначено частоту з якою проводять огляд відкритих інформаційних сайтів, що підлягають моніторингу}


{\footnotesize\bfseries
No: 2\\
Name: au\_\allowbreak{}13\_\allowbreak{}2\_\allowbreak{}01\\
Type: string\\
Default: nil
}

{\footnotesize проводиться огляд відкритих інформаційних сайтів, що підлягають моніторингу з <AU-13(02)\_\allowbreak{}ODP частотою>.}




\subsubsection{АВТОРИЗОВАНЕ КОПІЮВАННЯ ІНФОРМАЦІЇ (AU-13(3))}
Використовуйте методи виявлення, процеси та інструменти, щоб визначити, чи зовнішні суб'єкти копіюють організаційну інформацію неавторизованим способом.

\vspace{10pt}
{\footnotesize\bfseries
No: 1\\
Name: au\_\allowbreak{}13\_\allowbreak{}3\_\allowbreak{}01\\
Type: string\\
Default: nil
}

{\footnotesize застосовуються методи, процеси та інструменти виявлення, щоб визначити, чи не копіюють зовнішні суб'єкти інформацію організації в несанкціонований спосіб.}




\subsection{АУДИТ СЕСІЇ (AU-14)}
a. Надавати та реалізувати можливість для [Призначення: користувачів або ролей, визначених організацією] для [Вибору (одного або кількох): збору/запису або перегляду/прослуховування] вмісту сесії користувача в [Призначення: обставини, визначені організацією];\\ 
b. Розробляти, інтегрувати та використовувати діяльність з аудиту сесії, консультуючись із юрисконсультантом щодо її відповідності до чинних законів, розпоряджень, директив, нормативних актів, політик, стандартів і вказівок.

\vspace{10pt}
{\footnotesize\bfseries
No: 1\\
Name: au\_\allowbreak{}14\_\allowbreak{}odp\_\allowbreak{}1\\
Type: list\\
Default: [\textquotedbl{}admin\textquotedbl{}]
}

{\footnotesize визначено користувачів або ролі, які можуть перевіряти вміст сесії користувача;}


{\footnotesize\bfseries
No: 2\\
Name: au\_\allowbreak{}14\_\allowbreak{}odp\_\allowbreak{}2\\
Type: integer\\
Default: 30
}

{\footnotesize вибрано одне або декілька з наступних ЗНАЧЕНЬ ПАРАМЕТРІВ: \{збору/запису або перегляду/прослуховування\};}


{\footnotesize\bfseries
No: 3\\
Name: au\_\allowbreak{}14\_\allowbreak{}odp\_\allowbreak{}3\\
Type: list\\
Default: [\textquotedbl{}admin\textquotedbl{}]
}

{\footnotesize визначено обставини, за яких вміст сесії користувача може бути перевірено;}


{\footnotesize\bfseries
No: 4\\
Name: au\_\allowbreak{}14\_\allowbreak{}a\_\allowbreak{}1\\
Type: list\\
Default: [\textquotedbl{}admin\textquotedbl{}]
}

{\footnotesize <AU-14\_\allowbreak{}ODP[01] користувачам або ролям>  надається можливість}


{\footnotesize\bfseries
No: 5\\
Name: au\_\allowbreak{}14\_\allowbreak{}a\_\allowbreak{}2\\
Type: list\\
Default: [\textquotedbl{}admin\textquotedbl{}]
}

{\footnotesize реалізовано можливість для <AU-14\_\allowbreak{}ODP[01] користувачів або}


{\footnotesize\bfseries
No: 6\\
Name: au\_\allowbreak{}14\_\allowbreak{}b\_\allowbreak{}1\\
Type: string\\
Default: nil
}

{\footnotesize діяльність з аудиту сесій розробляється після консультацій з юристом та відповідно до чинних законів, розпоряджень, директив, нормативних актів, політик, стандартів та вказівок;}


{\footnotesize\bfseries
No: 7\\
Name: au\_\allowbreak{}14\_\allowbreak{}b\_\allowbreak{}2\\
Type: string\\
Default: nil
}

{\footnotesize діяльність з аудиту сесій інтегрується після консультацій з юристом та відповідно до чинних законів, розпоряджень, директив, нормативних актів, політик, стандартів та вказівок;}


{\footnotesize\bfseries
No: 8\\
Name: au\_\allowbreak{}14\_\allowbreak{}b\_\allowbreak{}3\\
Type: string\\
Default: nil
}

{\footnotesize діяльність з аудиту сесій використовується після консультацій з юристом та відповідно до чинних законів, розпоряджень, директив, нормативних актів, політик, стандартів та вказівок;}




\subsubsection{СИСТЕМА ЗАПУСКУ (AU-14(1))}


\vspace{10pt}
{\footnotesize\bfseries
No: 1\\
Name: au\_\allowbreak{}14\_\allowbreak{}1\_\allowbreak{}01\\
Type: string\\
Default: nil
}

{\footnotesize аудит сесії при запуску системи автоматично ініціюється}




\subsubsection{ЗАХОПЛЕННЯ ТА ЗАПИС ІНФОРМАЦІЇ (AU-14(2)) [Вилучено]}
[Вилучено: Включено до AU-14] [Вилучено: Включено до AU-14]

\vspace{10pt}
\textit{Немає параметрів для цього контролю.}
\vspace{10pt}


\subsubsection{ВІДДАЛЕНИЙ ПЕРЕГЛЯД ТА ПРОСЛУХОВУВАННЯ (AU-14(3))}
Забезпечити та реалізувати можливість авторизованих користувачів віддалено переглядати та прослуховувати вміст, пов'язаний із встановленою сесією користувача, у режимі реального часу.

\vspace{10pt}
{\footnotesize\bfseries
No: 1\\
Name: au\_\allowbreak{}14\_\allowbreak{}3\_\allowbreak{}1\\
Type: list\\
Default: [\textquotedbl{}admin\textquotedbl{}]
}

{\footnotesize забезпечується можливість для авторизованих користувачів віддалено переглядати та прослуховувати вміст, пов'язаний із встановленою сесією користувача, в режимі реального часу;}


{\footnotesize\bfseries
No: 2\\
Name: au\_\allowbreak{}14\_\allowbreak{}3\_\allowbreak{}2\\
Type: list\\
Default: [\textquotedbl{}admin\textquotedbl{}]
}

{\footnotesize реалізовано можливість для авторизованих користувачів віддалено переглядати та прослуховувати вміст, пов'язаний із встановленою сесією користувача, в режимі реального часу;}




\subsection{АЛЬТЕРНАТИВНА МОЖЛИВІСТЬ АУДИТУ (AU-15) [Вилучено]}
[Вилучено: Включено до AU-05(05)]

\vspace{10pt}
\textit{Немає параметрів для цього контролю.}
\vspace{10pt}


\subsection{МІЖОРГАНІЗАЦІЙНИЙ АУДИТ (AU-16)}
Використовувати [Призначення: визначені організацією методи] для координації [Призначення: визначеної організацією інформації] серед зовнішніх організацій, коли інформація аудиту передається за межі організації.

\vspace{10pt}
{\footnotesize\bfseries
No: 1\\
Name: au\_\allowbreak{}16\_\allowbreak{}odp\_\allowbreak{}1\\
Type: string\\
Default: nil
}

{\footnotesize визначено методи для координації інформації серед зовнішніх організацій, коли інформація аудиту передається через (за) межі організації (системи);}


{\footnotesize\bfseries
No: 2\\
Name: au\_\allowbreak{}16\_\allowbreak{}odp\_\allowbreak{}2\\
Type: string\\
Default: nil
}

{\footnotesize визначено інформацію для координації інформації серед зовнішніх організацій, коли інформація аудиту передається через (за) межі організації (системи);;}


{\footnotesize\bfseries
No: 3\\
Name: au\_\allowbreak{}16\_\allowbreak{}01\\
Type: string\\
Default: nil
}

{\footnotesize організація використовує <AU-16\_\allowbreak{}ODP[01] методи> для координації <AU-16\_\allowbreak{}ODP[02] інформації> серед зовнішніх організацій, коли інформація аудиту передається через (за) межі організації (системи).}




\subsubsection{ЗБЕРЕЖЕННЯ ІДЕНТИЧНОСТІ (AU-16(1))}


\vspace{10pt}
{\footnotesize\bfseries
No: 1\\
Name: au\_\allowbreak{}16\_\allowbreak{}1\_\allowbreak{}01\\
Type: string\\
Default: nil
}

{\footnotesize вимагається, щоб ідентичність особистості зберігалася в міжорганізаційних журналах аудиту.}




\subsubsection{ОБМІН ІНФОРМАЦІЄЮ АУДИТУ (AU-16(2))}
Надавати інформацію про міжорганізаційний аудит до [Призначення: організацій, визначених організацією] на основі [Призначення: визначеної організацією міжорганізаційної угоди про обмін].

\vspace{10pt}
{\footnotesize\bfseries
No: 1\\
Name: au\_\allowbreak{}16\_\allowbreak{}2\_\allowbreak{}odp\_\allowbreak{}1\\
Type: string\\
Default: nil
}

{\footnotesize визначено організації до яких надають інформацію про міжорганізаційний аудит}


{\footnotesize\bfseries
No: 2\\
Name: au\_\allowbreak{}16\_\allowbreak{}2\_\allowbreak{}odp\_\allowbreak{}2\\
Type: string\\
Default: nil
}

{\footnotesize визначено міжорганізаційну угоду про розподіл.}


{\footnotesize\bfseries
No: 3\\
Name: au\_\allowbreak{}16\_\allowbreak{}2\_\allowbreak{}01\\
Type: string\\
Default: nil
}

{\footnotesize надати інформацію про міжорганізаційний аудит до  <AU16(02)\_\allowbreak{}ODP[01] організацій>  організацією на основі  <AU16(02)\_\allowbreak{}ODP[02] угоди про розподіл>.}




\subsubsection{РОЗМЕЖУВАННЯ (AU-16(3))}
Запровадити [Призначення: заходи, визначені організацією], щоб розмежувати людей від інформації аудиту, що передається в межах організації.

\vspace{10pt}
{\footnotesize\bfseries
No: 1\\
Name: au\_\allowbreak{}16\_\allowbreak{}3\_\allowbreak{}odp\\
Type: string\\
Default: nil
}

{\footnotesize визначено заходи для розмежування окремих осіб від інформації аудиту, що передається в межах організації;}


{\footnotesize\bfseries
No: 2\\
Name: au\_\allowbreak{}16\_\allowbreak{}3\_\allowbreak{}01\\
Type: string\\
Default: nil
}

{\footnotesize <AU-16(03)\_\allowbreak{}ODP заходи> впроваджуються для того, щоб розмежування окремих осіб від інформації аудиту, що передається в межах організації;}




\section{CA}
Клас заходів захисту CA --- ОЦІНЮВАННЯ,

Цей клас регламентує процеси перевірки ефективності заходів захисту, авторизації систем та їх безперервного моніторингу.

\textbf{Перелік заходів захисту:} Політика і процедури оцінювання, акредитація та моніторинг безпеки (CA-1); Оцінювання (CA-2); Незалежні експерти (CA-2(1)); Спеціалізовані оцінки (CA-2(2)); Зовнішні організації (CA-2(3)); Взаємодія систем (CA-3); Незахищені з'єднання системи (CA-3(1)) [Вилучено]; Захищені з'єднання системи (CA-3(2)); Несекретні з'єднання системи безпеки, що не є національними (CA-3(3)) [Вилучено]; Підключення до загальнодоступних мереж (CA-3(4)) [Вилучено]; Обмеження зв'язку із зовнішніми системами (CA-3(5)) [Вилучено]; Передача дозволів (CA-3(6)); Транзитивний обмін інформацією (CA-3(7)); Сертифікація безпеки (CA-4) [Вилучено]; План усунення недоліків та контрольні показники (CA-5); Автоматизація підтримки задля точності та вживаності (CA-5(1)); Акредитація (CA-6); Спільна акредитація - одна і та сама організація (CA-6(1)); Спільна акредитація - різні організації (CA-6(2)); Безперервний моніторинг (CA-7); Незалежне оцінювання (CA-7(1)); Види оцінок (CA-7(2)) [Вилучено]; Аналіз тенденції (CA-7(3)); Моніторинг ризику (CA-7(4)); Узгоджений аналіз (CA-7(5)); Безперервний моніторингу (CA-7(6)); Тестування на проникнення (CA-8); Незалежна команда або агент на проникнення (CA-8(1)); Червона команда (CA-8(2)); Можливості перевірки на проникнення (CA-8(3)); Внутрішні системні зв'язки (CA-9); Відповідність заходів безпеки (CA-9(1)).

\subsection{ПОЛІТИКА І ПРОЦЕДУРИ ОЦІНЮВАННЯ, АКРЕДИТАЦІЯ ТА МОНІТОРИНГ БЕЗПЕКИ (CA-1)}
a. Розробити, задокументувати та поширити серед [Призначення: визначеного організацією персоналу або посад]:\\ 
1. [Вибір (один або декілька): Рівень організації; Рівень місії/бізнес-процесу; рівень системи] політика оцінювання, авторизації та моніторингу, яка:\\ 
(a) містить мету, сферу застосування, ролі, обов'язки, відповідальність керівництва, координацію між організаційними підрозділами та систему контролю відповідності (complaince);\\ 
(b) відповідає чинним законам, нормативним документам, наказам, положенням, політиці, стандартам і керівним принципам.\\ 
2. Процедури, що сприяють реалізації політики оцінювання, авторизації та моніторингу безпеки та приватності, а також пов'язаних з ними заходів оцінювання, авторизації та моніторингу безпеки та приватності.\\ 
b. Призначити [Призначення: посадова особа, визначена організацією] для управління розробкою, документуванням і розповсюдженням політики та процедур оцінювання, авторизації та моніторингу;\\ 
c. Переглядати та оновлювати поточне оцінювання, авторизацію та моніторинг:\\ 
1. Політику [Призначення: частота, визначена організацією] та наступне [Призначення: події, визначені організацією];\\ 
2. Процедури [Призначення: частота, визначена організацією] та наступні [Призначення: події, визначені організацією].

\vspace{10pt}
{\footnotesize\bfseries
No: 1\\
Name: ca\_\allowbreak{}1\_\allowbreak{}odp\_\allowbreak{}1\\
Type: list\\
Default: [\textquotedbl{}admin\textquotedbl{}]
}

{\footnotesize визначено персонал або ролі, серед яких має бути поширена політика оцінювання, авторизації та моніторингу;}


{\footnotesize\bfseries
No: 2\\
Name: ca\_\allowbreak{}1\_\allowbreak{}odp\_\allowbreak{}2\\
Type: list\\
Default: [\textquotedbl{}admin\textquotedbl{}]
}

{\footnotesize визначено персонал або ролі, серед яких мають бути поширені процедури оцінювання, авторизації та моніторингу;}


{\footnotesize\bfseries
No: 3\\
Name: ca\_\allowbreak{}1\_\allowbreak{}odp\_\allowbreak{}3\\
Type: integer\\
Default: 30
}

{\footnotesize вибрано одне або декілька з наступних ЗНА ЧЕНЬ ПАРА- МЕТРІВ: \{рівень організації; рівень місії/бізнес -процесу; рівень системи\};}


{\footnotesize\bfseries
No: 4\\
Name: ca\_\allowbreak{}1\_\allowbreak{}odp\_\allowbreak{}4\\
Type: string\\
Default: nil
}

{\footnotesize визначено посадову особу, яка управлятиме розробкою, документуванням і розповсюдженням політики та процедур оцінювання, авторизації та моніторингу;}


{\footnotesize\bfseries
No: 5\\
Name: ca\_\allowbreak{}1\_\allowbreak{}odp\_\allowbreak{}5\\
Type: integer\\
Default: 30
}

{\footnotesize визначено частоту, з якою переглядається та оновлюється поточна політика оцінювання, авторизації та моніторингу;}


{\footnotesize\bfseries
No: 6\\
Name: ca\_\allowbreak{}1\_\allowbreak{}odp\_\allowbreak{}6\\
Type: string\\
Default: nil
}

{\footnotesize визначено події, які потребують перегляду та оновлення поточної політики оцінки, авторизації та моніторингу;}


{\footnotesize\bfseries
No: 7\\
Name: ca\_\allowbreak{}1\_\allowbreak{}odp\_\allowbreak{}7\\
Type: integer\\
Default: 30
}

{\footnotesize визначено частоту, з якою переглядається та оновлюється поточні процедури оцінювання, авторизації та моніторингу;}


{\footnotesize\bfseries
No: 8\\
Name: ca\_\allowbreak{}1\_\allowbreak{}odp\_\allowbreak{}8\\
Type: string\\
Default: nil
}

{\footnotesize визначено події, які потребують перегляду та оновлення поточні процедури оцінки, авторизації та моніторингу;}


{\footnotesize\bfseries
No: 9\\
Name: ca\_\allowbreak{}1\_\allowbreak{}a\_\allowbreak{}1\\
Type: string\\
Default: nil
}

{\footnotesize розроблено та задокументовано політику оцінювання, авторизації та моніторингу;}


{\footnotesize\bfseries
No: 10\\
Name: ca\_\allowbreak{}1\_\allowbreak{}a\_\allowbreak{}2\\
Type: string\\
Default: nil
}

{\footnotesize політика оцінювання, авторизації та моніторингу поширюється}


{\footnotesize\bfseries
No: 11\\
Name: ca\_\allowbreak{}1\_\allowbreak{}a\_\allowbreak{}3\\
Type: list\\
Default: [\textquotedbl{}admin\textquotedbl{}]
}

{\footnotesize розроблені та задокументовані процедури оцінювання, авторизації та моніторингу, що сприяють впровадженню політики оцінювання, авторизації та моніторингу, а також пов'язані з ними засоби контролю оцінювання, авторизації та моніторингу;}


{\footnotesize\bfseries
No: 12\\
Name: ca\_\allowbreak{}1\_\allowbreak{}a\_\allowbreak{}4\\
Type: list\\
Default: [\textquotedbl{}admin\textquotedbl{}]
}

{\footnotesize процедури оцінювання, авторизації та  моніторингу поширюються серед <CA-01\_\allowbreak{}ODP[02] персоналу або ролей>; 156}


{\footnotesize\bfseries
No: 13\\
Name: ca\_\allowbreak{}1\_\allowbreak{}a\_\allowbreak{}1\_\allowbreak{}a\_\allowbreak{}1\\
Type: string\\
Default: nil
}

{\footnotesize політика <CA-01\_\allowbreak{}ODP[03] ВИБРАНЕ ЗНАЧЕННЯ ПАРАМЕТРА(ів)> оцінювання, авторизації та моніторингу містить мету;}


{\footnotesize\bfseries
No: 14\\
Name: ca\_\allowbreak{}1\_\allowbreak{}a\_\allowbreak{}1\_\allowbreak{}a\_\allowbreak{}2\\
Type: string\\
Default: nil
}

{\footnotesize політика <CA-01\_\allowbreak{}ODP[03] ВИБРАНЕ ЗНАЧЕННЯ ПАРАМЕТРА(ів)> оцінювання, авторизації та моніторингу містить сферу застосування;}


{\footnotesize\bfseries
No: 15\\
Name: ca\_\allowbreak{}1\_\allowbreak{}a\_\allowbreak{}1\_\allowbreak{}a\_\allowbreak{}3\\
Type: list\\
Default: [\textquotedbl{}admin\textquotedbl{}]
}

{\footnotesize політика <CA-01\_\allowbreak{}ODP[03] ВИБРАНЕ ЗНАЧЕННЯ ПАРАМЕТРА(ів)> оцінювання, авторизації та моніторингу містить ролі;}


{\footnotesize\bfseries
No: 16\\
Name: ca\_\allowbreak{}1\_\allowbreak{}a\_\allowbreak{}1\_\allowbreak{}a\_\allowbreak{}4\\
Type: string\\
Default: nil
}

{\footnotesize політика <CA-01\_\allowbreak{}ODP[03] ВИБРАНЕ ЗНАЧЕННЯ ПАРАМЕТРА(ів)> оцінювання, авторизації та моніторингу містить обов'язки;}


{\footnotesize\bfseries
No: 17\\
Name: ca\_\allowbreak{}1\_\allowbreak{}a\_\allowbreak{}1\_\allowbreak{}a\_\allowbreak{}5\\
Type: string\\
Default: nil
}

{\footnotesize політика <CA-01\_\allowbreak{}ODP[03] ВИБРАНЕ ЗНАЧЕННЯ ПАРАМЕТРА(ів)> оцінювання, авторизації та моніторингу містить відповідальність керівництва;}


{\footnotesize\bfseries
No: 18\\
Name: ca\_\allowbreak{}1\_\allowbreak{}a\_\allowbreak{}1\_\allowbreak{}a\_\allowbreak{}6\\
Type: string\\
Default: nil
}

{\footnotesize політика <CA-01\_\allowbreak{}ODP[03] ВИБРАНЕ ЗНАЧЕННЯ ПАРАМЕТРА(ів)> оцінювання, авторизації та моніторингу містить координацію між підрозділами організації;}


{\footnotesize\bfseries
No: 19\\
Name: ca\_\allowbreak{}1\_\allowbreak{}a\_\allowbreak{}1\_\allowbreak{}a\_\allowbreak{}7\\
Type: list\\
Default: [\textquotedbl{}admin\textquotedbl{}]
}

{\footnotesize політика <CA-01\_\allowbreak{}ODP[03] ВИБРАНЕ ЗНАЧЕННЯ ПАРАМЕТРА(ів)> оцінювання, авторизації та моніторингу містить систему контролю відповідності;}


{\footnotesize\bfseries
No: 20\\
Name: ca\_\allowbreak{}1\_\allowbreak{}a\_\allowbreak{}1\_\allowbreak{}b\\
Type: string\\
Default: nil
}

{\footnotesize <CA-01\_\allowbreak{}ODP[03] ВИБРАНЕ ЗНАЧЕННЯ ПАРАМЕТРА(ІВ)> політика оцінювання, надання дозволів та моніторингу відповідає чинним законам, нормативним документам, наказа м, положенням, політиці, стандартам і керівним принципам;}


{\footnotesize\bfseries
No: 21\\
Name: ca\_\allowbreak{}1\_\allowbreak{}b\\
Type: string\\
Default: nil
}

{\footnotesize <CA-01\_\allowbreak{}ODP[04] посадова особа > призначається для управління розробкою, документуванням та розповсюдженням політики та процедур оцінювання, авторизації та моніторингу;}


{\footnotesize\bfseries
No: 22\\
Name: ca\_\allowbreak{}1\_\allowbreak{}c\_\allowbreak{}1\_\allowbreak{}1\\
Type: string\\
Default: nil
}

{\footnotesize переглядається та оновлюється поточна політика оцінки, авторизації та моніторингу з <CA-01\_\allowbreak{}ODP[05] частотою>;}


{\footnotesize\bfseries
No: 23\\
Name: ca\_\allowbreak{}1\_\allowbreak{}c\_\allowbreak{}1\_\allowbreak{}2\\
Type: string\\
Default: nil
}

{\footnotesize переглядається та оновлюється поточна політика оцінки, авторизації та моніторингу після <CA-01\_\allowbreak{}ODP[06] подій>;}


{\footnotesize\bfseries
No: 24\\
Name: ca\_\allowbreak{}1\_\allowbreak{}c\_\allowbreak{}2\_\allowbreak{}1\\
Type: string\\
Default: nil
}

{\footnotesize переглядаються та оновлюються поточні процедури оцінки, авторизації та моніторингу з <CA-01\_\allowbreak{}ODP[07] частотою>;}


{\footnotesize\bfseries
No: 25\\
Name: ca\_\allowbreak{}1\_\allowbreak{}c\_\allowbreak{}2\_\allowbreak{}2\\
Type: string\\
Default: nil
}

{\footnotesize переглядаються та оновлюються поточні процедури оцінки, авторизації та моніторингу після <CA-01\_\allowbreak{}ODP[07] подій>; 157}




\subsection{ОЦІНЮВАННЯ (CA-2)}
a. Виберіть відповідного оцінювача або команду з оцінки для типу оцінювання, яке буде проводитися;\\ 
b. Розробіть план контрольної оцінки, який описує обсяг оцінки, в тому числі:\\ 
1. заходи захисту та посилені заходи, що підлягають оцінюванню;\\ 
2. процедури оцінювання, ефективності заходів; які використовуватимуться для визначення\\ 
3. середовище оцінювання, групу оцінювання, ролі й обов'язки з оцінювання.\\ 
c. Забезпечити розгляд і затвердження плану оцінювання уповноваженою офіційною особою або призначеним для проведення оцінювання представником;\\ 
d. Оцінити заходи захисту в системі та в її середовищі функціонування з [Призначення: визначеною організацією частотою] для визначення, наскільки коректно реалізовані заходи безпеки і чи працюють вони за призначенням і дають бажаний результат щодо дотримання встановлених вимог безпеки та приватності;\\ 
e. Підготовити звіт оцінювання безпеки, який документує результати оцінювання;\\ 
f. Надати результати оцінювання з безпеки [Призначення: особам або ролям, визначеним організацією].

\vspace{10pt}
{\footnotesize\bfseries
No: 1\\
Name: ca\_\allowbreak{}2\_\allowbreak{}odp\_\allowbreak{}1\\
Type: list\\
Default: [\textquotedbl{}admin\textquotedbl{}]
}

{\footnotesize визначено частоту, з якою слід оцінювати засоби контролю в системі та середовищі її функціонування;}


{\footnotesize\bfseries
No: 2\\
Name: ca\_\allowbreak{}2\_\allowbreak{}odp\_\allowbreak{}2\\
Type: list\\
Default: [\textquotedbl{}admin\textquotedbl{}]
}

{\footnotesize визначені особи або ролі, яким мають бути надані результати оцінювання з безпеки;}


{\footnotesize\bfseries
No: 3\\
Name: ca\_\allowbreak{}2\_\allowbreak{}a\\
Type: string\\
Default: nil
}

{\footnotesize обрано відповідного оцінювача або команду оцінювачів для проведення оцінювання;}


{\footnotesize\bfseries
No: 4\\
Name: ca\_\allowbreak{}2\_\allowbreak{}b\_\allowbreak{}1\\
Type: list\\
Default: [\textquotedbl{}admin\textquotedbl{}]
}

{\footnotesize розроблено план контрольної оцінки, який описує обсяг оцін ки, включаючи заходи захисту та посилені заходи, що підлягають оцінюванню.}


{\footnotesize\bfseries
No: 5\\
Name: ca\_\allowbreak{}2\_\allowbreak{}b\_\allowbreak{}2\\
Type: list\\
Default: [\textquotedbl{}admin\textquotedbl{}]
}

{\footnotesize розроблено план контрольної оцінки, який  описує обсяг оцінки, включаючи процедури оцінювання, які використовуватимуться для визначення ефективності заходів.}


{\footnotesize\bfseries
No: 6\\
Name: ca\_\allowbreak{}2\_\allowbreak{}b\_\allowbreak{}3\_\allowbreak{}1\\
Type: string\\
Default: nil
}

{\footnotesize розроблено план контро льної оцінки, який описує обсяг оцінки, включаючи середовище оцінювання.}


{\footnotesize\bfseries
No: 7\\
Name: ca\_\allowbreak{}2\_\allowbreak{}b\_\allowbreak{}3\_\allowbreak{}2\\
Type: list\\
Default: [\textquotedbl{}admin\textquotedbl{}]
}

{\footnotesize розроблено план контрольної оцінки, який описує обсяг оцінки, включаючи групу оцінювання.}


{\footnotesize\bfseries
No: 8\\
Name: ca\_\allowbreak{}2\_\allowbreak{}b\_\allowbreak{}3\_\allowbreak{}3\\
Type: list\\
Default: [\textquotedbl{}admin\textquotedbl{}]
}

{\footnotesize розроблено план контрольної оцінки, який описує обсяг оцінки, включаючи ролі й обов'язки з оцінювання.}


{\footnotesize\bfseries
No: 9\\
Name: ca\_\allowbreak{}2\_\allowbreak{}c\\
Type: string\\
Default: nil
}

{\footnotesize план оцінки заходів захисту розглядається та затверджується уповноваженою посадовою особою або призначеним представником перед проведенням оцінки;}


{\footnotesize\bfseries
No: 10\\
Name: ca\_\allowbreak{}2\_\allowbreak{}d\_\allowbreak{}1\\
Type: string\\
Default: nil
}

{\footnotesize заходи захисту оцінюються в системі та середовищі її визначити, наскільки коректно реалізовані заходи захисту, чи працюють вони за призначенням і чи дають бажаний результат щодо дотримання встановлених вимог до безпеки;}


{\footnotesize\bfseries
No: 11\\
Name: ca\_\allowbreak{}2\_\allowbreak{}d\_\allowbreak{}2\\
Type: string\\
Default: nil
}

{\footnotesize заходи захисту оцінюються в системі та середовищі її визначити, наскільки коректно реалізовані заходи захисту, чи працюють вони за призначенням і чи  дають бажаний результат щодо дотримання встановлених вимог до конфіденційності; 158}


{\footnotesize\bfseries
No: 12\\
Name: ca\_\allowbreak{}2\_\allowbreak{}e\\
Type: string\\
Default: nil
}

{\footnotesize готується звіт оцінювання , який документує результати оцінювання;}


{\footnotesize\bfseries
No: 13\\
Name: ca\_\allowbreak{}2\_\allowbreak{}f\\
Type: list\\
Default: [\textquotedbl{}admin\textquotedbl{}]
}

{\footnotesize результати оцінювання з безпеки надаються <CA-02\_\allowbreak{}ODP[02] особам або ролям>.}




\subsubsection{НЕЗАЛЕЖНІ ЕКСПЕРТИ (CA-2(1))}


\vspace{10pt}
{\footnotesize\bfseries
No: 1\\
Name: ca\_\allowbreak{}2\_\allowbreak{}1\_\allowbreak{}01\\
Type: list\\
Default: [\textquotedbl{}admin\textquotedbl{}]
}

{\footnotesize Для проведення контрольних оцінок залучаються незалежні експерти або групи експертів.}




\subsubsection{СПЕЦІАЛІЗОВАНІ ОЦІНКИ (CA-2(2))}
Ввести як частину оцінювання заходів безпеки та приватності, [Призначення: з визначеною організацією частотою], [Вибір: з попередженням; без попередження], [Вибір (один або кілька): поглиблений моніторинг; сканування уразливостей; тестування на шкідливих користувачів; оцінювання внутрішньої загрози; тестування продуктивності та навантаження; [Призначення: організаційно визначені інші форми оцінювання]].

\vspace{10pt}
{\footnotesize\bfseries
No: 1\\
Name: ca\_\allowbreak{}2\_\allowbreak{}2\_\allowbreak{}odp\_\allowbreak{}1\\
Type: integer\\
Default: 30
}

{\footnotesize визначено частоту, з якою слід включати спеціалізовані оцінки як частину оцінювання безпеки та конфіденційності;}


{\footnotesize\bfseries
No: 2\\
Name: ca\_\allowbreak{}2\_\allowbreak{}2\_\allowbreak{}odp\_\allowbreak{}2\\
Type: string\\
Default: nil
}

{\footnotesize вибрано одне з наступних ЗНАЧЕНЬ ПАРАМЕТРІВ: \{з попередженням; без попередження\};}


{\footnotesize\bfseries
No: 3\\
Name: ca\_\allowbreak{}2\_\allowbreak{}2\_\allowbreak{}odp\_\allowbreak{}3\\
Type: list\\
Default: [\textquotedbl{}admin\textquotedbl{}]
}

{\footnotesize вибрано одне або декілька з наступних ЗНАЧЕНЬ ПАРАМЕТРІВ: \{поглиблений моніторинг; сканування уразливостей; тестування на шкідливих користувачів; оцінювання внутрішньої загрози; тестування продуктивності та навантаження; \};}


{\footnotesize\bfseries
No: 4\\
Name: ca\_\allowbreak{}2\_\allowbreak{}2\_\allowbreak{}odp\_\allowbreak{}4\\
Type: string\\
Default: nil
}

{\footnotesize визначаються інші форми оцінювання (якщо вони були обрані); 159}


{\footnotesize\bfseries
No: 5\\
Name: ca\_\allowbreak{}2\_\allowbreak{}2\_\allowbreak{}01\\
Type: string\\
Default: nil
}

{\footnotesize <CA-02(02)\_\allowbreak{}ODP[01] періодичність оцінювання>  <CA02(02)\_\allowbreak{}ODP[02] ВИБРАНЕ ЗНАЧЕННЯ ПАРАМЕТРА> ПАРАМЕТРА(ів)> включаються як частина оцінювання заходів безпеки та конфіденційності;}




\subsubsection{ЗОВНІШНІ ОРГАНІЗАЦІЇ (CA-2(3))}
Використовуйте результати контрольного оцінювання, які виконує [Призначення: зовнішня організація, визначена організацією] на [Призначення: система, визначена організацією], коли оцінювання відповідає [Завдання: вимоги, визначені організацією].

\vspace{10pt}
{\footnotesize\bfseries
No: 1\\
Name: ca\_\allowbreak{}2\_\allowbreak{}3\_\allowbreak{}odp\_\allowbreak{}1\\
Type: string\\
Default: nil
}

{\footnotesize визначено організацію, яка надає результати оцінок заходів захисту інформації та персональних даних}


{\footnotesize\bfseries
No: 2\\
Name: ca\_\allowbreak{}2\_\allowbreak{}3\_\allowbreak{}odp\_\allowbreak{}2\\
Type: string\\
Default: nil
}

{\footnotesize визначено систему, яка приймає результати оцінок заходів захисту інформації та персональних даних}


{\footnotesize\bfseries
No: 3\\
Name: ca\_\allowbreak{}2\_\allowbreak{}3\_\allowbreak{}odp\_\allowbreak{}3\\
Type: string\\
Default: nil
}

{\footnotesize визначено вимоги до результатів оцінок заходів захисту інформації та персональних даних}


{\footnotesize\bfseries
No: 4\\
Name: ca\_\allowbreak{}2\_\allowbreak{}3\_\allowbreak{}01\\
Type: string\\
Default: nil
}

{\footnotesize прийняти результати оцінок заходів захисту інформації та}




\subsection{ВЗАЄМОДІЯ СИСТЕМ (CA-3)}
a. схвалити та керувати обміном інформацією між системою та іншими системами за допомогою [Вибір (один або кілька): угоди безпеки взаємозв'язку; договори безпеки обміну інформацією; меморандуми про взаєморозуміння; угоди про рівень обслуговування; угоди користувача; угоди про нерозголошення; [Доручення: тип договору, визначений організацією]];.\\ 
b. документувати, як частину угоди про обмін, характеристики інтерфейсу, вимоги до безпеки та приватності, засоби контролю та відповідальність для кожної системи, а також характер переданої інформації;\\ 
c. здійснювати перегляд та оновлення угод з [Призначення: визначеною організацією частотою].

\vspace{10pt}
{\footnotesize\bfseries
No: 1\\
Name: ca\_\allowbreak{}3\_\allowbreak{}odp\_\allowbreak{}1\\
Type: list\\
Default: [\textquotedbl{}admin\textquotedbl{}]
}

{\footnotesize вибрано одне або декілька з наступних ЗНАЧЕНЬ ПАРА- МЕТРІВ: \{): угоди безпеки взаємозв'язку; договори безпеки обміну інформацією; меморандуми про взаєморозуміння; 160 угоди про рівень обслуговування; угоди користувача; угоди}


{\footnotesize\bfseries
No: 2\\
Name: ca\_\allowbreak{}3\_\allowbreak{}odp\_\allowbreak{}2\\
Type: string\\
Default: nil
}

{\footnotesize визначено тип угоди, який використовується для схвалення та керування обміном інформацією (якщо вибрано);}


{\footnotesize\bfseries
No: 3\\
Name: ca\_\allowbreak{}3\_\allowbreak{}odp\_\allowbreak{}3\\
Type: integer\\
Default: 30
}

{\footnotesize визначено частоту, з якою необхідно переглядати та оновлювати угоди;}


{\footnotesize\bfseries
No: 4\\
Name: ca\_\allowbreak{}3\_\allowbreak{}a\\
Type: string\\
Default: nil
}

{\footnotesize обмін інформацією між системою та іншими системами схвалюється та керується за допомогою  <CA-03\_\allowbreak{}ODP[01] ВИБРАНЕ ЗНАЧЕННЯ ПАРАМЕТРА(ів)>;}


{\footnotesize\bfseries
No: 5\\
Name: ca\_\allowbreak{}3\_\allowbreak{}b\_\allowbreak{}1\\
Type: string\\
Default: nil
}

{\footnotesize характеристики інтерфейсу задокументовані як частина кожної угоди про обмін;}


{\footnotesize\bfseries
No: 6\\
Name: ca\_\allowbreak{}3\_\allowbreak{}b\_\allowbreak{}2\\
Type: string\\
Default: nil
}

{\footnotesize вимоги до безпеки задокументовані як частина кожної угоди про обмін;}


{\footnotesize\bfseries
No: 7\\
Name: ca\_\allowbreak{}3\_\allowbreak{}b\_\allowbreak{}3\\
Type: string\\
Default: nil
}

{\footnotesize вимоги щодо конфіденційності задокументовані як частина кожної угоди про обмін;}


{\footnotesize\bfseries
No: 8\\
Name: ca\_\allowbreak{}3\_\allowbreak{}b\_\allowbreak{}4\\
Type: string\\
Default: nil
}

{\footnotesize заходи захисту задокументовані як частина кожної угоди про обмін;}


{\footnotesize\bfseries
No: 9\\
Name: ca\_\allowbreak{}3\_\allowbreak{}b\_\allowbreak{}5\\
Type: string\\
Default: nil
}

{\footnotesize відповідальність за кожну сист ему задокументована як частина кожної угоди про обмін;}


{\footnotesize\bfseries
No: 10\\
Name: ca\_\allowbreak{}3\_\allowbreak{}b\_\allowbreak{}6\\
Type: string\\
Default: nil
}

{\footnotesize характер переданої інформації документується як частина кожної угоди про обмін;}


{\footnotesize\bfseries
No: 11\\
Name: ca\_\allowbreak{}3\_\allowbreak{}c\\
Type: string\\
Default: nil
}

{\footnotesize угоди переглядаються та оновлюються  <CA-03\_\allowbreak{}ODP[03] частота>.}




\subsubsection{НЕЗАХИЩЕНІ З'ЄДНАННЯ СИСТЕМИ (CA-3(1)) [Вилучено]}
[Вилучено: Включено до SC-07(25)]

\vspace{10pt}
\textit{Немає параметрів для цього контролю.}
\vspace{10pt}


\subsubsection{ЗАХИЩЕНІ З'ЄДНАННЯ СИСТЕМИ (CA-3(2))}


\vspace{10pt}
\textit{Немає параметрів для цього контролю.}
\vspace{10pt}


\subsubsection{НЕСЕКРЕТНІ З'ЄДНАННЯ СИСТЕМИ БЕЗПЕКИ, ЩО НЕ Є НАЦІОНАЛЬНИМИ (CA-3(3)) [Вилучено]}
[Вилучено: Включено до SC-07(27)]

\vspace{10pt}
\textit{Немає параметрів для цього контролю.}
\vspace{10pt}


\subsubsection{ПІДКЛЮЧЕННЯ ДО ЗАГАЛЬНОДОСТУПНИХ МЕРЕЖ (CA-3(4)) [Вилучено]}
[Вилучено: Включено до SC-07(28)]

\vspace{10pt}
\textit{Немає параметрів для цього контролю.}
\vspace{10pt}


\subsubsection{ОБМЕЖЕННЯ ЗВ'ЯЗКУ ІЗ ЗОВНІШНІМИ СИСТЕМАМИ (CA-3(5)) [Вилучено]}
[Вилучено: Включено до SC-07(05)]

\vspace{10pt}
\textit{Немає параметрів для цього контролю.}
\vspace{10pt}


\subsubsection{ПЕРЕДАЧА ДОЗВОЛІВ (CA-3(6))}
Переконатися, що особи або системи, які передають дані між взаємопов'язаними системами, мають необхідні повноваження (тобто дозволи на запис або привілеї), до прийняття таких даних.

\vspace{10pt}
{\footnotesize\bfseries
No: 1\\
Name: ca\_\allowbreak{}3\_\allowbreak{}6\_\allowbreak{}01\\
Type: string\\
Default: nil
}

{\footnotesize особи або системи, які передають дані між системами, що з'єднуються, мають необхідні повноваження (тобто дозволи на запис або привілеї) перед тим, як приймати такі дані.}




\subsubsection{ТРАНЗИТИВНИЙ ОБМІН ІНФОРМАЦІЄЮ (CA-3(7))}
a) Визначити транзитивний (низхідний) обмін інформацією з іншими системами через системи, визначені в CA-3a; b) Вжити заходів для забезпечення припинення транзитивного (низхідного) обміну інформацією, коли засоби контролю ідентифікованих транзитивних (низхідних) систем не можуть бути перевірені або підтверджені.

\vspace{10pt}
{\footnotesize\bfseries
No: 1\\
Name: ca\_\allowbreak{}3\_\allowbreak{}7\_\allowbreak{}a\\
Type: string\\
Default: nil
}

{\footnotesize визначено транзитний  обмін інформацією з іншими системами через системи, визначені в CA-03a;}


{\footnotesize\bfseries
No: 2\\
Name: ca\_\allowbreak{}3\_\allowbreak{}7\_\allowbreak{}b\\
Type: list\\
Default: [\textquotedbl{}admin\textquotedbl{}]
}

{\footnotesize вживаються заходи для забезпечення припинення транзитного  обміну інформацією, коли засоби контролю над ідентифікованими транзитними  системами не можуть бути перевірені або підтверджені.}




\subsection{СЕРТИФІКАЦІЯ БЕЗПЕКИ (CA-4) [Вилучено]}
[Вилучено: Включено до CA-02]

\vspace{10pt}
\textit{Немає параметрів для цього контролю.}
\vspace{10pt}


\subsection{ПЛАН УСУНЕННЯ НЕДОЛІКІВ ТА КОНТРОЛЬНІ ПОКАЗНИКИ (CA-5)}
a. Розробити для системи план усунення недоліків та контрольні показники з метою документування запланованих коригувальних дій організації для усунення недоліків і зауважень, які виявлені в ході оцінювання заходів захисту, а також для зменшення або усунення відомих вразливостей у системі.\\ 
b. Оновлювати чинний план усунення недоліків та контрольні показники з [Призначення: визначеною організацією частотою] на основі результатів оцінювання заходів, незалежних аудитів та постійного моніторингу.

\vspace{10pt}
{\footnotesize\bfseries
No: 1\\
Name: ca\_\allowbreak{}5\_\allowbreak{}odp\\
Type: list\\
Default: [\textquotedbl{}admin\textquotedbl{}]
}

{\footnotesize визначено частоту оновлення чинного плану усунення недоліків та контрольних показників на основі результатів оцінювання заходів захисту, незалежних аудитів та постійного моніторингу;}


{\footnotesize\bfseries
No: 2\\
Name: ca\_\allowbreak{}5\_\allowbreak{}a\\
Type: list\\
Default: [\textquotedbl{}admin\textquotedbl{}]
}

{\footnotesize розроблено план усунення недоліків та контрольні показники для системи, та задокументувано заплановані дії організації з коригування, спрямовані на усунення недоліків та зауважень, виявлених під час оцінки заходів захисту, а також на зменшення або усунення відомих вразливостей в системі;}


{\footnotesize\bfseries
No: 3\\
Name: ca\_\allowbreak{}5\_\allowbreak{}b\\
Type: string\\
Default: nil
}

{\footnotesize існуючий план оновлюються <CA-05\_\allowbreak{}ODP частота> на основі результатів оцінювання заходів, незалежних аудитів та постійного моніторингу.}




\subsubsection{АВТОМАТИЗАЦІЯ ПІДТРИМКИ ЗАДЛЯ ТОЧНОСТІ ТА ВЖИВАНОСТІ (CA-5(1))}


\vspace{10pt}
{\footnotesize\bfseries
No: 1\\
Name: ca\_\allowbreak{}5\_\allowbreak{}1\_\allowbreak{}odp\\
Type: string\\
Default: nil
}

{\footnotesize визначено автоматизовані механізми, які використовуються для забезпечення точності, актуальності та доступності плану усунення недоліків і основних етапів для системи;}


{\footnotesize\bfseries
No: 2\\
Name: ca\_\allowbreak{}5\_\allowbreak{}1\_\allowbreak{}01\\
Type: string\\
Default: nil
}

{\footnotesize <CA-05(01)\_\allowbreak{}ODP автоматизовані механізми>  використовуються для забезпечення точності, актуальності та доступності плану усунення недоліків і основних етапів для системи.}




\subsection{АКРЕДИТАЦІЯ (CA-6)}
a. Призначити старшого керівника, який відповідає за систему;\\ 
b. Призначити старшого керівника, відповідального за систему, та будь-які загальні заходи захисту, успадковані системою.\\ 
c. Переконатися перед початком функціонування системи, що посадова особа:\\ 
1. акредитує загальні заходи захисту, що успадковані системою;\\ 
2. акредитує систему на функціонування за призначенням.\\ 
d. Переконайтеся, що посадова особа, яка акредитує засоби захисту, дозволяє використання цих засобів захисту для успадкування організаційними системами;\\ 
e. Оновлювати акредитацію [Призначення: з визначеною організацією частотою].

\vspace{10pt}
{\footnotesize\bfseries
No: 1\\
Name: ca\_\allowbreak{}6\_\allowbreak{}odp\\
Type: integer\\
Default: 30
}

{\footnotesize визначено частоту, з якою потрібно оновлювати акредитації;}


{\footnotesize\bfseries
No: 2\\
Name: ca\_\allowbreak{}6\_\allowbreak{}a\\
Type: string\\
Default: nil
}

{\footnotesize призначено старшого керівника, який відповідає за систему; 164}


{\footnotesize\bfseries
No: 3\\
Name: ca\_\allowbreak{}6\_\allowbreak{}b\\
Type: string\\
Default: nil
}

{\footnotesize призначено старшого керівника, відповідального за систему, та будь-які загальні заходи захисту, успадковані системою;}


{\footnotesize\bfseries
No: 4\\
Name: ca\_\allowbreak{}6\_\allowbreak{}c\_\allowbreak{}1\\
Type: string\\
Default: nil
}

{\footnotesize перед початком функціонування системи посадова особа, яка відповідає за систему, акредитує загальні заходи захисту, що успадковані системою;}


{\footnotesize\bfseries
No: 5\\
Name: ca\_\allowbreak{}6\_\allowbreak{}c\_\allowbreak{}2\\
Type: string\\
Default: nil
}

{\footnotesize перед початком функціонування системи посадова особа, яка відповідає за систему, акредитує систему на функціонування за призначенням;}


{\footnotesize\bfseries
No: 6\\
Name: ca\_\allowbreak{}6\_\allowbreak{}d\\
Type: string\\
Default: nil
}

{\footnotesize посадова особа, яка акредитує заходи захисту, дозволяє використання цих заходів захисту для успадкування системами організації;}


{\footnotesize\bfseries
No: 7\\
Name: ca\_\allowbreak{}6\_\allowbreak{}e\\
Type: string\\
Default: nil
}

{\footnotesize акредитації оновлюються <CA-06\_\allowbreak{}ODP частота>.}




\subsubsection{СПІЛЬНА АКРЕДИТАЦІЯ - ОДНА І ТА САМА ОРГАНІЗАЦІЯ (CA-6(1))}


\vspace{10pt}
{\footnotesize\bfseries
No: 1\\
Name: ca\_\allowbreak{}6\_\allowbreak{}1\_\allowbreak{}1\\
Type: string\\
Default: nil
}

{\footnotesize для системи впроваджено спільний процес акредитації;}


{\footnotesize\bfseries
No: 2\\
Name: ca\_\allowbreak{}6\_\allowbreak{}1\_\allowbreak{}2\\
Type: string\\
Default: nil
}

{\footnotesize спільний процес акредитації, який використовується в системі, включає в себе кілька посадових осіб  з однієї організації , які надають акредитацію.}




\subsubsection{СПІЛЬНА АКРЕДИТАЦІЯ - РІЗНІ ОРГАНІЗАЦІЇ (CA-6(2))}
Впровадити спільний процес акредитації для системи, що має кількох уповноважених посадових осіб з принаймні однією уповноваженою посадовою особою з організації, яка є зовнішньою організацією, що здійснює акредитацію.

\vspace{10pt}
{\footnotesize\bfseries
No: 1\\
Name: ca\_\allowbreak{}6\_\allowbreak{}2\_\allowbreak{}1\\
Type: string\\
Default: nil
}

{\footnotesize для системи впроваджено спільний процес акредитації;}


{\footnotesize\bfseries
No: 2\\
Name: ca\_\allowbreak{}6\_\allowbreak{}2\_\allowbreak{}2\\
Type: string\\
Default: nil
}

{\footnotesize спільний процес акредитації, що використовується в системі, передбачає наявність кількох посадових осіб, які надають акредитації, принаймні одна з яких є посадовою особою з організації, що не належить до організації, яка здійснює акредитацію.}




\subsection{БЕЗПЕРЕРВНИЙ МОНІТОРИНГ (CA-7)}
Розробити стратегію безперервного моніторингу безпеки та приватності й упровадити програму безперервного моніторингу безпеки та приватності, яка охоплює:\\ 
a. встановлення показників безпеки та приватності, які необхідно відстежувати: [Призначення: визначені організацією метрики];\\ 
b. встановлення [Призначення: визначена організацією частота] для моніторингу та [Призначення: визначена організацією частота] для безперервного оцінювання ефективності заходів захисту;\\ 
c. поточні оцінювання заходів захисту відповідно до стратегії безперервного моніторингу організації;\\ 
d. постійний моніторинг стану безпеки та приватності відповідно до встановлених організацією метрик і відповідно до стратегії безперервного моніторингу організації;\\ 
e. зіставлення та аналіз інформації, отриманої в результаті оцінювання та моніторингу безпеки та приватності;\\ 
f. дії реагування за результатами аналізу інформації, пов'язаної з безпекою та приватністю;\\ 
g. повідомлення про статус безпеки та приватності системи [Призначення: визначені організацією персонал або ролі] з [Призначення: визначеною організацією частотою].

\vspace{10pt}
{\footnotesize\bfseries
No: 1\\
Name: ca\_\allowbreak{}7\_\allowbreak{}odp\_\allowbreak{}1\\
Type: string\\
Default: nil
}

{\footnotesize визначено метрики системного рівня, які підлягають моніторингу;}


{\footnotesize\bfseries
No: 2\\
Name: ca\_\allowbreak{}7\_\allowbreak{}odp\_\allowbreak{}2\\
Type: integer\\
Default: 30
}

{\footnotesize визначено частоту, з якою слід моніторити ефективність заходів захисту;}


{\footnotesize\bfseries
No: 3\\
Name: ca\_\allowbreak{}7\_\allowbreak{}odp\_\allowbreak{}3\\
Type: integer\\
Default: 30
}

{\footnotesize визначено частоту, з якою слід оцінювати ефективність заходів захисту;}


{\footnotesize\bfseries
No: 4\\
Name: ca\_\allowbreak{}7\_\allowbreak{}odp\_\allowbreak{}4\\
Type: list\\
Default: [\textquotedbl{}admin\textquotedbl{}]
}

{\footnotesize визначено персонал або ролі, яким повідомляється про стан безпеки системи;}


{\footnotesize\bfseries
No: 5\\
Name: ca\_\allowbreak{}7\_\allowbreak{}odp\_\allowbreak{}5\\
Type: integer\\
Default: 30
}

{\footnotesize визначено частоту, з якою повідомляється про стан безпеки системи;}


{\footnotesize\bfseries
No: 6\\
Name: ca\_\allowbreak{}7\_\allowbreak{}odp\_\allowbreak{}6\\
Type: list\\
Default: [\textquotedbl{}admin\textquotedbl{}]
}

{\footnotesize визначено персонал або ролі, яким повідомляється про стан конфіденційності системи;}


{\footnotesize\bfseries
No: 7\\
Name: ca\_\allowbreak{}7\_\allowbreak{}odp\_\allowbreak{}7\\
Type: integer\\
Default: 30
}

{\footnotesize визначено частоту, з якою повідомляється про стан конфіденційності системи;}


{\footnotesize\bfseries
No: 8\\
Name: ca\_\allowbreak{}7\_\allowbreak{}1\\
Type: string\\
Default: nil
}

{\footnotesize розроблено стратегію безперервного моніторингу на системному рівні; 166}


{\footnotesize\bfseries
No: 9\\
Name: ca\_\allowbreak{}7\_\allowbreak{}2\\
Type: string\\
Default: nil
}

{\footnotesize безперервний моніторинг на рівні системи здійснюється відповідно до стратегії безперервного моніторингу на рівні організації;}


{\footnotesize\bfseries
No: 10\\
Name: ca\_\allowbreak{}7\_\allowbreak{}a\\
Type: string\\
Default: nil
}

{\footnotesize безперервний моніторинг на рівні системи включає встановлення наступних метрик на рівні системи, які підлягають моніторингу:}


{\footnotesize\bfseries
No: 11\\
Name: ca\_\allowbreak{}7\_\allowbreak{}b\_\allowbreak{}1\\
Type: string\\
Default: nil
}

{\footnotesize безперервний моніторинг на рівні системи включає встановлені захисту;}


{\footnotesize\bfseries
No: 12\\
Name: ca\_\allowbreak{}7\_\allowbreak{}b\_\allowbreak{}2\\
Type: string\\
Default: nil
}

{\footnotesize безперервний моніторинг на рівні системи включає встановлені}


{\footnotesize\bfseries
No: 13\\
Name: ca\_\allowbreak{}7\_\allowbreak{}c\\
Type: list\\
Default: [\textquotedbl{}admin\textquotedbl{}]
}

{\footnotesize безперервний моніторинг на рівні системи включає поточні контрольні оцінки відповідно до стратегії безперервного моніторингу;}


{\footnotesize\bfseries
No: 14\\
Name: ca\_\allowbreak{}7\_\allowbreak{}d\\
Type: string\\
Default: nil
}

{\footnotesize безперервний моніторинг на рівні системи включає постійний моніторинг визначених системою та організацією показників відповідно до стратегії безперервного моніторингу;}


{\footnotesize\bfseries
No: 15\\
Name: ca\_\allowbreak{}7\_\allowbreak{}e\\
Type: string\\
Default: nil
}

{\footnotesize безперервний моніторинг на рівні системи включає зіставлення та аналіз інформації, отриманої в результаті оцінювання та моніторингу;}


{\footnotesize\bfseries
No: 16\\
Name: ca\_\allowbreak{}7\_\allowbreak{}f\\
Type: string\\
Default: nil
}

{\footnotesize безперервний моніторинг на рівні системи включає в себе дії з реагування на результати аналізу інформації, пов'язаної з безпекою та приватністю;}


{\footnotesize\bfseries
No: 17\\
Name: ca\_\allowbreak{}7\_\allowbreak{}g\_\allowbreak{}1\\
Type: string\\
Default: nil
}

{\footnotesize безперервний моніторинг на рівні системи включає повідомлення}


{\footnotesize\bfseries
No: 18\\
Name: ca\_\allowbreak{}7\_\allowbreak{}g\_\allowbreak{}2\\
Type: string\\
Default: nil
}

{\footnotesize безперервний моніторинг на рівні системи включає повідомлення}




\subsubsection{НЕЗАЛЕЖНЕ ОЦІНЮВАННЯ (CA-7(1))}


\vspace{10pt}
{\footnotesize\bfseries
No: 1\\
Name: ca\_\allowbreak{}7\_\allowbreak{}1\_\allowbreak{}01\\
Type: string\\
Default: nil
}

{\footnotesize для постійного моніторингу заходів захисту в системі залучаються незалежні експертів або групи з оцінювання.}




\subsubsection{ВИДИ ОЦІНОК (CA-7(2)) [Вилучено]}
[Вилучено: Включено до CA-02]

\vspace{10pt}
\textit{Немає параметрів для цього контролю.}
\vspace{10pt}


\subsubsection{АНАЛІЗ ТЕНДЕНЦІЇ (CA-7(3))}
Впровадити аналіз тенденцій, щоб визначити, чи потрібно змінювати реалізацію заходу захисту, частоту постійних моніторингових заходів і види діяльності, що використовуються в процесі безперервного моніторингу, на основі емпіричних даних.

\vspace{10pt}
{\footnotesize\bfseries
No: 1\\
Name: ca\_\allowbreak{}7\_\allowbreak{}3\_\allowbreak{}1\\
Type: string\\
Default: nil
}

{\footnotesize аналіз тенденцій використовується для визначення того, чи потрібно змінювати реалізацію заходів захисту, які використовуються в процесі безперервного  моніторингу, на основі емпіричних даних;}


{\footnotesize\bfseries
No: 2\\
Name: ca\_\allowbreak{}7\_\allowbreak{}3\_\allowbreak{}2\\
Type: string\\
Default: nil
}

{\footnotesize аналіз тенденцій застосовується для того, щоб на основі емпіричних даних визначити, чи потрібно змінювати частоту постійного моніторингу;}


{\footnotesize\bfseries
No: 3\\
Name: ca\_\allowbreak{}7\_\allowbreak{}3\_\allowbreak{}3\\
Type: string\\
Default: nil
}

{\footnotesize аналіз тенденцій застосовується для того, щоб на основі емпіричних даних визначити, чи потрібно змінювати види діяльності, які використовуються в процесі безперервного моніторингу.}




\subsubsection{МОНІТОРИНГ РИЗИКУ (CA-7(4))}
Забезпечити моніторинг ризиків, що є невід'ємною частиною стратегії постійного моніторингу та включає:\\ 
(a) моніторинг ефективності;\\ 
(b) моніторинг відповідності;\\ 
(c) моніторинг змін.

\vspace{10pt}
{\footnotesize\bfseries
No: 1\\
Name: ca\_\allowbreak{}7\_\allowbreak{}4\_\allowbreak{}01\\
Type: string\\
Default: nil
}

{\footnotesize моніторинг ризиків є невід'ємною частиною стратегії безперервного моніторингу;}


{\footnotesize\bfseries
No: 2\\
Name: ca\_\allowbreak{}7\_\allowbreak{}4\_\allowbreak{}a\\
Type: string\\
Default: nil
}

{\footnotesize моніторинг ефективності включено до моніторингу ризиків;}


{\footnotesize\bfseries
No: 3\\
Name: ca\_\allowbreak{}7\_\allowbreak{}4\_\allowbreak{}b\\
Type: string\\
Default: nil
}

{\footnotesize моніторинг відповідності включено до моніторингу ризиків;}


{\footnotesize\bfseries
No: 4\\
Name: ca\_\allowbreak{}7\_\allowbreak{}4\_\allowbreak{}c\\
Type: string\\
Default: nil
}

{\footnotesize моніторинг змін включений до моніторингу ризиків.}




\subsubsection{УЗГОДЖЕНИЙ АНАЛІЗ (CA-7(5))}
Застосуйте наступні дії, щоб перевірити, що політики встановлені, а запроваджені заходи захисту працюють узгоджено: [Призначення: дії, визначені організацією].

\vspace{10pt}
{\footnotesize\bfseries
No: 1\\
Name: ca\_\allowbreak{}7\_\allowbreak{}5\_\allowbreak{}odp\_\allowbreak{}1\\
Type: string\\
Default: nil
}

{\footnotesize визначені дії для підтвердження того, що політики встановлені;}


{\footnotesize\bfseries
No: 2\\
Name: ca\_\allowbreak{}7\_\allowbreak{}5\_\allowbreak{}odp\_\allowbreak{}2\\
Type: string\\
Default: nil
}

{\footnotesize визначені дії для підтвердження того, що впроваджені заходи захисту працюють узгоджено;}


{\footnotesize\bfseries
No: 3\\
Name: ca\_\allowbreak{}7\_\allowbreak{}5\_\allowbreak{}1\\
Type: string\\
Default: nil
}

{\footnotesize <CA-07(05)\_\allowbreak{}ODP[01] дії>  використовуються для перевірки того, що політики встановлено;}


{\footnotesize\bfseries
No: 4\\
Name: ca\_\allowbreak{}7\_\allowbreak{}5\_\allowbreak{}2\\
Type: string\\
Default: nil
}

{\footnotesize <CA-07(05)\_\allowbreak{}ODP[02] дії>  використовуються для перевірки того, що впроваджені заходи захисту працюють узгоджено}




\subsubsection{БЕЗПЕРЕРВНИЙ МОНІТОРИНГУ (CA-7(6))}
Забезпечити точність, актуальність і доступність результатів моніторингу для системи за допомогою [Завдання: автоматизовані механізми, визначені організацією]

\vspace{10pt}
{\footnotesize\bfseries
No: 1\\
Name: ca\_\allowbreak{}7\_\allowbreak{}6\_\allowbreak{}odp\\
Type: string\\
Default: nil
}

{\footnotesize визначено автоматизовані механізми забезпечення точності, актуальності та доступності результатів моніторингу системи;}


{\footnotesize\bfseries
No: 2\\
Name: ca\_\allowbreak{}7\_\allowbreak{}6\_\allowbreak{}01\\
Type: string\\
Default: nil
}

{\footnotesize <CA-07(06)\_\allowbreak{}ODP автоматизовані механізми > використовуються для забезпечення точності, актуальності та доступності результатів моніторингу системи.}




\subsection{ТЕСТУВАННЯ НА ПРОНИКНЕННЯ (CA-8)}
Проводити тестування на проникнення з [Призначення: визначеною організацією частотою] у [Призначення: визначеній організацією інформаційній системі чи системному компоненті].

\vspace{10pt}
{\footnotesize\bfseries
No: 1\\
Name: ca\_\allowbreak{}8\_\allowbreak{}odp\_\allowbreak{}1\\
Type: integer\\
Default: 30
}

{\footnotesize визначено частоту з якою проводить тестування на проникнення.}


{\footnotesize\bfseries
No: 2\\
Name: ca\_\allowbreak{}8\_\allowbreak{}odp\_\allowbreak{}2\\
Type: string\\
Default: nil
}

{\footnotesize визначено систему у якій проводить тестування на проникнення}


{\footnotesize\bfseries
No: 3\\
Name: ca\_\allowbreak{}8\_\allowbreak{}01\\
Type: string\\
Default: nil
}

{\footnotesize проводиться тестування на проникнення з <CA-08\_\allowbreak{}ODP[01] частотою> у <CA-08\_\allowbreak{}ODP[02] системі>.}




\subsubsection{НЕЗАЛЕЖНА КОМАНДА АБО АГЕНТ НА ПРОНИКНЕННЯ (CA-8(1))}


\vspace{10pt}
{\footnotesize\bfseries
No: 1\\
Name: ca\_\allowbreak{}8\_\allowbreak{}1\_\allowbreak{}01\\
Type: string\\
Default: nil
}

{\footnotesize для проведення тестування на проникнення в систему або компонентів системи залучається незалежний агент або команда з тестування на проникнення.}




\subsubsection{ЧЕРВОНА КОМАНДА (CA-8(2))}
Використовувати наступні вправи червоної команди, для імітації спроб супротивників скомпрометувати системи організації відповідно до прийнятих правил ведення бойових дій: [Завдання: визначені організацією вправи червоної команди].

\vspace{10pt}
{\footnotesize\bfseries
No: 1\\
Name: ca\_\allowbreak{}8\_\allowbreak{}2\_\allowbreak{}odp\\
Type: string\\
Default: nil
}

{\footnotesize визначено вправи  червоної команди для імітації спроби супротивників скомпрометувати системи організації;}


{\footnotesize\bfseries
No: 2\\
Name: ca\_\allowbreak{}8\_\allowbreak{}2\_\allowbreak{}01\\
Type: string\\
Default: nil
}

{\footnotesize залучити <CA-08(02)\_\allowbreak{}ODP вправи червоної команди>, щоб імітувати спроби супротивників скомпрометувати системи організації.}




\subsubsection{МОЖЛИВОСТІ ПЕРЕВІРКИ НА ПРОНИКНЕННЯ (CA-8(3))}
Впровадити процес тестування на проникнення, який охоплює [Призначення: визначену організацією частоту] [Вибір: з попередженням; без попередження] спроб обійти чи зламати заходи захисту, пов'язані з фізичними точками доступу до об'єкта.

\vspace{10pt}
{\footnotesize\bfseries
No: 1\\
Name: ca\_\allowbreak{}8\_\allowbreak{}3\_\allowbreak{}odp\_\allowbreak{}1\\
Type: integer\\
Default: 30
}

{\footnotesize визначено частоту спроби обійти або зламати заходи захисту, пов'язані з фізичними точками доступу до об'єкта в тестуванні на проникнення}


{\footnotesize\bfseries
No: 2\\
Name: ca\_\allowbreak{}8\_\allowbreak{}3\_\allowbreak{}odp\_\allowbreak{}2\\
Type: integer\\
Default: 30
}

{\footnotesize вибрано одне або декілька з наступних ЗНАЧЕНЬ ПАРАМЕТРІВ: \{з попередженням; без попередження\};}


{\footnotesize\bfseries
No: 3\\
Name: ca\_\allowbreak{}8\_\allowbreak{}3\_\allowbreak{}01\\
Type: string\\
Default: nil
}

{\footnotesize впроваджено процес тестування на проникнення, який}




\subsection{ВНУТРІШНІ СИСТЕМНІ ЗВ'ЯЗКИ (CA-9)}
a. Авторизувати внутрішні підключення [Призначення: системні компоненти або класи компонентів, що організація визначила] до системи;\\ 
b. Задокументувати, для кожного внутрішнього з'єднання, характеристики інтерфейсу, вимоги безпеки та приватності, а також характер переданої інформації;\\ 
c. Розірвати внутрішні системні підключення після [Призначення: умови, визначені організацією];\\ 
d. Переглядати [Призначення: частота, визначена організацією] постійну потребу в кожному внутрішньому з'єднанні.

\vspace{10pt}
{\footnotesize\bfseries
No: 1\\
Name: ca\_\allowbreak{}9\_\allowbreak{}odp\_\allowbreak{}1\\
Type: string\\
Default: nil
}

{\footnotesize визначено компоненти системи або класи компонентів, що потребують внутрішніх підключень до системи;}


{\footnotesize\bfseries
No: 2\\
Name: ca\_\allowbreak{}9\_\allowbreak{}odp\_\allowbreak{}2\\
Type: string\\
Default: nil
}

{\footnotesize визначено умови, за яких необхідно розірвати внутрішні підключення;}


{\footnotesize\bfseries
No: 3\\
Name: ca\_\allowbreak{}9\_\allowbreak{}odp\_\allowbreak{}3\\
Type: integer\\
Default: 30
}

{\footnotesize визначено частоту, з якою необхідно переглядати постійну потребу в кожному внутрішньому з'єднанні;}


{\footnotesize\bfseries
No: 4\\
Name: ca\_\allowbreak{}9\_\allowbreak{}a\\
Type: string\\
Default: nil
}

{\footnotesize внутрішні підключення <CA-09\_\allowbreak{}ODP[01] компонентів системи> 172 до системи є авторизовані;}


{\footnotesize\bfseries
No: 5\\
Name: ca\_\allowbreak{}9\_\allowbreak{}b\_\allowbreak{}1\\
Type: string\\
Default: nil
}

{\footnotesize для кожного внутрішнього з'єднання задокументовані характеристики інтерфейсу;}


{\footnotesize\bfseries
No: 6\\
Name: ca\_\allowbreak{}9\_\allowbreak{}b\_\allowbreak{}2\\
Type: string\\
Default: nil
}

{\footnotesize для кожного внутрішнього з'єднання задокументовані вимоги безпеки;}


{\footnotesize\bfseries
No: 7\\
Name: ca\_\allowbreak{}9\_\allowbreak{}b\_\allowbreak{}3\\
Type: string\\
Default: nil
}

{\footnotesize для кожного внутрішнього з'єднання задокументовані вимоги конфіденційності;}


{\footnotesize\bfseries
No: 8\\
Name: ca\_\allowbreak{}9\_\allowbreak{}b\_\allowbreak{}4\\
Type: string\\
Default: nil
}

{\footnotesize для кожного внутрішнього з'єднання задокументовані характер переданої інформації;}


{\footnotesize\bfseries
No: 9\\
Name: ca\_\allowbreak{}9\_\allowbreak{}c\\
Type: string\\
Default: nil
}

{\footnotesize внутрішні з'єднання системи розриваються після виконання <CA09\_\allowbreak{}ODP[02] умов>;}


{\footnotesize\bfseries
No: 10\\
Name: ca\_\allowbreak{}9\_\allowbreak{}d\\
Type: string\\
Default: nil
}

{\footnotesize переглядається подальша потреба у кожному внутрішньому з'єднанні <CA-09\_\allowbreak{}ODP[03] частота>.}




\subsubsection{ВІДПОВІДНІСТЬ ЗАХОДІВ БЕЗПЕКИ (CA-9(1))}
Забезпечується надання результатів внутрішніх підключень до компонентів системи.

\vspace{10pt}
{\footnotesize\bfseries
No: 1\\
Name: ca\_\allowbreak{}9\_\allowbreak{}1\_\allowbreak{}1\\
Type: string\\
Default: nil
}

{\footnotesize перед встановленням внутрішнього з'єднання виконується перевірка на  відповідність вимогам безпеки складових компонентів системи;}


{\footnotesize\bfseries
No: 2\\
Name: ca\_\allowbreak{}9\_\allowbreak{}1\_\allowbreak{}2\\
Type: string\\
Default: nil
}

{\footnotesize перед встановленням внутрішнього з'єднання виконується перевірка на відповідність вимогам конфіденційності складових компонентів системи;}




\section{CM}
Клас заходів захисту CM --- УПРАВЛІННЯ КОНФІГУРАЦІЄЮ

Цей клас гарантує створення та підтримання базових налаштувань системи, а також строгий контроль за змінами в апаратному та програмному забезпеченні.

\textbf{Перелік заходів захисту:} Політика та процедури управління конфігурацією (CM-1); Базова конфігурація (CM-2); Перегляд та оновлення (CM-2(1)) [Вилучено]; Автоматизація підтримки задля точності та вживаності (CM-2(2)); Зберігання попередніх версій конфігурацій (CM-2(3)); Неавторизоване програмне забезпечення (CM-2(4)); Авторизоване програмне забезпечення (CM-2(5)) [Вилучено]; Розробка та середовище тестування (CM-2(6)); Конфігурація систем та компонентів для сфер з високим ризиком (CM-2(7)); Управління змінами конфігурації (CM-3); Автоматизоване документування, повідомлення та заборона внесення змін (CM-3(1)); Тестування, валідація та документування змін (CM-3(2)); Автоматизована реалізація змін (CM-3(3)); Представник безпеки (CM-3(4)); Автоматичне реагування безпеки (CM-3(5)); Управління засобами криптографічного захисту (CM-3(6)); Перегляд змін у системі (CM-3(7)); Запобігання чи обмеження змін конфігурації (CM-3(8)); Аналіз впливу на безпеку та приватність (CM-4); Відокремлені випробувальні середовища (CM-4(1)); Верифікація функцій безпеки та приватності (CM-4(2)); Обмеження доступу до зміни (CM-5); Аудит і здійснення автоматичного доступу (CM-5(1)); Перегляд змін у системі (CM-5(2)) [Вилучено]; Підписані компоненти (CM-5(3)) [Вилучено]; Подвійна авторизація (CM-5(4)); Обмеження повноважень для виробництва та експлуатації (CM-5(5)); Обмеження повноважень для бібліотек (CM-5(6)); Обмеження доступу до зміни вадження заходів захисту (CM-5(7)) [Вилучено]; Налаштування конфігурації (CM-6); Автоматизоване управління, застосування та верифікація (CM-6(1)); Налаштування конфігурації санкціоновані зміни (CM-6(2)); Демонстрація відповідності (CM-6(4)) [Вилучено]; Мінімально необхідна функціональність (CM-7); Періодичний перегляд (CM-7(1)); Заборона виконання програми (CM-7(2)); Відповідність реєстрації (CM-7(3)); Неавторизоване програмне забезпечення - чорний список (CM-7(4)); Авторизоване програмне забезпечення -- білий список (CM-7(5)); Замкнуті середовища з обмеженими привілеями (CM-7(6)); Виконуваний код у захищеному середовищі (CM-7(7)); Бінарний або машинний виконуваний код (CM-7(8)); Заборона використання неавторизованого обладнання (CM-7(9)); Інвентаризація компонентів системи (CM-8); Оновлення під час встановлення та видалення (CM-8(1)); Автоматизована підтримка (CM-8(2)); Автоматизоване виявлення неавторизованих компонентів (CM-8(3)); Інформація про підзвітність (CM-8(4)); Інвентаризація компонентів системи дублювання компонентів обліку (CM-8(5)); Перевірені налаштування та затверджені відхилення (CM-8(6)); Централізоване сховище (CM-8(7)); Автоматизоване відстеження місця розташування (CM-8(8)); Призначення компонентів системам (CM-8(9)); План управління конфігурацією (CM-9); Встановлення відповідальності (CM-9(1)); Обмеження використання програмного забезпечення (CM-10); Обмеження використання програмного забезпечення програмне забезпечення з відкритим вихідним кодом (CM-10(1)); Встановлене користувачем програмне забезпечення (CM-11); Встановлене користувачем програмне забезпечення попередження про несанкціоновану інсталяцію (CM-11(1)); Встановлене користувачем програмне забезпечення встановлення програмного забезпечення з привілейованим статусом (CM-11(2)); Встановлене користувачем програмне забезпечення автоматичне виконання і моніторинг (CM-11(3)); Розташування інформації (CM-12); Автоматизовані інструменти підтримки розташування інформації (CM-12(1)); Відображення дій даних (CM-13); Підписані компоненти (CM-14).

\subsection{ПОЛІТИКА ТА ПРОЦЕДУРИ УПРАВЛІННЯ КОНФІГУРАЦІЄЮ (CM-1)}
a. Розробити, задокументувати та поширити серед [Призначення: визначених організацією персоналу або ролей]:\\ 
1. [Вибір (один або декілька): Рівень організації; Рівень місії/бізнес-процесу; рівень системи] політики управління конфігурацією, яка:\\ 
2.\\ 
(a) містить мету, сферу застосування, ролі, обов'язки, відповідальність керівництва, координацію між організаційними підрозділами та систему контролю відповідності (complaince);\\ 
(b) відповідає чинним законам, нормативним документам, наказам, положенням, політикам, стандартам і керівним принципам; процедури, що сприяють реалізації політики управління конфігурацією та пов'язаних з нею заходів управління конфігурацією.\\ 
b. Призначити [Призначення: посадова особа, визначена організацією] для управління розробкою, документуванням і розповсюдженням політики та процедур керування конфігурацією.\\ 
c. Переглядати та оновлювати поточну політику управління конфігурацією:\\ 
1. Політика [Призначення: частота, визначена організацією] та наступні [Призначення: події, визначені організацією];\\ 
2. Процедури [Призначення: частота, визначена організацією] та наступні [Призначення: події, визначені організацією].

\vspace{10pt}
{\footnotesize\bfseries
No: 1\\
Name: cm\_\allowbreak{}1\_\allowbreak{}odp\_\allowbreak{}1\\
Type: string\\
Default: nil
}

{\footnotesize визначено персонал або ролі, на яких поширюється політика управління конфігурацією;}


{\footnotesize\bfseries
No: 2\\
Name: cm\_\allowbreak{}1\_\allowbreak{}odp\_\allowbreak{}2\\
Type: string\\
Default: nil
}

{\footnotesize визначено персонал або ролі, на яких поширюється процедури управління конфігурацією;}


{\footnotesize\bfseries
No: 3\\
Name: cm\_\allowbreak{}1\_\allowbreak{}odp\_\allowbreak{}3\\
Type: string\\
Default: nil
}

{\footnotesize вибрано одне або декілька з наступних ЗНАЧЕНЬ ПАРА- МЕТРІВ: \{рівень організації; рівень місії/бізнес -процесу; рівень системи\};}


{\footnotesize\bfseries
No: 4\\
Name: cm\_\allowbreak{}1\_\allowbreak{}odp\_\allowbreak{}4\\
Type: string\\
Default: nil
}

{\footnotesize визначено посадову особу, яка управляє розробкою, документуванням і розповсюдженням політики та процедур керування конфігурацією;}


{\footnotesize\bfseries
No: 5\\
Name: cm\_\allowbreak{}1\_\allowbreak{}odp\_\allowbreak{}5\\
Type: string\\
Default: nil
}

{\footnotesize визначено частоту, з якою переглядається та оновлюється поточна політика управління конфігурацією;}


{\footnotesize\bfseries
No: 6\\
Name: cm\_\allowbreak{}1\_\allowbreak{}odp\_\allowbreak{}6\\
Type: string\\
Default: nil
}

{\footnotesize визначено події, після яких переглядається та оновлюється поточна політика управління конфігурацією;}


{\footnotesize\bfseries
No: 7\\
Name: cm\_\allowbreak{}1\_\allowbreak{}odp\_\allowbreak{}7\\
Type: string\\
Default: nil
}

{\footnotesize визначено частоту, з якою перегляда ються та оновлюються поточні процедури управління конфігурацією;}


{\footnotesize\bfseries
No: 8\\
Name: cm\_\allowbreak{}1\_\allowbreak{}odp\_\allowbreak{}8\\
Type: string\\
Default: nil
}

{\footnotesize визначено події, після яких переглядаються та оновлюються поточні процедури управління конфігурацією;}


{\footnotesize\bfseries
No: 9\\
Name: cm\_\allowbreak{}1\_\allowbreak{}a\_\allowbreak{}1\\
Type: string\\
Default: nil
}

{\footnotesize розроблено та задокументовано політику управління конфігурацією;}


{\footnotesize\bfseries
No: 10\\
Name: cm\_\allowbreak{}1\_\allowbreak{}a\_\allowbreak{}2\\
Type: string\\
Default: nil
}

{\footnotesize політика управління конфігурацією поширюється на  <CM01\_\allowbreak{}ODP[01] персонал або ролі>;}


{\footnotesize\bfseries
No: 11\\
Name: cm\_\allowbreak{}1\_\allowbreak{}a\_\allowbreak{}3\\
Type: string\\
Default: nil
}

{\footnotesize розроблені та задокументовані процедури управління , що сприяють реалізації політики управління конфігурацією та пов'язаних з нею заходів управління конфігурацією;}


{\footnotesize\bfseries
No: 12\\
Name: cm\_\allowbreak{}1\_\allowbreak{}a\_\allowbreak{}4\\
Type: string\\
Default: nil
}

{\footnotesize процедури управління конфігурацією поширюються на  <CM01\_\allowbreak{}ODP[02] персонал або ролі>;}


{\footnotesize\bfseries
No: 13\\
Name: cm\_\allowbreak{}1\_\allowbreak{}a\_\allowbreak{}1\_\allowbreak{}a\_\allowbreak{}1\\
Type: string\\
Default: nil
}

{\footnotesize <CM-01\_\allowbreak{}ODP[03] ВИБРАНЕ ЗНАЧЕННЯ ПАРАМЕТРА(ІВ)> політики управління конфігурацією містить мету; 175}


{\footnotesize\bfseries
No: 14\\
Name: cm\_\allowbreak{}1\_\allowbreak{}a\_\allowbreak{}1\_\allowbreak{}a\_\allowbreak{}2\\
Type: string\\
Default: nil
}

{\footnotesize <CM-01\_\allowbreak{}ODP[03] ВИБРАНЕ ЗНАЧЕННЯ ПАРАМЕТРА(ІВ)> політики управління конфігурацією містить сферу застосування;}


{\footnotesize\bfseries
No: 15\\
Name: cm\_\allowbreak{}1\_\allowbreak{}a\_\allowbreak{}1\_\allowbreak{}a\_\allowbreak{}3\\
Type: string\\
Default: nil
}

{\footnotesize <CM-01\_\allowbreak{}ODP[03] ВИБРАНЕ ЗНАЧЕННЯ ПАРАМЕТРА(ІВ)> політики управління конфігурацією містить ролі;}


{\footnotesize\bfseries
No: 16\\
Name: cm\_\allowbreak{}1\_\allowbreak{}a\_\allowbreak{}1\_\allowbreak{}a\_\allowbreak{}4\\
Type: string\\
Default: nil
}

{\footnotesize <CM-01\_\allowbreak{}ODP[03] ВИБРАНЕ ЗНАЧЕННЯ ПАРАМЕТРА(ІВ)> політики управління конфігурацією містить обов'язки;}


{\footnotesize\bfseries
No: 17\\
Name: cm\_\allowbreak{}1\_\allowbreak{}a\_\allowbreak{}1\_\allowbreak{}a\_\allowbreak{}5\\
Type: string\\
Default: nil
}

{\footnotesize <CM-01\_\allowbreak{}ODP[03] ВИБРАНЕ ЗНАЧЕННЯ ПАРАМЕТРА(ІВ)> політики управління конфігурацією містить відповідальність керівництва;}


{\footnotesize\bfseries
No: 18\\
Name: cm\_\allowbreak{}1\_\allowbreak{}a\_\allowbreak{}1\_\allowbreak{}a\_\allowbreak{}6\\
Type: string\\
Default: nil
}

{\footnotesize <CM-01\_\allowbreak{}ODP[03] ВИБРАНЕ ЗНАЧЕННЯ ПАРАМЕТРА(ІВ)> політики управління конфігурацією містить координацію між підрозділами організації;}


{\footnotesize\bfseries
No: 19\\
Name: cm\_\allowbreak{}1\_\allowbreak{}a\_\allowbreak{}1\_\allowbreak{}a\_\allowbreak{}7\\
Type: string\\
Default: nil
}

{\footnotesize <CM-01\_\allowbreak{}ODP[03] ВИБРАНЕ ЗНАЧЕННЯ ПАРАМЕТРА(ІВ)> політики управління конфігурацією містить систему контролю відповідності;}


{\footnotesize\bfseries
No: 20\\
Name: cm\_\allowbreak{}1\_\allowbreak{}a\_\allowbreak{}1\_\allowbreak{}b\\
Type: string\\
Default: nil
}

{\footnotesize політика управління конфігурацією відповідає чинним законам, нормативним документам, наказам, положенням, політикам, стандартам і керівним принципам;}


{\footnotesize\bfseries
No: 21\\
Name: cm\_\allowbreak{}1\_\allowbreak{}b\\
Type: string\\
Default: nil
}

{\footnotesize <CM-01\_\allowbreak{}ODP[04] посадова особа> призначається для управління розробкою, документуванням і розповсюдженням політики та процедур керування конфігурацією;}


{\footnotesize\bfseries
No: 22\\
Name: cm\_\allowbreak{}1\_\allowbreak{}c\_\allowbreak{}1\_\allowbreak{}1\\
Type: string\\
Default: nil
}

{\footnotesize переглядається та оновлюється поточна політика управління}


{\footnotesize\bfseries
No: 23\\
Name: cm\_\allowbreak{}1\_\allowbreak{}c\_\allowbreak{}1\_\allowbreak{}2\\
Type: string\\
Default: nil
}

{\footnotesize переглядається та оновлюється поточна політика управління}


{\footnotesize\bfseries
No: 24\\
Name: cm\_\allowbreak{}1\_\allowbreak{}c\_\allowbreak{}2\_\allowbreak{}1\\
Type: string\\
Default: nil
}

{\footnotesize переглядаються та оновлюються поточні процедури управління}


{\footnotesize\bfseries
No: 25\\
Name: cm\_\allowbreak{}1\_\allowbreak{}c\_\allowbreak{}2\_\allowbreak{}2\\
Type: string\\
Default: nil
}

{\footnotesize переглядаються та оновлюються поточні процедури управління}




\subsection{БАЗОВА КОНФІГУРАЦІЯ (CM-2)}
a. Розробити, задокументувати та підтримувати за допомогою заходів конфігурації поточні базові налаштування системи.\\ 
b. Переглядати та оновлювати базові налаштування системи:\\ 
1. з [Призначення: визначеною організацією частотою];\\ 
2. за потреби внаслідок [Призначення: визначених організацією обставин];\\ 
3. коли встановлені нові або оновлені компоненти системи.

\vspace{10pt}
{\footnotesize\bfseries
No: 1\\
Name: cm\_\allowbreak{}2\_\allowbreak{}odp\_\allowbreak{}1\\
Type: string\\
Default: nil
}

{\footnotesize визначено частоту перегляду та оновлення базових налаштувань;}


{\footnotesize\bfseries
No: 2\\
Name: cm\_\allowbreak{}2\_\allowbreak{}odp\_\allowbreak{}2\\
Type: string\\
Default: nil
}

{\footnotesize визначено обставини, що вимагають перегляду та оновлення базових налаштувань;}


{\footnotesize\bfseries
No: 3\\
Name: cm\_\allowbreak{}2\_\allowbreak{}a\_\allowbreak{}1\\
Type: string\\
Default: nil
}

{\footnotesize розроблено та задокументовано поточні базові налаштування системи;}


{\footnotesize\bfseries
No: 4\\
Name: cm\_\allowbreak{}2\_\allowbreak{}a\_\allowbreak{}2\\
Type: string\\
Default: nil
}

{\footnotesize поточні базові налаштування системи підтримуються за допомогою заходів конфігурації;}


{\footnotesize\bfseries
No: 5\\
Name: cm\_\allowbreak{}2\_\allowbreak{}b\_\allowbreak{}1\\
Type: string\\
Default: nil
}

{\footnotesize переглядаються та оновлюються базові налаштування системи}


{\footnotesize\bfseries
No: 6\\
Name: cm\_\allowbreak{}2\_\allowbreak{}b\_\allowbreak{}2\\
Type: string\\
Default: nil
}

{\footnotesize переглядаються та оновлюються базові налаштування системи}


{\footnotesize\bfseries
No: 7\\
Name: cm\_\allowbreak{}2\_\allowbreak{}b\_\allowbreak{}3\\
Type: string\\
Default: nil
}

{\footnotesize переглядаються та оновлюються базові налаштування системи коли встановлюються або модернізуються компоненти системи.}




\subsubsection{ПЕРЕГЛЯД ТА ОНОВЛЕННЯ (CM-2(1)) [Вилучено]}
[Вилучено: Включено до CM-02]

\vspace{10pt}
\textit{Немає параметрів для цього контролю.}
\vspace{10pt}


\subsubsection{АВТОМАТИЗАЦІЯ ПІДТРИМКИ ЗАДЛЯ ТОЧНОСТІ ТА ВЖИВАНОСТІ (CM-2(2))}
Підтримувати актуальність, повноту, точність і доступність базової конфігурації системи за допомогою [Призначення: автоматизовані механізми, визначені організацією].

\vspace{10pt}
{\footnotesize\bfseries
No: 1\\
Name: cm\_\allowbreak{}2\_\allowbreak{}2\_\allowbreak{}odp\\
Type: string\\
Default: nil
}

{\footnotesize визначено автоматизовані механізми підтримки базової конфігурації системи; 177}


{\footnotesize\bfseries
No: 2\\
Name: cm\_\allowbreak{}2\_\allowbreak{}2\_\allowbreak{}1\\
Type: string\\
Default: nil
}

{\footnotesize актуальність базової конфігурації системи підтримується за допомогою <CM-02(02)\_\allowbreak{}ODP автоматизовані механізми>;}


{\footnotesize\bfseries
No: 3\\
Name: cm\_\allowbreak{}2\_\allowbreak{}2\_\allowbreak{}2\\
Type: string\\
Default: nil
}

{\footnotesize повнота базової конфігурації системи підтримується за допомогою <CM-02(02)\_\allowbreak{}ODP автоматизовані механізми>;}


{\footnotesize\bfseries
No: 4\\
Name: cm\_\allowbreak{}2\_\allowbreak{}2\_\allowbreak{}3\\
Type: string\\
Default: nil
}

{\footnotesize точність базової конфігурації системи підтримується за допомогою <CM-02(02)\_\allowbreak{}ODP автоматизовані механізми>;}


{\footnotesize\bfseries
No: 5\\
Name: cm\_\allowbreak{}2\_\allowbreak{}2\_\allowbreak{}4\\
Type: string\\
Default: nil
}

{\footnotesize доступність базової конфігурації системи підтримується за допомогою <CM-02(02)\_\allowbreak{}ODP автоматизовані механізми>;}




\subsubsection{ЗБЕРІГАННЯ ПОПЕРЕДНІХ ВЕРСІЙ КОНФІГУРАЦІЙ (CM-2(3))}
Зберігати [Призначення: кількість, визначена організацією] попередніх версій базових конфігурацій системи для підтримки відкату.

\vspace{10pt}
{\footnotesize\bfseries
No: 1\\
Name: cm\_\allowbreak{}2\_\allowbreak{}3\_\allowbreak{}odp\\
Type: string\\
Default: nil
}

{\footnotesize визначено попередні версії базових конфігурацій системи необхідні для підтримки відкату}


{\footnotesize\bfseries
No: 2\\
Name: cm\_\allowbreak{}2\_\allowbreak{}3\_\allowbreak{}01\\
Type: string\\
Default: nil
}

{\footnotesize зберігати <CM-02(03)\_\allowbreak{}ODP попередні версії>  для підтримки відкату}




\subsubsection{НЕАВТОРИЗОВАНЕ ПРОГРАМНЕ ЗАБЕЗПЕЧЕННЯ (CM-2(4))}


\vspace{10pt}
\textit{Немає параметрів для цього контролю.}
\vspace{10pt}


\subsubsection{АВТОРИЗОВАНЕ ПРОГРАМНЕ ЗАБЕЗПЕЧЕННЯ (CM-2(5)) [Вилучено]}
[Вилучено: Включено до CM-07(05)]

\vspace{10pt}
\textit{Немає параметрів для цього контролю.}
\vspace{10pt}


\subsubsection{РОЗРОБКА ТА СЕРЕДОВИЩЕ ТЕСТУВАННЯ (CM-2(6))}
Підтримувати базову конфігурацію для розробки системи та тестових середовищ, які керуються окремо від робочої базової конфігурації.

\vspace{10pt}
{\footnotesize\bfseries
No: 1\\
Name: cm\_\allowbreak{}2\_\allowbreak{}6\_\allowbreak{}1\\
Type: string\\
Default: nil
}

{\footnotesize підтримується базова конфігурація для розробки системи, які керуються окремо від робочої базової конфігурації.}


{\footnotesize\bfseries
No: 2\\
Name: cm\_\allowbreak{}2\_\allowbreak{}6\_\allowbreak{}2\\
Type: string\\
Default: nil
}

{\footnotesize підтримується базова конфігурація для розробки тестових середовищ, які керуються окремо від робочої базової конфігурації.}




\subsubsection{КОНФІГУРАЦІЯ СИСТЕМ ТА КОМПОНЕНТІВ ДЛЯ СФЕР З ВИСОКИМ РИЗИКОМ (CM-2(7))}
(a) Видавати [Призначення: визначених організацією систем або компонентів систем] з [Призначенням: визначеними організацією конфігураціями] особам, що перебувають у місцях, які організація вважає місцями зі значним ризиком;\\ 
(b) Застосувати [Призначення: визначені організацією запобіжні заходи безпеки] до компонентів, коли особи повертаються з поїздки.

\vspace{10pt}
{\footnotesize\bfseries
No: 1\\
Name: cm\_\allowbreak{}2\_\allowbreak{}7\_\allowbreak{}odp\_\allowbreak{}1\\
Type: string\\
Default: nil
}

{\footnotesize визначено системи або компоненти систем, які мають видаватися особам, що перебувають у місцях зі значним ризиком;}


{\footnotesize\bfseries
No: 2\\
Name: cm\_\allowbreak{}2\_\allowbreak{}7\_\allowbreak{}odp\_\allowbreak{}1\\
Type: string\\
Default: nil
}

{\footnotesize визначено конфігурації систем або компонентів си179 стем, що видаються у місцях зі значним ризиком;}


{\footnotesize\bfseries
No: 3\\
Name: cm\_\allowbreak{}2\_\allowbreak{}7\_\allowbreak{}odp\_\allowbreak{}1\\
Type: string\\
Default: nil
}

{\footnotesize визначено заходи безпеки, які мають застосовуватися після повернення осіб з поїздки;}


{\footnotesize\bfseries
No: 4\\
Name: cm\_\allowbreak{}2\_\allowbreak{}7\_\allowbreak{}a\\
Type: string\\
Default: nil
}

{\footnotesize <CM-02(07)\_\allowbreak{}ODP[01] системи або компоненти системи> з <CM-02(07)\_\allowbreak{}ODP[02] конфігураціями> видаються особам, що перебувають у місцях, які організація вважає, становлять значний ризик;}


{\footnotesize\bfseries
No: 5\\
Name: cm\_\allowbreak{}2\_\allowbreak{}7\_\allowbreak{}b\\
Type: string\\
Default: nil
}

{\footnotesize <CM-02(07)\_\allowbreak{}ODP[03] заходи безпеки> застосовуються до систем або компонентів системи, коли особи повертаються з поїздки.}




\subsection{УПРАВЛІННЯ ЗМІНАМИ КОНФІГУРАЦІЇ (CM-3)}
a. Визначити типи змін у системі, які контролюються конфігурацією.\\ 
b. Переглядати запропоновані зміни в конфігурації, контрольовані системою, і схвалити або відхилити ці зміни з явним урахуванням аналізу наслідків безпеки.\\ 
c. Документувати рішення зі зміни конфігурації системи.\\ 
d. Впровадити схвалені зміни конфігурації в систему.\\ 
e. Зберігати записи змін конфігурації системі впродовж [Призначення: певного періоду часу, визначеного організацією].\\ 
f. Здійснювати моніторинг і аналіз дій, пов'язаних зі змінами конфігурації системи.\\ 
g. Координувати та впроваджувати нагляд за діяльністю з управління змінами конфігурації за допомогою [Призначення: елементу управління змінами конфігурації, визначеного організацією], який викликається [Вибір (один або кілька): [Призначення: з визначеною організацією частотою]; [Призначення: визначені організацією умови зміни конфігурації]].

\vspace{10pt}
{\footnotesize\bfseries
No: 1\\
Name: cm\_\allowbreak{}3\_\allowbreak{}odp\_\allowbreak{}1\\
Type: string\\
Default: nil
}

{\footnotesize визначено період часу, протягом якого зберігатимуться записи про зміни конфігурації;}


{\footnotesize\bfseries
No: 2\\
Name: cm\_\allowbreak{}3\_\allowbreak{}odp\_\allowbreak{}2\\
Type: string\\
Default: nil
}

{\footnotesize визначено елементи управління змінами конфігурації, відповідальні за координацію та нагляд за діяльністю з управління змінами;}


{\footnotesize\bfseries
No: 3\\
Name: cm\_\allowbreak{}3\_\allowbreak{}odp\_\allowbreak{}3\\
Type: string\\
Default: nil
}

{\footnotesize вибрано одне або декілька з наступних ЗНАЧЕНЬ}


{\footnotesize\bfseries
No: 4\\
Name: cm\_\allowbreak{}3\_\allowbreak{}odp\_\allowbreak{}4\\
Type: string\\
Default: nil
}

{\footnotesize визначено частоту, з якою викликаються елементи управління змінами конфігурації (якщо вибрано);}


{\footnotesize\bfseries
No: 5\\
Name: cm\_\allowbreak{}3\_\allowbreak{}odp\_\allowbreak{}5\\
Type: string\\
Default: nil
}

{\footnotesize визначено умови, за яких викликаються елементи управління змінами конфігурації (якщо вибрано); 180}


{\footnotesize\bfseries
No: 6\\
Name: cm\_\allowbreak{}3\_\allowbreak{}a\\
Type: string\\
Default: nil
}

{\footnotesize визначено та задокументовано типи змін до системи, які контролюються конфігурацією;}


{\footnotesize\bfseries
No: 7\\
Name: cm\_\allowbreak{}3\_\allowbreak{}b\_\allowbreak{}1\\
Type: string\\
Default: nil
}

{\footnotesize розглядаються запропоновані зміни в конфігурації, що контролюються системою;}


{\footnotesize\bfseries
No: 8\\
Name: cm\_\allowbreak{}3\_\allowbreak{}b\_\allowbreak{}2\\
Type: string\\
Default: nil
}

{\footnotesize запропоновані зміни в конфігурації, що контролюються системою, схвалюються або  відхиляються з урахуванням аналізу наслідків безпеки;}


{\footnotesize\bfseries
No: 9\\
Name: cm\_\allowbreak{}3\_\allowbreak{}c\\
Type: string\\
Default: nil
}

{\footnotesize рішення про зміну конфігурації системи документуються;}


{\footnotesize\bfseries
No: 10\\
Name: cm\_\allowbreak{}3\_\allowbreak{}d\\
Type: string\\
Default: nil
}

{\footnotesize впроваджуються схвалені зміни до конфігурації в систему;}


{\footnotesize\bfseries
No: 11\\
Name: cm\_\allowbreak{}3\_\allowbreak{}e\\
Type: string\\
Default: nil
}

{\footnotesize записи про зміни конфігурації у системі зберігаються протягом}


{\footnotesize\bfseries
No: 12\\
Name: cm\_\allowbreak{}3\_\allowbreak{}f\_\allowbreak{}1\\
Type: string\\
Default: nil
}

{\footnotesize здійснюється моніторинг дій, пов'язаних зі змінами конфігурації системи;}


{\footnotesize\bfseries
No: 13\\
Name: cm\_\allowbreak{}3\_\allowbreak{}f\_\allowbreak{}2\\
Type: string\\
Default: nil
}

{\footnotesize здійснюється аналіз дій, пов'язаних зі змінами конфігурації системи;}


{\footnotesize\bfseries
No: 14\\
Name: cm\_\allowbreak{}3\_\allowbreak{}g\_\allowbreak{}1\\
Type: string\\
Default: nil
}

{\footnotesize діяльність з управління змінами  конфігурації координується та}


{\footnotesize\bfseries
No: 15\\
Name: cm\_\allowbreak{}3\_\allowbreak{}g\_\allowbreak{}2\\
Type: string\\
Default: nil
}

{\footnotesize елемент управління зміною конфігурації викликається  <CM03\_\allowbreak{}ODP[03] ВИБРАНЕ ЗНАЧЕННЯ ПАРАМЕТРА(ів)>.}




\subsubsection{АВТОМАТИЗОВАНЕ ДОКУМЕНТУВАННЯ, ПОВІДОМЛЕННЯ ТА ЗАБОРОНА ВНЕСЕННЯ ЗМІН (CM-3(1))}


\vspace{10pt}
{\footnotesize\bfseries
No: 1\\
Name: cm\_\allowbreak{}3\_\allowbreak{}1\_\allowbreak{}odp\_\allowbreak{}1\\
Type: string\\
Default: nil
}

{\footnotesize визначено механізми, що використовуються для ав181 томатизації управління змінами конфігурації}


{\footnotesize\bfseries
No: 2\\
Name: cm\_\allowbreak{}3\_\allowbreak{}1\_\allowbreak{}odp\_\allowbreak{}2\\
Type: string\\
Default: nil
}

{\footnotesize визначені уповноважені органи, про які необхідно повідомляти та узгоджувати пропоновані зміни в системі;}


{\footnotesize\bfseries
No: 3\\
Name: cm\_\allowbreak{}3\_\allowbreak{}1\_\allowbreak{}odp\_\allowbreak{}3\\
Type: string\\
Default: nil
}

{\footnotesize визначено період часу, після якого слід виділяти зміни, які не були схвалені або відхилені;}


{\footnotesize\bfseries
No: 4\\
Name: cm\_\allowbreak{}3\_\allowbreak{}1\_\allowbreak{}odp\_\allowbreak{}4\\
Type: string\\
Default: nil
}

{\footnotesize визначено персонал, який буде повідомлений про завершення затверджених змін;}


{\footnotesize\bfseries
No: 5\\
Name: cm\_\allowbreak{}3\_\allowbreak{}1\_\allowbreak{}a\\
Type: string\\
Default: nil
}

{\footnotesize <CM-03(01)\_\allowbreak{}ODP[01] автоматизовані механізми>  використовуються для документування запропонованих змін до системи;}


{\footnotesize\bfseries
No: 6\\
Name: cm\_\allowbreak{}3\_\allowbreak{}1\_\allowbreak{}b\\
Type: string\\
Default: nil
}

{\footnotesize <CM-03(01)\_\allowbreak{}ODP[01] автоматизовані  механізми> використовуються для повідомлення <CM03(01)\_\allowbreak{}ODP[02] уповноважені органи>  про запропоновані зміни в системі та запиту на затвердження змін;}


{\footnotesize\bfseries
No: 7\\
Name: cm\_\allowbreak{}3\_\allowbreak{}1\_\allowbreak{}c\\
Type: string\\
Default: nil
}

{\footnotesize <CM-03(01)\_\allowbreak{}ODP[01] автоматизовані механізми>  використовуються для виділення запропонованих змін до системи, які не були схвалені або відхилені протягом}


{\footnotesize\bfseries
No: 8\\
Name: cm\_\allowbreak{}3\_\allowbreak{}1\_\allowbreak{}d\\
Type: string\\
Default: nil
}

{\footnotesize <CM-03(01)\_\allowbreak{}ODP[01] автоматизовані механізми>  використовуються для заборони внесення змін до системи до отримання відповідних погоджень;}


{\footnotesize\bfseries
No: 9\\
Name: cm\_\allowbreak{}3\_\allowbreak{}1\_\allowbreak{}e\\
Type: string\\
Default: nil
}

{\footnotesize <CM-03(01)\_\allowbreak{}ODP[01] автоматизовані механізми>  використовуються для документування всіх змін в системі;}


{\footnotesize\bfseries
No: 10\\
Name: cm\_\allowbreak{}3\_\allowbreak{}1\_\allowbreak{}f\\
Type: string\\
Default: nil
}

{\footnotesize <CM-03(01)\_\allowbreak{}ODP[01] автоматизовані механізми>  використовуються для повідомлення <CM03(01)\_\allowbreak{}ODP[04] персоналу>  про завершення погоджених змін у системі.}




\subsubsection{ТЕСТУВАННЯ, ВАЛІДАЦІЯ ТА ДОКУМЕНТУВАННЯ ЗМІН (CM-3(2))}
Тестувати, перевіряти та документувати зміни в системі до повної їх реалізації.

\vspace{10pt}
{\footnotesize\bfseries
No: 1\\
Name: cm\_\allowbreak{}3\_\allowbreak{}2\_\allowbreak{}1\\
Type: string\\
Default: nil
}

{\footnotesize зміни в системі тестуються перед повним впровадженням змін;}


{\footnotesize\bfseries
No: 2\\
Name: cm\_\allowbreak{}3\_\allowbreak{}2\_\allowbreak{}2\\
Type: string\\
Default: nil
}

{\footnotesize зміни в системі перевіряються перед повним впровадженням змін;}


{\footnotesize\bfseries
No: 3\\
Name: cm\_\allowbreak{}3\_\allowbreak{}2\_\allowbreak{}3\\
Type: string\\
Default: nil
}

{\footnotesize зміни в системі документуються перед повним впровадженням змін.}




\subsubsection{АВТОМАТИЗОВАНА РЕАЛІЗАЦІЯ ЗМІН (CM-3(3))}
Внести зміни в поточний базову план системи та розгорнути оновлений базовий план на встановленій базі за допомогою [Призначення: автоматизовані механізми, визначені організацією].

\vspace{10pt}
{\footnotesize\bfseries
No: 1\\
Name: cm\_\allowbreak{}3\_\allowbreak{}3\_\allowbreak{}odp\\
Type: string\\
Default: nil
}

{\footnotesize визначено автоматизовані механізми для внесення змін та розгортання оновленого базового плану по всій встановленій базі;}


{\footnotesize\bfseries
No: 2\\
Name: cm\_\allowbreak{}3\_\allowbreak{}3\_\allowbreak{}1\\
Type: string\\
Default: nil
}

{\footnotesize зміни до поточного базового плану системи реалізуються за допомогою <CM-03(03)\_\allowbreak{}ODP автоматизовані механізми>;}


{\footnotesize\bfseries
No: 3\\
Name: cm\_\allowbreak{}3\_\allowbreak{}3\_\allowbreak{}2\\
Type: string\\
Default: nil
}

{\footnotesize оновлений базовий план розгортається по всій встановленій базі за допомогою <CM-03(03)\_\allowbreak{}ODP автоматизовані механізми>.}




\subsubsection{ПРЕДСТАВНИК БЕЗПЕКИ (CM-3(4))}
Вимагати від [Призначення: визначеного організацією представника з інформаційної безпеки] бути членом [Призначення: визначеного організацією елементу керування зміною конфігурацій].

\vspace{10pt}
{\footnotesize\bfseries
No: 1\\
Name: cm\_\allowbreak{}3\_\allowbreak{}4\_\allowbreak{}odp\_\allowbreak{}1\\
Type: string\\
Default: nil
}

{\footnotesize визначено представника з безпеки, який має бути членом елементу керування змінами конфігурації;}


{\footnotesize\bfseries
No: 2\\
Name: cm\_\allowbreak{}3\_\allowbreak{}4\_\allowbreak{}odp\_\allowbreak{}2\\
Type: string\\
Default: nil
}

{\footnotesize визначено представника з  конфіденційності, який має бути членом елементу керування змінами конфігурації;}


{\footnotesize\bfseries
No: 3\\
Name: cm\_\allowbreak{}3\_\allowbreak{}4\_\allowbreak{}odp\_\allowbreak{}3\\
Type: string\\
Default: nil
}

{\footnotesize визначено елемент керування змінами конфігурації, членами якого мають бути представники безпеки та конфіденційності;}


{\footnotesize\bfseries
No: 4\\
Name: cm\_\allowbreak{}3\_\allowbreak{}4\_\allowbreak{}1\\
Type: string\\
Default: nil
}

{\footnotesize <CM-03(04)\_\allowbreak{}ODP[01] представники безпеки> повинні бути членами <CM-03(04)\_\allowbreak{}ODP[03] елемента керування змінами конфігурації>;}


{\footnotesize\bfseries
No: 5\\
Name: cm\_\allowbreak{}3\_\allowbreak{}4\_\allowbreak{}2\\
Type: string\\
Default: nil
}

{\footnotesize <CM-03(04)\_\allowbreak{}ODP[02] представники конфіденційності> повинні бути членами <CM-03(04)\_\allowbreak{}ODP[03] елемента керування змінами конфігурації>.}




\subsubsection{АВТОМАТИЧНЕ РЕАГУВАННЯ БЕЗПЕКИ (CM-3(5))}
Реалізувати автоматичне [Призначення: визначене організацією реагування безпеки], якщо базова конфігурація системи змінюється несанкціонованим чином.

\vspace{10pt}
{\footnotesize\bfseries
No: 1\\
Name: cm\_\allowbreak{}3\_\allowbreak{}5\_\allowbreak{}odp\\
Type: string\\
Default: nil
}

{\footnotesize визначено реагування безпеки, які мають бути застосовані автоматично;}


{\footnotesize\bfseries
No: 2\\
Name: cm\_\allowbreak{}3\_\allowbreak{}5\_\allowbreak{}01\\
Type: string\\
Default: nil
}

{\footnotesize <CM-03(05)\_\allowbreak{}ODP реагування безпеки> автоматично застосовуються, якщо базова конфігурація системи змінюється несанкціонованим чином.}




\subsubsection{УПРАВЛІННЯ ЗАСОБАМИ КРИПТОГРАФІЧНОГО ЗАХИСТУ (CM-3(6))}
Забезпечити, щоб криптографічні механізми, які використовуються для забезпечення відповідних заходів захисту перебували під управлінням конфігурацією [Призначення: визначених організацією заходів безпеки].

\vspace{10pt}
{\footnotesize\bfseries
No: 1\\
Name: cm\_\allowbreak{}3\_\allowbreak{}6\_\allowbreak{}odp\\
Type: string\\
Default: nil
}

{\footnotesize визначено заходи захисту;}


{\footnotesize\bfseries
No: 2\\
Name: cm\_\allowbreak{}3\_\allowbreak{}6\_\allowbreak{}01\\
Type: string\\
Default: nil
}

{\footnotesize криптографічні механізми, які використовуються для забезпечення відповідних заходів захисту перебува ють під управлінням конфігурацією <CM-03(06)\_\allowbreak{}ODP заходи захисту>.}




\subsubsection{ПЕРЕГЛЯД ЗМІН У СИСТЕМІ (CM-3(7))}
Перегляньте зміни в системі [Призначення: частота, визначена організацією] або коли [Призначення: обставини, визначені організацією], щоб визначити, чи відбулися неавторизовані зміни.

\vspace{10pt}
{\footnotesize\bfseries
No: 1\\
Name: cm\_\allowbreak{}3\_\allowbreak{}7\_\allowbreak{}odp\_\allowbreak{}1\\
Type: string\\
Default: nil
}

{\footnotesize визначено частоту, з якою необхідно переглядати зміни;}


{\footnotesize\bfseries
No: 2\\
Name: cm\_\allowbreak{}3\_\allowbreak{}7\_\allowbreak{}odp\_\allowbreak{}2\\
Type: string\\
Default: nil
}

{\footnotesize визначено обставини, за яких зміни мають бути переглянуті;}


{\footnotesize\bfseries
No: 3\\
Name: cm\_\allowbreak{}3\_\allowbreak{}7\_\allowbreak{}01\\
Type: string\\
Default: nil
}

{\footnotesize зміни в системі переглядаються < CM-03(07)\_\allowbreak{}ODP[01] визначити, чи відбулися неавторизовані зміни.}




\subsubsection{ЗАПОБІГАННЯ ЧИ ОБМЕЖЕННЯ ЗМІН КОНФІГУРАЦІЇ (CM-3(8))}
Запобігати або обмежити зміни конфігурації системи за таких обставин: [Призначення: обставини, визначені організацією].

\vspace{10pt}
{\footnotesize\bfseries
No: 1\\
Name: cm\_\allowbreak{}3\_\allowbreak{}8\_\allowbreak{}odp\\
Type: string\\
Default: nil
}

{\footnotesize визначено обставини, за яких зміни мають бути запобіжені або обмежені;}


{\footnotesize\bfseries
No: 2\\
Name: cm\_\allowbreak{}3\_\allowbreak{}8\_\allowbreak{}01\\
Type: string\\
Default: nil
}

{\footnotesize зміни конфігурації системи запобігають або обмежують за <CM-03(08)\_\allowbreak{}ODP обставин>.}




\subsection{АНАЛІЗ ВПЛИВУ НА БЕЗПЕКУ ТА ПРИВАТНІСТЬ (CM-4)}
Аналізувати зміни в системі, щоб визначити потенційну загрозу безпеці та приватності перед реалізацією змін.

\vspace{10pt}
{\footnotesize\bfseries
No: 1\\
Name: cm\_\allowbreak{}4\_\allowbreak{}1\\
Type: string\\
Default: nil
}

{\footnotesize аналізуються зміни в системі, щоб визначити потенційну загрозу безпеці перед реалізацією змін}


{\footnotesize\bfseries
No: 2\\
Name: cm\_\allowbreak{}4\_\allowbreak{}2\\
Type: string\\
Default: nil
}

{\footnotesize аналізуються зміни в системі, щоб визначити потенційну загрозу конфіденційності перед реалізацією змін}




\subsubsection{ВІДОКРЕМЛЕНІ ВИПРОБУВАЛЬНІ СЕРЕДОВИЩА (CM-4(1))}


\vspace{10pt}
{\footnotesize\bfseries
No: 1\\
Name: cm\_\allowbreak{}4\_\allowbreak{}1\_\allowbreak{}1\\
Type: string\\
Default: nil
}

{\footnotesize зміни в системі аналізуються в окремому тестовому середовищі перед впровадженням в операційному середовищі;}


{\footnotesize\bfseries
No: 2\\
Name: cm\_\allowbreak{}4\_\allowbreak{}1\_\allowbreak{}2\\
Type: string\\
Default: nil
}

{\footnotesize зміни в системі аналізуються на предмет впливу на безпеку через недоліки;}


{\footnotesize\bfseries
No: 3\\
Name: cm\_\allowbreak{}4\_\allowbreak{}1\_\allowbreak{}3\\
Type: string\\
Default: nil
}

{\footnotesize зміни в системі аналізуються на предмет впливу на конфіденційність через недоліки;}


{\footnotesize\bfseries
No: 4\\
Name: cm\_\allowbreak{}4\_\allowbreak{}1\_\allowbreak{}4\\
Type: string\\
Default: nil
}

{\footnotesize зміни в системі аналізуються на предмет впливу на безпеку через слабкості;}


{\footnotesize\bfseries
No: 5\\
Name: cm\_\allowbreak{}4\_\allowbreak{}1\_\allowbreak{}5\\
Type: string\\
Default: nil
}

{\footnotesize зміни в системі аналізуються на предмет впливу на конфіденційність через слабкості;}


{\footnotesize\bfseries
No: 6\\
Name: cm\_\allowbreak{}4\_\allowbreak{}1\_\allowbreak{}6\\
Type: string\\
Default: nil
}

{\footnotesize зміни в системі аналізуються на предмет впливу на безпеку через несумістність;}


{\footnotesize\bfseries
No: 7\\
Name: cm\_\allowbreak{}4\_\allowbreak{}1\_\allowbreak{}7\\
Type: string\\
Default: nil
}

{\footnotesize зміни в системі аналізуються на предмет впливу на конфіденційність через несумістність;}


{\footnotesize\bfseries
No: 8\\
Name: cm\_\allowbreak{}4\_\allowbreak{}1\_\allowbreak{}8\\
Type: string\\
Default: nil
}

{\footnotesize зміни в системі аналізуються на предмет впливу на безпеку 187 через навмисне спричинення шкоди;}


{\footnotesize\bfseries
No: 9\\
Name: cm\_\allowbreak{}4\_\allowbreak{}1\_\allowbreak{}9\\
Type: string\\
Default: nil
}

{\footnotesize зміни в системі аналізуються на предмет впливу на конфіденційність через навмисне спричинення шкоди;}




\subsubsection{ВЕРИФІКАЦІЯ ФУНКЦІЙ БЕЗПЕКИ ТА ПРИВАТНОСТІ (CM-4(2))}
Після змін у системі переконайтеся, що відповідні заходи захисту реалізовано правильно і вони функціонують належним чином та дають бажаний результат щодо дотримання вимог безпеки та приватності для системи.

\vspace{10pt}
{\footnotesize\bfseries
No: 1\\
Name: cm\_\allowbreak{}4\_\allowbreak{}2\_\allowbreak{}1\\
Type: string\\
Default: nil
}

{\footnotesize заходи захисту, на які було здійснено вплив, реалізовані правильно з точки зору відповідності вимогам безпеки системи після внесення змін до системи;}


{\footnotesize\bfseries
No: 2\\
Name: cm\_\allowbreak{}4\_\allowbreak{}2\_\allowbreak{}2\\
Type: string\\
Default: nil
}

{\footnotesize заходи захисту, на які було здійснено вплив, реалізовані правильно з точки зору відповідності вимогам конфіденційності системи після внесення змін до системи;}


{\footnotesize\bfseries
No: 3\\
Name: cm\_\allowbreak{}4\_\allowbreak{}2\_\allowbreak{}3\\
Type: string\\
Default: nil
}

{\footnotesize заходи захисту, на які було здійснено вплив, функціонують належним чином з точки зору відповідності вимогам безпеки системи після внесення змін до системи;}


{\footnotesize\bfseries
No: 4\\
Name: cm\_\allowbreak{}4\_\allowbreak{}2\_\allowbreak{}4\\
Type: string\\
Default: nil
}

{\footnotesize заходи захисту, на які було  здійснено вплив, функціонують належним чином з точки зору відповідності вимогам конфіденційності системи після внесення змін до системи;}


{\footnotesize\bfseries
No: 5\\
Name: cm\_\allowbreak{}4\_\allowbreak{}2\_\allowbreak{}5\\
Type: string\\
Default: nil
}

{\footnotesize заходи захисту, на які було здійснено вплив, дають бажаний результат з точки зору відповідності вимогам безпеки системи після внесення змін до системи;}


{\footnotesize\bfseries
No: 6\\
Name: cm\_\allowbreak{}4\_\allowbreak{}2\_\allowbreak{}6\\
Type: string\\
Default: nil
}

{\footnotesize заходи захисту, на які було здійснено вплив, дають бажаний результат з точки зору відповідності вимогам конфіденційності системи після внесення змін до системи; 188}




\subsection{ОБМЕЖЕННЯ ДОСТУПУ ДО ЗМІНИ (CM-5)}
Визначити, задокументувати, затвердити та забезпечити застосування фізичних і логічних обмежень доступу, пов'язаних зі змінами в системі.

\vspace{10pt}
{\footnotesize\bfseries
No: 1\\
Name: cm\_\allowbreak{}5\_\allowbreak{}1\\
Type: string\\
Default: nil
}

{\footnotesize визначені та задокументовані фізичні обмеження доступу, пов'язані зі змінами в системі;}


{\footnotesize\bfseries
No: 2\\
Name: cm\_\allowbreak{}5\_\allowbreak{}2\\
Type: string\\
Default: nil
}

{\footnotesize затверджені фізичні обмеження доступу, пов'язані зі змінами в системі;}


{\footnotesize\bfseries
No: 3\\
Name: cm\_\allowbreak{}5\_\allowbreak{}3\\
Type: string\\
Default: nil
}

{\footnotesize застосовуються фізичні обмеження доступу, пов'язані зі змінами в системі;}


{\footnotesize\bfseries
No: 4\\
Name: cm\_\allowbreak{}5\_\allowbreak{}4\\
Type: string\\
Default: nil
}

{\footnotesize визначені та задокументовані логічні обмеження доступу, пов'язані зі змінами в системі;}


{\footnotesize\bfseries
No: 5\\
Name: cm\_\allowbreak{}5\_\allowbreak{}5\\
Type: string\\
Default: nil
}

{\footnotesize затверджені логічні обмеження доступу, пов'язані зі змінами в системі;}


{\footnotesize\bfseries
No: 6\\
Name: cm\_\allowbreak{}5\_\allowbreak{}6\\
Type: string\\
Default: nil
}

{\footnotesize застосовуються логічні обмеження доступу, пов'язані зі змінами в системі;}




\subsubsection{АУДИТ І ЗДІЙСНЕННЯ АВТОМАТИЧНОГО ДОСТУПУ (CM-5(1))}


\vspace{10pt}
{\footnotesize\bfseries
No: 1\\
Name: cm\_\allowbreak{}5\_\allowbreak{}1\_\allowbreak{}odp\\
Type: string\\
Default: nil
}

{\footnotesize визначено механізми, що використовуються для автоматизації застосування обмежень доступу;}


{\footnotesize\bfseries
No: 2\\
Name: cm\_\allowbreak{}5\_\allowbreak{}1\_\allowbreak{}a\\
Type: string\\
Default: nil
}

{\footnotesize обмеження доступу до змін застосовуються за допомогою <CM-05(01)\_\allowbreak{}ODP автоматизованих механізмів>;}


{\footnotesize\bfseries
No: 3\\
Name: cm\_\allowbreak{}5\_\allowbreak{}1\_\allowbreak{}b\\
Type: string\\
Default: nil
}

{\footnotesize автоматично формуються записи аудиту для виконаних дій.}




\subsubsection{ПЕРЕГЛЯД ЗМІН У СИСТЕМІ (CM-5(2)) [Вилучено]}
[Вилучено: перенесено до СМ-03(07)]

\vspace{10pt}
\textit{Немає параметрів для цього контролю.}
\vspace{10pt}


\subsubsection{ПІДПИСАНІ КОМПОНЕНТИ (CM-5(3)) [Вилучено]}
[Вилучено: перенесено до СМ-14]

\vspace{10pt}
\textit{Немає параметрів для цього контролю.}
\vspace{10pt}


\subsubsection{ПОДВІЙНА АВТОРИЗАЦІЯ (CM-5(4))}
Здійснювати подвійну авторизацію для внесення змін до [Призначення: компонентів системи та інформації на рівні системи, визначених організацією].

\vspace{10pt}
{\footnotesize\bfseries
No: 1\\
Name: cm\_\allowbreak{}5\_\allowbreak{}4\_\allowbreak{}odp\_\allowbreak{}1\\
Type: string\\
Default: nil
}

{\footnotesize визначено компоненти системи, що потребують 190 подвійної авторизації для внесення змін;}


{\footnotesize\bfseries
No: 2\\
Name: cm\_\allowbreak{}5\_\allowbreak{}4\_\allowbreak{}odp\_\allowbreak{}2\\
Type: string\\
Default: nil
}

{\footnotesize визначено інформацію на рівні системи , що потребують подвійної авторизації для внесення змін;}


{\footnotesize\bfseries
No: 3\\
Name: cm\_\allowbreak{}5\_\allowbreak{}4\_\allowbreak{}1\\
Type: string\\
Default: nil
}

{\footnotesize запроваджено подвійну авторизацію для внесення змін}


{\footnotesize\bfseries
No: 4\\
Name: cm\_\allowbreak{}5\_\allowbreak{}4\_\allowbreak{}2\\
Type: string\\
Default: nil
}

{\footnotesize запроваджено подвійну авторизацію для внесення змін}




\subsubsection{ОБМЕЖЕННЯ ПОВНОВАЖЕНЬ ДЛЯ ВИРОБНИЦТВА ТА ЕКСПЛУАТАЦІЇ (CM-5(5))}
(a) обмежити повноваження для зміни компонентів системи та інформації, пов'язаної із системою, у виробничому або операційному середовищі;\\ 
(b) переглядати та переоцінювати повноваження [Призначення: визначеною організацією з частотою].

\vspace{10pt}
{\footnotesize\bfseries
No: 1\\
Name: cm\_\allowbreak{}5\_\allowbreak{}5\_\allowbreak{}odp\_\allowbreak{}1\\
Type: string\\
Default: nil
}

{\footnotesize визначено частоту перегляду повноважень;}


{\footnotesize\bfseries
No: 2\\
Name: cm\_\allowbreak{}5\_\allowbreak{}5\_\allowbreak{}odp\_\allowbreak{}2\\
Type: string\\
Default: nil
}

{\footnotesize визначено частоту переоцінення повноважень;}


{\footnotesize\bfseries
No: 3\\
Name: cm\_\allowbreak{}5\_\allowbreak{}5\_\allowbreak{}a\_\allowbreak{}1\\
Type: string\\
Default: nil
}

{\footnotesize повноваження для зміни компонентів системи у виробничому або операційному середовищі обмежені;}


{\footnotesize\bfseries
No: 4\\
Name: cm\_\allowbreak{}5\_\allowbreak{}5\_\allowbreak{}a\_\allowbreak{}2\\
Type: string\\
Default: nil
}

{\footnotesize повноваження для зміни інформації, пов'язаної із системою у виробничому або операційному середовищі обмежені;}


{\footnotesize\bfseries
No: 5\\
Name: cm\_\allowbreak{}5\_\allowbreak{}5\_\allowbreak{}b\_\allowbreak{}1\\
Type: string\\
Default: nil
}

{\footnotesize переглядаються повноваження <CM-05(05)\_\allowbreak{}ODP[01] частота>;}


{\footnotesize\bfseries
No: 6\\
Name: cm\_\allowbreak{}5\_\allowbreak{}5\_\allowbreak{}b\_\allowbreak{}2\\
Type: string\\
Default: nil
}

{\footnotesize переоцінюються повноваження <CM-05(05)\_\allowbreak{}ODP[02] частота>;}




\subsubsection{ОБМЕЖЕННЯ ПОВНОВАЖЕНЬ ДЛЯ БІБЛІОТЕК (CM-5(6))}
Обмежити повноваження для зміни програмного забезпечення, яке перебуває в бібліотеках програмного забезпечення.

\vspace{10pt}
{\footnotesize\bfseries
No: 1\\
Name: cm\_\allowbreak{}5\_\allowbreak{}6\_\allowbreak{}01\\
Type: string\\
Default: nil
}

{\footnotesize обмежено повноваження для зміни програмного забезпечення, яке перебуває в бібліотеках програмного забезпечення}




\subsubsection{ОБМЕЖЕННЯ ДОСТУПУ ДО ЗМІНИ ВАДЖЕННЯ ЗАХОДІВ ЗАХИСТУ (CM-5(7)) [Вилучено]}
[Вилучено: включено до SI-07]

\vspace{10pt}
\textit{Немає параметрів для цього контролю.}
\vspace{10pt}


\subsection{НАЛАШТУВАННЯ КОНФІГУРАЦІЇ (CM-6)}
a. Встановити та задокументувати параметри конфігурації компонентів, які застосовуються в системі, які відображають найбільш обмежений режим, що відповідає експлуатаційним вимогам, використовуючи [Призначення: визначені організацією загальні безпечні конфігурації].\\ 
b. Реалізувати конфігураційні установки.\\ 
c. Визначити, задокументувати та затвердити будь-які відхилення від встановлених конфігураційних параметрів конфігурації для [Призначення: визначених організацією компонентів системи] на основі [Призначення: визначених організацією експлуатаційних вимог].\\ 
d. Відстежувати та керувати змінами конфігураційних параметрів відповідно до організаційної політики та процедур.

\vspace{10pt}
{\footnotesize\bfseries
No: 1\\
Name: cm\_\allowbreak{}6\_\allowbreak{}odp\_\allowbreak{}1\\
Type: string\\
Default: nil
}

{\footnotesize визначено загальні безпечні конфігурації для встановлення та документування параметрів конфігурації компонентів, які застосовуються в системі; 192}


{\footnotesize\bfseries
No: 2\\
Name: cm\_\allowbreak{}6\_\allowbreak{}odp\_\allowbreak{}2\\
Type: string\\
Default: nil
}

{\footnotesize визначено компоненти системи, для яких необхідно затвердити відхилення;}


{\footnotesize\bfseries
No: 3\\
Name: cm\_\allowbreak{}6\_\allowbreak{}odp\_\allowbreak{}3\\
Type: string\\
Default: nil
}

{\footnotesize визначені експлуатаційні вимоги, що вимагають затвердження відхилень;}


{\footnotesize\bfseries
No: 4\\
Name: cm\_\allowbreak{}6\_\allowbreak{}a\\
Type: string\\
Default: nil
}

{\footnotesize налаштування конфігурації, які відображають найбільш обмежувальний режим, що відповідає експлуатаційним вимогам, встановлені та задокументовані для компонентів, що застосовуються}


{\footnotesize\bfseries
No: 5\\
Name: cm\_\allowbreak{}6\_\allowbreak{}b\\
Type: string\\
Default: nil
}

{\footnotesize реалізовано установки конфігурації, задокументовані в CM-06a;}


{\footnotesize\bfseries
No: 6\\
Name: cm\_\allowbreak{}6\_\allowbreak{}c\_\allowbreak{}1\\
Type: string\\
Default: nil
}

{\footnotesize будь-які відхилення від встановлених параметрів конфігурації вимог>;}


{\footnotesize\bfseries
No: 7\\
Name: cm\_\allowbreak{}6\_\allowbreak{}c\_\allowbreak{}2\\
Type: string\\
Default: nil
}

{\footnotesize будь-які відхилення від встановлених налаштувань конфігурації}


{\footnotesize\bfseries
No: 8\\
Name: cm\_\allowbreak{}6\_\allowbreak{}d\_\allowbreak{}1\\
Type: string\\
Default: nil
}

{\footnotesize зміни в налаштуваннях конфігурації відстежуються відповідно до політики та процедур організації;}


{\footnotesize\bfseries
No: 9\\
Name: cm\_\allowbreak{}6\_\allowbreak{}d\_\allowbreak{}2\\
Type: string\\
Default: nil
}

{\footnotesize зміни налаштувань конфігурації керуються відповідно до політики та процедур організації.}




\subsubsection{АВТОМАТИЗОВАНЕ УПРАВЛІННЯ, ЗАСТОСУВАННЯ ТА ВЕРИФІКАЦІЯ (CM-6(1))}


\vspace{10pt}
{\footnotesize\bfseries
No: 1\\
Name: cm\_\allowbreak{}6\_\allowbreak{}1\_\allowbreak{}odp\_\allowbreak{}1\\
Type: string\\
Default: nil
}

{\footnotesize визначено компоненти системи, для яких можна керувати, застосовувати та перевіряти налаштування конфігурації;}


{\footnotesize\bfseries
No: 2\\
Name: cm\_\allowbreak{}6\_\allowbreak{}1\_\allowbreak{}odp\_\allowbreak{}2\\
Type: string\\
Default: nil
}

{\footnotesize визначено автоматизовані механізми керування налаштуваннями конфігурації;}


{\footnotesize\bfseries
No: 3\\
Name: cm\_\allowbreak{}6\_\allowbreak{}1\_\allowbreak{}odp\_\allowbreak{}3\\
Type: string\\
Default: nil
}

{\footnotesize визначено автоматизовані механізми застосування налаштувань конфігурації;}


{\footnotesize\bfseries
No: 4\\
Name: cm\_\allowbreak{}6\_\allowbreak{}1\_\allowbreak{}odp\_\allowbreak{}4\\
Type: string\\
Default: nil
}

{\footnotesize визначено автоматизовані механізми перев ірки налаштувань конфігурації;}


{\footnotesize\bfseries
No: 5\\
Name: cm\_\allowbreak{}6\_\allowbreak{}1\_\allowbreak{}1\\
Type: string\\
Default: nil
}

{\footnotesize налаштування конфігурації для <CM-06(01)\_\allowbreak{}ODP[01]}


{\footnotesize\bfseries
No: 6\\
Name: cm\_\allowbreak{}6\_\allowbreak{}1\_\allowbreak{}2\\
Type: string\\
Default: nil
}

{\footnotesize налаштування конфігурації для <CM-06(01)\_\allowbreak{}ODP[01] компонентів системи>  застосовуються за допомогою}


{\footnotesize\bfseries
No: 7\\
Name: cm\_\allowbreak{}6\_\allowbreak{}1\_\allowbreak{}3\\
Type: string\\
Default: nil
}

{\footnotesize налаштування конфігурації для <CM-06(01)\_\allowbreak{}ODP[01] компонентів системи>  перевіряються за допомогою}




\subsubsection{НАЛАШТУВАННЯ КОНФІГУРАЦІЇ САНКЦІОНОВАНІ ЗМІНИ (CM-6(2))}
Виконайте такі дії у відповідь на неавторизовані зміни в [Призначення: параметри конфігурації, визначені організацією]: [Призначення: дії, визначені організацією].

\vspace{10pt}
{\footnotesize\bfseries
No: 1\\
Name: cm\_\allowbreak{}6\_\allowbreak{}2\_\allowbreak{}01\\
Type: string\\
Default: nil
}

{\footnotesize <CM-06(02)\_\allowbreak{}ODP[01] дії>  виконуються у відповідь на конфігурації>.}




\subsubsection{ДЕМОНСТРАЦІЯ ВІДПОВІДНОСТІ (CM-6(4)) [Вилучено]}
[Вилучено: Включено до CM-04]

\vspace{10pt}
\textit{Немає параметрів для цього контролю.}
\vspace{10pt}


\subsection{МІНІМАЛЬНО НЕОБХІДНА ФУНКЦІОНАЛЬНІСТЬ (CM-7)}
a. Налаштуйте систему для забезпечення лише [Призначення: основні функції, визначені організацією для місії];\\ 
b. Заборонити або обмежити використання таких функцій, портів, протоколів, програмного забезпечення та/або служб: [Призначення: визначені організацією заборонені або обмежені функції, системні порти, протоколи, програмне забезпечення та/або служби].

\vspace{10pt}
{\footnotesize\bfseries
No: 1\\
Name: cm\_\allowbreak{}7\_\allowbreak{}odp\_\allowbreak{}1\\
Type: string\\
Default: nil
}

{\footnotesize визначено основні функції системи, необхідні для виконання місії;}


{\footnotesize\bfseries
No: 2\\
Name: cm\_\allowbreak{}7\_\allowbreak{}odp\_\allowbreak{}2\\
Type: string\\
Default: nil
}

{\footnotesize визначено функції, які необхідно заборонити або обмежити;}


{\footnotesize\bfseries
No: 3\\
Name: cm\_\allowbreak{}7\_\allowbreak{}odp\_\allowbreak{}3\\
Type: string\\
Default: nil
}

{\footnotesize визначено системні порти, які необхідно заборонити або обмежити;}


{\footnotesize\bfseries
No: 4\\
Name: cm\_\allowbreak{}7\_\allowbreak{}odp\_\allowbreak{}4\\
Type: string\\
Default: nil
}

{\footnotesize визначено протоколи, які необхідно заборонити або обмежити; 195}


{\footnotesize\bfseries
No: 5\\
Name: cm\_\allowbreak{}7\_\allowbreak{}odp\_\allowbreak{}5\\
Type: string\\
Default: nil
}

{\footnotesize визначено програмне забезпечення, яке необхідно заборонити або обмежити;}


{\footnotesize\bfseries
No: 6\\
Name: cm\_\allowbreak{}7\_\allowbreak{}odp\_\allowbreak{}6\\
Type: string\\
Default: nil
}

{\footnotesize визначено служби, які необхідно заборонити або обмежити;}


{\footnotesize\bfseries
No: 7\\
Name: cm\_\allowbreak{}7\_\allowbreak{}a\\
Type: string\\
Default: nil
}

{\footnotesize система налаштована на забезпечення лише  <CM-07\_\allowbreak{}ODP[01] основні функції системи >;}


{\footnotesize\bfseries
No: 8\\
Name: cm\_\allowbreak{}7\_\allowbreak{}b\_\allowbreak{}1\\
Type: string\\
Default: nil
}

{\footnotesize використання <CM-07\_\allowbreak{}ODP[02] функцій>  заборонено або обмежено;}


{\footnotesize\bfseries
No: 9\\
Name: cm\_\allowbreak{}7\_\allowbreak{}b\_\allowbreak{}2\\
Type: string\\
Default: nil
}

{\footnotesize використання <CM-07\_\allowbreak{}ODP[03] порти> заборонено або обмежено;}


{\footnotesize\bfseries
No: 10\\
Name: cm\_\allowbreak{}7\_\allowbreak{}b\_\allowbreak{}3\\
Type: string\\
Default: nil
}

{\footnotesize використання <CM-07\_\allowbreak{}ODP[04] протоколи> заборонено або обмежено;}


{\footnotesize\bfseries
No: 11\\
Name: cm\_\allowbreak{}7\_\allowbreak{}b\_\allowbreak{}4\\
Type: string\\
Default: nil
}

{\footnotesize використання <CM-07\_\allowbreak{}ODP[05] програмне забезпечення > заборонено або обмежено;}


{\footnotesize\bfseries
No: 12\\
Name: cm\_\allowbreak{}7\_\allowbreak{}b\_\allowbreak{}5\\
Type: string\\
Default: nil
}

{\footnotesize використання <CM-07\_\allowbreak{}ODP[06] служби> заборонено або обмежено;}




\subsubsection{ПЕРІОДИЧНИЙ ПЕРЕГЛЯД (CM-7(1))}


\vspace{10pt}
{\footnotesize\bfseries
No: 1\\
Name: cm\_\allowbreak{}7\_\allowbreak{}1\_\allowbreak{}odp\_\allowbreak{}1\\
Type: string\\
Default: nil
}

{\footnotesize визначено частоту, з якою слід переглядати систему для виявлення непотрібних та/або незахищених функцій, портів, протоколів і послуг;}


{\footnotesize\bfseries
No: 2\\
Name: cm\_\allowbreak{}7\_\allowbreak{}1\_\allowbreak{}odp\_\allowbreak{}2\\
Type: string\\
Default: nil
}

{\footnotesize визначено функції, які слід вимкнути, якщо вони вважаються непотрібними або незахищеними; 196}


{\footnotesize\bfseries
No: 3\\
Name: cm\_\allowbreak{}7\_\allowbreak{}1\_\allowbreak{}odp\_\allowbreak{}3\\
Type: string\\
Default: nil
}

{\footnotesize визначено порти, які слід вимкнути, якщо вони вважаються непотрібними або незахищеними;}


{\footnotesize\bfseries
No: 4\\
Name: cm\_\allowbreak{}7\_\allowbreak{}1\_\allowbreak{}odp\_\allowbreak{}4\\
Type: string\\
Default: nil
}

{\footnotesize визначено протоколи, які слід вимкнути, якщо вони вважаються непотрібними або незахищеними;}


{\footnotesize\bfseries
No: 5\\
Name: cm\_\allowbreak{}7\_\allowbreak{}1\_\allowbreak{}odp\_\allowbreak{}5\\
Type: string\\
Default: nil
}

{\footnotesize визначено послуги, які слід вимкнути, якщо вони вважаються непотрібними або незахищеними;}


{\footnotesize\bfseries
No: 6\\
Name: cm\_\allowbreak{}7\_\allowbreak{}1\_\allowbreak{}a\\
Type: string\\
Default: nil
}

{\footnotesize система переглядається <CM-07(01)\_\allowbreak{}ODP[01] частота> для виявлення непотрібних та/або незахищених функцій, портів, протоколів і послуг;}


{\footnotesize\bfseries
No: 7\\
Name: cm\_\allowbreak{}7\_\allowbreak{}1\_\allowbreak{}b\_\allowbreak{}1\\
Type: string\\
Default: nil
}

{\footnotesize <CM-07(01)\_\allowbreak{}ODP[02] функції> , які вважаються непотрібними та/або незахищеними, вимкнено;}


{\footnotesize\bfseries
No: 8\\
Name: cm\_\allowbreak{}7\_\allowbreak{}1\_\allowbreak{}b\_\allowbreak{}2\\
Type: string\\
Default: nil
}

{\footnotesize <CM-07(01)\_\allowbreak{}ODP[03] порти> , які вважаються непотрібними та/або незахищеними, вимкнено;}


{\footnotesize\bfseries
No: 9\\
Name: cm\_\allowbreak{}7\_\allowbreak{}1\_\allowbreak{}b\_\allowbreak{}3\\
Type: string\\
Default: nil
}

{\footnotesize <CM-07(01)\_\allowbreak{}ODP[04] протоколи> , які вважаються непотрібними та/або незахищеними, вимкнено;}


{\footnotesize\bfseries
No: 10\\
Name: cm\_\allowbreak{}7\_\allowbreak{}1\_\allowbreak{}b\_\allowbreak{}4\\
Type: string\\
Default: nil
}

{\footnotesize <CM-07(01)\_\allowbreak{}ODP[05] послуги>, які вважаються непотрібними та/або незахищеними, вимкнено;}




\subsubsection{ЗАБОРОНА ВИКОНАННЯ ПРОГРАМИ (CM-7(2))}
Заборонити виконання програми відповідно до [Вибір (один або кілька): [Призначення: визначеної організацією політики, правил поведінки та/або угод про доступ щодо використання програмного забезпечення та обмежень]; правил, що встановлюють терміни та умови використання програмного забезпечення].

\vspace{10pt}
{\footnotesize\bfseries
No: 1\\
Name: cm\_\allowbreak{}7\_\allowbreak{}2\_\allowbreak{}odp\_\allowbreak{}1\\
Type: string\\
Default: nil
}

{\footnotesize вибрано одне або декілька з наступних ЗНАЧЕНЬ стання програмного забезпечення та обмежень>; правила, що встановлюють терміни та умови використання програмного забезпечення\};}


{\footnotesize\bfseries
No: 2\\
Name: cm\_\allowbreak{}7\_\allowbreak{}2\_\allowbreak{}odp\_\allowbreak{}2\\
Type: string\\
Default: nil
}

{\footnotesize визначені політики, правил поведінки та/або угод про доступ щодо використання програмного забезпечення та обмежень (якщо вибрано);}


{\footnotesize\bfseries
No: 3\\
Name: cm\_\allowbreak{}7\_\allowbreak{}2\_\allowbreak{}01\\
Type: string\\
Default: nil
}

{\footnotesize виконання програми заборонено відповідно до <CM07(02)\_\allowbreak{}ODP[01] ВИБРАНЕ ЗНАЧЕННЯ ПАРАМЕТРА(ів)>.}




\subsubsection{ВІДПОВІДНІСТЬ РЕЄСТРАЦІЇ (CM-7(3))}
Забезпечити відповідність [Призначення: визначеним організацією вимогам до реєстрації для функцій, портів, протоколів і послуг].

\vspace{10pt}
{\footnotesize\bfseries
No: 1\\
Name: cm\_\allowbreak{}7\_\allowbreak{}3\_\allowbreak{}odp\\
Type: string\\
Default: nil
}

{\footnotesize визначено вимоги до реєстрації функцій, портів, протоколів та сервісів;}


{\footnotesize\bfseries
No: 2\\
Name: cm\_\allowbreak{}7\_\allowbreak{}3\_\allowbreak{}01\\
Type: string\\
Default: nil
}

{\footnotesize <CM-07(03)\_\allowbreak{}ODP вимоги до реєстрації> дотримано.}




\subsubsection{НЕАВТОРИЗОВАНЕ ПРОГРАМНЕ ЗАБЕЗПЕЧЕННЯ - ЧОРНИЙ СПИСОК (CM-7(4))}
(a) Визначити [Призначення: визначені організацією програмне забезпечення, що не має дозволу виконуватися в системі].\\ 
(b) Впровадити політику «дозволу всього, за винятком деяких» для заборони виконання неавторизованих програм у системі\\ 
(c) Переглядати та оновлювати список неавторизованих програм [Призначення: з визначеною організацією частотою].

\vspace{10pt}
{\footnotesize\bfseries
No: 1\\
Name: cm\_\allowbreak{}7\_\allowbreak{}4\_\allowbreak{}odp\_\allowbreak{}1\\
Type: string\\
Default: nil
}

{\footnotesize визначено програмне забезпечення, яке не має дозволу на виконання в системі;}


{\footnotesize\bfseries
No: 2\\
Name: cm\_\allowbreak{}7\_\allowbreak{}4\_\allowbreak{}odp\_\allowbreak{}2\\
Type: string\\
Default: nil
}

{\footnotesize визначено частоту, з якою слід переглядати та оновлювати список неавторизованих програм;}


{\footnotesize\bfseries
No: 3\\
Name: cm\_\allowbreak{}7\_\allowbreak{}4\_\allowbreak{}a\\
Type: string\\
Default: nil
}

{\footnotesize визначено <CM-07(04)\_\allowbreak{}ODP[01] програмне забезпечення>;}


{\footnotesize\bfseries
No: 4\\
Name: cm\_\allowbreak{}7\_\allowbreak{}4\_\allowbreak{}b\\
Type: string\\
Default: nil
}

{\footnotesize політика \textquotedbl{}дозволу всього, за винятком деяких\textquotedbl{} застосовується для заборони виконання неавторизованих програм у системі;}


{\footnotesize\bfseries
No: 5\\
Name: cm\_\allowbreak{}7\_\allowbreak{}4\_\allowbreak{}c\\
Type: string\\
Default: nil
}

{\footnotesize переглядається та оновлюється список неавторизованих}




\subsubsection{АВТОРИЗОВАНЕ ПРОГРАМНЕ ЗАБЕЗПЕЧЕННЯ -- БІЛИЙ СПИСОК (CM-7(5))}
(a) Визначити [Призначення: визначені організацією програмне забезпечення, яке авторизовано виконується в системі]\\ 
(b) Впровадити політику «заборони всього, за винятком деяких», щоб дозволити виконання авторизованих програм у системі.\\ 
(c) Переглядати та оновлювати список авторизованих програм [Призначення: з визначеною організацією частотою].

\vspace{10pt}
{\footnotesize\bfseries
No: 1\\
Name: cm\_\allowbreak{}7\_\allowbreak{}5\_\allowbreak{}odp\_\allowbreak{}1\\
Type: string\\
Default: nil
}

{\footnotesize визначено програмне забезпечення, яке авторизовано 199 для виконання в системі;}


{\footnotesize\bfseries
No: 2\\
Name: cm\_\allowbreak{}7\_\allowbreak{}5\_\allowbreak{}odp\_\allowbreak{}2\\
Type: string\\
Default: nil
}

{\footnotesize визначено частоту перегляду та оновлення списку авторизованих програм;}


{\footnotesize\bfseries
No: 3\\
Name: cm\_\allowbreak{}7\_\allowbreak{}5\_\allowbreak{}a\\
Type: string\\
Default: nil
}

{\footnotesize визначено <CM-07(05)\_\allowbreak{}ODP[01] програмне забезпечення>;}


{\footnotesize\bfseries
No: 4\\
Name: cm\_\allowbreak{}7\_\allowbreak{}5\_\allowbreak{}b\\
Type: string\\
Default: nil
}

{\footnotesize політика \textquotedbl{}заборонити все, дозволити за винятком\textquotedbl{} застосовуэться, щоб дозволити виконання авторизованих програм у системі;}


{\footnotesize\bfseries
No: 5\\
Name: cm\_\allowbreak{}7\_\allowbreak{}5\_\allowbreak{}c\\
Type: string\\
Default: nil
}

{\footnotesize переглядається та оновлюється список  авторизованих}




\subsubsection{ЗАМКНУТІ СЕРЕДОВИЩА З ОБМЕЖЕНИМИ ПРИВІЛЕЯМИ (CM-7(6))}
Вимагайте, щоб визначене програмне забезпечення, встановлене користувачем, виконувалося в обмеженому середовищі фізичної або віртуальної машини з обмеженими привілеями: [Призначення: програмне забезпечення, встановлене користувачем, визначене організацією]

\vspace{10pt}
{\footnotesize\bfseries
No: 1\\
Name: cm\_\allowbreak{}7\_\allowbreak{}6\_\allowbreak{}odp\\
Type: string\\
Default: nil
}

{\footnotesize визначено встановлене користувачем програмне забезпечення, яке потрібно виконувати в обмеженому середовищі;}


{\footnotesize\bfseries
No: 2\\
Name: cm\_\allowbreak{}7\_\allowbreak{}6\_\allowbreak{}01\\
Type: string\\
Default: nil
}

{\footnotesize <CM-07(06)\_\allowbreak{}ODP програмне забезпечення, встановлене користувачем> має виконуватися в обмеженому середовищі фізичної або віртуальної машини з обмеженими привілеями.}




\subsubsection{ВИКОНУВАНИЙ КОД У ЗАХИЩЕНОМУ СЕРЕДОВИЩІ (CM-7(7))}
Дозволити виконання двійкового або машинно-виконуваного коду лише в обмеженому фізичному або віртуальному машинному середовищі та за явного дозволу [Призначення: персонал або ролі, визначені організацією], якщо такий код: a) Отримано з джерел з обмеженою гарантією або без неї; та/або b) Без надання вихідного коду.

\vspace{10pt}
{\footnotesize\bfseries
No: 1\\
Name: cm\_\allowbreak{}7\_\allowbreak{}7\_\allowbreak{}odp\\
Type: string\\
Default: nil
}

{\footnotesize визначено персонал або ролі для явного дозволу на виконання двійкового або машинно-виконуваного коду;}


{\footnotesize\bfseries
No: 2\\
Name: cm\_\allowbreak{}7\_\allowbreak{}7\_\allowbreak{}01\\
Type: string\\
Default: nil
}

{\footnotesize виконання двійкового або машинного коду дозволено лише в обмеженому середовищі фізичної або віртуальної машини;}


{\footnotesize\bfseries
No: 3\\
Name: cm\_\allowbreak{}7\_\allowbreak{}7\_\allowbreak{}a\\
Type: string\\
Default: nil
}

{\footnotesize виконання двійкового або машинного коду, отриманого з джерел з обмеженою гарантією або без неї, дозволяється лише з явного дозволу  <CM-07(07)\_\allowbreak{}ODP персонал або ролі>;}


{\footnotesize\bfseries
No: 4\\
Name: cm\_\allowbreak{}7\_\allowbreak{}7\_\allowbreak{}b\\
Type: string\\
Default: nil
}

{\footnotesize виконання двійкового або машинного коду без надання вихідного коду дозволяється лише з явного дозволу < CM07(07)\_\allowbreak{}ODP персонал або ролі>.}




\subsubsection{БІНАРНИЙ АБО МАШИННИЙ ВИКОНУВАНИЙ КОД (CM-7(8))}
a) Заборонити використання двійкового або машинно-виконуваного коду з джерел з обмеженою гарантією або без неї або без надання вихідного коду; і b) Дозволяти винятки лише для обов'язкових місій або оперативних вимог і за погодженням з уповноваженою посадовою особою.

\vspace{10pt}
{\footnotesize\bfseries
No: 1\\
Name: cm\_\allowbreak{}7\_\allowbreak{}8\_\allowbreak{}a\\
Type: string\\
Default: nil
}

{\footnotesize використання двійкового або машинного коду заборонено, якщо він походить з джерел з обмеженою гарантією або без неї, або без надання вихідного коду;}


{\footnotesize\bfseries
No: 2\\
Name: cm\_\allowbreak{}7\_\allowbreak{}8\_\allowbreak{}b\_\allowbreak{}1\\
Type: string\\
Default: nil
}

{\footnotesize винятки із заборони на використання двійкового або машинного коду з джерел з обмеженою гарантією або без неї, або без надання вихідного коду допускаються лише для обов'язкових місій або оперативних вимог;}


{\footnotesize\bfseries
No: 3\\
Name: cm\_\allowbreak{}7\_\allowbreak{}8\_\allowbreak{}b\_\allowbreak{}2\\
Type: string\\
Default: nil
}

{\footnotesize винятки із заборони на використання двійкового або машинного коду з джерел з обмеженою гарантією або без неї, або без надання вихідного коду допускаються лише у разі погодження з уповноваженою посадовою особою;}




\subsubsection{ЗАБОРОНА ВИКОРИСТАННЯ НЕАВТОРИЗОВАНОГО ОБЛАДНАННЯ (CM-7(9))}
a.\\ 
b. c. ЗАБОРОНА Визначити [Призначення: апаратні компоненти, визначені організацією, авторизовані для використання в системі]; Заборонити використання або підключення неавторизованих апаратних компонентів; Перегляд та оновлення списку авторизованих апаратних компонентів [Призначення: частота, визначена організацією].

\vspace{10pt}
{\footnotesize\bfseries
No: 1\\
Name: cm\_\allowbreak{}7\_\allowbreak{}9\_\allowbreak{}odp\_\allowbreak{}1\\
Type: string\\
Default: nil
}

{\footnotesize визначено апаратні компоненти, дозволені для використання в системі;}


{\footnotesize\bfseries
No: 2\\
Name: cm\_\allowbreak{}7\_\allowbreak{}9\_\allowbreak{}a\\
Type: string\\
Default: nil
}

{\footnotesize ідентифіковано < CM-07(09)\_\allowbreak{}ODP[01] апаратні компоненти>;}


{\footnotesize\bfseries
No: 3\\
Name: cm\_\allowbreak{}7\_\allowbreak{}9\_\allowbreak{}b\\
Type: string\\
Default: nil
}

{\footnotesize використання або підключення несанкціонованих апаратних компонентів заборонено;}


{\footnotesize\bfseries
No: 4\\
Name: cm\_\allowbreak{}7\_\allowbreak{}9\_\allowbreak{}c\\
Type: string\\
Default: nil
}

{\footnotesize список дозволених апаратних компонентів переглядається та оновлюється з <CM-07(09)\_\allowbreak{}ODP[02] частота>;.}




\subsection{ІНВЕНТАРИЗАЦІЯ КОМПОНЕНТІВ СИСТЕМИ (CM-8)}
a.\\ 
b. Розробити та задокументувати процес інвентаризації компонентів системи, який:\\ 
1. точно описує поточну систему;\\ 
2. охоплює всі компоненти в межах акредитації системи;\\ 
3. не включає повторний облік компонентів або компонентів, будь-якої іншої системи;\\ 
4. визначає рівень деталізації, який є необхідним для відстеження та звітування;\\ 
5. включає інформацію для досягнення підзвітності компонентів системи: [Призначення: визначена організацією інформація, необхідна для досягнення ефективної підзвітності компонентів системи]. Переглядати та оновлювати опис компонентів системи з [Призначення: визначеною організацією частотою].

\vspace{10pt}
{\footnotesize\bfseries
No: 1\\
Name: cm\_\allowbreak{}8\_\allowbreak{}odp\_\allowbreak{}1\\
Type: string\\
Default: nil
}

{\footnotesize визначено інформацію, яка вважається необхідною для досягнення ефективної підзвітності компонентів системи;}


{\footnotesize\bfseries
No: 2\\
Name: cm\_\allowbreak{}8\_\allowbreak{}odp\_\allowbreak{}2\\
Type: string\\
Default: nil
}

{\footnotesize визначено частоту перегляду та оновлення опису компонентів системи;}


{\footnotesize\bfseries
No: 3\\
Name: cm\_\allowbreak{}8\_\allowbreak{}a\_\allowbreak{}1\\
Type: string\\
Default: nil
}

{\footnotesize розроблено та задокументовано процес інвентаризації компонентів системи, який точно описує поточну систему;}


{\footnotesize\bfseries
No: 4\\
Name: cm\_\allowbreak{}8\_\allowbreak{}a\_\allowbreak{}2\\
Type: string\\
Default: nil
}

{\footnotesize розроблено та задокументовано процес інвентаризації компонентів системи, який охоплює всі компоненти в межах акредитації системи;}


{\footnotesize\bfseries
No: 5\\
Name: cm\_\allowbreak{}8\_\allowbreak{}a\_\allowbreak{}3\\
Type: string\\
Default: nil
}

{\footnotesize розроблено та задокументовано процес інвентаризації компонентів системи, який не включає повторний облік компонентів 203 або компонентів, будь-якої іншої системи;}


{\footnotesize\bfseries
No: 6\\
Name: cm\_\allowbreak{}8\_\allowbreak{}a\_\allowbreak{}4\\
Type: string\\
Default: nil
}

{\footnotesize розроблено та задокументовано процес інвентаризації компонентів системи, який визначає рівень деталізації, який є необхідним для відстеження та звітування;}


{\footnotesize\bfseries
No: 7\\
Name: cm\_\allowbreak{}8\_\allowbreak{}a\_\allowbreak{}5\\
Type: string\\
Default: nil
}

{\footnotesize розроблено та задокументовано процес інвентаризації компонентів системи, який включає <CM-08\_\allowbreak{}ODP[01] інформацію>;}


{\footnotesize\bfseries
No: 8\\
Name: cm\_\allowbreak{}8\_\allowbreak{}b\\
Type: string\\
Default: nil
}

{\footnotesize переглядається та оновлюється опис компонентів системи}




\subsubsection{ОНОВЛЕННЯ ПІД ЧАС ВСТАНОВЛЕННЯ ТА ВИДАЛЕННЯ (CM-8(1))}


\vspace{10pt}
{\footnotesize\bfseries
No: 1\\
Name: cm\_\allowbreak{}8\_\allowbreak{}1\_\allowbreak{}1\\
Type: string\\
Default: nil
}

{\footnotesize інвентаризація компонентів системи оновлюється в рамках інсталяцій компонентів системи;}


{\footnotesize\bfseries
No: 2\\
Name: cm\_\allowbreak{}8\_\allowbreak{}1\_\allowbreak{}2\\
Type: string\\
Default: nil
}

{\footnotesize інвентаризація компонентів системи оновлюється в рамках видалення компонентів системи;}


{\footnotesize\bfseries
No: 3\\
Name: cm\_\allowbreak{}8\_\allowbreak{}1\_\allowbreak{}3\\
Type: string\\
Default: nil
}

{\footnotesize інвентаризація компонентів системи оновлюється в рамках оновлення компонентів системи;}




\subsubsection{АВТОМАТИЗОВАНА ПІДТРИМКА (CM-8(2))}


\vspace{10pt}
{\footnotesize\bfseries
No: 1\\
Name: cm\_\allowbreak{}8\_\allowbreak{}2\_\allowbreak{}odp\_\allowbreak{}1\\
Type: string\\
Default: nil
}

{\footnotesize визначено автоматизовані механізми підтримки актуальності інвентаризації компонентів системи;}


{\footnotesize\bfseries
No: 2\\
Name: cm\_\allowbreak{}8\_\allowbreak{}2\_\allowbreak{}odp\_\allowbreak{}2\\
Type: string\\
Default: nil
}

{\footnotesize визначено автоматизовані механізми підтримки повноти інвентаризації компонентів системи;}


{\footnotesize\bfseries
No: 3\\
Name: cm\_\allowbreak{}8\_\allowbreak{}2\_\allowbreak{}odp\_\allowbreak{}3\\
Type: string\\
Default: nil
}

{\footnotesize визначено автоматизовані механізми підтримки точності інвентаризації компонентів системи;}


{\footnotesize\bfseries
No: 4\\
Name: cm\_\allowbreak{}8\_\allowbreak{}2\_\allowbreak{}odp\_\allowbreak{}4\\
Type: string\\
Default: nil
}

{\footnotesize визначено автоматизовані механізми підтримки доступності інвентаризації компонентів системи;}


{\footnotesize\bfseries
No: 5\\
Name: cm\_\allowbreak{}8\_\allowbreak{}2\_\allowbreak{}1\\
Type: string\\
Default: nil
}

{\footnotesize <CM-08(02)\_\allowbreak{}ODP[01] автоматизовані механізми>  використовуються для підтримки актуальності інвентаризації компонентів системи;}


{\footnotesize\bfseries
No: 6\\
Name: cm\_\allowbreak{}8\_\allowbreak{}2\_\allowbreak{}2\\
Type: string\\
Default: nil
}

{\footnotesize <CM-08(02)\_\allowbreak{}ODP[02] автоматизовані механізми>  використовуються для підтримки повноти інвентаризації компонентів системи;}


{\footnotesize\bfseries
No: 7\\
Name: cm\_\allowbreak{}8\_\allowbreak{}2\_\allowbreak{}3\\
Type: string\\
Default: nil
}

{\footnotesize <CM-08(02)\_\allowbreak{}ODP[03] автоматизовані механізми> використовуються для підтримки точності інвентаризації компонентів системи;}


{\footnotesize\bfseries
No: 8\\
Name: cm\_\allowbreak{}8\_\allowbreak{}2\_\allowbreak{}4\\
Type: string\\
Default: nil
}

{\footnotesize <CM-08(02)\_\allowbreak{}ODP[04] автоматизовані механізми>  використовуються для підтримки доступності інвентаризації компонентів системи;}




\subsubsection{АВТОМАТИЗОВАНЕ ВИЯВЛЕННЯ НЕАВТОРИЗОВАНИХ КОМПОНЕНТІВ (CM-8(3))}


\vspace{10pt}
{\footnotesize\bfseries
No: 1\\
Name: cm\_\allowbreak{}8\_\allowbreak{}3\_\allowbreak{}odp\_\allowbreak{}1\\
Type: string\\
Default: nil
}

{\footnotesize визначено автоматизовані механізми, що використовуються для виявлення наявності несанкціонованого обладнання в системі;}


{\footnotesize\bfseries
No: 2\\
Name: cm\_\allowbreak{}8\_\allowbreak{}3\_\allowbreak{}odp\_\allowbreak{}2\\
Type: string\\
Default: nil
}

{\footnotesize визначено автоматизовані механізми, що використовуються для виявлення наявності несанкціонованого програмного забезпечення в системі;}


{\footnotesize\bfseries
No: 3\\
Name: cm\_\allowbreak{}8\_\allowbreak{}3\_\allowbreak{}odp\_\allowbreak{}3\\
Type: string\\
Default: nil
}

{\footnotesize визначено автоматизовані механізми, що використовуються для виявлення наявності несанкціонованих мікропрограмних компонентів в системі;}


{\footnotesize\bfseries
No: 4\\
Name: cm\_\allowbreak{}8\_\allowbreak{}3\_\allowbreak{}odp\_\allowbreak{}4\\
Type: string\\
Default: nil
}

{\footnotesize визначено частоту, з якою використовуються автоматизовані механізми для виявлення присутності несанкціонованих компонентів системи в системі;}


{\footnotesize\bfseries
No: 5\\
Name: cm\_\allowbreak{}8\_\allowbreak{}3\_\allowbreak{}odp\_\allowbreak{}5\\
Type: string\\
Default: nil
}

{\footnotesize вибрано одне або декілька з наступних ЗНАЧЕНЬ}


{\footnotesize\bfseries
No: 6\\
Name: cm\_\allowbreak{}8\_\allowbreak{}3\_\allowbreak{}odp\_\allowbreak{}6\\
Type: string\\
Default: nil
}

{\footnotesize визначено персонал або ролі, я кі мають бути повідомлені при виявленні несанкціонованих компонентів (якщо вибрано);}


{\footnotesize\bfseries
No: 7\\
Name: cm\_\allowbreak{}8\_\allowbreak{}3\_\allowbreak{}a\_\allowbreak{}1\\
Type: string\\
Default: nil
}

{\footnotesize наявність несанкціонованого обладнання в системі виявляється за допомогою  <CM-08(03)\_\allowbreak{}ODP[01] частота>;}


{\footnotesize\bfseries
No: 8\\
Name: cm\_\allowbreak{}8\_\allowbreak{}3\_\allowbreak{}a\_\allowbreak{}2\\
Type: string\\
Default: nil
}

{\footnotesize наявність несанкціонованого програмного забезпечення}


{\footnotesize\bfseries
No: 9\\
Name: cm\_\allowbreak{}8\_\allowbreak{}3\_\allowbreak{}a\_\allowbreak{}3\\
Type: string\\
Default: nil
}

{\footnotesize наявність несанкціонованих мікропрограмних компонентів в  системі виявляється за допомогою  <CM08(03)\_\allowbreak{}ODP[03] автоматизованих механізмів> <CM 08(03)\_\allowbreak{}ODP[04] частота>;}


{\footnotesize\bfseries
No: 10\\
Name: cm\_\allowbreak{}8\_\allowbreak{}3\_\allowbreak{}b\_\allowbreak{}1\\
Type: string\\
Default: nil
}

{\footnotesize <CM-08(03)\_\allowbreak{}ODP[05] ВИБРАНЕ ЗНАЧЕННЯ ПАРАМЕТРА(ів)> приймаються при виявленні несанкціонованого обладнання;}


{\footnotesize\bfseries
No: 11\\
Name: cm\_\allowbreak{}8\_\allowbreak{}3\_\allowbreak{}b\_\allowbreak{}2\\
Type: string\\
Default: nil
}

{\footnotesize <CM-08(03)\_\allowbreak{}ODP[05] ВИБРАНЕ ЗНАЧЕННЯ ПА206 РАМЕТРА(ів)> приймаються при виявленні несанкціонованого програмного забезпечення;}


{\footnotesize\bfseries
No: 12\\
Name: cm\_\allowbreak{}8\_\allowbreak{}3\_\allowbreak{}b\_\allowbreak{}3\\
Type: string\\
Default: nil
}

{\footnotesize <CM-08(03)\_\allowbreak{}ODP[05] ВИБРАНЕ ЗНАЧЕННЯ ПАРАМЕТРА(ів)> приймаються при виявленні несанкціонованих мікропрограмних компонентів;}




\subsubsection{ІНФОРМАЦІЯ ПРО ПІДЗВІТНІСТЬ (CM-8(4))}


\vspace{10pt}
{\footnotesize\bfseries
No: 1\\
Name: cm\_\allowbreak{}8\_\allowbreak{}4\_\allowbreak{}odp\\
Type: string\\
Default: nil
}

{\footnotesize вибрано одне або декілька з наступних ЗНАЧЕНЬ ПАРАМЕТРІВ: \{ім'я; позиція; роль\};}


{\footnotesize\bfseries
No: 2\\
Name: cm\_\allowbreak{}8\_\allowbreak{}4\_\allowbreak{}01\\
Type: string\\
Default: nil
}

{\footnotesize впровадженно в інвентаризацію компонент ів системи засіб для ідентифікації  <CM-08(04)\_\allowbreak{}ODP ВИБРАНЕ ЗНАЧЕННЯ ПАРАМЕТРА(ів)> осіб, відповідальних і підзвітних за управління цими компонентами.}




\subsubsection{ІНВЕНТАРИЗАЦІЯ КОМПОНЕНТІВ СИСТЕМИ ДУБЛЮВАННЯ КОМПОНЕНТІВ ОБЛІКУ (CM-8(5))}


\vspace{10pt}
\textit{Немає параметрів для цього контролю.}
\vspace{10pt}


\subsubsection{ПЕРЕВІРЕНІ НАЛАШТУВАННЯ ТА ЗАТВЕРДЖЕНІ ВІДХИЛЕННЯ (CM-8(6))}


\vspace{10pt}
{\footnotesize\bfseries
No: 1\\
Name: cm\_\allowbreak{}8\_\allowbreak{}6\_\allowbreak{}1\\
Type: string\\
Default: nil
}

{\footnotesize включено перевірені налаштування компонентів до поточних розгорнутих конфігурацій в інвентаризаційний облік компонентів системи;}


{\footnotesize\bfseries
No: 2\\
Name: cm\_\allowbreak{}8\_\allowbreak{}6\_\allowbreak{}2\\
Type: string\\
Default: nil
}

{\footnotesize включено будь-які затверджені відхилення до поточних розгорнутих конфігурацій в інвентаризаційний облік компонентів системи.}




\subsubsection{ЦЕНТРАЛІЗОВАНЕ СХОВИЩЕ (CM-8(7))}
Впровадити централізоване компонентів системи. сховище СИСТЕМИ для ЦЕНТРАЛІЗОВАНЕ інвентаризаційного обліку

\vspace{10pt}
{\footnotesize\bfseries
No: 1\\
Name: cm\_\allowbreak{}8\_\allowbreak{}7\_\allowbreak{}01\\
Type: string\\
Default: nil
}

{\footnotesize впроваджено централізоване сховище для інвентаризаційного обліку компонентів системи.}




\subsubsection{АВТОМАТИЗОВАНЕ ВІДСТЕЖЕННЯ МІСЦЯ РОЗТАШУВАННЯ (CM-8(8))}


\vspace{10pt}
{\footnotesize\bfseries
No: 1\\
Name: cm\_\allowbreak{}8\_\allowbreak{}8\_\allowbreak{}odp\\
Type: string\\
Default: nil
}

{\footnotesize визначено автоматизовані механізми відстеження компонентів;}


{\footnotesize\bfseries
No: 2\\
Name: cm\_\allowbreak{}8\_\allowbreak{}8\_\allowbreak{}01\\
Type: string\\
Default: nil
}

{\footnotesize використовуютья <CM-08(08)\_\allowbreak{}ODP автоматизовані механізми> для підтримки відстеження компонентів системи за географічним розташуванням}




\subsubsection{ПРИЗНАЧЕННЯ КОМПОНЕНТІВ СИСТЕМАМ (CM-8(9))}


\vspace{10pt}
{\footnotesize\bfseries
No: 1\\
Name: cm\_\allowbreak{}8\_\allowbreak{}9\_\allowbreak{}odp\\
Type: string\\
Default: nil
}

{\footnotesize визначено персонал або ролі, від яких слід отримувати підтвердження;}


{\footnotesize\bfseries
No: 2\\
Name: cm\_\allowbreak{}8\_\allowbreak{}9\_\allowbreak{}a\\
Type: string\\
Default: nil
}

{\footnotesize компоненти системи призначаються системі;}


{\footnotesize\bfseries
No: 3\\
Name: cm\_\allowbreak{}8\_\allowbreak{}9\_\allowbreak{}b\\
Type: string\\
Default: nil
}

{\footnotesize отримано підтвердження призначення компонента від <CM-08(09)\_\allowbreak{}ODP персонал або ролі>.}




\subsection{ПЛАН УПРАВЛІННЯ КОНФІГУРАЦІЄЮ (CM-9)}
Унікально ідентифікувати та автентифікувати користувачів або процеси, що діють від імені користувачів.

\vspace{10pt}
{\footnotesize\bfseries
No: 1\\
Name: cm\_\allowbreak{}9\_\allowbreak{}odp\\
Type: string\\
Default: nil
}

{\footnotesize визначено персонал або ролі для розгляду та затвердження плану управління конфігурацією;}


{\footnotesize\bfseries
No: 2\\
Name: cm\_\allowbreak{}9\_\allowbreak{}1\\
Type: string\\
Default: nil
}

{\footnotesize розроблено та задокументовано план управління конфігурацією системи;}


{\footnotesize\bfseries
No: 3\\
Name: cm\_\allowbreak{}9\_\allowbreak{}1\\
Type: string\\
Default: nil
}

{\footnotesize реалізовано план управління конфігурацією системи;}


{\footnotesize\bfseries
No: 4\\
Name: cm\_\allowbreak{}9\_\allowbreak{}a\_\allowbreak{}1\\
Type: string\\
Default: nil
}

{\footnotesize план управління конфігурацією описує ролі;}


{\footnotesize\bfseries
No: 5\\
Name: cm\_\allowbreak{}9\_\allowbreak{}a\_\allowbreak{}2\\
Type: string\\
Default: nil
}

{\footnotesize план управління конфігурацією описує відповідальність;}


{\footnotesize\bfseries
No: 6\\
Name: cm\_\allowbreak{}9\_\allowbreak{}a\_\allowbreak{}3\\
Type: string\\
Default: nil
}

{\footnotesize план управління конфігурацією описує процеси та процедури управління конфігурацією;}


{\footnotesize\bfseries
No: 7\\
Name: cm\_\allowbreak{}9\_\allowbreak{}b\_\allowbreak{}1\\
Type: string\\
Default: nil
}

{\footnotesize план управління конфігурацією встановлює процес ідентифікації елементів конфігурації протягом всього життєвого циклу розробки системи;}


{\footnotesize\bfseries
No: 8\\
Name: cm\_\allowbreak{}9\_\allowbreak{}b\_\allowbreak{}2\\
Type: string\\
Default: nil
}

{\footnotesize план управління конфігурацією встановлює процес управління конфігурацією елементів;}


{\footnotesize\bfseries
No: 9\\
Name: cm\_\allowbreak{}9\_\allowbreak{}c\_\allowbreak{}1\\
Type: string\\
Default: nil
}

{\footnotesize план управління конфігурацією визначає елементи конфігурації системи;}


{\footnotesize\bfseries
No: 10\\
Name: cm\_\allowbreak{}9\_\allowbreak{}c\_\allowbreak{}2\\
Type: string\\
Default: nil
}

{\footnotesize план управління конфігурацією розміщує елементи конфігурації під управлінням конфігурацією;}


{\footnotesize\bfseries
No: 11\\
Name: cm\_\allowbreak{}9\_\allowbreak{}d\\
Type: string\\
Default: nil
}

{\footnotesize план управління конфігурацією розглянуто та затверджено  <CM09\_\allowbreak{}ODP персонал або ролі>; 210}


{\footnotesize\bfseries
No: 12\\
Name: cm\_\allowbreak{}9\_\allowbreak{}e\_\allowbreak{}1\\
Type: string\\
Default: nil
}

{\footnotesize план управління конфігурацією захищений від несанкціонованого розкриття;}


{\footnotesize\bfseries
No: 13\\
Name: cm\_\allowbreak{}9\_\allowbreak{}e\_\allowbreak{}2\\
Type: string\\
Default: nil
}

{\footnotesize план управління конфігурацією захищений від несанкціонованої модифікації.}




\subsubsection{ВСТАНОВЛЕННЯ ВІДПОВІДАЛЬНОСТІ (CM-9(1))}


\vspace{10pt}
{\footnotesize\bfseries
No: 1\\
Name: cm\_\allowbreak{}9\_\allowbreak{}1\_\allowbreak{}01\\
Type: string\\
Default: nil
}

{\footnotesize встановлено відповідальність за реалізацію процесу управління конфігурацією персоналу, який безпосередньо не бере участь у розробці системи.}




\subsection{ОБМЕЖЕННЯ ВИКОРИСТАННЯ ПРОГРАМНОГО ЗАБЕЗПЕЧЕННЯ (CM-10)}


\vspace{10pt}
{\footnotesize\bfseries
No: 1\\
Name: cm\_\allowbreak{}10\_\allowbreak{}a\\
Type: string\\
Default: nil
}

{\footnotesize програмне забезпечення та супутня  документація використовуються відповідно до договірних угод та законів про авторські права;}


{\footnotesize\bfseries
No: 2\\
Name: cm\_\allowbreak{}10\_\allowbreak{}b\\
Type: string\\
Default: nil
}

{\footnotesize використання програмного забезпечення та пов'язаної з ним документації, захищеної ліцензіями, відстежується для контролю за копіюванням та розповсюдженням;}


{\footnotesize\bfseries
No: 3\\
Name: cm\_\allowbreak{}10\_\allowbreak{}c\\
Type: string\\
Default: nil
}

{\footnotesize використання технології однорангового обміну файлами контролюється та документується, щоб гарантувати, що одноранговий обмін файлами не використовується для несанкціонованого розповсюдження, відображення, виконання або відтворен ня програмного забезпечення, захищеного авторським правом.}




\subsubsection{ОБМЕЖЕННЯ ВИКОРИСТАННЯ ПРОГРАМНОГО ЗАБЕЗПЕЧЕННЯ ПРОГРАМНЕ ЗАБЕЗПЕЧЕННЯ З ВІДКРИТИМ ВИХІДНИМ КОДОМ (CM-10(1))}


\vspace{10pt}
{\footnotesize\bfseries
No: 1\\
Name: cm\_\allowbreak{}10\_\allowbreak{}1\_\allowbreak{}odp\\
Type: string\\
Default: nil
}

{\footnotesize визначено обмеження на використання програмного забезпечення з відкритим вихідним кодом}


{\footnotesize\bfseries
No: 2\\
Name: cm\_\allowbreak{}10\_\allowbreak{}1\_\allowbreak{}01\\
Type: string\\
Default: nil
}

{\footnotesize встановлено <CM-10(01)\_\allowbreak{}ODP обмеження> на використання програмного забезпечення з відкритим вихідним кодом.}




\subsection{ВСТАНОВЛЕНЕ КОРИСТУВАЧЕМ ПРОГРАМНЕ ЗАБЕЗПЕЧЕННЯ (CM-11)}


\vspace{10pt}
{\footnotesize\bfseries
No: 1\\
Name: cm\_\allowbreak{}11\_\allowbreak{}odp\_\allowbreak{}1\\
Type: string\\
Default: nil
}

{\footnotesize визначено правила (політики), що регулюють встановлення програмного забезпечення користувачами;}


{\footnotesize\bfseries
No: 2\\
Name: cm\_\allowbreak{}11\_\allowbreak{}odp\_\allowbreak{}2\\
Type: string\\
Default: nil
}

{\footnotesize визначено методи, що використовуються для забезпечення дотримання правил (політик) встановлення програмного забезпечення;}


{\footnotesize\bfseries
No: 3\\
Name: cm\_\allowbreak{}11\_\allowbreak{}odp\_\allowbreak{}3\\
Type: string\\
Default: nil
}

{\footnotesize визначено частоту, з якою слід контролювати відповідність правил (політик);}


{\footnotesize\bfseries
No: 4\\
Name: cm\_\allowbreak{}11\_\allowbreak{}a\\
Type: string\\
Default: nil
}

{\footnotesize встановлено <CM-11\_\allowbreak{}ODP[01] правила (політики)>, що регулюють встановлення програмного забезпечення користувачами;}


{\footnotesize\bfseries
No: 5\\
Name: cm\_\allowbreak{}11\_\allowbreak{}b\\
Type: string\\
Default: nil
}

{\footnotesize правила (політики) встановлення програмного забезпечення}


{\footnotesize\bfseries
No: 6\\
Name: cm\_\allowbreak{}11\_\allowbreak{}c\\
Type: string\\
Default: nil
}

{\footnotesize дотримання <CM-11\_\allowbreak{}ODP[01] політик> контролюється <CM11\_\allowbreak{}ODP[03] частота>.}




\subsubsection{ВСТАНОВЛЕНЕ КОРИСТУВАЧЕМ ПРОГРАМНЕ ЗАБЕЗПЕЧЕННЯ ПОПЕРЕДЖЕННЯ ПРО НЕСАНКЦІОНОВАНУ ІНСТАЛЯЦІЮ (CM-11(1))}


\vspace{10pt}
\textit{Немає параметрів для цього контролю.}
\vspace{10pt}


\subsubsection{ВСТАНОВЛЕНЕ КОРИСТУВАЧЕМ ПРОГРАМНЕ ЗАБЕЗПЕЧЕННЯ ВСТАНОВЛЕННЯ ПРОГРАМНОГО ЗАБЕЗПЕЧЕННЯ З ПРИВІЛЕЙОВАНИМ СТАТУСОМ (CM-11(2))}


\vspace{10pt}
{\footnotesize\bfseries
No: 1\\
Name: cm\_\allowbreak{}11\_\allowbreak{}2\_\allowbreak{}01\\
Type: string\\
Default: nil
}

{\footnotesize встановлювати програмне забезпечення дозволено користувачеві лише при наявності привілейованого статусу.}




\subsubsection{ВСТАНОВЛЕНЕ КОРИСТУВАЧЕМ ПРОГРАМНЕ ЗАБЕЗПЕЧЕННЯ АВТОМАТИЧНЕ ВИКОНАННЯ І МОНІТОРИНГ (CM-11(3))}


\vspace{10pt}
{\footnotesize\bfseries
No: 1\\
Name: cm\_\allowbreak{}11\_\allowbreak{}3\_\allowbreak{}odp\_\allowbreak{}1\\
Type: string\\
Default: nil
}

{\footnotesize визначено автоматизовані механізми для забезпечення дотримання політик виконання програмного забезпечення;}


{\footnotesize\bfseries
No: 2\\
Name: cm\_\allowbreak{}11\_\allowbreak{}3\_\allowbreak{}odp\_\allowbreak{}2\\
Type: string\\
Default: nil
}

{\footnotesize визначено автоматизовані механізми для забезпечення дотримання політик контролю програмного забезпечення;}


{\footnotesize\bfseries
No: 3\\
Name: cm\_\allowbreak{}11\_\allowbreak{}3\_\allowbreak{}1\\
Type: string\\
Default: nil
}

{\footnotesize дотримання політик виконання програмного забезпечення забезпечується за допомогою  <CM11(03)\_\allowbreak{}ODP[01] автоматизованих механізмів>; 214}


{\footnotesize\bfseries
No: 4\\
Name: cm\_\allowbreak{}11\_\allowbreak{}3\_\allowbreak{}2\\
Type: string\\
Default: nil
}

{\footnotesize дотримання політик контролю програмного забезпечення контролюється за допомогою  <CM-11(03)\_\allowbreak{}ODP[02] автоматизованих механізмів>.}




\subsection{РОЗТАШУВАННЯ ІНФОРМАЦІЇ (CM-12)}


\vspace{10pt}
{\footnotesize\bfseries
No: 1\\
Name: cm\_\allowbreak{}12\_\allowbreak{}odp\\
Type: string\\
Default: nil
}

{\footnotesize визначено інформацію, місцезнаходження якої має бути визначено та задокументовано;}


{\footnotesize\bfseries
No: 2\\
Name: cm\_\allowbreak{}12\_\allowbreak{}a\_\allowbreak{}1\\
Type: string\\
Default: nil
}

{\footnotesize місцезнаходження <CM-12\_\allowbreak{}ODP інформація> визначено та задокументовано;}


{\footnotesize\bfseries
No: 3\\
Name: cm\_\allowbreak{}12\_\allowbreak{}a\_\allowbreak{}2\\
Type: string\\
Default: nil
}

{\footnotesize визначено та задокументовано конкретні компоненти системи, на яких обробляється <CM-12\_\allowbreak{}ODP інформація>;}


{\footnotesize\bfseries
No: 4\\
Name: cm\_\allowbreak{}12\_\allowbreak{}a\_\allowbreak{}3\\
Type: string\\
Default: nil
}

{\footnotesize визначено та задокументовано конкретні компоненти системи, на яких зберігається <CM-12\_\allowbreak{}ODP інформація>;}


{\footnotesize\bfseries
No: 5\\
Name: cm\_\allowbreak{}12\_\allowbreak{}b\_\allowbreak{}1\\
Type: string\\
Default: nil
}

{\footnotesize визначено та задокументовано користувачів, які мають доступ до системи та компонентів системи, де обробляється <CM-12\_\allowbreak{}ODP інформація>;}


{\footnotesize\bfseries
No: 6\\
Name: cm\_\allowbreak{}12\_\allowbreak{}b\_\allowbreak{}2\\
Type: string\\
Default: nil
}

{\footnotesize визначено та задокументовано користувачів, які мають доступ до системи та компонентів системи, де зберігає ться <CM-12\_\allowbreak{}ODP інформація>; 215}


{\footnotesize\bfseries
No: 7\\
Name: cm\_\allowbreak{}12\_\allowbreak{}c\_\allowbreak{}1\\
Type: string\\
Default: nil
}

{\footnotesize задокументовано зміни розташування ( наприклад, системи або компонентів системи), де обробляється <CM-12\_\allowbreak{}ODP інформація>;}


{\footnotesize\bfseries
No: 8\\
Name: cm\_\allowbreak{}12\_\allowbreak{}c\_\allowbreak{}2\\
Type: string\\
Default: nil
}

{\footnotesize задокументовано зміни розташування ( наприклад, системи або компонентів сис теми), де зберігається <CM-12\_\allowbreak{}ODP інформація>;}




\subsubsection{АВТОМАТИЗОВАНІ ІНСТРУМЕНТИ ПІДТРИМКИ РОЗТАШУВАННЯ ІНФОРМАЦІЇ (CM-12(1))}


\vspace{10pt}
{\footnotesize\bfseries
No: 1\\
Name: cm\_\allowbreak{}12\_\allowbreak{}1\_\allowbreak{}odp\_\allowbreak{}1\\
Type: string\\
Default: nil
}

{\footnotesize інформація визначена за типом;}


{\footnotesize\bfseries
No: 2\\
Name: cm\_\allowbreak{}12\_\allowbreak{}1\_\allowbreak{}odp\_\allowbreak{}2\\
Type: string\\
Default: nil
}

{\footnotesize визначено компоненти системи, де знаходиться інформація;}


{\footnotesize\bfseries
No: 3\\
Name: cm\_\allowbreak{}12\_\allowbreak{}1\_\allowbreak{}01\\
Type: string\\
Default: nil
}

{\footnotesize автоматизовані інструменти використовуються для ідентифікації <CM-12(01)\_\allowbreak{}ODP[01] інформації за типом> забезпечити впровадження належних за ходів захисту щодо інформації про організацію і персональних даних.}




\subsection{ВІДОБРАЖЕННЯ ДІЙ ДАНИХ (CM-13)}


\vspace{10pt}
{\footnotesize\bfseries
No: 1\\
Name: cm\_\allowbreak{}13\_\allowbreak{}01\\
Type: string\\
Default: nil
}

{\footnotesize розроблено та задокументовано карту дій з даними системи.}




\subsection{ПІДПИСАНІ КОМПОНЕНТИ (CM-14)}


\vspace{10pt}
{\footnotesize\bfseries
No: 1\\
Name: cm\_\allowbreak{}14\_\allowbreak{}odp\_\allowbreak{}1\\
Type: string\\
Default: nil
}

{\footnotesize визначено програмне забезпечення, яке потребує перевірки сертифікату з цифровим підписом перед встановленням;}


{\footnotesize\bfseries
No: 2\\
Name: cm\_\allowbreak{}14\_\allowbreak{}odp\_\allowbreak{}2\\
Type: string\\
Default: nil
}

{\footnotesize визначено мікропрограмні компоненти, які потребує перевірки сертифікату з цифровим підписом перед встановленням;}


{\footnotesize\bfseries
No: 3\\
Name: cm\_\allowbreak{}14\_\allowbreak{}1\\
Type: string\\
Default: nil
}

{\footnotesize інсталяція < CM-14\_\allowbreak{}ODP[01] програмне забезпечення > попередньо запобігається, якщо не буде перевірено, що програмне забезпечення було  підписано цифровим підписом за допомогою сертифіката, визнаного та затвердженого організацією;}


{\footnotesize\bfseries
No: 4\\
Name: cm\_\allowbreak{}14\_\allowbreak{}2\\
Type: string\\
Default: nil
}

{\footnotesize інсталяція < CM-14\_\allowbreak{}ODP[02] мікропрограмні компоненти > попередньо запобігається, якщо не буде перевірено, що мікропрограмні компоненти були підписані  цифровим підписом за допомогою сертифіката, визнаного та затвердженого організацією;}




\section{CP}
Клас заходів захисту CP --- ПЛАНУВАННЯ БЕЗПЕРЕРВНОЇ РОБОТИ

Цей клас описує заходи для відновлення функціонування інформаційної системи та забезпечення її безперебійної роботи у разі надзвичайних ситуацій.

\textbf{Перелік заходів захисту:} Політика та процедури планування безперервної роботи (CP-1); План забезпечення безперервної роботи та відновлення функціонування (CP-2); Координація з пов'язаними планами (CP-2(1)); Навчання із забезпечення безперервної роботи (CP-3); Тестування плану забезпечення безперервної роботи та відновлення функціонування (CP-4); Альтернативна платформа тестування (CP-4(2)); Автоматичне тестування (CP-4(3)); Повне відновлення (CP-4(4)); Застосовуються для порушення та негативного впливу на (CP-4(5)); Оновлення плану забезпечення безперервної роботи та відновлення функціонування (CP-5) [Вилучено]; Альтернативне місце зберігання (CP-6); Відділення від первинного сховища (CP-6(1)); Час відновлення та встановлення цілей відновлення (CP-6(2)); Доступність (CP-6(3)); Альтернативний майданчик роботи (CP-7); Відділення від основного майданчика (CP-7(1)); Доступність (CP-7(2)); Пріоритет обслуговування (CP-7(3)); Підготовка для використання (CP-7(4)); Нездатність повернутися на основний майданчик (CP-7(6)); Комунікаційні послуги (CP-8); Пріоритет постачання послуг (CP-8(1)); Єдині точки відмови (CP-8(2)); Відділення основних та альтернативних провайдерів (CP-8(3)); План забезпечення безперервної роботи постачальника комунікаційних послуг (CP-8(4)); Тестування альтернативних комунікаційних послуг (CP-8(5)); Резервне копіювання (CP-9); Випробування на надійність та цілісність (CP-9(1)); Тестування відновлення з використанням зразків (CP-9(2)); Відокремлене сховище критичної інформації (CP-9(3)); Передача на альтернативне сховище зберігання (CP-9(5)); Надлишкова вторинна система (CP-9(6)); Подвійна авторизація (CP-9(7)); Криптографічний захист (CP-9(8)); Відновлення та відтворення системи (CP-10); Відновлення транзакцій (CP-10(2)); Відновлення в межах часового періоду (CP-10(4)); Здатність відмовостійкості (CP-10(5)) [Вилучено]; Альтернативні протоколи зв'язку (CP-11); Безпечний режим (CP-12); Альтернативні механізми безпеки (CP-13).

\subsection{ПОЛІТИКА ТА ПРОЦЕДУРИ ПЛАНУВАННЯ БЕЗПЕРЕРВНОЇ РОБОТИ (CP-1)}


\vspace{10pt}
{\footnotesize\bfseries
No: 1\\
Name: cp\_\allowbreak{}1\_\allowbreak{}odp\_\allowbreak{}1\\
Type: string\\
Default: nil
}

{\footnotesize визначено персонал або посади, на які поширюється політика планування безперервної роботи на випадок надзвичайних ситуацій;}


{\footnotesize\bfseries
No: 2\\
Name: cp\_\allowbreak{}1\_\allowbreak{}odp\_\allowbreak{}2\\
Type: string\\
Default: nil
}

{\footnotesize визначено персонал або посади, на які поширюється процедури планування безперервної роботи на випадок надзвичайних ситуацій;}


{\footnotesize\bfseries
No: 3\\
Name: cp\_\allowbreak{}1\_\allowbreak{}odp\_\allowbreak{}3\\
Type: string\\
Default: nil
}

{\footnotesize вибрано одне або декілька з наступних ЗНАЧЕНЬ ПАРА- МЕТРІВ: \{рівень організації; рівень місії/бізнес -процесу; рівень системи\};}


{\footnotesize\bfseries
No: 4\\
Name: cp\_\allowbreak{}1\_\allowbreak{}odp\_\allowbreak{}4\\
Type: string\\
Default: nil
}

{\footnotesize визначено посадову особу, яка кер уватиме політикою та процедурами планування безперервної роботи на випадок надзвичайних ситуацій;}


{\footnotesize\bfseries
No: 5\\
Name: cp\_\allowbreak{}1\_\allowbreak{}odp\_\allowbreak{}5\\
Type: string\\
Default: nil
}

{\footnotesize визначено частоту, з якою переглядається та оновлюється поточна політика;}


{\footnotesize\bfseries
No: 6\\
Name: cp\_\allowbreak{}1\_\allowbreak{}odp\_\allowbreak{}6\\
Type: string\\
Default: nil
}

{\footnotesize визначено події, після яких переглядається та оновлюється поточна політика;}


{\footnotesize\bfseries
No: 7\\
Name: cp\_\allowbreak{}1\_\allowbreak{}odp\_\allowbreak{}7\\
Type: string\\
Default: nil
}

{\footnotesize визначено частоту, з якою переглядаються та оновлюються поточні процедури;}


{\footnotesize\bfseries
No: 8\\
Name: cp\_\allowbreak{}1\_\allowbreak{}odp\_\allowbreak{}8\\
Type: string\\
Default: nil
}

{\footnotesize визначено події, після яких переглядаються та оновлюються поточні процедури;}


{\footnotesize\bfseries
No: 9\\
Name: cp\_\allowbreak{}1\_\allowbreak{}a\_\allowbreak{}1\\
Type: string\\
Default: nil
}

{\footnotesize розроблено та задокументовано політику планування безперервної роботи на випадок надзвичайних ситуацій;}


{\footnotesize\bfseries
No: 10\\
Name: cp\_\allowbreak{}1\_\allowbreak{}a\_\allowbreak{}2\\
Type: string\\
Default: nil
}

{\footnotesize політика планування безперервної роботи на випадок надзвичайних ситуацій поширюється на  <CP-01\_\allowbreak{}ODP[01] персонал або посади>;}


{\footnotesize\bfseries
No: 11\\
Name: cp\_\allowbreak{}1\_\allowbreak{}a\_\allowbreak{}3\\
Type: string\\
Default: nil
}

{\footnotesize розроблені та задокументовані процедури планування безперервної роботи на випадок надзвичайних ситуацій, що сприяють впровадженню політики планування безперервної роботи на випадок надзвичайних ситуацій та пов'язаних з нею заходів захисту на випадок надзвичайних ситуацій; 219}


{\footnotesize\bfseries
No: 12\\
Name: cp\_\allowbreak{}1\_\allowbreak{}a\_\allowbreak{}4\\
Type: string\\
Default: nil
}

{\footnotesize процедури планування безперервної роботи на випадок надзвичайних ситуацій поширюються на <CP-01\_\allowbreak{}ODP[02] персонал або посади>;}


{\footnotesize\bfseries
No: 13\\
Name: cp\_\allowbreak{}1\_\allowbreak{}a\_\allowbreak{}1\_\allowbreak{}a\_\allowbreak{}1\\
Type: string\\
Default: nil
}

{\footnotesize політика планування безперервної роботи на випадок надзвичайних ситуацій  <CP-01\_\allowbreak{}ODP[03] ВИБРАНЕ ЗНАЧЕННЯ ПАРАМЕТРА(ів)> містить мету;}


{\footnotesize\bfseries
No: 14\\
Name: cp\_\allowbreak{}1\_\allowbreak{}a\_\allowbreak{}1\_\allowbreak{}a\_\allowbreak{}2\\
Type: string\\
Default: nil
}

{\footnotesize політика планування безперервної роботи на випадок надзвичайних ситуацій  <CP-01\_\allowbreak{}ODP[03] ВИБРАНЕ ЗНАЧЕННЯ ПАРАМЕТРА(ів)> містить сферу застосування;}


{\footnotesize\bfseries
No: 15\\
Name: cp\_\allowbreak{}1\_\allowbreak{}a\_\allowbreak{}1\_\allowbreak{}a\_\allowbreak{}3\\
Type: string\\
Default: nil
}

{\footnotesize політика план ування безперервної роботи на випадок надзвичайних ситуацій  <CP-01\_\allowbreak{}ODP[03] ВИБРАНЕ ЗНАЧЕННЯ ПАРАМЕТРА(ів)> містить ролі;}


{\footnotesize\bfseries
No: 16\\
Name: cp\_\allowbreak{}1\_\allowbreak{}a\_\allowbreak{}1\_\allowbreak{}a\_\allowbreak{}4\\
Type: string\\
Default: nil
}

{\footnotesize політика планування безперервної роботи на випадок надзвичайних ситуацій <CP-01\_\allowbreak{}ODP[03] ВИБРАНЕ ЗНАЧЕННЯ ПАРАМЕТРА(ів)> містить обов'язки;}


{\footnotesize\bfseries
No: 17\\
Name: cp\_\allowbreak{}1\_\allowbreak{}a\_\allowbreak{}1\_\allowbreak{}a\_\allowbreak{}5\\
Type: string\\
Default: nil
}

{\footnotesize політика планування безперервної роботи на випадок надзвичайних ситуацій <CP-01\_\allowbreak{}ODP[03] ВИБРАНЕ ЗНАЧЕННЯ ПАРАМЕТРА(ів)> містить відповідальність керівництва;}


{\footnotesize\bfseries
No: 18\\
Name: cp\_\allowbreak{}1\_\allowbreak{}a\_\allowbreak{}1\_\allowbreak{}a\_\allowbreak{}6\\
Type: string\\
Default: nil
}

{\footnotesize політика планування безперервної роботи на випадок надзвичайних ситуацій <CP-01\_\allowbreak{}ODP[03] ВИБРАНЕ ЗНАЧЕННЯ ПАРАМЕТРА(ів)> містить координацію між підрозділами організації;}


{\footnotesize\bfseries
No: 19\\
Name: cp\_\allowbreak{}1\_\allowbreak{}a\_\allowbreak{}1\_\allowbreak{}a\_\allowbreak{}7\\
Type: string\\
Default: nil
}

{\footnotesize політика планування безперервної роботи на випадок надзвичайних ситуацій  <CP-01\_\allowbreak{}ODP[03] ВИБРАНЕ ЗНАЧЕННЯ ПАРАМЕТРА(ів)> містить систему контролю відповідності;}


{\footnotesize\bfseries
No: 20\\
Name: cp\_\allowbreak{}1\_\allowbreak{}a\_\allowbreak{}1\_\allowbreak{}b\\
Type: string\\
Default: nil
}

{\footnotesize <CP-01\_\allowbreak{}ODP[03] ВИБРАНЕ ЗНАЧЕННЯ ПАРАМЕТРА(ів)> політика планування безперервної роботи на випадок надзвичайних ситуацій в ідповідає чинним законам, виконавчим розпорядженням, директивам, положенням, політикам, стандартам і керівним принципам;}


{\footnotesize\bfseries
No: 21\\
Name: cp\_\allowbreak{}1\_\allowbreak{}b\\
Type: string\\
Default: nil
}

{\footnotesize <CP-01\_\allowbreak{}ODP[04] посадова особа> призначається для управління , документуванням та розповсюдженням політики та процедур планування безперервної роботи на випадок надзвичайних ситуацій;}


{\footnotesize\bfseries
No: 22\\
Name: cp\_\allowbreak{}1\_\allowbreak{}c\_\allowbreak{}1\_\allowbreak{}1\\
Type: string\\
Default: nil
}

{\footnotesize переглядається та оновлюється поточна політика планування}


{\footnotesize\bfseries
No: 23\\
Name: cp\_\allowbreak{}1\_\allowbreak{}c\_\allowbreak{}1\_\allowbreak{}2\\
Type: string\\
Default: nil
}

{\footnotesize переглядається т а оновлюється поточна політика планування безперервної роботи  на випадок надзвичайних ситуацій після 220}


{\footnotesize\bfseries
No: 24\\
Name: cp\_\allowbreak{}1\_\allowbreak{}c\_\allowbreak{}2\_\allowbreak{}1\\
Type: string\\
Default: nil
}

{\footnotesize переглядаються та оновлюються поточні процедури планування безперервної роботи  на випадок надзвичайних ситуацій}


{\footnotesize\bfseries
No: 25\\
Name: cp\_\allowbreak{}1\_\allowbreak{}c\_\allowbreak{}2\_\allowbreak{}2\\
Type: string\\
Default: nil
}

{\footnotesize переглядаються та оновлюються поточні процедури планування безперервної роботи  на випадок надзвичайних ситуацій}




\subsection{ПЛАН ЗАБЕЗПЕЧЕННЯ БЕЗПЕРЕРВНОЇ РОБОТИ ТА ВІДНОВЛЕННЯ ФУНКЦІОНУВАННЯ (CP-2)}


\vspace{10pt}
{\footnotesize\bfseries
No: 1\\
Name: cp\_\allowbreak{}2\_\allowbreak{}odp\_\allowbreak{}1\\
Type: string\\
Default: nil
}

{\footnotesize визначено персонал або ролі для перегляду плану забезпечення безперервної роботи у надзвичайних ситуаціях;}


{\footnotesize\bfseries
No: 2\\
Name: cp\_\allowbreak{}2\_\allowbreak{}odp\_\allowbreak{}2\\
Type: string\\
Default: nil
}

{\footnotesize визначено персонал або ролі для затвердження плану забезпечення безперервної роботи у надзвичайних ситуаціях;}


{\footnotesize\bfseries
No: 3\\
Name: cp\_\allowbreak{}2\_\allowbreak{}odp\_\allowbreak{}3\\
Type: string\\
Default: nil
}

{\footnotesize визначено ключовий резервний персонал (ідентифікований за іменами та/або за ролями), якому поширюються копії плану забезпечення безперервної роботи на випадок надзвичайних ситуацій;}


{\footnotesize\bfseries
No: 4\\
Name: cp\_\allowbreak{}2\_\allowbreak{}odp\_\allowbreak{}4\\
Type: string\\
Default: nil
}

{\footnotesize визначено ключові елементи, на які поширюються копії плану забезпечення безперервної роботи на випадок надзвичайних ситуацій;}


{\footnotesize\bfseries
No: 5\\
Name: cp\_\allowbreak{}2\_\allowbreak{}odp\_\allowbreak{}5\\
Type: string\\
Default: nil
}

{\footnotesize визначено періодичність перегляду плану забезпечення безперервної роботи  у надзвичайних ситуаціях;}


{\footnotesize\bfseries
No: 6\\
Name: cp\_\allowbreak{}2\_\allowbreak{}odp\_\allowbreak{}6\\
Type: string\\
Default: nil
}

{\footnotesize визначено ключовий резервний персонал (ідентифікований за іменами та/або ролями), якому необхідно повідомити про зміни;}


{\footnotesize\bfseries
No: 7\\
Name: cp\_\allowbreak{}2\_\allowbreak{}odp\_\allowbreak{}7\\
Type: string\\
Default: nil
}

{\footnotesize визначено ключові елементи організації, і які необхідно повідомити про зміни;}


{\footnotesize\bfseries
No: 8\\
Name: cp\_\allowbreak{}2\_\allowbreak{}a\_\allowbreak{}1\\
Type: string\\
Default: nil
}

{\footnotesize розроблено план забезпечення безперервної роботи у надзвичайних ситуаціях для системи, який визначає основні завдання, функції та пов'язані з ними вимоги щодо безперервної роботи; 221}


{\footnotesize\bfseries
No: 9\\
Name: cp\_\allowbreak{}2\_\allowbreak{}a\_\allowbreak{}2\_\allowbreak{}1\\
Type: string\\
Default: nil
}

{\footnotesize розроблено план забезпечення безперервної роботи у надзвичайних ситуаціях для системи, який забезпечує цілі;}


{\footnotesize\bfseries
No: 10\\
Name: cp\_\allowbreak{}2\_\allowbreak{}a\_\allowbreak{}2\_\allowbreak{}2\\
Type: string\\
Default: nil
}

{\footnotesize розроблено план забезпечення безперервної роботи у надзвичайних ситуаціях для системи, який забезпечує пріорітети;}


{\footnotesize\bfseries
No: 11\\
Name: cp\_\allowbreak{}2\_\allowbreak{}a\_\allowbreak{}2\_\allowbreak{}3\\
Type: string\\
Default: nil
}

{\footnotesize розроблено план забезпечення безперервної роботи у надзвичайних ситуаціях для системи, який забезпечує відповідні показники;}


{\footnotesize\bfseries
No: 12\\
Name: cp\_\allowbreak{}2\_\allowbreak{}a\_\allowbreak{}3\_\allowbreak{}1\\
Type: string\\
Default: nil
}

{\footnotesize розроблено план забезпечення безперервної роботи на випадок надзвичайних ситуацій для системи, в якому визначено ролі;}


{\footnotesize\bfseries
No: 13\\
Name: cp\_\allowbreak{}2\_\allowbreak{}a\_\allowbreak{}3\_\allowbreak{}2\\
Type: string\\
Default: nil
}

{\footnotesize розроблено план забезпечення безперервної роботи на випадок надзвичайних ситуацій для системи, в якому визначено обов'язки;}


{\footnotesize\bfseries
No: 14\\
Name: cp\_\allowbreak{}2\_\allowbreak{}a\_\allowbreak{}3\_\allowbreak{}3\\
Type: string\\
Default: nil
}

{\footnotesize розроблено план забезпечення безперервної роботи  на випадок надзвичайних ситуацій для системи, в якому визначено відповідальних осіб з контактною інформацією;}


{\footnotesize\bfseries
No: 15\\
Name: cp\_\allowbreak{}2\_\allowbreak{}a\_\allowbreak{}4\\
Type: string\\
Default: nil
}

{\footnotesize розроблено план забезпечення безперервної роботи на випадок непередбачених обставин для системи, який спрямований на підтримку основних завдань і функції, попри системні збої, компрометації або помилки;}


{\footnotesize\bfseries
No: 16\\
Name: cp\_\allowbreak{}2\_\allowbreak{}a\_\allowbreak{}5\\
Type: string\\
Default: nil
}

{\footnotesize розроблено план забезпечення безперервної роботи у надзвичайних ситуаціях для системи, який спрямований на повне відновлення функціонування системи без погіршення запланованих і реалізованих заходів захисту інформації та персональних даних;}


{\footnotesize\bfseries
No: 17\\
Name: cp\_\allowbreak{}2\_\allowbreak{}a\_\allowbreak{}6\\
Type: string\\
Default: nil
}

{\footnotesize розроблено план забезпечення безперервної роботи на випадок надзвичайних ситуацій для системи, який вирішує питання обміну інформацією про надзвичайні ситуації;}


{\footnotesize\bfseries
No: 18\\
Name: cp\_\allowbreak{}2\_\allowbreak{}a\_\allowbreak{}7\_\allowbreak{}1\\
Type: string\\
Default: nil
}

{\footnotesize розроблено план забезпечення безперервної роботи у надзвичайних ситуаціях для системи, який переглядається <CP-02\_\allowbreak{}ODP[01] персоналом або ролями>;}


{\footnotesize\bfseries
No: 19\\
Name: cp\_\allowbreak{}2\_\allowbreak{}a\_\allowbreak{}7\_\allowbreak{}2\\
Type: string\\
Default: nil
}

{\footnotesize розроблено план забезпечення безперервної роботи у надзвичайних ситуаціях для системи, який затверджено <CP-02\_\allowbreak{}ODP[02] персоналом або ролями>;}


{\footnotesize\bfseries
No: 20\\
Name: cp\_\allowbreak{}2\_\allowbreak{}b\_\allowbreak{}1\\
Type: string\\
Default: nil
}

{\footnotesize копії плану забезпечення безперервної роботи на випадок надзвичайних ситуацій розповсюджуються серед <CP-02\_\allowbreak{}ODP[03] персоналу>;}


{\footnotesize\bfseries
No: 21\\
Name: cp\_\allowbreak{}2\_\allowbreak{}b\_\allowbreak{}2\\
Type: string\\
Default: nil
}

{\footnotesize копії плану забезпечення безперервної роботи на випадок надзвичайних ситуацій розповсюджуються серед <CP-02\_\allowbreak{}ODP[04] елементів>;}


{\footnotesize\bfseries
No: 22\\
Name: cp\_\allowbreak{}2\_\allowbreak{}c\\
Type: string\\
Default: nil
}

{\footnotesize діяльність з планування безперервної роботи координується з діяльністю із заходами по усуненню інцидентів; 222}


{\footnotesize\bfseries
No: 23\\
Name: cp\_\allowbreak{}2\_\allowbreak{}d\\
Type: string\\
Default: nil
}

{\footnotesize переглядається план  забезпечення безперервної роботи у надзвичайних ситуаціях для системи <CP-02\_\allowbreak{}ODP[05] частота>;}


{\footnotesize\bfseries
No: 24\\
Name: cp\_\allowbreak{}2\_\allowbreak{}e\_\allowbreak{}1\\
Type: string\\
Default: nil
}

{\footnotesize план забезпечення безперервної роботи на випадок надзвичайних ситуацій оновлюється з урахуванням змін в організації, системі або середовищі функціонування;}


{\footnotesize\bfseries
No: 25\\
Name: cp\_\allowbreak{}2\_\allowbreak{}e\_\allowbreak{}2\\
Type: string\\
Default: nil
}

{\footnotesize план забезпечення безперервної роботи  на випадок надзвичайних ситуацій оновлюється для вирішення проблем, що виникають під час впровадження, виконання або тестування плану дій на випадок надзвичайних ситуацій;}


{\footnotesize\bfseries
No: 26\\
Name: cp\_\allowbreak{}2\_\allowbreak{}f\_\allowbreak{}1\\
Type: string\\
Default: nil
}

{\footnotesize зміни в плані забезпечення безперервної роботи на випадок}


{\footnotesize\bfseries
No: 27\\
Name: cp\_\allowbreak{}2\_\allowbreak{}f\_\allowbreak{}2\\
Type: string\\
Default: nil
}

{\footnotesize зміни в плані забезпечення безперервної роботи на випадок}


{\footnotesize\bfseries
No: 28\\
Name: cp\_\allowbreak{}2\_\allowbreak{}g\_\allowbreak{}1\\
Type: string\\
Default: nil
}

{\footnotesize уроки, отримані під час тестування планів забезпечення безперервної роботи  у надзвичайних ситуаціях або фактичних дій у надзвичайних ситуаціях, включаються в навчання;}


{\footnotesize\bfseries
No: 29\\
Name: cp\_\allowbreak{}2\_\allowbreak{}g\_\allowbreak{}2\\
Type: string\\
Default: nil
}

{\footnotesize уроки, отримані під час тестування планів забезпечення безперервної роботи  у надзвичайних ситуаціях або фактичних дій у надзвичайних ситуаціях, включаються в тестування;}


{\footnotesize\bfseries
No: 30\\
Name: cp\_\allowbreak{}2\_\allowbreak{}h\_\allowbreak{}1\\
Type: string\\
Default: nil
}

{\footnotesize план забезпечення безперервної роботи у надзвичайних ситуаціях захищений від несанкціонованого доступу;}


{\footnotesize\bfseries
No: 31\\
Name: cp\_\allowbreak{}2\_\allowbreak{}h\_\allowbreak{}2\\
Type: string\\
Default: nil
}

{\footnotesize план забезпечення безперервної роботи у надзвичайних ситуаціях захищений від несанкціонованих змін;}




\subsubsection{КООРДИНАЦІЯ З ПОВ'ЯЗАНИМИ ПЛАНАМИ (CP-2(1))}


\vspace{10pt}
{\footnotesize\bfseries
No: 1\\
Name: cp\_\allowbreak{}2\_\allowbreak{}1\_\allowbreak{}01\\
Type: string\\
Default: nil
}

{\footnotesize розробка плану забезпечення безперервної роботи у надзвичайних ситуаціях координується зі структурними підрозділами, які відповідають за розробку та реалізацію пов'язаних планів.}




\subsection{НАВЧАННЯ ІЗ ЗАБЕЗПЕЧЕННЯ БЕЗПЕРЕРВНОЇ РОБОТИ (CP-3)}


\vspace{10pt}
{\footnotesize\bfseries
No: 1\\
Name: cp\_\allowbreak{}3\_\allowbreak{}odp\_\allowbreak{}1\\
Type: string\\
Default: nil
}

{\footnotesize визначено період часу, протягом якого необхідно провести тренінг з підготовки до дій в умовах надзвичайних ситуацій після прийняття на себе ролі або відповідальності в умовах надзвичайних ситуацій;}


{\footnotesize\bfseries
No: 2\\
Name: cp\_\allowbreak{}3\_\allowbreak{}odp\_\allowbreak{}2\\
Type: string\\
Default: nil
}

{\footnotesize визначено частоту проведення тренінгів для користувачів системи, які виконують непередбачувану роль або несуть відповідальність;}


{\footnotesize\bfseries
No: 3\\
Name: cp\_\allowbreak{}3\_\allowbreak{}odp\_\allowbreak{}3\\
Type: string\\
Default: nil
}

{\footnotesize визначено частоту, з якою необхідно переглядати та оновлювати зміст тренувань на випадок надзвичайних ситуацій;}


{\footnotesize\bfseries
No: 4\\
Name: cp\_\allowbreak{}3\_\allowbreak{}odp\_\allowbreak{}4\\
Type: string\\
Default: nil
}

{\footnotesize визначено події, які потребують перегля ду та оновлення тренувань на випадок надзвичайних ситуацій;}


{\footnotesize\bfseries
No: 5\\
Name: cp\_\allowbreak{}3\_\allowbreak{}a\_\allowbreak{}1\\
Type: string\\
Default: nil
}

{\footnotesize підготовка на випадок надзвичайних ситуацій надається користувачам системи відповідно до призначених ролей та обов'язків протягом <CP-03\_\allowbreak{}ODP[01] періоду часу>  з моменту прийняття на себе надзвичайної ролі або обов'язку;}


{\footnotesize\bfseries
No: 6\\
Name: cp\_\allowbreak{}3\_\allowbreak{}a\_\allowbreak{}2\\
Type: string\\
Default: nil
}

{\footnotesize навчання на випадок надзвичайних ситуацій проводиться для користувачів системи відповідно до призначених ролей та обов'язків, якщо цього вимагають зміни в системі;}


{\footnotesize\bfseries
No: 7\\
Name: cp\_\allowbreak{}3\_\allowbreak{}a\_\allowbreak{}3\\
Type: string\\
Default: nil
}

{\footnotesize користувачам системи надається навчання на випадок надзвичайних ситуацій відповідно до призначених ролей та обов'язків}


{\footnotesize\bfseries
No: 8\\
Name: cp\_\allowbreak{}3\_\allowbreak{}b\_\allowbreak{}1\\
Type: string\\
Default: nil
}

{\footnotesize переглядається та оновлюється зміст тренувань за планом реагування на надзвичайні ситуації <CP-03\_\allowbreak{}ODP[03] частота>;}


{\footnotesize\bfseries
No: 9\\
Name: cp\_\allowbreak{}3\_\allowbreak{}b\_\allowbreak{}2\\
Type: string\\
Default: nil
}

{\footnotesize зміст тренувань за планом реагування на надзвичайні ситуації}




\subsection{ТЕСТУВАННЯ ПЛАНУ ЗАБЕЗПЕЧЕННЯ БЕЗПЕРЕРВНОЇ РОБОТИ ТА ВІДНОВЛЕННЯ ФУНКЦІОНУВАННЯ (CP-4)}


\vspace{10pt}
{\footnotesize\bfseries
No: 1\\
Name: cp\_\allowbreak{}4\_\allowbreak{}odp\_\allowbreak{}1\\
Type: string\\
Default: nil
}

{\footnotesize визначено частоту тестування плану забезпечення безперервної роботи у надзвичайних ситуаціях для системи;}


{\footnotesize\bfseries
No: 2\\
Name: cp\_\allowbreak{}4\_\allowbreak{}odp\_\allowbreak{}2\\
Type: string\\
Default: nil
}

{\footnotesize визначено тести для визначення ефективності плану забезпечення безперервної роботи у надзвичайних ситуаціях;}


{\footnotesize\bfseries
No: 3\\
Name: cp\_\allowbreak{}4\_\allowbreak{}odp\_\allowbreak{}3\\
Type: string\\
Default: nil
}

{\footnotesize визначені тести для визначення готовності до виконання плану забезпечення безперервної роботи у надзвичайних ситуаціях;}


{\footnotesize\bfseries
No: 4\\
Name: cp\_\allowbreak{}4\_\allowbreak{}a\_\allowbreak{}1\\
Type: string\\
Default: nil
}

{\footnotesize тестується план забезпечення безперервної робо ти у надзвичайних ситуаціях для системи <CP-04\_\allowbreak{}ODP[01] частота>;}


{\footnotesize\bfseries
No: 5\\
Name: cp\_\allowbreak{}4\_\allowbreak{}a\_\allowbreak{}2\\
Type: string\\
Default: nil
}

{\footnotesize <CP-04\_\allowbreak{}ODP[02] тести>  використовуються для визначення ефективності плану;}


{\footnotesize\bfseries
No: 6\\
Name: cp\_\allowbreak{}4\_\allowbreak{}a\_\allowbreak{}3\\
Type: string\\
Default: nil
}

{\footnotesize <CP-04\_\allowbreak{}ODP[03] тести>  використовуються для визначення готовності до виконання плану;}


{\footnotesize\bfseries
No: 7\\
Name: cp\_\allowbreak{}4\_\allowbreak{}b\\
Type: string\\
Default: nil
}

{\footnotesize переглядаються результати тестування плану забезпечення безперервної роботи у надзвичайних ситуаціях;}


{\footnotesize\bfseries
No: 8\\
Name: cp\_\allowbreak{}4\_\allowbreak{}c\\
Type: string\\
Default: nil
}

{\footnotesize за необхідності ініціюються коригувальні дії.}




\subsubsection{АЛЬТЕРНАТИВНА ПЛАТФОРМА ТЕСТУВАННЯ (CP-4(2))}


\vspace{10pt}
{\footnotesize\bfseries
No: 1\\
Name: cp\_\allowbreak{}4\_\allowbreak{}2\_\allowbreak{}a\\
Type: string\\
Default: nil
}

{\footnotesize план забезпечення безперервної роботи у надзвичайних ситуаціях тестується на альтернативній платформі для ознайомлення персоналу з об'єктом та наявними ресурсами;}


{\footnotesize\bfseries
No: 2\\
Name: cp\_\allowbreak{}4\_\allowbreak{}2\_\allowbreak{}b\\
Type: string\\
Default: nil
}

{\footnotesize план забезпечення безперервної роботи у надзвичайних ситуаціях тестується на альтернативній платформі для оцінки можливостей альтернативної платформи;}




\subsubsection{АВТОМАТИЧНЕ ТЕСТУВАННЯ (CP-4(3))}


\vspace{10pt}
{\footnotesize\bfseries
No: 1\\
Name: cp\_\allowbreak{}4\_\allowbreak{}3\_\allowbreak{}odp\\
Type: string\\
Default: nil
}

{\footnotesize визначено автоматизовані механізми тестування планів забезпечення безперервної роботи у надзвичайних ситуаціях;}


{\footnotesize\bfseries
No: 2\\
Name: cp\_\allowbreak{}4\_\allowbreak{}3\_\allowbreak{}01\\
Type: string\\
Default: nil
}

{\footnotesize план забезпечення безперервної роботи  на випадок надзвичайних ситуацій тестується за допомогою <CP-04(03)\_\allowbreak{}ODP автоматизованих механізмів>.}




\subsubsection{ПОВНЕ ВІДНОВЛЕННЯ (CP-4(4))}


\vspace{10pt}
{\footnotesize\bfseries
No: 1\\
Name: cp\_\allowbreak{}4\_\allowbreak{}4\_\allowbreak{}1\\
Type: string\\
Default: nil
}

{\footnotesize включено повне відновлення системи до відомого стану як частину тестування плану забезпечення безперервної роботи та відновлення функціонування}


{\footnotesize\bfseries
No: 2\\
Name: cp\_\allowbreak{}4\_\allowbreak{}4\_\allowbreak{}2\\
Type: string\\
Default: nil
}

{\footnotesize включено повне повернення системи до відомого стану як частину тестування плану забезпечення безперервної роботи та відновлення функціонування}




\subsubsection{ЗАСТОСОВУЮТЬСЯ ДЛЯ ПОРУШЕННЯ ТА НЕГАТИВНОГО ВПЛИВУ НА (CP-4(5))}


\vspace{10pt}
{\footnotesize\bfseries
No: 1\\
Name: cp\_\allowbreak{}4\_\allowbreak{}5\_\allowbreak{}odp\_\allowbreak{}1\\
Type: string\\
Default: nil
}

{\footnotesize визначено механізми, що застосовуються для порушення та негативного впливу на систему або на компонент системи;}


{\footnotesize\bfseries
No: 2\\
Name: cp\_\allowbreak{}4\_\allowbreak{}5\_\allowbreak{}odp\_\allowbreak{}2\\
Type: string\\
Default: nil
}

{\footnotesize визначено систему або компонент системи, до яких застосовуються механізми порушення та негативного впливу}


{\footnotesize\bfseries
No: 3\\
Name: cp\_\allowbreak{}4\_\allowbreak{}5\_\allowbreak{}01\\
Type: string\\
Default: nil
}

{\footnotesize <CP-04(05)\_\allowbreak{}ODP[01] механізми> застосовуються для порушення та негативного впливу на  <CP-04(05)\_\allowbreak{}ODP[02] систему або компонент системи>.}




\subsection{ОНОВЛЕННЯ ПЛАНУ ЗАБЕЗПЕЧЕННЯ БЕЗПЕРЕРВНОЇ РОБОТИ ТА ВІДНОВЛЕННЯ ФУНКЦІОНУВАННЯ (CP-5) [Вилучено]}
[Вилучено: Включено до СР-02]

\vspace{10pt}
\textit{Немає параметрів для цього контролю.}
\vspace{10pt}


\subsection{АЛЬТЕРНАТИВНЕ МІСЦЕ ЗБЕРІГАННЯ (CP-6)}


\vspace{10pt}
{\footnotesize\bfseries
No: 1\\
Name: cp\_\allowbreak{}6\_\allowbreak{}a\_\allowbreak{}1\\
Type: string\\
Default: nil
}

{\footnotesize створено альтернативне місце зберігання;}


{\footnotesize\bfseries
No: 2\\
Name: cp\_\allowbreak{}6\_\allowbreak{}a\_\allowbreak{}2\\
Type: string\\
Default: nil
}

{\footnotesize створення альтернативного місця зберігання включає в себе необхідні угоди, що дозволяють зберігати та видавати інформацію резервного копіювання системи;}


{\footnotesize\bfseries
No: 3\\
Name: cp\_\allowbreak{}6\_\allowbreak{}b\\
Type: string\\
Default: nil
}

{\footnotesize в альтернативному місці зберігання впроваджені заходи захисту, аналогічні заходам захисту основної локації.}




\subsubsection{ВІДДІЛЕННЯ ВІД ПЕРВИННОГО СХОВИЩА (CP-6(1))}


\vspace{10pt}
{\footnotesize\bfseries
No: 1\\
Name: cp\_\allowbreak{}6\_\allowbreak{}1\_\allowbreak{}01\\
Type: string\\
Default: nil
}

{\footnotesize визначено альтернативне місце зберігання, яке відокремлено від основного місця зберігання, щоб зменшити сприйнятливість до тих самих загроз.}




\subsubsection{ЧАС ВІДНОВЛЕННЯ ТА ВСТАНОВЛЕННЯ ЦІЛЕЙ ВІДНОВЛЕННЯ (CP-6(2))}


\vspace{10pt}
{\footnotesize\bfseries
No: 1\\
Name: cp\_\allowbreak{}6\_\allowbreak{}2\_\allowbreak{}1\\
Type: string\\
Default: nil
}

{\footnotesize налаштувати альтернативне місце зберігання для полегшення операцій відновлення відповідно до часу відновлення;}


{\footnotesize\bfseries
No: 2\\
Name: cp\_\allowbreak{}6\_\allowbreak{}2\_\allowbreak{}2\\
Type: string\\
Default: nil
}

{\footnotesize налаштувати альтернативне місце зберігання для полегшення операцій відновлення відповідно до встановлених цілей відновлення.}




\subsubsection{ДОСТУПНІСТЬ (CP-6(3))}


\vspace{10pt}
{\footnotesize\bfseries
No: 1\\
Name: cp\_\allowbreak{}6\_\allowbreak{}3\_\allowbreak{}1\\
Type: string\\
Default: nil
}

{\footnotesize визначено потенційні проблеми доступності для альтернативного місця зберігання в разі збоїв або  стихійних лих по всьому регіоні;}


{\footnotesize\bfseries
No: 2\\
Name: cp\_\allowbreak{}6\_\allowbreak{}3\_\allowbreak{}2\\
Type: string\\
Default: nil
}

{\footnotesize в загальних рисах окреслено дії щодо пом'якшення наслідків}




\subsection{АЛЬТЕРНАТИВНИЙ МАЙДАНЧИК РОБОТИ (CP-7)}


\vspace{10pt}
{\footnotesize\bfseries
No: 1\\
Name: cp\_\allowbreak{}7\_\allowbreak{}odp\_\allowbreak{}1\\
Type: string\\
Default: nil
}

{\footnotesize визначені операції системи для основних завдань і функцій;}


{\footnotesize\bfseries
No: 2\\
Name: cp\_\allowbreak{}7\_\allowbreak{}odp\_\allowbreak{}2\\
Type: string\\
Default: nil
}

{\footnotesize визначено період часу, відповідно термінам відновлення та встановленим цілям відновлення;}


{\footnotesize\bfseries
No: 3\\
Name: cp\_\allowbreak{}7\_\allowbreak{}a\\
Type: string\\
Default: nil
}

{\footnotesize альтернативний майданчик для роботи, включно з необхідними}


{\footnotesize\bfseries
No: 4\\
Name: cp\_\allowbreak{}7\_\allowbreak{}b\_\allowbreak{}1\\
Type: string\\
Default: nil
}

{\footnotesize обладнання та прилади, необхідні для передачі, доступні на альтернативному місці роботи або якщо укладені контракти на періоду часу> для передачі;}


{\footnotesize\bfseries
No: 5\\
Name: cp\_\allowbreak{}7\_\allowbreak{}b\_\allowbreak{}2\\
Type: string\\
Default: nil
}

{\footnotesize обладнання та прилади, необхідні для відновлення, доступні на альтернативному місці роботи або якщо укладені контракти на періоду часу> для передачі;}


{\footnotesize\bfseries
No: 6\\
Name: cp\_\allowbreak{}7\_\allowbreak{}c\\
Type: string\\
Default: nil
}

{\footnotesize впроваджено на альтернативному майданчику роботи заходи захисту, еквівалентні тим, що впровадженні на основному майданчику.}




\subsubsection{ВІДДІЛЕННЯ ВІД ОСНОВНОГО МАЙДАНЧИКА (CP-7(1))}


\vspace{10pt}
{\footnotesize\bfseries
No: 1\\
Name: cp\_\allowbreak{}7\_\allowbreak{}1\_\allowbreak{}01\\
Type: string\\
Default: nil
}

{\footnotesize визначено альтернативний майданчик для роботи, який відокремлений від основного майданчика, з метою зменшення сприйнятливості до тих самих загроз}




\subsubsection{ДОСТУПНІСТЬ (CP-7(2))}


\vspace{10pt}
{\footnotesize\bfseries
No: 1\\
Name: cp\_\allowbreak{}7\_\allowbreak{}2\_\allowbreak{}1\\
Type: string\\
Default: nil
}

{\footnotesize визначено потенційні проблеми доступності для альтернативного майданчика для роботи в разі збоїв або катастрофи по всьому регіону}


{\footnotesize\bfseries
No: 2\\
Name: cp\_\allowbreak{}7\_\allowbreak{}2\_\allowbreak{}2\\
Type: string\\
Default: nil
}

{\footnotesize окреслено чіткі заходи щодо пом'якшення наслідків}




\subsubsection{ПРІОРИТЕТ ОБСЛУГОВУВАННЯ (CP-7(3))}


\vspace{10pt}
{\footnotesize\bfseries
No: 1\\
Name: cp\_\allowbreak{}7\_\allowbreak{}3\_\allowbreak{}01\\
Type: string\\
Default: nil
}

{\footnotesize розроблено угоди про альтернативний майданчик для роботи, які містять положення щодо пріоритету обслуговування відповідно до вимог стосовно доступності (включно з вимогами щодо часу відно237 влення).}




\subsubsection{ПІДГОТОВКА ДЛЯ ВИКОРИСТАННЯ (CP-7(4))}


\vspace{10pt}
{\footnotesize\bfseries
No: 1\\
Name: cp\_\allowbreak{}7\_\allowbreak{}4\_\allowbreak{}01\\
Type: string\\
Default: nil
}

{\footnotesize підготовлено альтернативний майданчик для роботи таким чином, щоб майданчик був готовий до використання як оперативний майданчик, що підтримує виконання основних завдань та функцій}




\subsubsection{НЕЗДАТНІСТЬ ПОВЕРНУТИСЯ НА ОСНОВНИЙ МАЙДАНЧИК (CP-7(6))}


\vspace{10pt}
{\footnotesize\bfseries
No: 1\\
Name: cp\_\allowbreak{}7\_\allowbreak{}6\_\allowbreak{}1\\
Type: string\\
Default: nil
}

{\footnotesize розроблено план до обставин, які виключають повернення на основне місце роботи;}


{\footnotesize\bfseries
No: 2\\
Name: cp\_\allowbreak{}7\_\allowbreak{}6\_\allowbreak{}2\\
Type: string\\
Default: nil
}

{\footnotesize підготувалися до обставин, які виключають повернення на основне місце роботи;}




\subsection{КОМУНІКАЦІЙНІ ПОСЛУГИ (CP-8)}


\vspace{10pt}
{\footnotesize\bfseries
No: 1\\
Name: cp\_\allowbreak{}8\_\allowbreak{}odp\_\allowbreak{}1\\
Type: string\\
Default: nil
}

{\footnotesize визначено операції системи, які необхідно відновити для виконання основних завдань та функцій;}


{\footnotesize\bfseries
No: 2\\
Name: cp\_\allowbreak{}8\_\allowbreak{}odp\_\allowbreak{}2\\
Type: string\\
Default: nil
}

{\footnotesize визначено період часу, протягом якого необхідно відновити основні завдання та функції, коли основні комунікаційні можливості недоступні;}


{\footnotesize\bfseries
No: 3\\
Name: cp\_\allowbreak{}8\_\allowbreak{}01\\
Type: string\\
Default: nil
}

{\footnotesize альтернативні комунікаційні послуги, включно з необхідними угодами, що дозволяють відновити <CP-08\_\allowbreak{}ODP[01] операції системи>, створюються для основних завдань та функцій протягом <CP08\_\allowbreak{}ODP[02] періоду часу>, коли основні комунікаційні можливості недоступн і на основному або альтернативному місцях роботи або зберігання.}




\subsubsection{ПРІОРИТЕТ ПОСТАЧАННЯ ПОСЛУГ (CP-8(1))}


\vspace{10pt}
{\footnotesize\bfseries
No: 1\\
Name: cp\_\allowbreak{}8\_\allowbreak{}1\_\allowbreak{}a\_\allowbreak{}1\\
Type: string\\
Default: nil
}

{\footnotesize розроблено угоди про надання основних комунікаційних послуг, які містять пріоритетні положення про надання послуг відповідно до вимог щодо доступності (включно з вимогами щодо часу відновлення);}


{\footnotesize\bfseries
No: 2\\
Name: cp\_\allowbreak{}8\_\allowbreak{}1\_\allowbreak{}a\_\allowbreak{}2\\
Type: string\\
Default: nil
}

{\footnotesize розроблено альтернативні угоди про надання комунікаційних послуг, які містять положення про пріоритетність надання послуг відповідно до вимог доступності (включно з вимогами щодо часу відновлення);}


{\footnotesize\bfseries
No: 3\\
Name: cp\_\allowbreak{}8\_\allowbreak{}1\_\allowbreak{}b\\
Type: string\\
Default: nil
}

{\footnotesize надсилається запит про пріоритети комуніка ційних послуг для всіх комунікаційних послуг, що використовуються для забезпечення безперервності роботи, якщо основні та/або альтернативні комунікаційні послуги надаються загальним оператором.}




\subsubsection{ЄДИНІ ТОЧКИ ВІДМОВИ (CP-8(2))}


\vspace{10pt}
{\footnotesize\bfseries
No: 1\\
Name: cp\_\allowbreak{}8\_\allowbreak{}2\_\allowbreak{}01\\
Type: string\\
Default: nil
}

{\footnotesize отримано альтернативні комунікаційні послуги з метою зменшення ймовірності спільного використання єдиної точки відмови з основними комунікаційними послугами.}




\subsubsection{ВІДДІЛЕННЯ ОСНОВНИХ ТА АЛЬТЕРНАТИВНИХ ПРОВАЙДЕРІВ (CP-8(3))}


\vspace{10pt}
{\footnotesize\bfseries
No: 1\\
Name: cp\_\allowbreak{}8\_\allowbreak{}3\_\allowbreak{}01\\
Type: string\\
Default: nil
}

{\footnotesize отримуються альтернативні комунікаційні послуги від постачальників, які відокремлені від основних постачальників послуг, щоб зменшити сприйнятливості до тих самих загроз.}




\subsubsection{ПЛАН ЗАБЕЗПЕЧЕННЯ БЕЗПЕРЕРВНОЇ РОБОТИ ПОСТАЧАЛЬНИКА КОМУНІКАЦІЙНИХ ПОСЛУГ (CP-8(4))}


\vspace{10pt}
{\footnotesize\bfseries
No: 1\\
Name: cp\_\allowbreak{}8\_\allowbreak{}4\_\allowbreak{}odp\_\allowbreak{}1\\
Type: string\\
Default: nil
}

{\footnotesize визначено частоту, з якою постачальники послуг повинні надавати свідчення про тестування планів забезпечення безперервної роботи;}


{\footnotesize\bfseries
No: 2\\
Name: cp\_\allowbreak{}8\_\allowbreak{}4\_\allowbreak{}odp\_\allowbreak{}2\\
Type: string\\
Default: nil
}

{\footnotesize визначено частоту, з якою постачальники послуг повинні надавати свідчення про тренування з планів забезпечення безперервної роботи;}


{\footnotesize\bfseries
No: 3\\
Name: cp\_\allowbreak{}8\_\allowbreak{}4\_\allowbreak{}a\_\allowbreak{}1\\
Type: string\\
Default: nil
}

{\footnotesize постачальники основних комунікаційних послуг зобов'язані мати плани забезпечення безперервної роботи;}


{\footnotesize\bfseries
No: 4\\
Name: cp\_\allowbreak{}8\_\allowbreak{}4\_\allowbreak{}a\_\allowbreak{}2\\
Type: string\\
Default: nil
}

{\footnotesize постачальники альтернативних комунікаційних послуг зобов'язані мати плани забезпечення безперервної роботи;}


{\footnotesize\bfseries
No: 5\\
Name: cp\_\allowbreak{}8\_\allowbreak{}4\_\allowbreak{}b\\
Type: string\\
Default: nil
}

{\footnotesize переглядаються плани забезпечення безперервної роботи постачальників комунікаційних послуг для забезпечення 241 відповідності планам забе зпечення безперервної роботи організації;}


{\footnotesize\bfseries
No: 6\\
Name: cp\_\allowbreak{}8\_\allowbreak{}4\_\allowbreak{}c\_\allowbreak{}1\\
Type: string\\
Default: nil
}

{\footnotesize отримано свідчення про тестування планів забезпечення}


{\footnotesize\bfseries
No: 7\\
Name: cp\_\allowbreak{}8\_\allowbreak{}4\_\allowbreak{}c\_\allowbreak{}2\\
Type: string\\
Default: nil
}

{\footnotesize отримано свідчення про тренування з планів забезпечення}




\subsubsection{ТЕСТУВАННЯ АЛЬТЕРНАТИВНИХ КОМУНІКАЦІЙНИХ ПОСЛУГ (CP-8(5))}


\vspace{10pt}
{\footnotesize\bfseries
No: 1\\
Name: cp\_\allowbreak{}8\_\allowbreak{}5\_\allowbreak{}odp\\
Type: string\\
Default: nil
}

{\footnotesize визначено частоту з якою необхідно тестувати надання альтернативних комунікаційних послуг;}


{\footnotesize\bfseries
No: 2\\
Name: cp\_\allowbreak{}8\_\allowbreak{}5\_\allowbreak{}01\\
Type: string\\
Default: nil
}

{\footnotesize тестування надання альтернативних комунікаційних послуг з <CP-08(05)\_\allowbreak{}ODP частотою>.}




\subsection{РЕЗЕРВНЕ КОПІЮВАННЯ (CP-9)}


\vspace{10pt}
{\footnotesize\bfseries
No: 1\\
Name: cp\_\allowbreak{}9\_\allowbreak{}odp\_\allowbreak{}1\\
Type: string\\
Default: nil
}

{\footnotesize визначено компоненти системи, для яких необхідно проводити резервне копіювання інформації користувачів;}


{\footnotesize\bfseries
No: 2\\
Name: cp\_\allowbreak{}9\_\allowbreak{}odp\_\allowbreak{}2\\
Type: string\\
Default: nil
}

{\footnotesize визначено частоту, з якою слід проводити резервне копіювання інформації користувача відповідно до часу відновлення та цілей відновлення;}


{\footnotesize\bfseries
No: 3\\
Name: cp\_\allowbreak{}9\_\allowbreak{}odp\_\allowbreak{}3\\
Type: string\\
Default: nil
}

{\footnotesize визначено частоту проведення резервного копіювання інформації системи, що відповідає завдань відновлення і встановлених цілей відновлення;}


{\footnotesize\bfseries
No: 4\\
Name: cp\_\allowbreak{}9\_\allowbreak{}odp\_\allowbreak{}4\\
Type: string\\
Default: nil
}

{\footnotesize визначено частоту, з якою слід проводити резервне копіювання документації системи відповідно до часу відновлення та цілей точки відновлення;}


{\footnotesize\bfseries
No: 5\\
Name: cp\_\allowbreak{}9\_\allowbreak{}a\\
Type: string\\
Default: nil
}

{\footnotesize резервне копіювання інформації кор истувача, що міститься в}


{\footnotesize\bfseries
No: 6\\
Name: cp\_\allowbreak{}9\_\allowbreak{}b\\
Type: string\\
Default: nil
}

{\footnotesize виконується резервне копіювання інформації системи, що CP-09(с) створюються резервні копії документації системи, включаючи}


{\footnotesize\bfseries
No: 7\\
Name: cp\_\allowbreak{}9\_\allowbreak{}d\_\allowbreak{}1\\
Type: string\\
Default: nil
}

{\footnotesize конфіденційність резервних копій інформації захищена;}


{\footnotesize\bfseries
No: 8\\
Name: cp\_\allowbreak{}9\_\allowbreak{}d\_\allowbreak{}2\\
Type: string\\
Default: nil
}

{\footnotesize цілісність резервних копій інформації захищена;}


{\footnotesize\bfseries
No: 9\\
Name: cp\_\allowbreak{}9\_\allowbreak{}d\_\allowbreak{}3\\
Type: string\\
Default: nil
}

{\footnotesize доступність резервних копій інформації захищена.}




\subsubsection{ВИПРОБУВАННЯ НА НАДІЙНІСТЬ ТА ЦІЛІСНІСТЬ (CP-9(1))}


\vspace{10pt}
{\footnotesize\bfseries
No: 1\\
Name: cp\_\allowbreak{}9\_\allowbreak{}1\_\allowbreak{}odp\_\allowbreak{}1\\
Type: string\\
Default: nil
}

{\footnotesize визначено частоту тестування на надійність носіїв резервних копій інформації;}


{\footnotesize\bfseries
No: 2\\
Name: cp\_\allowbreak{}9\_\allowbreak{}1\_\allowbreak{}odp\_\allowbreak{}2\\
Type: string\\
Default: nil
}

{\footnotesize визначено частоту тестування на цілісність носіїв резервних копій інформації;}


{\footnotesize\bfseries
No: 3\\
Name: cp\_\allowbreak{}9\_\allowbreak{}1\_\allowbreak{}1\\
Type: string\\
Default: nil
}

{\footnotesize носії резервних копій інформації тестується <CP09(01)\_\allowbreak{}ODP[01] частота> для перевірки надійності;}


{\footnotesize\bfseries
No: 4\\
Name: cp\_\allowbreak{}9\_\allowbreak{}1\_\allowbreak{}2\\
Type: string\\
Default: nil
}

{\footnotesize носії резервних копій інформації тестується <CP09(01)\_\allowbreak{}ODP[02] частота> для перевірки цілісності;}




\subsubsection{ТЕСТУВАННЯ ВІДНОВЛЕННЯ З ВИКОРИСТАННЯМ ЗРАЗКІВ (CP-9(2))}


\vspace{10pt}
{\footnotesize\bfseries
No: 1\\
Name: cp\_\allowbreak{}9\_\allowbreak{}2\_\allowbreak{}01\\
Type: string\\
Default: nil
}

{\footnotesize використовується зразок резервної копії інформації при відновленні вибраних функцій системи як частину тестування плану забезпечення безперервної роботи та відновлення функціонування}




\subsubsection{ВІДОКРЕМЛЕНЕ СХОВИЩЕ КРИТИЧНОЇ ІНФОРМАЦІЇ (CP-9(3))}


\vspace{10pt}
{\footnotesize\bfseries
No: 1\\
Name: cp\_\allowbreak{}9\_\allowbreak{}3\_\allowbreak{}odp\\
Type: string\\
Default: nil
}

{\footnotesize визначено критичного системного програмного забезпечення та іншої інформації, пов'язаної з безпекою яке має зберігатися в окремому сховищі або у вогнетривкому контейнері;}


{\footnotesize\bfseries
No: 2\\
Name: cp\_\allowbreak{}9\_\allowbreak{}3\_\allowbreak{}01\\
Type: string\\
Default: nil
}

{\footnotesize резервні копії <CP-09(03)\_\allowbreak{}ODP критичного системного програмного забезпечення та іншої інформації, пов'язаної з безпекою> зберігаються в окремому сховищі або у вогнетривкому контейнері, не пов'язаному з системою.}




\subsubsection{ПЕРЕДАЧА НА АЛЬТЕРНАТИВНЕ СХОВИЩЕ ЗБЕРІГАННЯ (CP-9(5))}


\vspace{10pt}
{\footnotesize\bfseries
No: 1\\
Name: cp\_\allowbreak{}9\_\allowbreak{}5\_\allowbreak{}odp\_\allowbreak{}1\\
Type: string\\
Default: nil
}

{\footnotesize визначено період часу, що відповідає часу відновлення та цілям відновлення;}


{\footnotesize\bfseries
No: 2\\
Name: cp\_\allowbreak{}9\_\allowbreak{}5\_\allowbreak{}odp\_\allowbreak{}2\\
Type: string\\
Default: nil
}

{\footnotesize визначено швидкість передачі даних, що відповідає часу відновлення та цілям відновлення;}


{\footnotesize\bfseries
No: 3\\
Name: cp\_\allowbreak{}9\_\allowbreak{}5\_\allowbreak{}1\\
Type: string\\
Default: nil
}

{\footnotesize інформація резервної копії системи передається до альтернативного сховища протягом <CP-09(05)\_\allowbreak{}ODP[01] періоду часу>; 245}


{\footnotesize\bfseries
No: 4\\
Name: cp\_\allowbreak{}9\_\allowbreak{}5\_\allowbreak{}2\\
Type: string\\
Default: nil
}

{\footnotesize інформація резервної копії системи передається до альтернативного сховища з <CP-09(05)\_\allowbreak{}ODP[02] швидкість передачі >;}




\subsubsection{НАДЛИШКОВА ВТОРИННА СИСТЕМА (CP-9(6))}


\vspace{10pt}
{\footnotesize\bfseries
No: 1\\
Name: cp\_\allowbreak{}9\_\allowbreak{}6\_\allowbreak{}1\\
Type: string\\
Default: nil
}

{\footnotesize резервне копіювання системи здійснюється шляхом підтримки надлишкової вторинної системи, яка не пов'язана з первинною системою;}


{\footnotesize\bfseries
No: 2\\
Name: cp\_\allowbreak{}9\_\allowbreak{}6\_\allowbreak{}2\\
Type: string\\
Default: nil
}

{\footnotesize резервне копіювання системи здійснюється шляхом підтримки резервної вторинної системи, яка може бути активована без втрати інформації або порушення роботи.}




\subsubsection{ПОДВІЙНА АВТОРИЗАЦІЯ (CP-9(7))}


\vspace{10pt}
{\footnotesize\bfseries
No: 1\\
Name: cp\_\allowbreak{}9\_\allowbreak{}7\_\allowbreak{}odp\\
Type: string\\
Default: nil
}

{\footnotesize визначено резервну інформацію, для якої необхідно застосувати подвійну авторизацію з метою видалення або знищення;}


{\footnotesize\bfseries
No: 2\\
Name: cp\_\allowbreak{}9\_\allowbreak{}7\_\allowbreak{}01\\
Type: string\\
Default: nil
}

{\footnotesize застосовано подвійну авторизацію для видалення або знищення <CP-09(07)\_\allowbreak{}ODP резервної інформації>.}




\subsubsection{КРИПТОГРАФІЧНИЙ ЗАХИСТ (CP-9(8))}


\vspace{10pt}
{\footnotesize\bfseries
No: 1\\
Name: cp\_\allowbreak{}9\_\allowbreak{}8\_\allowbreak{}odp\\
Type: string\\
Default: nil
}

{\footnotesize визначено резервні копії інформації для захисту від несанкціонованого розкриття та змін;}


{\footnotesize\bfseries
No: 2\\
Name: cp\_\allowbreak{}9\_\allowbreak{}8\_\allowbreak{}01\\
Type: string\\
Default: nil
}

{\footnotesize реалізовано криптографічні механізми для запобігання несанкціонованому розкриттю та  зміні <CP-09(08)\_\allowbreak{}ODP резервної інформації>.}




\subsection{ВІДНОВЛЕННЯ ТА ВІДТВОРЕННЯ СИСТЕМИ (CP-10)}


\vspace{10pt}
{\footnotesize\bfseries
No: 1\\
Name: cp\_\allowbreak{}10\_\allowbreak{}odp\_\allowbreak{}1\\
Type: string\\
Default: nil
}

{\footnotesize визначено період часу для відновлення, часу та цілям відновлення системи;}


{\footnotesize\bfseries
No: 2\\
Name: cp\_\allowbreak{}10\_\allowbreak{}odp\_\allowbreak{}2\\
Type: string\\
Default: nil
}

{\footnotesize визначено період часу для відтворення, часу та цілям відновлення системи;}


{\footnotesize\bfseries
No: 3\\
Name: cp\_\allowbreak{}10\_\allowbreak{}1\\
Type: string\\
Default: nil
}

{\footnotesize відновлення системи до відомого стану забезпечується протягом}


{\footnotesize\bfseries
No: 4\\
Name: cp\_\allowbreak{}10\_\allowbreak{}2\\
Type: string\\
Default: nil
}

{\footnotesize відтворення системи до відомого стану забезпечується протягом}




\subsubsection{ВІДНОВЛЕННЯ ТРАНЗАКЦІЙ (CP-10(2))}


\vspace{10pt}
{\footnotesize\bfseries
No: 1\\
Name: cp\_\allowbreak{}10\_\allowbreak{}2\_\allowbreak{}01\\
Type: string\\
Default: nil
}

{\footnotesize реалізовано відновлення транзакцій для систем, що базуються на транзакціях}




\subsubsection{ВІДНОВЛЕННЯ В МЕЖАХ ЧАСОВОГО ПЕРІОДУ (CP-10(4))}


\vspace{10pt}
{\footnotesize\bfseries
No: 1\\
Name: cp\_\allowbreak{}10\_\allowbreak{}4\_\allowbreak{}odp\\
Type: string\\
Default: nil
}

{\footnotesize визначено період часу відновлення, протягом якого компоненти системи відновлюються до відомого, робочого стану;}


{\footnotesize\bfseries
No: 2\\
Name: cp\_\allowbreak{}10\_\allowbreak{}4\_\allowbreak{}01\\
Type: string\\
Default: nil
}

{\footnotesize забезпечено можливість відновлення компонентів системи протягом <CP-10(04)\_\allowbreak{}ODP період часу відновлення>  з інформації управління конфігурацією та захищеною цілісністю, яка описує відомий робочий стан компонентів.}




\subsubsection{ЗДАТНІСТЬ ВІДМОВОСТІЙКОСТІ (CP-10(5)) [Вилучено]}
[Вилучено: Включено до SI-13]

\vspace{10pt}
\textit{Немає параметрів для цього контролю.}
\vspace{10pt}


\subsection{АЛЬТЕРНАТИВНІ ПРОТОКОЛИ ЗВ'ЯЗКУ (CP-11)}


\vspace{10pt}
{\footnotesize\bfseries
No: 1\\
Name: cp\_\allowbreak{}11\_\allowbreak{}odp\\
Type: string\\
Default: nil
}

{\footnotesize організація визначає альтернативні протоколи зв'язку для підтримки збереження безперервності функціонування}


{\footnotesize\bfseries
No: 2\\
Name: cp\_\allowbreak{}11\_\allowbreak{}01\\
Type: string\\
Default: nil
}

{\footnotesize організація забезпечує можливість застосування <CP-11\_\allowbreak{}ODP альтернативних протоколів зв'язку > для підтримки збереження безперервності функціонування}




\subsection{БЕЗПЕЧНИЙ РЕЖИМ (CP-12)}


\vspace{10pt}
{\footnotesize\bfseries
No: 1\\
Name: cp\_\allowbreak{}12\_\allowbreak{}odp\_\allowbreak{}1\\
Type: string\\
Default: nil
}

{\footnotesize визначено умови, за яких організація вводить безпечний режим роботи;}


{\footnotesize\bfseries
No: 2\\
Name: cp\_\allowbreak{}12\_\allowbreak{}odp\_\allowbreak{}2\\
Type: string\\
Default: nil
}

{\footnotesize визначено обмеження в безпечному режимі роботи;}


{\footnotesize\bfseries
No: 3\\
Name: cp\_\allowbreak{}12\_\allowbreak{}01\\
Type: string\\
Default: nil
}

{\footnotesize при виявлені <CP-12\_\allowbreak{}ODP[01] умови>, вводиться безпечний режим роботи з <CP-12\_\allowbreak{}ODP[02] обмеженням>.}




\subsection{АЛЬТЕРНАТИВНІ МЕХАНІЗМИ БЕЗПЕКИ (CP-13)}


\vspace{10pt}
{\footnotesize\bfseries
No: 1\\
Name: cp\_\allowbreak{}13\_\allowbreak{}odp\_\allowbreak{}1\\
Type: string\\
Default: nil
}

{\footnotesize визначені альтернативні або додаткові механізми безпеки;}


{\footnotesize\bfseries
No: 2\\
Name: cp\_\allowbreak{}13\_\allowbreak{}odp\_\allowbreak{}2\\
Type: string\\
Default: nil
}

{\footnotesize визначені функції безпеки;}


{\footnotesize\bfseries
No: 3\\
Name: cp\_\allowbreak{}13\_\allowbreak{}01\\
Type: string\\
Default: nil
}

{\footnotesize <CP-13\_\allowbreak{}ODP[01] альтернативні або додаткові механізми безпеки> використовуються для реалізації <CP-13\_\allowbreak{}ODP[02] функцій безпеки> , коли основні засоби реалізації функцій безпеки недоступні або скомпрометовані.}




\section{IA}
Клас заходів захисту IA --- ІДЕНТИФІКАЦІЯ ТА АВТЕНТИФІКАЦІЯ

Цей клас відповідає за однозначне розпізнавання користувачів або пристроїв та підтвердження їхньої справжності перед наданням доступу до системи.

\textbf{Перелік заходів захисту:} Політика та процедури ідентифікації та автентифікації (IA-1); ІДЕНТИФІКАЦІЯ ТА АВТЕНТИФІКАЦІЯ (КОРИСТУВАЧІВ ОРГАНІЗАЦІЇ) (IA-2); Багатофакторна автентифікація привілейованих облікових записів (IA-2(1)); Багатофакторна автентифікація непривілейованих (IA-2(2)); Локальний доступ до привілейованих облікових записів (IA-2(3)) [Вилучено]; Локальний доступ до непривілейованих облікових записів (IA-2(4)) [Вилучено]; Індивідуальна автентифікація з груповою автентифікацією (IA-2(5)); Мережевий доступ до непривілейованих облікових записів -- окремий пристрій (IA-2(7)) [Вилучено]; Доступ до облікових записів -- стійкість до відтворення (IA-2(8)); Доступ до непривілейованих облікових записів -- стійкість до відтворення (IA-2(9)) [Вилучено]; Єдина точка входу (IA-2(10)); Віддалений доступ - окремий пристрій (IA-2(11)) [Вилучено]; ПРИЙНЯТТЯ ПОВНОВАЖЕНЬ ДЛЯ ВЕРИФІКАЦІЇ ОСОБИСТОЇ ІНФОРМАЦІЇ (PIV CREDENTIALS) (IA-2(12)); Автентифікація по зовнішньому каналу (IA-2(13)); Ідентифікація та автентифікація пристроїв (IA-3); Криптографічна двобічна автентифікація (IA-3(1)); Криптографічний двобічна мережа автентифікація [виключено: включено до іа-03(01)]. (IA-3(2)); Динамічний розподіл адреси (IA-3(3)); Ідентифікація та автентифікація пристрою виконується на основі атестації за допомогою (IA-3(4)); Управління ідентифікацією (IA-4); Заборона використання ідентифікаторів облікових записів таки самих, як й публічні ідентифікатори (IA-4(1)); Авторизація супервайзера [виключено: включено до іа-12(01)]. (IA-4(2)); Множинні форми сертифікації (IA-4(3)); Ідентифікація статусу користувача (IA-4(4)); Динамічне управління (IA-4(5)); Крос-організаційне управління (IA-4(6)); Особиста реєстрація [виключено: включено до іа-12(04)]. (IA-4(7)); Попарні псевдонімні ідентифікатори (IA-4(8)); Попарні псевдонімні ідентифікатори (IA-4(9)); Управління автентифікатором (IA-5); Автентифікація на основі пароля (IA-5(1)); Автентифікація на основі відкритого ключа (IA-5(2)); Особиста або довірча автентифікація зовнішньої сторони [виключено: включено до іа-12(04)]. (IA-5(3)); Автоматизована підтримка для визначення міцності пароля [виключено: включено до іа-05(01)]. (IA-5(4)); Зміна автентифікаторів до доставки (IA-5(5)); Захист автентифікаторів (IA-5(6)); Відсутність вбудованих незашифрованих статичних автентифікаторів (IA-5(7)); Багатосистемні облікові записи (IA-5(8)); Управління об'єднанням автентифікаторів (IA-5(9)); Динамічне зв'язування мандатів (IA-5(10)); Автентифікація на основі апаратних токенів (IA-5(11)) [Вилучено]; Ефективність біометричної автентифікації (IA-5(12)); Закінчення терміну шування автентифікаторів (IA-5(13)); Управління змістом довірчих сховищ інфраструктури відкритих ключів (IA-5(14)); Продукти та послуги, затверджені уповноваженим органом (IA-5(15)); Передача особистої або довірчої автентифікації зовнішньої сторони (IA-5(16)); Автоматизовані засоби виявлення атак із використанням біометричних автентифікаторів (IA-5(17)); Менеджер паролів (IA-5(18)); ПРИХОВУВАННЯ ЗВОРОТНОГО ЗВ'ЯЗКУ АВТЕНТИФІКАТОРА; Автентифікація криптографічного модуля (IA-7); ІДЕНТИФІКАЦІЯ ТА АВТЕНТИФІКАЦІЯ (КОРИСТУВАЧІ, ЩО НЕ НАЛЕЖАТЬ ДО ОРГАНІЗАЦІЇ) (IA-8); Використання затверджених продуктів (IA-8(3)) [Вилучено]; ВИЗНАННЯ ПОСВІДЧЕНЬ ОСОБИ (PIV-I) (IA-8(5)); Розмежування (IA-8(6)); Послуги ідентифікації та автентифікації (IA-9); Обмін інфо- (IA-9(1)); Передача рішень [виключено: перенесено до ia-09] (IA-9(2)); Адаптивна автентифікація (IA-10); Повторна автентифікація (IA-11); ПЕРЕВІРКА СПРАВЖНОСТІ (ІДЕНТИЧНОСТІ) (IA-12); Авторизація супервайзера (IA-12(1)); Посвідчення особи (IA-12(2)); Очна перевірка та верифікація (IA-12(4)); Підтвердження адреси (IA-12(5)); Прийняття ідентифікацій схвалених третьою стороною (IA-12(6)).

\subsection{ПОЛІТИКА ТА ПРОЦЕДУРИ ІДЕНТИФІКАЦІЇ ТА АВТЕНТИФІКАЦІЇ (IA-1)}


\vspace{10pt}
{\footnotesize\bfseries
No: 1\\
Name: ia\_\allowbreak{}1\_\allowbreak{}odp\_\allowbreak{}1\\
Type: string\\
Default: nil
}

{\footnotesize визначено персонал або ролі, на які поширюється політика ідентифікації та автентифікації;}


{\footnotesize\bfseries
No: 2\\
Name: ia\_\allowbreak{}1\_\allowbreak{}odp\_\allowbreak{}2\\
Type: string\\
Default: nil
}

{\footnotesize визначено персонал або ролі, на які поширюються процедури ідентифікації та автентифікації;}


{\footnotesize\bfseries
No: 3\\
Name: ia\_\allowbreak{}1\_\allowbreak{}odp\_\allowbreak{}3\\
Type: string\\
Default: nil
}

{\footnotesize вибрано одне або декілька з наступних ЗНАЧЕНЬ ПАРА- МЕТРІВ: \{рівень організації; рівень місії/бізнес -процесів; рівень системи\};}


{\footnotesize\bfseries
No: 4\\
Name: ia\_\allowbreak{}1\_\allowbreak{}odp\_\allowbreak{}4\\
Type: string\\
Default: nil
}

{\footnotesize визначено посадову особу, яка управляє політикою та процедурами ідентифікації та автентифікації;}


{\footnotesize\bfseries
No: 5\\
Name: ia\_\allowbreak{}1\_\allowbreak{}odp\_\allowbreak{}5\\
Type: string\\
Default: nil
}

{\footnotesize визначено частоту, з якою переглядається та оновлюється поточна політика ідентифікації та автентифікації;}


{\footnotesize\bfseries
No: 6\\
Name: ia\_\allowbreak{}1\_\allowbreak{}odp\_\allowbreak{}6\\
Type: string\\
Default: nil
}

{\footnotesize визначено події, які потребують перегляду та оновлення поточної політики ідентифікації та автентифікації;}


{\footnotesize\bfseries
No: 7\\
Name: ia\_\allowbreak{}1\_\allowbreak{}odp\_\allowbreak{}7\\
Type: string\\
Default: nil
}

{\footnotesize визначено частоту, з якою переглядаються та оновлюються поточні процедури ідентифікації та автентифікації;}


{\footnotesize\bfseries
No: 8\\
Name: ia\_\allowbreak{}1\_\allowbreak{}odp\_\allowbreak{}8\\
Type: string\\
Default: nil
}

{\footnotesize визначено події, які потребують перегляду та оновлення процедур ідентифікації та автентифікації;}


{\footnotesize\bfseries
No: 9\\
Name: ia\_\allowbreak{}1\_\allowbreak{}a\_\allowbreak{}1\\
Type: string\\
Default: nil
}

{\footnotesize розроблено та задокументовано політику ідентифікації та автентифікації;}


{\footnotesize\bfseries
No: 10\\
Name: ia\_\allowbreak{}1\_\allowbreak{}a\_\allowbreak{}2\\
Type: string\\
Default: nil
}

{\footnotesize політика ідентифікації та автентифікації поширюється на  <IA01\_\allowbreak{}ODP[01] персонал або ролі>;}


{\footnotesize\bfseries
No: 11\\
Name: ia\_\allowbreak{}1\_\allowbreak{}a\_\allowbreak{}3\\
Type: string\\
Default: nil
}

{\footnotesize розроблені та задокументовані процедури ідентифікації та автентифікації, що сприяють впровадженню політики ідентифікації та автентифікації, а також відповідні заходи ідентифікації та перевірки автентичності;}


{\footnotesize\bfseries
No: 12\\
Name: ia\_\allowbreak{}1\_\allowbreak{}a\_\allowbreak{}4\\
Type: string\\
Default: nil
}

{\footnotesize процедури ідентифікації та автентифікації поширюються на  <IA01\_\allowbreak{}ODP[02] персонал або ролі>;}


{\footnotesize\bfseries
No: 13\\
Name: ia\_\allowbreak{}1\_\allowbreak{}a\_\allowbreak{}1\_\allowbreak{}a\_\allowbreak{}1\\
Type: string\\
Default: nil
}

{\footnotesize <IA-01\_\allowbreak{}ODP[03] ВИБРАНЕ ЗНАЧЕННЯ ПАРАМЕТРА(ів)> політики ідентифікації та автентифікації містить мету; 253}


{\footnotesize\bfseries
No: 14\\
Name: ia\_\allowbreak{}1\_\allowbreak{}a\_\allowbreak{}1\_\allowbreak{}a\_\allowbreak{}2\\
Type: string\\
Default: nil
}

{\footnotesize <IA-01\_\allowbreak{}ODP[03] ВИБРАНЕ ЗНАЧЕННЯ ПАРАМЕТРА(ів)> політики ідентифікації та автентифікації містить сферу застосування;}


{\footnotesize\bfseries
No: 15\\
Name: ia\_\allowbreak{}1\_\allowbreak{}a\_\allowbreak{}1\_\allowbreak{}a\_\allowbreak{}3\\
Type: string\\
Default: nil
}

{\footnotesize <IA-01\_\allowbreak{}ODP[03] ВИБРАНЕ ЗНАЧЕННЯ ПАРАМЕТРА(ів)> політики ідентифікації та автентифікації містить ролі;}


{\footnotesize\bfseries
No: 16\\
Name: ia\_\allowbreak{}1\_\allowbreak{}a\_\allowbreak{}1\_\allowbreak{}a\_\allowbreak{}4\\
Type: string\\
Default: nil
}

{\footnotesize <IA-01\_\allowbreak{}ODP[03] ВИБРАНЕ ЗНАЧЕННЯ ПАРАМЕТРА(ів)> політики ідентифікації та автентифікації містить обов'язки;}


{\footnotesize\bfseries
No: 17\\
Name: ia\_\allowbreak{}1\_\allowbreak{}a\_\allowbreak{}1\_\allowbreak{}a\_\allowbreak{}5\\
Type: string\\
Default: nil
}

{\footnotesize <IA-01\_\allowbreak{}ODP[03] ВИБРАНЕ ЗНАЧЕННЯ ПАРАМЕТРА(ів)> політики ідентифікації та автентифікації містить відповідальність керівництва;}


{\footnotesize\bfseries
No: 18\\
Name: ia\_\allowbreak{}1\_\allowbreak{}a\_\allowbreak{}1\_\allowbreak{}a\_\allowbreak{}6\\
Type: string\\
Default: nil
}

{\footnotesize <IA-01\_\allowbreak{}ODP[03] ВИБРАНЕ ЗНАЧЕННЯ ПАРАМЕТРА(ів)> політики ідентифікації та автентифікації містить координацію між підрозділами організації;}


{\footnotesize\bfseries
No: 19\\
Name: ia\_\allowbreak{}1\_\allowbreak{}a\_\allowbreak{}1\_\allowbreak{}a\_\allowbreak{}7\\
Type: string\\
Default: nil
}

{\footnotesize <IA-01\_\allowbreak{}ODP[03] ВИБРАНЕ ЗНАЧЕННЯ ПАРАМЕТРА(ів)> політики ідентифікації та автентифікації містить систему контролю відповідності;}


{\footnotesize\bfseries
No: 20\\
Name: ia\_\allowbreak{}1\_\allowbreak{}a\_\allowbreak{}1\_\allowbreak{}b\\
Type: string\\
Default: nil
}

{\footnotesize політика ідентифікації та автентифікації  <IA-01\_\allowbreak{}ODP[03] ВИБРАНЕ ЗНАЧЕННЯ ПАРАМЕТРА(ів)> відповідає чинному законодавству, виконавчим розпорядженням, директивам, положенням, політиці, стандартам і керівним принципам;}


{\footnotesize\bfseries
No: 21\\
Name: ia\_\allowbreak{}1\_\allowbreak{}b\\
Type: string\\
Default: nil
}

{\footnotesize <IA-01\_\allowbreak{}ODP[04] посадова особа> призначається для управління політикою та процедурами ідентифікації та автентифікації;}


{\footnotesize\bfseries
No: 22\\
Name: ia\_\allowbreak{}1\_\allowbreak{}c\_\allowbreak{}1\_\allowbreak{}1\\
Type: string\\
Default: nil
}

{\footnotesize переглядається та оновлюється поточна політика ідентифікації та}


{\footnotesize\bfseries
No: 23\\
Name: ia\_\allowbreak{}1\_\allowbreak{}c\_\allowbreak{}1\_\allowbreak{}2\\
Type: string\\
Default: nil
}

{\footnotesize переглядається та оновлюється поточна політика ідентифікації та}


{\footnotesize\bfseries
No: 24\\
Name: ia\_\allowbreak{}1\_\allowbreak{}c\_\allowbreak{}2\_\allowbreak{}1\\
Type: string\\
Default: nil
}

{\footnotesize переглядаються та оновлю ються поточні процедури ідентифікації}


{\footnotesize\bfseries
No: 25\\
Name: ia\_\allowbreak{}1\_\allowbreak{}c\_\allowbreak{}2\_\allowbreak{}2\\
Type: string\\
Default: nil
}

{\footnotesize переглядаються та оновлю ються поточні процедури ідентифікації}




\subsection{ІДЕНТИФІКАЦІЯ ТА АВТЕНТИФІКАЦІЯ (КОРИСТУВАЧІВ ОРГАНІЗАЦІЇ) (IA-2)}


\vspace{10pt}
{\footnotesize\bfseries
No: 1\\
Name: ia\_\allowbreak{}2\_\allowbreak{}1\\
Type: string\\
Default: nil
}

{\footnotesize користувачі унікально ідентифіковані та автентифіковані;}


{\footnotesize\bfseries
No: 2\\
Name: ia\_\allowbreak{}2\_\allowbreak{}2\\
Type: string\\
Default: nil
}

{\footnotesize процеси що діють від імені користувачів унікально ідентифіковані та автентифіковані;}




\subsubsection{БАГАТОФАКТОРНА АВТЕНТИФІКАЦІЯ ПРИВІЛЕЙОВАНИХ ОБЛІКОВИХ ЗАПИСІВ (IA-2(1))}


\vspace{10pt}
{\footnotesize\bfseries
No: 1\\
Name: ia\_\allowbreak{}2\_\allowbreak{}1\_\allowbreak{}01\\
Type: string\\
Default: nil
}

{\footnotesize реалізувано багатофакторну автентифікацію для доступу до привілейованих облікових записів}




\subsubsection{БАГАТОФАКТОРНА АВТЕНТИФІКАЦІЯ НЕПРИВІЛЕЙОВАНИХ (IA-2(2))}


\vspace{10pt}
{\footnotesize\bfseries
No: 1\\
Name: ia\_\allowbreak{}2\_\allowbreak{}2\_\allowbreak{}01\\
Type: string\\
Default: nil
}

{\footnotesize реалізовано багатофакторну автентифікацію для доступу до непривілейованих облікових записів}




\subsubsection{ЛОКАЛЬНИЙ ДОСТУП ДО ПРИВІЛЕЙОВАНИХ ОБЛІКОВИХ ЗАПИСІВ (IA-2(3)) [Вилучено]}
[Вилучено: Включено до ІА-02(01)]

\vspace{10pt}
\textit{Немає параметрів для цього контролю.}
\vspace{10pt}


\subsubsection{ЛОКАЛЬНИЙ ДОСТУП ДО НЕПРИВІЛЕЙОВАНИХ ОБЛІКОВИХ ЗАПИСІВ (IA-2(4)) [Вилучено]}
[Вилучено: Включено до ІА-02(02)]

\vspace{10pt}
\textit{Немає параметрів для цього контролю.}
\vspace{10pt}


\subsubsection{ІНДИВІДУАЛЬНА АВТЕНТИФІКАЦІЯ З ГРУПОВОЮ АВТЕНТИФІКАЦІЄЮ (IA-2(5))}


\vspace{10pt}
{\footnotesize\bfseries
No: 1\\
Name: ia\_\allowbreak{}2\_\allowbreak{}5\_\allowbreak{}01\\
Type: string\\
Default: nil
}

{\footnotesize користувачі повинні пройти індивідуальну автентифікацію перед наданням доступу до спільних облікових записів або ресурсів, якщо використовуються спільні облікові записи або автентифікатори.}




\subsubsection{МЕРЕЖЕВИЙ ДОСТУП ДО НЕПРИВІЛЕЙОВАНИХ ОБЛІКОВИХ ЗАПИСІВ -- ОКРЕМИЙ ПРИСТРІЙ (IA-2(7)) [Вилучено]}
[Вилучено: Включено до ІА-02(01)]

\vspace{10pt}
\textit{Немає параметрів для цього контролю.}
\vspace{10pt}


\subsubsection{ДОСТУП ДО ОБЛІКОВИХ ЗАПИСІВ -- СТІЙКІСТЬ ДО ВІДТВОРЕННЯ (IA-2(8))}


\vspace{10pt}
{\footnotesize\bfseries
No: 1\\
Name: ia\_\allowbreak{}2\_\allowbreak{}8\_\allowbreak{}odp\\
Type: string\\
Default: nil
}

{\footnotesize вибрано одне або декілька з наступних ЗНАЧЕНЬ ПАРА- МЕТРІВ: \{привілейовані облікові записи; непривілейовані облікові записи\};}


{\footnotesize\bfseries
No: 2\\
Name: ia\_\allowbreak{}2\_\allowbreak{}8\_\allowbreak{}01\\
Type: string\\
Default: nil
}

{\footnotesize реалізовано стійкі до повторного відтворення механізми автентифікації для доступу до <IA-02(08)\_\allowbreak{}ODP ВИБРАНЕ ЗНАЧЕННЯ ПАРАМЕТРА(ів)>.}




\subsubsection{ДОСТУП ДО НЕПРИВІЛЕЙОВАНИХ ОБЛІКОВИХ ЗАПИСІВ -- СТІЙКІСТЬ ДО ВІДТВОРЕННЯ (IA-2(9)) [Вилучено]}
[Вилучено: Включено до ІА-02(08)]

\vspace{10pt}
\textit{Немає параметрів для цього контролю.}
\vspace{10pt}


\subsubsection{ЄДИНА ТОЧКА ВХОДУ (IA-2(10))}


\vspace{10pt}
{\footnotesize\bfseries
No: 1\\
Name: ia\_\allowbreak{}2\_\allowbreak{}10\_\allowbreak{}odp\\
Type: string\\
Default: nil
}

{\footnotesize визначено облікові записи та послуги системи, для яких має бути забезпечена можливість єдиного входу;}


{\footnotesize\bfseries
No: 2\\
Name: ia\_\allowbreak{}2\_\allowbreak{}10\_\allowbreak{}01\\
Type: string\\
Default: nil
}

{\footnotesize забезпечено можливість єдиного входу для <IA-02(10)\_\allowbreak{}ODP облікових записів і послуг системи>.}




\subsubsection{ВІДДАЛЕНИЙ ДОСТУП - ОКРЕМИЙ ПРИСТРІЙ (IA-2(11)) [Вилучено]}
[Вилучено: Включено до ІА-02(06)]

\vspace{10pt}
\textit{Немає параметрів для цього контролю.}
\vspace{10pt}


\subsubsection{ПРИЙНЯТТЯ ПОВНОВАЖЕНЬ ДЛЯ ВЕРИФІКАЦІЇ ОСОБИСТОЇ ІНФОРМАЦІЇ (PIV CREDENTIALS) (IA-2(12))}


\vspace{10pt}
{\footnotesize\bfseries
No: 1\\
Name: ia\_\allowbreak{}2\_\allowbreak{}12\_\allowbreak{}01\\
Type: string\\
Default: nil
}

{\footnotesize приймаються та електронним шляхом підтверджуються повноваження облікових даних особистої ідентифікації.}




\subsubsection{АВТЕНТИФІКАЦІЯ ПО ЗОВНІШНЬОМУ КАНАЛУ (IA-2(13))}


\vspace{10pt}
{\footnotesize\bfseries
No: 1\\
Name: ia\_\allowbreak{}2\_\allowbreak{}13\_\allowbreak{}odp\_\allowbreak{}1\\
Type: string\\
Default: nil
}

{\footnotesize визначено механізми зовнішньої автентифікації;}


{\footnotesize\bfseries
No: 2\\
Name: ia\_\allowbreak{}2\_\allowbreak{}13\_\allowbreak{}odp\_\allowbreak{}2\\
Type: string\\
Default: nil
}

{\footnotesize визначено умови, за яких має бути реалізована зовнішня автентифікація;}


{\footnotesize\bfseries
No: 3\\
Name: ia\_\allowbreak{}2\_\allowbreak{}13\_\allowbreak{}01\\
Type: string\\
Default: nil
}

{\footnotesize <IA-02(13)\_\allowbreak{}ODP[01] механізми зовнішньої автентифікації> застосовано за <IA-02(13)\_\allowbreak{}ODP[02] умов>.}




\subsection{ІДЕНТИФІКАЦІЯ ТА АВТЕНТИФІКАЦІЯ ПРИСТРОЇВ (IA-3)}


\vspace{10pt}
{\footnotesize\bfseries
No: 1\\
Name: ia\_\allowbreak{}3\_\allowbreak{}odp\_\allowbreak{}1\\
Type: string\\
Default: nil
}

{\footnotesize визначені пристрої та/або типи пристроїв, які повинні бути унікально ідентифіковані та автентифіковані перед установкою підключення;}


{\footnotesize\bfseries
No: 2\\
Name: ia\_\allowbreak{}3\_\allowbreak{}odp\_\allowbreak{}2\\
Type: string\\
Default: nil
}

{\footnotesize вибрано одне або декілька з наступних ЗНАЧЕНЬ ПАРА- МЕТРІВ: \{локальний; віддалений; мережевий\}; 260}


{\footnotesize\bfseries
No: 3\\
Name: ia\_\allowbreak{}3\_\allowbreak{}01\\
Type: string\\
Default: nil
}

{\footnotesize <IA-03\_\allowbreak{}ODP[01] пристрої та/або типи пристроїв>  унікально ідентифіковані та автентифіковано перед встановленням підключення}




\subsubsection{КРИПТОГРАФІЧНА ДВОБІЧНА АВТЕНТИФІКАЦІЯ (IA-3(1))}


\vspace{10pt}
{\footnotesize\bfseries
No: 1\\
Name: ia\_\allowbreak{}3\_\allowbreak{}1\_\allowbreak{}odp\_\allowbreak{}1\\
Type: string\\
Default: nil
}

{\footnotesize визначено пристрої та/або типи пристроїв, які потребують використання двобічної автентифікації яка заснована на криптографічних механізмах для автентифікації перед встановленням підключення;}


{\footnotesize\bfseries
No: 2\\
Name: ia\_\allowbreak{}3\_\allowbreak{}1\_\allowbreak{}odp\_\allowbreak{}2\\
Type: string\\
Default: nil
}

{\footnotesize вибрано одне або декілька з наступних ЗНАЧЕНЬ ПАРАМЕТРІВ: \{локальний; віддалений; мережевий\};}


{\footnotesize\bfseries
No: 3\\
Name: ia\_\allowbreak{}3\_\allowbreak{}1\_\allowbreak{}01\\
Type: string\\
Default: nil
}

{\footnotesize <IA-03(01)\_\allowbreak{}ODP[01] пристрої та/або типи пристроїв> автентифікуються перед встановленням підключення  <IA03(01)\_\allowbreak{}ODP[02] ВИБРАНЕ ЗНАЧЕННЯ ПАРАМЕТРА(ів)> за допомогою двобічної автентифікації, яка заснована на криптографічних механізмах.}




\subsubsection{КРИПТОГРАФІЧНИЙ ДВОБІЧНА МЕРЕЖА АВТЕНТИФІКАЦІЯ [Виключено: включено до ІА-03(01)]. (IA-3(2))}


\vspace{10pt}
\textit{Немає параметрів для цього контролю.}
\vspace{10pt}


\subsubsection{ДИНАМІЧНИЙ РОЗПОДІЛ АДРЕСИ (IA-3(3))}


\vspace{10pt}
{\footnotesize\bfseries
No: 1\\
Name: ia\_\allowbreak{}3\_\allowbreak{}3\_\allowbreak{}odp\_\allowbreak{}1\\
Type: string\\
Default: nil
}

{\footnotesize визначено інформацію про оренду, яку буде використано для стандартизації динамічного розподілу адрес для пристроїв;}


{\footnotesize\bfseries
No: 2\\
Name: ia\_\allowbreak{}3\_\allowbreak{}3\_\allowbreak{}odp\_\allowbreak{}2\\
Type: string\\
Default: nil
}

{\footnotesize визначено тривалість оренди, яка буде використовуватися для стандартизації динамічного виділення адрес для пристроїв;}


{\footnotesize\bfseries
No: 3\\
Name: ia\_\allowbreak{}3\_\allowbreak{}3\_\allowbreak{}a\_\allowbreak{}1\\
Type: string\\
Default: nil
}

{\footnotesize інформація про оренду динамічного розподілу адрес, що призначається пристроям з динамічним розподілом адрес,}


{\footnotesize\bfseries
No: 4\\
Name: ia\_\allowbreak{}3\_\allowbreak{}3\_\allowbreak{}a\_\allowbreak{}2\\
Type: string\\
Default: nil
}

{\footnotesize тривалість оренди динамічного виділення адреси, що призначається пристроям з динамічним виділенням адреси, стандартизовано відповідно до <IA-03(03)\_\allowbreak{}ODP[02] тривалість оренди>;}


{\footnotesize\bfseries
No: 5\\
Name: ia\_\allowbreak{}3\_\allowbreak{}3\_\allowbreak{}b\\
Type: string\\
Default: nil
}

{\footnotesize інформація про оренду  перевіряється, коли її призначено пристрою.}




\subsubsection{ІДЕНТИФІКАЦІЯ ТА АВТЕНТИФІКАЦІЯ ПРИСТРОЮ ВИКОНУЄТЬСЯ НА ОСНОВІ АТЕСТАЦІЇ ЗА ДОПОМОГОЮ (IA-3(4))}


\vspace{10pt}
{\footnotesize\bfseries
No: 1\\
Name: ia\_\allowbreak{}3\_\allowbreak{}4\_\allowbreak{}odp\\
Type: string\\
Default: nil
}

{\footnotesize визначено процес управління конфігурацією, який буде використовуватися для ідентифікації та автентифікації пристроїв на основі атестації;}


{\footnotesize\bfseries
No: 2\\
Name: ia\_\allowbreak{}3\_\allowbreak{}4\_\allowbreak{}01\\
Type: string\\
Default: nil
}

{\footnotesize ідентифікація та автентифікація пристрою виконується на основі атестації за допомогою  <IA-03(04)\_\allowbreak{}ODP процес управління конфігурацією>.}




\subsection{УПРАВЛІННЯ ІДЕНТИФІКАЦІЄЮ (IA-4)}


\vspace{10pt}
{\footnotesize\bfseries
No: 1\\
Name: ia\_\allowbreak{}4\_\allowbreak{}odp\_\allowbreak{}1\\
Type: string\\
Default: nil
}

{\footnotesize визначено персонал або ролі, від яких необхідно отримати дозвіл на призначення ідентифікатора;}


{\footnotesize\bfseries
No: 2\\
Name: ia\_\allowbreak{}4\_\allowbreak{}odp\_\allowbreak{}2\\
Type: string\\
Default: nil
}

{\footnotesize визначено період часу для запобігання повторному використанню ідентифікаторів;}


{\footnotesize\bfseries
No: 3\\
Name: ia\_\allowbreak{}4\_\allowbreak{}a\\
Type: string\\
Default: nil
}

{\footnotesize управління ідентифікаторами здійснюється шляхом отримання дозволу від <IA-04\_\allowbreak{}ODP[01] персоналу або ролей>  на призначення ідентифікатора особі, групі, ролі або пристрою;}


{\footnotesize\bfseries
No: 4\\
Name: ia\_\allowbreak{}4\_\allowbreak{}b\\
Type: string\\
Default: nil
}

{\footnotesize управління ідентифікаторами здійснюється шляхом вибору ідентифікатора, який ідентифікує окрему особу, групу, ролі або пристрій;}


{\footnotesize\bfseries
No: 5\\
Name: ia\_\allowbreak{}4\_\allowbreak{}c\\
Type: string\\
Default: nil
}

{\footnotesize управління ідентифікаторами здійснюється шляхом призначення ідентифікатора особі, групі, ролі або пристрою; 263}


{\footnotesize\bfseries
No: 6\\
Name: ia\_\allowbreak{}4\_\allowbreak{}d\\
Type: string\\
Default: nil
}

{\footnotesize ідентифікатори управляються шляхом запобігання повторному використанню ідентифікаторів впродовж <IA-04\_\allowbreak{}ODP[02] період часу>.}




\subsubsection{ЗАБОРОНА ВИКОРИСТАННЯ ІДЕНТИФІКАТОРІВ ОБЛІКОВИХ ЗАПИСІВ ТАКИ САМИХ, ЯК Й ПУБЛІЧНІ ІДЕНТИФІКАТОРИ (IA-4(1))}


\vspace{10pt}
{\footnotesize\bfseries
No: 1\\
Name: ia\_\allowbreak{}4\_\allowbreak{}1\_\allowbreak{}01\\
Type: string\\
Default: nil
}

{\footnotesize заборонено використання ідентифікаторів облікових записів системи, які збігаються із загальнодоступними ідентифікаторами для індивідуальних облікових записів.}




\subsubsection{АВТОРИЗАЦІЯ СУПЕРВАЙЗЕРА [Виключено: Включено до ІА-12(01)]. (IA-4(2))}


\vspace{10pt}
{\footnotesize\bfseries
No: 1\\
Name: ia\_\allowbreak{}4\_\allowbreak{}2\_\allowbreak{}01\\
Type: string\\
Default: nil
}

{\footnotesize УПРАВЛІННЯ ІДЕНТИФІКАЦІЄЮ - АВТОРИЗАЦІЯ СУПЕРВАЙЗЕРА [Виключено: Включено до ІА-12(01)]}




\subsubsection{МНОЖИННІ ФОРМИ СЕРТИФІКАЦІЇ (IA-4(3))}


\vspace{10pt}
{\footnotesize\bfseries
No: 1\\
Name: ia\_\allowbreak{}4\_\allowbreak{}3\_\allowbreak{}01\\
Type: string\\
Default: nil
}

{\footnotesize УПРАВЛІННЯ ІДЕНТИФІКАЦІЄЮ - МНОЖИННІ ФОРМИ СЕРТИФІКАЦІЇ [Виключено: Включено до ІА-12(02)]}




\subsubsection{ІДЕНТИФІКАЦІЯ СТАТУСУ КОРИСТУВАЧА (IA-4(4))}


\vspace{10pt}
{\footnotesize\bfseries
No: 1\\
Name: ia\_\allowbreak{}4\_\allowbreak{}4\_\allowbreak{}odp\\
Type: string\\
Default: nil
}

{\footnotesize визначено ознаку, що ідентифікує індивідуальний статус;}


{\footnotesize\bfseries
No: 2\\
Name: ia\_\allowbreak{}4\_\allowbreak{}4\_\allowbreak{}01\\
Type: string\\
Default: nil
}

{\footnotesize управління індивідуальними ідентифікаторами, однозначно ідентифікуючи кожного індивідуума <IA-04(04)\_\allowbreak{}ODP ознака>, що ідентифікує індивідуальний статус.}




\subsubsection{ДИНАМІЧНЕ УПРАВЛІННЯ (IA-4(5))}


\vspace{10pt}
{\footnotesize\bfseries
No: 1\\
Name: ia\_\allowbreak{}4\_\allowbreak{}5\_\allowbreak{}odp\\
Type: string\\
Default: nil
}

{\footnotesize визначено політику динамічних ідентифікаторів для управління індивідуальними ідентифікаторами;}


{\footnotesize\bfseries
No: 2\\
Name: ia\_\allowbreak{}4\_\allowbreak{}5\_\allowbreak{}01\\
Type: string\\
Default: nil
}

{\footnotesize індивідуальні ідентифікатори динамічно управляються відповідно до <IA-04(05)\_\allowbreak{}ODP політика динамічних ідентифікаторів>.}




\subsubsection{КРОС-ОРГАНІЗАЦІЙНЕ УПРАВЛІННЯ (IA-4(6))}


\vspace{10pt}
{\footnotesize\bfseries
No: 1\\
Name: ia\_\allowbreak{}4\_\allowbreak{}6\_\allowbreak{}odp\\
Type: string\\
Default: nil
}

{\footnotesize визначено зовнішні організації з якими необхідно здійснювати координацію для крос-організаційного управління ідентифікаторами;}


{\footnotesize\bfseries
No: 2\\
Name: ia\_\allowbreak{}4\_\allowbreak{}6\_\allowbreak{}01\\
Type: string\\
Default: nil
}

{\footnotesize здійснюється координація з  <IA-04(06)\_\allowbreak{}ODP зовнішні організації> для крос-організаційного управління ідентифікаторами.}




\subsubsection{ОСОБИСТА РЕЄСТРАЦІЯ [Виключено: Включено до ІА-12(04)]. (IA-4(7))}


\vspace{10pt}
{\footnotesize\bfseries
No: 1\\
Name: ia\_\allowbreak{}4\_\allowbreak{}7\_\allowbreak{}01\\
Type: list\\
Default: [\textquotedbl{}admin\textquotedbl{}, \textquotedbl{}security\_\allowbreak{}officer\textquotedbl{}]
}

{\footnotesize УПРАВЛІННЯ ІДЕНТИФІКАЦІЄЮ - ОСОБИСТА РЕЄСТРАЦІЯ [Виключено: Включено до ІА-12(04)]}




\subsubsection{ПОПАРНІ ПСЕВДОНІМНІ ІДЕНТИФІКАТОРИ (IA-4(8))}


\vspace{10pt}
{\footnotesize\bfseries
No: 1\\
Name: ia\_\allowbreak{}4\_\allowbreak{}8\_\allowbreak{}01\\
Type: string\\
Default: nil
}

{\footnotesize створено попарні псевдонімні ідентифікатори.}




\subsubsection{ПОПАРНІ ПСЕВДОНІМНІ ІДЕНТИФІКАТОРИ (IA-4(9))}


\vspace{10pt}
{\footnotesize\bfseries
No: 1\\
Name: ia\_\allowbreak{}4\_\allowbreak{}9\_\allowbreak{}odp\\
Type: string\\
Default: nil
}

{\footnotesize визначено захищене центральне сховище, яке використовується для зберігання атрибутів для кожної унікально ідентифікованої особи, пристрою або служби;}


{\footnotesize\bfseries
No: 2\\
Name: ia\_\allowbreak{}4\_\allowbreak{}9\_\allowbreak{}01\\
Type: string\\
Default: nil
}

{\footnotesize атрибути для кожної унікально ідентифікованої особи, пристрою або служби зберігаються у  <IA-04(09)\_\allowbreak{}ODP захищеному центральному сховищі>.}




\subsection{УПРАВЛІННЯ АВТЕНТИФІКАТОРОМ (IA-5)}


\vspace{10pt}
{\footnotesize\bfseries
No: 1\\
Name: ia\_\allowbreak{}5\_\allowbreak{}odp\_\allowbreak{}1\\
Type: string\\
Default: nil
}

{\footnotesize визначено період часу для зміни або оновлення автентифікаторів за типом автентифікатора;}


{\footnotesize\bfseries
No: 2\\
Name: ia\_\allowbreak{}5\_\allowbreak{}odp\_\allowbreak{}2\\
Type: string\\
Default: nil
}

{\footnotesize визначено події, які викликають зміну або оновлення автентифікаторів;}


{\footnotesize\bfseries
No: 3\\
Name: ia\_\allowbreak{}5\_\allowbreak{}a\\
Type: string\\
Default: nil
}

{\footnotesize управління системними автентифікаторами здійснюється шляхом перевірки, як частини початкового розподілу автентифікатора, особи, групи, ролі або пристрою, який отримує автентифікатор;}


{\footnotesize\bfseries
No: 4\\
Name: ia\_\allowbreak{}5\_\allowbreak{}b\\
Type: string\\
Default: nil
}

{\footnotesize управління системними автентифікаторами здійснюється шляхом створення вихідного вмісту автентифікатора для будь -яких автентифікаторів, виданих організацією;}


{\footnotesize\bfseries
No: 5\\
Name: ia\_\allowbreak{}5\_\allowbreak{}c\\
Type: string\\
Default: nil
}

{\footnotesize управління системними автентифікаторами здійснюється шляхом забезпечення того, щоб автентифікатори мали достатню стійкість механізму для їх використання за призначенням; 267}


{\footnotesize\bfseries
No: 6\\
Name: ia\_\allowbreak{}5\_\allowbreak{}d\\
Type: string\\
Default: nil
}

{\footnotesize управління системними автентифікаторами здійснюється шляхом створення та реалізація адміністративних процедур для первинного розповсюдження автентифікаторів, для втрачених/ск омпрометованих або пошкоджених автентифікаторів, а також для відкликання автентифікаторів;}


{\footnotesize\bfseries
No: 7\\
Name: ia\_\allowbreak{}5\_\allowbreak{}e\\
Type: string\\
Default: nil
}

{\footnotesize управління системними автентифікаторами здійснюється шляхом зміни типових автентифікаторів перед першим використанням;}


{\footnotesize\bfseries
No: 8\\
Name: ia\_\allowbreak{}5\_\allowbreak{}f\\
Type: string\\
Default: nil
}

{\footnotesize управління системними ав тентифікаторами здійснюється шляхом}


{\footnotesize\bfseries
No: 9\\
Name: ia\_\allowbreak{}5\_\allowbreak{}g\\
Type: string\\
Default: nil
}

{\footnotesize управління системними автентифікаторами здійснюється шляхом захисту вмісту автентифікатора від несанкціонованого розкриття та модифікацій;}


{\footnotesize\bfseries
No: 10\\
Name: ia\_\allowbreak{}5\_\allowbreak{}h\\
Type: string\\
Default: nil
}

{\footnotesize управління системними автентифікаторами здійснюється шляхом вимоги до осіб, які використовують пристрої, використовувати спеціальні заходи безпеки для захисту автентифікаторів;}


{\footnotesize\bfseries
No: 11\\
Name: ia\_\allowbreak{}5\_\allowbreak{}i\\
Type: string\\
Default: nil
}

{\footnotesize управління системними автентифікаторами здійснюється шляхом вимоги змінювати автентифікатори для облікових записів груп/ролей при зміні членства в цих облікових записах.}




\subsubsection{АВТЕНТИФІКАЦІЯ НА ОСНОВІ ПАРОЛЯ (IA-5(1))}


\vspace{10pt}
{\footnotesize\bfseries
No: 1\\
Name: ia\_\allowbreak{}5\_\allowbreak{}1\_\allowbreak{}odp\_\allowbreak{}1\\
Type: string\\
Default: nil
}

{\footnotesize визначено частоту оновлення списку часто використовуваних, очікуваних або скомпрометованих паролів;}


{\footnotesize\bfseries
No: 2\\
Name: ia\_\allowbreak{}5\_\allowbreak{}1\_\allowbreak{}odp\_\allowbreak{}2\\
Type: string\\
Default: nil
}

{\footnotesize визначено склад та правила складності автентифікатора;}


{\footnotesize\bfseries
No: 3\\
Name: ia\_\allowbreak{}5\_\allowbreak{}1\_\allowbreak{}a\\
Type: string\\
Default: nil
}

{\footnotesize для автентифікації на основі паролів підтримується та оновлюється список часто використовуваних, очікуваних або 268 також коли є підозра, що паролі організації були скомпрометовані прямо чи опосередковано;}


{\footnotesize\bfseries
No: 4\\
Name: ia\_\allowbreak{}5\_\allowbreak{}1\_\allowbreak{}b\\
Type: string\\
Default: nil
}

{\footnotesize для автентифікації на основі паролів, коли паролі створюються або оновлюються користувачами, паролі перевіряються на відсутність у списку загальновживаних, очікуваних або скомпрометованих паролів в IA-05(01)(a);}


{\footnotesize\bfseries
No: 5\\
Name: ia\_\allowbreak{}5\_\allowbreak{}1\_\allowbreak{}c\\
Type: string\\
Default: nil
}

{\footnotesize для автентифікації на основі паролів, паролі передаються лише криптографічно захищеними каналами;}


{\footnotesize\bfseries
No: 6\\
Name: ia\_\allowbreak{}5\_\allowbreak{}1\_\allowbreak{}d\\
Type: string\\
Default: nil
}

{\footnotesize для автентифікації на основі паролів паролі зберігаються за допомогою затвердженого алгоритму гешування, переважно використовуючи ключову геш-функцію;}


{\footnotesize\bfseries
No: 7\\
Name: ia\_\allowbreak{}5\_\allowbreak{}1\_\allowbreak{}e\\
Type: string\\
Default: nil
}

{\footnotesize для автентифікації на основі пароля після відновлення облікового запису потрібно негайно вибрати новий пароль;}


{\footnotesize\bfseries
No: 8\\
Name: ia\_\allowbreak{}5\_\allowbreak{}1\_\allowbreak{}f\\
Type: string\\
Default: nil
}

{\footnotesize для автентифікації на основі пароля дозволяється вибір користувачем довгих паролів і фраз, що включають пробіли та всі друковані символи;}


{\footnotesize\bfseries
No: 9\\
Name: ia\_\allowbreak{}5\_\allowbreak{}1\_\allowbreak{}g\\
Type: string\\
Default: nil
}

{\footnotesize для автентифікації на основі пароля використовуються автоматизовані інструменти, які допомагають користувачеві у виборі надійних автентифікаторів паролів;}


{\footnotesize\bfseries
No: 10\\
Name: ia\_\allowbreak{}5\_\allowbreak{}1\_\allowbreak{}h\\
Type: string\\
Default: nil
}

{\footnotesize для автентифікації на основі пароля застосовуються <IA05(01)\_\allowbreak{}ODP[02] склад та правила складності>.}




\subsubsection{АВТЕНТИФІКАЦІЯ НА ОСНОВІ ВІДКРИТОГО КЛЮЧА (IA-5(2))}


\vspace{10pt}
\textit{Немає параметрів для цього контролю.}
\vspace{10pt}


\subsubsection{ОСОБИСТА АБО ДОВІРЧА АВТЕНТИФІКАЦІЯ ЗОВНІШНЬОЇ СТОРОНИ [Виключено: Включено до ІА-12(04)]. (IA-5(3))}


\vspace{10pt}
{\footnotesize\bfseries
No: 1\\
Name: ia\_\allowbreak{}5\_\allowbreak{}3\_\allowbreak{}01\\
Type: list\\
Default: [\textquotedbl{}admin\textquotedbl{}, \textquotedbl{}security\_\allowbreak{}officer\textquotedbl{}]
}

{\footnotesize УПРАВЛІННЯ АВТЕНТИФІКАТОРОМ - ОСОБИСТА АБО ДОВІРЧА АВТЕНТИФІКАЦІЯ ЗОВНІШНЬОЇ СТОРОНИ [Виключено: Включено до ІА-12(04)]}




\subsubsection{АВТОМАТИЗОВАНА ПІДТРИМКА ДЛЯ ВИЗНАЧЕННЯ МІЦНОСТІ ПАРОЛЯ [Виключено: Включено до ІА-05(01)]. (IA-5(4))}


\vspace{10pt}
{\footnotesize\bfseries
No: 1\\
Name: ia\_\allowbreak{}5\_\allowbreak{}4\_\allowbreak{}01\\
Type: string\\
Default: nil
}

{\footnotesize УПРАВЛІННЯ АВТЕНТИФІКАТОРОМ - АВТОМАТИЗОВАНА ПІДТРИМКА ДЛЯ ВИЗНАЧЕННЯ МІЦНОСТІ ПАРОЛЯ [Виключено: Включено до ІА-05(01)]}




\subsubsection{ЗМІНА АВТЕНТИФІКАТОРІВ ДО ДОСТАВКИ (IA-5(5))}


\vspace{10pt}
{\footnotesize\bfseries
No: 1\\
Name: ia\_\allowbreak{}5\_\allowbreak{}5\_\allowbreak{}01\\
Type: string\\
Default: nil
}

{\footnotesize розробники та інсталятори компонентів системи зобов'язані надавати унікальні автентифікатори або змінювати автентифікатори за замовчуванням до доставки та встановлення.}




\subsubsection{ЗАХИСТ АВТЕНТИФІКАТОРІВ (IA-5(6))}


\vspace{10pt}
{\footnotesize\bfseries
No: 1\\
Name: ia\_\allowbreak{}5\_\allowbreak{}6\_\allowbreak{}01\\
Type: string\\
Default: nil
}

{\footnotesize автентифікатори захищені відповідно до категорії безпеки інформації, до якої дозволяє доступ використання автентифікатора.}




\subsubsection{ВІДСУТНІСТЬ ВБУДОВАНИХ НЕЗАШИФРОВАНИХ СТАТИЧНИХ АВТЕНТИФІКАТОРІВ (IA-5(7))}


\vspace{10pt}
{\footnotesize\bfseries
No: 1\\
Name: ia\_\allowbreak{}5\_\allowbreak{}7\_\allowbreak{}01\\
Type: string\\
Default: nil
}

{\footnotesize незашифровані статичні автентифікатори не вбудовуються в застосунки або сценарії доступу та не збережені на функціональній клавіші.}




\subsubsection{БАГАТОСИСТЕМНІ ОБЛІКОВІ ЗАПИСИ (IA-5(8))}


\vspace{10pt}
{\footnotesize\bfseries
No: 1\\
Name: ia\_\allowbreak{}5\_\allowbreak{}8\_\allowbreak{}odp\\
Type: string\\
Default: nil
}

{\footnotesize визначено заходи захисту, впроваджені для управління ризиком компрометації через те, що користувачі мають облікові записи в декількох системах.}


{\footnotesize\bfseries
No: 2\\
Name: ia\_\allowbreak{}5\_\allowbreak{}8\_\allowbreak{}01\\
Type: string\\
Default: nil
}

{\footnotesize <IA-05(08)\_\allowbreak{}ODP заходи захисту> впроваджені для управління ризиком компрометації через наявність облікових записів у декількох системах.}




\subsubsection{УПРАВЛІННЯ ОБ'ЄДНАННЯМ АВТЕНТИФІКАТОРІВ (IA-5(9))}


\vspace{10pt}
{\footnotesize\bfseries
No: 1\\
Name: ia\_\allowbreak{}5\_\allowbreak{}9\_\allowbreak{}odp\\
Type: string\\
Default: nil
}

{\footnotesize визначено зовнішні організації, які будуть використовуватися для об'єднання автентифікаторів;}


{\footnotesize\bfseries
No: 2\\
Name: ia\_\allowbreak{}5\_\allowbreak{}9\_\allowbreak{}01\\
Type: string\\
Default: nil
}

{\footnotesize <IA-05(09)\_\allowbreak{}ODP зовнішні організації> використовуються для об'єднання автентифікаторів.}




\subsubsection{ДИНАМІЧНЕ ЗВ'ЯЗУВАННЯ МАНДАТІВ (IA-5(10))}


\vspace{10pt}
{\footnotesize\bfseries
No: 1\\
Name: ia\_\allowbreak{}5\_\allowbreak{}10\_\allowbreak{}odp\\
Type: string\\
Default: nil
}

{\footnotesize визначено правила для динамічного зв'язування ідентифікаторів та автентифікаторів;}


{\footnotesize\bfseries
No: 2\\
Name: ia\_\allowbreak{}5\_\allowbreak{}10\_\allowbreak{}01\\
Type: string\\
Default: nil
}

{\footnotesize ідентифікатори та автентифікатори динамічно зв'язуються за допомогою <IA-05(10)\_\allowbreak{}ODP правила для динамічного зв'язування >.}




\subsubsection{АВТЕНТИФІКАЦІЯ НА ОСНОВІ АПАРАТНИХ ТОКЕНІВ (IA-5(11)) [Вилучено]}
[Вилучено: Включено до ІА-02(01), ІА-02(02)]

\vspace{10pt}
\textit{Немає параметрів для цього контролю.}
\vspace{10pt}


\subsubsection{ЕФЕКТИВНІСТЬ БІОМЕТРИЧНОЇ АВТЕНТИФІКАЦІЇ (IA-5(12))}


\vspace{10pt}
{\footnotesize\bfseries
No: 1\\
Name: ia\_\allowbreak{}5\_\allowbreak{}12\_\allowbreak{}odp\\
Type: string\\
Default: nil
}

{\footnotesize визначено вимоги до якості біометрії; 273}


{\footnotesize\bfseries
No: 2\\
Name: ia\_\allowbreak{}5\_\allowbreak{}12\_\allowbreak{}01\\
Type: string\\
Default: nil
}

{\footnotesize для біометричної автентифікації використовувати механізми, які задовольняють  <IA-05(12)\_\allowbreak{}ODP вимоги>.}




\subsubsection{ЗАКІНЧЕННЯ ТЕРМІНУ ШУВАННЯ АВТЕНТИФІКАТОРІВ (IA-5(13))}


\vspace{10pt}
{\footnotesize\bfseries
No: 1\\
Name: ia\_\allowbreak{}5\_\allowbreak{}13\_\allowbreak{}odp\\
Type: string\\
Default: nil
}

{\footnotesize визначено періоду часу після якого необхідно заборонити використання кешованих автентифікаторів}


{\footnotesize\bfseries
No: 2\\
Name: ia\_\allowbreak{}5\_\allowbreak{}13\_\allowbreak{}01\\
Type: string\\
Default: nil
}

{\footnotesize забороняється використання кешованих автентифікаторів після <IA-05(13)\_\allowbreak{}ODP періоду часу>.}




\subsubsection{УПРАВЛІННЯ ЗМІСТОМ ДОВІРЧИХ СХОВИЩ ІНФРАСТРУКТУРИ ВІДКРИТИХ КЛЮЧІВ (IA-5(14))}


\vspace{10pt}
{\footnotesize\bfseries
No: 1\\
Name: ia\_\allowbreak{}5\_\allowbreak{}14\_\allowbreak{}01\\
Type: string\\
Default: nil
}

{\footnotesize для автентифікації на основі інфраструктури з відкритим ключем використовується загальноорганізаційна методологія управління 274 вмістом довірених сховищ інфраструктури відкритого ключа, встановлених на всіх платформах, включно з мережами, операційними системами, браузерами та застосунками.}




\subsubsection{ПРОДУКТИ ТА ПОСЛУГИ, ЗАТВЕРДЖЕНІ УПОВНОВАЖЕНИМ ОРГАНОМ (IA-5(15))}


\vspace{10pt}
{\footnotesize\bfseries
No: 1\\
Name: ia\_\allowbreak{}5\_\allowbreak{}15\_\allowbreak{}01\\
Type: string\\
Default: nil
}

{\footnotesize використовуються лише схвалені та затверджені уповноваженим органом продукти та послуги.}




\subsubsection{ПЕРЕДАЧА ОСОБИСТОЇ АБО ДОВІРЧОЇ АВТЕНТИФІКАЦІЇ ЗОВНІШНЬОЇ СТОРОНИ (IA-5(16))}


\vspace{10pt}
{\footnotesize\bfseries
No: 1\\
Name: ia\_\allowbreak{}5\_\allowbreak{}16\_\allowbreak{}odp\_\allowbreak{}1\\
Type: string\\
Default: nil
}

{\footnotesize визначено типи та/або конкретні автентифікатори, які будуть передаватися;}


{\footnotesize\bfseries
No: 2\\
Name: ia\_\allowbreak{}5\_\allowbreak{}16\_\allowbreak{}odp\_\allowbreak{}2\\
Type: string\\
Default: nil
}

{\footnotesize вибрано одне з наступних ЗНАЧЕНЬ ПАРАМЕТРІВ: \{особисто; довіреною зовнішньою стороною\};}


{\footnotesize\bfseries
No: 3\\
Name: ia\_\allowbreak{}5\_\allowbreak{}16\_\allowbreak{}odp\_\allowbreak{}3\\
Type: string\\
Default: nil
}

{\footnotesize визначено зареєстрований орган, який приймає автентифікатори;}


{\footnotesize\bfseries
No: 4\\
Name: ia\_\allowbreak{}5\_\allowbreak{}16\_\allowbreak{}odp\_\allowbreak{}4\\
Type: string\\
Default: nil
}

{\footnotesize визначено персонал або ролі, які уповноважують передачу автентифікаторів;}


{\footnotesize\bfseries
No: 5\\
Name: ia\_\allowbreak{}5\_\allowbreak{}16\_\allowbreak{}01\\
Type: string\\
Default: nil
}

{\footnotesize передача <IA-05(16)\_\allowbreak{}ODP[01] типів та/або конкретних}




\subsubsection{АВТОМАТИЗОВАНІ ЗАСОБИ ВИЯВЛЕННЯ АТАК ІЗ ВИКОРИСТАННЯМ БІОМЕТРИЧНИХ АВТЕНТИФІКАТОРІВ (IA-5(17))}


\vspace{10pt}
{\footnotesize\bfseries
No: 1\\
Name: ia\_\allowbreak{}5\_\allowbreak{}17\_\allowbreak{}01\\
Type: string\\
Default: nil
}

{\footnotesize використовуються механізми виявлення атак із використанням штучно виготовлених артефактів для біометричних автентифікаторів.}




\subsubsection{МЕНЕДЖЕР ПАРОЛІВ (IA-5(18))}


\vspace{10pt}
{\footnotesize\bfseries
No: 1\\
Name: ia\_\allowbreak{}5\_\allowbreak{}18\_\allowbreak{}odp\_\allowbreak{}1\\
Type: string\\
Default: nil
}

{\footnotesize визначено менеджери паролів, які використовуються для створення та керування паролями;}


{\footnotesize\bfseries
No: 2\\
Name: ia\_\allowbreak{}5\_\allowbreak{}18\_\allowbreak{}odp\_\allowbreak{}2\\
Type: string\\
Default: nil
}

{\footnotesize визначено елементи керування для захисту паролів;}


{\footnotesize\bfseries
No: 3\\
Name: ia\_\allowbreak{}5\_\allowbreak{}18\_\allowbreak{}a\\
Type: string\\
Default: nil
}

{\footnotesize <IA-05(18)\_\allowbreak{}ODP[01] менеджери паролів> використовуються для створення та керування паролями;}


{\footnotesize\bfseries
No: 4\\
Name: ia\_\allowbreak{}5\_\allowbreak{}18\_\allowbreak{}b\\
Type: string\\
Default: nil
}

{\footnotesize паролі зах ищені за допомогою  <IA-05(18)\_\allowbreak{}ODP[02] елементи керування >}




\subsection{ПРИХОВУВАННЯ ЗВОРОТНОГО ЗВ'ЯЗКУ АВТЕНТИФІКАТОРА}


\vspace{10pt}
{\footnotesize\bfseries
No: 1\\
Name: ia\_\allowbreak{}6\_\allowbreak{}01\\
Type: string\\
Default: nil
}

{\footnotesize забезпечино приховану зворотну передачу інформації автентифікації в про277 цесі автентифікації для забезпечення захисту інформації від можливої експлуатації та використання неавторизованими особами.}




\subsection{АВТЕНТИФІКАЦІЯ КРИПТОГРАФІЧНОГО МОДУЛЯ (IA-7)}


\vspace{10pt}
{\footnotesize\bfseries
No: 1\\
Name: ia\_\allowbreak{}7\_\allowbreak{}01\\
Type: string\\
Default: nil
}

{\footnotesize впроваджено механізми автентифікації в криптографічний модуль, який відповідає вимогам чинних законів, виконавчих розпоряджень, директив, політик, правил, стандартів та рекомендацій для такої автентифікації.}




\subsection{ІДЕНТИФІКАЦІЯ ТА АВТЕНТИФІКАЦІЯ (КОРИСТУВАЧІ, ЩО НЕ НАЛЕЖАТЬ ДО ОРГАНІЗАЦІЇ) (IA-8)}


\vspace{10pt}
{\footnotesize\bfseries
No: 1\\
Name: ia\_\allowbreak{}8\_\allowbreak{}01\\
Type: string\\
Default: nil
}

{\footnotesize унікально ідентифіку ються та автентифікуються користувачі, що не належать до організації або процеси (що не належать організації), які діють від імені користувачів.}




\subsubsection{ВИКОРИСТАННЯ ЗАТВЕРДЖЕНИХ ПРОДУКТІВ (IA-8(3)) [Вилучено]}
[Вилучено: Включено до ІА-08(02)]

\vspace{10pt}
\textit{Немає параметрів для цього контролю.}
\vspace{10pt}


\subsubsection{ВИЗНАННЯ ПОСВІДЧЕНЬ ОСОБИ (PIV-I) (IA-8(5))}


\vspace{10pt}
{\footnotesize\bfseries
No: 1\\
Name: ia\_\allowbreak{}8\_\allowbreak{}5\_\allowbreak{}odp\\
Type: string\\
Default: nil
}

{\footnotesize визначено політику використання ативних облікових даних або облікових даних PKI;}


{\footnotesize\bfseries
No: 2\\
Name: ia\_\allowbreak{}8\_\allowbreak{}5\_\allowbreak{}1\\
Type: string\\
Default: nil
}

{\footnotesize приймаються облікові дані  або дані PKI, які відповідають <IA-08(05)\_\allowbreak{}ODP політика>;}


{\footnotesize\bfseries
No: 3\\
Name: ia\_\allowbreak{}8\_\allowbreak{}5\_\allowbreak{}2\\
Type: string\\
Default: nil
}

{\footnotesize підтверджуються облікові дані  або дані PKI, які відповідають <IA-08(05)\_\allowbreak{}ODP політика>.}




\subsubsection{РОЗМЕЖУВАННЯ (IA-8(6))}


\vspace{10pt}
{\footnotesize\bfseries
No: 1\\
Name: ia\_\allowbreak{}8\_\allowbreak{}6\_\allowbreak{}odp\\
Type: string\\
Default: nil
}

{\footnotesize визначено заходи, щоб розмежувати атрибути користувача або зв'язки підтвердження ідентифікатора між окремими особами, постачальниками облікових даних і довіреними 281 сторонами;}


{\footnotesize\bfseries
No: 2\\
Name: ia\_\allowbreak{}8\_\allowbreak{}6\_\allowbreak{}01\\
Type: string\\
Default: nil
}

{\footnotesize впроваджено <IA-08(06)\_\allowbreak{}ODP заходи> щоб розмежувати атрибути користува ча або зв'язки підтвердження ідентифікатора між окремими особами, постачальниками облікових даних і довіреними сторонами.}




\subsection{ПОСЛУГИ ІДЕНТИФІКАЦІЇ ТА АВТЕНТИФІКАЦІЇ (IA-9)}


\vspace{10pt}
{\footnotesize\bfseries
No: 1\\
Name: ia\_\allowbreak{}9\_\allowbreak{}odp\\
Type: string\\
Default: nil
}

{\footnotesize визначено системні служби та застосунки, які мають бути унікально ідентифіковані та автентифіковані;}


{\footnotesize\bfseries
No: 2\\
Name: ia\_\allowbreak{}9\_\allowbreak{}01\\
Type: string\\
Default: nil
}

{\footnotesize <IA-09\_\allowbreak{}ODP системні служби та застосунки> унікально ідентифікуються та автентифікуються перед встановленням зв'язку з пристроями, користувачами або іншими службами чи застосунками.}




\subsubsection{ОБМІН ІНФО- (IA-9(1))}


\vspace{10pt}
{\footnotesize\bfseries
No: 1\\
Name: ia\_\allowbreak{}9\_\allowbreak{}1\_\allowbreak{}01\\
Type: string\\
Default: nil
}

{\footnotesize ПОСЛУГИ ІДЕНТИФІКАЦІЇ ТА АВТЕНТИФІКАЦІЇ - ОБМІН ІНФО- РМАЦІЄЮ [Виключено: перенесено до IA-09]}




\subsubsection{ПЕРЕДАЧА РІШЕНЬ [Виключено: перенесено до IA-09] (IA-9(2))}


\vspace{10pt}
{\footnotesize\bfseries
No: 1\\
Name: ia\_\allowbreak{}9\_\allowbreak{}2\_\allowbreak{}01\\
Type: string\\
Default: nil
}

{\footnotesize ПОСЛУГИ ІДЕНТИФІКАЦІЇ ТА АВТЕНТИФІКАЦІЇ - ПЕРЕДАЧА РІШЕНЬ [Виключено: перенесено до IA-09]}




\subsection{АДАПТИВНА АВТЕНТИФІКАЦІЯ (IA-10)}


\vspace{10pt}
{\footnotesize\bfseries
No: 1\\
Name: ia\_\allowbreak{}10\_\allowbreak{}odp\_\allowbreak{}1\\
Type: string\\
Default: nil
}

{\footnotesize визначені додаткові методи або механізми автентифікації, які будуть застосовуватися при доступі до системи за певних обставин або ситуацій;}


{\footnotesize\bfseries
No: 2\\
Name: ia\_\allowbreak{}10\_\allowbreak{}odp\_\allowbreak{}1\\
Type: string\\
Default: nil
}

{\footnotesize визначені обставини або ситуації, які вимагають від осіб, що отримують доступ до системи, використання додаткових методів або механізмів автентифікації;}


{\footnotesize\bfseries
No: 3\\
Name: ia\_\allowbreak{}10\_\allowbreak{}01\\
Type: string\\
Default: nil
}

{\footnotesize особи, які отримують доступ до системи, повинні використовувати ситуацій>.}




\subsection{ПОВТОРНА АВТЕНТИФІКАЦІЯ (IA-11)}


\vspace{10pt}
{\footnotesize\bfseries
No: 1\\
Name: ia\_\allowbreak{}11\_\allowbreak{}odp\\
Type: string\\
Default: nil
}

{\footnotesize визначено обставини або ситуації, що вимагають повторної автен283 тифікації;}


{\footnotesize\bfseries
No: 2\\
Name: ia\_\allowbreak{}11\_\allowbreak{}01\\
Type: string\\
Default: nil
}

{\footnotesize користувачі повинні повторно автентифікуватися, коли  <IA-11\_\allowbreak{}ODP обставини або ситуації>.}




\subsection{ПЕРЕВІРКА СПРАВЖНОСТІ (ІДЕНТИЧНОСТІ) (IA-12)}


\vspace{10pt}
{\footnotesize\bfseries
No: 1\\
Name: ia\_\allowbreak{}12\_\allowbreak{}a\\
Type: string\\
Default: nil
}

{\footnotesize користувачі, яким потрібні облікові записи для логічного доступу до систем на основі вимог гарантій відповідного рівня, як це зазначено у відповідних стандартах і рекомендаціях, мають підтверджену ідентичність;}


{\footnotesize\bfseries
No: 2\\
Name: ia\_\allowbreak{}12\_\allowbreak{}b\\
Type: string\\
Default: nil
}

{\footnotesize встановлені ідентифікатори користувачів унікальні для особи;}


{\footnotesize\bfseries
No: 3\\
Name: ia\_\allowbreak{}12\_\allowbreak{}c\_\allowbreak{}1\\
Type: string\\
Default: nil
}

{\footnotesize збираються докази ідентичності особи;}


{\footnotesize\bfseries
No: 4\\
Name: ia\_\allowbreak{}12\_\allowbreak{}c\_\allowbreak{}2\\
Type: string\\
Default: nil
}

{\footnotesize затверджуються докази ідентичності особи;}


{\footnotesize\bfseries
No: 5\\
Name: ia\_\allowbreak{}12\_\allowbreak{}c\_\allowbreak{}3\\
Type: string\\
Default: nil
}

{\footnotesize перевіряються докази ідентичності особи;}




\subsubsection{АВТОРИЗАЦІЯ СУПЕРВАЙЗЕРА (IA-12(1))}


\vspace{10pt}
{\footnotesize\bfseries
No: 1\\
Name: ia\_\allowbreak{}12\_\allowbreak{}1\_\allowbreak{}01\\
Type: string\\
Default: nil
}

{\footnotesize процес реєстрації для отримання облікового запису для логічного доступу включає авторизацію супервайзера.}




\subsubsection{ПОСВІДЧЕННЯ ОСОБИ (IA-12(2))}


\vspace{10pt}
{\footnotesize\bfseries
No: 1\\
Name: ia\_\allowbreak{}12\_\allowbreak{}2\_\allowbreak{}01\\
Type: string\\
Default: nil
}

{\footnotesize документи, що посвідчують особу пред'являються до реєстраційного органу.}




\subsubsection{ОЧНА ПЕРЕВІРКА ТА ВЕРИФІКАЦІЯ (IA-12(4))}


\vspace{10pt}
{\footnotesize\bfseries
No: 1\\
Name: ia\_\allowbreak{}12\_\allowbreak{}4\_\allowbreak{}01\\
Type: string\\
Default: nil
}

{\footnotesize підтвердження та перевірка посвідчення особи проводиться особисто в призначеному органі реєстрації.}




\subsubsection{ПІДТВЕРДЖЕННЯ АДРЕСИ (IA-12(5))}


\vspace{10pt}
{\footnotesize\bfseries
No: 1\\
Name: ia\_\allowbreak{}12\_\allowbreak{}5\_\allowbreak{}odp\\
Type: string\\
Default: nil
}

{\footnotesize вибрано одне з наступних ЗНАЧЕНЬ ПАРАМЕТРІВ: \{реєстраційний код; повідомлення про перевірку\}; доставляється через зовнішні канали для перевірки адреси (фізичної або цифрової) користувача.}




\subsubsection{ПРИЙНЯТТЯ ІДЕНТИФІКАЦІЙ СХВАЛЕНИХ ТРЕТЬОЮ СТОРОНОЮ (IA-12(6))}


\vspace{10pt}
{\footnotesize\bfseries
No: 1\\
Name: ia\_\allowbreak{}12\_\allowbreak{}6\_\allowbreak{}odp\\
Type: string\\
Default: nil
}

{\footnotesize визначено рівень гарантії ідентичності для прийняття зовнішньо підтверджених ідентифікаторів;}


{\footnotesize\bfseries
No: 2\\
Name: ia\_\allowbreak{}12\_\allowbreak{}6\_\allowbreak{}01\\
Type: string\\
Default: nil
}

{\footnotesize приймаються зовнішньо підтверджені ідентифікатори <IA12(06)\_\allowbreak{}ODP рівень гарантії ідентичності>.}




\section{IR}
Клас заходів захисту IR --- РЕАГУВАННЯ НА ІНЦИДЕНТИ

Цей клас визначає процедури виявлення, аналізу, локалізації та усунення наслідків інцидентів інформаційної безпеки.

\textbf{Перелік заходів захисту:} Політика та процедури реагування на інциденти (IR-1); Навчання з реагування на інциденти (IR-2); Моделювання подій (IR-2(1)); Злам (IR-2(3)); Перевірка реагувань на інциденти (IR-3); Координація з пов'язаними планами (IR-3(2)); Постійне покращення (IR-3(3)); Обробка інциденту (IR-4); Динамічна реконфігурація (IR-4(2)); Безперервність операцій (IR-4(3)); Інформаційна кореляція (IR-4(4)); Внутрішні загрози - особливі можливості (IR-4(6)); Координація з зовнішніми організаціями (IR-4(8)); Здатність динамічного реагування (IR-4(9)); Координація ланцюга постачання (IR-4(10)); Інтегрована група реагування на інцеденти (IR-4(11)); Зловмисний код та криміналістичний аналіз (IR-4(12)); Аналіз поведінки (IR-4(13)); Центр безпеки (IR-4(14)); Зв'язки з громадкістю та відновлення репутації (IR-4(15)); Моніторинг інциденту (IR-5); Автоматизоване відстеження, збір даних і аналіз (IR-5(1)); Звітність про інциденти (IR-6); Автоматичне звітування (IR-6(1)); Координація ланцюжка постачання (IR-6(3)); Підтримка реагування на інциденти (IR-7); Автоматизація підтримки для забезпечення доступності інформації та підтримки (IR-7(1)); Координація з зовнішніми постачальниками (IR-7(2)); План реагування на інциденти (IR-8); Обробка персональних даних (IR-8(1)); Реагування на витік інформації (IR-9); Відповідальний персонал (IR-9(1)) [Вилучено]; Тренування (IR-9(2)); Робота після витоку (IR-9(3)); Викриття неавторизованого персоналу (IR-9(4)); Інтегрована команда аналізу інформаційної безпеки (IR-10) [Вилучено].

\subsection{ПОЛІТИКА ТА ПРОЦЕДУРИ РЕАГУВАННЯ НА ІНЦИДЕНТИ (IR-1)}


\vspace{10pt}
{\footnotesize\bfseries
No: 1\\
Name: ir\_\allowbreak{}1\_\allowbreak{}odp\_\allowbreak{}1\\
Type: string\\
Default: nil
}

{\footnotesize визначено персонал або ролі, до яких має бути доведена політика реагування на інциденти;}


{\footnotesize\bfseries
No: 2\\
Name: ir\_\allowbreak{}1\_\allowbreak{}odp\_\allowbreak{}2\\
Type: string\\
Default: nil
}

{\footnotesize визначено персонал або ролі, до яких мають бути доведені процедури реагування на інциденти;}


{\footnotesize\bfseries
No: 3\\
Name: ir\_\allowbreak{}1\_\allowbreak{}odp\_\allowbreak{}3\\
Type: string\\
Default: nil
}

{\footnotesize вибрано одне або декілька з наступних ЗНАЧЕНЬ ПАРА- МЕТРІВ: \{рівень організації; рівень місії/бізнес -процесу; рівень системи\};}


{\footnotesize\bfseries
No: 4\\
Name: ir\_\allowbreak{}1\_\allowbreak{}odp\_\allowbreak{}4\\
Type: string\\
Default: nil
}

{\footnotesize визначено посадову особу, яка керуватиме політикою та процедурами реагування на інциденти;}


{\footnotesize\bfseries
No: 5\\
Name: ir\_\allowbreak{}1\_\allowbreak{}odp\_\allowbreak{}5\\
Type: string\\
Default: nil
}

{\footnotesize визначено частоту, з якою переглядається та оновлюється поточна політика реагування на інциденти;}


{\footnotesize\bfseries
No: 6\\
Name: ir\_\allowbreak{}1\_\allowbreak{}odp\_\allowbreak{}6\\
Type: string\\
Default: nil
}

{\footnotesize визначаються події, які потребують перегляду та оновлення поточної політики реагування на інциденти;}


{\footnotesize\bfseries
No: 7\\
Name: ir\_\allowbreak{}1\_\allowbreak{}odp\_\allowbreak{}7\\
Type: string\\
Default: nil
}

{\footnotesize визначено частоту, з якою переглядаються та оновлюються поточні процедури реагування на інциденти;}


{\footnotesize\bfseries
No: 8\\
Name: ir\_\allowbreak{}1\_\allowbreak{}odp\_\allowbreak{}8\\
Type: string\\
Default: nil
}

{\footnotesize визначено події, які пот ребують перегляду та оновлення процедур реагування на інциденти;}


{\footnotesize\bfseries
No: 9\\
Name: ir\_\allowbreak{}1\_\allowbreak{}a\_\allowbreak{}1\\
Type: string\\
Default: nil
}

{\footnotesize розроблено та задокументовано політику реагування на інциденти;}


{\footnotesize\bfseries
No: 10\\
Name: ir\_\allowbreak{}1\_\allowbreak{}a\_\allowbreak{}2\\
Type: string\\
Default: nil
}

{\footnotesize політика реагування на інциденти поширюється серед  <IR01\_\allowbreak{}ODP[01] персоналу або ролей>;}


{\footnotesize\bfseries
No: 11\\
Name: ir\_\allowbreak{}1\_\allowbreak{}a\_\allowbreak{}3\\
Type: string\\
Default: nil
}

{\footnotesize розроблені та задокументовані процедури реагування на інциденти, що сприяють впровадженню політики реагування на інциденти та пов'язаних з нею заходів захисту з реагування на інциденти;}


{\footnotesize\bfseries
No: 12\\
Name: ir\_\allowbreak{}1\_\allowbreak{}a\_\allowbreak{}4\\
Type: string\\
Default: nil
}

{\footnotesize процедури реагування на інциденти поширюються серед <IR01\_\allowbreak{}ODP[02] персоналу або ролей>;}


{\footnotesize\bfseries
No: 13\\
Name: ir\_\allowbreak{}1\_\allowbreak{}a\_\allowbreak{}1\_\allowbreak{}a\_\allowbreak{}1\\
Type: string\\
Default: nil
}

{\footnotesize політика реагування на інциденти  <IR-01\_\allowbreak{}ODP[03] ВИБРАНЕ ЗНАЧЕННЯ(Я) ПАРАМЕТРА(ів)> містить мету;}


{\footnotesize\bfseries
No: 14\\
Name: ir\_\allowbreak{}1\_\allowbreak{}a\_\allowbreak{}1\_\allowbreak{}a\_\allowbreak{}2\\
Type: string\\
Default: nil
}

{\footnotesize політика реагування на інциденти <IR-01\_\allowbreak{}ODP[03] ВИБРАНЕ ЗНАЧЕННЯ(Я) ПАРАМЕТРА(ів)> містить сферу застосування; 288}


{\footnotesize\bfseries
No: 15\\
Name: ir\_\allowbreak{}1\_\allowbreak{}a\_\allowbreak{}1\_\allowbreak{}a\_\allowbreak{}3\\
Type: string\\
Default: nil
}

{\footnotesize політика реагування на інциденти  <IR-01\_\allowbreak{}ODP[03] ВИБРАНЕ ЗНАЧЕННЯ(Я) ПАРАМЕТРА(ів)> містить ролі;}


{\footnotesize\bfseries
No: 16\\
Name: ir\_\allowbreak{}1\_\allowbreak{}a\_\allowbreak{}1\_\allowbreak{}a\_\allowbreak{}4\\
Type: string\\
Default: nil
}

{\footnotesize політика реагування на інциденти  <IR-01\_\allowbreak{}ODP[03] ВИБРАНЕ ЗНАЧЕННЯ(Я) ПАРАМЕТРА(ів)> містить обов'язки;}


{\footnotesize\bfseries
No: 17\\
Name: ir\_\allowbreak{}1\_\allowbreak{}a\_\allowbreak{}1\_\allowbreak{}a\_\allowbreak{}5\\
Type: string\\
Default: nil
}

{\footnotesize політика реагування на інциденти  <IR-01\_\allowbreak{}ODP[03] ВИБРАНЕ ЗНАЧЕННЯ(Я) ПАРАМЕТРА(ів)> містить відповідальність керівництва;}


{\footnotesize\bfseries
No: 18\\
Name: ir\_\allowbreak{}1\_\allowbreak{}a\_\allowbreak{}1\_\allowbreak{}a\_\allowbreak{}6\\
Type: string\\
Default: nil
}

{\footnotesize політика реагування на інциденти  <IR-01\_\allowbreak{}ODP[03] ВИБРАНЕ ЗНАЧЕННЯ(Я) ПАРАМЕТРА(ів)> містить координацію між підрозділами організації;}


{\footnotesize\bfseries
No: 19\\
Name: ir\_\allowbreak{}1\_\allowbreak{}a\_\allowbreak{}1\_\allowbreak{}a\_\allowbreak{}7\\
Type: string\\
Default: nil
}

{\footnotesize політика реагування на інциденти  <IR-01\_\allowbreak{}ODP[03] ВИБРАНЕ ЗНАЧЕННЯ(Я) ПАРАМЕТРА(ів)> містить систему контролю відповідності;}


{\footnotesize\bfseries
No: 20\\
Name: ir\_\allowbreak{}1\_\allowbreak{}a\_\allowbreak{}1\_\allowbreak{}b\\
Type: string\\
Default: nil
}

{\footnotesize політика реагування на інциденти  <IR-01\_\allowbreak{}ODP[03] ВИБРАНЕ ЗНАЧЕННЯ ПАРАМЕТРА(ів)> відповідає чинному законодавству, виконавчим наказам, директивам, нормам, політикам, стандартам та керівним принципам;}


{\footnotesize\bfseries
No: 21\\
Name: ir\_\allowbreak{}1\_\allowbreak{}b\\
Type: string\\
Default: nil
}

{\footnotesize <IR-01\_\allowbreak{}ODP[04] посадова особа> призначається для управління розробкою, документуванням та розповсюдженням політики та процедур реагування на інциденти;}


{\footnotesize\bfseries
No: 22\\
Name: ir\_\allowbreak{}1\_\allowbreak{}c\_\allowbreak{}1\_\allowbreak{}1\\
Type: string\\
Default: nil
}

{\footnotesize переглядається та оновлюється поточна політика реагування на}


{\footnotesize\bfseries
No: 23\\
Name: ir\_\allowbreak{}1\_\allowbreak{}c\_\allowbreak{}1\_\allowbreak{}2\\
Type: string\\
Default: nil
}

{\footnotesize поточна політика реагування на інциденти переглядається та}


{\footnotesize\bfseries
No: 24\\
Name: ir\_\allowbreak{}1\_\allowbreak{}c\_\allowbreak{}2\_\allowbreak{}1\\
Type: string\\
Default: nil
}

{\footnotesize переглядаються та оновлюються поточні процедури реагування на}


{\footnotesize\bfseries
No: 25\\
Name: ir\_\allowbreak{}1\_\allowbreak{}c\_\allowbreak{}2\_\allowbreak{}2\\
Type: string\\
Default: nil
}

{\footnotesize поточні процедури реагування на інциденти переглядаються та}




\subsection{НАВЧАННЯ З РЕАГУВАННЯ НА ІНЦИДЕНТИ (IR-2)}


\vspace{10pt}
{\footnotesize\bfseries
No: 1\\
Name: ir\_\allowbreak{}2\_\allowbreak{}odp\_\allowbreak{}1\\
Type: string\\
Default: nil
}

{\footnotesize визначено період часу, протягом якого має бути проведено навчання з реагування на інциденти для користувачів системи, які беруть на себе роль або відповідальність за реагування на інциденти;}


{\footnotesize\bfseries
No: 2\\
Name: ir\_\allowbreak{}2\_\allowbreak{}odp\_\allowbreak{}2\\
Type: string\\
Default: nil
}

{\footnotesize визначено частоту, з якою користувачі повинні п роходити навчання з реагування на інциденти;}


{\footnotesize\bfseries
No: 3\\
Name: ir\_\allowbreak{}2\_\allowbreak{}odp\_\allowbreak{}3\\
Type: string\\
Default: nil
}

{\footnotesize визначено частоту перегляду та оновлення змісту навчання з реагування на інциденти;}


{\footnotesize\bfseries
No: 4\\
Name: ir\_\allowbreak{}2\_\allowbreak{}odp\_\allowbreak{}4\\
Type: string\\
Default: nil
}

{\footnotesize визначено події, які ініціюють перегляд змісту навчання з реагування на інциденти;}


{\footnotesize\bfseries
No: 5\\
Name: ir\_\allowbreak{}2\_\allowbreak{}a\_\allowbreak{}1\\
Type: string\\
Default: nil
}

{\footnotesize навчання з реагування на інциденти надається користувачам системи відповідно до призначених ролей та обов'язків протягом або обов'язків з реагування на інциденти або отримання доступу до системи;}


{\footnotesize\bfseries
No: 6\\
Name: ir\_\allowbreak{}2\_\allowbreak{}a\_\allowbreak{}2\\
Type: string\\
Default: nil
}

{\footnotesize навчання з реагування на інциденти надається користувачам системи відповідно до призначених ролей та обов'язків, коли цього вимагають зміни в системі;}


{\footnotesize\bfseries
No: 7\\
Name: ir\_\allowbreak{}2\_\allowbreak{}a\_\allowbreak{}3\\
Type: string\\
Default: nil
}

{\footnotesize користувачам системи надається навчання з реагування на інциденти відповідно до призначених ролей та обов'язків  <IR02\_\allowbreak{}ODP[02] частота>;}


{\footnotesize\bfseries
No: 8\\
Name: ir\_\allowbreak{}2\_\allowbreak{}b\_\allowbreak{}1\\
Type: string\\
Default: nil
}

{\footnotesize зміст навчання з реагування на інциденти переглядається та}


{\footnotesize\bfseries
No: 9\\
Name: ir\_\allowbreak{}2\_\allowbreak{}b\_\allowbreak{}2\\
Type: string\\
Default: nil
}

{\footnotesize зміст навчання з реагування на інциденти переглядається та}




\subsubsection{МОДЕЛЮВАННЯ ПОДІЙ (IR-2(1))}


\vspace{10pt}
{\footnotesize\bfseries
No: 1\\
Name: ir\_\allowbreak{}2\_\allowbreak{}1\_\allowbreak{}01\\
Type: string\\
Default: nil
}

{\footnotesize моделювання подій включається в процес навчання з реагування на інциденти для забезпечення ефективного реагування персоналу в кризових ситуаціях.}




\subsubsection{ЗЛАМ (IR-2(3))}


\vspace{10pt}
{\footnotesize\bfseries
No: 1\\
Name: ir\_\allowbreak{}2\_\allowbreak{}3\_\allowbreak{}1\\
Type: string\\
Default: nil
}

{\footnotesize проводиться навчання з реагування на інциденти щодо виявлення та реагування на порушення; 291}


{\footnotesize\bfseries
No: 2\\
Name: ir\_\allowbreak{}2\_\allowbreak{}3\_\allowbreak{}2\\
Type: string\\
Default: nil
}

{\footnotesize проводиться навчання з реагування на інциденти щодо процесу повідомлення про порушення в організації.}




\subsection{ПЕРЕВІРКА РЕАГУВАНЬ НА ІНЦИДЕНТИ (IR-3)}


\vspace{10pt}
{\footnotesize\bfseries
No: 1\\
Name: ir\_\allowbreak{}3\_\allowbreak{}odp\_\allowbreak{}1\\
Type: string\\
Default: nil
}

{\footnotesize визначено частоту, з якою необхідно перевіряти ефективність реагування системи на інциденти;}


{\footnotesize\bfseries
No: 2\\
Name: ir\_\allowbreak{}3\_\allowbreak{}odp\_\allowbreak{}2\\
Type: string\\
Default: nil
}

{\footnotesize визначено тести, що використовуються для перевірки ефективності реагування на інциденти в системі;}


{\footnotesize\bfseries
No: 3\\
Name: ir\_\allowbreak{}3\_\allowbreak{}01\\
Type: string\\
Default: nil
}

{\footnotesize ефективність реагування системи на інциденти перевіряється тестів>.}




\subsubsection{КООРДИНАЦІЯ З ПОВ'ЯЗАНИМИ ПЛАНАМИ (IR-3(2))}


\vspace{10pt}
{\footnotesize\bfseries
No: 1\\
Name: ir\_\allowbreak{}3\_\allowbreak{}2\_\allowbreak{}01\\
Type: string\\
Default: nil
}

{\footnotesize Тестування реагування на інциденти координується з елементами організації, відповідальними за пов'язані плани}




\subsubsection{ПОСТІЙНЕ ПОКРАЩЕННЯ (IR-3(3))}


\vspace{10pt}
{\footnotesize\bfseries
No: 1\\
Name: ir\_\allowbreak{}3\_\allowbreak{}3\_\allowbreak{}a\_\allowbreak{}1\\
Type: string\\
Default: nil
}

{\footnotesize якісні дані тестування використовуються для визначення ефективності процесів реагування на інциденти; 293}


{\footnotesize\bfseries
No: 2\\
Name: ir\_\allowbreak{}3\_\allowbreak{}3\_\allowbreak{}a\_\allowbreak{}2\\
Type: string\\
Default: nil
}

{\footnotesize кількісні дані тестування використовуються для визначення ефективності процесів реагування на інциденти;}


{\footnotesize\bfseries
No: 3\\
Name: ir\_\allowbreak{}3\_\allowbreak{}3\_\allowbreak{}b\_\allowbreak{}1\\
Type: string\\
Default: nil
}

{\footnotesize якісні дані тестування використовуються для постійного вдосконалення процесів реагування на інциденти;}


{\footnotesize\bfseries
No: 4\\
Name: ir\_\allowbreak{}3\_\allowbreak{}3\_\allowbreak{}b\_\allowbreak{}2\\
Type: string\\
Default: nil
}

{\footnotesize кількісні дані тестування використовуються для постійного вдосконалення процесів реагування на інциденти;}


{\footnotesize\bfseries
No: 5\\
Name: ir\_\allowbreak{}3\_\allowbreak{}3\_\allowbreak{}c\_\allowbreak{}1\\
Type: string\\
Default: nil
}

{\footnotesize якісні дані, отримані за результатами тестування , використовуються для забезпечення точних показників та метрик реагування на інциденти;}


{\footnotesize\bfseries
No: 6\\
Name: ir\_\allowbreak{}3\_\allowbreak{}3\_\allowbreak{}c\_\allowbreak{}2\\
Type: string\\
Default: nil
}

{\footnotesize кількісні дані, отримані за результатами тестування , використовуються для забезпечення точних показників та метрик реагування на інциденти;}


{\footnotesize\bfseries
No: 7\\
Name: ir\_\allowbreak{}3\_\allowbreak{}3\_\allowbreak{}c\_\allowbreak{}3\\
Type: string\\
Default: nil
}

{\footnotesize якісні дані, отримані за результатами тестування , використовуються для забезпе чення послідовності показників та метрик реагування на інциденти;}


{\footnotesize\bfseries
No: 8\\
Name: ir\_\allowbreak{}3\_\allowbreak{}3\_\allowbreak{}c\_\allowbreak{}4\\
Type: string\\
Default: nil
}

{\footnotesize кількісні дані, отримані за результатами тестування , використовуються для забезпечення послідовності показників та метрик реагування на інциденти;}


{\footnotesize\bfseries
No: 9\\
Name: ir\_\allowbreak{}3\_\allowbreak{}3\_\allowbreak{}c\_\allowbreak{}5\\
Type: string\\
Default: nil
}

{\footnotesize якісні дані, отримані за результатами тестування , використовуються для забезпечення відтворюваносні показників та метрик реагування на інциденти;}


{\footnotesize\bfseries
No: 10\\
Name: ir\_\allowbreak{}3\_\allowbreak{}3\_\allowbreak{}c\_\allowbreak{}6\\
Type: string\\
Default: nil
}

{\footnotesize кількусні дані, отримані за результатами тестування , використовуються для забезпечення ві дтворюваносні показників та метрик реагування на інциденти;}




\subsection{ОБРОБКА ІНЦИДЕНТУ (IR-4)}


\vspace{10pt}
{\footnotesize\bfseries
No: 1\\
Name: ir\_\allowbreak{}4\_\allowbreak{}a\_\allowbreak{}1\\
Type: string\\
Default: nil
}

{\footnotesize впроваджено можливість обробки інцидентів безпеки включно з підготовкою;}


{\footnotesize\bfseries
No: 2\\
Name: ir\_\allowbreak{}4\_\allowbreak{}a\_\allowbreak{}2\\
Type: string\\
Default: nil
}

{\footnotesize впроваджено можливість обробки інцидентів безпеки включно з виявленням;}


{\footnotesize\bfseries
No: 3\\
Name: ir\_\allowbreak{}4\_\allowbreak{}a\_\allowbreak{}3\\
Type: string\\
Default: nil
}

{\footnotesize впроваджено можливість обробки інцидентів безпеки включно з аналізом;}


{\footnotesize\bfseries
No: 4\\
Name: ir\_\allowbreak{}4\_\allowbreak{}a\_\allowbreak{}4\\
Type: string\\
Default: nil
}

{\footnotesize впроваджено можливість обробки інцидентів безпеки включно з локалізацією;}


{\footnotesize\bfseries
No: 5\\
Name: ir\_\allowbreak{}4\_\allowbreak{}a\_\allowbreak{}5\\
Type: string\\
Default: nil
}

{\footnotesize впроваджено можливість обробки інцидентів безпеки включно з ліквідацією;}


{\footnotesize\bfseries
No: 6\\
Name: ir\_\allowbreak{}4\_\allowbreak{}a\_\allowbreak{}6\\
Type: string\\
Default: nil
}

{\footnotesize впроваджено можливість обробки інцидентів безпеки включно з відновленням;}


{\footnotesize\bfseries
No: 7\\
Name: ir\_\allowbreak{}4\_\allowbreak{}b\\
Type: string\\
Default: nil
}

{\footnotesize діяльність з обробки інцидентів координується із заходами із забезпечення безперервності функціонування;}


{\footnotesize\bfseries
No: 8\\
Name: ir\_\allowbreak{}4\_\allowbreak{}c\_\allowbreak{}1\\
Type: string\\
Default: nil
}

{\footnotesize уроки, отримані з поточних дій з обробки інцидентів, включаються в процедури реагування на інциденти, навчання та тестування;}


{\footnotesize\bfseries
No: 9\\
Name: ir\_\allowbreak{}4\_\allowbreak{}c\_\allowbreak{}2\\
Type: string\\
Default: nil
}

{\footnotesize зміни, що випливають з отриманих уроків, впроваджуються відповідним чином;}


{\footnotesize\bfseries
No: 10\\
Name: ir\_\allowbreak{}4\_\allowbreak{}d\_\allowbreak{}1\\
Type: string\\
Default: nil
}

{\footnotesize строгість заходів з обробки інцидентів є порівнянною та передбачуваною в межах всієї організації;}


{\footnotesize\bfseries
No: 11\\
Name: ir\_\allowbreak{}4\_\allowbreak{}d\_\allowbreak{}2\\
Type: string\\
Default: nil
}

{\footnotesize інтенсивність заходів з обробки інцидентів є порівнянною та передбачуваною в межах всієї організації;}


{\footnotesize\bfseries
No: 12\\
Name: ir\_\allowbreak{}4\_\allowbreak{}d\_\allowbreak{}3\\
Type: string\\
Default: nil
}

{\footnotesize обсяг заходів з обробки інцидентів є порівнянною та передбачуваною в межах всієї організації;}


{\footnotesize\bfseries
No: 13\\
Name: ir\_\allowbreak{}4\_\allowbreak{}d\_\allowbreak{}4\\
Type: string\\
Default: nil
}

{\footnotesize результати діяльності  заходів з обробки інцидентів є порівнянною та передбачуваною в межах всієї організації;}




\subsubsection{ДИНАМІЧНА РЕКОНФІГУРАЦІЯ (IR-4(2))}


\vspace{10pt}
{\footnotesize\bfseries
No: 1\\
Name: ir\_\allowbreak{}4\_\allowbreak{}2\_\allowbreak{}odp\_\allowbreak{}1\\
Type: string\\
Default: nil
}

{\footnotesize визначено типи динамічної реконфігурації для компонентів системи;}


{\footnotesize\bfseries
No: 2\\
Name: ir\_\allowbreak{}4\_\allowbreak{}2\_\allowbreak{}odp\_\allowbreak{}2\\
Type: string\\
Default: nil
}

{\footnotesize визначено компоненти системи, які потребують динамічної реконфігурації;}


{\footnotesize\bfseries
No: 3\\
Name: ir\_\allowbreak{}4\_\allowbreak{}2\_\allowbreak{}01\\
Type: string\\
Default: nil
}

{\footnotesize <IR-04(02)\_\allowbreak{}ODP[01] типи динамічної реконфігурації> для частина здатності реагування на інциденти.}




\subsubsection{БЕЗПЕРЕРВНІСТЬ ОПЕРАЦІЙ (IR-4(3))}


\vspace{10pt}
{\footnotesize\bfseries
No: 1\\
Name: ir\_\allowbreak{}4\_\allowbreak{}3\_\allowbreak{}odp\_\allowbreak{}1\\
Type: string\\
Default: nil
}

{\footnotesize визначено класи інцидентів, що вимагають вживання}


{\footnotesize\bfseries
No: 2\\
Name: ir\_\allowbreak{}4\_\allowbreak{}3\_\allowbreak{}odp\_\allowbreak{}2\\
Type: string\\
Default: nil
}

{\footnotesize визначено дії, які необхідно вжити у відповідь на визначені організацією класи інцидентів;}


{\footnotesize\bfseries
No: 3\\
Name: ir\_\allowbreak{}4\_\allowbreak{}3\_\allowbreak{}1\\
Type: string\\
Default: nil
}

{\footnotesize ідентифіковано <IR-04(03)\_\allowbreak{}ODP[01] класи інцидентів>;}


{\footnotesize\bfseries
No: 4\\
Name: ir\_\allowbreak{}4\_\allowbreak{}3\_\allowbreak{}2\\
Type: string\\
Default: nil
}

{\footnotesize <IR-04(03)\_\allowbreak{}ODP[02] дії> вживаються у відповідь на ці інциденти (визначені в IR-04(03)\_\allowbreak{}ODP[01]) для забезпечення продовження виконання завдань та функцій організації.}




\subsubsection{ІНФОРМАЦІЙНА КОРЕЛЯЦІЯ (IR-4(4))}


\vspace{10pt}
{\footnotesize\bfseries
No: 1\\
Name: ir\_\allowbreak{}4\_\allowbreak{}4\_\allowbreak{}01\\
Type: string\\
Default: nil
}

{\footnotesize інформація про інциденти та індивідуальне реагування на інциденти зіставляється з метою досягнення загальноорганізаційного бачення на обізнаність про інциденти та реагування на них.}




\subsubsection{ВНУТРІШНІ ЗАГРОЗИ - ОСОБЛИВІ МОЖЛИВОСТІ (IR-4(6))}


\vspace{10pt}
{\footnotesize\bfseries
No: 1\\
Name: ir\_\allowbreak{}4\_\allowbreak{}6\_\allowbreak{}01\\
Type: string\\
Default: nil
}

{\footnotesize реалізовано можливість обробки інцидентів, пов'язаних з внутрішніми загрозами}




\subsubsection{КООРДИНАЦІЯ З ЗОВНІШНІМИ ОРГАНІЗАЦІЯМИ (IR-4(8))}


\vspace{10pt}
{\footnotesize\bfseries
No: 1\\
Name: ir\_\allowbreak{}4\_\allowbreak{}8\_\allowbreak{}odp\_\allowbreak{}1\\
Type: string\\
Default: nil
}

{\footnotesize визначено зовнішні організації, з якими необхідно координувати та обмінюватися інформацією про інциденти в організації;}


{\footnotesize\bfseries
No: 2\\
Name: ir\_\allowbreak{}4\_\allowbreak{}8\_\allowbreak{}odp\_\allowbreak{}2\\
Type: string\\
Default: nil
}

{\footnotesize визначено інформацію про інциденти, яку необхідно зіставляти та поширювати із зовнішніми організаціями;}


{\footnotesize\bfseries
No: 3\\
Name: ir\_\allowbreak{}4\_\allowbreak{}8\_\allowbreak{}01\\
Type: string\\
Default: nil
}

{\footnotesize здійснюється координація з  <IR-04(08)\_\allowbreak{}ODP[01] нення міжорганізаційного бачення щодо обізнаності про інциденти та більш ефективного реагування на інциденти.}




\subsubsection{ЗДАТНІСТЬ ДИНАМІЧНОГО РЕАГУВАННЯ (IR-4(9))}


\vspace{10pt}
{\footnotesize\bfseries
No: 1\\
Name: ir\_\allowbreak{}4\_\allowbreak{}9\_\allowbreak{}odp\\
Type: string\\
Default: nil
}

{\footnotesize визначено можливості динамічного реагування для ефективного реагування на інциденти безпеки.}


{\footnotesize\bfseries
No: 2\\
Name: ir\_\allowbreak{}4\_\allowbreak{}9\_\allowbreak{}01\\
Type: string\\
Default: nil
}

{\footnotesize використуються <IR-04(09)\_\allowbreak{}ODP можливості динамічного реагування> для ефективного реагування на інциденти безпеки.}




\subsubsection{КООРДИНАЦІЯ ЛАНЦЮГА ПОСТАЧАННЯ (IR-4(10))}


\vspace{10pt}
{\footnotesize\bfseries
No: 1\\
Name: ir\_\allowbreak{}4\_\allowbreak{}10\_\allowbreak{}01\\
Type: string\\
Default: nil
}

{\footnotesize координується діяльність з обробки інцидентів, пов'язана з подіями ланцюжка постачання, з іншими організаціями, що беруть участь у 300 ланцюжку постачання.}




\subsubsection{ІНТЕГРОВАНА ГРУПА РЕАГУВАННЯ НА ІНЦЕДЕНТИ (IR-4(11))}


\vspace{10pt}
{\footnotesize\bfseries
No: 1\\
Name: ir\_\allowbreak{}4\_\allowbreak{}11\_\allowbreak{}odp\\
Type: string\\
Default: nil
}

{\footnotesize визначено період часу, протягом якого може бути розгорнута інтегрована група реагування на інцидент;}


{\footnotesize\bfseries
No: 2\\
Name: ir\_\allowbreak{}4\_\allowbreak{}11\_\allowbreak{}1\\
Type: string\\
Default: nil
}

{\footnotesize створена та підтримується інтегрована група реагування на інциденти;}


{\footnotesize\bfseries
No: 3\\
Name: ir\_\allowbreak{}4\_\allowbreak{}11\_\allowbreak{}2\\
Type: string\\
Default: nil
}

{\footnotesize інтегрована група реагування на інциденти може бути розгорнута в будь -якому місці, визначеному організацією протягом <IR04(11)\_\allowbreak{}ODP часового періоду>.}




\subsubsection{ЗЛОВМИСНИЙ КОД ТА КРИМІНАЛІСТИЧНИЙ АНАЛІЗ (IR-4(12))}


\vspace{10pt}
{\footnotesize\bfseries
No: 1\\
Name: ir\_\allowbreak{}4\_\allowbreak{}12\_\allowbreak{}1\\
Type: string\\
Default: nil
}

{\footnotesize шкідливий код, що залишився в системі, аналізується після інциденту;}


{\footnotesize\bfseries
No: 2\\
Name: ir\_\allowbreak{}4\_\allowbreak{}12\_\allowbreak{}2\\
Type: string\\
Default: nil
}

{\footnotesize інші залишкові артефакти, що залишилися в системі (якщо такі є), аналізуються після інциденту. 301}




\subsubsection{АНАЛІЗ ПОВЕДІНКИ (IR-4(13))}


\vspace{10pt}
{\footnotesize\bfseries
No: 1\\
Name: ir\_\allowbreak{}4\_\allowbreak{}13\_\allowbreak{}odp\\
Type: string\\
Default: nil
}

{\footnotesize визначаються середовища або ресурси, які можуть містити або можуть бути пов'язані з аномальною або підозрілою ворожою поведінкою;}


{\footnotesize\bfseries
No: 2\\
Name: ir\_\allowbreak{}4\_\allowbreak{}13\_\allowbreak{}01\\
Type: string\\
Default: nil
}

{\footnotesize аналізується аномальна або підозрювана ворожа поведінка в <IR-04(13)\_\allowbreak{}ODP середовищах або ресурсах > або пов'язана з ними.}




\subsubsection{ЦЕНТР БЕЗПЕКИ (IR-4(14))}


\vspace{10pt}
{\footnotesize\bfseries
No: 1\\
Name: ir\_\allowbreak{}4\_\allowbreak{}14\_\allowbreak{}1\\
Type: string\\
Default: nil
}

{\footnotesize створено оперативний центр безпеки;}


{\footnotesize\bfseries
No: 2\\
Name: ir\_\allowbreak{}4\_\allowbreak{}14\_\allowbreak{}2\\
Type: string\\
Default: nil
}

{\footnotesize підтримується оперативний центр безпеки;}




\subsubsection{ЗВ'ЯЗКИ З ГРОМАДКІСТЮ ТА ВІДНОВЛЕННЯ РЕПУТАЦІЇ (IR-4(15))}


\vspace{10pt}
\textit{Немає параметрів для цього контролю.}
\vspace{10pt}


\subsection{МОНІТОРИНГ ІНЦИДЕНТУ (IR-5)}


\vspace{10pt}
{\footnotesize\bfseries
No: 1\\
Name: ir\_\allowbreak{}5\_\allowbreak{}1\\
Type: string\\
Default: nil
}

{\footnotesize відстежуються інциденти безпеки та приватності;}


{\footnotesize\bfseries
No: 2\\
Name: ir\_\allowbreak{}5\_\allowbreak{}2\\
Type: string\\
Default: nil
}

{\footnotesize документуються інциденти безпеки та приватності. 303}




\subsubsection{АВТОМАТИЗОВАНЕ ВІДСТЕЖЕННЯ, ЗБІР ДАНИХ І АНАЛІЗ (IR-5(1))}


\vspace{10pt}
{\footnotesize\bfseries
No: 1\\
Name: ir\_\allowbreak{}5\_\allowbreak{}1\_\allowbreak{}odp\_\allowbreak{}1\\
Type: string\\
Default: nil
}

{\footnotesize визначено автоматизовані механізми відстеження інцидентів;}


{\footnotesize\bfseries
No: 2\\
Name: ir\_\allowbreak{}5\_\allowbreak{}1\_\allowbreak{}odp\_\allowbreak{}2\\
Type: string\\
Default: nil
}

{\footnotesize визначено автоматизовані механізми збору інформації про інциденти;}


{\footnotesize\bfseries
No: 3\\
Name: ir\_\allowbreak{}5\_\allowbreak{}1\_\allowbreak{}odp\_\allowbreak{}3\\
Type: string\\
Default: nil
}

{\footnotesize визначено автоматизовані механізми аналізу інформації про інциденти;}


{\footnotesize\bfseries
No: 4\\
Name: ir\_\allowbreak{}5\_\allowbreak{}1\_\allowbreak{}1\\
Type: string\\
Default: nil
}

{\footnotesize інциденти відстежуються за допомогою <IR05(01)\_\allowbreak{}ODP[01] автоматизованих механізмів>;}


{\footnotesize\bfseries
No: 5\\
Name: ir\_\allowbreak{}5\_\allowbreak{}1\_\allowbreak{}2\\
Type: string\\
Default: nil
}

{\footnotesize інформація про інциденти збирається за допомогою <IR05(01)\_\allowbreak{}ODP[02] автоматизованих механізмів>;}


{\footnotesize\bfseries
No: 6\\
Name: ir\_\allowbreak{}5\_\allowbreak{}1\_\allowbreak{}3\\
Type: string\\
Default: nil
}

{\footnotesize інформація про інциденти аналізується за допомогою <IR05(01)\_\allowbreak{}ODP[03] автоматизованих механізмів>.}




\subsection{ЗВІТНІСТЬ ПРО ІНЦИДЕНТИ (IR-6)}


\vspace{10pt}
{\footnotesize\bfseries
No: 1\\
Name: ir\_\allowbreak{}6\_\allowbreak{}odp\_\allowbreak{}1\\
Type: string\\
Default: nil
}

{\footnotesize визначено період часу, протягом якого персонал повинен повідомляти про підозрілі інциденти до уповноваженого органу;}


{\footnotesize\bfseries
No: 2\\
Name: ir\_\allowbreak{}6\_\allowbreak{}odp\_\allowbreak{}2\\
Type: string\\
Default: nil
}

{\footnotesize визначені органи, до яких слід повідомляти інформацію про інцидент;}


{\footnotesize\bfseries
No: 3\\
Name: ir\_\allowbreak{}6\_\allowbreak{}a\\
Type: string\\
Default: nil
}

{\footnotesize персонал зобов'язаний повідомляти про підозрілі інциде нти протягом <IR-06\_\allowbreak{}ODP[01] періоду часу>;}


{\footnotesize\bfseries
No: 4\\
Name: ir\_\allowbreak{}6\_\allowbreak{}b\\
Type: string\\
Default: nil
}

{\footnotesize інформацію про інцидент повідомляється  <IR-06\_\allowbreak{}ODP[02] органам>.}




\subsubsection{АВТОМАТИЧНЕ ЗВІТУВАННЯ (IR-6(1))}


\vspace{10pt}
{\footnotesize\bfseries
No: 1\\
Name: ir\_\allowbreak{}6\_\allowbreak{}1\_\allowbreak{}odp\\
Type: string\\
Default: nil
}

{\footnotesize визначені автоматичні механізми звітування про інциденти;}


{\footnotesize\bfseries
No: 2\\
Name: ir\_\allowbreak{}6\_\allowbreak{}1\_\allowbreak{}01\\
Type: string\\
Default: nil
}

{\footnotesize використовуються <IR-06(01)\_\allowbreak{}ODP автоматичні механізми> звітування про інциденти.}




\subsubsection{КООРДИНАЦІЯ ЛАНЦЮЖКА ПОСТАЧАННЯ (IR-6(3))}


\vspace{10pt}
{\footnotesize\bfseries
No: 1\\
Name: ir\_\allowbreak{}6\_\allowbreak{}3\_\allowbreak{}01\\
Type: string\\
Default: nil
}

{\footnotesize інформація про інцидент надається постачальнику продукту або послуги та іншим організаціям, які беруть участь у ланцюжку постачання систем або компонентів системи, пов'язаних з інцидентом.}




\subsection{ПІДТРИМКА РЕАГУВАННЯ НА ІНЦИДЕНТИ (IR-7)}


\vspace{10pt}
{\footnotesize\bfseries
No: 1\\
Name: ir\_\allowbreak{}7\_\allowbreak{}1\\
Type: string\\
Default: nil
}

{\footnotesize надається ресурс підтримки реагування на інциденти, що є невід'ємною частиною спроможності організації реагувати на інциденти;}


{\footnotesize\bfseries
No: 2\\
Name: ir\_\allowbreak{}7\_\allowbreak{}1\\
Type: string\\
Default: nil
}

{\footnotesize ресурс підтримки реагування на інциденти містить поради та допомогу користувачам системи для обробки та формування звітності про інциденти.}




\subsubsection{АВТОМАТИЗАЦІЯ ПІДТРИМКИ ДЛЯ ЗАБЕЗПЕЧЕННЯ ДОСТУПНОСТІ ІНФОРМАЦІЇ ТА ПІДТРИМКИ (IR-7(1))}


\vspace{10pt}
{\footnotesize\bfseries
No: 1\\
Name: ir\_\allowbreak{}7\_\allowbreak{}1\_\allowbreak{}odp\\
Type: string\\
Default: nil
}

{\footnotesize визначено автоматизовані механізми, що використовуються для збільшення доступності інформації та підтримки при реагуванні на інциденти;}


{\footnotesize\bfseries
No: 2\\
Name: ir\_\allowbreak{}7\_\allowbreak{}1\_\allowbreak{}01\\
Type: string\\
Default: nil
}

{\footnotesize підвищено доступність інформації та підтримки реагування на інциденти з використанням <IR-07(01)\_\allowbreak{}ODP автоматизованих механізмів>.}




\subsubsection{КООРДИНАЦІЯ З ЗОВНІШНІМИ ПОСТАЧАЛЬНИКАМИ (IR-7(2))}


\vspace{10pt}
{\footnotesize\bfseries
No: 1\\
Name: ir\_\allowbreak{}7\_\allowbreak{}2\_\allowbreak{}a\\
Type: string\\
Default: nil
}

{\footnotesize встановлено прямі відносини кооперації між здатністю реагування на інциденти та зовнішніми постачальниками можливостей захисту системи.}


{\footnotesize\bfseries
No: 2\\
Name: ir\_\allowbreak{}7\_\allowbreak{}2\_\allowbreak{}b\\
Type: string\\
Default: nil
}

{\footnotesize визначено членів команди реагування на інциденти в організації для зовнішніх постачальників послуг.}




\subsection{ПЛАН РЕАГУВАННЯ НА ІНЦИДЕНТИ (IR-8)}


\vspace{10pt}
{\footnotesize\bfseries
No: 1\\
Name: ir\_\allowbreak{}8\_\allowbreak{}odp\_\allowbreak{}1\\
Type: string\\
Default: nil
}

{\footnotesize визначено персонал або ролі, які переглядають та затверджують план реагування на інциденти;}


{\footnotesize\bfseries
No: 2\\
Name: ir\_\allowbreak{}8\_\allowbreak{}odp\_\allowbreak{}2\\
Type: string\\
Default: nil
}

{\footnotesize визначено періодичність перегляду та затвердження плану реагування на інциденти;}


{\footnotesize\bfseries
No: 3\\
Name: ir\_\allowbreak{}8\_\allowbreak{}odp\_\allowbreak{}3\\
Type: string\\
Default: nil
}

{\footnotesize визначені організації, персонал або ролі, які несуть відповідальність за реагування на інциденти; 308}


{\footnotesize\bfseries
No: 4\\
Name: ir\_\allowbreak{}8\_\allowbreak{}odp\_\allowbreak{}4\\
Type: string\\
Default: nil
}

{\footnotesize визначено персонал з реагування на інцидент (ідентифікований за іменами та/або за ролями), якому мають бути роздані копії плану реагування на інцидент;}


{\footnotesize\bfseries
No: 5\\
Name: ir\_\allowbreak{}8\_\allowbreak{}odp\_\allowbreak{}5\\
Type: string\\
Default: nil
}

{\footnotesize визначено елементи організації, серед яких мають бути розповсюджені копії плану реагування на інцидент;}


{\footnotesize\bfseries
No: 6\\
Name: ir\_\allowbreak{}8\_\allowbreak{}odp\_\allowbreak{}6\\
Type: string\\
Default: nil
}

{\footnotesize визначено персонал з реагування на інцидент (ідентифікований за іменами та/або ролями), якому повідомляються зміни до плану реагування на інцидент;}


{\footnotesize\bfseries
No: 7\\
Name: ir\_\allowbreak{}8\_\allowbreak{}odp\_\allowbreak{}7\\
Type: string\\
Default: nil
}

{\footnotesize визначено елементи організації, яким повідомляється про зміни в плані реагування на інцидент;}


{\footnotesize\bfseries
No: 8\\
Name: ir\_\allowbreak{}8\_\allowbreak{}a\_\allowbreak{}1\\
Type: string\\
Default: nil
}

{\footnotesize розроблено план реагування на інциденти, який надає організації дорожню карту для впровадження її можливостей реагування на інциденти;}


{\footnotesize\bfseries
No: 9\\
Name: ir\_\allowbreak{}8\_\allowbreak{}a\_\allowbreak{}2\\
Type: string\\
Default: nil
}

{\footnotesize розроблено план реагування на інциденти, який  описує структуру та організацію спроможності реагування на інциденти;}


{\footnotesize\bfseries
No: 10\\
Name: ir\_\allowbreak{}8\_\allowbreak{}a\_\allowbreak{}3\\
Type: string\\
Default: nil
}

{\footnotesize розроблено план реагування на інциденти, який надає високорівневий підхід до того, як здатність реагування на інциденти вписується в загальну практику організації;}


{\footnotesize\bfseries
No: 11\\
Name: ir\_\allowbreak{}8\_\allowbreak{}a\_\allowbreak{}4\\
Type: string\\
Default: nil
}

{\footnotesize розроблено план реагування на інциденти, який відповідає унікальним вимогам організації, які пов'язані із завданнями, розміром, структурою і функціями;}


{\footnotesize\bfseries
No: 12\\
Name: ir\_\allowbreak{}8\_\allowbreak{}a\_\allowbreak{}5\\
Type: string\\
Default: nil
}

{\footnotesize розроблено план реагування на інциденти, який визначає підзвітні інциденти;}


{\footnotesize\bfseries
No: 13\\
Name: ir\_\allowbreak{}8\_\allowbreak{}a\_\allowbreak{}6\\
Type: string\\
Default: nil
}

{\footnotesize розроблено план реагування на інциденти, який надає показники для вимірювання можливостей реагування на інциденти всередині організації;}


{\footnotesize\bfseries
No: 14\\
Name: ir\_\allowbreak{}8\_\allowbreak{}a\_\allowbreak{}7\\
Type: string\\
Default: nil
}

{\footnotesize розроблено план реагування на інциденти, який визначає ресурси та управлінську підтримку, необхідну для ефективної підтримки та розвитку можливостей реагування на інциденти;}


{\footnotesize\bfseries
No: 15\\
Name: ir\_\allowbreak{}8\_\allowbreak{}a\_\allowbreak{}8\\
Type: string\\
Default: nil
}

{\footnotesize розроблено план реагування на інциденти, який вирішує питання обміну інформацією про інциденти;}


{\footnotesize\bfseries
No: 16\\
Name: ir\_\allowbreak{}8\_\allowbreak{}a\_\allowbreak{}9\\
Type: string\\
Default: nil
}

{\footnotesize розроблено план реагування на інцидент, який розглядається та затверджується <IR-08\_\allowbreak{}ODP[01] персоналом або ролями> < IR08\_\allowbreak{}ODP[02] частота>;}


{\footnotesize\bfseries
No: 17\\
Name: ir\_\allowbreak{}8\_\allowbreak{}a\_\allowbreak{}10\\
Type: string\\
Default: nil
}

{\footnotesize розроблено план реагування на інциденти, в якому чітко визначено організацій, персоналу або ролей>.}


{\footnotesize\bfseries
No: 18\\
Name: ir\_\allowbreak{}8\_\allowbreak{}b\_\allowbreak{}1\\
Type: string\\
Default: nil
}

{\footnotesize копії плану реагування на інцидент розповсюджуються серед  <IR309}


{\footnotesize\bfseries
No: 19\\
Name: ir\_\allowbreak{}8\_\allowbreak{}b\_\allowbreak{}2\\
Type: string\\
Default: nil
}

{\footnotesize копії плану реагування на інцидент розповсюджуються серед  <IR08\_\allowbreak{}ODP[05] елементів організації>;}


{\footnotesize\bfseries
No: 20\\
Name: ir\_\allowbreak{}8\_\allowbreak{}c\\
Type: string\\
Default: nil
}

{\footnotesize план реагування на інциденти оновлюється з урахуванням змін у системі та організації або проблем, що виникають під час впровадження, виконання або тестування плану;}


{\footnotesize\bfseries
No: 21\\
Name: ir\_\allowbreak{}8\_\allowbreak{}d\_\allowbreak{}1\\
Type: string\\
Default: nil
}

{\footnotesize зміни в плані реагування на інцидент повідомляються <IR08\_\allowbreak{}ODP[06] персоналу з реагування на інциденти>;}


{\footnotesize\bfseries
No: 22\\
Name: ir\_\allowbreak{}8\_\allowbreak{}d\_\allowbreak{}2\\
Type: string\\
Default: nil
}

{\footnotesize зміни в плані реагування на інциденти надсилаються до  <IR08\_\allowbreak{}ODP[07] елементів організації>;}


{\footnotesize\bfseries
No: 23\\
Name: ir\_\allowbreak{}8\_\allowbreak{}e\_\allowbreak{}1\\
Type: string\\
Default: nil
}

{\footnotesize план реагування на інциденти захищений від несанкціонованого розкриття;}


{\footnotesize\bfseries
No: 24\\
Name: ir\_\allowbreak{}8\_\allowbreak{}e\_\allowbreak{}2\\
Type: string\\
Default: nil
}

{\footnotesize план реагування на інциденти захищений від несанкціонованої модифікації.}




\subsubsection{ОБРОБКА ПЕРСОНАЛЬНИХ ДАНИХ (IR-8(1))}


\vspace{10pt}
{\footnotesize\bfseries
No: 1\\
Name: ir\_\allowbreak{}8\_\allowbreak{}1\_\allowbreak{}a\\
Type: string\\
Default: nil
}

{\footnotesize план реагування на інциденти для інцидентів, пов'язаних з персональними даними, включає процес визначення доцільності повідомлення наглядових організацій і надання такого повідомлення, якщо це доречно;}


{\footnotesize\bfseries
No: 2\\
Name: ir\_\allowbreak{}8\_\allowbreak{}1\_\allowbreak{}b\\
Type: string\\
Default: nil
}

{\footnotesize план реагування на інциденти для інцидентів, пов'язаних з персональними даними, включає процес оцінювання для визначення ступеня шкоди, труднощів, незручностей або несправедливості щодо постраждалих осіб та будь-які механізми пом'якшення такої шкоди;}


{\footnotesize\bfseries
No: 3\\
Name: ir\_\allowbreak{}8\_\allowbreak{}1\_\allowbreak{}c\\
Type: string\\
Default: nil
}

{\footnotesize план реагування на інциденти для інцидентів, пов'язаних з персональними даними, включає ідентифікацію застосовних вимог щодо конфіденційності. 310}




\subsection{РЕАГУВАННЯ НА ВИТІК ІНФОРМАЦІЇ (IR-9)}


\vspace{10pt}
{\footnotesize\bfseries
No: 1\\
Name: ir\_\allowbreak{}9\_\allowbreak{}odp\_\allowbreak{}1\\
Type: string\\
Default: nil
}

{\footnotesize визначено персонал або ролі, на які покладено відповідальність за реагування на витоки інформації;}


{\footnotesize\bfseries
No: 2\\
Name: ir\_\allowbreak{}9\_\allowbreak{}odp\_\allowbreak{}2\\
Type: string\\
Default: nil
}

{\footnotesize визначено персонал або ролі, які мають бути сповіщені про витік інформації за допомогою методу зв'язку, не пов'язаного з витоком;}


{\footnotesize\bfseries
No: 3\\
Name: ir\_\allowbreak{}9\_\allowbreak{}odp\_\allowbreak{}3\\
Type: string\\
Default: nil
}

{\footnotesize визначені дії, які необхідно виконати;}


{\footnotesize\bfseries
No: 4\\
Name: ir\_\allowbreak{}9\_\allowbreak{}a\\
Type: string\\
Default: nil
}

{\footnotesize <IR-09\_\allowbreak{}ODP[01] персонал або ролі>  призначено відповідальним за реагування на витоки інформації;}


{\footnotesize\bfseries
No: 5\\
Name: ir\_\allowbreak{}9\_\allowbreak{}b\\
Type: string\\
Default: nil
}

{\footnotesize у відповідь на витік інформації  визначається конкретна інформація, пов'язана з джерелом витоку в системі;}


{\footnotesize\bfseries
No: 6\\
Name: ir\_\allowbreak{}9\_\allowbreak{}c\\
Type: string\\
Default: nil
}

{\footnotesize <IR-09\_\allowbreak{}ODP[02] персонал або ролі>  попереджається про витік інформації за допомогою методу зв'язку, не пов'язаного з витоком;}


{\footnotesize\bfseries
No: 7\\
Name: ir\_\allowbreak{}9\_\allowbreak{}d\\
Type: string\\
Default: nil
}

{\footnotesize ізоляцюється система або компонент  системи де відбувся витік інформації;}


{\footnotesize\bfseries
No: 8\\
Name: ir\_\allowbreak{}9\_\allowbreak{}e\\
Type: string\\
Default: nil
}

{\footnotesize інформація видаляється із зараженої системи або компонента у відповідь на витік інформації;}


{\footnotesize\bfseries
No: 9\\
Name: ir\_\allowbreak{}9\_\allowbreak{}f\\
Type: string\\
Default: nil
}

{\footnotesize у відповідь на витік інформації визначаються інші системи або компоненти системи, які могли бути згодом  джерелом витоку інформації;}


{\footnotesize\bfseries
No: 10\\
Name: ir\_\allowbreak{}9\_\allowbreak{}g\\
Type: string\\
Default: nil
}

{\footnotesize <IR-09\_\allowbreak{}ODP[03] дії> виконуються у відповідь на витік інформації.}




\subsubsection{ВІДПОВІДАЛЬНИЙ ПЕРСОНАЛ (IR-9(1)) [Вилучено]}
[Вилучено: включено до IR-09]

\vspace{10pt}
\textit{Немає параметрів для цього контролю.}
\vspace{10pt}


\subsubsection{ТРЕНУВАННЯ (IR-9(2))}


\vspace{10pt}
{\footnotesize\bfseries
No: 1\\
Name: ir\_\allowbreak{}9\_\allowbreak{}2\_\allowbreak{}odp\\
Type: string\\
Default: nil
}

{\footnotesize визначено частоту навчання з реагування на витік інформації;}


{\footnotesize\bfseries
No: 2\\
Name: ir\_\allowbreak{}9\_\allowbreak{}2\_\allowbreak{}01\\
Type: string\\
Default: nil
}

{\footnotesize забезпечено навчання з реагування на витік інформації <IR09(02)\_\allowbreak{}ODP частота>.}




\subsubsection{РОБОТА ПІСЛЯ ВИТОКУ (IR-9(3))}


\vspace{10pt}
{\footnotesize\bfseries
No: 1\\
Name: ir\_\allowbreak{}9\_\allowbreak{}3\_\allowbreak{}odp\\
Type: string\\
Default: nil
}

{\footnotesize визначено процедури з метою забезпечення спроможності персоналу організації, на який впливає витік інформації, продовжувати виконувати поставлені завдання, у той час, як постраждалі системи зазнають коригувальних дій.}


{\footnotesize\bfseries
No: 2\\
Name: ir\_\allowbreak{}9\_\allowbreak{}3\_\allowbreak{}01\\
Type: string\\
Default: nil
}

{\footnotesize реалізувано <IR-09(03)\_\allowbreak{}ODP процедури>, з метою забезпечення спроможності для персоналу організації, на який впливає витік інформації, продовжувати виконувати поставлені завдання, у той час, як постраждалі системи зазнають коригу312 вальних дій.}




\subsubsection{ВИКРИТТЯ НЕАВТОРИЗОВАНОГО ПЕРСОНАЛУ (IR-9(4))}


\vspace{10pt}
{\footnotesize\bfseries
No: 1\\
Name: ir\_\allowbreak{}9\_\allowbreak{}4\_\allowbreak{}odp\\
Type: string\\
Default: nil
}

{\footnotesize визначено механізми захисту  для персоналу, що має доступ до інформації, яка не відповідає призначеним правам доступу;}


{\footnotesize\bfseries
No: 2\\
Name: ir\_\allowbreak{}9\_\allowbreak{}4\_\allowbreak{}01\\
Type: string\\
Default: nil
}

{\footnotesize застосовуються <IR-09(04)\_\allowbreak{}ODP механізми захисту > для персоналу, що має доступ до інформації, яка не відпов ідає призначеним правам доступу.}




\subsection{ІНТЕГРОВАНА КОМАНДА АНАЛІЗУ ІНФОРМАЦІЙНОЇ БЕЗПЕКИ (IR-10) [Вилучено]}
[Вилучено: перенесено до IR-04(11)]

\vspace{10pt}
\textit{Немає параметрів для цього контролю.}
\vspace{10pt}


\section{MA}
Клас заходів захисту MA --- ТЕХНІЧНЕ ОБСЛУГОВУВАННЯ

Цей клас регулює процеси планового та позапланового технічного обслуговування компонентів системи для запобігання збоям.

\textbf{Перелік заходів захисту:} Політика та процедури технічного обслуговування (MA-1); Зміст запису (MA-2(1)) [Вилучено]; Автоматизована технічна діяльність (MA-2(2)); Інструменти для обслуговування (MA-3); Перевірка інструментів (MA-3(1)); Перевірка носіїв інформації (MA-3(2)); Запобігання несанкціонованому переміщенню (MA-3(3)); Обмеження використання інструмента (MA-3(4)); Привілейоване виконання (MA-3(5)); Оновлення програм- (MA-3(6)); Віддалене обслуговування (MA-4); Аудит та огляд (MA-4(1)); Документування віддаленого обслуговування (MA-4(2)) [Вилучено]; Порівняльна безпека і очищення (MA-4(3)); Схвалення та повідомлення (MA-4(5)); Віддалене обслуговування роз'єднання (MA-4(7)); Технічний персонал (MA-5); Особи без належного доступу (MA-5(1)); Оформлення допуску для систем, що обробляють інформацію з обмеженим доступом (MA-5(2)); Вимоги до громадянства (MA-5(3)); Іноземні громадяни (MA-5(4)); Несистемне обслуговування (MA-5(5)); Своєчасне обслуговування (MA-6); Профілактичне обслуговування (MA-6(1)); Планове технічне обслуговування (MA-6(2)); Автоматизована підтримка планового технічного обслуговування (MA-6(3)); Технічне обслуговування в польових умовах (MA-7).

\subsection{ПОЛІТИКА ТА ПРОЦЕДУРИ ТЕХНІЧНОГО ОБСЛУГОВУВАННЯ (MA-1)}
a. Планувати, документувати та переглядати записи з технічного обслуговування, ремонту або заміни компонентів системи відповідно до вимог виробника та постачальників та/або вимог організації.\\ 
b. Затвердити та здійснювати моніторинг усіх заходів з технічного обслуговування, незалежно від того, виконуються вони на місці або віддалено, а також чи обслуговуються системи або системні компоненти на місці, чи переміщуються в інше місце.\\ 
c. Вимагати, щоб [Призначення: визначені організацією персонал чи ролі] явно схвалили видалення системи або компоненту системи з організаційного обладнання для технічного обслуговування, ремонту чи заміни поза об'єктами експлуатації.\\ 
d. Очищати обладнання з погляду видалення всієї інформації з носіїв до вилучення обладнання організації для технічного обслуговування, ремонту чи заміни поза об'єктами експлуатації.\\ 
e. Перевірити всі потенційно порушені заходи захисту, щоб переконатися, що вони, як і раніше, працюють належним чином після дій з обслуговування, ремонту або заміни.\\ 
f. Вносити [Призначення: визначену організацією інформацію, пов'язану з технічним обслуговуванням] до записів з технічного обслуговування.

\vspace{10pt}
{\footnotesize\bfseries
No: 1\\
Name: ma\_\allowbreak{}1\_\allowbreak{}odp\_\allowbreak{}1\\
Type: string\\
Default: nil
}

{\footnotesize визначено персонал або ролі, на які поширюється політика технічного обслуговування;}


{\footnotesize\bfseries
No: 2\\
Name: ma\_\allowbreak{}1\_\allowbreak{}odp\_\allowbreak{}2\\
Type: string\\
Default: nil
}

{\footnotesize визначено персонал або ролі, на які поширюються процедури технічного обслуговування;}


{\footnotesize\bfseries
No: 3\\
Name: ma\_\allowbreak{}1\_\allowbreak{}odp\_\allowbreak{}3\\
Type: string\\
Default: nil
}

{\footnotesize вибрано одне або декілька з наступних ЗНАЧЕНЬ ПАРА- МЕТРІВ: \{рівень організації; рівень місії/бізнес -процесу; рівень системи\};}


{\footnotesize\bfseries
No: 4\\
Name: ma\_\allowbreak{}1\_\allowbreak{}odp\_\allowbreak{}4\\
Type: string\\
Default: nil
}

{\footnotesize визначено посадову особу, яка керуватиме політикою та процедурами технічного обслуговування;}


{\footnotesize\bfseries
No: 5\\
Name: ma\_\allowbreak{}1\_\allowbreak{}odp\_\allowbreak{}5\\
Type: string\\
Default: nil
}

{\footnotesize визначено частоту, з якою переглядається та оновлюється поточна політика технічного обслуговування;}


{\footnotesize\bfseries
No: 6\\
Name: ma\_\allowbreak{}1\_\allowbreak{}odp\_\allowbreak{}6\\
Type: string\\
Default: nil
}

{\footnotesize визначено події, які потребують перегляду та оновлення поточної політики технічного обслуговування;}


{\footnotesize\bfseries
No: 7\\
Name: ma\_\allowbreak{}1\_\allowbreak{}odp\_\allowbreak{}7\\
Type: string\\
Default: nil
}

{\footnotesize визначено частоту, з якою переглядаються та оновлюються поточні процедури технічного обслуговування;}


{\footnotesize\bfseries
No: 8\\
Name: ma\_\allowbreak{}1\_\allowbreak{}odp\_\allowbreak{}8\\
Type: string\\
Default: nil
}

{\footnotesize визначено події, які потребують перегляду та оновлення процедур технічного обслуговування;}


{\footnotesize\bfseries
No: 9\\
Name: ma\_\allowbreak{}1\_\allowbreak{}a\_\allowbreak{}1\\
Type: string\\
Default: nil
}

{\footnotesize розроблено та задокументовано політику технічного обслуговування;}


{\footnotesize\bfseries
No: 10\\
Name: ma\_\allowbreak{}1\_\allowbreak{}a\_\allowbreak{}2\\
Type: string\\
Default: nil
}

{\footnotesize політика технічного обслуговування поширюється на  <MA01\_\allowbreak{}ODP[01] персонал або ролі>;}


{\footnotesize\bfseries
No: 11\\
Name: ma\_\allowbreak{}1\_\allowbreak{}a\_\allowbreak{}3\\
Type: string\\
Default: nil
}

{\footnotesize розроблені та задокументовані процедури технічного обслуговування, що сприяють впровадженню політики технічного обслуговування та пов'язаних з нею заходів технічного обслуговування;}


{\footnotesize\bfseries
No: 12\\
Name: ma\_\allowbreak{}1\_\allowbreak{}a\_\allowbreak{}4\\
Type: string\\
Default: nil
}

{\footnotesize процедури технічного обслуговування розповсюджуються серед <MA-01\_\allowbreak{}ODP[02] персоналу або ролей>;}


{\footnotesize\bfseries
No: 13\\
Name: ma\_\allowbreak{}1\_\allowbreak{}a\_\allowbreak{}1\_\allowbreak{}a\_\allowbreak{}1\\
Type: string\\
Default: nil
}

{\footnotesize політика технічного обслуговування  <MA-01\_\allowbreak{}ODP[03] ВИБРАНЕ ЗНАЧЕННЯ ПАРАМЕТРА(ів)> містить мету; 314}


{\footnotesize\bfseries
No: 14\\
Name: ma\_\allowbreak{}1\_\allowbreak{}a\_\allowbreak{}1\_\allowbreak{}a\_\allowbreak{}2\\
Type: string\\
Default: nil
}

{\footnotesize політика технічного обслуговування  <MA-01\_\allowbreak{}ODP[03] ВИБРАНЕ ЗНАЧЕННЯ ПАРАМЕТРА(ів)> містить сферу застосування;}


{\footnotesize\bfseries
No: 15\\
Name: ma\_\allowbreak{}1\_\allowbreak{}a\_\allowbreak{}1\_\allowbreak{}a\_\allowbreak{}3\\
Type: string\\
Default: nil
}

{\footnotesize політика технічного обслуговування  <MA-01\_\allowbreak{}ODP[03] ВИБРАНЕ ЗНАЧЕННЯ ПАРАМЕТРА(ів)> містить ролі;}


{\footnotesize\bfseries
No: 16\\
Name: ma\_\allowbreak{}1\_\allowbreak{}a\_\allowbreak{}1\_\allowbreak{}a\_\allowbreak{}4\\
Type: string\\
Default: nil
}

{\footnotesize політика технічного обслуговування  <MA-01\_\allowbreak{}ODP[03] ВИБРАНЕ ЗНАЧЕННЯ ПАРАМЕТРА(ів)> містить обов'язки;}


{\footnotesize\bfseries
No: 17\\
Name: ma\_\allowbreak{}1\_\allowbreak{}a\_\allowbreak{}1\_\allowbreak{}a\_\allowbreak{}5\\
Type: string\\
Default: nil
}

{\footnotesize політика технічного обслуговування <MA-01\_\allowbreak{}ODP[03] ВИБРАНЕ ЗНАЧЕННЯ ПАРАМЕТРА(ів)> містить відповідальність керівництва;}


{\footnotesize\bfseries
No: 18\\
Name: ma\_\allowbreak{}1\_\allowbreak{}a\_\allowbreak{}1\_\allowbreak{}a\_\allowbreak{}6\\
Type: string\\
Default: nil
}

{\footnotesize політика технічного обслуговування <MA-01\_\allowbreak{}ODP[03] ВИБРАНЕ ЗНАЧЕННЯ ПАРАМЕТРА(ів)> містить координацію між підрозділами організації;}


{\footnotesize\bfseries
No: 19\\
Name: ma\_\allowbreak{}1\_\allowbreak{}a\_\allowbreak{}1\_\allowbreak{}a\_\allowbreak{}7\\
Type: string\\
Default: nil
}

{\footnotesize політика технічного обслуговування  <MA-01\_\allowbreak{}ODP[03] ВИБРАНЕ ЗНАЧЕННЯ ПАРАМЕТРА(ів)> містить систему контролю відповідності;}


{\footnotesize\bfseries
No: 20\\
Name: ma\_\allowbreak{}1\_\allowbreak{}a\_\allowbreak{}1\_\allowbreak{}b\\
Type: string\\
Default: nil
}

{\footnotesize політика технічного обслуговування  <MA-01\_\allowbreak{}ODP[03] ВИБРАНЕ ЗНАЧЕННЯ ПАРАМЕТРА(ів)> відповідає чинному законодавству, виконавчим наказам, директивам, нормам, політикам, стандартам і керівним принципам;}


{\footnotesize\bfseries
No: 21\\
Name: ma\_\allowbreak{}1\_\allowbreak{}b\\
Type: string\\
Default: nil
}

{\footnotesize <MA-01\_\allowbreak{}ODP[04] посадова особа> призначається для управління розробкою, документуванням та розповсюдженням політики та процедур технічного обслуговування;}


{\footnotesize\bfseries
No: 22\\
Name: ma\_\allowbreak{}1\_\allowbreak{}c\_\allowbreak{}1\_\allowbreak{}1\\
Type: string\\
Default: nil
}

{\footnotesize переглядається та оновлюється поточна політика обслуговування <MA-01\_\allowbreak{}ODP[05] частота>;}


{\footnotesize\bfseries
No: 23\\
Name: ma\_\allowbreak{}1\_\allowbreak{}c\_\allowbreak{}1\_\allowbreak{}2\\
Type: string\\
Default: nil
}

{\footnotesize переглядається та оновлюється поточна політика обслуговування після <MA-01\_\allowbreak{}ODP[06] подій>;}


{\footnotesize\bfseries
No: 24\\
Name: ma\_\allowbreak{}1\_\allowbreak{}c\_\allowbreak{}2\_\allowbreak{}1\\
Type: string\\
Default: nil
}

{\footnotesize переглядаються та оновлюються поточні процедури технічного}


{\footnotesize\bfseries
No: 25\\
Name: ma\_\allowbreak{}1\_\allowbreak{}c\_\allowbreak{}2\_\allowbreak{}2\\
Type: string\\
Default: nil
}

{\footnotesize переглядаються та оновлюються поточні процедури технічного}




\subsubsection{ЗМІСТ ЗАПИСУ (MA-2(1)) [Вилучено]}
[Вилучено: Включено до МА-02]

\vspace{10pt}
\textit{Немає параметрів для цього контролю.}
\vspace{10pt}


\subsubsection{АВТОМАТИЗОВАНА ТЕХНІЧНА ДІЯЛЬНІСТЬ (MA-2(2))}


\vspace{10pt}
{\footnotesize\bfseries
No: 1\\
Name: ma\_\allowbreak{}2\_\allowbreak{}2\_\allowbreak{}odp\_\allowbreak{}1\\
Type: string\\
Default: nil
}

{\footnotesize визначено автоматизовані механізми, що використовуються для планування дій з технічного обслуговування, ремонту та заміни системи;}


{\footnotesize\bfseries
No: 2\\
Name: ma\_\allowbreak{}2\_\allowbreak{}2\_\allowbreak{}odp\_\allowbreak{}2\\
Type: string\\
Default: nil
}

{\footnotesize визначено автоматизовані механізми, що використовуються для проведення дій з технічного обслуговування, ремонту та заміни системи;}


{\footnotesize\bfseries
No: 3\\
Name: ma\_\allowbreak{}2\_\allowbreak{}2\_\allowbreak{}odp\_\allowbreak{}3\\
Type: string\\
Default: nil
}

{\footnotesize визначено автоматизовані механізми, що використовуються для документування дій з технічного обслуговування, ремонту та заміни системи;}


{\footnotesize\bfseries
No: 4\\
Name: ma\_\allowbreak{}2\_\allowbreak{}2\_\allowbreak{}a\_\allowbreak{}1\\
Type: string\\
Default: nil
}

{\footnotesize <MA-02(02)\_\allowbreak{}ODP[01] автоматизовані механізми> використовуються для планування дій з технічного обслуговування, ремонту та заміни системи;}


{\footnotesize\bfseries
No: 5\\
Name: ma\_\allowbreak{}2\_\allowbreak{}2\_\allowbreak{}a\_\allowbreak{}2\\
Type: string\\
Default: nil
}

{\footnotesize <MA-02(02)\_\allowbreak{}ODP[02] автоматизовані механізми> використовуються для проведення дій з технічного обслуговування, ремонту та заміни системи; 317}


{\footnotesize\bfseries
No: 6\\
Name: ma\_\allowbreak{}2\_\allowbreak{}2\_\allowbreak{}a\_\allowbreak{}3\\
Type: string\\
Default: nil
}

{\footnotesize <MA-02(02)\_\allowbreak{}ODP[03] автоматизовані механізми> використовуються для документування дій з технічного обслуговування, ремонту та заміни системи;}


{\footnotesize\bfseries
No: 7\\
Name: ma\_\allowbreak{}2\_\allowbreak{}2\_\allowbreak{}b\_\allowbreak{}1\\
Type: string\\
Default: nil
}

{\footnotesize надаються актуальні, точні та повні записи про всі замовлені, заплановані, виконувані та завершені дії з технічного обслуговування;}


{\footnotesize\bfseries
No: 8\\
Name: ma\_\allowbreak{}2\_\allowbreak{}2\_\allowbreak{}b\_\allowbreak{}2\\
Type: string\\
Default: nil
}

{\footnotesize надаються актуальні, точні та повні записи про всі замовлені, заплановані, виконувані та завершені дії ремонту;}


{\footnotesize\bfseries
No: 9\\
Name: ma\_\allowbreak{}2\_\allowbreak{}2\_\allowbreak{}b\_\allowbreak{}3\\
Type: string\\
Default: nil
}

{\footnotesize надаються актуальні, точні та повні записи про всі замовлені, заплановані, виконувані та завершені дії з заміни;}




\subsection{ІНСТРУМЕНТИ ДЛЯ ОБСЛУГОВУВАННЯ (MA-3)}
a. Затвердити, контролювати та відстежувати використання засобів технічного обслуговування.\\ 
b. Переглядати раніше затверджені інструменти технічного [Призначення: з частотою, визначеною організацією]. обслуговування

\vspace{10pt}
{\footnotesize\bfseries
No: 1\\
Name: ma\_\allowbreak{}3\_\allowbreak{}odp\\
Type: string\\
Default: nil
}

{\footnotesize визначено частоту, з якою слід переглядати раніше затверджені інструменти технічного обслуговування;}


{\footnotesize\bfseries
No: 2\\
Name: ma\_\allowbreak{}3\_\allowbreak{}a\_\allowbreak{}1\\
Type: string\\
Default: nil
}

{\footnotesize використання засобів технічного обслуговування затверджено;}


{\footnotesize\bfseries
No: 3\\
Name: ma\_\allowbreak{}3\_\allowbreak{}a\_\allowbreak{}2\\
Type: string\\
Default: nil
}

{\footnotesize використання засобів технічного обслуговування контролюється;}


{\footnotesize\bfseries
No: 4\\
Name: ma\_\allowbreak{}3\_\allowbreak{}a\_\allowbreak{}3\\
Type: string\\
Default: nil
}

{\footnotesize використання засобів технічного обслуговування відстажуються;}


{\footnotesize\bfseries
No: 5\\
Name: ma\_\allowbreak{}3\_\allowbreak{}b\\
Type: string\\
Default: nil
}

{\footnotesize переглядаються раніше затверджені інструменти технічного обслуговування <MA-03\_\allowbreak{}ODP частота>.}




\subsubsection{ПЕРЕВІРКА ІНСТРУМЕНТІВ (MA-3(1))}


\vspace{10pt}
{\footnotesize\bfseries
No: 1\\
Name: ma\_\allowbreak{}3\_\allowbreak{}1\_\allowbreak{}01\\
Type: string\\
Default: nil
}

{\footnotesize оглядаються інструменти для технічного обслуговування, які доставлені на об'єкт обслуговуючим персоналом, на предмет неправильних або несанкціонованих модифікацій.}




\subsubsection{ПЕРЕВІРКА НОСІЇВ ІНФОРМАЦІЇ (MA-3(2))}
Перед використанням носіїв у системі перевірити носії, що містять діагностичні та тестові програми на наявність шкідливого коду.

\vspace{10pt}
{\footnotesize\bfseries
No: 1\\
Name: ma\_\allowbreak{}3\_\allowbreak{}2\_\allowbreak{}01\\
Type: string\\
Default: nil
}

{\footnotesize перед використанням носіїв у системі перевіряються носії, що містять діагностичні  та тестові програми на наявність шкідливого коду.}




\subsubsection{ЗАПОБІГАННЯ НЕСАНКЦІОНОВАНОМУ ПЕРЕМІЩЕННЮ (MA-3(3))}
Запобігти переміщенню обладнання для технічного обслуговування, що містить організаційну інформацію, шляхом:\\ 
(a) перевірки відсутності організаційної інформації, розміщеної на обладнанні;\\ 
(b) очищення або знищення обладнання;\\ 
(c) утримання обладнання на об'єкті;\\ 
(d) отримання дозволу від [Призначення: визначених організацією персоналу чи ролей], які явно дозволяють переміщення обладнання з об'єкта.

\vspace{10pt}
{\footnotesize\bfseries
No: 1\\
Name: ma\_\allowbreak{}3\_\allowbreak{}3\_\allowbreak{}odp\\
Type: string\\
Default: nil
}

{\footnotesize визначено персонал або ролі, які можуть надавати дозвіл на переміщення обладнання з об'єкту;}


{\footnotesize\bfseries
No: 2\\
Name: ma\_\allowbreak{}3\_\allowbreak{}3\_\allowbreak{}a\\
Type: string\\
Default: nil
}

{\footnotesize переміщення обладнання для технічного обслуговування, що містить інформацію організації, запобігається шляхом перевірки того, що на обладнанні не міститься ніякої інформації організації; або}


{\footnotesize\bfseries
No: 3\\
Name: ma\_\allowbreak{}3\_\allowbreak{}3\_\allowbreak{}b\\
Type: string\\
Default: nil
}

{\footnotesize переміщення обладнання для технічного обслуговування, що містить інформацію організації, запобігається шляхом очищення або знищення обладнання; або}


{\footnotesize\bfseries
No: 4\\
Name: ma\_\allowbreak{}3\_\allowbreak{}3\_\allowbreak{}c\\
Type: string\\
Default: nil
}

{\footnotesize переміщення обладнання для технічного обслуговування, що містить інформацію організації, запобігається шляхом утримання обладнання на об'єкті; або}


{\footnotesize\bfseries
No: 5\\
Name: ma\_\allowbreak{}3\_\allowbreak{}3\_\allowbreak{}d\\
Type: string\\
Default: nil
}

{\footnotesize переміщення обладнання для технічного обслуговування, що містить інформацію організації, запобігається шляхом отримання дозволу від <MA-03(03)\_\allowbreak{}ODP персонал або ролі>.}




\subsubsection{ОБМЕЖЕННЯ ВИКОРИСТАННЯ ІНСТРУМЕНТА (MA-3(4))}
Обмежити використання інструментів авторизованим персоналом. технічного ОБМЕЖЕННЯ обслуговування лише

\vspace{10pt}
{\footnotesize\bfseries
No: 1\\
Name: ma\_\allowbreak{}3\_\allowbreak{}4\_\allowbreak{}01\\
Type: string\\
Default: nil
}

{\footnotesize обмежено використання інструментів технічного обслуговування лише авторизованим персоналом.}




\subsubsection{ПРИВІЛЕЙОВАНЕ ВИКОНАННЯ (MA-3(5))}
(a) вимагати схвалення кожного віддаленого сеансу технічного обслуговування [Призначення: персоналом або роллю, що визначила організація];\\ 
(b) повідомити [Призначення: персонал або ролі, що визначила організація] про дату та час запланованого віддаленого обслуговування.

\vspace{10pt}
{\footnotesize\bfseries
No: 1\\
Name: ma\_\allowbreak{}3\_\allowbreak{}5\_\allowbreak{}01\\
Type: string\\
Default: nil
}

{\footnotesize відстежується використання інструментів обслуговування, які виконуються з підвищеними привілеями.}




\subsubsection{ОНОВЛЕННЯ ПРОГРАМ- (MA-3(6))}
Запровадити криптографічні механізми для захисту цілісності та конфіденційності віддаленого обслуговування та діагностичних комунікацій.

\vspace{10pt}
{\footnotesize\bfseries
No: 1\\
Name: ma\_\allowbreak{}3\_\allowbreak{}6\_\allowbreak{}01\\
Type: string\\
Default: nil
}

{\footnotesize інструменти технічного обслуговування перевіряються, щоб переконатися, що встановлені найновіші оновлення програмного забезпечення та патчі.}




\subsection{ВІДДАЛЕНЕ ОБСЛУГОВУВАННЯ (MA-4)}


\vspace{10pt}
{\footnotesize\bfseries
No: 1\\
Name: ma\_\allowbreak{}4\_\allowbreak{}a\_\allowbreak{}1\\
Type: string\\
Default: nil
}

{\footnotesize впроваджено віддалені дії з обслуговування та діагностики;}


{\footnotesize\bfseries
No: 2\\
Name: ma\_\allowbreak{}4\_\allowbreak{}a\_\allowbreak{}2\\
Type: string\\
Default: nil
}

{\footnotesize відстежуються віддалені дії з обслуговування та діагностики;}


{\footnotesize\bfseries
No: 3\\
Name: ma\_\allowbreak{}4\_\allowbreak{}b\_\allowbreak{}1\\
Type: string\\
Default: nil
}

{\footnotesize використання віддалених засобів технічного обслуговування та діагностики дозволено лише відповідно до політики організації;}


{\footnotesize\bfseries
No: 4\\
Name: ma\_\allowbreak{}4\_\allowbreak{}b\_\allowbreak{}2\\
Type: string\\
Default: nil
}

{\footnotesize використання віддалених засобів технічного обслуговування та діагностики задокументовано в плані захисту інформації;}


{\footnotesize\bfseries
No: 5\\
Name: ma\_\allowbreak{}4\_\allowbreak{}c\\
Type: string\\
Default: nil
}

{\footnotesize надійна автентифікація використовується при встановленні віддалених технічних та діагностичних сеансів;}


{\footnotesize\bfseries
No: 6\\
Name: ma\_\allowbreak{}4\_\allowbreak{}d\\
Type: string\\
Default: nil
}

{\footnotesize ведеться облік віддалених дій з обслуговування та діагностики;}


{\footnotesize\bfseries
No: 7\\
Name: ma\_\allowbreak{}4\_\allowbreak{}e\_\allowbreak{}1\\
Type: string\\
Default: nil
}

{\footnotesize сесія припиняється, коли завершено віддалене обслуговування}


{\footnotesize\bfseries
No: 8\\
Name: ma\_\allowbreak{}4\_\allowbreak{}e\_\allowbreak{}2\\
Type: string\\
Default: nil
}

{\footnotesize мережеве з'єднання припиняється, коли завершено віддалене 322 обслуговування}




\subsubsection{АУДИТ ТА ОГЛЯД (MA-4(1))}


\vspace{10pt}
{\footnotesize\bfseries
No: 1\\
Name: ma\_\allowbreak{}4\_\allowbreak{}1\_\allowbreak{}odp\_\allowbreak{}1\\
Type: string\\
Default: nil
}

{\footnotesize визначено події аудиту, які слід журналювати  для віддалених сеансів обслуговування;}


{\footnotesize\bfseries
No: 2\\
Name: ma\_\allowbreak{}4\_\allowbreak{}1\_\allowbreak{}odp\_\allowbreak{}2\\
Type: string\\
Default: nil
}

{\footnotesize визначено події аудиту, які слід журналювати  для віддалених сеансів діагностики;}


{\footnotesize\bfseries
No: 3\\
Name: ma\_\allowbreak{}4\_\allowbreak{}1\_\allowbreak{}a\_\allowbreak{}1\\
Type: string\\
Default: nil
}

{\footnotesize <MA-04(01)\_\allowbreak{}ODP[01] події аудиту>  журналюються для віддалених сеансів обслуговування;}


{\footnotesize\bfseries
No: 4\\
Name: ma\_\allowbreak{}4\_\allowbreak{}1\_\allowbreak{}a\_\allowbreak{}2\\
Type: string\\
Default: nil
}

{\footnotesize <MA-04(01)\_\allowbreak{}ODP[02] події аудиту>  журналюються для віддалених сеансів діагностики;}


{\footnotesize\bfseries
No: 5\\
Name: ma\_\allowbreak{}4\_\allowbreak{}1\_\allowbreak{}b\_\allowbreak{}1\\
Type: string\\
Default: nil
}

{\footnotesize здійснюється огляд записів про сеанси віддаленого обслуговування;}


{\footnotesize\bfseries
No: 6\\
Name: ma\_\allowbreak{}4\_\allowbreak{}1\_\allowbreak{}b\_\allowbreak{}2\\
Type: string\\
Default: nil
}

{\footnotesize здійснюється огляд записів про сеанси віддаленої діагностики.}




\subsubsection{ДОКУМЕНТУВАННЯ ВІДДАЛЕНОГО ОБСЛУГОВУВАННЯ (MA-4(2)) [Вилучено]}
[Вилучено: включено до MA-01 та МА-04]

\vspace{10pt}
\textit{Немає параметрів для цього контролю.}
\vspace{10pt}


\subsubsection{ПОРІВНЯЛЬНА БЕЗПЕКА І ОЧИЩЕННЯ (MA-4(3))}


\vspace{10pt}
{\footnotesize\bfseries
No: 1\\
Name: ma\_\allowbreak{}4\_\allowbreak{}3\_\allowbreak{}a\_\allowbreak{}1\\
Type: string\\
Default: nil
}

{\footnotesize віддалені послуги з технічного обслуговування повинні виконуватися з системи, яка реалізує заходи захисту, співставні з заходами захисту, реалізованими в системі, що обслуговується;}


{\footnotesize\bfseries
No: 2\\
Name: ma\_\allowbreak{}4\_\allowbreak{}3\_\allowbreak{}a\_\allowbreak{}2\\
Type: string\\
Default: nil
}

{\footnotesize віддалені послуги з діагностики повинні виконуватися з системи, яка реалізує заходи захисту, співставні з заходами захисту, реалізованими в системі, що обслуговується;}


{\footnotesize\bfseries
No: 3\\
Name: ma\_\allowbreak{}4\_\allowbreak{}3\_\allowbreak{}b\_\allowbreak{}1\\
Type: string\\
Default: nil
}

{\footnotesize компонент, що підлягає обслуговуванню, видаляється з системи перед проведенням віддаленого технічного обслуговування або діагностики;}


{\footnotesize\bfseries
No: 4\\
Name: ma\_\allowbreak{}4\_\allowbreak{}3\_\allowbreak{}b\_\allowbreak{}2\\
Type: string\\
Default: nil
}

{\footnotesize компонент, що підлягає обслуговуванню, пройшов процедуру очищення (від інф ормації  організації);}


{\footnotesize\bfseries
No: 5\\
Name: ma\_\allowbreak{}4\_\allowbreak{}3\_\allowbreak{}b\_\allowbreak{}3\\
Type: string\\
Default: nil
}

{\footnotesize компонент перевіряється та очищується (на наявність потенційно шкідливого програмного забезпечення) після виконання послуги та перед повторним підключенням компонента до системи.}




\subsubsection{СХВАЛЕННЯ ТА ПОВІДОМЛЕННЯ (MA-4(5))}


\vspace{10pt}
{\footnotesize\bfseries
No: 1\\
Name: ma\_\allowbreak{}4\_\allowbreak{}5\_\allowbreak{}odp\_\allowbreak{}1\\
Type: string\\
Default: nil
}

{\footnotesize визначено персонал або ролі, необхідні для затвердження кожного віддаленого сеансу технічного обслуговування; 325}


{\footnotesize\bfseries
No: 2\\
Name: ma\_\allowbreak{}4\_\allowbreak{}5\_\allowbreak{}odp\_\allowbreak{}2\\
Type: string\\
Default: nil
}

{\footnotesize визначено персонал та ролі, які мають бути повідомлені про дату та час запланованого віддаленого технічного обслуговування;}


{\footnotesize\bfseries
No: 3\\
Name: ma\_\allowbreak{}4\_\allowbreak{}5\_\allowbreak{}a\\
Type: string\\
Default: nil
}

{\footnotesize схвалення кожного віддаленого сеансу обслуговування}


{\footnotesize\bfseries
No: 4\\
Name: ma\_\allowbreak{}4\_\allowbreak{}5\_\allowbreak{}b\\
Type: string\\
Default: nil
}

{\footnotesize <MA-04(05)\_\allowbreak{}ODP[02] персонал і ролі> повідомлено про дату і час запланованого віддаленого технічного обслуговування.}




\subsubsection{ВІДДАЛЕНЕ ОБСЛУГОВУВАННЯ РОЗ'ЄДНАННЯ (MA-4(7))}


\vspace{10pt}
{\footnotesize\bfseries
No: 1\\
Name: ma\_\allowbreak{}4\_\allowbreak{}7\_\allowbreak{}1\\
Type: string\\
Default: nil
}

{\footnotesize реалізувано перевірку роз'єднання у разі припинення віддалених сеансів обслуговування}


{\footnotesize\bfseries
No: 2\\
Name: ma\_\allowbreak{}4\_\allowbreak{}7\_\allowbreak{}2\\
Type: string\\
Default: nil
}

{\footnotesize реалізувано перевірку роз'єднання у разі припинення віддалених сеансів діагностики}




\subsection{ТЕХНІЧНИЙ ПЕРСОНАЛ (MA-5)}
a. Встановити процедуру авторизації технічного персоналу та вести перелік авторизованих організацій технічного обслуговування або персоналу.\\ 
b. Перевіряти, що персонал, який не супроводжується та виконує технічне обслуговування в системі, має необхідні дозволи на доступ.\\ 
c. Визначити персонал організації з необхідними повноваженнями щодо доступу та технічною компетенцією для нагляду за персоналом з технічного обслуговування, який не має необхідних дозволів на доступ.

\vspace{10pt}
{\footnotesize\bfseries
No: 1\\
Name: ma\_\allowbreak{}5\_\allowbreak{}a\_\allowbreak{}1\\
Type: string\\
Default: nil
}

{\footnotesize запроваджено процес авторизації технічного персоналу;}


{\footnotesize\bfseries
No: 2\\
Name: ma\_\allowbreak{}5\_\allowbreak{}a\_\allowbreak{}2\\
Type: string\\
Default: nil
}

{\footnotesize ведеться перелік авторизованих організацій або персоналу з технічного обслуговування;}


{\footnotesize\bfseries
No: 3\\
Name: ma\_\allowbreak{}5\_\allowbreak{}b\\
Type: string\\
Default: nil
}

{\footnotesize персонал без супроводу, який виконує технічне обслуговування системи, має необхідні дозволи на доступ;}


{\footnotesize\bfseries
No: 4\\
Name: ma\_\allowbreak{}5\_\allowbreak{}c\\
Type: string\\
Default: nil
}

{\footnotesize персонал організації з необхідними повноваженнями доступу та технічною компетентністю призначений/призначені для нагляду за 327 діяльністю з технічного обслуговування персоналу, який не має необхідних дозволів на доступ.}




\subsubsection{ОСОБИ БЕЗ НАЛЕЖНОГО ДОСТУПУ (MA-5(1))}


\vspace{10pt}
{\footnotesize\bfseries
No: 1\\
Name: ma\_\allowbreak{}5\_\allowbreak{}1\_\allowbreak{}odp\\
Type: string\\
Default: nil
}

{\footnotesize визначені альтернативні заходи захисту, які мають бути розроблені та впроваджені на випадок, якщо компонент системи не може бути очищений, вилучений або відключений від системи;}


{\footnotesize\bfseries
No: 2\\
Name: ma\_\allowbreak{}5\_\allowbreak{}1\_\allowbreak{}a\_\allowbreak{}1\\
Type: string\\
Default: nil
}

{\footnotesize впроваджені процедури залучення персоналу з технічного обслуговування, який не має відповідних дозволів (допуску) або не є громадянами України, містять вимогу: обслуговуючий персонал, що не має необхідних пра в доступу, рівня допуску, або офіційного затвердженого доступу,повинен супроводжуватися та бути під наглядом уповноваженого організацією персоналу, з необхідним рівнем допуску, а також мати відповідну технічну кваліфікацію для виконання технічного обслугов ування та діагностичних заходів у системі;}


{\footnotesize\bfseries
No: 3\\
Name: ma\_\allowbreak{}5\_\allowbreak{}1\_\allowbreak{}a\_\allowbreak{}2\\
Type: string\\
Default: nil
}

{\footnotesize впроваджені процедури залучення персоналу з технічного обслуговування, який не має відповідних дозволів (допуску) або не є громадянами України, містять вимогу: перед тим, як розпочати технічне обслуговування або діагностику персоналом, який не має необхідних прав допуску, рівня допуску або офіційного затвердженого доступу,упевнитися, що всі компоненти енергонезалежного зберігання інформації в системі очищуються, а всі енергонезалежні носії видал яються або фізично відключаються від системи та надійно захищаються;}


{\footnotesize\bfseries
No: 4\\
Name: ma\_\allowbreak{}5\_\allowbreak{}1\_\allowbreak{}b\\
Type: string\\
Default: nil
}

{\footnotesize <MA-05(01)\_\allowbreak{}ODP альтернативні заходи захисту> розробляються і впроваджуються у випадку, якщо систему неможливо очистити, вилучити або відключити від систе328 ми.}




\subsubsection{ОФОРМЛЕННЯ ДОПУСКУ ДЛЯ СИСТЕМ, ЩО ОБРОБЛЯЮТЬ ІНФОРМАЦІЮ З ОБМЕЖЕНИМ ДОСТУПОМ (MA-5(2))}
Переконатися, що персонал, який виконує технічне обслуговування та діагностику в системі, що обробляє, зберігає або передає інформацію з обмеженим доступом, має рівень допуску та офіційне схвалення на доступ для найвищого рівня секретності та для всієї інформації в системі.

\vspace{10pt}
{\footnotesize\bfseries
No: 1\\
Name: ma\_\allowbreak{}5\_\allowbreak{}2\_\allowbreak{}1\\
Type: string\\
Default: nil
}

{\footnotesize персонал, який виконує роботи з технічного обслуговування та діагностики в системі, що обробляє, зберігає або передає інформацію з обмеженим доступом, має відповідний рівень допуску;}


{\footnotesize\bfseries
No: 2\\
Name: ma\_\allowbreak{}5\_\allowbreak{}2\_\allowbreak{}2\\
Type: string\\
Default: nil
}

{\footnotesize персонал, який виконує роботи з технічного обслуговування та діагностики в системі, що обробляє, зберігає або передає інформацію з обмеженим доступом, має офіційне схвалення на допуск;}




\subsubsection{ВИМОГИ ДО ГРОМАДЯНСТВА (MA-5(3))}
Переконатися, що працівники, які виконують технічне обслуговування та діагностичні заходи з обробки, зберігання або передачі таємної інформації, є громадянами України.

\vspace{10pt}
{\footnotesize\bfseries
No: 1\\
Name: ma\_\allowbreak{}5\_\allowbreak{}3\_\allowbreak{}01\\
Type: string\\
Default: nil
}

{\footnotesize працівники, які виконують технічне обслуговування та діагностичні заходи з обробки, зберігання або передачі інформації з обмеженим доступом, є громадянами України.}




\subsubsection{ІНОЗЕМНІ ГРОМАДЯНИ (MA-5(4))}
Переконайтеся, що:\\ 
(a) іноземні громадяни з відповідним рівнем допуску залучаються для проведення технічного обслуговування та діагностичних робіт у системах, що обробляють інформацію з обмеженим доступом тільки тоді, коли ці системи спільно належать і експлуатуються урядами України та закордонних союзників, або належать та експлуатуються виключно іноземними союзними урядами;\\ 
(b) схвалення, згоди та додаткові умови експлуатації, що стосуються залучення іноземних громадян для проведення робіт з технічного обслуговування та діагностики систем, що обробляють інформацію з обмеженим доступом, повністю задокументовані в Меморандумі про угоду.

\vspace{10pt}
{\footnotesize\bfseries
No: 1\\
Name: ma\_\allowbreak{}5\_\allowbreak{}4\_\allowbreak{}a\\
Type: string\\
Default: nil
}

{\footnotesize іноземні громадяни з відповідним рівнем допуску залучаються для проведення технічного обслуговування та діагностичних робіт у системах, що обробляють інформацію з обмеженим доступом тільки тоді, коли ці системи спільно належать і експлуатуються урядами України та закордонних союзників, або належать та е ксплуатуються виключно іноземними союзними урядами;}


{\footnotesize\bfseries
No: 2\\
Name: ma\_\allowbreak{}5\_\allowbreak{}4\_\allowbreak{}b\_\allowbreak{}1\\
Type: string\\
Default: nil
}

{\footnotesize схвалення, що стосуються залучення іноземних громадян для проведення робіт з технічного обслуговування та діагностики систем, що обробляють інформацію з обмеженим доступом, повністю задокументовані в Меморандумі про угоду.}


{\footnotesize\bfseries
No: 3\\
Name: ma\_\allowbreak{}5\_\allowbreak{}4\_\allowbreak{}b\_\allowbreak{}2\\
Type: string\\
Default: nil
}

{\footnotesize згоди, що стосуються залучення іноземних громадян для проведення робіт з технічного обслуговування та діагностики систем, що обробляють інформацію з обмеженим доступом, повністю задокументовані в Меморандумі про угоду.}


{\footnotesize\bfseries
No: 4\\
Name: ma\_\allowbreak{}5\_\allowbreak{}4\_\allowbreak{}b\_\allowbreak{}3\\
Type: string\\
Default: nil
}

{\footnotesize додаткові умови експлуатації, що стосуються залучення іноземних громадян для проведення робіт з технічного обслуговування та діагностики систем, що обробляють інформацію з обмеженим доступом, повністю задокументова ні в Меморандумі про угоду.}




\subsubsection{НЕСИСТЕМНЕ ОБСЛУГОВУВАННЯ (MA-5(5))}
Переконатися, що персонал, який не супроводжується та здійснює ремонтні роботи, не пов'язаний безпосередньо із системою, але перебуває фізично близько від системи, має необхідні дозволи на доступ.

\vspace{10pt}
{\footnotesize\bfseries
No: 1\\
Name: ma\_\allowbreak{}5\_\allowbreak{}5\_\allowbreak{}01\\
Type: string\\
Default: nil
}

{\footnotesize персонал, який не супроводжується, що здійснює ремонтні роботи, не пов'язаний безпосередньо з системою, але знаходиться фізично близько від системи, має необхідні дозволи на доступ.}




\subsection{СВОЄЧАСНЕ ОБСЛУГОВУВАННЯ (MA-6)}
Отримати технічну підтримку та/або запасні частини для [Призначення: визначених організацією компонентів системи] в межах [Призначення: визначеного організацією періоду часу] у разі відмови.

\vspace{10pt}
{\footnotesize\bfseries
No: 1\\
Name: ma\_\allowbreak{}6\_\allowbreak{}odp\_\allowbreak{}1\\
Type: string\\
Default: nil
}

{\footnotesize визначено компоненти системи, для яких отримується технічна підтримка та/або запасні частини;}


{\footnotesize\bfseries
No: 2\\
Name: ma\_\allowbreak{}6\_\allowbreak{}odp\_\allowbreak{}2\\
Type: string\\
Default: nil
}

{\footnotesize визначено період часу, протягом якого можна отримати технічну підтримку та/або запасні частини у разі відмови;}


{\footnotesize\bfseries
No: 3\\
Name: ma\_\allowbreak{}6\_\allowbreak{}01\\
Type: string\\
Default: nil
}

{\footnotesize технічна підтримка та/або запасні частини отримуються для 331}




\subsubsection{ПРОФІЛАКТИЧНЕ ОБСЛУГОВУВАННЯ (MA-6(1))}


\vspace{10pt}
{\footnotesize\bfseries
No: 1\\
Name: ma\_\allowbreak{}6\_\allowbreak{}1\_\allowbreak{}odp\_\allowbreak{}1\\
Type: string\\
Default: nil
}

{\footnotesize визначено компоненти системи яким необхідно здійснювати профілактичне обслуговування;}


{\footnotesize\bfseries
No: 2\\
Name: ma\_\allowbreak{}6\_\allowbreak{}1\_\allowbreak{}odp\_\allowbreak{}2\\
Type: string\\
Default: nil
}

{\footnotesize визначено часові інтервали з якими необхідно здійснювати профілактичне обслуговування визначеним компонентам системи;}


{\footnotesize\bfseries
No: 3\\
Name: ma\_\allowbreak{}6\_\allowbreak{}1\_\allowbreak{}01\\
Type: string\\
Default: nil
}

{\footnotesize здійснюється профілактичне обслуговування <MA06(01)\_\allowbreak{}ODP[01]  компонентів системи>  у  <MA06(01)\_\allowbreak{}ODP[02] часові інтервали>.}




\subsubsection{ПЛАНОВЕ ТЕХНІЧНЕ ОБСЛУГОВУВАННЯ (MA-6(2))}
Здійснювати планове технічне обслуговування [Призначення: визначених організацією компонентів системи] у [Призначення: визначені організацією часові інтервали].

\vspace{10pt}
{\footnotesize\bfseries
No: 1\\
Name: ma\_\allowbreak{}6\_\allowbreak{}2\_\allowbreak{}odp\_\allowbreak{}1\\
Type: string\\
Default: nil
}

{\footnotesize визначено компоненти системи яким необхідно здійснювати планове  технічне обслуговування;}


{\footnotesize\bfseries
No: 2\\
Name: ma\_\allowbreak{}6\_\allowbreak{}2\_\allowbreak{}odp\_\allowbreak{}2\\
Type: string\\
Default: nil
}

{\footnotesize визначено часові інтервали з якими необхідно здійснювати планове  технічне обслуговування;}


{\footnotesize\bfseries
No: 3\\
Name: ma\_\allowbreak{}6\_\allowbreak{}2\_\allowbreak{}01\\
Type: string\\
Default: nil
}

{\footnotesize здійснюється планове  технічне обслуговування <MA06(02)\_\allowbreak{}ODP[01]  компонентів системи>  у  <MA06(02)\_\allowbreak{}ODP[02] часові інтервали>.}




\subsubsection{АВТОМАТИЗОВАНА ПІДТРИМКА ПЛАНОВОГО ТЕХНІЧНОГО ОБСЛУГОВУВАННЯ (MA-6(3))}
Використовувати автоматизовані механізми для передачі даних планового технічного обслуговування до комп'ютеризованої системи управління обслуговуванням [Призначення: автоматизовані засоби визначені організацією].

\vspace{10pt}
{\footnotesize\bfseries
No: 1\\
Name: ma\_\allowbreak{}6\_\allowbreak{}3\_\allowbreak{}odp\\
Type: string\\
Default: nil
}

{\footnotesize визначено автоматизовані механізми для передачі даних планового технічного обслуговування до комп'ютеризованої системи управління обслуговуванням;}


{\footnotesize\bfseries
No: 2\\
Name: ma\_\allowbreak{}6\_\allowbreak{}3\_\allowbreak{}01\\
Type: string\\
Default: nil
}

{\footnotesize використовуються <MA-06(03)\_\allowbreak{}ODP автоматизовані механізми> для передачі даних планового технічного обслуговування до комп'ютеризованої системи управління обслуговуванням.}




\subsection{ТЕХНІЧНЕ ОБСЛУГОВУВАННЯ В ПОЛЬОВИХ УМОВАХ (MA-7)}


\vspace{10pt}
{\footnotesize\bfseries
No: 1\\
Name: ma\_\allowbreak{}7\_\allowbreak{}odp\_\allowbreak{}1\\
Type: string\\
Default: nil
}

{\footnotesize визначені системи або компоненти системи, на яких технічне обслуговування в польових умовах обмежене або заборонене}


{\footnotesize\bfseries
No: 2\\
Name: ma\_\allowbreak{}7\_\allowbreak{}odp\_\allowbreak{}2\\
Type: string\\
Default: nil
}

{\footnotesize визначено довірені засоби технічного обслуговуваннятехнічного обслуговування}


{\footnotesize\bfseries
No: 3\\
Name: ma\_\allowbreak{}7\_\allowbreak{}01\\
Type: string\\
Default: nil
}

{\footnotesize технічне обслуговування в польових умовах < MA-07\_\allowbreak{}ODP[01] систем або компонентів системи > обмежене або заборонене}




\section{MP}
Клас заходів захисту MP --- ЗАХИСТ НОСІЇВ ІНФОРМАЦІЇ

Цей клас спрямований на безпечне зберігання, транспортування, використання та знищення як цифрових, так і паперових носіїв інформації.

\textbf{Перелік заходів захисту:} Політика та процедури щодо захисту носіїв інформації (MP-1); Доступ до носіїв інформації (MP-2); Автоматизований обмеженний доступ (MP-2(1)) [Вилучено]; Криптографічний захист (MP-2(2)) [Вилучено]; Маркування носіїв інформації (MP-3); Зберігання носіїв інформації (MP-4); Криптографічний захист (MP-4(1)) [Вилучено]; Автоматизований обмежений доступ (MP-4(2)); Транспортування носіїв інформації (MP-5); Захист поза контрольованими зонами (MP-5(1)) [Вилучено]; Документування дій (MP-5(2)) [Вилучено]; Зберігачі (MP-5(3)); Криптографічний захист (MP-5(4)) [Вилучено]; Знищення інформації на носіях інформації (MP-6); Переглядати, затвердження, відстеження, документування та перевірка (MP-6(1)); Перевірка обладнання (MP-6(2)); Неруйнівні методи (MP-6(3)); Керована несекретна інформація (MP-6(4)) [Вилучено]; Секретна інформація (MP-6(5)) [Вилучено]; Знищення носіїв інформації (MP-6(6)) [Вилучено]; Подвійна авторизація (MP-6(7)); Віддалене очищення або стирання інформації (MP-6(8)); Використання носіїв інформації (MP-7); Заборона використання без визначеного власника (MP-7(1)) [Вилучено]; Заборона використання (MP-7(2)); Зниження категорії безпеки носіїв інформації (MP-8); Документування процесу (MP-8(1)); Перевірка обладнання (MP-8(2)); Критична інформація (MP-8(3)); Таємна інформація (MP-8(4)).

\subsection{ПОЛІТИКА ТА ПРОЦЕДУРИ ЩОДО ЗАХИСТУ НОСІЇВ ІНФОРМАЦІЇ (MP-1)}
a. Розробити, задокументувати та поширити серед [Призначення: визначеного організацією персоналу або посад]:\\ 
1. 2. політику захисту носіїв інформації, яка:\\ 
(a) містить мету, сферу застосування, ролі, обов'язки, відповідальність керівництва, координацію між організаційними підрозділами та систему контролю відповідності (complaince);\\ 
(b) відповідає чинному законодавству, виконавчим наказам, директивам, нормам, політикам, стандартам та керівним принципам; процедури, які сприяють здійсненню політики та заходів захисту носіїв інформації.\\ 
b. Призначити [Призначення: визначену організацією посадову особу] для управління розробкою, документування, та розповсюдження політики та процедурами захисту носіїв інформації.\\ 
c. Переглядати та оновлювати чинну систему захисту носіїв інформації:\\ 
1. поточну політику захисту носіїв інформації [Призначення: з визначеною організацією частотою];\\ 
2. поточні процедури захисту носіїв інформації [Призначення: з визначеною організацією частотою].

\vspace{10pt}
{\footnotesize\bfseries
No: 1\\
Name: mp\_\allowbreak{}1\_\allowbreak{}odp\_\allowbreak{}1\\
Type: string\\
Default: nil
}

{\footnotesize визначено персонал або ролі, серед яких має бути поширена політика захисту носіїв інформації;}


{\footnotesize\bfseries
No: 2\\
Name: mp\_\allowbreak{}1\_\allowbreak{}odp\_\allowbreak{}2\\
Type: string\\
Default: nil
}

{\footnotesize визначено персонал або ролі, серед яких мають бути поширені процедури захисту носіїв інформації;}


{\footnotesize\bfseries
No: 3\\
Name: mp\_\allowbreak{}1\_\allowbreak{}odp\_\allowbreak{}3\\
Type: string\\
Default: nil
}

{\footnotesize вибрано одне або декілька з наступних ЗНАЧЕНЬ ПАРА- МЕТРІВ: \{рівень організації; рівень місії/бізнес -процесу; рівень системи\};}


{\footnotesize\bfseries
No: 4\\
Name: mp\_\allowbreak{}1\_\allowbreak{}odp\_\allowbreak{}4\\
Type: string\\
Default: nil
}

{\footnotesize визначено посадову особу, яка керуватиме політикою та процедурами захисту носіїв інформації;}


{\footnotesize\bfseries
No: 5\\
Name: mp\_\allowbreak{}1\_\allowbreak{}odp\_\allowbreak{}5\\
Type: string\\
Default: nil
}

{\footnotesize визначено частоту, з якою переглядається та оновлюється поточна політика захисту носіїв інформації;}


{\footnotesize\bfseries
No: 6\\
Name: mp\_\allowbreak{}1\_\allowbreak{}odp\_\allowbreak{}6\\
Type: string\\
Default: nil
}

{\footnotesize визначено події, які потребують перегляду та оновлення чинної політики захисту носіїв інформації;}


{\footnotesize\bfseries
No: 7\\
Name: mp\_\allowbreak{}1\_\allowbreak{}odp\_\allowbreak{}7\\
Type: string\\
Default: nil
}

{\footnotesize визначено частоту, з якою переглядаються та оновлюються чинні процедури захисту носіїв інформації;}


{\footnotesize\bfseries
No: 8\\
Name: mp\_\allowbreak{}1\_\allowbreak{}odp\_\allowbreak{}8\\
Type: string\\
Default: nil
}

{\footnotesize визначено події, які потребують перегляду та оновлення процедур захисту носіїв інформації;}


{\footnotesize\bfseries
No: 9\\
Name: mp\_\allowbreak{}1\_\allowbreak{}a\_\allowbreak{}1\\
Type: string\\
Default: nil
}

{\footnotesize розроблено та задокументовано політику захисту носіїв інформації;}


{\footnotesize\bfseries
No: 10\\
Name: mp\_\allowbreak{}1\_\allowbreak{}a\_\allowbreak{}2\\
Type: string\\
Default: nil
}

{\footnotesize політика захисту носіїв інформації поширюється на  <MP01\_\allowbreak{}ODP[01] персонал або ролі>;}


{\footnotesize\bfseries
No: 11\\
Name: mp\_\allowbreak{}1\_\allowbreak{}a\_\allowbreak{}3\\
Type: string\\
Default: nil
}

{\footnotesize розроблено та задокументовано процедури захисту носіїв інформації, що сприятимуть реалізації політики захисту носіїв інформації та заходів захисту носіїв інформації;}


{\footnotesize\bfseries
No: 12\\
Name: mp\_\allowbreak{}1\_\allowbreak{}a\_\allowbreak{}4\\
Type: string\\
Default: nil
}

{\footnotesize процедури захисту носіїв інформації поширюються на <MP01\_\allowbreak{}ODP[02] персонал або ролі>;}


{\footnotesize\bfseries
No: 13\\
Name: mp\_\allowbreak{}1\_\allowbreak{}a\_\allowbreak{}1\_\allowbreak{}a\_\allowbreak{}1\\
Type: string\\
Default: nil
}

{\footnotesize <MP-01\_\allowbreak{}ODP[03] ВИБРАНЕ ЗНАЧЕННЯ ПАРАМЕТРА(ів)> політика захисту носіїв інформації містить мету; 335}


{\footnotesize\bfseries
No: 14\\
Name: mp\_\allowbreak{}1\_\allowbreak{}a\_\allowbreak{}1\_\allowbreak{}a\_\allowbreak{}2\\
Type: string\\
Default: nil
}

{\footnotesize <MP-01\_\allowbreak{}ODP[03] ВИБРАНЕ ЗНАЧЕННЯ ПАРАМЕТРА(ів)> політика захисту носіїв інформації містить сферу застосування;}


{\footnotesize\bfseries
No: 15\\
Name: mp\_\allowbreak{}1\_\allowbreak{}a\_\allowbreak{}1\_\allowbreak{}a\_\allowbreak{}3\\
Type: string\\
Default: nil
}

{\footnotesize <MP-01\_\allowbreak{}ODP[03] ВИБРАНЕ ЗНАЧЕННЯ ПАРАМЕТРА(ів)> політика захисту носіїв інформації містить ролі;}


{\footnotesize\bfseries
No: 16\\
Name: mp\_\allowbreak{}1\_\allowbreak{}a\_\allowbreak{}1\_\allowbreak{}a\_\allowbreak{}4\\
Type: string\\
Default: nil
}

{\footnotesize <MP-01\_\allowbreak{}ODP[03] ВИБРАНЕ ЗНАЧЕННЯ ПАРАМЕТРА(ів)> політика захисту носіїв інформації містить обов'язки;}


{\footnotesize\bfseries
No: 17\\
Name: mp\_\allowbreak{}1\_\allowbreak{}a\_\allowbreak{}1\_\allowbreak{}a\_\allowbreak{}5\\
Type: string\\
Default: nil
}

{\footnotesize <MP-01\_\allowbreak{}ODP[03] ВИБРАНЕ ЗНАЧЕННЯ ПАРАМЕТРА(ів)> політика захисту носіїв інформації містить відповідальність керівництва;}


{\footnotesize\bfseries
No: 18\\
Name: mp\_\allowbreak{}1\_\allowbreak{}a\_\allowbreak{}1\_\allowbreak{}a\_\allowbreak{}6\\
Type: string\\
Default: nil
}

{\footnotesize <MP-01\_\allowbreak{}ODP[03] ВИБРАНЕ ЗНАЧЕННЯ ПАРАМЕТРА(ів)> політика захисту носіїв інформації містить координацію між підрозділами організації;}


{\footnotesize\bfseries
No: 19\\
Name: mp\_\allowbreak{}1\_\allowbreak{}a\_\allowbreak{}1\_\allowbreak{}a\_\allowbreak{}7\\
Type: string\\
Default: nil
}

{\footnotesize <MP-01\_\allowbreak{}ODP[03] ВИБРАНЕ ЗНАЧЕННЯ ПАРАМЕТРА(ів)> політика захисту носіїв інформації містить систему контролю відповідності;}


{\footnotesize\bfseries
No: 20\\
Name: mp\_\allowbreak{}1\_\allowbreak{}a\_\allowbreak{}1\_\allowbreak{}b\\
Type: string\\
Default: nil
}

{\footnotesize політика захисту носіїв інформації відповідає чинному законодавству, виконавчим наказам, директивам, нормам, політикам, стандартам та керівним принципам;}


{\footnotesize\bfseries
No: 21\\
Name: mp\_\allowbreak{}1\_\allowbreak{}b\\
Type: string\\
Default: nil
}

{\footnotesize <MP-01\_\allowbreak{}ODP[04] посадова особа> призначається для управління розробкою, документуванням та розповсюдженням політики та процедур захисту носіїв інформації.}


{\footnotesize\bfseries
No: 22\\
Name: mp\_\allowbreak{}1\_\allowbreak{}c\_\allowbreak{}1\_\allowbreak{}1\\
Type: string\\
Default: nil
}

{\footnotesize переглядається та оновлюється поточна політика захисту носіїв}


{\footnotesize\bfseries
No: 23\\
Name: mp\_\allowbreak{}1\_\allowbreak{}c\_\allowbreak{}1\_\allowbreak{}2\\
Type: string\\
Default: nil
}

{\footnotesize переглядається та оновлюється поточна політика захисту носіїв}


{\footnotesize\bfseries
No: 24\\
Name: mp\_\allowbreak{}1\_\allowbreak{}c\_\allowbreak{}2\_\allowbreak{}1\\
Type: string\\
Default: nil
}

{\footnotesize переглядаються та оновлюються поточні процедури захисту}


{\footnotesize\bfseries
No: 25\\
Name: mp\_\allowbreak{}1\_\allowbreak{}c\_\allowbreak{}2\_\allowbreak{}2\\
Type: string\\
Default: nil
}

{\footnotesize переглядаються та оновлюються поточні процедури захисту}




\subsection{ДОСТУП ДО НОСІЇВ ІНФОРМАЦІЇ (MP-2)}
Обмежити доступ до [Призначення: визначених організацією типів цифрових та/або нецифрових носіїв інформації] [Призначення: визначеним організацією персоналом або ролями].

\vspace{10pt}
{\footnotesize\bfseries
No: 1\\
Name: mp\_\allowbreak{}2\_\allowbreak{}odp\_\allowbreak{}1\\
Type: string\\
Default: nil
}

{\footnotesize визначено типи цифрових носіїв інформації, доступ до  яких обмежено;}


{\footnotesize\bfseries
No: 2\\
Name: mp\_\allowbreak{}2\_\allowbreak{}odp\_\allowbreak{}2\\
Type: string\\
Default: nil
}

{\footnotesize визначено персонал або ролі, уповноважені на доступ до цифрових носіїв інформації;}


{\footnotesize\bfseries
No: 3\\
Name: mp\_\allowbreak{}2\_\allowbreak{}odp\_\allowbreak{}3\\
Type: string\\
Default: nil
}

{\footnotesize визначено типи нецифрових носіїв інформації, доступ до яких обмежено;}


{\footnotesize\bfseries
No: 4\\
Name: mp\_\allowbreak{}2\_\allowbreak{}odp\_\allowbreak{}4\\
Type: string\\
Default: nil
}

{\footnotesize визначено персонал або ролі, уповноважені на доступ до нецифрових носіїв інформації;}


{\footnotesize\bfseries
No: 5\\
Name: mp\_\allowbreak{}2\_\allowbreak{}1\\
Type: string\\
Default: nil
}

{\footnotesize доступ до  <MP-02\_\allowbreak{}ODP[01] типів цифрових носіїв інформації> обмежено для <MP-02\_\allowbreak{}ODP[02] персоналу або ролей>;}


{\footnotesize\bfseries
No: 6\\
Name: mp\_\allowbreak{}2\_\allowbreak{}2\\
Type: string\\
Default: nil
}

{\footnotesize доступ до  <MP-02\_\allowbreak{}ODP[03] типів нецифрових носіїв інформації> обмежено для <MP-02\_\allowbreak{}ODP[04] персоналу або ролей>;}




\subsubsection{АВТОМАТИЗОВАНИЙ ОБМЕЖЕННИЙ ДОСТУП (MP-2(1)) [Вилучено]}
[Вилучено: Включено до MP-04(02)]

\vspace{10pt}
\textit{Немає параметрів для цього контролю.}
\vspace{10pt}


\subsubsection{КРИПТОГРАФІЧНИЙ ЗАХИСТ (MP-2(2)) [Вилучено]}
[Вилучено: Включено до SC-28(01)]

\vspace{10pt}
\textit{Немає параметрів для цього контролю.}
\vspace{10pt}


\subsection{МАРКУВАННЯ НОСІЇВ ІНФОРМАЦІЇ (MP-3)}
a. Наносити на носії інформації маркування, що вказують на обмеження поширення, обробки, а також застереження та відповідні мітки безпеки (якщо такі є) інформації.\\ 
b. Звільнити [Призначення: визначені організацією типи носіїв системи] від маркування, якщо носії залишаються в межах [Призначення: визначених організацією контрольованих зон].

\vspace{10pt}
{\footnotesize\bfseries
No: 1\\
Name: mp\_\allowbreak{}3\_\allowbreak{}odp\_\allowbreak{}1\\
Type: string\\
Default: nil
}

{\footnotesize визначено типи носіїв  інформації, які звільняються від маркування під час перебування на контрольованих зонах;}


{\footnotesize\bfseries
No: 2\\
Name: mp\_\allowbreak{}3\_\allowbreak{}odp\_\allowbreak{}2\\
Type: string\\
Default: nil
}

{\footnotesize визначено контрольовані зони, де носії інформації звільняються від маркування;}


{\footnotesize\bfseries
No: 3\\
Name: mp\_\allowbreak{}3\_\allowbreak{}a\\
Type: string\\
Default: nil
}

{\footnotesize носії інформації маркуються, щоб вказати на обмеження поширення, обробки, а також застереження та відповідні мітки безпеки (якщо такі є) інформації;}


{\footnotesize\bfseries
No: 4\\
Name: mp\_\allowbreak{}3\_\allowbreak{}b\\
Type: string\\
Default: nil
}

{\footnotesize <MP-03\_\allowbreak{}ODP[01] типи носіїв інформації> залишаються в}




\subsection{ЗБЕРІГАННЯ НОСІЇВ ІНФОРМАЦІЇ (MP-4)}
a. Фізично контролювати та безпечно зберігати [Призначення: визначені організацією типи цифрових та/або нецифрових носіїв інформації] в межах [Призначення: визначених організацією контрольованих зон].\\ 
b. Захищати системні носії, які визначені в МР-4 до того часу, як носії знищуються або очищаються, з використанням затвердженого обладнання, методів та процедур.

\vspace{10pt}
{\footnotesize\bfseries
No: 1\\
Name: mp\_\allowbreak{}4\_\allowbreak{}odp\_\allowbreak{}1\\
Type: string\\
Default: nil
}

{\footnotesize визначено типи цифрових носіїв інформації, які підлягають фізичному контролю (якщо вибрано);}


{\footnotesize\bfseries
No: 2\\
Name: mp\_\allowbreak{}4\_\allowbreak{}odp\_\allowbreak{}2\\
Type: string\\
Default: nil
}

{\footnotesize визначено типи нецифрових носіїв інформації, які підлягають фізичному контролю (якщо вибрано);}


{\footnotesize\bfseries
No: 3\\
Name: mp\_\allowbreak{}4\_\allowbreak{}odp\_\allowbreak{}3\\
Type: string\\
Default: nil
}

{\footnotesize визначено типи цифрових носіїв інформації для безпечного зберігання (якщо вибрано);}


{\footnotesize\bfseries
No: 4\\
Name: mp\_\allowbreak{}4\_\allowbreak{}odp\_\allowbreak{}4\\
Type: string\\
Default: nil
}

{\footnotesize визначено типи нецифрових носіїв інформації для безпечного зберігання (якщо вибрано);}


{\footnotesize\bfseries
No: 5\\
Name: mp\_\allowbreak{}4\_\allowbreak{}odp\_\allowbreak{}5\\
Type: string\\
Default: nil
}

{\footnotesize визначено контрольовані зони, в яких можна безпечно зберігати цифрові носії інформації; 338}


{\footnotesize\bfseries
No: 6\\
Name: mp\_\allowbreak{}4\_\allowbreak{}odp\_\allowbreak{}6\\
Type: string\\
Default: nil
}

{\footnotesize визначено контрольовані зони, в яких можна безпечно зберігати нецифрові носії інформації;}


{\footnotesize\bfseries
No: 7\\
Name: mp\_\allowbreak{}4\_\allowbreak{}a\_\allowbreak{}1\\
Type: string\\
Default: nil
}

{\footnotesize <MP-04\_\allowbreak{}ODP[01] типи цифрових носіїв> контролюються фізично;}


{\footnotesize\bfseries
No: 8\\
Name: mp\_\allowbreak{}4\_\allowbreak{}a\_\allowbreak{}2\\
Type: string\\
Default: nil
}

{\footnotesize <<MP-04\_\allowbreak{}ODP[02] типи нецифрових носіїв> контролюються фізично;}


{\footnotesize\bfseries
No: 9\\
Name: mp\_\allowbreak{}4\_\allowbreak{}a\_\allowbreak{}3\\
Type: string\\
Default: nil
}

{\footnotesize <MP-04\_\allowbreak{}ODP[03] типи цифрових носіїв> безпечно зберігаються}


{\footnotesize\bfseries
No: 10\\
Name: mp\_\allowbreak{}4\_\allowbreak{}a\_\allowbreak{}4\\
Type: string\\
Default: nil
}

{\footnotesize <MP-04\_\allowbreak{}ODP[04] типи нецифрових носіїв> безпечно}


{\footnotesize\bfseries
No: 11\\
Name: mp\_\allowbreak{}4\_\allowbreak{}b\\
Type: string\\
Default: nil
}

{\footnotesize типи носіїв інформації (визначені в MP -04\_\allowbreak{}ODP[01], MP 04\_\allowbreak{}ODP[02], MP -04\_\allowbreak{}ODP[03], MP -04\_\allowbreak{}ODP[04]) захищені доти, доки носії інформації не будуть знищені або очищені за допомогою визначеного обладнання, методик та процедур.}




\subsubsection{КРИПТОГРАФІЧНИЙ ЗАХИСТ (MP-4(1)) [Вилучено]}
[Вилучено: Включено до SC-28(01)]

\vspace{10pt}
\textit{Немає параметрів для цього контролю.}
\vspace{10pt}


\subsubsection{АВТОМАТИЗОВАНИЙ ОБМЕЖЕНИЙ ДОСТУП (MP-4(2))}


\vspace{10pt}
{\footnotesize\bfseries
No: 1\\
Name: mp\_\allowbreak{}4\_\allowbreak{}2\_\allowbreak{}odp\_\allowbreak{}1\\
Type: string\\
Default: nil
}

{\footnotesize визначено автоматизовані механізми обмеження доступу до зон зберігання носіїв інформації;}


{\footnotesize\bfseries
No: 2\\
Name: mp\_\allowbreak{}4\_\allowbreak{}2\_\allowbreak{}odp\_\allowbreak{}2\\
Type: string\\
Default: nil
}

{\footnotesize визначено автоматизовані механізми реєстрації спроб доступу до зон зберігання носіїв інформації; 339}


{\footnotesize\bfseries
No: 3\\
Name: mp\_\allowbreak{}4\_\allowbreak{}2\_\allowbreak{}odp\_\allowbreak{}3\\
Type: string\\
Default: nil
}

{\footnotesize визначено автоматизовані механізми реєстрації доступу, наданого до зон зберігання носіїв інформації;}


{\footnotesize\bfseries
No: 4\\
Name: mp\_\allowbreak{}4\_\allowbreak{}2\_\allowbreak{}1\\
Type: string\\
Default: nil
}

{\footnotesize доступ до зон зберігання носіїв інформації обмежено за}


{\footnotesize\bfseries
No: 5\\
Name: mp\_\allowbreak{}4\_\allowbreak{}2\_\allowbreak{}2\\
Type: string\\
Default: nil
}

{\footnotesize спроби доступу до зон зберігання носіїв інфор мації автоматизованих механізмів>;}


{\footnotesize\bfseries
No: 6\\
Name: mp\_\allowbreak{}4\_\allowbreak{}2\_\allowbreak{}3\\
Type: string\\
Default: nil
}

{\footnotesize доступ, наданий до зон зберігання носіїв, реєструється за}




\subsection{ТРАНСПОРТУВАННЯ НОСІЇВ ІНФОРМАЦІЇ (MP-5)}
a. Захищати та контролювати [Призначення: визначені організацією типи носіїв системи] під час транспортування за межами контрольованих зон, використовуючи [Призначення: визначені організацією заходи безпеки].\\ 
b. Вести облік носіїв системи інформації під час транспортування за межами контрольованих зон.\\ 
c. Документувати дії, пов'язані з транспортуванням носіїв системи.\\ 
d. Обмежити діяльність уповноваженого персоналу, пов'язану з транспортуванням носіїв системи.

\vspace{10pt}
{\footnotesize\bfseries
No: 1\\
Name: mp\_\allowbreak{}5\_\allowbreak{}odp\_\allowbreak{}1\\
Type: string\\
Default: nil
}

{\footnotesize визначено типи носіїв інформації системи для захисту та контролю під час транспортування за межі контрольованих зон;}


{\footnotesize\bfseries
No: 2\\
Name: mp\_\allowbreak{}5\_\allowbreak{}odp\_\allowbreak{}2\\
Type: string\\
Default: nil
}

{\footnotesize визначено заходи безпеки, що використовуються для захисту носіїв інформації системи поза контрольованими зонами;}


{\footnotesize\bfseries
No: 3\\
Name: mp\_\allowbreak{}5\_\allowbreak{}odp\_\allowbreak{}3\\
Type: string\\
Default: nil
}

{\footnotesize визначено заходи безпеки, що використовуються для контролю носіїв інформації системи за межами контрольованих зон;}


{\footnotesize\bfseries
No: 4\\
Name: mp\_\allowbreak{}5\_\allowbreak{}a\_\allowbreak{}1\\
Type: string\\
Default: nil
}

{\footnotesize <MP-05\_\allowbreak{}ODP[01] типи носіїв інформації системи>  захищаються під час транспортування за межі контрольованих зон за допомогою}


{\footnotesize\bfseries
No: 5\\
Name: mp\_\allowbreak{}5\_\allowbreak{}a\_\allowbreak{}2\\
Type: string\\
Default: nil
}

{\footnotesize <MP-05\_\allowbreak{}ODP[01] типи носіїв інформації системи> 340 контролюються під час транспортування за межі контрольованих}


{\footnotesize\bfseries
No: 6\\
Name: mp\_\allowbreak{}5\_\allowbreak{}b\\
Type: string\\
Default: nil
}

{\footnotesize під час транспортування за межі контрольованих зон ведеться облік носіїв інформації системи;}


{\footnotesize\bfseries
No: 7\\
Name: mp\_\allowbreak{}5\_\allowbreak{}c\\
Type: string\\
Default: nil
}

{\footnotesize діяльність, пов'язана з транспортуванням носіїв інформації системи, задокументована;}


{\footnotesize\bfseries
No: 8\\
Name: mp\_\allowbreak{}5\_\allowbreak{}d\_\allowbreak{}1\\
Type: string\\
Default: nil
}

{\footnotesize визначено персонал, уповноважений здійснювати діяльність з транспортування носіїв інформації;}


{\footnotesize\bfseries
No: 9\\
Name: mp\_\allowbreak{}5\_\allowbreak{}d\_\allowbreak{}2\\
Type: string\\
Default: nil
}

{\footnotesize діяльність, пов'язана з транспортуванням носіїв інформації системи, обмежується визначеним уповноваженим персоналом.}




\subsubsection{ЗАХИСТ ПОЗА КОНТРОЛЬОВАНИМИ ЗОНАМИ (MP-5(1)) [Вилучено]}
[Вилучено: Включено до MP-05]

\vspace{10pt}
\textit{Немає параметрів для цього контролю.}
\vspace{10pt}


\subsubsection{ДОКУМЕНТУВАННЯ ДІЙ (MP-5(2)) [Вилучено]}
[Вилучено: Включено до MP-05]

\vspace{10pt}
\textit{Немає параметрів для цього контролю.}
\vspace{10pt}


\subsubsection{ЗБЕРІГАЧІ (MP-5(3))}


\vspace{10pt}
{\footnotesize\bfseries
No: 1\\
Name: mp\_\allowbreak{}5\_\allowbreak{}3\_\allowbreak{}1\\
Type: string\\
Default: nil
}

{\footnotesize визначено зберігачів інформації під час транспортування носіїв інформації системи за межі контрольованих зон.}


{\footnotesize\bfseries
No: 2\\
Name: mp\_\allowbreak{}5\_\allowbreak{}3\_\allowbreak{}2\\
Type: string\\
Default: nil
}

{\footnotesize залучено визначених зберігачів інформації під час транспорту341 вання носіїв інформації системи за межі контрольованих зон.}




\subsubsection{КРИПТОГРАФІЧНИЙ ЗАХИСТ (MP-5(4)) [Вилучено]}
[Вилучено: Включено до SC-28(01)]

\vspace{10pt}
\textit{Немає параметрів для цього контролю.}
\vspace{10pt}


\subsection{ЗНИЩЕННЯ ІНФОРМАЦІЇ НА НОСІЯХ ІНФОРМАЦІЇ (MP-6)}
a. Очищувати [Призначення: визначені організацією системні носії] перед утилізацією, випуском за межі організаційного контролю, або перед повторним використанням [Призначення: методами та процедурами очищення, визначеними організацією].\\ 
b. Використовувати механізми очищення зі стійкістю та цілісністю, що відповідає категорії безпеки або рівню секретності інформації.

\vspace{10pt}
{\footnotesize\bfseries
No: 1\\
Name: mp\_\allowbreak{}6\_\allowbreak{}odp\_\allowbreak{}1\\
Type: string\\
Default: nil
}

{\footnotesize визначено носії інформації системи, які підлягають очищенню перед утилізацією;}


{\footnotesize\bfseries
No: 2\\
Name: mp\_\allowbreak{}6\_\allowbreak{}odp\_\allowbreak{}2\\
Type: string\\
Default: nil
}

{\footnotesize визначено носії інформації системи, які підлягають очищенню перед випуском за межі контрольованої зони;}


{\footnotesize\bfseries
No: 3\\
Name: mp\_\allowbreak{}6\_\allowbreak{}odp\_\allowbreak{}3\\
Type: string\\
Default: nil
}

{\footnotesize визначені носії інформації системи, що підлягають очищенню перед повторним використанням;}


{\footnotesize\bfseries
No: 4\\
Name: mp\_\allowbreak{}6\_\allowbreak{}odp\_\allowbreak{}4\\
Type: string\\
Default: nil
}

{\footnotesize визначено методи та процедури очищення, які слід використовувати для очищення перед утилізацією;}


{\footnotesize\bfseries
No: 5\\
Name: mp\_\allowbreak{}6\_\allowbreak{}odp\_\allowbreak{}5\\
Type: string\\
Default: nil
}

{\footnotesize визначено методи та процедури очищення, які слід використовувати для очищення перед вип уском за межі контрольованої зони;}


{\footnotesize\bfseries
No: 6\\
Name: mp\_\allowbreak{}6\_\allowbreak{}odp\_\allowbreak{}6\\
Type: string\\
Default: nil
}

{\footnotesize визначено методи та процедури очищення, які слід використовувати для очищення перед повторним використанням;}


{\footnotesize\bfseries
No: 7\\
Name: mp\_\allowbreak{}6\_\allowbreak{}a\_\allowbreak{}1\\
Type: string\\
Default: nil
}

{\footnotesize <MP-06\_\allowbreak{}ODP[01] носії інформації системи>  перед утилізацією та процедур очищення>;}


{\footnotesize\bfseries
No: 8\\
Name: mp\_\allowbreak{}6\_\allowbreak{}a\_\allowbreak{}2\\
Type: string\\
Default: nil
}

{\footnotesize <MP-06\_\allowbreak{}ODP[02] носії інформації системи>  очищаються за перед випуском за межі контрольованої зони;}


{\footnotesize\bfseries
No: 9\\
Name: mp\_\allowbreak{}6\_\allowbreak{}a\_\allowbreak{}3\\
Type: string\\
Default: nil
}

{\footnotesize <MP-06\_\allowbreak{}ODP[03] носії інформації системи>  очищуються за 342 перед повторним використанням;}


{\footnotesize\bfseries
No: 10\\
Name: mp\_\allowbreak{}6\_\allowbreak{}b\\
Type: string\\
Default: nil
}

{\footnotesize застосовуються механізми очищення, надійність і цілісність яких відповідає категорії безпеки або рівню секретності інформації.}




\subsubsection{ПЕРЕГЛЯДАТИ, ЗАТВЕРДЖЕННЯ, ВІДСТЕЖЕННЯ, ДОКУМЕНТУВАННЯ ТА ПЕРЕВІРКА (MP-6(1))}


\vspace{10pt}
{\footnotesize\bfseries
No: 1\\
Name: mp\_\allowbreak{}6\_\allowbreak{}1\_\allowbreak{}1\\
Type: string\\
Default: nil
}

{\footnotesize переглядаються заходи з очищення та утилізації носіїв інформації;}


{\footnotesize\bfseries
No: 2\\
Name: mp\_\allowbreak{}6\_\allowbreak{}1\_\allowbreak{}2\\
Type: string\\
Default: nil
}

{\footnotesize затверджуються заходи з очищення та утилізації носіїв інформації;}


{\footnotesize\bfseries
No: 3\\
Name: mp\_\allowbreak{}6\_\allowbreak{}1\_\allowbreak{}3\\
Type: string\\
Default: nil
}

{\footnotesize відстежуються заходи з очищення та утилізації носіїв інформації;}


{\footnotesize\bfseries
No: 4\\
Name: mp\_\allowbreak{}6\_\allowbreak{}1\_\allowbreak{}4\\
Type: string\\
Default: nil
}

{\footnotesize документуються заходи з очищення та утилізації носіїв інформації;}


{\footnotesize\bfseries
No: 5\\
Name: mp\_\allowbreak{}6\_\allowbreak{}1\_\allowbreak{}5\\
Type: string\\
Default: nil
}

{\footnotesize перевіряються заходи з очищення та утилізації носіїв інформації;}




\subsubsection{ПЕРЕВІРКА ОБЛАДНАННЯ (MP-6(2))}


\vspace{10pt}
{\footnotesize\bfseries
No: 1\\
Name: mp\_\allowbreak{}6\_\allowbreak{}2\_\allowbreak{}odp\_\allowbreak{}1\\
Type: string\\
Default: nil
}

{\footnotesize визначена частота, з якою проводиться перевірка обладнання для очищення;}


{\footnotesize\bfseries
No: 2\\
Name: mp\_\allowbreak{}6\_\allowbreak{}2\_\allowbreak{}odp\_\allowbreak{}2\\
Type: string\\
Default: nil
}

{\footnotesize визначено частоту, з якою потрібно перевіряти процедури очищення;}


{\footnotesize\bfseries
No: 3\\
Name: mp\_\allowbreak{}6\_\allowbreak{}2\_\allowbreak{}1\\
Type: string\\
Default: nil
}

{\footnotesize обладнання для очищення тестується  <MP06(02)\_\allowbreak{}ODP[01] частота> , щоб переконатися в досягненні запланованого очищення;}


{\footnotesize\bfseries
No: 4\\
Name: mp\_\allowbreak{}6\_\allowbreak{}2\_\allowbreak{}2\\
Type: string\\
Default: nil
}

{\footnotesize процедури санітарної обробки тестуються <MP06(02)\_\allowbreak{}ODP[02] частота> , щоб переконатися в досягненні запланованого очищення.}




\subsubsection{НЕРУЙНІВНІ МЕТОДИ (MP-6(3))}


\vspace{10pt}
{\footnotesize\bfseries
No: 1\\
Name: mp\_\allowbreak{}6\_\allowbreak{}3\_\allowbreak{}odp\\
Type: string\\
Default: nil
}

{\footnotesize визначено умови, які  вимагають очищення зовнішніх носіїв інформації;}


{\footnotesize\bfseries
No: 2\\
Name: mp\_\allowbreak{}6\_\allowbreak{}3\_\allowbreak{}01\\
Type: string\\
Default: nil
}

{\footnotesize застосовуються методи неруйнівного очищення до зовнішніх носіїв інформації перед підключенням таких пристроїв до 344 системи при <MP-06(03)\_\allowbreak{}ODP умовах>, що вимагають очищення зовнішніх носіїв інформації.}




\subsubsection{КЕРОВАНА НЕСЕКРЕТНА ІНФОРМАЦІЯ (MP-6(4)) [Вилучено]}
[Вилучено: Включено до MP-06]

\vspace{10pt}
\textit{Немає параметрів для цього контролю.}
\vspace{10pt}


\subsubsection{СЕКРЕТНА ІНФОРМАЦІЯ (MP-6(5)) [Вилучено]}
[Вилучено: Включено до MP-06]

\vspace{10pt}
\textit{Немає параметрів для цього контролю.}
\vspace{10pt}


\subsubsection{ЗНИЩЕННЯ НОСІЇВ ІНФОРМАЦІЇ (MP-6(6)) [Вилучено]}
[Вилучено: Включено до MP-06]

\vspace{10pt}
\textit{Немає параметрів для цього контролю.}
\vspace{10pt}


\subsubsection{ПОДВІЙНА АВТОРИЗАЦІЯ (MP-6(7))}


\vspace{10pt}
{\footnotesize\bfseries
No: 1\\
Name: mp\_\allowbreak{}6\_\allowbreak{}7\_\allowbreak{}odp\\
Type: string\\
Default: nil
}

{\footnotesize визначено носії інформації для яких необхідно здійснювати подвійну авторизацію для очищення;}


{\footnotesize\bfseries
No: 2\\
Name: mp\_\allowbreak{}6\_\allowbreak{}7\_\allowbreak{}01\\
Type: string\\
Default: nil
}

{\footnotesize здійснюється подвійна авторизація для очищення <MP06(07)\_\allowbreak{}ODP носії інформації>.}




\subsubsection{ВІДДАЛЕНЕ ОЧИЩЕННЯ АБО СТИРАННЯ ІНФОРМАЦІЇ (MP-6(8))}


\vspace{10pt}
{\footnotesize\bfseries
No: 1\\
Name: mp\_\allowbreak{}6\_\allowbreak{}8\_\allowbreak{}odp\_\allowbreak{}1\\
Type: string\\
Default: nil
}

{\footnotesize визначено системи або компоненти системи для очищення або стирання інформації віддалено або за певних умов;}


{\footnotesize\bfseries
No: 2\\
Name: mp\_\allowbreak{}6\_\allowbreak{}8\_\allowbreak{}odp\_\allowbreak{}2\\
Type: string\\
Default: nil
}

{\footnotesize вибрано одне з наступних ЗНАЧЕНЬ ПАРАМЕТРІВ:}


{\footnotesize\bfseries
No: 3\\
Name: mp\_\allowbreak{}6\_\allowbreak{}8\_\allowbreak{}odp\_\allowbreak{}3\\
Type: string\\
Default: nil
}

{\footnotesize визначаються умови, за яких інформація підлягає очищенню або стиранню (якщо вибрано);}


{\footnotesize\bfseries
No: 4\\
Name: mp\_\allowbreak{}6\_\allowbreak{}8\_\allowbreak{}01\\
Type: string\\
Default: nil
}

{\footnotesize передбачено можливість очищення або стирання інформації з <MP-06(08) \_\allowbreak{} ODP[01] систем або компонентів}




\subsection{ВИКОРИСТАННЯ НОСІЇВ ІНФОРМАЦІЇ (MP-7)}
a. [Вибір: обмежити; заборонити] використання [Призначення: визначених організацією типів носіїв системи] на [Призначення: визначені організацією системи або компоненти системи], використовуючи [Призначення: визначені організацією заходи безпеки].\\ 
b. Заборонити використання портативних пристроїв зберігання даних в системах організації, якщо такі пристрої не мають визначеного власника.

\vspace{10pt}
{\footnotesize\bfseries
No: 1\\
Name: mp\_\allowbreak{}7\_\allowbreak{}b\\
Type: string\\
Default: nil
}

{\footnotesize використання зовнішніх  носіїв інформації в системах організації заборонено, якщо такі пристрої не мають власника, якого можна ідентифікувати.}




\subsubsection{ЗАБОРОНА ВИКОРИСТАННЯ БЕЗ ВИЗНАЧЕНОГО ВЛАСНИКА (MP-7(1)) [Вилучено]}
[Вилучено: Включено до MP-07]

\vspace{10pt}
\textit{Немає параметрів для цього контролю.}
\vspace{10pt}


\subsubsection{ЗАБОРОНА ВИКОРИСТАННЯ (MP-7(2))}


\vspace{10pt}
{\footnotesize\bfseries
No: 1\\
Name: mp\_\allowbreak{}7\_\allowbreak{}2\_\allowbreak{}1\\
Type: string\\
Default: nil
}

{\footnotesize ідентифіковано стійкі до очищення носії інформації;}


{\footnotesize\bfseries
No: 2\\
Name: mp\_\allowbreak{}7\_\allowbreak{}2\_\allowbreak{}2\\
Type: string\\
Default: nil
}

{\footnotesize використання стійких до очищення носіїв інформації в системах організації заборонено.}




\subsection{ЗНИЖЕННЯ КАТЕГОРІЇ БЕЗПЕКИ НОСІЇВ ІНФОРМАЦІЇ (MP-8)}
a. сплануйте розташування або ділянку об'єкта, де знаходиться система, враховуючи фізичні та екологічні ризики;\\ 
b. для існуючих об'єктів врахуйте фізичні та екологічні ризики в організаційній стратегії управління ризиками.

\vspace{10pt}
{\footnotesize\bfseries
No: 1\\
Name: mp\_\allowbreak{}8\_\allowbreak{}odp\_\allowbreak{}1\\
Type: string\\
Default: nil
}

{\footnotesize визначено процес зниження категорії безпеки носіїв інформації;}


{\footnotesize\bfseries
No: 2\\
Name: mp\_\allowbreak{}8\_\allowbreak{}odp\_\allowbreak{}2\\
Type: string\\
Default: nil
}

{\footnotesize визначено носії інформації системи, що вимагають зниження категорії безпеки;}


{\footnotesize\bfseries
No: 3\\
Name: mp\_\allowbreak{}8\_\allowbreak{}a\_\allowbreak{}1\\
Type: string\\
Default: nil
}

{\footnotesize встановлено <MP-08\_\allowbreak{}ODP[01] процес зниження категорії безпеки носіїв інформаці>;}


{\footnotesize\bfseries
No: 4\\
Name: mp\_\allowbreak{}8\_\allowbreak{}a\_\allowbreak{}2\\
Type: string\\
Default: nil
}

{\footnotesize <MP-08\_\allowbreak{}ODP[01] процес зниження категорії безпеки носіїв інформаці> охоплює використання механізмів зниження грифа секретності носіїв інформації за стійкістю та цілісністю, що відповідає категорії безпеки або рівню секретності інформації;}


{\footnotesize\bfseries
No: 5\\
Name: mp\_\allowbreak{}8\_\allowbreak{}b\_\allowbreak{}1\\
Type: string\\
Default: nil
}

{\footnotesize здійснюється перевірка того, що процес зниження категорії безпеки носіїв інформації відповідає категорії безпеки та/або рівню секретності інформації, що підлягає видаленню;}


{\footnotesize\bfseries
No: 6\\
Name: mp\_\allowbreak{}8\_\allowbreak{}b\_\allowbreak{}2\\
Type: string\\
Default: nil
}

{\footnotesize здійснюється перевірка того, що процес зниження категорії безпеки носіїв інформації співмірний з правами доступу потенційних одержувачів пониженої інформації; 348}


{\footnotesize\bfseries
No: 7\\
Name: mp\_\allowbreak{}8\_\allowbreak{}c\\
Type: string\\
Default: nil
}

{\footnotesize визначено <MP-08\_\allowbreak{}ODP[02] носії інформації системи, що потребують пониження статусу>;}


{\footnotesize\bfseries
No: 8\\
Name: mp\_\allowbreak{}8\_\allowbreak{}d\\
Type: string\\
Default: nil
}

{\footnotesize визначений носій інформації понижено у категорії безпеки за безпеки носіїв інформації>.}




\subsubsection{ДОКУМЕНТУВАННЯ ПРОЦЕСУ (MP-8(1))}


\vspace{10pt}
{\footnotesize\bfseries
No: 1\\
Name: mp\_\allowbreak{}8\_\allowbreak{}1\_\allowbreak{}01\\
Type: string\\
Default: nil
}

{\footnotesize документуються дії зі зниження категорії безпеки носіїв інформації.}




\subsubsection{ПЕРЕВІРКА ОБЛАДНАННЯ (MP-8(2))}


\vspace{10pt}
{\footnotesize\bfseries
No: 1\\
Name: mp\_\allowbreak{}8\_\allowbreak{}2\_\allowbreak{}odp\_\allowbreak{}1\\
Type: string\\
Default: nil
}

{\footnotesize визначено частоту, з якою потрібно перевіряти обладнання для заниження категорії безпеки ;}


{\footnotesize\bfseries
No: 2\\
Name: mp\_\allowbreak{}8\_\allowbreak{}2\_\allowbreak{}odp\_\allowbreak{}2\\
Type: string\\
Default: nil
}

{\footnotesize визначено частоту, з якою потрібно перевіряти процедури для заниження категорії безпеки ;}


{\footnotesize\bfseries
No: 3\\
Name: mp\_\allowbreak{}8\_\allowbreak{}2\_\allowbreak{}1\\
Type: string\\
Default: nil
}

{\footnotesize обладнання для заниження категорії безпеки перевіряється}


{\footnotesize\bfseries
No: 4\\
Name: mp\_\allowbreak{}8\_\allowbreak{}2\_\allowbreak{}2\\
Type: string\\
Default: nil
}

{\footnotesize процедури для заниження категорії безпеки перевіряється}




\subsubsection{КРИТИЧНА ІНФОРМАЦІЯ (MP-8(3))}


\vspace{10pt}
{\footnotesize\bfseries
No: 1\\
Name: mp\_\allowbreak{}8\_\allowbreak{}3\_\allowbreak{}1\\
Type: string\\
Default: nil
}

{\footnotesize визначено критичну інформацію за наявності якої на носії інформації, знижують категорію безпеки до рівня публічного доступу.}


{\footnotesize\bfseries
No: 2\\
Name: mp\_\allowbreak{}8\_\allowbreak{}3\_\allowbreak{}2\\
Type: string\\
Default: nil
}

{\footnotesize знижується категорія безпеки носіїв  інформації, що містять визначену критичну інформацію до рівня публічного доступу.}




\subsubsection{ТАЄМНА ІНФОРМАЦІЯ (MP-8(4))}


\vspace{10pt}
{\footnotesize\bfseries
No: 1\\
Name: mp\_\allowbreak{}8\_\allowbreak{}4\_\allowbreak{}1\\
Type: string\\
Default: nil
}

{\footnotesize ідентифіковано носії інформації, що містять інформацію з обмеженим доступом;}


{\footnotesize\bfseries
No: 2\\
Name: mp\_\allowbreak{}8\_\allowbreak{}4\_\allowbreak{}2\\
Type: string\\
Default: nil
}

{\footnotesize носії інформації, що містять інформацію з обмеженим доступом, знижуються в класі перед передачею особам, які не мають необхідних дозволів на доступ.}




\section{PE}
Клас заходів захисту PE --- ФІЗИЧНИЙ ЗАХИСТ І ЗАХИСТ РОБОЧОГО

Цей клас охоплює заходи контролю фізичного доступу до об'єктів організації та захисту обладнання від загроз навколишнього середовища.

\textbf{Перелік заходів захисту:} Політика та процедури фізичного захисту та захисту робочого середовища (PE-1); Авторизація фізичного доступу (PE-2); Доступ на основі посади або ролі (PE-2(1)); Дві форми ідентифікації (PE-2(2)); Безперервна охорона (PE-3(3)); Шафи з блокуванням (PE-3(4)); Керування фізичним доступом (PE-3); Доступ до системи (PE-3(1)); Межі об'єкту та системи (PE-3(2)); Керування фізичним доступом --- захист від злому (PE-3(5)); Тестування на можливість проникнення (PE-3(6)); Фізичні перешкоди (PE-3(7)); КОНТРОЛЬ ДОСТУПУ У ВЕСТИБЮЛІ (ХОЛІ) (PE-3(8)); Ліній електроживлення (PE-4); Контроль доступу в приміщення для відображення інформації (PE-5); Доступ до вихідних даних уповноваженими особами (PE-5(1)) [Вилучено]; Доступ до вихідних даних фізичними особами (PE-5(2)); Маркування пристроїв виведення інформації (PE-5(3)) [Вилучено]; Моніторинг фізичного доступу (PE-6); Охоронна сигналізація та обладнання для спостереження (PE-6(1)); Моніторинг фізичного доступу --- автоматичні розпізнавання вторгнень і відповідна реакція (PE-6(2)); Відеоспостереження (PE-6(3)); Моніторинг фізичного доступу до системи (PE-6(4)); Контроль відвідувачів (PE-7); Реєстр доступу відвідувачів (PE-8); Реєстри доступу відвідувачів --- обмеження інформації, що ідентифікує особу (PE-8(3)); Енергетичне обладнання та кабелі (PE-9); Резервні кабелі (PE-9(1)); Автоматичне керування напругою (PE-9(2)); Аварійне відключення (PE-10); Випадкова і несанкціонована активація (PE-10(1)) [Вилучено]; Аварійне енергозабезпечення (PE-11); Довгострокове альтернативне джерело живлення - мінімальні експлуатаційні можливості (PE-11(1)); Довгострокове альтернативне джерело живлення -- автономне живлення (PE-11(2)); Аварійне освітлення (PE-12); Основні завдання та функції (PE-12(1)); Протипожежний захист (PE-13); Пристрої та системи виявлення (PE-13(1)); Пристрої та системи автоматичного пожежогасіння (PE-13(2)); Автоматичне пожежогасіння (PE-13(3)) [Вилучено]; Перевірки (PE-13(4)); Контроль температури та вологості (PE-14); Автоматичний контроль (PE-14(1)); Моніторинг за допомогою сигналізацій та сповіщень (PE-14(2)); Захист від пошкодження водою (PE-15); Автоматична підтримка (PE-15(1)); Доставка та видалення (PE-16); Альтернативне робоче місце (PE-17); Розташування компонентів системи (PE-18); Місце розміщення об'єкта (PE-18(1)) [Вилучено]; Витік інформації (PE-19); Національні політики та процедури щодо пемв (PE-19(1)); Моніторинг і відстеження активів (PE-20); Захист від електромагнітного імпульсу (PE-21); Маркування компонентів (PE-22); Розташування об'єкта (PE-23).

\subsection{ПОЛІТИКА ТА ПРОЦЕДУРИ ФІЗИЧНОГО ЗАХИСТУ ТА ЗАХИСТУ РОБОЧОГО СЕРЕДОВИЩА (PE-1)}


\vspace{10pt}
{\footnotesize\bfseries
No: 1\\
Name: pe\_\allowbreak{}1\_\allowbreak{}odp\_\allowbreak{}1\\
Type: string\\
Default: nil
}

{\footnotesize визначено персонал або ролі, до яких має бути доведена політика фізичного захисту та захисту робочого середовища;}


{\footnotesize\bfseries
No: 2\\
Name: pe\_\allowbreak{}1\_\allowbreak{}odp\_\allowbreak{}2\\
Type: string\\
Default: nil
}

{\footnotesize визначено персонал або ролі, на які поширюються процедури фізичного захисту та захисту робочого середовища;}


{\footnotesize\bfseries
No: 3\\
Name: pe\_\allowbreak{}1\_\allowbreak{}odp\_\allowbreak{}3\\
Type: string\\
Default: nil
}

{\footnotesize вибрано одне або декілька з наступних ЗНАЧЕНЬ ПАРА- МЕТРІВ: \{рівень організації; рівень місії/бізнес-процесу; рівень системи\};}


{\footnotesize\bfseries
No: 4\\
Name: pe\_\allowbreak{}1\_\allowbreak{}odp\_\allowbreak{}4\\
Type: string\\
Default: nil
}

{\footnotesize визначено посадову особу, яка керуватиме політикою та процедурами фізичного захисту та захисту робочого середовища;}


{\footnotesize\bfseries
No: 5\\
Name: pe\_\allowbreak{}1\_\allowbreak{}odp\_\allowbreak{}5\\
Type: string\\
Default: nil
}

{\footnotesize визначено частоту, з якою переглядається та оновлюється поточна політика фізичного захисту та захисту робочого середовища;}


{\footnotesize\bfseries
No: 6\\
Name: pe\_\allowbreak{}1\_\allowbreak{}odp\_\allowbreak{}6\\
Type: string\\
Default: nil
}

{\footnotesize визначено події, які потребують перегляду та оновлення поточної політики фізичного захисту та захисту робочого середовища;}


{\footnotesize\bfseries
No: 7\\
Name: pe\_\allowbreak{}1\_\allowbreak{}odp\_\allowbreak{}7\\
Type: string\\
Default: nil
}

{\footnotesize визначено частоту, з якою переглядаються  та оновлюються поточні процедури фізичного захисту та захисту робочого середовища;}


{\footnotesize\bfseries
No: 8\\
Name: pe\_\allowbreak{}1\_\allowbreak{}odp\_\allowbreak{}8\\
Type: string\\
Default: nil
}

{\footnotesize визначено події, які потребують перегляду та оновлення процедур фізичного захисту та захисту робочого середовища;}


{\footnotesize\bfseries
No: 9\\
Name: pe\_\allowbreak{}1\_\allowbreak{}a\_\allowbreak{}1\\
Type: string\\
Default: nil
}

{\footnotesize розроблено та задокументо вано політику фізичного захисту та захисту робочого середовища;}


{\footnotesize\bfseries
No: 10\\
Name: pe\_\allowbreak{}1\_\allowbreak{}a\_\allowbreak{}2\\
Type: string\\
Default: nil
}

{\footnotesize політика фізичного захисту та захисту робочого середовища}


{\footnotesize\bfseries
No: 11\\
Name: pe\_\allowbreak{}1\_\allowbreak{}a\_\allowbreak{}3\\
Type: string\\
Default: nil
}

{\footnotesize розроблені та задокументовані процедури фізичного захисту та захисту робочого середовища, що сприяють впровадженню політики фізичного захисту та захисту робочого середовища, а також пов'язані з ними заходи захисту; 352}


{\footnotesize\bfseries
No: 12\\
Name: pe\_\allowbreak{}1\_\allowbreak{}a\_\allowbreak{}4\\
Type: string\\
Default: nil
}

{\footnotesize процедури фізичного захисту та захисту робочого середовища}


{\footnotesize\bfseries
No: 13\\
Name: pe\_\allowbreak{}1\_\allowbreak{}a\_\allowbreak{}1\_\allowbreak{}a\_\allowbreak{}1\\
Type: string\\
Default: nil
}

{\footnotesize <PE-01\_\allowbreak{}ODP[03] ВИБРАНЕ ЗНАЧЕННЯ ПАРАМЕТРА(ів)> політика фізичного захисту та захисту робочого середовища містить мету;}


{\footnotesize\bfseries
No: 14\\
Name: pe\_\allowbreak{}1\_\allowbreak{}a\_\allowbreak{}1\_\allowbreak{}a\_\allowbreak{}2\\
Type: string\\
Default: nil
}

{\footnotesize <PE-01\_\allowbreak{}ODP[03] ВИБРАНЕ ЗНАЧЕННЯ ПАРАМЕТРА(ів)> політика фізичного захисту та захисту робочого середовища містить сферу застосування;}


{\footnotesize\bfseries
No: 15\\
Name: pe\_\allowbreak{}1\_\allowbreak{}a\_\allowbreak{}1\_\allowbreak{}a\_\allowbreak{}3\\
Type: string\\
Default: nil
}

{\footnotesize <PE-01\_\allowbreak{}ODP[03] ВИБРАНЕ ЗНАЧЕННЯ ПАРАМЕТРА(ів)> політика фізичного захисту та захисту робочого середовища містить ролі;}


{\footnotesize\bfseries
No: 16\\
Name: pe\_\allowbreak{}1\_\allowbreak{}a\_\allowbreak{}1\_\allowbreak{}a\_\allowbreak{}4\\
Type: string\\
Default: nil
}

{\footnotesize <PE-01\_\allowbreak{}ODP[03] ВИБРАНЕ ЗНАЧЕННЯ ПАРАМЕТРА(ів)> політика фізичного захисту та захисту робочого середовища містить обов'язки;}


{\footnotesize\bfseries
No: 17\\
Name: pe\_\allowbreak{}1\_\allowbreak{}a\_\allowbreak{}1\_\allowbreak{}a\_\allowbreak{}5\\
Type: string\\
Default: nil
}

{\footnotesize <PE-01\_\allowbreak{}ODP[03] ВИБРАНЕ ЗНАЧЕННЯ ПАРАМЕТРА(ів)> політика фізичного захисту та захисту робочого середовища містить відповідальність керівництва;}


{\footnotesize\bfseries
No: 18\\
Name: pe\_\allowbreak{}1\_\allowbreak{}a\_\allowbreak{}1\_\allowbreak{}a\_\allowbreak{}6\\
Type: string\\
Default: nil
}

{\footnotesize <PE-01\_\allowbreak{}ODP[03] ВИБРАНЕ ЗНАЧЕННЯ ПАРАМЕТРА(ів)> політика фізичного захисту та захисту робочого середовища містить координацію між підрозділами організації;}


{\footnotesize\bfseries
No: 19\\
Name: pe\_\allowbreak{}1\_\allowbreak{}a\_\allowbreak{}1\_\allowbreak{}a\_\allowbreak{}7\\
Type: string\\
Default: nil
}

{\footnotesize <PE-01\_\allowbreak{}ODP[03] ВИБРАНЕ ЗНАЧЕННЯ ПАРАМЕТРА(ів)> політика фізичного захисту та захисту робочого середовища містить систему контролю відповідності;}


{\footnotesize\bfseries
No: 20\\
Name: pe\_\allowbreak{}1\_\allowbreak{}a\_\allowbreak{}1\_\allowbreak{}b\\
Type: string\\
Default: nil
}

{\footnotesize <PE-01\_\allowbreak{}ODP[03] ВИБРАНЕ ЗНАЧЕННЯ ПАРАМЕТРА(ів)> політика фізичного захисту та захисту робочого середовища відповідає чинному законодавству, виконавчим наказам, директивам, нормам, політикам, стандартам і керівним принципам;}


{\footnotesize\bfseries
No: 21\\
Name: pe\_\allowbreak{}1\_\allowbreak{}b\\
Type: string\\
Default: nil
}

{\footnotesize <PE-01\_\allowbreak{}ODP[04] посадова особа>  призначається для управління розробкою, документуванням та розповсюдженням політики та процедур фізичного захисту та захисту робочого середовища;}


{\footnotesize\bfseries
No: 22\\
Name: pe\_\allowbreak{}1\_\allowbreak{}c\_\allowbreak{}1\_\allowbreak{}1\\
Type: string\\
Default: nil
}

{\footnotesize переглядається та оновлюється поточна політика фізичного захисту та захисту робочого середовища <PE-01\_\allowbreak{}ODP[05] частота>;}


{\footnotesize\bfseries
No: 23\\
Name: pe\_\allowbreak{}1\_\allowbreak{}c\_\allowbreak{}1\_\allowbreak{}2\\
Type: string\\
Default: nil
}

{\footnotesize поточна політика фізичного захисту та захисту робочого середовища переглядається та оновлюється після <PE-01\_\allowbreak{}ODP[06] подій>;}


{\footnotesize\bfseries
No: 24\\
Name: pe\_\allowbreak{}1\_\allowbreak{}c\_\allowbreak{}2\_\allowbreak{}1\\
Type: string\\
Default: nil
}

{\footnotesize переглядаються та оновлюються поточні процедури фізичного 353}


{\footnotesize\bfseries
No: 25\\
Name: pe\_\allowbreak{}1\_\allowbreak{}c\_\allowbreak{}2\_\allowbreak{}2\\
Type: string\\
Default: nil
}

{\footnotesize поточні процедури фізичного та екологічного захисту переглядаються та оновлюються після <PE-01\_\allowbreak{}ODP[08] подій>.}




\subsection{АВТОРИЗАЦІЯ ФІЗИЧНОГО ДОСТУПУ (PE-2)}


\vspace{10pt}
{\footnotesize\bfseries
No: 1\\
Name: pe\_\allowbreak{}2\_\allowbreak{}odp\\
Type: string\\
Default: nil
}

{\footnotesize визначено періодичність перегляду списку доступу, у якому закріплений перелік персоналу або ролей, яким дозволений санкціонований доступ до об'єкта;}


{\footnotesize\bfseries
No: 2\\
Name: pe\_\allowbreak{}2\_\allowbreak{}a\_\allowbreak{}1\\
Type: string\\
Default: nil
}

{\footnotesize розроблено перелік осіб, які мають право авторизованого доступу до об'єкта, де перебуває система;}


{\footnotesize\bfseries
No: 3\\
Name: pe\_\allowbreak{}2\_\allowbreak{}a\_\allowbreak{}2\\
Type: string\\
Default: nil
}

{\footnotesize затверджено перелік осіб, які мають право авторизованого доступу до об'єкта, де перебуває система;}


{\footnotesize\bfseries
No: 4\\
Name: pe\_\allowbreak{}2\_\allowbreak{}a\_\allowbreak{}3\\
Type: string\\
Default: nil
}

{\footnotesize ведеться перелік осіб, які мають право  авторизованого доступу до об'єкта, де перебуває система;}


{\footnotesize\bfseries
No: 5\\
Name: pe\_\allowbreak{}2\_\allowbreak{}b\\
Type: string\\
Default: nil
}

{\footnotesize для доступу до об'єкта надаються повноваження;}


{\footnotesize\bfseries
No: 6\\
Name: pe\_\allowbreak{}2\_\allowbreak{}c\\
Type: string\\
Default: nil
}

{\footnotesize переглядається список доступу, , у якому закріплений перелік персоналу або ролей, яким дозволений санкціонований доступ до об'єкта <PE-02\_\allowbreak{}ODP частота>;}


{\footnotesize\bfseries
No: 7\\
Name: pe\_\allowbreak{}2\_\allowbreak{}d\\
Type: string\\
Default: nil
}

{\footnotesize особи видаляються зі списку доступу до об'єкта, коли доступ більше не потрібен.}




\subsubsection{ДОСТУП НА ОСНОВІ ПОСАДИ АБО РОЛІ (PE-2(1))}


\vspace{10pt}
\textit{Немає параметрів для цього контролю.}
\vspace{10pt}


\subsubsection{ДВІ ФОРМИ ІДЕНТИФІКАЦІЇ (PE-2(2))}


\vspace{10pt}
{\footnotesize\bfseries
No: 1\\
Name: pe\_\allowbreak{}2\_\allowbreak{}2\_\allowbreak{}odp\\
Type: string\\
Default: nil
}

{\footnotesize визначено список прийнятних форм ідентифікації}


{\footnotesize\bfseries
No: 2\\
Name: pe\_\allowbreak{}2\_\allowbreak{}2\_\allowbreak{}01\\
Type: string\\
Default: nil
}

{\footnotesize вимагається дві форми ідентифікації від <PE-02(02)\_\allowbreak{}ODP списку прийнятних форм ідентифікації>  для доступу відвідувачів до об'єкта, де знаходиться система.}




\subsection{КЕРУВАННЯ ФІЗИЧНИМ ДОСТУПОМ (PE-3)}


\vspace{10pt}
{\footnotesize\bfseries
No: 1\\
Name: pe\_\allowbreak{}3\_\allowbreak{}odp\_\allowbreak{}1\\
Type: string\\
Default: nil
}

{\footnotesize визначено точки входу та виходу в об'єкт, в якому знаходиться система;}


{\footnotesize\bfseries
No: 2\\
Name: pe\_\allowbreak{}3\_\allowbreak{}odp\_\allowbreak{}2\\
Type: string\\
Default: nil
}

{\footnotesize вибрано одне або декілька з наступних ЗНАЧЕНЬ ПАРА-}


{\footnotesize\bfseries
No: 3\\
Name: pe\_\allowbreak{}3\_\allowbreak{}odp\_\allowbreak{}3\\
Type: string\\
Default: nil
}

{\footnotesize визначено фізичні системи або пристрої контролю доступу, що використовуються для контролю входу та виходу на об'єкт (якщо вибрано); 356}


{\footnotesize\bfseries
No: 4\\
Name: pe\_\allowbreak{}3\_\allowbreak{}odp\_\allowbreak{}4\\
Type: string\\
Default: nil
}

{\footnotesize визначено точки входу або виходу, для яких ведуться журнали контролю фізичного доступу;}


{\footnotesize\bfseries
No: 5\\
Name: pe\_\allowbreak{}3\_\allowbreak{}odp\_\allowbreak{}5\\
Type: string\\
Default: nil
}

{\footnotesize визначено заходи захисту для контролю доступу в зони всередині об'єкту, позначені як загальнодоступні;}


{\footnotesize\bfseries
No: 6\\
Name: pe\_\allowbreak{}3\_\allowbreak{}odp\_\allowbreak{}6\\
Type: string\\
Default: nil
}

{\footnotesize визначено умови, що вимагають супроводу відвідувачів та моніторингу активності відвідувачів;}


{\footnotesize\bfseries
No: 7\\
Name: pe\_\allowbreak{}3\_\allowbreak{}odp\_\allowbreak{}7\\
Type: string\\
Default: nil
}

{\footnotesize визначені пристрої фізичного доступу, що підлягають інвентаризації;}


{\footnotesize\bfseries
No: 8\\
Name: pe\_\allowbreak{}3\_\allowbreak{}odp\_\allowbreak{}8\\
Type: string\\
Default: nil
}

{\footnotesize визначено частоту проведення інвентаризації пристроїв фізичного доступу;}


{\footnotesize\bfseries
No: 9\\
Name: pe\_\allowbreak{}3\_\allowbreak{}odp\_\allowbreak{}9\\
Type: string\\
Default: nil
}

{\footnotesize визначено частоту, з якою потрібно змінювати коди доступу;}


{\footnotesize\bfseries
No: 10\\
Name: pe\_\allowbreak{}3\_\allowbreak{}odp\_\allowbreak{}10\\
Type: string\\
Default: nil
}

{\footnotesize визначено частоту, з якою потрібно змінювати ключі;}


{\footnotesize\bfseries
No: 11\\
Name: pe\_\allowbreak{}3\_\allowbreak{}a\_\allowbreak{}1\\
Type: string\\
Default: nil
}

{\footnotesize авторизація фізичного доступу забезпечується в <PE-03\_\allowbreak{}ODP[01] пунктах входу і виходу>  шляхом перевірки індивідуальних дозволів доступу;}


{\footnotesize\bfseries
No: 12\\
Name: pe\_\allowbreak{}3\_\allowbreak{}a\_\allowbreak{}2\\
Type: string\\
Default: nil
}

{\footnotesize авторизація фізичного доступу здійснюється у <PE-03\_\allowbreak{}ODP[01] точках входу та виходу>  шляхом управління входом та виходом ЗНАЧЕННЯ ПАРАМЕТРА(ів)>;}


{\footnotesize\bfseries
No: 13\\
Name: pe\_\allowbreak{}3\_\allowbreak{}b\\
Type: string\\
Default: nil
}

{\footnotesize журнали контролю фізичного доступу ведуться для <PE03\_\allowbreak{}ODP[04] точок входу або виходу>;}


{\footnotesize\bfseries
No: 14\\
Name: pe\_\allowbreak{}3\_\allowbreak{}c\\
Type: string\\
Default: nil
}

{\footnotesize доступ в зони всередині об'єкту, визначені як загальнодоступні, захисту>;}


{\footnotesize\bfseries
No: 15\\
Name: pe\_\allowbreak{}3\_\allowbreak{}d\_\allowbreak{}1\\
Type: string\\
Default: nil
}

{\footnotesize відвідувачів супроводжують;}


{\footnotesize\bfseries
No: 16\\
Name: pe\_\allowbreak{}3\_\allowbreak{}d\_\allowbreak{}2\\
Type: string\\
Default: nil
}

{\footnotesize активність відвідувачів контролюється <PE-03\_\allowbreak{}ODP[06] умови>;}


{\footnotesize\bfseries
No: 17\\
Name: pe\_\allowbreak{}3\_\allowbreak{}e\_\allowbreak{}1\\
Type: string\\
Default: nil
}

{\footnotesize ключі захищені;}


{\footnotesize\bfseries
No: 18\\
Name: pe\_\allowbreak{}3\_\allowbreak{}e\_\allowbreak{}2\\
Type: string\\
Default: nil
}

{\footnotesize коди доступу захищені;}


{\footnotesize\bfseries
No: 19\\
Name: pe\_\allowbreak{}3\_\allowbreak{}e\_\allowbreak{}3\\
Type: string\\
Default: nil
}

{\footnotesize інші пристрої фізичного доступу захищені;}


{\footnotesize\bfseries
No: 20\\
Name: pe\_\allowbreak{}3\_\allowbreak{}f\\
Type: string\\
Default: nil
}

{\footnotesize <PE-03\_\allowbreak{}ODP[07] пристрої фізичного доступу>  інвентаризуються}


{\footnotesize\bfseries
No: 21\\
Name: pe\_\allowbreak{}3\_\allowbreak{}g\_\allowbreak{}1\\
Type: string\\
Default: nil
}

{\footnotesize коди доступу змінюється <PE-03\_\allowbreak{}ODP[09] частота> , коли код скомпрометовано, або коли особи, які володіють кодом, переводяться або звільняються; 357}


{\footnotesize\bfseries
No: 22\\
Name: pe\_\allowbreak{}3\_\allowbreak{}g\_\allowbreak{}2\\
Type: string\\
Default: nil
}

{\footnotesize ключі змінюються <PE-03\_\allowbreak{}ODP[10] частота>, коли ключі втрачено, або коли особи, що володіють ключами, переводяться або звільняються.}




\subsubsection{ДОСТУП ДО СИСТЕМИ (PE-3(1))}


\vspace{10pt}
{\footnotesize\bfseries
No: 1\\
Name: pe\_\allowbreak{}3\_\allowbreak{}1\_\allowbreak{}odp\\
Type: string\\
Default: nil
}

{\footnotesize визначено фізичні приміщення, що містять один або декілька компонентів системи;}


{\footnotesize\bfseries
No: 2\\
Name: pe\_\allowbreak{}3\_\allowbreak{}1\_\allowbreak{}1\\
Type: string\\
Default: nil
}

{\footnotesize фізичні авторизації доступу до системи є обов'язковими;}


{\footnotesize\bfseries
No: 3\\
Name: pe\_\allowbreak{}3\_\allowbreak{}1\_\allowbreak{}2\\
Type: string\\
Default: nil
}

{\footnotesize для об'єкта застосовуються засоби контролю фізичного доступу в <PE-03(01)\_\allowbreak{}ODP фізичні приміщення>.}




\subsubsection{МЕЖІ ОБ'ЄКТУ ТА СИСТЕМИ (PE-3(2))}


\vspace{10pt}
{\footnotesize\bfseries
No: 1\\
Name: pe\_\allowbreak{}3\_\allowbreak{}2\_\allowbreak{}odp\\
Type: string\\
Default: nil
}

{\footnotesize не визначено частоту проведення перевірок безпеки на фізичній межі об'єкта або системи на предмет витоку інформації або вилучення компонентів системи;}


{\footnotesize\bfseries
No: 2\\
Name: pe\_\allowbreak{}3\_\allowbreak{}2\_\allowbreak{}01\\
Type: string\\
Default: nil
}

{\footnotesize перевірки безпеки проводяться  <PE-03(02)\_\allowbreak{}ODP частота> на фізичному периметрі об'єкта або системи на предмет витоку інформації або вилучення компонентів системи.}




\subsubsection{БЕЗПЕРЕРВНА ОХОРОНА (PE-3(3))}


\vspace{10pt}
{\footnotesize\bfseries
No: 1\\
Name: pe\_\allowbreak{}2\_\allowbreak{}3\_\allowbreak{}odp\_\allowbreak{}1\\
Type: string\\
Default: nil
}

{\footnotesize вибрано одне або декілька з наступних ЗНАЧЕНЬ ПАРАМЕТРІВ: \{рівень допуску для всієї інформації, що міститься в системі; авторизація офіційного доступу до всієї інформації, що міститься в системі; необхідність}


{\footnotesize\bfseries
No: 2\\
Name: pe\_\allowbreak{}2\_\allowbreak{}3\_\allowbreak{}odp\_\allowbreak{}2\\
Type: string\\
Default: nil
}

{\footnotesize визначено повноваження фізичного доступу для доступу без супроводу до об'єкта, де знаходиться система (якщо вибрано);}


{\footnotesize\bfseries
No: 3\\
Name: pe\_\allowbreak{}2\_\allowbreak{}3\_\allowbreak{}01\\
Type: string\\
Default: nil
}

{\footnotesize доступ без супроводу до приміщення, де знаходиться система, обмежено для персоналу з <PE-02(03)\_\allowbreak{}ODP[01] ВИБРАНИМ ЗНАЧЕННЯМ ПАРАМЕТРА(ів)>.}




\subsubsection{ШАФИ З БЛОКУВАННЯМ (PE-3(4))}


\vspace{10pt}
\textit{Немає параметрів для цього контролю.}
\vspace{10pt}


\subsubsection{КЕРУВАННЯ ФІЗИЧНИМ ДОСТУПОМ --- ЗАХИСТ ВІД ЗЛОМУ (PE-3(5))}


\vspace{10pt}
{\footnotesize\bfseries
No: 1\\
Name: pe\_\allowbreak{}3\_\allowbreak{}5\_\allowbreak{}odp\_\allowbreak{}1\\
Type: string\\
Default: nil
}

{\footnotesize визначено заходи захисту від фізичної підробки або підміни;}


{\footnotesize\bfseries
No: 2\\
Name: pe\_\allowbreak{}3\_\allowbreak{}5\_\allowbreak{}odp\_\allowbreak{}2\\
Type: string\\
Default: nil
}

{\footnotesize вибрано одне або декілька з наступних ЗНАЧЕНЬ ПАРАМЕТРІВ: \{виявлення; запобігання\};}


{\footnotesize\bfseries
No: 3\\
Name: pe\_\allowbreak{}3\_\allowbreak{}5\_\allowbreak{}odp\_\allowbreak{}3\\
Type: string\\
Default: nil
}

{\footnotesize визначено апаратні компоненти, які мають бути захищені від фізичної підробки або підміни;}


{\footnotesize\bfseries
No: 4\\
Name: pe\_\allowbreak{}3\_\allowbreak{}5\_\allowbreak{}01\\
Type: string\\
Default: nil
}

{\footnotesize <PE-03(05)\_\allowbreak{}ODP[01] заходи захисту > застосовуються для системи.}




\subsubsection{ТЕСТУВАННЯ НА МОЖЛИВІСТЬ ПРОНИКНЕННЯ (PE-3(6))}


\vspace{10pt}
\textit{Немає параметрів для цього контролю.}
\vspace{10pt}


\subsubsection{ФІЗИЧНІ ПЕРЕШКОДИ (PE-3(7))}


\vspace{10pt}
\textit{Немає параметрів для цього контролю.}
\vspace{10pt}


\subsubsection{КОНТРОЛЬ ДОСТУПУ У ВЕСТИБЮЛІ (ХОЛІ) (PE-3(8))}


\vspace{10pt}
{\footnotesize\bfseries
No: 1\\
Name: pe\_\allowbreak{}3\_\allowbreak{}8\_\allowbreak{}odp\\
Type: string\\
Default: nil
}

{\footnotesize визначено місця на об'єкті, де необхідний контроль доступу;}


{\footnotesize\bfseries
No: 2\\
Name: pe\_\allowbreak{}3\_\allowbreak{}8\_\allowbreak{}01\\
Type: string\\
Default: nil
}

{\footnotesize контроль доступу використовуються в < PE-03(08)\_\allowbreak{}ODP місцях>.}




\subsection{ЛІНІЙ ЕЛЕКТРОЖИВЛЕННЯ (PE-4)}


\vspace{10pt}
{\footnotesize\bfseries
No: 1\\
Name: pe\_\allowbreak{}4\_\allowbreak{}odp\_\allowbreak{}1\\
Type: string\\
Default: nil
}

{\footnotesize визначені системи розподілу та постачання живлення, які потребують фізичного контролю доступу;}


{\footnotesize\bfseries
No: 2\\
Name: pe\_\allowbreak{}4\_\allowbreak{}odp\_\allowbreak{}2\\
Type: string\\
Default: nil
}

{\footnotesize визначено заходи захисту, які необхідно впровадити для контролю фізичного доступу до систем розподілу та постачання живлення в межах об'єкту організації;}


{\footnotesize\bfseries
No: 3\\
Name: pe\_\allowbreak{}4\_\allowbreak{}01\\
Type: string\\
Default: nil
}

{\footnotesize фізичний доступ до  <PE-04\_\allowbreak{}ODP[01] систем розподілу та постачання живлення> в межах об'єктів організації контролюється}




\subsection{КОНТРОЛЬ ДОСТУПУ В ПРИМІЩЕННЯ ДЛЯ ВІДОБРАЖЕННЯ ІНФОРМАЦІЇ (PE-5)}


\vspace{10pt}
{\footnotesize\bfseries
No: 1\\
Name: pe\_\allowbreak{}5\_\allowbreak{}odp\\
Type: string\\
Default: nil
}

{\footnotesize визначено пристрої для виведення інформації над якими необхідний контроль над фізичним доступом до вихідних даних; 362}


{\footnotesize\bfseries
No: 2\\
Name: pe\_\allowbreak{}5\_\allowbreak{}01\\
Type: string\\
Default: nil
}

{\footnotesize керування фізичним доступом до вихідних даних здійснюється з <PE-05\_\allowbreak{}ODP пристроїв для виведення інформації >, для запобігання несанкціонованого отримання користувачами вихідних даних.}




\subsubsection{ДОСТУП ДО ВИХІДНИХ ДАНИХ УПОВНОВАЖЕНИМИ ОСОБАМИ (PE-5(1)) [Вилучено]}
[Вилучено: включено до PE-05].

\vspace{10pt}
\textit{Немає параметрів для цього контролю.}
\vspace{10pt}


\subsubsection{ДОСТУП ДО ВИХІДНИХ ДАНИХ ФІЗИЧНИМИ ОСОБАМИ (PE-5(2))}
пов'язуються дані про цифрову ідентичність з підтвердженням отримання даних від вихідних пристроїв;

\vspace{10pt}
{\footnotesize\bfseries
No: 1\\
Name: pe\_\allowbreak{}5\_\allowbreak{}2\\
Type: string\\
Default: nil
}

{\footnotesize пов'язуються дані про цифрову ідентичність з підтвердженням отримання даних від вихідних пристроїв;}




\subsubsection{МАРКУВАННЯ ПРИСТРОЇВ ВИВЕДЕННЯ ІНФОРМАЦІЇ (PE-5(3)) [Вилучено]}
[Вилучено: включено до PE-22].

\vspace{10pt}
\textit{Немає параметрів для цього контролю.}
\vspace{10pt}


\subsection{МОНІТОРИНГ ФІЗИЧНОГО ДОСТУПУ (PE-6)}


\vspace{10pt}
{\footnotesize\bfseries
No: 1\\
Name: pe\_\allowbreak{}6\_\allowbreak{}odp\_\allowbreak{}1\\
Type: string\\
Default: nil
}

{\footnotesize визначено частоту перегляду журналів фізичного доступу;}


{\footnotesize\bfseries
No: 2\\
Name: pe\_\allowbreak{}6\_\allowbreak{}odp\_\allowbreak{}2\\
Type: string\\
Default: nil
}

{\footnotesize визначено події або потенційні ознаки подій, що вимагають перегляду журналів фізичного доступу;}


{\footnotesize\bfseries
No: 3\\
Name: pe\_\allowbreak{}6\_\allowbreak{}a\\
Type: string\\
Default: nil
}

{\footnotesize фізичний доступ до об'єкту, де знаходиться система, моніториться з метою виявлення та реагування на інциденти фізичної безпеки;}


{\footnotesize\bfseries
No: 4\\
Name: pe\_\allowbreak{}6\_\allowbreak{}b\_\allowbreak{}1\\
Type: string\\
Default: nil
}

{\footnotesize переглядаються журнали фізичного доступу  <PE-06\_\allowbreak{}ODP[01] частота>;}


{\footnotesize\bfseries
No: 5\\
Name: pe\_\allowbreak{}6\_\allowbreak{}b\_\allowbreak{}2\\
Type: string\\
Default: nil
}

{\footnotesize журнали фізичного доступу переглядаються при виникненні  <PE06\_\allowbreak{}ODP[02] подій>;}


{\footnotesize\bfseries
No: 6\\
Name: pe\_\allowbreak{}6\_\allowbreak{}c\_\allowbreak{}1\\
Type: string\\
Default: nil
}

{\footnotesize результати переглядів узгоджуються з можливостями організації щодо реагування на інциденти;}


{\footnotesize\bfseries
No: 7\\
Name: pe\_\allowbreak{}6\_\allowbreak{}c\_\allowbreak{}2\\
Type: string\\
Default: nil
}

{\footnotesize результати розслідувань узгоджуються з можливостями організації щодо реагування на інциденти;}




\subsubsection{ОХОРОННА СИГНАЛІЗАЦІЯ ТА ОБЛАДНАННЯ ДЛЯ СПОСТЕРЕЖЕННЯ (PE-6(1))}


\vspace{10pt}
\textit{Немає параметрів для цього контролю.}
\vspace{10pt}


\subsubsection{МОНІТОРИНГ ФІЗИЧНОГО ДОСТУПУ --- АВТОМАТИЧНІ РОЗПІЗНАВАННЯ ВТОРГНЕНЬ І ВІДПОВІДНА РЕАКЦІЯ (PE-6(2))}


\vspace{10pt}
{\footnotesize\bfseries
No: 1\\
Name: pe\_\allowbreak{}6\_\allowbreak{}2\_\allowbreak{}odp\_\allowbreak{}1\\
Type: string\\
Default: nil
}

{\footnotesize визначено класи або типи вторгнень, які мають розпізнаватися автоматизованими механізмами;}


{\footnotesize\bfseries
No: 2\\
Name: pe\_\allowbreak{}6\_\allowbreak{}2\_\allowbreak{}odp\_\allowbreak{}2\\
Type: string\\
Default: nil
}

{\footnotesize визначено реакції, які мають ініціюватися автоматизованими механізмами при розпізнаванні визначених організацією класів або типів вторгнень;}


{\footnotesize\bfseries
No: 3\\
Name: pe\_\allowbreak{}6\_\allowbreak{}2\_\allowbreak{}odp\_\allowbreak{}3\\
Type: string\\
Default: nil
}

{\footnotesize визначено автоматизовані механізми, що використовуються для розпізнавання класів або типів вторгнень та ініціювання дій реагування (визначені в PE06(02)\_\allowbreak{}ODP);}


{\footnotesize\bfseries
No: 4\\
Name: pe\_\allowbreak{}6\_\allowbreak{}2\_\allowbreak{}1\\
Type: string\\
Default: nil
}

{\footnotesize розпізнаються <PE-06(02)\_\allowbreak{}ODP[01] класи або типи вторгнень>;}


{\footnotesize\bfseries
No: 5\\
Name: pe\_\allowbreak{}6\_\allowbreak{}2\_\allowbreak{}2\\
Type: string\\
Default: nil
}

{\footnotesize <PE-06(02)\_\allowbreak{}ODP[02] реакції> ініціюються за допомогою}




\subsubsection{ВІДЕОСПОСТЕРЕЖЕННЯ (PE-6(3))}
ведеться відеоспостереження за <PE-06(03)\_\allowbreak{}ODP[01] зонами>;\\ 
відеозаписи переглядаються <PE-06(03)\_\allowbreak{}ODP[02] частота>;\\ 
відеозаписи зберігаються протягом <PE-06(03)\_\allowbreak{}ODP[03] періоду часу>.

\vspace{10pt}
{\footnotesize\bfseries
No: 1\\
Name: pe\_\allowbreak{}6\_\allowbreak{}3\_\allowbreak{}odp\_\allowbreak{}1\\
Type: string\\
Default: nil
}

{\footnotesize визначено зони, де буде застосовуватися відеоспостереження;}


{\footnotesize\bfseries
No: 2\\
Name: pe\_\allowbreak{}6\_\allowbreak{}3\_\allowbreak{}odp\_\allowbreak{}2\\
Type: string\\
Default: nil
}

{\footnotesize визначено частоту перегляду відеозаписів;}


{\footnotesize\bfseries
No: 3\\
Name: pe\_\allowbreak{}6\_\allowbreak{}3\_\allowbreak{}odp\_\allowbreak{}3\\
Type: string\\
Default: nil
}

{\footnotesize визначено період часу, протягом якого необхідно зберігати відеозаписи;}


{\footnotesize\bfseries
No: 4\\
Name: pe\_\allowbreak{}6\_\allowbreak{}3\_\allowbreak{}a\\
Type: string\\
Default: nil
}

{\footnotesize ведеться відеоспостереження за <PE-06(03)\_\allowbreak{}ODP[01] зонами>;}


{\footnotesize\bfseries
No: 5\\
Name: pe\_\allowbreak{}6\_\allowbreak{}3\_\allowbreak{}b\\
Type: string\\
Default: nil
}

{\footnotesize відеозаписи переглядаються <PE-06(03)\_\allowbreak{}ODP[02] частота>;}


{\footnotesize\bfseries
No: 6\\
Name: pe\_\allowbreak{}6\_\allowbreak{}3\_\allowbreak{}c\\
Type: string\\
Default: nil
}

{\footnotesize відеозаписи зберігаються протягом <PE-06(03)\_\allowbreak{}ODP[03] періоду часу>.}




\subsubsection{МОНІТОРИНГ ФІЗИЧНОГО ДОСТУПУ ДО СИСТЕМИ (PE-6(4))}
моніторинг фізичного доступу до системи здійснюється на додаток до моніторингу фізичного доступу до об'єкта в <PE-06(04)\_\allowbreak{}ODP фізичні приміщення>.

\vspace{10pt}
{\footnotesize\bfseries
No: 1\\
Name: pe\_\allowbreak{}6\_\allowbreak{}4\_\allowbreak{}odp\\
Type: string\\
Default: nil
}

{\footnotesize визначено фізичні приміщення, що містять один або більше компонентів системи;}


{\footnotesize\bfseries
No: 2\\
Name: pe\_\allowbreak{}6\_\allowbreak{}4\_\allowbreak{}01\\
Type: string\\
Default: nil
}

{\footnotesize моніторинг фізичного доступу до системи здійснюється на додаток до моніторингу фізичного доступу до об'єкта в <PE-06(04)\_\allowbreak{}ODP фізичні приміщення>.}




\subsection{КОНТРОЛЬ ВІДВІДУВАЧІВ (PE-7)}


\vspace{10pt}
\textit{Немає параметрів для цього контролю.}
\vspace{10pt}


\subsection{РЕЄСТР ДОСТУПУ ВІДВІДУВАЧІВ (PE-8)}


\vspace{10pt}
{\footnotesize\bfseries
No: 1\\
Name: pe\_\allowbreak{}8\_\allowbreak{}odp\_\allowbreak{}1\\
Type: string\\
Default: nil
}

{\footnotesize визначено період часу, протягом якого зберігатимуться записи про доступ відвідувачів до об'єкта, на якому перебуває система;}


{\footnotesize\bfseries
No: 2\\
Name: pe\_\allowbreak{}8\_\allowbreak{}odp\_\allowbreak{}2\\
Type: string\\
Default: nil
}

{\footnotesize визначено частоту перегляду записів про доступ відвідувачів;}


{\footnotesize\bfseries
No: 3\\
Name: pe\_\allowbreak{}8\_\allowbreak{}odp\_\allowbreak{}3\\
Type: string\\
Default: nil
}

{\footnotesize визначено персонал, якому повідомляється про порушення записів про доступ відвідувачів;}


{\footnotesize\bfseries
No: 4\\
Name: pe\_\allowbreak{}8\_\allowbreak{}a\\
Type: string\\
Default: nil
}

{\footnotesize записи про доступ відвідувачів для об'єкта, на якому знаходиться}


{\footnotesize\bfseries
No: 5\\
Name: pe\_\allowbreak{}8\_\allowbreak{}b\\
Type: string\\
Default: nil
}

{\footnotesize переглядаються записи про доступ відвідувачів  <PE-08\_\allowbreak{}ODP[02] частота>;}


{\footnotesize\bfseries
No: 6\\
Name: pe\_\allowbreak{}8\_\allowbreak{}c\\
Type: string\\
Default: nil
}

{\footnotesize про порушення записів про доступ відвідувачів повідомляється}




\subsubsection{РЕЄСТРИ ДОСТУПУ ВІДВІДУВАЧІВ --- ОБМЕЖЕННЯ ІНФОРМАЦІЇ, ЩО ІДЕНТИФІКУЄ ОСОБУ (PE-8(3))}


\vspace{10pt}
{\footnotesize\bfseries
No: 1\\
Name: pe\_\allowbreak{}8\_\allowbreak{}3\_\allowbreak{}odp\\
Type: string\\
Default: nil
}

{\footnotesize в оцінці ризиків для конфіденційності визначено елементи, що обмежуються в реєстрі відвідувачів}


{\footnotesize\bfseries
No: 2\\
Name: pe\_\allowbreak{}8\_\allowbreak{}3\_\allowbreak{}01\\
Type: string\\
Default: nil
}

{\footnotesize інформація, що ідентифікує особу, яка міститься в реєстрах доступу відвідувачів, обмежується < PE-08(03)\_\allowbreak{}ODP елементи>, визначеними в оцінці ризиків для конфіденційності.}




\subsection{ЕНЕРГЕТИЧНЕ ОБЛАДНАННЯ ТА КАБЕЛІ (PE-9)}


\vspace{10pt}
\textit{Немає параметрів для цього контролю.}
\vspace{10pt}


\subsubsection{РЕЗЕРВНІ КАБЕЛІ (PE-9(1))}
використовувати резервні силові кабельні системи, які фізично відокремлені на <PE-09(01)\_\allowbreak{}ODP відстань>.

\vspace{10pt}
{\footnotesize\bfseries
No: 1\\
Name: pe\_\allowbreak{}9\_\allowbreak{}1\_\allowbreak{}odp\\
Type: string\\
Default: nil
}

{\footnotesize визначено відстань, на яку повинні бути відокремлені резервні силові кабельні системи;}


{\footnotesize\bfseries
No: 2\\
Name: pe\_\allowbreak{}9\_\allowbreak{}1\_\allowbreak{}01\\
Type: string\\
Default: nil
}

{\footnotesize використовувати резервні силові кабельні системи, які фізично відокремлені на <PE-09(01)\_\allowbreak{}ODP відстань>.}




\subsubsection{АВТОМАТИЧНЕ КЕРУВАННЯ НАПРУГОЮ (PE-9(2))}
впроваджено механізми автоматичного керування напругою для <PE-09(02)\_\allowbreak{}ODP критичних компонентів системи>.

\vspace{10pt}
{\footnotesize\bfseries
No: 1\\
Name: pe\_\allowbreak{}9\_\allowbreak{}2\_\allowbreak{}odp\\
Type: string\\
Default: nil
}

{\footnotesize визначено критичні компоненти системи для яких необхідно впровадити механізми автоматичного керування напругою;}


{\footnotesize\bfseries
No: 2\\
Name: pe\_\allowbreak{}9\_\allowbreak{}2\_\allowbreak{}01\\
Type: string\\
Default: nil
}

{\footnotesize впроваджено механізми автоматичного керування напругою для <PE-09(02)\_\allowbreak{}ODP критичних компонентів системи>.}




\subsection{АВАРІЙНЕ ВІДКЛЮЧЕННЯ (PE-10)}


\vspace{10pt}
{\footnotesize\bfseries
No: 1\\
Name: pe\_\allowbreak{}10\_\allowbreak{}odp\_\allowbreak{}1\\
Type: string\\
Default: nil
}

{\footnotesize визначено систему або окремі компоненти системи, які потребують можливості вимкнення живлення в надзвичайних ситуаціях;}


{\footnotesize\bfseries
No: 2\\
Name: pe\_\allowbreak{}10\_\allowbreak{}odp\_\allowbreak{}2\\
Type: string\\
Default: nil
}

{\footnotesize визначено розташування перемикачів або пристроїв аварійного вимкнення в системі або компоненті системи;}


{\footnotesize\bfseries
No: 3\\
Name: pe\_\allowbreak{}10\_\allowbreak{}a\\
Type: string\\
Default: nil
}

{\footnotesize передбачена можливість відключення живлення <PE-10\_\allowbreak{}ODP[01] системи або окремих компонентів системи>  в надзвичайних ситуаціях;}


{\footnotesize\bfseries
No: 4\\
Name: pe\_\allowbreak{}10\_\allowbreak{}b\\
Type: string\\
Default: nil
}

{\footnotesize аварійні перемикачі або пристрої вимкнення розміщені в < PE10\_\allowbreak{}ODP[02] розташування>, щоб забезпечити доступ для персоналу;}


{\footnotesize\bfseries
No: 5\\
Name: pe\_\allowbreak{}10\_\allowbreak{}c\\
Type: string\\
Default: nil
}

{\footnotesize можливість аварійного вимкнення живлення захищена від несанкціонованої активації.}




\subsubsection{ВИПАДКОВА І НЕСАНКЦІОНОВАНА АКТИВАЦІЯ (PE-10(1)) [Вилучено]}
[Виключено: включено до PE-10].

\vspace{10pt}
\textit{Немає параметрів для цього контролю.}
\vspace{10pt}


\subsection{АВАРІЙНЕ ЕНЕРГОЗАБЕЗПЕЧЕННЯ (PE-11)}


\vspace{10pt}
{\footnotesize\bfseries
No: 1\\
Name: pe\_\allowbreak{}11\_\allowbreak{}odp\\
Type: string\\
Default: nil
}

{\footnotesize вибрано одне з наступних ЗНАЧЕНЬ ПАРАМЕТРІВ: \{впорядковане виключення системи; перехід системи на довгострокову альтернативну систему живлення\};): 371}


{\footnotesize\bfseries
No: 2\\
Name: pe\_\allowbreak{}11\_\allowbreak{}01\\
Type: string\\
Default: nil
}

{\footnotesize передбачено джерело безперебійного живлення для полегшення <PE-11\_\allowbreak{}ODP ВИБРАНЕ ЗНАЧЕННЯ ПАРАМЕТРА> у випадку втрати основного джерела живлення.}


{\footnotesize\bfseries
No: 3\\
Name: pe\_\allowbreak{}11\_\allowbreak{}a\\
Type: string\\
Default: nil
}

{\footnotesize альтернативне джерело живлення, передбачене для системи, є автономним;}


{\footnotesize\bfseries
No: 4\\
Name: pe\_\allowbreak{}11\_\allowbreak{}b\\
Type: string\\
Default: nil
}

{\footnotesize альтернативне джерело живлення, передбачене для системи, не залежить від зовнішнього постачання енергії;}


{\footnotesize\bfseries
No: 5\\
Name: pe\_\allowbreak{}11\_\allowbreak{}c\\
Type: string\\
Default: nil
}

{\footnotesize альтернативне джерело живлення, передбачене для системи, здатне підтримувати <PE-11(02)\_\allowbreak{}ODP[02] ВИБРАНЕ ЗНАЧЕННЯ ПАРАМЕТРА> у разі тривалої втрати основного джерела живлення.}




\subsubsection{ДОВГОСТРОКОВЕ АЛЬТЕРНАТИВНЕ ДЖЕРЕЛО ЖИВЛЕННЯ - МІНІМАЛЬНІ ЕКСПЛУАТАЦІЙНІ МОЖЛИВОСТІ (PE-11(1))}
активується альтернативне джерело живлення, передбачене для системи <PE-11(01)\_\allowbreak{}ODP ВИБРАНЕ ЗНАЧЕННЯ ПАРАМЕТРА>;\\ 
альтернативне джерело живлення, передбачене для системи, може підтримувати мінімально необхідну працездатність у разі тривалої втрати основного джерела живлення.

\vspace{10pt}
{\footnotesize\bfseries
No: 1\\
Name: pe\_\allowbreak{}11\_\allowbreak{}1\_\allowbreak{}odp\\
Type: string\\
Default: nil
}

{\footnotesize вибрано одне з наступних ЗНАЧЕНЬ ПАРАМЕТРА: \{вручну; автоматично\};}


{\footnotesize\bfseries
No: 2\\
Name: pe\_\allowbreak{}11\_\allowbreak{}1\_\allowbreak{}01\\
Type: string\\
Default: nil
}

{\footnotesize активується альтернативне джерело живлення, передбачене для системи <PE-11(01)\_\allowbreak{}ODP ВИБРАНЕ ЗНАЧЕННЯ ПАРАМЕТРА>;}


{\footnotesize\bfseries
No: 3\\
Name: pe\_\allowbreak{}11\_\allowbreak{}1\_\allowbreak{}02\\
Type: string\\
Default: nil
}

{\footnotesize альтернативне джерело живлення, передбачене для системи, може підтримувати мінімально необхідну працездатність у разі тривалої втрати основного джерела живлення.}




\subsubsection{ДОВГОСТРОКОВЕ АЛЬТЕРНАТИВНЕ ДЖЕРЕЛО ЖИВЛЕННЯ -- АВТОНОМНЕ ЖИВЛЕННЯ (PE-11(2))}
активується альтернативне джерело живлення, передбачене для системи <PE-11(02)\_\allowbreak{}ODP[01] ВИБРАНЕ ЗНАЧЕННЯ ПАРАМЕТРА>;\\ 
альтернативне джерело живлення, передбачене для системи, є автономним;\\ 
альтернативне джерело живлення, передбачене для системи, не залежить від зовнішнього постачання енергії;\\ 
альтернативне джерело живлення, передбачене для системи, здатне підтримувати <PE-11(02)\_\allowbreak{}ODP[02] ВИБРАНЕ ЗНАЧЕННЯ ПАРАМЕТРА> у разі тривалої втрати основного джерела живлення.

\vspace{10pt}
{\footnotesize\bfseries
No: 1\\
Name: pe\_\allowbreak{}11\_\allowbreak{}2\_\allowbreak{}odp\_\allowbreak{}1\\
Type: string\\
Default: nil
}

{\footnotesize вибрано одне з наступних ЗНАЧЕНЬ ПАРАМЕТРІВ: \{вручну; автоматично\};}


{\footnotesize\bfseries
No: 2\\
Name: pe\_\allowbreak{}11\_\allowbreak{}2\_\allowbreak{}odp\_\allowbreak{}2\\
Type: string\\
Default: nil
}

{\footnotesize вибрано одне з наступних ЗНАЧЕНЬ ПАРАМЕТРІВ: \{мінімально необхідні операційні можливості; повна експлуатаційна здатність\};}


{\footnotesize\bfseries
No: 3\\
Name: pe\_\allowbreak{}11\_\allowbreak{}2\_\allowbreak{}01\\
Type: string\\
Default: nil
}

{\footnotesize активується альтернативне джерело живлення, передбачене для системи <PE-11(02)\_\allowbreak{}ODP[01] ВИБРАНЕ ЗНАЧЕННЯ ПАРАМЕТРА>;}


{\footnotesize\bfseries
No: 4\\
Name: pe\_\allowbreak{}11\_\allowbreak{}2\_\allowbreak{}a\\
Type: string\\
Default: nil
}

{\footnotesize альтернативне джерело живлення, передбачене для системи, є автономним;}


{\footnotesize\bfseries
No: 5\\
Name: pe\_\allowbreak{}11\_\allowbreak{}2\_\allowbreak{}b\\
Type: string\\
Default: nil
}

{\footnotesize альтернативне джерело живлення, передбачене для системи, не залежить від зовнішнього постачання енергії;}


{\footnotesize\bfseries
No: 6\\
Name: pe\_\allowbreak{}11\_\allowbreak{}2\_\allowbreak{}c\\
Type: string\\
Default: nil
}

{\footnotesize альтернативне джерело живлення, передбачене для системи, здатне підтримувати <PE-11(02)\_\allowbreak{}ODP[02] ВИБРАНЕ ЗНАЧЕННЯ ПАРАМЕТРА> у разі тривалої втрати основного джерела живлення.}




\subsection{АВАРІЙНЕ ОСВІТЛЕННЯ (PE-12)}


\vspace{10pt}
{\footnotesize\bfseries
No: 1\\
Name: pe\_\allowbreak{}12\_\allowbreak{}1\\
Type: string\\
Default: nil
}

{\footnotesize автоматичне аварійне освітлення, яке вмикається в разі відключення або збою в електропостачанні;}


{\footnotesize\bfseries
No: 2\\
Name: pe\_\allowbreak{}12\_\allowbreak{}2\\
Type: string\\
Default: nil
}

{\footnotesize підтримується автоматичне аварійне освітлення, яке вмикається в разі відключення або збою в електропостачанні;}


{\footnotesize\bfseries
No: 3\\
Name: pe\_\allowbreak{}12\_\allowbreak{}3\\
Type: string\\
Default: nil
}

{\footnotesize автоматичне аварійне освітлення системи освітлює евакуаційні виходи в межах об'єкта;}


{\footnotesize\bfseries
No: 4\\
Name: pe\_\allowbreak{}12\_\allowbreak{}4\\
Type: string\\
Default: nil
}

{\footnotesize автоматичне аварійне освітлення системи освітлює шляхи евакуації в 373 межах об'єкта;}




\subsubsection{ОСНОВНІ ЗАВДАННЯ ТА ФУНКЦІЇ (PE-12(1))}
аварійне освітлення передбачено для всіх зон, що підтримують виконання основних завдань і функцій.

\vspace{10pt}
{\footnotesize\bfseries
No: 1\\
Name: pe\_\allowbreak{}12\_\allowbreak{}1\\
Type: string\\
Default: nil
}

{\footnotesize аварійне освітлення передбачено для всіх зон, що підтримують виконання основних завдань і функцій.}




\subsection{ПРОТИПОЖЕЖНИЙ ЗАХИСТ (PE-13)}
застосовуються системи пожежної сигналізації;\\ 
системи пожежної сигналізації підтримуються незалежним джерелом енергії;\\ 
підтримуються в робочому стані системи пожежної сигналізації;\\ 
застосовуються системи пожежогасіння;\\ 
системи пожежогасіння підтримуються незалежним джерелом енергії;\\ 
підтримуються в робочому стані системи пожежогасіння.

\vspace{10pt}
{\footnotesize\bfseries
No: 1\\
Name: pe\_\allowbreak{}13\_\allowbreak{}01\\
Type: string\\
Default: nil
}

{\footnotesize застосовуються системи пожежної сигналізації;}


{\footnotesize\bfseries
No: 2\\
Name: pe\_\allowbreak{}13\_\allowbreak{}02\\
Type: string\\
Default: nil
}

{\footnotesize системи пожежної сигналізації підтримуються незалежним джерелом енергії;}


{\footnotesize\bfseries
No: 3\\
Name: pe\_\allowbreak{}13\_\allowbreak{}03\\
Type: string\\
Default: nil
}

{\footnotesize підтримуються в робочому стані системи пожежної сигналізації;}


{\footnotesize\bfseries
No: 4\\
Name: pe\_\allowbreak{}13\_\allowbreak{}04\\
Type: string\\
Default: nil
}

{\footnotesize застосовуються системи пожежогасіння;}


{\footnotesize\bfseries
No: 5\\
Name: pe\_\allowbreak{}13\_\allowbreak{}05\\
Type: string\\
Default: nil
}

{\footnotesize системи пожежогасіння підтримуються незалежним джерелом енергії;}


{\footnotesize\bfseries
No: 6\\
Name: pe\_\allowbreak{}13\_\allowbreak{}06\\
Type: string\\
Default: nil
}

{\footnotesize підтримуються в робочому стані системи пожежогасіння.}




\subsubsection{ПРИСТРОЇ ТА СИСТЕМИ ВИЯВЛЕННЯ (PE-13(1))}
використовуються системи пожежної сигналізації, які автоматично спрацьовують у разі пожежі;\\ 
використовуються системи виявлення пожежі, які автоматично сповіщають персонал або ролі у разі виникнення пожежі;\\ 
використовуються системи виявлення пожежі, які автоматично сповіщають аварійні команди у разі виникнення пожежі.

\vspace{10pt}
{\footnotesize\bfseries
No: 1\\
Name: pe\_\allowbreak{}13\_\allowbreak{}1\_\allowbreak{}odp\_\allowbreak{}1\\
Type: string\\
Default: nil
}

{\footnotesize визначено персонал або ролі, які мають бути повідомлені у випадку пожежі;}


{\footnotesize\bfseries
No: 2\\
Name: pe\_\allowbreak{}13\_\allowbreak{}1\_\allowbreak{}odp\_\allowbreak{}2\\
Type: string\\
Default: nil
}

{\footnotesize визначено аварійні команди, які мають бути сповіщені у випадку пожежі;}


{\footnotesize\bfseries
No: 3\\
Name: pe\_\allowbreak{}13\_\allowbreak{}1\_\allowbreak{}01\\
Type: string\\
Default: nil
}

{\footnotesize використовуються системи пожежної сигналізації, які автоматично спрацьовують у разі пожежі;}


{\footnotesize\bfseries
No: 4\\
Name: pe\_\allowbreak{}13\_\allowbreak{}1\_\allowbreak{}02\\
Type: string\\
Default: nil
}

{\footnotesize використовуються системи виявлення пожежі, які автоматично сповіщають <PE-13(01)\_\allowbreak{}ODP[01] персонал або ролі> у разі виникнення пожежі;}


{\footnotesize\bfseries
No: 5\\
Name: pe\_\allowbreak{}13\_\allowbreak{}1\_\allowbreak{}03\\
Type: string\\
Default: nil
}

{\footnotesize використовуються системи виявлення пожежі, які автоматично сповіщають <PE-13(01)\_\allowbreak{}ODP[02] аварійні команди> у разі виникнення пожежі.}




\subsubsection{ПРИСТРОЇ ТА СИСТЕМИ АВТОМАТИЧНОГО ПОЖЕЖОГАСІННЯ (PE-13(2))}
застосовуються системи пожежогасіння, які активуються автоматично;\\ 
використовуються системи пожежогасіння, які автоматично сповіщають персонал або ролі;\\ 
використовуються системи пожежогасіння, які автоматично сповіщають аварійні команди;\\ 
використовується автоматична система пожежогасіння, коли об'єкт не укомплектований відповідним персоналом на постійній основі.

\vspace{10pt}
{\footnotesize\bfseries
No: 1\\
Name: pe\_\allowbreak{}13\_\allowbreak{}2\_\allowbreak{}odp\_\allowbreak{}1\\
Type: string\\
Default: nil
}

{\footnotesize визначено персонал або ролі, які мають бути сповіщені у випадку пожежі;}


{\footnotesize\bfseries
No: 2\\
Name: pe\_\allowbreak{}13\_\allowbreak{}2\_\allowbreak{}odp\_\allowbreak{}2\\
Type: string\\
Default: nil
}

{\footnotesize визначені аварійні команди, які повинні бути сповіщені в разі пожежі;}


{\footnotesize\bfseries
No: 3\\
Name: pe\_\allowbreak{}13\_\allowbreak{}2\_\allowbreak{}a\_\allowbreak{}01\\
Type: string\\
Default: nil
}

{\footnotesize застосовуються системи пожежогасіння, які активуються автоматично;}


{\footnotesize\bfseries
No: 4\\
Name: pe\_\allowbreak{}13\_\allowbreak{}2\_\allowbreak{}a\_\allowbreak{}02\\
Type: string\\
Default: nil
}

{\footnotesize використовуються системи пожежогасіння, які автоматично сповіщають <PE-13(02)\_\allowbreak{}ODP[01] персонал або ролі>;}


{\footnotesize\bfseries
No: 5\\
Name: pe\_\allowbreak{}13\_\allowbreak{}2\_\allowbreak{}a\_\allowbreak{}03\\
Type: string\\
Default: nil
}

{\footnotesize використовуються системи пожежогасіння, які автоматично сповіщають <PE-13(02)\_\allowbreak{}ODP[02] аварійні команди>;}


{\footnotesize\bfseries
No: 6\\
Name: pe\_\allowbreak{}13\_\allowbreak{}2\_\allowbreak{}b\\
Type: string\\
Default: nil
}

{\footnotesize використовується автоматична система пожежогасіння, коли об'єкт не укомплектований відповідним персоналом на постійній основі.}




\subsubsection{АВТОМАТИЧНЕ ПОЖЕЖОГАСІННЯ (PE-13(3)) [Вилучено]}
[Виключено: включено до РЕ-13(02)].

\vspace{10pt}
\textit{Немає параметрів для цього контролю.}
\vspace{10pt}


\subsubsection{ПЕРЕВІРКИ (PE-13(4))}
об'єкт проходить перевірки пожежної безпеки уповноваженими та кваліфікованими інспекторами;\\ 
виявлені недоліки за результатами перевірок пожежної безпеки усуваються у визначений термін.

\vspace{10pt}
{\footnotesize\bfseries
No: 1\\
Name: pe\_\allowbreak{}13\_\allowbreak{}4\_\allowbreak{}odp\_\allowbreak{}1\\
Type: string\\
Default: nil
}

{\footnotesize визначено частоту проведення перевірок пожежної безпеки на об'єкті;}


{\footnotesize\bfseries
No: 2\\
Name: pe\_\allowbreak{}13\_\allowbreak{}4\_\allowbreak{}odp\_\allowbreak{}2\\
Type: string\\
Default: nil
}

{\footnotesize визначено термін для усунення недоліків, виявлених перевірками пожежного нагляду;}


{\footnotesize\bfseries
No: 3\\
Name: pe\_\allowbreak{}13\_\allowbreak{}4\_\allowbreak{}01\\
Type: string\\
Default: nil
}

{\footnotesize об'єкт проходить перевірки пожежної безпеки <PE-13(04)\_\allowbreak{}ODP[01] частота> уповноваженими та кваліфікованими інспекторами;}


{\footnotesize\bfseries
No: 4\\
Name: pe\_\allowbreak{}13\_\allowbreak{}4\_\allowbreak{}02\\
Type: string\\
Default: nil
}

{\footnotesize виявлені недоліки за результатами перевірок пожежної безпеки усуваються у <PE-13(04)\_\allowbreak{}ODP[02] термін>.}




\subsection{КОНТРОЛЬ ТЕМПЕРАТУРИ ТА ВОЛОГОСТІ (PE-14)}


\vspace{10pt}
{\footnotesize\bfseries
No: 1\\
Name: pe\_\allowbreak{}14\_\allowbreak{}odp\_\allowbreak{}1\\
Type: string\\
Default: nil
}

{\footnotesize визначено прийнятні рівні для температури та вологості;}


{\footnotesize\bfseries
No: 2\\
Name: pe\_\allowbreak{}14\_\allowbreak{}odp\_\allowbreak{}2\\
Type: string\\
Default: nil
}

{\footnotesize визначено частоту моніторингу рівнів температури та вологості;}


{\footnotesize\bfseries
No: 3\\
Name: pe\_\allowbreak{}14\_\allowbreak{}a\\
Type: string\\
Default: nil
}

{\footnotesize температура та вологость підтримуються на <PE-14\_\allowbreak{}ODP[01] рівні> у приміщенні, де знаходиться система;}


{\footnotesize\bfseries
No: 4\\
Name: pe\_\allowbreak{}14\_\allowbreak{}b\\
Type: string\\
Default: nil
}

{\footnotesize контролюються рівні температури та вологості  <PE-14\_\allowbreak{}ODP[02] частота>.}




\subsubsection{АВТОМАТИЧНИЙ КОНТРОЛЬ (PE-14(1))}
впроваджено <PE-14(01)\_\allowbreak{}ODP механізми> на об'єкті для запобігання потенційно шкідливим для системи коливанням.

\vspace{10pt}
{\footnotesize\bfseries
No: 1\\
Name: pe\_\allowbreak{}14\_\allowbreak{}1\_\allowbreak{}odp\\
Type: string\\
Default: nil
}

{\footnotesize визначено механізми автоматичного регулювання температури та вологості;}


{\footnotesize\bfseries
No: 2\\
Name: pe\_\allowbreak{}14\_\allowbreak{}1\_\allowbreak{}01\\
Type: string\\
Default: nil
}

{\footnotesize впроваджено <PE-14(01)\_\allowbreak{}ODP механізми> на об'єкті для запобігання потенційно шкідливим для системи коливанням.}




\subsubsection{МОНІТОРИНГ ЗА ДОПОМОГОЮ СИГНАЛІЗАЦІЙ ТА СПОВІЩЕНЬ (PE-14(2))}
застосовується моніторинг температури та вологості;\\ 
функція моніторингу температури та вологості надає сигнал тривоги або повідомлення <PE-14(02)\_\allowbreak{}ODP персоналу або ролям>, коли зміни є потенційно шкідливими для персоналу або обладнання.

\vspace{10pt}
{\footnotesize\bfseries
No: 1\\
Name: pe\_\allowbreak{}14\_\allowbreak{}2\_\allowbreak{}odp\\
Type: string\\
Default: nil
}

{\footnotesize визначено персонал або ролі, які необхідно повідомляти в рамках моніторингу температури та вологості, коли зміни температури та вологості є потенційно шкідливими для персоналу або обладнання;}


{\footnotesize\bfseries
No: 2\\
Name: pe\_\allowbreak{}14\_\allowbreak{}2\_\allowbreak{}01\\
Type: string\\
Default: nil
}

{\footnotesize застосовується моніторинг температури та вологості;}


{\footnotesize\bfseries
No: 3\\
Name: pe\_\allowbreak{}14\_\allowbreak{}2\_\allowbreak{}02\\
Type: string\\
Default: nil
}

{\footnotesize функція моніторингу температури та вологості надає сигнал тривоги або повідомлення <PE-14(02)\_\allowbreak{}ODP персоналу або ролям>, коли зміни є потенційно шкідливими для персоналу або обладнання.}




\subsection{ЗАХИСТ ВІД ПОШКОДЖЕННЯ ВОДОЮ (PE-15)}


\vspace{10pt}
\textit{Немає параметрів для цього контролю.}
\vspace{10pt}


\subsubsection{АВТОМАТИЧНА ПІДТРИМКА (PE-15(1))}
наявність води поблизу системи можна виявити автоматично;\\ 
<PE-15(01)\_\allowbreak{}ODP[01] персонал або ролі> оповіщаються за допомогою <PE-15(01)\_\allowbreak{}ODP[02] автоматизованих механізмів>.

\vspace{10pt}
{\footnotesize\bfseries
No: 1\\
Name: pe\_\allowbreak{}15\_\allowbreak{}1\_\allowbreak{}odp\_\allowbreak{}1\\
Type: string\\
Default: nil
}

{\footnotesize визначено персонал або ролі, які слід сповіщати, коли біля системи виявлено присутність води;}


{\footnotesize\bfseries
No: 2\\
Name: pe\_\allowbreak{}15\_\allowbreak{}1\_\allowbreak{}odp\_\allowbreak{}2\\
Type: string\\
Default: nil
}

{\footnotesize визначено автоматизовані механізми, що використовуються для виявлення присутності води поблизу системи;}


{\footnotesize\bfseries
No: 3\\
Name: pe\_\allowbreak{}15\_\allowbreak{}1\_\allowbreak{}01\\
Type: string\\
Default: nil
}

{\footnotesize наявність води поблизу системи можна виявити автоматично;}


{\footnotesize\bfseries
No: 4\\
Name: pe\_\allowbreak{}15\_\allowbreak{}1\_\allowbreak{}02\\
Type: string\\
Default: nil
}

{\footnotesize <PE-15(01)\_\allowbreak{}ODP[01] персонал або ролі> оповіщаються за допомогою <PE-15(01)\_\allowbreak{}ODP[02] автоматизованих механізмів>.}




\subsection{ДОСТАВКА ТА ВИДАЛЕННЯ (PE-16)}


\vspace{10pt}
{\footnotesize\bfseries
No: 1\\
Name: pe\_\allowbreak{}16\_\allowbreak{}odp\_\allowbreak{}1\\
Type: string\\
Default: nil
}

{\footnotesize визначено типи компонентів системи, які підлягають авторизації та контролю при вході на об'єкт;}


{\footnotesize\bfseries
No: 2\\
Name: pe\_\allowbreak{}16\_\allowbreak{}odp\_\allowbreak{}2\\
Type: string\\
Default: nil
}

{\footnotesize визначено типи компонентів системи, які підлягають авторизації та контролю при виході з об'єкта;}


{\footnotesize\bfseries
No: 3\\
Name: pe\_\allowbreak{}16\_\allowbreak{}a\_\allowbreak{}1\\
Type: string\\
Default: nil
}

{\footnotesize <PE-16\_\allowbreak{}ODP[01] типи компонентів системи>  авторизуються при вході на об'єкт;}


{\footnotesize\bfseries
No: 4\\
Name: pe\_\allowbreak{}16\_\allowbreak{}a\_\allowbreak{}2\\
Type: string\\
Default: nil
}

{\footnotesize <PE-16\_\allowbreak{}ODP[01] типи компонентів системи>  контролюються при вході на об'єкт;}


{\footnotesize\bfseries
No: 5\\
Name: pe\_\allowbreak{}16\_\allowbreak{}a\_\allowbreak{}3\\
Type: string\\
Default: nil
}

{\footnotesize <PE-16\_\allowbreak{}ODP[02] типи компонентів системи>  авторизуються при виході з об'єкта;}


{\footnotesize\bfseries
No: 6\\
Name: pe\_\allowbreak{}16\_\allowbreak{}a\_\allowbreak{}4\\
Type: string\\
Default: nil
}

{\footnotesize <PE-16\_\allowbreak{}ODP[02] типи компонентів системи>  контролюються при виході з об'єкта;}


{\footnotesize\bfseries
No: 7\\
Name: pe\_\allowbreak{}16\_\allowbreak{}b\\
Type: string\\
Default: nil
}

{\footnotesize ведеться облік компонентів системи, зазначених вище}




\subsection{АЛЬТЕРНАТИВНЕ РОБОЧЕ МІСЦЕ (PE-17)}


\vspace{10pt}
{\footnotesize\bfseries
No: 1\\
Name: pe\_\allowbreak{}17\_\allowbreak{}odp\_\allowbreak{}1\\
Type: string\\
Default: nil
}

{\footnotesize визначені альтернативні робочі місця, дозволені для використання працівниками;}


{\footnotesize\bfseries
No: 2\\
Name: pe\_\allowbreak{}17\_\allowbreak{}odp\_\allowbreak{}2\\
Type: string\\
Default: nil
}

{\footnotesize визначаються заходи захисту, які будуть застосовуватися на альтернативних робочих місцях; РЕ-17(c) оцінюється ефективність заходів захисту на альтернативних робочих місцях; РЕ-17(d) працівникам надаються засоби комунікації з персоналом служби інформаційної безпеки на випадок інцидентів.}




\subsection{РОЗТАШУВАННЯ КОМПОНЕНТІВ СИСТЕМИ (PE-18)}


\vspace{10pt}
{\footnotesize\bfseries
No: 1\\
Name: pe\_\allowbreak{}18\_\allowbreak{}odp\\
Type: string\\
Default: nil
}

{\footnotesize визначено фізичні та екологічні небезпеки, які можуть призвести до потенційного пошкодження компонентів системи на об'єкті;}


{\footnotesize\bfseries
No: 2\\
Name: pe\_\allowbreak{}18\_\allowbreak{}01\\
Type: string\\
Default: nil
}

{\footnotesize компоненти системи розміщені в межах об'єкта так, щоб мінімізувати потенційну шкоду від  <PE-18\_\allowbreak{}ODP фізичні та екологічні небезпеки> і звести до мінімуму можливість несанкціонованого доступу.}




\subsubsection{МІСЦЕ РОЗМІЩЕННЯ ОБ'ЄКТА (PE-18(1)) [Вилучено]}
[Виключено: включено до PE-23].

\vspace{10pt}
\textit{Немає параметрів для цього контролю.}
\vspace{10pt}


\subsection{ВИТІК ІНФОРМАЦІЇ (PE-19)}


\vspace{10pt}
\textit{Немає параметрів для цього контролю.}
\vspace{10pt}


\subsubsection{НАЦІОНАЛЬНІ ПОЛІТИКИ ТА ПРОЦЕДУРИ ЩОДО ПЕМВ (PE-19(1))}
компоненти системи захищені відповідно до національних політик і процедур щодо ПЕМВ;\\ 
передача даних захищається відповідно до національних політик і процедур щодо ПЕМВ;\\ 
мережі захищені відповідно до національних політик і процедур щодо ПЕМВ.

\vspace{10pt}
{\footnotesize\bfseries
No: 1\\
Name: pe\_\allowbreak{}19\_\allowbreak{}1\_\allowbreak{}01\\
Type: string\\
Default: nil
}

{\footnotesize компоненти системи захищені відповідно до національних політик і процедур щодо ПЕМВ;}


{\footnotesize\bfseries
No: 2\\
Name: pe\_\allowbreak{}19\_\allowbreak{}1\_\allowbreak{}02\\
Type: string\\
Default: nil
}

{\footnotesize передача даних захищається відповідно до національних політик і процедур щодо ПЕМВ;}


{\footnotesize\bfseries
No: 3\\
Name: pe\_\allowbreak{}19\_\allowbreak{}1\_\allowbreak{}03\\
Type: string\\
Default: nil
}

{\footnotesize мережі захищені відповідно до національних політик і процедур щодо ПЕМВ.}




\subsection{МОНІТОРИНГ І ВІДСТЕЖЕННЯ АКТИВІВ (PE-20)}


\vspace{10pt}
{\footnotesize\bfseries
No: 1\\
Name: pe\_\allowbreak{}20\_\allowbreak{}odp\_\allowbreak{}1\\
Type: string\\
Default: nil
}

{\footnotesize визначено технології, які будуть використовуватися для відстеження та моніторингу місцезнаходження та переміщення активів;}


{\footnotesize\bfseries
No: 2\\
Name: pe\_\allowbreak{}20\_\allowbreak{}odp\_\allowbreak{}2\\
Type: string\\
Default: nil
}

{\footnotesize визначено активи, місцезнаходження та переміщення яких необхідно відстежувати та моніторити;}


{\footnotesize\bfseries
No: 3\\
Name: pe\_\allowbreak{}20\_\allowbreak{}odp\_\allowbreak{}3\\
Type: string\\
Default: nil
}

{\footnotesize визначено контрольовані зони, в межах яких місцезнаходження та переміщення активів підлягають відстеженню та моніторингу;}


{\footnotesize\bfseries
No: 4\\
Name: pe\_\allowbreak{}20\_\allowbreak{}01\\
Type: string\\
Default: nil
}

{\footnotesize <PE-20\_\allowbreak{}ODP[01] технології> використовуються для відстеження та моніторингу місцезнаходження і переміщення <PE383 зони>.}




\subsection{ЗАХИСТ ВІД ЕЛЕКТРОМАГНІТНОГО ІМПУЛЬСУ (PE-21)}


\vspace{10pt}
{\footnotesize\bfseries
No: 1\\
Name: pe\_\allowbreak{}21\_\allowbreak{}odp\_\allowbreak{}1\\
Type: string\\
Default: nil
}

{\footnotesize визначено заходи захисту від пошкодження електромагнітними імпульсами;}


{\footnotesize\bfseries
No: 2\\
Name: pe\_\allowbreak{}21\_\allowbreak{}odp\_\allowbreak{}2\\
Type: string\\
Default: nil
}

{\footnotesize визначено систему та компоненти системи, що потребують захисту від пошкодження електромагнітними імпульсами;}


{\footnotesize\bfseries
No: 3\\
Name: pe\_\allowbreak{}21\_\allowbreak{}01\\
Type: string\\
Default: nil
}

{\footnotesize <PE-21\_\allowbreak{}ODP[01] заходи захисту>  застосовуються проти пошкодження електромагнітними імпульсами для <PE-21\_\allowbreak{}ODP[02] системи та компонентів системи>.}




\subsection{МАРКУВАННЯ КОМПОНЕНТІВ (PE-22)}


\vspace{10pt}
{\footnotesize\bfseries
No: 1\\
Name: pe\_\allowbreak{}22\_\allowbreak{}odp\\
Type: string\\
Default: nil
}

{\footnotesize визначено апаратні компоненти системи, які підлягають маркуванню із зазначенням рівня впливу або класифікації інформації, яку дозволено обробляти, зберігати або передавати за допомогою апаратного компонента;}


{\footnotesize\bfseries
No: 2\\
Name: pe\_\allowbreak{}22\_\allowbreak{}01\\
Type: string\\
Default: nil
}

{\footnotesize <PE-22\_\allowbreak{}ODP апаратні компоненти системи>  позначаються із зазначенням рівня впливу або класифікації інформації, яку дозволено обробляти, зберігати або передавати за допомогою апаратного компонента.}




\subsection{РОЗТАШУВАННЯ ОБ'ЄКТА (PE-23)}
a. Місце розташування або ділянка об'єкта, де знаходиться система, планується з урахуванням фізичних та екологічних небезпек;\\ 
b. Для існуючих об'єктів фізичні та екологічні небезпеки враховуються в стратегії управління ризиками організації.

\vspace{10pt}
{\footnotesize\bfseries
No: 1\\
Name: pe\_\allowbreak{}23\_\allowbreak{}a\\
Type: string\\
Default: nil
}

{\footnotesize місце розташування або ділянка об'єкта, де знаходиться система, планується з урахуванням фізичних та екологічних небезпек;}


{\footnotesize\bfseries
No: 2\\
Name: pe\_\allowbreak{}23\_\allowbreak{}b\\
Type: string\\
Default: nil
}

{\footnotesize для існуючих об'єктів фізичні та екологічні небезпеки враховуються в стратегії управління ризиками організації.}




\section{PL}
Клас заходів захисту PL --- ПЛАНУВАННЯ БЕЗПЕКИ

Цей клас регламентує розробку, документування та регулярне оновлення планів захисту інформації, архітектури безпеки та приватності.

\textbf{Перелік заходів захисту:} Політики та процедури планування безпеки (PL-1); Плани захисту інформації та персональних даних (PL-2); Керування ризиками ланцюга диверсифікація постачальників (PL-2(1)) [Вилучено]; Плани захисту інформації функціональна архітектура (PL-2(2)) [Вилучено]; Оновлення планів захисту інформації та персональних даних (PL-3) [Вилучено]; Правила поведінки (PL-4); Обмеження на соціальні медіа та мережу (PL-4(1)); Оцінка впливу на приватність (PL-5) [Вилучено]; Планування діяльності, пов'язаної з безпекою (PL-6) [Вилучено]; Концепція експлуатації (PL-7); Архітектура безпеки та приватності (PL-8); «глибока оборона» (PL-8(1)); Різноманітність постачальників (PL-8(2)); Централізоване управління (PL-9); Вибір базового профілю безпеки (PL-10); Налаштування базового профілю безпеки (PL-11).

\subsection{ПОЛІТИКИ ТА ПРОЦЕДУРИ ПЛАНУВАННЯ БЕЗПЕКИ (PL-1)}
a. Розробити, задокументувати та поширити серед [Призначення: визначеного організацією персоналу або ролей]:\\ 
1. 2. [вибір (один або декілька): рівень організації; рівень місії/бізнес процесу; рівень системи] політику планування, яка:\\ 
(a) містить мету, сферу застосування, ролі, обов'язки, відповідальність керівництва, координацію між організаційними підрозділами та системою контролю (complaince);\\ 
(b) відповідає чинному законодавству, виконавчим наказам, директивам, нормам, політикам, стандартам та керівним принципам; процедури, що полегшують здійснення планування політики безпеки та приватності й пов'язані з ними заходи.\\ 
b. Призначити [Призначення: визначену організацією посадову особу] для управління політикою та процедурами планування політики безпеки та приватності.\\ 
c. Переглядати та оновлювати поточне планування:\\ 
1. політики планування безпеки та приватності [Призначення:з визначеною організацією частотою] та наступні [Призначення: визначені організацією події];\\ 
2. поточні процедури планування політики [Призначення: з визначеною організацією [Призначення: визначені організацією події]. безпеки та приватності частотою] та наступні

\vspace{10pt}
{\footnotesize\bfseries
No: 1\\
Name: pl\_\allowbreak{}1\_\allowbreak{}odp\_\allowbreak{}01\\
Type: list\\
Default: [\textquotedbl{}admin\textquotedbl{}, \textquotedbl{}security\_\allowbreak{}officer\textquotedbl{}]
}

{\footnotesize Визначено персонал або ролі, до яких має бути доведена політика планування безпеки}


{\footnotesize\bfseries
No: 2\\
Name: pl\_\allowbreak{}1\_\allowbreak{}odp\_\allowbreak{}02\\
Type: list\\
Default: [\textquotedbl{}admin\textquotedbl{}, \textquotedbl{}security\_\allowbreak{}officer\textquotedbl{}]
}

{\footnotesize Визначено персонал або ролі, на які поширюватимуться процедури планування безпеки}


{\footnotesize\bfseries
No: 3\\
Name: pl\_\allowbreak{}1\_\allowbreak{}odp\_\allowbreak{}03\\
Type: string\\
Default: nil
}

{\footnotesize Вибрано одне або декілька з наступних ЗНАЧЕНЬ ПАРАМЕТРІВ: \{рівень організації; рівень місії/бізнеспроцесу; рівень системи\}}


{\footnotesize\bfseries
No: 4\\
Name: pl\_\allowbreak{}1\_\allowbreak{}odp\_\allowbreak{}04\\
Type: list\\
Default: [\textquotedbl{}admin\textquotedbl{}, \textquotedbl{}security\_\allowbreak{}officer\textquotedbl{}]
}

{\footnotesize Визначено посадову особу, яка керуватиме політикою та процедурами планування безпеки}


{\footnotesize\bfseries
No: 5\\
Name: pl\_\allowbreak{}1\_\allowbreak{}odp\_\allowbreak{}05\\
Type: list\\
Default: [\textquotedbl{}default\_\allowbreak{}deny\_\allowbreak{}rule\textquotedbl{}, \textquotedbl{}abac\_\allowbreak{}rule\_\allowbreak{}1\textquotedbl{}]
}

{\footnotesize Визначено періодичність перегляду та оновлення поточної політики планування безпеки}


{\footnotesize\bfseries
No: 6\\
Name: pl\_\allowbreak{}1\_\allowbreak{}odp\_\allowbreak{}06\\
Type: list\\
Default: [\textquotedbl{}default\_\allowbreak{}deny\_\allowbreak{}rule\textquotedbl{}, \textquotedbl{}abac\_\allowbreak{}rule\_\allowbreak{}1\textquotedbl{}]
}

{\footnotesize Є події, які потребують перегляду та оновлення поточної політики планування безпеки}


{\footnotesize\bfseries
No: 7\\
Name: pl\_\allowbreak{}1\_\allowbreak{}odp\_\allowbreak{}07\\
Type: string\\
Default: \textquotedbl{}щорічно\textquotedbl{}
}

{\footnotesize Визначена частота, з якою переглядаються та оновлюються поточні процедури планування безпеки}




\subsection{ПЛАНИ ЗАХИСТУ ІНФОРМАЦІЇ ТА ПЕРСОНАЛЬНИХ ДАНИХ (PL-2)}
a. Розробити план захисту інформації та персональних даних для інформаційної системи, який:\\ 
1. узгоджується з архітектурою підприємства організації;\\ 
2. чітко визначає складові компоненти системи;\\ 
3. описує оперативний контекст інформаційної системи з точки зору завдань та процесів;\\ 
4. визначає осіб, які виконують системні ролі та обов'язки;\\ 
5. визначає тип інформації, яка обробляється, зберігається та передається системою;\\ 
6. надає огляд вимог безпеки та приватності інформаційної системи;\\ 
7. описує будь які конкретні загрози системи, які викликають стурбованість організації;\\ 
8. надає результати оцінки ризику конфіденційності для систем, в яких обробляються персональні дані;\\ 
9. описує робоче середовище інформаційної системи та будь які залежності від систем або компонентів систем або підключень до таких систем та їх компонентів;\\ 
10. надає огляд вимог безпеки та конфіденційності системи;\\ 
11. визначає будь які відповідні контрольні базові рівні або накладання, якщо вони застосовуються;\\ 
12. описує чинні або заплановані заходи щодо забезпечення безпеки та приватності, включно з обґрунтуванням рішень щодо налаштування\\ 
13. включає виявлення ризиків для архітектури безпеки і приватності, а також проєктних рішень;\\ 
14. включає дії, пов'язані з безпекою та конфіденційністю, які впливають на систему, виконання яких вимагає планування та координацію з [Призначення: визначені організацією окремі особи або групи];\\ 
15. розглядається та затверджується уповноваженою посадовою особою або призначеним представником до початку реалізації плану.\\ 
b. Поширити копії планів захисту інформації та персональних даних і повідомляти про подальші зміни планів серед [Призначення:визначеного організацією персоналу або ролей].\\ 
c. Переглядати плани захисту інформації та персональних даних [Призначення: з визначеною організацією частотою].\\ 
d. Оновлювати плани захисту інформації та персональних даних для врахування змін в інформаційній системі й робочому середовищі або проблем, виявлених у ході реалізації або оцінювання заходів безпеки та приватності.\\ 
e. забезпечити захист планів захисту інформації та персональних даних від несанкціонованого розголошення та змін.

\vspace{10pt}
{\footnotesize\bfseries
No: 1\\
Name: pl\_\allowbreak{}2\_\allowbreak{}odp\_\allowbreak{}01\\
Type: list\\
Default: [\textquotedbl{}admin\textquotedbl{}, \textquotedbl{}security\_\allowbreak{}officer\textquotedbl{}]
}

{\footnotesize Призначені особи або групи, з якими пов'язана діяльність з безпекою та конфіденційністю, що впливає на систему, яка потребує планування та координації}


{\footnotesize\bfseries
No: 2\\
Name: pl\_\allowbreak{}2\_\allowbreak{}odp\_\allowbreak{}02\\
Type: list\\
Default: [\textquotedbl{}admin\textquotedbl{}, \textquotedbl{}security\_\allowbreak{}officer\textquotedbl{}]
}

{\footnotesize Призначено персонал або ролі для отримання розповсюджуваних копій планів захисту інформації та конфіденційності системи}


{\footnotesize\bfseries
No: 3\\
Name: pl\_\allowbreak{}2\_\allowbreak{}odp\_\allowbreak{}03\\
Type: list\\
Default: [\textquotedbl{}admin\textquotedbl{}, \textquotedbl{}security\_\allowbreak{}officer\textquotedbl{}]
}

{\footnotesize Визначено періодичність перегляду інформації та конфіденційності системи; планів захисту PL-02a.01[01] розроблено план захисту інформації, архітектурі підприємства організації; який відповідає PL-02a.01[02] розроблено план конфіденційності для системи, відповідає архітектурі підприємства організації; який PL-02a.02[01] розроблено план захисту інформації, який чітко визначає складові компоненти системи; PL-02a.02[02] розроблено для системи план забезпечення конфіденційності, який чітко визначає складові компоненти системи; PL-02a.03[01] розроблено план захисту інформації системи, який описує операційний контекст системи з точки зору місії та бізнеспроцесів; PL-02a.03[02] розроблено для системи план забезпечення конфіденційності, який описує операційний контекст системи з точки зору місії та бізнес-процесів; PL-02a.04[01] розроблено план захисту інформації, який визначає осіб, що виконують системні ролі та обов'язки; PL-02a.04[02] розроблено для системи план забезпечення конфіденційності, який визначає осіб, що виконують ролі та обов'язки в системі; PL-02a.05[01] розроблено план захисту інформації, який визначає типи інформації, що обробляється, зберігається та передається системою; PL-02a.05[02] розроблено план конфіденційності для системи, який визначає типи інформації, що обробляється, зберігається та передається системою; PL-02a.06[01] розроблено план захисту інформації, який передбачає категоризацію безпеки системи, включаючи відповідне обґрунтування; PL-02a.06[02] розроблено для системи план забезпечення конфіденційності, який передбачає категоризацію системи за рівнем безпеки, включаючи обґрунтування; PL-02a.07[01] розроблено план захисту інформації, який описує будь-які конкретні загрози для системи, що викликають потенційні ризики для рганізації; PL-02a.07[02] розроблено план конфіденційності для системи, який описує будь-які конкретні загрози для системи, що викликають потенційні ризики для організації; PL-02a.08[01] розроблено план захисту інформації, який містить результати оцінки ризиків конфіденційності для систем, що обробляють інформацію, яка ідентифікує особу; PL-02a.08[02] розроблено для системи план забезпечення конфіденційності, який містить результати оцінки ризиків конфіденційності для систем, що обробляють інформацію, яка ідентифікує особу; PL-02a.09[01] розроблено план захисту інформації, який описує операційне середовище системи та будь-які залежності або зв'язки з іншими системами чи компонентами системи; PL-02a.09[02] розроблено для системи план забезпечення конфіденційності, який описує робоче середовище системи та будь-які залежності або зв'язки з іншими системами чи компонентами системи; PL-02a.10[01] розроблено план захисту інформації, який містить огляд вимог до безпеки системи; PL-02a.10[02] розроблено для системи план забезпечення конфіденційності, який містить огляд вимог до конфіденційності системи; PL-02a.11[01] розроблено план захисту інформації, який визначає будь-які відповідні базові рівні контролю або обмеження, якщо такі є; PL-02a.11[02] розроблено для системи план забезпечення конфіденційності, який визначає будь-які відповідні базові рівні контролю або обмеження, якщо такі є; PL-02a.12[01] розроблено план захисту інформації, який описує наявні або заплановані засоби контролю для виконання вимог безпеки, включаючи обґрунтування будь-яких рішень щодо адаптації; PL-02a.12[02] розроблено для системи план забезпечення конфіденційності, який описує наявні або заплановані засоби контролю для виконання вимог щодо конфіденційності, включаючи обґрунтування будь-яких рішень, пов'язаних з адаптацією; PL-02a.13[01] розроблено план захисту інформації, який включає визначення ризиків для архітектури безпеки та проектних рішень; PL-02a.13[02] розроблено план забезпечення конфіденційності для системи, який включає визначення ризиків для архітектури конфіденційності та проектних рішень; PL-02a.14[01] розроблено захисту інформації, який включає діяльність, пов'язану з безпекою, що впливає на систему і потребує планування та координації з окремими особами або групами; PL-02a.14[02] розроблено для системи план забезпечення конфіденційності, який включає діяльність, пов'язану з конфіденційністю, що впливає на систему і потребує планування та координації з окремими особами або групами; PL-02a.15[01] розроблено план захисту інформації, який розглядається та затверджується уповноваженою посадовою особою або призначеним представником до початку реалізації плану; PL-02a.15[02] розроблено план забезпечення конфіденційності для системи, який перевіряється та затверджується уповноваженою посадовою особою або призначеним представником перед впровадженням плану. PL-02b.[01] розповсюджуються копії персонал або ролі>; PL-02b.[02] повідомляються наступні зміни до планів персонал або ролі; PL-02c. переглядаються плани частота; планів серед <PL-02\_\allowbreak{}ODP[02] PL-02d.[01] оновлюються плани відповідно до змін у системі та середовищі діяльності; PL-02d.[02] оновлюються плани для вирішення проблем, виявлених під час реалізації плану; PL-02d.[03] оновлюються плани для вирішення проблем, виявлених під час контрольних оцінок; PL-02e.[01] захищені плани від несанкціонованого розголошення; PL-02e.[02] захищені плани від несанкціонованої модифікації}




\subsubsection{КЕРУВАННЯ РИЗИКАМИ ЛАНЦЮГА ДИВЕРСИФІКАЦІЯ ПОСТАЧАЛЬНИКІВ (PL-2(1)) [Вилучено]}
[Вилучено: включено до PL-7]

\vspace{10pt}
\textit{Немає параметрів для цього контролю.}
\vspace{10pt}


\subsubsection{ПЛАНИ ЗАХИСТУ ІНФОРМАЦІЇ ФУНКЦІОНАЛЬНА АРХІТЕКТУРА (PL-2(2)) [Вилучено]}
[Вилучено: включено до PL-8]

\vspace{10pt}
\textit{Немає параметрів для цього контролю.}
\vspace{10pt}


\subsection{ОНОВЛЕННЯ ПЛАНІВ ЗАХИСТУ ІНФОРМАЦІЇ ТА ПЕРСОНАЛЬНИХ ДАНИХ (PL-3) [Вилучено]}
[Вилучено: включено до PL-2]

\vspace{10pt}
\textit{Немає параметрів для цього контролю.}
\vspace{10pt}


\subsection{ПРАВИЛА ПОВЕДІНКИ (PL-4)}
a. Створити та надати особам, що потребують доступу до інформаційної системи, правила, які описують їхні обов'язки й очікувану поведінку щодо інформації та використання інформаційної системи, безпеки та приватності.\\ 
b. Отримати документальне підтвердження від таких осіб про те, що вони прочитали, зрозуміли та погодилися дотримуватися правил поведінки, перш ніж дозволяти доступ до інформації та інформаційної системи.\\ 
c. Переглядати й оновлювати правила поведінки [Призначення: з визначеною організацією частотою].\\ 
d. Вимагати від осіб, які підписали попередню версію правил поведінки, перечитати та повторно підписати правила [Вибір (один або декілька): [Призначення: з визначеною організацією частотою]; коли правила переглядаються чи оновлюються].

\vspace{10pt}
{\footnotesize\bfseries
No: 1\\
Name: pl\_\allowbreak{}4\_\allowbreak{}odp\_\allowbreak{}01\\
Type: string\\
Default: \textquotedbl{}щорічно\textquotedbl{}
}

{\footnotesize Визначено періодичність перегляду та оновлення правил поведінки}


{\footnotesize\bfseries
No: 2\\
Name: pl\_\allowbreak{}4\_\allowbreak{}odp\_\allowbreak{}02\\
Type: list\\
Default: [\textquotedbl{}default\_\allowbreak{}deny\_\allowbreak{}rule\textquotedbl{}, \textquotedbl{}abac\_\allowbreak{}rule\_\allowbreak{}1\textquotedbl{}]
}

{\footnotesize Вибрано одне або декілька з наступних ЗНАЧЕНЬ ПАРАМЕТРА:\{частота; коли правила переглядаються або оновлюються\}}


{\footnotesize\bfseries
No: 3\\
Name: pl\_\allowbreak{}4\_\allowbreak{}odp\_\allowbreak{}03\\
Type: list\\
Default: [\textquotedbl{}admin\textquotedbl{}, \textquotedbl{}security\_\allowbreak{}officer\textquotedbl{}]
}

{\footnotesize Визначена періодичність перегляду та повторного підтвердження правил поведінки (якщо вибрано); PL-04a.[01] встановлені правила, які описують обов'язки та очікувану поведінку щодо використання інформації та системи, безпеки та конфіденційності для осіб, яким потрібен доступ до системи; PL-04a.[02] надаються правила, які описують обов'язки та очікувану поведінку щодо використання інформації та системи, безпеки та конфіденційності, особам, які отримують доступ до системи; PL-04b. отримано перед наданням доступу до інформації та системи задокументоване підтвердження від таких осіб про те, що вони прочитали, зрозуміли та згодні дотримуватися правил поведінки; PL-04c. переглядаються та оновлюються правила поведінки частота; PL-04d. потрібно особам, які визнали попередню версію правил поведінки, прочитати та повторно визнати ЗНАЧЕННЯ ВИБРАНОГО ПАРАМЕТРА(ів)}




\subsubsection{ОБМЕЖЕННЯ НА СОЦІАЛЬНІ МЕДІА ТА МЕРЕЖУ (PL-4(1))}


\vspace{10pt}
{\footnotesize\bfseries
No: 1\\
Name: pl\_\allowbreak{}4\_\allowbreak{}1\_\allowbreak{}a\\
Type: list\\
Default: [\textquotedbl{}default\_\allowbreak{}deny\_\allowbreak{}rule\textquotedbl{}, \textquotedbl{}abac\_\allowbreak{}rule\_\allowbreak{}1\textquotedbl{}]
}

{\footnotesize Включають правила поведінки обмеження на використання соціальних медіа, соціальних мереж та зовнішніх сайтів/додатків}


{\footnotesize\bfseries
No: 2\\
Name: pl\_\allowbreak{}4\_\allowbreak{}1\_\allowbreak{}b\\
Type: list\\
Default: [\textquotedbl{}default\_\allowbreak{}deny\_\allowbreak{}rule\textquotedbl{}, \textquotedbl{}abac\_\allowbreak{}rule\_\allowbreak{}1\textquotedbl{}]
}

{\footnotesize Включають правила поведінки обмеження на розміщення інформації про організацію на публічних веб-сайтах}


{\footnotesize\bfseries
No: 3\\
Name: pl\_\allowbreak{}4\_\allowbreak{}1\_\allowbreak{}c\\
Type: list\\
Default: [\textquotedbl{}admin\textquotedbl{}, \textquotedbl{}security\_\allowbreak{}officer\textquotedbl{}]
}

{\footnotesize Включають правила поведінки обмеження на використання наданих організацією ідентифікаторів (наприклад, адрес електронної пошти) та секретів автентифікації (наприклад, паролів) для створення облікових записів на зовнішніх сайтах/додатках. ПОТЕНЦІЙНІ МЕТОДИ ТА ОБ'ЄКТИ ОЦІНКИ:}




\subsection{ОЦІНКА ВПЛИВУ НА ПРИВАТНІСТЬ (PL-5) [Вилучено]}
[Вилучено: включено до RА-8]

\vspace{10pt}
\textit{Немає параметрів для цього контролю.}
\vspace{10pt}


\subsection{ПЛАНУВАННЯ ДІЯЛЬНОСТІ, ПОВ'ЯЗАНОЇ З БЕЗПЕКОЮ (PL-6) [Вилучено]}
[Вилучено: включено до PL-2]

\vspace{10pt}
\textit{Немає параметрів для цього контролю.}
\vspace{10pt}


\subsection{КОНЦЕПЦІЯ ЕКСПЛУАТАЦІЇ (PL-7)}
a. Розробити концепцію експлуатації інформаційної системи, яка описує, як організація має намір керувати системою з погляду забезпечення безпеки та приватності інформації.\\ 
b. Переглядати й оновлювати концепцію експлуатації [Призначення: з визначеною організацією частотою].

\vspace{10pt}
{\footnotesize\bfseries
No: 1\\
Name: pl\_\allowbreak{}7\_\allowbreak{}odp\\
Type: string\\
Default: \textquotedbl{}щорічно\textquotedbl{}
}

{\footnotesize Визначена періодичність перегляду концепцію експлуатації системи; та оновлення PL-07a. розроблено концепцію експлуатації системи, що описує, як організація має намір експлуатувати систему з точки зору інформаційної безпеки та конфіденційності; PL-07b. переглядається та оновлюється концепція експлуатації системи частота}




\subsection{АРХІТЕКТУРА БЕЗПЕКИ ТА ПРИВАТНОСТІ (PL-8)}
a.\\ 
b. Розробити архітектуру безпеки та приватності для інформаційної системи, яка:\\ 
1. характеризує методологію, вимоги та підходи, які слід вживати для забезпечення конфіденційності, цілісності та доступності інформації, що циркулює в системі;\\ 
2. характеризує методологію, вимоги та підхід до обробки персональних даних для мінімізації ризику їх втрати;\\ 
3. характеризує, як архітектури безпеки та приватності інтегруються в архітектуру підприємства;\\ 
4. характеризує будь-які припущення, що пов'язані з безпекою та приватністю, щодо зовнішніх служб і залежності від них. Переглядати й оновлювати архітектуру безпеки та приватності [Призначення:з визначеною організацією частотою], щоб відображати оновлення в архітектурі підприємства.\\ 
c. Відображати заплановані зміни архітектури плану безпеки та приватності, концепції експлуатації інформаційної системи, аналізу критичності, організаційних заходах, постачань та закупівель.

\vspace{10pt}
{\footnotesize\bfseries
No: 1\\
Name: pl\_\allowbreak{}8\_\allowbreak{}01\\
Type: string\\
Default: nil
}

{\footnotesize АРХІТЕКТУРА БЕЗПЕКИ ТА ПРИВАТНОСТІ МЕТА ОЦІНКИ: Визначити, чи:}


{\footnotesize\bfseries
No: 2\\
Name: pl\_\allowbreak{}8\_\allowbreak{}odp\\
Type: list\\
Default: [\textquotedbl{}admin\textquotedbl{}, \textquotedbl{}security\_\allowbreak{}officer\textquotedbl{}]
}

{\footnotesize Потрібно переглядати та оновлювати частоту, відображати зміни в архітектурі підприємства; PL-08a.01 описує архітектура безпеки системи вимоги та підходи до захисту конфіденційності, цілісності та доступності організаційної інформації; PL-08a.02 описує архітектура конфіденційності вимоги та підходи до обробки персональних даних з метою мінімізації ризиків для приватного життя людей; PL-08a.03[01] описує архітектура безпеки системи те, як вона інтегрована в архітектуру підприємства та підтримує її; PL-08a.03[02] описує архітектура конфіденційності системи те, як вона інтегрована в архітектуру підприємства та підтримує її; PL-08a.04[01] описує архітектура безпеки системи будь-які припущення та залежності від зовнішніх систем та сервісів; PL-08a.04[02] описує архітектура конфіденційності системи будь-які припущення та залежності від зовнішніх систем і сервісів; PL-08b. переглядаються та оновлюються зміни в архітектурі підприємства частота для відображення змін в архітектурі підприємства; PL-08c.[01] заплановані зміни в архітектурі відображені в плані безпеки; PL-08c.[02] відображені заплановані конфіденційності; PL-08c.[03] заплановані зміни архітектури відображені діяльності концепції експлуатації системи; PL-08c.[04] заплановані зміни критичності; PL-08c.[05] відображені заплановані зміни архітектури в організаційних процедурах; PL-08c.[06] заплановані зміни в архітектурі відображаються на закупівлях та придбанні. в зміни архітектурі в архітектурі в щоб плані в Концепції відображені в аналізі}




\subsubsection{«ГЛИБОКА ОБОРОНА» (PL-8(1))}


\vspace{10pt}
{\footnotesize\bfseries
No: 1\\
Name: pl\_\allowbreak{}8\_\allowbreak{}1\_\allowbreak{}a\_\allowbreak{}01\\
Type: string\\
Default: nil
}

{\footnotesize Архітектура безпеки системи розроблена з використанням підходу «глибокої оборони», який розподіляє елементи керування за місцями та архітектурними рівнями}


{\footnotesize\bfseries
No: 2\\
Name: pl\_\allowbreak{}8\_\allowbreak{}1\_\allowbreak{}a\_\allowbreak{}02\\
Type: string\\
Default: nil
}

{\footnotesize Архітектура конфіденційності системи розроблена з використанням підходу глибокого захисту, який розподіляє елементи керування за місцями та архітектурними рівнями}


{\footnotesize\bfseries
No: 3\\
Name: pl\_\allowbreak{}8\_\allowbreak{}1\_\allowbreak{}b\_\allowbreak{}01\\
Type: string\\
Default: \textquotedbl{}автоматизований засіб моніторингу\textquotedbl{}
}

{\footnotesize Архітектура безпеки системи розроблена з використанням підходу «глибокої оборони», який гарантує, що виділені засоби контролю працюють скоординовано і взаємно підсилюють один одного}


{\footnotesize\bfseries
No: 4\\
Name: pl\_\allowbreak{}8\_\allowbreak{}1\_\allowbreak{}b\_\allowbreak{}02\\
Type: string\\
Default: \textquotedbl{}автоматизований засіб моніторингу\textquotedbl{}
}

{\footnotesize Архітектура конфіденційності системи розроблена з використанням підходу «глибокої оборони», який гарантує, що виділені засоби контролю працюють скоординовано і взаємно підкріплюють один одного}




\subsubsection{РІЗНОМАНІТНІСТЬ ПОСТАЧАЛЬНИКІВ (PL-8(2))}


\vspace{10pt}
{\footnotesize\bfseries
No: 1\\
Name: pl\_\allowbreak{}8\_\allowbreak{}2\_\allowbreak{}01\\
Type: string\\
Default: nil
}

{\footnotesize Потрібно отримувати елементи керування, призначені для локацій та архітектурних рівнів, від різних постачальників}




\subsection{ЦЕНТРАЛІЗОВАНЕ УПРАВЛІННЯ (PL-9)}
Централізовано управляти [Призначення: визначеними організацією організаційними заходами захисту та пов'язаними з ними процесами].

\vspace{10pt}
{\footnotesize\bfseries
No: 1\\
Name: pl\_\allowbreak{}9\_\allowbreak{}01\\
Type: string\\
Default: nil
}

{\footnotesize Здійснюється централізоване управління контролями та пов'язаними з ними процесами}


{\footnotesize\bfseries
No: 2\\
Name: pl\_\allowbreak{}9\_\allowbreak{}odp\\
Type: string\\
Default: \textquotedbl{}автоматизований засіб моніторингу\textquotedbl{}
}

{\footnotesize Визначені засоби контролю безпеки та конфіденційності і пов'язані з ними процеси, якими слід централізовано керувати}




\subsection{ВИБІР БАЗОВОГО ПРОФІЛЮ БЕЗПЕКИ (PL-10)}
Вибрати базовий профіль безпеки для інформаційної системи.

\vspace{10pt}
{\footnotesize\bfseries
No: 1\\
Name: pl\_\allowbreak{}10\_\allowbreak{}01\\
Type: string\\
Default: nil
}

{\footnotesize Вибрано базовий профіль безпеки для системи}




\subsection{НАЛАШТУВАННЯ БАЗОВОГО ПРОФІЛЮ БЕЗПЕКИ (PL-11)}
a. Розробити та поширити на організаційному рівні план програми (концепцію) з інформаційної безпеки, яка:\\ 
1. містить огляд вимог до програми (концепції) безпеки й описує заходи управління програмою інформаційної безпеки та загальних заходів безпеки, які використовуються або плануються для виконання цих вимог;\\ 
2. містить визначення та розподіл ролей, обов'язків, відповідальності керівництва, заходи з координації діяльності організації і забезпечення відповідності вимогам законодавства та іншим нормативним документам;\\ 
3. відображає координацію між організаційними елементами, що відповідають за інформаційну безпеку;\\ 
4. затверджена вищою посадовою особою, що відповідає та підзвітна за управління ризиками, пов'язаними з організаційними операціями (включно з завданнями (місією), функціями, іміджем і репутацією організації), організаційні активи, фізичних осіб, інші організації та державу.\\ 
b. Переглядати та оновлювати загальноорганізаційний план програми (концепцію) інформаційної безпеки [Призначення: з визначеною організацією частотою] та у випадку [Призначення: визначені організацією випадки].\\ 
c. Забезпечити захист плану програми (концепції) інформаційної безпеки від несанкціонованого розкриття та зміни.

\vspace{10pt}
{\footnotesize\bfseries
No: 1\\
Name: pl\_\allowbreak{}11\_\allowbreak{}01\\
Type: list\\
Default: [\textquotedbl{}login\textquotedbl{}, \textquotedbl{}logout\textquotedbl{}, \textquotedbl{}failed\_\allowbreak{}attempt\textquotedbl{}]
}

{\footnotesize Налаштуватовано вибраний базовий профіль застосовуючи вказані дії для налаштування. безпеки,}




\section{PS}
Клас заходів захисту PS --- КАДРОВА БЕЗПЕКА

Цей клас встановлює вимоги до перевірки, надання повноважень та звільнення персоналу з метою мінімізації ризиків інсайдерських загроз.

\textbf{Перелік заходів захисту:} Політика та процедури кадрової безпеки (PS-1); Визначення посадового ризику (PS-2); Перевірка персоналу (PS-3); Звільнення персоналу (PS-4); Переведення персоналу (PS-5); Угоди про доступ (PS-6); Інформація, що вимагає спеціального захисту (PS-6(1)) [Вилучено]; Безпека зовнішнього персоналу (PS-7); Кадрові санкції (PS-8); Опис позицій (PS-9).

\subsection{Політика та процедури кадрової безпеки (PS-1)}


\vspace{10pt}
{\footnotesize\bfseries
No: 1\\
Name: ps\_\allowbreak{}1\_\allowbreak{}odp\_\allowbreak{}01\\
Type: list\\
Default: [\textquotedbl{}admin\textquotedbl{}, \textquotedbl{}security\_\allowbreak{}officer\textquotedbl{}]
}

{\footnotesize Визначено персонал або ролі, на які поширюється політика кадрової безпеки}


{\footnotesize\bfseries
No: 2\\
Name: ps\_\allowbreak{}1\_\allowbreak{}odp\_\allowbreak{}02\\
Type: list\\
Default: [\textquotedbl{}admin\textquotedbl{}, \textquotedbl{}security\_\allowbreak{}officer\textquotedbl{}]
}

{\footnotesize Визначено персонал або ролі, на які поширюються процедури кадрової безпеки}


{\footnotesize\bfseries
No: 3\\
Name: ps\_\allowbreak{}1\_\allowbreak{}odp\_\allowbreak{}03\\
Type: string\\
Default: nil
}

{\footnotesize Вибрано одне або декілька з наступних ЗНАЧЕНЬ ПАРАМЕТРІВ: \{рівень організації; рівень місії/бізнеспроцесу; рівень системи\}}


{\footnotesize\bfseries
No: 4\\
Name: ps\_\allowbreak{}1\_\allowbreak{}odp\_\allowbreak{}04\\
Type: list\\
Default: [\textquotedbl{}admin\textquotedbl{}, \textquotedbl{}security\_\allowbreak{}officer\textquotedbl{}]
}

{\footnotesize Визначено посадову особу, яка керуватиме політикою та процедурами кадрової безпеки}


{\footnotesize\bfseries
No: 5\\
Name: ps\_\allowbreak{}1\_\allowbreak{}odp\_\allowbreak{}05\\
Type: list\\
Default: [\textquotedbl{}default\_\allowbreak{}deny\_\allowbreak{}rule\textquotedbl{}, \textquotedbl{}abac\_\allowbreak{}rule\_\allowbreak{}1\textquotedbl{}]
}

{\footnotesize Визначена періодичність перегляду та оновлення поточної політики кадрової безпеки}


{\footnotesize\bfseries
No: 6\\
Name: ps\_\allowbreak{}1\_\allowbreak{}odp\_\allowbreak{}06\\
Type: list\\
Default: [\textquotedbl{}default\_\allowbreak{}deny\_\allowbreak{}rule\textquotedbl{}, \textquotedbl{}abac\_\allowbreak{}rule\_\allowbreak{}1\textquotedbl{}]
}

{\footnotesize Є події, які вимагають перегляду та оновлення поточної політики кадрової безпеки}


{\footnotesize\bfseries
No: 7\\
Name: ps\_\allowbreak{}1\_\allowbreak{}odp\_\allowbreak{}07\\
Type: string\\
Default: \textquotedbl{}щорічно\textquotedbl{}
}

{\footnotesize Визначено періодичність перегляду та оновлення поточних процедур кадрової безпеки}




\subsection{Визначення посадового ризику (PS-2)}


\vspace{10pt}
{\footnotesize\bfseries
No: 1\\
Name: ps\_\allowbreak{}2\_\allowbreak{}odp\\
Type: list\\
Default: [\textquotedbl{}admin\textquotedbl{}, \textquotedbl{}security\_\allowbreak{}officer\textquotedbl{}]
}

{\footnotesize Визначено періодичність перегляду ідентифікаторів посадових ризиків; та оновлення PS-02a. всім посадам в організації присвоєно ідентифікатор ризику; PS-02b. вибрано одне або декілька з наступних ЗНАЧЕНЬ ПАРАМЕТРІВ: \{рівень організації; рівень місії/бізнес-процесу; рівень системи\}; PS-02c. встановлені критерії відбору для осіб, які обіймають посади в організації}




\subsection{Перевірка персоналу (PS-3)}


\vspace{10pt}
{\footnotesize\bfseries
No: 1\\
Name: ps\_\allowbreak{}3\_\allowbreak{}odp\_\allowbreak{}01\\
Type: list\\
Default: []
}

{\footnotesize Визначені умови, що вимагають повторної перевірки осіб}


{\footnotesize\bfseries
No: 2\\
Name: ps\_\allowbreak{}3\_\allowbreak{}odp\_\allowbreak{}02\\
Type: list\\
Default: [\textquotedbl{}admin\textquotedbl{}, \textquotedbl{}security\_\allowbreak{}officer\textquotedbl{}]
}

{\footnotesize Визначена частота повторної перевірки осіб, для яких це показано; PS-03a. проходять особи перевірку перед тим, як надати їм доступ до системи; PS-03b.[01] проходять особи повторну перевірку відповідно до умови, що вимагають повторної перевірки; PS-03b.[02] проводиться повторна перевірка у випадках, коли це зазначено, частота}




\subsection{Звільнення персоналу (PS-4)}


\vspace{10pt}
{\footnotesize\bfseries
No: 1\\
Name: ps\_\allowbreak{}4\_\allowbreak{}odp\_\allowbreak{}01\\
Type: integer\\
Default: 30
}

{\footnotesize Визначено період часу, протягом якого забороняється доступ до системи}


{\footnotesize\bfseries
No: 2\\
Name: ps\_\allowbreak{}4\_\allowbreak{}odp\_\allowbreak{}02\\
Type: list\\
Default: [\textquotedbl{}admin\textquotedbl{}, \textquotedbl{}security\_\allowbreak{}officer\textquotedbl{}]
}

{\footnotesize Визначені теми інформаційної безпеки для обговорення під час проведення співбесід; PS-04a. при звільненні працівника доступ до системи вимикається протягом часового періоду; PS-04b. припиняють дію або анулюють будь-які автентифікатори та облікові дані після припинення трудових відносин з окремими особами; PS-04c. проводяться при звільненні окремих працівників співбесіди, які включають обговорення питань інформаційної безпеки; PS-04d. отримується після звільнення особи все майно, пов'язане з безпекою організаційної системи; PS-04e. зберігається доступ до організаційної інформації та систем, які раніше перебували під контролем особи, що звільняється, після припинення нею трудових відносин}




\subsection{Переведення персоналу (PS-5)}


\vspace{10pt}
{\footnotesize\bfseries
No: 1\\
Name: ps\_\allowbreak{}5\_\allowbreak{}odp\_\allowbreak{}01\\
Type: list\\
Default: [\textquotedbl{}login\textquotedbl{}, \textquotedbl{}logout\textquotedbl{}, \textquotedbl{}failed\_\allowbreak{}attempt\textquotedbl{}]
}

{\footnotesize Визначені дії, які мають бути ініційовані після переведення або перепризначення}


{\footnotesize\bfseries
No: 2\\
Name: ps\_\allowbreak{}5\_\allowbreak{}odp\_\allowbreak{}02\\
Type: list\\
Default: [\textquotedbl{}login\textquotedbl{}, \textquotedbl{}logout\textquotedbl{}, \textquotedbl{}failed\_\allowbreak{}attempt\textquotedbl{}]
}

{\footnotesize Визначено період часу, протягом якого мають бути здійснені дії з переведення або перепризначення після переведення або перепризначення}


{\footnotesize\bfseries
No: 3\\
Name: ps\_\allowbreak{}5\_\allowbreak{}odp\_\allowbreak{}03\\
Type: list\\
Default: [\textquotedbl{}admin\textquotedbl{}, \textquotedbl{}security\_\allowbreak{}officer\textquotedbl{}]
}

{\footnotesize Визначено персонал або ролі, про які необхідно повідомляти, коли осіб призначають на інші посади або переводять на інші посади в організації}


{\footnotesize\bfseries
No: 4\\
Name: ps\_\allowbreak{}5\_\allowbreak{}odp\_\allowbreak{}04\\
Type: list\\
Default: [\textquotedbl{}admin\textquotedbl{}, \textquotedbl{}security\_\allowbreak{}officer\textquotedbl{}]
}

{\footnotesize Визначено період часу, протягом якого необхідно повідомляти визначений організацією персонал або ролі, коли осіб перепризначають або переводять на інші посади в межах організації; PS-05a. переглядаються та підтверджуються поточні потреби в логічних та фізичних дозволах на доступ до систем та об'єктів при перепризначенні або переведенні осіб на інші посади в організації; PS-05b. були ініційовані дії з переведення або перепризначення протягом періоду часу після формальної дії з переведення; PS-05c. змінюється авторизація доступу за необхідності, щоб відповідати будь-яким змінам в оперативних потребах у зв'язку з перепризначенням або переведенням; PS-05d. було повідомлено персонал або ролі протягом часового періоду}




\subsection{Угоди про доступ (PS-6)}


\vspace{10pt}
{\footnotesize\bfseries
No: 1\\
Name: ps\_\allowbreak{}6\_\allowbreak{}odp\_\allowbreak{}01\\
Type: string\\
Default: \textquotedbl{}щорічно\textquotedbl{}
}

{\footnotesize Визначено періодичність перегляду та оновлення угод про доступ}


{\footnotesize\bfseries
No: 2\\
Name: ps\_\allowbreak{}6\_\allowbreak{}odp\_\allowbreak{}02\\
Type: list\\
Default: [\textquotedbl{}admin\textquotedbl{}, \textquotedbl{}security\_\allowbreak{}officer\textquotedbl{}]
}

{\footnotesize Визначена періодичність перепідписання угод про доступ для збереження доступу до інформації організації; PS-06a. розроблені та задокументовані угоди про доступ до систем організації; PS-06b. переглядаються та оновлюються угоди про доступ частота; PS-06c.01 підписують особи, яким потрібен доступ до інформації та систем організації, відповідні угоди про доступ до того, як їм буде надано доступ; PS-06c.02 перепідписують особи, яким потрібен доступ до інформації та систем організації, угоди про доступ для збереження доступу до систем організації, коли угоди про доступ були оновлені чи як частота}




\subsubsection{ІНФОРМАЦІЯ, ЩО ВИМАГАЄ СПЕЦІАЛЬНОГО ЗАХИСТУ (PS-6(1)) [Вилучено]}
[Вилучено: включено до PS-3]

\vspace{10pt}
\textit{Немає параметрів для цього контролю.}
\vspace{10pt}


\subsection{Безпека зовнішнього персоналу (PS-7)}


\vspace{10pt}
{\footnotesize\bfseries
No: 1\\
Name: ps\_\allowbreak{}7\_\allowbreak{}odp\_\allowbreak{}01\\
Type: list\\
Default: [\textquotedbl{}admin\textquotedbl{}, \textquotedbl{}security\_\allowbreak{}officer\textquotedbl{}]
}

{\footnotesize Визначено персонал або ролі, які мають бути повідомлені про будь-які кадрові переведення або звільнення зовнішнього персоналу, який володіє організаційними повноваженнями та/або бейджами, або має системні привілеї}


{\footnotesize\bfseries
No: 2\\
Name: ps\_\allowbreak{}7\_\allowbreak{}odp\_\allowbreak{}02\\
Type: list\\
Default: [\textquotedbl{}admin\textquotedbl{}, \textquotedbl{}security\_\allowbreak{}officer\textquotedbl{}]
}

{\footnotesize Визначено період часу, протягом якого сторонні провайдери повинні повідомляти визначений організацією персонал або ролі про будь-які кадрові переведення або звільнення зовнішнього персоналу, який володіє організаційними повноваженнями та/або бейджами або має системні привілеї; PS-07a. встановлені вимоги до безпеки персоналу, включаючи ролі та обов'язки зовнішніх постачальників послуг у сфері безпеки; PS-07b. зобов'язані зовнішні провайдери дотримуватися політики та процедур кадрової безпеки, встановлених організацією; PS-07c. задокументовані вимоги до безпеки персоналу; PS-07d. зобов'язані зовнішні провайдери повідомляти персонал або ролі про будь-які кадрові переведення або звільнення зовнішнього персоналу, який володіє організаційними повноваженнями та/або бейджами або має системні привілеї протягом часового періоду; PS-07e. контролюється дотримання провайдером вимог щодо безпеки персоналу}




\subsection{Кадрові санкції (PS-8)}


\vspace{10pt}
{\footnotesize\bfseries
No: 1\\
Name: ps\_\allowbreak{}8\_\allowbreak{}odp\_\allowbreak{}01\\
Type: list\\
Default: [\textquotedbl{}admin\textquotedbl{}, \textquotedbl{}security\_\allowbreak{}officer\textquotedbl{}]
}

{\footnotesize Визначено персонал або ролі, про які необхідно повідомляти, коли ініціюється офіційний процес санкцій щодо працівників}




\subsection{ОПИС ПОЗИЦІЙ (PS-9)}


\vspace{10pt}
{\footnotesize\bfseries
No: 1\\
Name: ps\_\allowbreak{}9\_\allowbreak{}02\\
Type: list\\
Default: [\textquotedbl{}admin\textquotedbl{}, \textquotedbl{}security\_\allowbreak{}officer\textquotedbl{}]
}

{\footnotesize Включені функції та обов'язки з безпеки в описи посадових осіб в організації; включені ролі та обов'язки щодо конфіденційності в описи посадових осіб в організації}




\section{RA}
Клас заходів захисту RA --- ОЦІНЮВАННЯ РИЗИКУ

Цей клас забезпечує систематичний підхід до виявлення, аналізу та реагування на ризики, пов'язані з обробкою інформації та функціонуванням системи.

\textbf{Перелік заходів захисту:} Політика та процедури оцінювання ризику (RA-1); Категоріювання безпеки (RA-2); Категоріювання другого рівня (RA-2(1)); Оцінювання ризику (RA-3); Оцінювання постачання (RA-3(1)); Використання інформації з усіх доступних джерел (RA-3(2)); Усвідомлення динамічних загроз (RA-3(3)); Прогностична кібераналітика (RA-3(4)); Оновлення оцінювання ризику (RA-4) [Вилучено]; Сканування вразливостей (RA-5); Сканування вразливостей інструментів (RA-5(1)) [Вилучено]; Оновлення за частотою, перед новим скануванням або при ідентифікації (RA-5(2)); Широта та глибина покриття (RA-5(3)); Виявна інформація (RA-5(4)); Привілейований доступ (RA-5(5)); Сканування тенденцій (RA-5(6)); Автоматизоване виявлення та сповіщення про неавторизовані компоненти (RA-5(7)) [Вилучено]; Огляд журналів аудиту за минулі періоди (RA-5(8)); Сканування вразливостей проникнення (RA-5(9)) [Вилучено]; Заходи протидії технічній розвідці (RA-6); Реагування на ризик (RA-7); Оцінка впливу на приватність (RA-8); Аналіз критичності (RA-9); Активний пошук загроз (RA-10).

\subsection{ПОЛІТИКА ТА ПРОЦЕДУРИ ОЦІНЮВАННЯ РИЗИКУ (RA-1)}
a. Розробити, задокументувати та поширити серед [Призначення: визначеного організацією персоналу або посадових осіб]:\\ 
1. 2. Політику оцінювання ризику, яка:\\ 
(a) містить мету, сферу застосування, ролі, обов'язки, відповідальність керівництва, координацію між організаційними підрозділами та систему контролю відповідності (complaince);\\ 
(b) відповідає чинному законодавству, виконавчим наказам, директивам, нормам, політикам, стандартам і керівним принципам. Процедури, що сприяють здійсненню політики оцінювання ризику та пов'язаних з ними заходів оцінювання ризику.\\ 
b. Призначити [Призначення: визначена організацією посадову особу] для управління політикою та процедурами оцінювання ризику.\\ 
c. Переглядати й оновлювати:\\ 
1. Поточну політику оцінювання організацією частотою]. ризику [Призначення: з визначеною\\ 
2. Поточні процедури оцінювання ризику [Призначення: з визначеною організацією частотою].

\vspace{10pt}
{\footnotesize\bfseries
No: 1\\
Name: ra\_\allowbreak{}1\_\allowbreak{}odp\_\allowbreak{}01\\
Type: list\\
Default: [\textquotedbl{}admin\textquotedbl{}, \textquotedbl{}security\_\allowbreak{}officer\textquotedbl{}]
}

{\footnotesize Визначено персонал або ролі, на які поширюється політика оцінювання ризику}


{\footnotesize\bfseries
No: 2\\
Name: ra\_\allowbreak{}1\_\allowbreak{}odp\_\allowbreak{}02\\
Type: list\\
Default: [\textquotedbl{}admin\textquotedbl{}, \textquotedbl{}security\_\allowbreak{}officer\textquotedbl{}]
}

{\footnotesize Визначено персонал або ролі, на які поширюються процедури, що сприяють здійсненню політики оцінювання ризику}


{\footnotesize\bfseries
No: 3\\
Name: ra\_\allowbreak{}1\_\allowbreak{}odp\_\allowbreak{}03\\
Type: string\\
Default: nil
}

{\footnotesize Вибрано одне або декілька з наступних ЗНАЧЕНЬ ПАРАМЕТРІВ: \{рівень організації; рівень місії/бізнеспроцесу; рівень системи\}}


{\footnotesize\bfseries
No: 4\\
Name: ra\_\allowbreak{}1\_\allowbreak{}odp\_\allowbreak{}04\\
Type: list\\
Default: [\textquotedbl{}admin\textquotedbl{}, \textquotedbl{}security\_\allowbreak{}officer\textquotedbl{}]
}

{\footnotesize Визначена посадова особа, відповідальна за управління політикою та процедурами оцінювання ризику}


{\footnotesize\bfseries
No: 5\\
Name: ra\_\allowbreak{}1\_\allowbreak{}odp\_\allowbreak{}05\\
Type: list\\
Default: [\textquotedbl{}default\_\allowbreak{}deny\_\allowbreak{}rule\textquotedbl{}, \textquotedbl{}abac\_\allowbreak{}rule\_\allowbreak{}1\textquotedbl{}]
}

{\footnotesize Визначена періодичність перегляду та оновлення поточної політики оцінювання ризику}


{\footnotesize\bfseries
No: 6\\
Name: ra\_\allowbreak{}1\_\allowbreak{}odp\_\allowbreak{}06\\
Type: list\\
Default: [\textquotedbl{}default\_\allowbreak{}deny\_\allowbreak{}rule\textquotedbl{}, \textquotedbl{}abac\_\allowbreak{}rule\_\allowbreak{}1\textquotedbl{}]
}

{\footnotesize Є події, які вимагають перегляду та оновлення поточної політики оцінювання ризику}


{\footnotesize\bfseries
No: 7\\
Name: ra\_\allowbreak{}1\_\allowbreak{}odp\_\allowbreak{}07\\
Type: string\\
Default: \textquotedbl{}щорічно\textquotedbl{}
}

{\footnotesize Визначено періодичність перегляду поточних процедур оцінювання ризику}




\subsection{КАТЕГОРІЮВАННЯ БЕЗПЕКИ (RA-2)}
a.\\ 
b. Проводити оцінювання ризику, включно з вірогідністю й величиною шкоди від:\\ 
1. несанкціонованого доступу, використання, розголошення, руйнування, модифікації або знищення інформаційної системи, інформації, яку вона обробляє, зберігає та передає; а також будь-якої пов'язаної інформації;\\ 
2. проблем, пов'язаних з приватністю фізичних осіб, що виникають у результаті обробки персональних даних. Інтегрувати результати оцінювання ризику та рішення з управління ризиками на рівні організації та завдань/процесів з оцінюванням ризиків на рівні інформаційної системи.\\ 
c. Задокументувати результати оцінювання ризику до [Вибір: планів безпеки та приватності; звіту про оцінювання ризику; [Призначення: визначеного організацією документа]].\\ 
d. Переглядати результати оцінювання ризиків [Призначення: з визначеною організацією частотою].\\ 
e. Поширити результати оцінювання ризику серед [Призначення: визначеного організацією персоналу або посадових осіб].\\ 
f. Оновлювати оцінювання ризику [Призначення: з визначеною організацією частотою] або коли є суттєві зміни в інформаційній системі, її робочому середовищі чи інших умовах, які можуть вплинути на стан безпеки або приватність інформаційної системи.

\vspace{10pt}
\textit{Немає параметрів для цього контролю.}
\vspace{10pt}


\subsubsection{КАТЕГОРІЮВАННЯ ДРУГОГО РІВНЯ (RA-2(1))}


\vspace{10pt}
{\footnotesize\bfseries
No: 1\\
Name: ra\_\allowbreak{}2\_\allowbreak{}1\_\allowbreak{}01\\
Type: string\\
Default: nil
}

{\footnotesize Проводиться категоріювання другого рівня для інформаційних систем організації з метою отримання додаткової деталізації рівнів критичності системи}




\subsection{ОЦІНЮВАННЯ РИЗИКУ (RA-3)}


\vspace{10pt}
{\footnotesize\bfseries
No: 1\\
Name: ra\_\allowbreak{}3\_\allowbreak{}odp\_\allowbreak{}01\\
Type: string\\
Default: nil
}

{\footnotesize Вибрано одне з наступних ЗНАЧЕНЬ ПАРАМЕТРІВ: \{плани безпеки та приватності; звіт про оцінювання ризику; документ\}}


{\footnotesize\bfseries
No: 2\\
Name: ra\_\allowbreak{}3\_\allowbreak{}odp\_\allowbreak{}02\\
Type: string\\
Default: nil
}

{\footnotesize Визначено документ, в якому мають бути задокументовані результати оцінювання ризику (якщо вони не задокументовані в планах безпеки та приватності або в звіті про оцінювання ризику) (якщо вибрано)}


{\footnotesize\bfseries
No: 3\\
Name: ra\_\allowbreak{}3\_\allowbreak{}odp\_\allowbreak{}03\\
Type: string\\
Default: \textquotedbl{}щорічно\textquotedbl{}
}

{\footnotesize Визначено періодичність оцінювання ризику}


{\footnotesize\bfseries
No: 4\\
Name: ra\_\allowbreak{}3\_\allowbreak{}odp\_\allowbreak{}04\\
Type: list\\
Default: [\textquotedbl{}admin\textquotedbl{}, \textquotedbl{}security\_\allowbreak{}officer\textquotedbl{}]
}

{\footnotesize Визначено персонал або ролі, до яких мають бути доведені результати оцінювання ризику}


{\footnotesize\bfseries
No: 5\\
Name: ra\_\allowbreak{}3\_\allowbreak{}odp\_\allowbreak{}05\\
Type: list\\
Default: [\textquotedbl{}admin\textquotedbl{}, \textquotedbl{}security\_\allowbreak{}officer\textquotedbl{}]
}

{\footnotesize Визначена періодичність оновлення оцінювання ризику; перегляду результатів RA-03a.01 проводиться оцінювання ризику для виявлення загроз і вразливостей у системі; RA-03a.02 проводиться оцінювання ризику для визначення ймовірності та розміру шкоди від несанкціонованого доступу, використання, розкриття, порушення, модифікації або знищення системи; інформації, яку вона обробляє, зберігає або передає; а також будь-якої пов'язаної з нею інформації; RA-03a.03 проводиться оцінювання ризику для визначення ймовірності та впливу несприятливих наслідків для фізичних осіб, що виникають у зв'язку з обробкою інформації, яка ідентифікує особу; RA-03b. інтегровані результати оцінювання ризику та рішення з управління ризиками з точки зору організації та місії або бізнес-процесів з оцінкою ризиків на системному рівні; RA-03c. результати оцінювання ризику задокументовані в ЗНАЧЕННЯ ВИБРАНОГО ПАРАМЕТРА; RA-03d. переглядаються результати 03\_\allowbreak{}ODP[03] частота>; RA-03e. поширюються результати оцінювання ризику серед персоналу або ролей; RA-03f. оновлюється оцінювання ризику частота або коли відбуваються значні зміни в системі, середовищі її функціонування або інших умовах, які можуть вплинути на стан безпеки або приватності системи оцінювання ризику <RA-}




\subsubsection{ОЦІНЮВАННЯ ПОСТАЧАННЯ (RA-3(1))}


\vspace{10pt}
{\footnotesize\bfseries
No: 1\\
Name: ra\_\allowbreak{}3\_\allowbreak{}1\_\allowbreak{}a\\
Type: string\\
Default: nil
}

{\footnotesize Оцінювання ризику ланцюга постачання, пов'язані з системами, системними компонентами та системними послугами}


{\footnotesize\bfseries
No: 2\\
Name: ra\_\allowbreak{}3\_\allowbreak{}1\_\allowbreak{}b\\
Type: string\\
Default: nil
}

{\footnotesize Потрібно оновлювати оцінювання ризику ланцюга постачання, коли відбуваються значні зміни у відповідному ланцюгу постачання, або коли зміни в системі, середовищі функціонування чи інших умовах можуть вимагати змін у ланцюгу постачання}




\subsubsection{ВИКОРИСТАННЯ ІНФОРМАЦІЇ З УСІХ ДОСТУПНИХ ДЖЕРЕЛ (RA-3(2))}


\vspace{10pt}
{\footnotesize\bfseries
No: 1\\
Name: ra\_\allowbreak{}3\_\allowbreak{}2\_\allowbreak{}01\\
Type: string\\
Default: nil
}

{\footnotesize Використовується інформація з усіх доступних джерел для аналізу ризиків}




\subsubsection{УСВІДОМЛЕННЯ ДИНАМІЧНИХ ЗАГРОЗ (RA-3(3))}


\vspace{10pt}
{\footnotesize\bfseries
No: 1\\
Name: ra\_\allowbreak{}3\_\allowbreak{}3\_\allowbreak{}01\\
Type: string\\
Default: \textquotedbl{}автоматизований засіб моніторингу\textquotedbl{}
}

{\footnotesize Визначається поточне середовище кіберзагроз на постійній основі за допомогою засоби}


{\footnotesize\bfseries
No: 2\\
Name: ra\_\allowbreak{}3\_\allowbreak{}3\_\allowbreak{}odp\\
Type: string\\
Default: \textquotedbl{}автоматизований засіб моніторингу\textquotedbl{}
}

{\footnotesize Є засоби для постійного визначення поточного стану кіберзагроз}




\subsubsection{ПРОГНОСТИЧНА КІБЕРАНАЛІТИКА (RA-3(4))}


\vspace{10pt}
{\footnotesize\bfseries
No: 1\\
Name: ra\_\allowbreak{}3\_\allowbreak{}4\_\allowbreak{}01\\
Type: string\\
Default: nil
}

{\footnotesize Використовуються розширені можливості автоматизації для прогнозування та виявлення ризику для систем або компонентів системи}


{\footnotesize\bfseries
No: 2\\
Name: ra\_\allowbreak{}3\_\allowbreak{}4\_\allowbreak{}02\\
Type: string\\
Default: nil
}

{\footnotesize Застосовуються розширені аналітичні можливості для прогнозування та виявлення ризику для систем або компонентів системи. аналітики}




\subsection{ОНОВЛЕННЯ ОЦІНЮВАННЯ РИЗИКУ (RA-4) [Вилучено]}
[Вилучено: включено до RA-3]

\vspace{10pt}
\textit{Немає параметрів для цього контролю.}
\vspace{10pt}


\subsection{СКАНУВАННЯ ВРАЗЛИВОСТЕЙ (RA-5)}
Використовувати заходи ПДТР за [Призначення: визначені організацією місця] [Вибір (один або кілька): [Призначення: з визначеною організацією частотою]; [Призначення: за визначеними організацією подіями або показниками]].

\vspace{10pt}
{\footnotesize\bfseries
No: 1\\
Name: ra\_\allowbreak{}5\_\allowbreak{}odp\_\allowbreak{}01\\
Type: string\\
Default: nil
}

{\footnotesize Визначена необхідність моніторингу систем та розміщених застосунків на наявність вразливостей}


{\footnotesize\bfseries
No: 2\\
Name: ra\_\allowbreak{}5\_\allowbreak{}odp\_\allowbreak{}02\\
Type: string\\
Default: \textquotedbl{}щорічно\textquotedbl{}
}

{\footnotesize Визначена періодичність перевірки систем та розміщених на них застосунків на наявність вразливостей}


{\footnotesize\bfseries
No: 3\\
Name: ra\_\allowbreak{}5\_\allowbreak{}odp\_\allowbreak{}03\\
Type: integer\\
Default: 30
}

{\footnotesize Визначено час реагування на усунення законних вразливостей відповідно до організаційної оцінення ризику}


{\footnotesize\bfseries
No: 4\\
Name: ra\_\allowbreak{}5\_\allowbreak{}odp\_\allowbreak{}04\\
Type: list\\
Default: [\textquotedbl{}admin\textquotedbl{}, \textquotedbl{}security\_\allowbreak{}officer\textquotedbl{}]
}

{\footnotesize Потрібно ділитися інформацією, отриманою в процесі сканування вразливостей та оцінок контролю, з персоналом або ролями, з якими потрібно ділитися; RA-05a.[01] здійснюється моніторинг систем та розміщених застосунків на наявність вразливостей частота та/або випадковість відповідно до визначеного організацією процесу, а також коли виявляються та повідомляються нові вразливості, що потенційно можуть вплинути на систему; RA-05a.[02] перевіряються системи та розміщені застосунки на наявність вразливостей частота та/або випадковим чином відповідно до визначеного організацією процесу, а також коли виявляються та повідомляються нові вразливості, що потенційно можуть вплинути на систему; RA-05b. застосовуються інструменти та методи моніторингу вразливостей для забезпечення сумісності між інструментами; RA-05b.01 застосовуються інструменти та методи моніторингу вразливостей для автоматизації частини процесу управління вразливостями, використовуючи стандарти для переліку платформ, недоліків програмного забезпечення та неправильних конфігурацій; RA-05b.02 застосовуються інструменти та методи моніторингу вразливостей для полегшення взаємодії між інструментами та автоматизації частини процесу управління вразливостями шляхом використання стандартів для формування контрольних списків та процедур тестування; RA-05b.03 застосовуються інструменти та методи моніторингу вразливостей для полегшення взаємодії між інструментами та автоматизації частин процесу управління вразливостями шляхом використання стандартів для вимірювання впливу вразливостей; RA-05c. аналізуються звіти про сканування вразливостей та результати моніторингу вразливостей; RA-05d. усуваються легітимні вразливості час реагування відповідно до організаційної оцінки ризиків; RA-05e. надається інформація, отримана в процесі моніторингу вразливостей та оцінки контролю, персонал або ролі, щоб допомогти усунути подібні вразливості в інших системах; RA-05f. використовуються інструменти моніторингу вразливостей, які передбачають можливість швидкого оновлення вразливостей, що підлягають скануванню}




\subsubsection{СКАНУВАННЯ ВРАЗЛИВОСТЕЙ ІНСТРУМЕНТІВ (RA-5(1)) [Вилучено]}
[Вилучено: включено до RA-5]

\vspace{10pt}
\textit{Немає параметрів для цього контролю.}
\vspace{10pt}


\subsubsection{ОНОВЛЕННЯ ЗА ЧАСТОТОЮ, ПЕРЕД НОВИМ СКАНУВАННЯМ АБО ПРИ ІДЕНТИФІКАЦІЇ (RA-5(2))}


\vspace{10pt}
{\footnotesize\bfseries
No: 1\\
Name: ra\_\allowbreak{}5\_\allowbreak{}2\_\allowbreak{}01\\
Type: integer\\
Default: 30
}

{\footnotesize Визначено час реагування на усунення законних вразливостей відповідно до організаційної оцінення ризику}




\subsubsection{ШИРОТА ТА ГЛИБИНА ПОКРИТТЯ (RA-5(3))}


\vspace{10pt}
{\footnotesize\bfseries
No: 1\\
Name: ra\_\allowbreak{}5\_\allowbreak{}3\_\allowbreak{}01\\
Type: string\\
Default: nil
}

{\footnotesize Визначено ширину вразливостей. та глибину охоплення сканування}




\subsubsection{ВИЯВНА ІНФОРМАЦІЯ (RA-5(4))}


\vspace{10pt}
{\footnotesize\bfseries
No: 1\\
Name: ra\_\allowbreak{}5\_\allowbreak{}4\_\allowbreak{}01\\
Type: list\\
Default: [\textquotedbl{}login\textquotedbl{}, \textquotedbl{}logout\textquotedbl{}, \textquotedbl{}failed\_\allowbreak{}attempt\textquotedbl{}]
}

{\footnotesize Є інформація про систему відкритою; RA-05(04)[02] вживаються коригувальні дії, коли інформація про систему підтверджується як така, що може бути виявлена}


{\footnotesize\bfseries
No: 2\\
Name: ra\_\allowbreak{}5\_\allowbreak{}4\_\allowbreak{}odp\\
Type: list\\
Default: [\textquotedbl{}login\textquotedbl{}, \textquotedbl{}logout\textquotedbl{}, \textquotedbl{}failed\_\allowbreak{}attempt\textquotedbl{}]
}

{\footnotesize Визначені коригувальні дії, які необхідно вжити, якщо інформація про систему буде виявлена}




\subsubsection{ПРИВІЛЕЙОВАНИЙ ДОСТУП (RA-5(5))}


\vspace{10pt}
{\footnotesize\bfseries
No: 1\\
Name: ra\_\allowbreak{}5\_\allowbreak{}5\_\allowbreak{}01\\
Type: string\\
Default: nil
}

{\footnotesize Реалізовано авторизацію привілейованого доступу до компоненти системи для діяльность зі сканування вразливостей}




\subsubsection{СКАНУВАННЯ ТЕНДЕНЦІЙ (RA-5(6))}


\vspace{10pt}
{\footnotesize\bfseries
No: 1\\
Name: ra\_\allowbreak{}5\_\allowbreak{}6\_\allowbreak{}01\\
Type: string\\
Default: \textquotedbl{}автоматизований засіб моніторингу\textquotedbl{}
}

{\footnotesize Порівнюються результати багаторазового сканування вразливостей за допомогою автоматизовані механізми}


{\footnotesize\bfseries
No: 2\\
Name: ra\_\allowbreak{}5\_\allowbreak{}6\_\allowbreak{}odp\\
Type: string\\
Default: \textquotedbl{}автоматизований засіб моніторингу\textquotedbl{}
}

{\footnotesize Визначені автоматизовані механізми для порівняння результатів багаторазового сканування вразливостей}




\subsubsection{АВТОМАТИЗОВАНЕ ВИЯВЛЕННЯ ТА СПОВІЩЕННЯ ПРО НЕАВТОРИЗОВАНІ КОМПОНЕНТИ (RA-5(7)) [Вилучено]}
[Вилучено: включено до CM-8]

\vspace{10pt}
\textit{Немає параметрів для цього контролю.}
\vspace{10pt}


\subsubsection{ОГЛЯД ЖУРНАЛІВ АУДИТУ ЗА МИНУЛІ ПЕРІОДИ (RA-5(8))}


\vspace{10pt}
\textit{Немає параметрів для цього контролю.}
\vspace{10pt}


\subsubsection{СКАНУВАННЯ ВРАЗЛИВОСТЕЙ ПРОНИКНЕННЯ (RA-5(9)) [Вилучено]}
[Вилучено: включено до CA-8]

\vspace{10pt}
\textit{Немає параметрів для цього контролю.}
\vspace{10pt}


\subsection{ЗАХОДИ ПРОТИДІЇ ТЕХНІЧНІЙ РОЗВІДЦІ (RA-6)}


\vspace{10pt}
{\footnotesize\bfseries
No: 1\\
Name: ra\_\allowbreak{}6\_\allowbreak{}odp\_\allowbreak{}01\\
Type: string\\
Default: nil
}

{\footnotesize Визначені місця для використання заходів ПДТР}


{\footnotesize\bfseries
No: 2\\
Name: ra\_\allowbreak{}6\_\allowbreak{}odp\_\allowbreak{}02\\
Type: list\\
Default: [\textquotedbl{}login\textquotedbl{}, \textquotedbl{}logout\textquotedbl{}, \textquotedbl{}failed\_\allowbreak{}attempt\textquotedbl{}]
}

{\footnotesize Вибрано одне або декілька з наступних ЗНАЧЕНЬ ПАРАМЕТРІВ: \{частота; коли події або показники\}}


{\footnotesize\bfseries
No: 3\\
Name: ra\_\allowbreak{}6\_\allowbreak{}odp\_\allowbreak{}03\\
Type: integer\\
Default: 30
}

{\footnotesize Визначено частоту, з якою слід проводити заходи ПДТР (якщо обрано)}


{\footnotesize\bfseries
No: 4\\
Name: ra\_\allowbreak{}6\_\allowbreak{}odp\_\allowbreak{}04\\
Type: list\\
Default: [\textquotedbl{}login\textquotedbl{}, \textquotedbl{}logout\textquotedbl{}, \textquotedbl{}failed\_\allowbreak{}attempt\textquotedbl{}]
}

{\footnotesize Визначені події або показники, які, у разі їх виникнення, спричиняють проведення заходів ПДТР (якщо вони були обрані)}




\subsection{РЕАГУВАННЯ НА РИЗИК (RA-7)}
Реагувати на результати оцінювання, моніторингу й аудиту безпеки та приватності.

\vspace{10pt}
{\footnotesize\bfseries
No: 1\\
Name: ra\_\allowbreak{}7\_\allowbreak{}01\\
Type: string\\
Default: nil
}

{\footnotesize Вживаються заходи реагування на результати оцінок безпеки відповідно до організаційної толерантності до ризиків}


{\footnotesize\bfseries
No: 2\\
Name: ra\_\allowbreak{}7\_\allowbreak{}02\\
Type: string\\
Default: nil
}

{\footnotesize Вживаються заходи реагування на результати оцінювання приватності відповідно до організаційної толерантності до ризиків}


{\footnotesize\bfseries
No: 3\\
Name: ra\_\allowbreak{}7\_\allowbreak{}03\\
Type: string\\
Default: nil
}

{\footnotesize Вживаються заходи реагування на результати моніторингу відповідно до організаційної толерантності до ризиків}


{\footnotesize\bfseries
No: 4\\
Name: ra\_\allowbreak{}7\_\allowbreak{}04\\
Type: string\\
Default: nil
}

{\footnotesize Вживаються заходи реагування на висновки аудиту відповідно до організаційної толерантності до ризиків}




\subsection{ОЦІНКА ВПЛИВУ НА ПРИВАТНІСТЬ (RA-8)}
Проводити оцінювання впливу на приватність інформаційних систем, програм або інших заходів, які становлять ризик приватності перед тим, як:\\ 
a. розробити або закупити інформаційні технології, які збирають, підтримують чи поширюють критичну інформацію;\\ 
b. ініціювати створення нових архівів інформації, яка:\\ 
1. буде зібрана, збережена чи розповсюджена за допомогою інформаційних технологій;\\ 
2. містить персональні дані, які дозволяють встановити фізичне або онлайнз'єднання з конкретною особою.

\vspace{10pt}
\textit{Немає параметрів для цього контролю.}
\vspace{10pt}


\subsection{АНАЛІЗ КРИТИЧНОСТІ (RA-9)}
Визначити критичні компоненти інформаційної системи та функції, виконавши аналіз критичності для [Призначення: визначених організацією систем, компонентів системи або послуг для системи] в [Призначення: визначенні організацією точки ухвалення рішень у життєвому циклі розробки системи].

\vspace{10pt}
{\footnotesize\bfseries
No: 1\\
Name: ra\_\allowbreak{}9\_\allowbreak{}01\\
Type: string\\
Default: nil
}

{\footnotesize Визначені критичні компоненти та функції системи шляхом проведення аналізу критичності для систем, системних компонентів або системних служб в точках прийняття рішень в життєвому циклі розробки системи}


{\footnotesize\bfseries
No: 2\\
Name: ra\_\allowbreak{}9\_\allowbreak{}odp\_\allowbreak{}01\\
Type: string\\
Default: nil
}

{\footnotesize Визначені системи, компоненти системи або системні сервіси, що підлягають аналізу на предмет критичності}


{\footnotesize\bfseries
No: 3\\
Name: ra\_\allowbreak{}9\_\allowbreak{}odp\_\allowbreak{}02\\
Type: string\\
Default: nil
}

{\footnotesize Визначені точки прийняття рішень в життєвому циклі розробки системи, коли необхідно проводити аналіз критичності}




\subsection{АКТИВНИЙ ПОШУК ЗАГРОЗ (RA-10)}
a. Створити та підтримувати можливості активного пошуку кіберзагроз:\\ 
1. пошук індикаторів компроментації в системах організації;\\ 
2. виявлення, відстеження та знищення загроз, які можуть обходити існуючі засоби контролю безпеки.\\ 
b. використовуйте можливості активного пошуку загроз [Призначення: частота, визначена організацією].

\vspace{10pt}
{\footnotesize\bfseries
No: 1\\
Name: ra\_\allowbreak{}10\_\allowbreak{}odp\\
Type: list\\
Default: [\textquotedbl{}login\textquotedbl{}, \textquotedbl{}logout\textquotedbl{}, \textquotedbl{}failed\_\allowbreak{}attempt\textquotedbl{}]
}

{\footnotesize Визначена частота, з якою можливість виявлення загроз; слід використовувати RA-10a.01 створена та підтримується спроможність протидії кіберзагрозам для пошуку індикаторів компрометації в організаційних системах; RA-10a.02 створена та підтримується спроможність виявляти, відслідковувати та знешкоджувати кіберзагрози, які не піддаються існуючому контролю; RA-10b. використовується функція відстеження загроз частота}




\section{SA}
Клас заходів захисту SA --- ПРИДБАННЯ СИСТЕМИ ТА ПОСЛУГ

Цей клас зосереджується на інтеграції вимог безпеки на всіх етапах життєвого циклу розробки та закупівлі ІТ-продуктів чи послуг.

\textbf{Перелік заходів захисту:} Політики та процедури придбання систем та послуг (SA-1); Розподіл ресурсів (SA-2); Життєвий цикл розробки системи (SA-3); Середовищем розробки (SA-3(1)); Життєвий цикл реальних даних (SA-3(2)); Життєвий цикл технологій (SA-3(3)); Процес закупівель (SA-4); Функціональні властивості заходів (SA-4(1)); Розробка та впровадження інформації для заходів (SA-4(2)); Процес закупівель розробки (SA-4(3)); Віднесення компонентів до систем (SA-4(4)) [Вилучено]; Конфігурації системи, компонента та системної служби (SA-4(5)); Процес закупівель інформації (SA-4(6)); Затверджені профілі захищеності (SA-4(7)); План безперервного моніторингу заходів безпеки (SA-4(8)); Функції, послуги, що використовуються (SA-4(9)); Системна документація (SA-5); Функціональні властивості заходів безпеки (SA-5(1)) [Вилучено]; Зовнішні системні інтерфейси, (SA-5(2)) [Вилучено]; АРХІТЕКТУРА (ПРОЄКТ) ВИСОКОГО РІВНЯ (SA-5(3)) [Вилучено]; АРХІТЕКТУРА (ПРОЄКТ) НИЗЬКОГО РІВНЯ (SA-5(4)) [Вилучено]; Вихідний код (SA-5(5)) [Вилучено]; Обмеження щодо використання забезпечення (SA-6) [Вилучено]; Встановлене користувачем програмне забезпечення (SA-7) [Вилучено]; Безпека та приватність принципів інжинірингу (SA-8); Чітка абстракція (SA-8(1)); Безпека та приватність принципів найменш поширений механізм (SA-8(2)); Безпека та приватність принципів модульність і багаторівневість (SA-8(3)); Безпека та приватність принципів частково впорядковані залежності (SA-8(4)); Безпека та приватність принципів ефективний опосередкований доступ (SA-8(5)); Безпека та приватність мінімізований обмін (SA-8(6)); Знижена складність (SA-8(7)); Безпека та приватність принципів еволюція безпеки в системі (SA-8(8)); Довірені компоненти системи (SA-8(9)); Реалізують принцип ієрархічної довіри в дизайні безпеки (SA-8(10)); Реалізують принцип проектування безпеки зворотного порогу модифікації (SA-8(11)); Реалізують принцип побудови ієрархічного захисту (SA-8(12)); Реалізують принцип проектування безпеки з мінімізацією елементів безпеки (SA-8(13)); Реалізують принцип побудови безпеки за принципом найменших привілеїв (SA-8(14)); Реалізують принцип проектування безпеки попереднього дозволу (SA-8(15)); Реалізують принцип проектування безпеки самодостатньої надійності (SA-8(16)); Реалізують принцип проектування безпеки захищеного розподіленого вмісту (SA-8(17)); Реалізують принцип побудови безпеки довірених каналів зв'язку (SA-8(18)); Реалізують принцип проектування безпеки довготривалого (SA-8(19)); Реалізують принцип безпечного управління метаданими (SA-8(20)); Реалізують принцип проектування безпеки на основі самоаналізу (SA-8(21)); Реалізують принцип проектування безпечних налаштувань за замовчуванням (SA-8(23)); <sa-08(25)\_\allowbreak{}odp системи або компоненти реалізують принцип безпеки економіки. (SA-8(25)); <sa-08(26)\_\allowbreak{}odp системи або компоненти реалізують принцип безпеки продуктивності. (SA-8(26)); Реалізують принцип безпеки з урахуванням людського фактору (SA-8(27)); <sa-08(28)\_\allowbreak{}odp системи або компоненти реалізують принцип прийнятного рівня безпеки. (SA-8(28)); Реалізують принцип безпечної модифікації систем (SA-8(31)); Реалізують принцип достатньої допускної документації (SA-8(32)); <sa-08(33)\_\allowbreak{}odp системи або компоненти реалізують принцип мінімізації конфіденційності. (SA-8(33)); Зовнішні послуги для системи (SA-9); Зовнішні послуги для системи- оцінювання ризиків та організаційні погодження (SA-9(1)); Зовнішні послуги для системи- визначення функцій, портів, протоколів та служб (SA-9(2)); Створення та підтримка довірчих відносин з постачальниками (SA-9(3)); Зовнішні системні служби споживачів і постачальників (SA-9(4)); Зовнішні послуги для системи- місце обробки, зберігання та обслуговування (SA-9(5)); Зовнішні послуги для системи- криптографічні ключі, керовані організацією (SA-9(6)); Зовнішні послуги для системи- перевірка цілісності, що контролюється організацією (SA-9(7)); Зовнішні послуги для системизберігання -- юрисдикція україни (SA-9(8)); Управління конфігурацією розробника (SA-10); Управління конфігурацією розробника --- перевірка цілісності програмного забезпечення та мікропрограм (SA-10(1)); Управління конфігурацією розробника --- довірене постачання (SA-10(6)); Представники безпеки та приватності (SA-11); Тестування та оцінювання розробника --- моделювання загроз і аналіз вразливостей (SA-11(2)); Ate\_\allowbreak{}ind.2 eal2 незалежне тестування eal3 незалежне тестування --- eal4 eal5 eal6 зразок (SA-11(3)); Тестування та оцінювання розробника --- тестування на проникнення (SA-11(5)); Покриття аналіз покриття (SA-11(7)); Динамічний аналіз коду (SA-11(8)); Керування ризиками ланцюга постачання (SA-12) [Вилучено]; Довірчість (SA-13) [Вилучено]; Аналіз критичності (SA-14) [Вилучено]; Процеси, стандарти та інструменти розробки (SA-15); Зменшення поверхні атаки (SA-15(5)); Постійне вдосконалення (SA-15(6)); Аналіз вразливостей автоматизований (SA-15(7)); Повторне використання інформації про загрози та вразливості (SA-15(8)); Процес, стандарти та інструменти використання реальних даних (SA-15(9)) [Вилучено]; Навчання, що надається розробниками (SA-16); Проєкт та архітектура розробника (SA-17); Проєкт і архітектура безпеки розробника --- формальна модель політики (SA-17(1)); Проєкт і архітектура безпеки розробника --- компоненти, що необхідні для забезпечення безпеки (SA-17(2)); Проєкт і архітектура безпеки розробника --- формальна відповідність (SA-17(3)); Проєкт і архітектура безпеки розробника --- неформальна відповідність (SA-17(4)); Захист та виявлення підробки (SA-18) [Вилучено]; Справжність компонента (SA-19) [Вилучено]; Індивідуальна розробка критичних компонентів (SA-20); Перевірка розробника (SA-21); Компоненти системи, що не підтримуються (SA-22); Спеціалізація (SA-23).

\subsection{ПОЛІТИКИ ТА ПРОЦЕДУРИ ПРИДБАННЯ СИСТЕМ ТА ПОСЛУГ (SA-1)}
a. Розробити, задокументувати та поширити серед [Призначення: визначеного організацією персоналу або ролей]: 1. 2.\\ 
b. [Вибір (один або декілька): рівень організації; рівень місії/бізнес-процесу; рівень системи] політики придбання систем і послуг, яка:\\ 
(a) містить мету, сферу застосування, ролі, обов'язки, відповідальність керівництва, координацію між організаційними підрозділами та систему контролю (complaince);\\ 
(b) відповідає чинному законодавству, виконавчим наказам, директивам, нормам, політикам, стандартам і рекомендаціям.\\ 
Процедури, що полегшують впровадження політики та заходів придбання систем і послуг.\\ 
Призначити [Призначення: визначену організацією посадову особу] для управління політикою та процедурами придбання системи та послуг.\\ 
c. Переглядати й оновлювати поточні політику та процедури придбання систем та послуг:\\ 
1. Поточну політику придбання системи та послуг [Призначення: з визначеною організацією частотою].\\ 
2. Поточні процедури придбання системи та послуг [Призначення: з визначеною організацією частотою] та наступні [Призначення: події, визначені організацією].\\ 


\vspace{10pt}
{\footnotesize\bfseries
No: 1\\
Name: sa\_\allowbreak{}1\_\allowbreak{}odp\_\allowbreak{}01\\
Type: list\\
Default: [\textquotedbl{}admin\textquotedbl{}, \textquotedbl{}security\_\allowbreak{}officer\textquotedbl{}]
}

{\footnotesize Визначено персонал або ролі, на які поширюватиметься політика придбання систем і послуг}


{\footnotesize\bfseries
No: 2\\
Name: sa\_\allowbreak{}1\_\allowbreak{}odp\_\allowbreak{}02\\
Type: list\\
Default: [\textquotedbl{}admin\textquotedbl{}, \textquotedbl{}security\_\allowbreak{}officer\textquotedbl{}]
}

{\footnotesize Визначено персонал або ролі, на які поширюються процедури придбання систем і послуг}


{\footnotesize\bfseries
No: 3\\
Name: sa\_\allowbreak{}1\_\allowbreak{}odp\_\allowbreak{}03\\
Type: string\\
Default: nil
}

{\footnotesize Вибрано одне або декілька з наступних ЗНАЧЕНЬ ПАРАМЕТРІВ: \{рівень організації; рівень місії/бізнеспроцесу; рівень системи\}}


{\footnotesize\bfseries
No: 4\\
Name: sa\_\allowbreak{}1\_\allowbreak{}odp\_\allowbreak{}04\\
Type: list\\
Default: [\textquotedbl{}admin\textquotedbl{}, \textquotedbl{}security\_\allowbreak{}officer\textquotedbl{}]
}

{\footnotesize Визначено посадову особу, яка керуватиме політикою та процедурами придбання систем і послуг}


{\footnotesize\bfseries
No: 5\\
Name: sa\_\allowbreak{}1\_\allowbreak{}odp\_\allowbreak{}05\\
Type: list\\
Default: [\textquotedbl{}default\_\allowbreak{}deny\_\allowbreak{}rule\textquotedbl{}, \textquotedbl{}abac\_\allowbreak{}rule\_\allowbreak{}1\textquotedbl{}]
}

{\footnotesize Визначено періодичність перегляду та оновлення поточної політики придбання систем і послуг}


{\footnotesize\bfseries
No: 6\\
Name: sa\_\allowbreak{}1\_\allowbreak{}odp\_\allowbreak{}06\\
Type: list\\
Default: [\textquotedbl{}default\_\allowbreak{}deny\_\allowbreak{}rule\textquotedbl{}, \textquotedbl{}abac\_\allowbreak{}rule\_\allowbreak{}1\textquotedbl{}]
}

{\footnotesize Є події, які вимагають перегляду та оновлення поточної політики придбання систем і послуг}


{\footnotesize\bfseries
No: 7\\
Name: sa\_\allowbreak{}1\_\allowbreak{}odp\_\allowbreak{}07\\
Type: integer\\
Default: 30
}

{\footnotesize Визначено частоту, з якою переглядаються та оновлюються поточні процедури придбання систем і послуг}


{\footnotesize\bfseries
No: 8\\
Name: sa\_\allowbreak{}1\_\allowbreak{}odp\_\allowbreak{}08\\
Type: integer\\
Default: 30
}

{\footnotesize Є події, які вимагають перегляду та оновлення процедур придбання систем і послуг}


{\footnotesize\bfseries
No: 9\\
Name: sa\_\allowbreak{}1\_\allowbreak{}a\_\allowbreak{}01\\
Type: integer\\
Default: 30
}

{\footnotesize Розроблена та задокументована політика придбання систем і послуг}


{\footnotesize\bfseries
No: 10\\
Name: sa\_\allowbreak{}1\_\allowbreak{}a\_\allowbreak{}02\\
Type: integer\\
Default: 30
}

{\footnotesize Поширюється політика придбання систем і послуг на SA-1-ODP-01 персонал або ролі}


{\footnotesize\bfseries
No: 11\\
Name: sa\_\allowbreak{}1\_\allowbreak{}a\_\allowbreak{}03\\
Type: integer\\
Default: 30
}

{\footnotesize Розроблені та задокументовані процедури придбання систем і послуг, що сприяють впровадженню політики придбання систем та послуг, а також відповідні засоби контролю за придбанням систем та послуг}


{\footnotesize\bfseries
No: 12\\
Name: sa\_\allowbreak{}1\_\allowbreak{}a\_\allowbreak{}03\\
Type: integer\\
Default: 30
}

{\footnotesize Розроблені та задокументовані процедури придбання систем і послуг, що сприяють впровадженню політики придбання систем та послуг, а також відповідні засоби контролю за придбанням систем та послуг}


{\footnotesize\bfseries
No: 13\\
Name: sa\_\allowbreak{}1\_\allowbreak{}a\_\allowbreak{}04\\
Type: integer\\
Default: 30
}

{\footnotesize Поширюються процедури придбання системи і послуг на SA-1-ODP-02 персонал або ролі}


{\footnotesize\bfseries
No: 14\\
Name: sa\_\allowbreak{}1\_\allowbreak{}a\_\allowbreak{}04\_\allowbreak{}1\\
Type: integer\\
Default: 30
}

{\footnotesize Відповідає політика придбання системи і послуг <SA01\_\allowbreak{}ODP[03] ВИБІРКОВЕ ЗНАЧЕННЯ ПАРАМЕТРА(ів)> поставленій меті}


{\footnotesize\bfseries
No: 15\\
Name: sa\_\allowbreak{}1\_\allowbreak{}a\_\allowbreak{}04\_\allowbreak{}2\\
Type: integer\\
Default: 30
}

{\footnotesize Політика придбання систем і послуг <SA-01\_\allowbreak{}ODP[03] ВИБІРКОВЕ ЗНАЧЕННЯ ПАРАМЕТРА(S)> враховує сферу застосування}


{\footnotesize\bfseries
No: 16\\
Name: sa\_\allowbreak{}1\_\allowbreak{}a\_\allowbreak{}04\_\allowbreak{}3\\
Type: integer\\
Default: 30
}

{\footnotesize Політика придбання системи і послуг <SA-01\_\allowbreak{}ODP[03] ВИБІРКОВЕ ЗНАЧЕННЯ ПАРАМЕТРА(S)> стосується ролей}


{\footnotesize\bfseries
No: 17\\
Name: sa\_\allowbreak{}1\_\allowbreak{}a\_\allowbreak{}04\_\allowbreak{}4\\
Type: integer\\
Default: 30
}

{\footnotesize Політика придбання систем і послуг <SA-01\_\allowbreak{}ODP[03] ВИБІРКОВЕ ЗНАЧЕННЯ ПАРАМЕТРА(ів)> враховує відповідальність}


{\footnotesize\bfseries
No: 18\\
Name: sa\_\allowbreak{}1\_\allowbreak{}a\_\allowbreak{}04\_\allowbreak{}5\\
Type: integer\\
Default: 30
}

{\footnotesize Враховує політика придбання систем і послуг <SA01\_\allowbreak{}ODP[03] ВИБІРКОВЕ ЗНАЧЕННЯ ПАРАМЕТРА(ів)> зобов'язання керівництва}


{\footnotesize\bfseries
No: 19\\
Name: sa\_\allowbreak{}1\_\allowbreak{}a\_\allowbreak{}04\_\allowbreak{}6\\
Type: integer\\
Default: 30
}

{\footnotesize Політика придбання систем і послуг <SA-01\_\allowbreak{}ODP[03] ВИБІРКОВЕ ЗНАЧЕННЯ ПАРАМЕТРА(ів)> передбачає координацію між організаціями}


{\footnotesize\bfseries
No: 20\\
Name: sa\_\allowbreak{}1\_\allowbreak{}a\_\allowbreak{}04\_\allowbreak{}7\\
Type: integer\\
Default: 30
}

{\footnotesize Політика придбання систем і послуг <SA-01\_\allowbreak{}ODP[03] ВИБІРКОВЕ ЗНАЧЕННЯ ПАРАМЕТРА(ів)> враховує відповідність вимогам}


{\footnotesize\bfseries
No: 21\\
Name: sa\_\allowbreak{}1\_\allowbreak{}a\_\allowbreak{}04\_\allowbreak{}8\\
Type: integer\\
Default: 30
}

{\footnotesize Відповідає <SA-01\_\allowbreak{}ODP[03] ВИБІРКОВЕ ЗНАЧЕННЯ ПАРАМЕТРА(ів)> політика придбання системи і послуг чинному законодавству, виконавчим наказам, директивам, положенням, політикам, стандартам та настановам}


{\footnotesize\bfseries
No: 22\\
Name: sa\_\allowbreak{}1\_\allowbreak{}b\\
Type: integer\\
Default: 30
}

{\footnotesize Призначена <SA-01\_\allowbreak{}ODP[04] посадова особа> для управління розробкою, документуванням та розповсюдженням політики та процедур придбання системи і послуг}


{\footnotesize\bfseries
No: 23\\
Name: sa\_\allowbreak{}1\_\allowbreak{}c\_\allowbreak{}01\_\allowbreak{}1\\
Type: integer\\
Default: 30
}

{\footnotesize Переглядаються та оновлюються поточна політика придбання систем і послуг <SA-01\_\allowbreak{}ODP[07] частота>}


{\footnotesize\bfseries
No: 24\\
Name: sa\_\allowbreak{}1\_\allowbreak{}c\_\allowbreak{}01\_\allowbreak{}2\\
Type: integer\\
Default: 30
}

{\footnotesize Переглядається та оновлюється поточна політика придбання систем та послуг після <SA-01\_\allowbreak{}ODP[06] подій>}


{\footnotesize\bfseries
No: 25\\
Name: sa\_\allowbreak{}1\_\allowbreak{}c\_\allowbreak{}02\_\allowbreak{}1\\
Type: integer\\
Default: 30
}

{\footnotesize Переглядаються та оновлюються поточні процедури придбання систем і послуг <SA-01\_\allowbreak{}ODP[07] частота>}


{\footnotesize\bfseries
No: 26\\
Name: sa\_\allowbreak{}1\_\allowbreak{}c\_\allowbreak{}02\_\allowbreak{}2\\
Type: integer\\
Default: 30
}

{\footnotesize Переглядаються та оновлюються поточні процедури закупівлі систем та послуг після <SA-01\_\allowbreak{}ODP[08] подій>}




\subsection{РОЗПОДІЛ РЕСУРСІВ (SA-2)}
a. Визначити вимоги щодо інформаційної безпеки та приватності для систем або послуг для системи при плануванні завдань та процесів.\\ 
b. Визначити, задокументувати та розподілити ресурси, які необхідні для захисту систем або послуг для системи у рамках процесу фінансового планування в організації та управління інвестиціями.\\ 
c. Створити окрему позицію бюджету для фінансування заходів із забезпечення інформаційної безпеки та приватності.\\ 


\vspace{10pt}
{\footnotesize\bfseries
No: 1\\
Name: sa\_\allowbreak{}2\_\allowbreak{}a\_\allowbreak{}01\\
Type: integer\\
Default: 30
}

{\footnotesize Визначені вимоги до інформаційної безпеки високого рівня для систем або послуг для системи при плануванні місії та бізнес-процесів}


{\footnotesize\bfseries
No: 2\\
Name: sa\_\allowbreak{}2\_\allowbreak{}a\_\allowbreak{}02\\
Type: integer\\
Default: 30
}

{\footnotesize Визначені вимоги до приватності високого рівня для систем або послуг для системи при плануванні місії та бізнес-процесів}


{\footnotesize\bfseries
No: 3\\
Name: sa\_\allowbreak{}2\_\allowbreak{}b\_\allowbreak{}01\\
Type: integer\\
Default: 30
}

{\footnotesize Визначені та задокументовані ресурси, необхідні для захисту систем або послуг для системи, як частина процесу планування організаційного капіталу та контролю за інвестиціями}


{\footnotesize\bfseries
No: 4\\
Name: sa\_\allowbreak{}2\_\allowbreak{}b\_\allowbreak{}02\\
Type: integer\\
Default: 30
}

{\footnotesize Виділені ресурси, необхідні для захисту систем або послуг для системи, в рамках процесу планування організаційного капіталу та контролю інвестицій}


{\footnotesize\bfseries
No: 5\\
Name: sa\_\allowbreak{}2\_\allowbreak{}c\_\allowbreak{}01\\
Type: integer\\
Default: 30
}

{\footnotesize Передбачено окрему статтю витрат на інформаційну безпеку в програмній та бюджетній документації організації}


{\footnotesize\bfseries
No: 6\\
Name: sa\_\allowbreak{}2\_\allowbreak{}c\_\allowbreak{}02\\
Type: integer\\
Default: 30
}

{\footnotesize Передбачена окрема стаття для захисту приватності в програмній та бюджетній документації організації}




\subsection{ЖИТТЄВИЙ ЦИКЛ РОЗРОБКИ СИСТЕМИ (SA-3)}
a. Придбати, розробити та керувати системою, використовуючи [Призначення: визначений організацією життєвий цикл розробки], який охоплює питання захисту інформації та приватності.\\ 
b. Визначити та задокументувати роль і обов'язки із забезпечення безпеки та приватності інформації протягом усього життєвого циклу розробки системи.\\ 
c. Визначити осіб, які мають повноваження та обов'язки в області інформаційної безпеки та приватності.\\ 
d. Інтегрувати процес управління інформаційною безпекою та приватністю в процеси життєвого циклу розробки системи\\ 


\vspace{10pt}
{\footnotesize\bfseries
No: 1\\
Name: sa\_\allowbreak{}3\_\allowbreak{}odp\\
Type: list\\
Default: [\textquotedbl{}admin\textquotedbl{}, \textquotedbl{}security\_\allowbreak{}officer\textquotedbl{}]
}

{\footnotesize Визначено життєвий цикл розробки системи; SA-03a.[01] система придбана, розроблена та керується з використанням життєвого циклу життєвий цикл розробки системи , який охоплює інформаційну безпеку; SA-03a.[02] система придбана, розроблена та керується з використанням життєвий цикл розробки системи , який охоплює приватність; SA-03b.[01] визначені та задокументовані ролі та обов'язки з інформаційної безпеки протягом усього життєвого циклу розробки системи; SA-03b.[02] визначені та задокументовані ролі та обов'язки щодо приватності протягом усього життєвого циклу розробки системи; SA-03c.[01] визначені особи, які виконують функції та обов'язки з інформаційної безпеки; SA-03c.[02] визначені особи з функціями та обов'язками, пов'язаними з приватністю; SA-03d.[01] інтегровані процеси управління інформаційною безпекою організації в діяльність життєвого циклу розробки системи; SA-03d.[02] інтегровані процеси управління приватністю в життєвого циклу розробки системи}


{\footnotesize\bfseries
No: 2\\
Name: sa\_\allowbreak{}3\_\allowbreak{}a\_\allowbreak{}1\\
Type: list\\
Default: [\textquotedbl{}admin\textquotedbl{}, \textquotedbl{}security\_\allowbreak{}officer\textquotedbl{}]
}

{\footnotesize Система придбана, розроблена та керується з використанням життєвого циклу життєвий цикл розробки системи , який охоплює інформаційну безпеку}


{\footnotesize\bfseries
No: 3\\
Name: sa\_\allowbreak{}3\_\allowbreak{}a\_\allowbreak{}2\\
Type: list\\
Default: [\textquotedbl{}admin\textquotedbl{}, \textquotedbl{}security\_\allowbreak{}officer\textquotedbl{}]
}

{\footnotesize Система придбана, розроблена та керується з використанням життєвий цикл розробки системи , який охоплює приватність}


{\footnotesize\bfseries
No: 4\\
Name: sa\_\allowbreak{}3\_\allowbreak{}b\_\allowbreak{}1\\
Type: list\\
Default: [\textquotedbl{}admin\textquotedbl{}, \textquotedbl{}security\_\allowbreak{}officer\textquotedbl{}]
}

{\footnotesize Визначені та задокументовані ролі та обов'язки з інформаційної безпеки протягом усього життєвого циклу розробки системи}


{\footnotesize\bfseries
No: 5\\
Name: sa\_\allowbreak{}3\_\allowbreak{}b\_\allowbreak{}2\\
Type: list\\
Default: [\textquotedbl{}admin\textquotedbl{}, \textquotedbl{}security\_\allowbreak{}officer\textquotedbl{}]
}

{\footnotesize Визначені та задокументовані ролі та обов'язки щодо приватності протягом усього життєвого циклу розробки системи}


{\footnotesize\bfseries
No: 6\\
Name: sa\_\allowbreak{}3\_\allowbreak{}c\_\allowbreak{}1\\
Type: list\\
Default: [\textquotedbl{}admin\textquotedbl{}, \textquotedbl{}security\_\allowbreak{}officer\textquotedbl{}]
}

{\footnotesize Визначені особи, які виконують функції та обов'язки з інформаційної безпеки}


{\footnotesize\bfseries
No: 7\\
Name: sa\_\allowbreak{}3\_\allowbreak{}c\_\allowbreak{}2\\
Type: list\\
Default: [\textquotedbl{}admin\textquotedbl{}, \textquotedbl{}security\_\allowbreak{}officer\textquotedbl{}]
}

{\footnotesize Визначені особи з функціями та обов'язками, пов'язаними з приватністю}


{\footnotesize\bfseries
No: 8\\
Name: sa\_\allowbreak{}3\_\allowbreak{}d\_\allowbreak{}1\\
Type: list\\
Default: [\textquotedbl{}admin\textquotedbl{}, \textquotedbl{}security\_\allowbreak{}officer\textquotedbl{}]
}

{\footnotesize Інтегровані процеси управління інформаційною безпекою організації в діяльність життєвого циклу розробки системи}


{\footnotesize\bfseries
No: 9\\
Name: sa\_\allowbreak{}3\_\allowbreak{}d\_\allowbreak{}2\\
Type: list\\
Default: [\textquotedbl{}admin\textquotedbl{}, \textquotedbl{}security\_\allowbreak{}officer\textquotedbl{}]
}

{\footnotesize Інтегровані процеси управління приватністю в життєвого циклу розробки системи}




\subsubsection{СЕРЕДОВИЩЕМ РОЗРОБКИ (SA-3(1))}


\vspace{10pt}
{\footnotesize\bfseries
No: 1\\
Name: sa\_\allowbreak{}3\_\allowbreak{}1\_\allowbreak{}01\\
Type: string\\
Default: nil
}

{\footnotesize Захищене середовище розробки системи, відповідно до ризиків протягом усього життєвого циклу розробки системи для системи, компонентів системи або служб.}




\subsubsection{ЖИТТЄВИЙ ЦИКЛ РЕАЛЬНИХ ДАНИХ (SA-3(2))}


\vspace{10pt}
\textit{Немає параметрів для цього контролю.}
\vspace{10pt}


\subsubsection{ЖИТТЄВИЙ ЦИКЛ ТЕХНОЛОГІЙ (SA-3(3))}


\vspace{10pt}
{\footnotesize\bfseries
No: 1\\
Name: sa\_\allowbreak{}3\_\allowbreak{}3\_\allowbreak{}01\\
Type: string\\
Default: nil
}

{\footnotesize Планується оновлення технологій для підтримки системи протягом усього життєвого циклу розробки системи}


{\footnotesize\bfseries
No: 2\\
Name: sa\_\allowbreak{}3\_\allowbreak{}3\_\allowbreak{}02\\
Type: string\\
Default: nil
}

{\footnotesize Впроваджено графік оновлення технологій для підтримки системи протягом усього життєвого циклу розробки системи}




\subsection{ПРОЦЕС ЗАКУПІВЕЛЬ (SA-4)}
Включіть такі вимоги, описи та критерії, явно або за допомогою посилання, використовуючи [Вибір (один або більше): стандартні пункти контракту; [Призначення: пункти контракту, визначені організацією]] в контракті про придбання системи, системного компонента або системної послуги:\\ 
a. функціональні вимоги безпеки та приватності;\\ 
b. вимоги до стійкості механізму;\\ 
c. вимоги до забезпечення безпеки та приватності;\\ 
d. заходи захисту для забезпечення вимог безпеки та приватності;\\ 
e. вимоги до захисту документації з безпеки та приватності;\\ 
f. опис середовища розробки системи та середовища, у якому система призначена для роботи;\\ 
g. розподіл відповідальності або визначення сторін, відповідальних за управління інформаційною безпекою, приватністю та управлінням ланцюгами постачання;\\ 
h. критерії прийнятності.

\vspace{10pt}
{\footnotesize\bfseries
No: 1\\
Name: sa\_\allowbreak{}4\_\allowbreak{}odp\_\allowbreak{}01\\
Type: string\\
Default: nil
}

{\footnotesize Вибрано одне або декілька з наступних ЗНАЧЕНЬ ПАРАМЕТРА: \{ стандартні пункти контракту; пункт контракту\}}




\subsubsection{ФУНКЦІОНАЛЬНІ ВЛАСТИВОСТІ ЗАХОДІВ (SA-4(1))}


\vspace{10pt}
{\footnotesize\bfseries
No: 1\\
Name: sa\_\allowbreak{}4\_\allowbreak{}1\_\allowbreak{}01\\
Type: string\\
Default: nil
}

{\footnotesize Планується оновлення технологій для системи протягом життєвого циклу розробки системи}




\subsubsection{РОЗРОБКА ТА ВПРОВАДЖЕННЯ ІНФОРМАЦІЇ ДЛЯ ЗАХОДІВ (SA-4(2))}


\vspace{10pt}
\textit{Немає параметрів для цього контролю.}
\vspace{10pt}


\subsubsection{ПРОЦЕС ЗАКУПІВЕЛЬ РОЗРОБКИ (SA-4(3))}


\vspace{10pt}
{\footnotesize\bfseries
No: 1\\
Name: sa\_\allowbreak{}4\_\allowbreak{}3\_\allowbreak{}a\\
Type: string\\
Default: nil
}

{\footnotesize Повинен розробник системи, системного компонента або системної послуги демонструвати використання процесу життєвого циклу розробки системи, який включає методи системної інженерії}




\subsubsection{ВІДНЕСЕННЯ КОМПОНЕНТІВ ДО СИСТЕМ (SA-4(4)) [Вилучено]}
[Вилучено: включено до CM-8(9)]

\vspace{10pt}
\textit{Немає параметрів для цього контролю.}
\vspace{10pt}


\subsubsection{КОНФІГУРАЦІЇ СИСТЕМИ, КОМПОНЕНТА ТА СИСТЕМНОЇ СЛУЖБИ (SA-4(5))}


\vspace{10pt}
{\footnotesize\bfseries
No: 1\\
Name: sa\_\allowbreak{}4\_\allowbreak{}5\_\allowbreak{}a\\
Type: string\\
Default: nil
}

{\footnotesize Повинен розробник системи, компонента системи або системної служби постачати систему, компонент або службу із впровадженими конфігураціями безпеки}


{\footnotesize\bfseries
No: 2\\
Name: sa\_\allowbreak{}4\_\allowbreak{}5\_\allowbreak{}b\\
Type: string\\
Default: nil
}

{\footnotesize Будуть конфігурації використовуватися за замовчуванням для будь-якої наступної переінсталяції або оновлення системи, компонента чи служби}


{\footnotesize\bfseries
No: 3\\
Name: sa\_\allowbreak{}4\_\allowbreak{}5\_\allowbreak{}odp\\
Type: string\\
Default: nil
}

{\footnotesize Визначено конфігурації безпеки для системи, компонента або служби}




\subsubsection{ПРОЦЕС ЗАКУПІВЕЛЬ ІНФОРМАЦІЇ (SA-4(6))}


\vspace{10pt}
{\footnotesize\bfseries
No: 1\\
Name: sa\_\allowbreak{}4\_\allowbreak{}6\_\allowbreak{}a\\
Type: string\\
Default: \textquotedbl{}AES-256-GCM\textquotedbl{}
}

{\footnotesize Використовуються лише засоби захисту інформації, які пройшли державну експертизу або сертифікацію, створені для технічного та криптографічного захисту інформації}


{\footnotesize\bfseries
No: 2\\
Name: sa\_\allowbreak{}4\_\allowbreak{}6\_\allowbreak{}b\\
Type: string\\
Default: \textquotedbl{}автоматизований засіб моніторингу\textquotedbl{}
}

{\footnotesize Були ці засоби захисту мають позитивний експертний висновок або сертифікат відповідності, а також відповідні дозволи для використання для захисту критичної інформації}




\subsubsection{ЗАТВЕРДЖЕНІ ПРОФІЛІ ЗАХИЩЕНОСТІ (SA-4(7))}


\vspace{10pt}
{\footnotesize\bfseries
No: 1\\
Name: sa\_\allowbreak{}4\_\allowbreak{}7\_\allowbreak{}a\\
Type: string\\
Default: nil
}

{\footnotesize Обмежується використання комерційної готової до використання технічної продукції, створеної для захисту інформації та з функцією підтримки забезпечення безпеки інформації, до тих продуктів, які були успішно оцінені відповідно до профілю захищеності для конкретного типу технології, затвердженого уповноваженим державним органом, якщо такий профіль наявний}


{\footnotesize\bfseries
No: 2\\
Name: sa\_\allowbreak{}4\_\allowbreak{}7\_\allowbreak{}b\\
Type: list\\
Default: [\textquotedbl{}default\_\allowbreak{}deny\_\allowbreak{}rule\textquotedbl{}, \textquotedbl{}abac\_\allowbreak{}rule\_\allowbreak{}1\textquotedbl{}]
}

{\footnotesize Якщо немає профілю захищеності для певного типу технологій, затвердженого уповноваженим органом, але забезпечення політики безпеки продукту, що надається на комерційній основі, залежить від криптографічних функцій, --- вимагати, щоб криптографічний модуль пройшов державну експертизу, мав позитивний експертний висновок і був рекомендований до використання уповноваженим органом}




\subsubsection{ПЛАН БЕЗПЕРЕРВНОГО МОНІТОРИНГУ ЗАХОДІВ БЕЗПЕКИ (SA-4(8))}


\vspace{10pt}
{\footnotesize\bfseries
No: 1\\
Name: sa\_\allowbreak{}4\_\allowbreak{}8\_\allowbreak{}01\\
Type: string\\
Default: nil
}

{\footnotesize Розробник системи, системного компонента або системної служби створив план безперервного моніторингу ефективності заходів безпеки та приватності, який узгоджується з відповідним планом постійного моніторингу організації}




\subsubsection{ФУНКЦІЇ, ПОСЛУГИ, ЩО ВИКОРИСТОВУЮТЬСЯ (SA-4(9))}


\vspace{10pt}
{\footnotesize\bfseries
No: 1\\
Name: sa\_\allowbreak{}4\_\allowbreak{}9\_\allowbreak{}01\\
Type: string\\
Default: nil
}

{\footnotesize Зобов'язаний розробник системи, системного компонента або системного сервісу визначити функції, призначені для використання в організації}


{\footnotesize\bfseries
No: 2\\
Name: sa\_\allowbreak{}4\_\allowbreak{}9\_\allowbreak{}02\\
Type: string\\
Default: nil
}

{\footnotesize Зобов'язаний розробник системи, системного компонента або системної служби визначити порти, призначені для використання в організації}


{\footnotesize\bfseries
No: 3\\
Name: sa\_\allowbreak{}4\_\allowbreak{}9\_\allowbreak{}03\\
Type: string\\
Default: \textquotedbl{}TLS 1.3\textquotedbl{}
}

{\footnotesize Зобов'язаний розробник системи, системного компонента або системної служби визначити протоколи, призначені для використання в організації}


{\footnotesize\bfseries
No: 4\\
Name: sa\_\allowbreak{}4\_\allowbreak{}9\_\allowbreak{}04\\
Type: string\\
Default: nil
}

{\footnotesize Зобов'язаний розробник системи, системного компонента або системної послуги визначити послуги, призначені для використання в організації}




\subsection{СИСТЕМНА ДОКУМЕНТАЦІЯ (SA-5)}
a.\\ 
b. Отримати або розробити документацію адміністратора для системи, системного компонента або системної служби, яка описує:\\ 
1. безпечне налаштування, установку та роботу системи, компонента або служби;\\ 
2. ефективне використання, підтримку функцій та механізмів безпеки та приватності;\\ 
3. відомі вразливості щодо конфігурації та використання адміністративних або привілейованих функцій. Отримати або розробити документацію користувача для системи, системного компонента або системної служби, яка описує:\\ 
1. функції та механізми безпеки та приватності та способи ефективного використання цих функцій і механізмів;\\ 
2. методи взаємодії з користувачем, що дозволяють окремим особам використовувати систему, компонент або службу безпечнішим чином та захищати індивідуальну приватність;\\ 
3. обов'язки користувача щодо забезпечення безпеки системи, компонента або служби та приватності окремих осіб.\\ 
c. Документувати спроби отримати доступ до документації системи, системного компонента чи системної служби, коли така документація недоступна або ж відсутня, і вжити [Призначення: визначені організацією заходи] у відповідь.\\ 
d. Поширити документацію серед [Призначення: визначеного організацією персоналу або посадових осіб].

\vspace{10pt}
{\footnotesize\bfseries
No: 1\\
Name: sa\_\allowbreak{}5\_\allowbreak{}odp\_\allowbreak{}01\\
Type: list\\
Default: [\textquotedbl{}login\textquotedbl{}, \textquotedbl{}logout\textquotedbl{}, \textquotedbl{}failed\_\allowbreak{}attempt\textquotedbl{}]
}

{\footnotesize Визначені дії, яких слід вжити, коли документація на систему, системний компонент або системне обслуговування недоступна або відсутня}


{\footnotesize\bfseries
No: 2\\
Name: sa\_\allowbreak{}5\_\allowbreak{}odp\_\allowbreak{}02\\
Type: list\\
Default: [\textquotedbl{}admin\textquotedbl{}, \textquotedbl{}security\_\allowbreak{}officer\textquotedbl{}]
}

{\footnotesize Визначено персонал або системної документації; SA-05a.01[01] отримана або розроблена документація для адміністратора системи, системного компонента або системної служби, яка описує безпечну конфігурацію системи, компонента або служби; SA-05a.01[02] була отримана або розроблена документація для адміністратора системи, системного компонента або системної служби, яка описує безпечне встановлення системи, компонента або служби; SA-05a.01[03] отримана або розроблена ролі для розповсюдження документація адміністратора системи, системного компонента або системної служби, яка описує безпечну роботу системи, компонента або служби; SA-05a.02[01] отримана або розроблена адміністраторська документація системи, системного компонента або системної служби, яка описує ефективне використання функцій та механізмів безпеки; SA-05a.02[02] була отримана або розроблена документація адміністратора системи, системного компонента або системної служби, яка описує безпечне встановлення системи, компонента або служби; SA-05a.02[03] отримана або розроблена документація для адміністратора системи, системного компонента або системної служби, яка описує ефективне використання функцій і механізмів забезпечення конфіденційності; SA-05a.02[04] отримана або розроблена документація для адміністратора системи, системного компонента або системної служби, яка описує ефективну підтримку функцій і механізмів забезпечення конфіденційності; SA-05a.03[01] була отримана або розроблена документація адміністратора системи, системного компонента або системної служби, яка описує відомі вразливості, що стосуються конфігурації адміністративних або привілейованих функцій; SA-05a.03[02] була отримана або розроблена документація адміністратора системи, системного компонента або системної служби, яка описує відомі вразливості, пов'язані з використанням адміністративних або привілейованих функцій; SA-05b.01[01] отримана або розроблена користувацька документація системи, системного компонента або системної служби, яка описує доступні користувачеві функції та механізми безпеки; SA-05b.01[02] отримана або розроблена користувацька документація системи, системного компоненту або системної служби, яка описує, як ефективно використовувати ці (доступні користувачеві) функції та механізми безпеки; SA-05b.01[03] отримана або розроблена користувацька документація системи, системного компонента або системної служби, яка описує доступні користувачеві функції та механізми забезпечення конфіденційності; SA-05b.01[04] отримана або розроблена користувацька документація системи, системного компонента або системної служби, яка описує, як ефективно використовувати ці (доступні користувачеві) функції та механізми захисту приватності; SA-05b.02[01] отримана або розроблена користувацька документація системи, системного компонента або системної послуги, яка описує методи взаємодії з користувачем, що дозволяють особам використовувати систему, компонент або послугу в більш безпечний спосіб; SA-05b.02[02] отримана або розроблена користувацька документація системи, системного компонента або системної служби, яка описує методи взаємодії з користувачем, що дозволяють особам використовувати систему, компонент або службу для захисту особистої приватності; SA-05b.03[01] отримана або розроблена користувацька документація системи, системного компоненту або системного сервісу, яка описує обов'язки користувача щодо підтримання безпеки системи, компоненту або сервісу; SA-05b.03[02] отримана або розроблена користувацька документація системи, системного компонента або системної служби, яка описує обов'язки користувачів щодо збереження приватності приватних осіб; SA-05c.[01] були задокументовані спроби отримати документацію на систему, системний компонент або системне обслуговування, коли така документація або недоступна, або взагалі не існує; SA-05c.[02] після спроб отримати документацію системи, системного компонента або системної служби, коли така документація недоступна або не існує, у відповідь виконуються дії; SA-05d. розповсюджується документація персоналу або ролей>. серед <SA-05\_\allowbreak{}ODP[02]}




\subsubsection{ФУНКЦІОНАЛЬНІ ВЛАСТИВОСТІ ЗАХОДІВ БЕЗПЕКИ (SA-5(1)) [Вилучено]}
[Вилучено: включено до SA-4(1)]

\vspace{10pt}
\textit{Немає параметрів для цього контролю.}
\vspace{10pt}


\subsubsection{ЗОВНІШНІ СИСТЕМНІ ІНТЕРФЕЙСИ, (SA-5(2)) [Вилучено]}
[Вилучено: включено до SA-4(2)]

\vspace{10pt}
\textit{Немає параметрів для цього контролю.}
\vspace{10pt}


\subsubsection{АРХІТЕКТУРА (ПРОЄКТ) ВИСОКОГО РІВНЯ (SA-5(3)) [Вилучено]}
[Вилучено: включено до SA-4(2)]

\vspace{10pt}
\textit{Немає параметрів для цього контролю.}
\vspace{10pt}


\subsubsection{АРХІТЕКТУРА (ПРОЄКТ) НИЗЬКОГО РІВНЯ (SA-5(4)) [Вилучено]}
[Вилучено: включено до SA-4(2)]

\vspace{10pt}
\textit{Немає параметрів для цього контролю.}
\vspace{10pt}


\subsubsection{ВИХІДНИЙ КОД (SA-5(5)) [Вилучено]}
[Вилучено: включено до SA-4(2)]

\vspace{10pt}
\textit{Немає параметрів для цього контролю.}
\vspace{10pt}


\subsection{ОБМЕЖЕННЯ ЩОДО ВИКОРИСТАННЯ ЗАБЕЗПЕЧЕННЯ (SA-6) [Вилучено]}
[Вилучено: включено до CM-10 та SI-7]

\vspace{10pt}
\textit{Немає параметрів для цього контролю.}
\vspace{10pt}


\subsection{ВСТАНОВЛЕНЕ КОРИСТУВАЧЕМ ПРОГРАМНЕ ЗАБЕЗПЕЧЕННЯ (SA-7) [Вилучено]}
[Вилучено: включено до CM-11 та SI-7]

\vspace{10pt}
\textit{Немає параметрів для цього контролю.}
\vspace{10pt}


\subsection{БЕЗПЕКА ТА ПРИВАТНІСТЬ ПРИНЦИПІВ ІНЖИНІРИНГУ (SA-8)}
Застосовувати [Призначення: визначені організацією принципи інжинірингу безпеки та конфіденційності системи] в специфікації, проєктуванні, розробці, впровадженні та зміні системи й компонентів системи.

\vspace{10pt}
{\footnotesize\bfseries
No: 1\\
Name: sa\_\allowbreak{}8\_\allowbreak{}01\\
Type: string\\
Default: nil
}

{\footnotesize Застосовуються принципи інжинірингу безпеки систем у специфікації системи та компонентів системи}


{\footnotesize\bfseries
No: 2\\
Name: sa\_\allowbreak{}8\_\allowbreak{}02\\
Type: string\\
Default: nil
}

{\footnotesize Застосовуються принципи інжинірингу безпеки систем при розробці системи та її компонентів}


{\footnotesize\bfseries
No: 3\\
Name: sa\_\allowbreak{}8\_\allowbreak{}03\\
Type: string\\
Default: nil
}

{\footnotesize Були застосовані принципи інжинірингу безпеки систем при розробці системи та компонентів систем}


{\footnotesize\bfseries
No: 4\\
Name: sa\_\allowbreak{}8\_\allowbreak{}04\\
Type: string\\
Default: nil
}

{\footnotesize Застосовуються принципи інжинірингу безпеки систем при реалізації системи та компонентів систем}


{\footnotesize\bfseries
No: 5\\
Name: sa\_\allowbreak{}8\_\allowbreak{}05\\
Type: string\\
Default: nil
}

{\footnotesize Застосовуються принципи інжинірингу безпеки систем при модифікації системи та компонентів систем}


{\footnotesize\bfseries
No: 6\\
Name: sa\_\allowbreak{}8\_\allowbreak{}06\\
Type: string\\
Default: nil
}

{\footnotesize Застосовуються принципи інжинірингу конфіденційності у специфікації системи та компонентів систем; принципи інжинірингу конфіденційності SA-08[07] застосовуються принципи інжинірингу конфіденційності при розробці системи та компонентів систем}


{\footnotesize\bfseries
No: 7\\
Name: sa\_\allowbreak{}8\_\allowbreak{}08\\
Type: string\\
Default: nil
}

{\footnotesize Застосовуються принципи інжинірингу конфіденційності при розробці системи та компонентів систем}


{\footnotesize\bfseries
No: 8\\
Name: sa\_\allowbreak{}8\_\allowbreak{}09\\
Type: string\\
Default: nil
}

{\footnotesize Застосовуються принципи інжинірингу конфіденційності при реалізації системи та компонентів систем}


{\footnotesize\bfseries
No: 9\\
Name: sa\_\allowbreak{}8\_\allowbreak{}10\\
Type: string\\
Default: nil
}

{\footnotesize Застосовуються принципи інжинірингу конфіденційності при модифікації системи та компонентів систем}


{\footnotesize\bfseries
No: 10\\
Name: sa\_\allowbreak{}8\_\allowbreak{}odp\_\allowbreak{}01\\
Type: string\\
Default: nil
}

{\footnotesize Визначені принципи інжинірингу безпеки систем}


{\footnotesize\bfseries
No: 11\\
Name: sa\_\allowbreak{}8\_\allowbreak{}odp\_\allowbreak{}02\\
Type: string\\
Default: nil
}

{\footnotesize Визначені системи}




\subsubsection{ЧІТКА АБСТРАКЦІЯ (SA-8(1))}


\vspace{10pt}
{\footnotesize\bfseries
No: 1\\
Name: sa\_\allowbreak{}8\_\allowbreak{}1\_\allowbreak{}01\\
Type: string\\
Default: nil
}

{\footnotesize Реалізовано принцип проектування безпеки чітких абстракцій}




\subsubsection{БЕЗПЕКА ТА ПРИВАТНІСТЬ ПРИНЦИПІВ НАЙМЕНШ ПОШИРЕНИЙ МЕХАНІЗМ (SA-8(2))}


\vspace{10pt}
{\footnotesize\bfseries
No: 1\\
Name: sa\_\allowbreak{}8\_\allowbreak{}2\_\allowbreak{}01\\
Type: string\\
Default: nil
}

{\footnotesize Реалізують системи або компоненти системи принцип побудови безпеки за принципом найменш поширеного механізму}


{\footnotesize\bfseries
No: 2\\
Name: sa\_\allowbreak{}8\_\allowbreak{}2\_\allowbreak{}odp\\
Type: string\\
Default: nil
}

{\footnotesize Реалізовано абстракцій. принцип проектування безпеки чітких}




\subsubsection{БЕЗПЕКА ТА ПРИВАТНІСТЬ ПРИНЦИПІВ МОДУЛЬНІСТЬ І БАГАТОРІВНЕВІСТЬ (SA-8(3))}


\vspace{10pt}
{\footnotesize\bfseries
No: 1\\
Name: sa\_\allowbreak{}8\_\allowbreak{}3\_\allowbreak{}01\\
Type: string\\
Default: nil
}

{\footnotesize БЕЗПЕКА ТА ПРИВАТНІСТЬ ПРИНЦИПІВ МОДУЛЬНІСТЬ І БАГАТОРІВНЕВІСТЬ ІНЖИНІРИНГУ - МЕТА ОЦІНКИ: Визначити, чи: SA08(03)\_\allowbreak{}ODP[01] визначені системи або компоненти системи, реалізують принцип модульності дизайну безпеки; які SA08(03)\_\allowbreak{}ODP[02] визначені системи або компоненти реалізують принцип багаторівневого безпеки}


{\footnotesize\bfseries
No: 2\\
Name: sa\_\allowbreak{}8\_\allowbreak{}3\_\allowbreak{}02\\
Type: string\\
Default: nil
}

{\footnotesize Системи або компоненти системи реалізують принцип багаторівневого проектування безпеки. системи, які проектування}




\subsubsection{БЕЗПЕКА ТА ПРИВАТНІСТЬ ПРИНЦИПІВ ЧАСТКОВО ВПОРЯДКОВАНІ ЗАЛЕЖНОСТІ (SA-8(4))}


\vspace{10pt}
{\footnotesize\bfseries
No: 1\\
Name: sa\_\allowbreak{}8\_\allowbreak{}4\_\allowbreak{}01\\
Type: integer\\
Default: 30
}

{\footnotesize Системи або компоненти системи реалізують принцип проектування безпеки частково впорядкованих залежностей}


{\footnotesize\bfseries
No: 2\\
Name: sa\_\allowbreak{}8\_\allowbreak{}4\_\allowbreak{}odp\\
Type: integer\\
Default: 30
}

{\footnotesize Визначені системи або компоненти системи, які реалізують принцип проектування безпеки частково впорядкованих залежностей}




\subsubsection{БЕЗПЕКА ТА ПРИВАТНІСТЬ ПРИНЦИПІВ ЕФЕКТИВНИЙ ОПОСЕРЕДКОВАНИЙ ДОСТУП (SA-8(5))}


\vspace{10pt}
{\footnotesize\bfseries
No: 1\\
Name: sa\_\allowbreak{}8\_\allowbreak{}5\_\allowbreak{}01\\
Type: string\\
Default: nil
}

{\footnotesize Системи або компоненти системи реалізують принцип проектування безпеки ефективного опосередкованого доступу}


{\footnotesize\bfseries
No: 2\\
Name: sa\_\allowbreak{}8\_\allowbreak{}5\_\allowbreak{}odp\\
Type: string\\
Default: nil
}

{\footnotesize Визначені системи або її компоненти, які реалізують принцип проектування безпеки ефективного опосередкованого доступу}




\subsubsection{БЕЗПЕКА ТА ПРИВАТНІСТЬ МІНІМІЗОВАНИЙ ОБМІН (SA-8(6))}


\vspace{10pt}
{\footnotesize\bfseries
No: 1\\
Name: sa\_\allowbreak{}8\_\allowbreak{}6\_\allowbreak{}01\\
Type: string\\
Default: nil
}

{\footnotesize Системи або компоненти системи реалізують принцип побудови безпеки за принципом мінімізації спільного використання}


{\footnotesize\bfseries
No: 2\\
Name: sa\_\allowbreak{}8\_\allowbreak{}6\_\allowbreak{}odp\\
Type: string\\
Default: nil
}

{\footnotesize Визначені системи або компоненти системи, які реалізують принцип проектування безпеки, що полягає в мінімізації спільного використання}




\subsubsection{ЗНИЖЕНА СКЛАДНІСТЬ (SA-8(7))}


\vspace{10pt}
{\footnotesize\bfseries
No: 1\\
Name: sa\_\allowbreak{}8\_\allowbreak{}7\_\allowbreak{}01\\
Type: string\\
Default: nil
}

{\footnotesize Системи або компоненти системи реалізують принцип проектування безпеки за принципом зниженої складності. ПОТЕНЦІЙНІ МЕТОДИ ТА ОБ'ЄКТИ ОЦІНКИ:}


{\footnotesize\bfseries
No: 2\\
Name: sa\_\allowbreak{}8\_\allowbreak{}7\_\allowbreak{}odp\\
Type: string\\
Default: nil
}

{\footnotesize Визначені системи або компоненти системи, які реалізують принцип проектування безпеки за принципом зниженої складності}




\subsubsection{БЕЗПЕКА ТА ПРИВАТНІСТЬ ПРИНЦИПІВ ЕВОЛЮЦІЯ БЕЗПЕКИ В СИСТЕМІ (SA-8(8))}


\vspace{10pt}
{\footnotesize\bfseries
No: 1\\
Name: sa\_\allowbreak{}8\_\allowbreak{}8\_\allowbreak{}01\\
Type: string\\
Default: nil
}

{\footnotesize Системи або компоненти системи реалізують принцип безпечного проектування безпечної еволюційності}


{\footnotesize\bfseries
No: 2\\
Name: sa\_\allowbreak{}8\_\allowbreak{}8\_\allowbreak{}odp\\
Type: string\\
Default: nil
}

{\footnotesize Визначені системи або компоненти системи, які реалізують принцип безпечного проектування безпечної еволюційності}




\subsubsection{ДОВІРЕНІ КОМПОНЕНТИ СИСТЕМИ (SA-8(9))}


\vspace{10pt}
{\footnotesize\bfseries
No: 1\\
Name: sa\_\allowbreak{}8\_\allowbreak{}9\_\allowbreak{}01\\
Type: string\\
Default: nil
}

{\footnotesize Системи або компоненти реалізують принцип побудови безпеки компонентів. системи довірених}


{\footnotesize\bfseries
No: 2\\
Name: sa\_\allowbreak{}8\_\allowbreak{}9\_\allowbreak{}odp\\
Type: string\\
Default: nil
}

{\footnotesize Визначені системи або компоненти системи, які реалізують принцип побудови безпеки довірених компонентів}




\subsubsection{РЕАЛІЗУЮТЬ ПРИНЦИП ІЄРАРХІЧНОЇ ДОВІРИ В ДИЗАЙНІ БЕЗПЕКИ (SA-8(10))}


\vspace{10pt}
{\footnotesize\bfseries
No: 1\\
Name: sa\_\allowbreak{}8\_\allowbreak{}10\_\allowbreak{}01\\
Type: string\\
Default: nil
}

{\footnotesize Системи або компоненти системи реалізують принцип ієрархічної довіри в дизайні безпеки. ПОТЕНЦІЙНІ МЕТОДИ ТА ОБ'ЄКТИ ОЦІНКИ:}


{\footnotesize\bfseries
No: 2\\
Name: sa\_\allowbreak{}8\_\allowbreak{}10\_\allowbreak{}odp\\
Type: string\\
Default: nil
}

{\footnotesize Визначені системи або компоненти системи, які реалізують принцип ієрархічної довіри при проектуванні безпеки}




\subsubsection{РЕАЛІЗУЮТЬ ПРИНЦИП ПРОЕКТУВАННЯ БЕЗПЕКИ ЗВОРОТНОГО ПОРОГУ МОДИФІКАЦІЇ (SA-8(11))}


\vspace{10pt}
{\footnotesize\bfseries
No: 1\\
Name: sa\_\allowbreak{}8\_\allowbreak{}11\_\allowbreak{}01\\
Type: string\\
Default: nil
}

{\footnotesize Системи або компоненти системи реалізують принцип проектування безпеки зворотного порогу модифікації}


{\footnotesize\bfseries
No: 2\\
Name: sa\_\allowbreak{}8\_\allowbreak{}11\_\allowbreak{}odp\\
Type: string\\
Default: nil
}

{\footnotesize Визначені системи або компоненти системи, які реалізують принцип ієрархічної довіри при проектуванні безпеки}




\subsubsection{РЕАЛІЗУЮТЬ ПРИНЦИП ПОБУДОВИ ІЄРАРХІЧНОГО ЗАХИСТУ (SA-8(12))}


\vspace{10pt}
{\footnotesize\bfseries
No: 1\\
Name: sa\_\allowbreak{}8\_\allowbreak{}12\_\allowbreak{}01\\
Type: string\\
Default: nil
}

{\footnotesize Системи або компоненти системи реалізують принцип побудови ієрархічного захисту}


{\footnotesize\bfseries
No: 2\\
Name: sa\_\allowbreak{}8\_\allowbreak{}12\_\allowbreak{}odp\\
Type: string\\
Default: nil
}

{\footnotesize Визначені системи або компоненти системи, реалізують принцип ієрархічного дизайну безпеки; які}




\subsubsection{РЕАЛІЗУЮТЬ ПРИНЦИП ПРОЕКТУВАННЯ БЕЗПЕКИ З МІНІМІЗАЦІЄЮ ЕЛЕМЕНТІВ БЕЗПЕКИ (SA-8(13))}


\vspace{10pt}
{\footnotesize\bfseries
No: 1\\
Name: sa\_\allowbreak{}8\_\allowbreak{}13\_\allowbreak{}01\\
Type: string\\
Default: nil
}

{\footnotesize Системи або компоненти системи реалізують принцип проектування безпеки з мінімізацією елементів безпеки. ПОТЕНЦІЙНІ МЕТОДИ ТА ОБ'ЄКТИ ОЦІНКИ:}


{\footnotesize\bfseries
No: 2\\
Name: sa\_\allowbreak{}8\_\allowbreak{}13\_\allowbreak{}odp\\
Type: string\\
Default: nil
}

{\footnotesize Визначено системи або компоненти системи, які реалізують принцип проектування безпеки за принципом мінімізації елементів безпеки}




\subsubsection{РЕАЛІЗУЮТЬ ПРИНЦИП ПОБУДОВИ БЕЗПЕКИ ЗА ПРИНЦИПОМ НАЙМЕНШИХ ПРИВІЛЕЇВ (SA-8(14))}


\vspace{10pt}
{\footnotesize\bfseries
No: 1\\
Name: sa\_\allowbreak{}8\_\allowbreak{}14\_\allowbreak{}01\\
Type: string\\
Default: nil
}

{\footnotesize Системи або компоненти системи реалізують принцип побудови безпеки за принципом найменших привілеїв}


{\footnotesize\bfseries
No: 2\\
Name: sa\_\allowbreak{}8\_\allowbreak{}14\_\allowbreak{}odp\\
Type: string\\
Default: nil
}

{\footnotesize Визначено системи або компоненти системи, які реалізують принцип побудови безпеки за принципом найменших привілеїв}




\subsubsection{РЕАЛІЗУЮТЬ ПРИНЦИП ПРОЕКТУВАННЯ БЕЗПЕКИ ПОПЕРЕДНЬОГО ДОЗВОЛУ (SA-8(15))}


\vspace{10pt}
{\footnotesize\bfseries
No: 1\\
Name: sa\_\allowbreak{}8\_\allowbreak{}15\_\allowbreak{}01\\
Type: string\\
Default: nil
}

{\footnotesize Системи або компоненти системи реалізують принцип проектування безпеки попереднього дозволу}


{\footnotesize\bfseries
No: 2\\
Name: sa\_\allowbreak{}8\_\allowbreak{}15\_\allowbreak{}odp\\
Type: string\\
Default: nil
}

{\footnotesize Визначено системи або компоненти системи, які реалізують принцип проектування безпеки попереднього дозволу}




\subsubsection{РЕАЛІЗУЮТЬ ПРИНЦИП ПРОЕКТУВАННЯ БЕЗПЕКИ САМОДОСТАТНЬОЇ НАДІЙНОСТІ (SA-8(16))}


\vspace{10pt}
{\footnotesize\bfseries
No: 1\\
Name: sa\_\allowbreak{}8\_\allowbreak{}16\_\allowbreak{}01\\
Type: string\\
Default: nil
}

{\footnotesize Системи або компоненти системи реалізують принцип проектування безпеки самодостатньої надійності. ПОТЕНЦІЙНІ МЕТОДИ ТА ОБ'ЄКТИ ОЦІНКИ:}


{\footnotesize\bfseries
No: 2\\
Name: sa\_\allowbreak{}8\_\allowbreak{}16\_\allowbreak{}odp\\
Type: string\\
Default: nil
}

{\footnotesize Визначено системи або компоненти реалізують принцип проектування ґрунтується на самодостатній надійності; системи, безпеки, які що}




\subsubsection{РЕАЛІЗУЮТЬ ПРИНЦИП ПРОЕКТУВАННЯ БЕЗПЕКИ ЗАХИЩЕНОГО РОЗПОДІЛЕНОГО ВМІСТУ (SA-8(17))}


\vspace{10pt}
{\footnotesize\bfseries
No: 1\\
Name: sa\_\allowbreak{}8\_\allowbreak{}17\_\allowbreak{}01\\
Type: string\\
Default: nil
}

{\footnotesize Системи або компоненти системи реалізують принцип проектування безпеки захищеного розподіленого вмісту}


{\footnotesize\bfseries
No: 2\\
Name: sa\_\allowbreak{}8\_\allowbreak{}17\_\allowbreak{}odp\\
Type: string\\
Default: nil
}

{\footnotesize Визначено системи або компоненти системи, які реалізують принцип безпечного проектування захищеного розподіленого вмісту}




\subsubsection{РЕАЛІЗУЮТЬ ПРИНЦИП ПОБУДОВИ БЕЗПЕКИ ДОВІРЕНИХ КАНАЛІВ ЗВ'ЯЗКУ (SA-8(18))}


\vspace{10pt}
{\footnotesize\bfseries
No: 1\\
Name: sa\_\allowbreak{}8\_\allowbreak{}18\_\allowbreak{}01\\
Type: string\\
Default: nil
}

{\footnotesize Системи або компоненти системи реалізують принцип побудови безпеки довірених каналів зв'язку}


{\footnotesize\bfseries
No: 2\\
Name: sa\_\allowbreak{}8\_\allowbreak{}18\_\allowbreak{}odp\\
Type: string\\
Default: nil
}

{\footnotesize Визначено системи або компоненти системи, які реалізують принцип побудови безпеки довірених каналів зв'язку}




\subsubsection{РЕАЛІЗУЮТЬ ПРИНЦИП ПРОЕКТУВАННЯ БЕЗПЕКИ ДОВГОТРИВАЛОГО (SA-8(19))}


\vspace{10pt}
{\footnotesize\bfseries
No: 1\\
Name: sa\_\allowbreak{}8\_\allowbreak{}19\_\allowbreak{}01\\
Type: string\\
Default: nil
}

{\footnotesize Системи або компоненти системи реалізують принцип проектування безпеки довготривалого захисту}


{\footnotesize\bfseries
No: 2\\
Name: sa\_\allowbreak{}8\_\allowbreak{}19\_\allowbreak{}odp\\
Type: string\\
Default: nil
}

{\footnotesize Визначено системи або компоненти системи, які втілюють принцип проектування безпеки довготривалого захисту}




\subsubsection{РЕАЛІЗУЮТЬ ПРИНЦИП БЕЗПЕЧНОГО УПРАВЛІННЯ МЕТАДАНИМИ (SA-8(20))}


\vspace{10pt}
{\footnotesize\bfseries
No: 1\\
Name: sa\_\allowbreak{}8\_\allowbreak{}20\_\allowbreak{}01\\
Type: string\\
Default: nil
}

{\footnotesize Системи або компоненти системи реалізують принцип безпечного управління метаданими}


{\footnotesize\bfseries
No: 2\\
Name: sa\_\allowbreak{}8\_\allowbreak{}20\_\allowbreak{}odp\\
Type: string\\
Default: nil
}

{\footnotesize Визначено системи або компоненти системи, які реалізують принцип безпечного управління метаданими}




\subsubsection{РЕАЛІЗУЮТЬ ПРИНЦИП ПРОЕКТУВАННЯ БЕЗПЕКИ НА ОСНОВІ САМОАНАЛІЗУ (SA-8(21))}


\vspace{10pt}
{\footnotesize\bfseries
No: 1\\
Name: sa\_\allowbreak{}8\_\allowbreak{}21\_\allowbreak{}01\\
Type: string\\
Default: nil
}

{\footnotesize Системи або компоненти системи реалізують принцип проектування безпеки на основі самоаналізу}


{\footnotesize\bfseries
No: 2\\
Name: sa\_\allowbreak{}8\_\allowbreak{}21\_\allowbreak{}odp\\
Type: string\\
Default: nil
}

{\footnotesize Визначено системи або компоненти системи, які реалізують принцип самоаналізу при проектуванні безпеки}




\subsubsection{РЕАЛІЗУЮТЬ ПРИНЦИП ПРОЕКТУВАННЯ БЕЗПЕЧНИХ НАЛАШТУВАНЬ ЗА ЗАМОВЧУВАННЯМ (SA-8(23))}


\vspace{10pt}
{\footnotesize\bfseries
No: 1\\
Name: sa\_\allowbreak{}8\_\allowbreak{}23\_\allowbreak{}odp\\
Type: string\\
Default: nil
}

{\footnotesize Визначено системи або компоненти системи, реалізують принцип безпечних налаштувань замовчуванням; які за SA-08(23) системи або компоненти системи реалізують принцип проектування безпечних налаштувань за замовчуванням}




\subsubsection{<SA-08(25)\_\allowbreak{}ODP системи або компоненти реалізують принцип безпеки економіки. (SA-8(25))}


\vspace{10pt}
{\footnotesize\bfseries
No: 1\\
Name: sa\_\allowbreak{}8\_\allowbreak{}25\_\allowbreak{}01\\
Type: string\\
Default: nil
}

{\footnotesize Системи або компоненти реалізують принцип безпеки економіки. які системи}


{\footnotesize\bfseries
No: 2\\
Name: sa\_\allowbreak{}8\_\allowbreak{}25\_\allowbreak{}odp\\
Type: string\\
Default: nil
}

{\footnotesize Визначено системи або компоненти реалізують принцип безпеки економіки; системи,}




\subsubsection{<SA-08(26)\_\allowbreak{}ODP системи або компоненти реалізують принцип безпеки продуктивності. (SA-8(26))}


\vspace{10pt}
{\footnotesize\bfseries
No: 1\\
Name: sa\_\allowbreak{}8\_\allowbreak{}26\_\allowbreak{}01\\
Type: string\\
Default: nil
}

{\footnotesize Системи або компоненти реалізують принцип безпеки продуктивності. які системи}


{\footnotesize\bfseries
No: 2\\
Name: sa\_\allowbreak{}8\_\allowbreak{}26\_\allowbreak{}odp\\
Type: string\\
Default: nil
}

{\footnotesize Визначено системи або компоненти системи, реалізують принцип безпеки продуктивності}




\subsubsection{РЕАЛІЗУЮТЬ ПРИНЦИП БЕЗПЕКИ З УРАХУВАННЯМ ЛЮДСЬКОГО ФАКТОРУ (SA-8(27))}


\vspace{10pt}
{\footnotesize\bfseries
No: 1\\
Name: sa\_\allowbreak{}8\_\allowbreak{}27\_\allowbreak{}01\\
Type: string\\
Default: nil
}

{\footnotesize Системи або компоненти системи реалізують принцип безпеки з урахуванням людського фактору}


{\footnotesize\bfseries
No: 2\\
Name: sa\_\allowbreak{}8\_\allowbreak{}27\_\allowbreak{}odp\\
Type: string\\
Default: nil
}

{\footnotesize Визначено системи або компоненти системи, які реалізують принцип безпеки з урахуванням людського фактору}




\subsubsection{<SA-08(28)\_\allowbreak{}ODP системи або компоненти реалізують принцип прийнятного рівня безпеки. (SA-8(28))}


\vspace{10pt}
{\footnotesize\bfseries
No: 1\\
Name: sa\_\allowbreak{}8\_\allowbreak{}28\_\allowbreak{}01\\
Type: string\\
Default: nil
}

{\footnotesize Системи або компоненти системи реалізують принцип повторюваних та задокументованих процедур}


{\footnotesize\bfseries
No: 2\\
Name: sa\_\allowbreak{}8\_\allowbreak{}28\_\allowbreak{}odp\\
Type: string\\
Default: nil
}

{\footnotesize Визначено системи або компоненти системи, реалізують принцип прийнятного рівня безпеки}




\subsubsection{РЕАЛІЗУЮТЬ ПРИНЦИП БЕЗПЕЧНОЇ МОДИФІКАЦІЇ СИСТЕМ (SA-8(31))}


\vspace{10pt}
{\footnotesize\bfseries
No: 1\\
Name: sa\_\allowbreak{}8\_\allowbreak{}31\_\allowbreak{}01\\
Type: string\\
Default: nil
}

{\footnotesize Системи або компоненти системи реалізують принцип безпечної модифікації систем}


{\footnotesize\bfseries
No: 2\\
Name: sa\_\allowbreak{}8\_\allowbreak{}31\_\allowbreak{}odp\\
Type: string\\
Default: nil
}

{\footnotesize Визначено системи або компоненти системи, реалізують принцип безпечної модифікації систем; які}




\subsubsection{РЕАЛІЗУЮТЬ ПРИНЦИП ДОСТАТНЬОЇ ДОПУСКНОЇ ДОКУМЕНТАЦІЇ (SA-8(32))}


\vspace{10pt}
{\footnotesize\bfseries
No: 1\\
Name: sa\_\allowbreak{}8\_\allowbreak{}32\_\allowbreak{}01\\
Type: string\\
Default: nil
}

{\footnotesize Системи або компоненти системи реалізують принцип достатньої допускної документації}


{\footnotesize\bfseries
No: 2\\
Name: sa\_\allowbreak{}8\_\allowbreak{}32\_\allowbreak{}odp\\
Type: string\\
Default: nil
}

{\footnotesize Визначено системи або компоненти системи, які реалізують принцип достатньої допускної документації}




\subsubsection{<SA-08(33)\_\allowbreak{}ODP системи або компоненти реалізують принцип мінімізації конфіденційності. (SA-8(33))}


\vspace{10pt}
{\footnotesize\bfseries
No: 1\\
Name: sa\_\allowbreak{}8\_\allowbreak{}33\_\allowbreak{}odp\\
Type: string\\
Default: nil
}

{\footnotesize Визначено системи або компоненти системи, реалізують принцип мінімізації конфіденційності; які SA-08(33) системи або компоненти реалізують принцип мінімізації конфіденційності. системи}




\subsection{ЗОВНІШНІ ПОСЛУГИ ДЛЯ СИСТЕМИ (SA-9)}
a. Вимагати, щоб постачальники зовнішніх послуг для системи відповідали вимогам безпеки та приватності в організації та застосовували такі заходи захисту [Призначення: встановлені організацією заходи безпеки та приватності].\\ 
b. Визначити та задокументувати нагляд організаціїповноваження та обов'язки користувачів щодо зовнішніх послуг для системи.\\ 
c. Використовувати наступні процеси, методи та техніки для постійного моніторингу дотримання контролю зовнішніми постачальниками послуг: [Призначення: визначені організацією процеси, методи та техніки].

\vspace{10pt}
{\footnotesize\bfseries
No: 1\\
Name: sa\_\allowbreak{}9\_\allowbreak{}odp\_\allowbreak{}01\\
Type: string\\
Default: \textquotedbl{}автоматизований засіб моніторингу\textquotedbl{}
}

{\footnotesize Визначені засоби контролю, які будуть застосовуватися зовнішніми постачальниками системних послуг}


{\footnotesize\bfseries
No: 2\\
Name: sa\_\allowbreak{}9\_\allowbreak{}odp\_\allowbreak{}02\\
Type: list\\
Default: [\textquotedbl{}admin\textquotedbl{}, \textquotedbl{}security\_\allowbreak{}officer\textquotedbl{}]
}

{\footnotesize Визначені процеси, методи та техніки, що застосовуються для моніторингу дотримання вимог контролю зовнішніми постачальниками послуг; SA-09a.[01] дотримуються постачальники зовнішніх системних послуг вимог організаційної безпеки; SA-09a.[02] дотримуються постачальники зовнішніх системних послуг вимог приватності організації; SA-09a.[03] використовують постачальники зовнішніх системних послуг елементи керування; SA-09b.[01] визначено та задокументовано організаційний нагляд за зовнішніми системними послугами; SA-09b.[02] визначені та задокументовані ролі та обов'язки користувачів щодо зовнішніх системних сервісів; SA-09c. визначені та задокументовані ролі та обов'язки користувачів щодо зовнішніх системних послуг}




\subsubsection{ЗОВНІШНІ ПОСЛУГИ ДЛЯ СИСТЕМИ- ОЦІНЮВАННЯ РИЗИКІВ ТА ОРГАНІЗАЦІЙНІ ПОГОДЖЕННЯ (SA-9(1))}
(a) Проводити організаційне оцінювання ризиків перед придбанням або переданням послуг інформаційної безпеки служб інформаційної безпеки.\\ 
(b) Переконатися, що придбання або передача спеціалізованих служб інформаційної безпеки погоджені [Призначення: визначеним організацією персоналом або посадовими особами].

\vspace{10pt}
{\footnotesize\bfseries
No: 1\\
Name: sa\_\allowbreak{}9\_\allowbreak{}1\_\allowbreak{}a\\
Type: string\\
Default: nil
}

{\footnotesize Проводиться організаційна оцінка ризиків перед придбанням або передачею послуг з інформаційної безпеки}


{\footnotesize\bfseries
No: 2\\
Name: sa\_\allowbreak{}9\_\allowbreak{}1\_\allowbreak{}b\\
Type: list\\
Default: [\textquotedbl{}admin\textquotedbl{}, \textquotedbl{}security\_\allowbreak{}officer\textquotedbl{}]
}

{\footnotesize Схвалюють персонал або ролі придбання або передачу спеціальних послуг з інформаційної безпеки}


{\footnotesize\bfseries
No: 3\\
Name: sa\_\allowbreak{}9\_\allowbreak{}1\_\allowbreak{}odp\\
Type: list\\
Default: [\textquotedbl{}admin\textquotedbl{}, \textquotedbl{}security\_\allowbreak{}officer\textquotedbl{}]
}

{\footnotesize Визначено персонал або ролі, які схвалюють придбання або передачу спеціальних послуг з інформаційної безпеки}




\subsubsection{ЗОВНІШНІ ПОСЛУГИ ДЛЯ СИСТЕМИ- ВИЗНАЧЕННЯ ФУНКЦІЙ, ПОРТІВ, ПРОТОКОЛІВ ТА СЛУЖБ (SA-9(2))}
Вимагати від постачальників наведених нижче зовнішніх послуг для системи [Призначення: визначених організацією зовнішніх послуг для системи] визначити функції, порти, протоколи та інші служби, необхідні для використання таких служб.

\vspace{10pt}
{\footnotesize\bfseries
No: 1\\
Name: sa\_\allowbreak{}9\_\allowbreak{}2\_\allowbreak{}01\\
Type: string\\
Default: \textquotedbl{}TLS 1.3\textquotedbl{}
}

{\footnotesize Необхідно постачальникам зовнішніх системних послуг ідентифікувати функції, порти, протоколи та інші послуги, необхідні для використання таких послуг}


{\footnotesize\bfseries
No: 2\\
Name: sa\_\allowbreak{}9\_\allowbreak{}2\_\allowbreak{}odp\\
Type: string\\
Default: \textquotedbl{}TLS 1.3\textquotedbl{}
}

{\footnotesize Визначені зовнішні системні сервіси, які потребують ідентифікації функцій, портів, протоколів та інших сервісів}




\subsubsection{СТВОРЕННЯ ТА ПІДТРИМКА ДОВІРЧИХ ВІДНОСИН З ПОСТАЧАЛЬНИКАМИ (SA-9(3))}
Створити, задокументувати та підтримувати довірчі відносини із зовнішніми постачальниками послуг на основі таких вимог, властивостей, факторів або умов: [Призначення: визначених організацією вимог, властивостей, факторів або умов щодо безпеки та приватності, що визначають прийнятні довірчі відносини].

\vspace{10pt}
{\footnotesize\bfseries
No: 1\\
Name: sa\_\allowbreak{}9\_\allowbreak{}3\_\allowbreak{}01\\
Type: string\\
Default: nil
}

{\footnotesize Встановлені та задокументовані довірчі відносини з зовнішніми постачальниками послуг на основі вимог, властивостей, факторів або умов безпеки}


{\footnotesize\bfseries
No: 2\\
Name: sa\_\allowbreak{}9\_\allowbreak{}3\_\allowbreak{}02\\
Type: string\\
Default: nil
}

{\footnotesize Підтримуються довірчі відносини з зовнішніми постачальниками послуг на основі вимог, властивостей, факторів або умов безпеки}


{\footnotesize\bfseries
No: 3\\
Name: sa\_\allowbreak{}9\_\allowbreak{}3\_\allowbreak{}03\\
Type: string\\
Default: nil
}

{\footnotesize Встановлені та задокументовані довірчі відносини із зовнішніми постачальниками послуг на основі вимог, властивостей, факторів або умов конфіденційності}


{\footnotesize\bfseries
No: 4\\
Name: sa\_\allowbreak{}9\_\allowbreak{}3\_\allowbreak{}04\\
Type: string\\
Default: nil
}

{\footnotesize Підтримуються довірчі відносини із зовнішніми постачальниками послуг на основі вимог, властивостей, факторів або умов конфіденційності}




\subsubsection{ЗОВНІШНІ СИСТЕМНІ СЛУЖБИ СПОЖИВАЧІВ І ПОСТАЧАЛЬНИКІВ (SA-9(4))}
Виконайте такі дії, щоб переконатися, що інтереси [Призначення: визначених організацією зовнішніх постачальників послуг] узгоджуються з інтересами організації та відображають їх: [Призначення: дії, визначені організацією].

\vspace{10pt}
{\footnotesize\bfseries
No: 1\\
Name: sa\_\allowbreak{}9\_\allowbreak{}4\_\allowbreak{}01\\
Type: list\\
Default: [\textquotedbl{}login\textquotedbl{}, \textquotedbl{}logout\textquotedbl{}, \textquotedbl{}failed\_\allowbreak{}attempt\textquotedbl{}]
}

{\footnotesize Вживаються дії для перевірки того, що інтереси зовнішніх постачальників послуг узгоджуються з інтересами організації та відображають їх}




\subsubsection{ЗОВНІШНІ ПОСЛУГИ ДЛЯ СИСТЕМИ- МІСЦЕ ОБРОБКИ, ЗБЕРІГАННЯ ТА ОБСЛУГОВУВАННЯ (SA-9(5))}
Обмежити розташування [Вибір (один або більше): обробка інформації; інформація або дані; системні служби] до [Призначення: визначені організацією місця] на основі [Призначення: визначених організацією вимог або умов].

\vspace{10pt}
\textit{Немає параметрів для цього контролю.}
\vspace{10pt}


\subsubsection{ЗОВНІШНІ ПОСЛУГИ ДЛЯ СИСТЕМИ- КРИПТОГРАФІЧНІ КЛЮЧІ, КЕРОВАНІ ОРГАНІЗАЦІЄЮ (SA-9(6))}
Зберігати контроль над криптографічними ключами для зашифрованої інформації, яка зберігається або передається через зовнішню систему.

\vspace{10pt}
{\footnotesize\bfseries
No: 1\\
Name: sa\_\allowbreak{}9\_\allowbreak{}6\_\allowbreak{}01\\
Type: string\\
Default: \textquotedbl{}AES-256-GCM\textquotedbl{}
}

{\footnotesize Зберігається ексклюзивний контроль над криптографічними ключами для зашифрованих матеріалів, що зберігаються або передаються через зовнішню систему}




\subsubsection{ЗОВНІШНІ ПОСЛУГИ ДЛЯ СИСТЕМИ- ПЕРЕВІРКА ЦІЛІСНОСТІ, ЩО КОНТРОЛЮЄТЬСЯ ОРГАНІЗАЦІЄЮ (SA-9(7))}
Забезпечити можливість перевірки цілісності інформації в організації під час її перебування в зовнішній системі.

\vspace{10pt}
{\footnotesize\bfseries
No: 1\\
Name: sa\_\allowbreak{}9\_\allowbreak{}7\_\allowbreak{}01\\
Type: integer\\
Default: 30
}

{\footnotesize Передбачена можливість перевірки цілісності інформації під час її перебування у зовнішній системі}




\subsubsection{ЗОВНІШНІ ПОСЛУГИ ДЛЯ СИСТЕМИЗБЕРІГАННЯ -- ЮРИСДИКЦІЯ УКРАЇНИ (SA-9(8))}
Обмежити географічне розміщення обробки та зберігання даних об'єктами, розташованими в межах юридичної юрисдикції України.

\vspace{10pt}
{\footnotesize\bfseries
No: 1\\
Name: sa\_\allowbreak{}9\_\allowbreak{}8\_\allowbreak{}01\\
Type: string\\
Default: nil
}

{\footnotesize Обмежується географічне розміщення обробки інформації та зберігання даних об'єктами, розташованими в межах юридичної юрисдикції України}




\subsection{УПРАВЛІННЯ КОНФІГУРАЦІЄЮ РОЗРОБНИКА (SA-10)}
Вимагати від розробника системи, системного компонента або системної служби:\\ 
a. виконання управління конфігурацією під час [Вибір (один або кілька): проєктування; розробки; реалізації; експлуатації; видалення] системи, компонента або служби;\\ 
b. документувати, керувати та контролювати цілісність змін у [Призначення: визначених організацією елементах конфігурації при управлінні конфігурацією];\\ 
c. впроваджувати тільки схвалені організацією зміни в системі, компоненті або службі;\\ 
d. документувати зміни в системі, компоненті або службі та можливі наслідки таких змін для безпеки та приватності;\\ 
e. відстежувати недоліки безпеки та усунення дефектів у системі, компоненті або службі та повідомляти про результати [Призначення: визначений організацією персонал].

\vspace{10pt}
{\footnotesize\bfseries
No: 1\\
Name: sa\_\allowbreak{}10\_\allowbreak{}odp\_\allowbreak{}02\\
Type: string\\
Default: nil
}

{\footnotesize Визначено елементи конфігурації під керуванням}




\subsubsection{Управління конфігурацією розробника --- Перевірка цілісності програмного забезпечення та мікропрограм (SA-10(1))}


\vspace{10pt}
{\footnotesize\bfseries
No: 1\\
Name: sa\_\allowbreak{}10\_\allowbreak{}1\_\allowbreak{}01\\
Type: string\\
Default: nil
}

{\footnotesize Зобов'язаний розробник системи, системного компонента або системного служби забезпечити перевірку цілісності програмного забезпечення та компонентів програмного забезпечення}




\subsubsection{Управління конфігурацією розробника --- Довірене постачання (SA-10(6))}


\vspace{10pt}
{\footnotesize\bfseries
No: 1\\
Name: sa\_\allowbreak{}10\_\allowbreak{}6\_\allowbreak{}01\\
Type: string\\
Default: \textquotedbl{}автоматизований засіб моніторингу\textquotedbl{}
}

{\footnotesize Повинен розробник системи, системного компонента або системної служби виконувати процедури для забезпечення того, щоб апаратні засоби, програмне забезпечення й оновлення прошивки, що стосуються безпеки й оновлюються в організації, точно відповідали оригінальним копіям}




\subsection{УПРАВЛІННЯ КОНФІГУРАЦІЄЮ РОЗРОБНИКА - ПРЕДСТАВНИКИ БЕЗПЕКИ ТА ПРИВАТНОСТІ (SA-11)}
Вимагати від розробника системи, системного компонента або системної служби на всіх етапах проєктування та життєвого циклу розробки системи:\\ 
a. створити та впровадити план з оцінювання безпеки та приватності;\\ 
b. виконати [Вибір (один або кілька): одиниця; інтеграція; система; регресія] тестування/оцінювання [Призначення: з визначеною організацією частотою] з [Призначення: визначена організацією глибиною та охопленням];\\ 
c. надати докази (свідчення) виконання плану оцінювання та результати тестування й оцінювань;\\ 
d. впровадити перевірку процесу виправлення недоліків;\\ 
e. виправити дефекти, виявлені під час тестування та оцінювання.

\vspace{10pt}
{\footnotesize\bfseries
No: 1\\
Name: sa\_\allowbreak{}11\_\allowbreak{}odp\_\allowbreak{}01\\
Type: string\\
Default: nil
}

{\footnotesize Вибрано одне або декілька з наступних ЗНАЧЕНЬ ПАРАМЕТРІВ: \{одиниця; інтеграція; система; регресія\}}




\subsubsection{Тестування та оцінювання розробника --- Моделювання загроз і аналіз вразливостей (SA-11(2))}


\vspace{10pt}
{\footnotesize\bfseries
No: 1\\
Name: sa\_\allowbreak{}11\_\allowbreak{}2\_\allowbreak{}a\_\allowbreak{}01\\
Type: integer\\
Default: 30
}

{\footnotesize Повинен розробник системи, компонента системи або системного сервісу виконувати моделювання загроз під час розробки системи, компонента або сервісу, що використовує <SA-11(02)\_\allowbreak{}ODP[01]\_\allowbreak{}інформацію>}


{\footnotesize\bfseries
No: 2\\
Name: sa\_\allowbreak{}11\_\allowbreak{}2\_\allowbreak{}a\_\allowbreak{}02\\
Type: integer\\
Default: 30
}

{\footnotesize Повинен розробник системи, системного компонента або системної служби виконувати аналіз вразливостей під час розробки системи, компонента або служби, які використовують <SA-11(02)\_\allowbreak{}ODP[01]\_\allowbreak{}інформацію>}


{\footnotesize\bfseries
No: 3\\
Name: sa\_\allowbreak{}11\_\allowbreak{}2\_\allowbreak{}a\_\allowbreak{}03\\
Type: integer\\
Default: 30
}

{\footnotesize Повинен розробник системи, компонента системи або системного сервісу виконувати моделювання загроз під час подальшого тестування та оцінювання системи, компонента або сервісу, що використовує <SA- 11(02)\_\allowbreak{}ODP[01]\_\allowbreak{}інформацію>}


{\footnotesize\bfseries
No: 4\\
Name: sa\_\allowbreak{}11\_\allowbreak{}2\_\allowbreak{}a\_\allowbreak{}04\\
Type: integer\\
Default: 30
}

{\footnotesize Повинен розробник системи, системного компонента або системного сервісу виконувати аналіз вразливостей під час подальшого тестування та оцінювання системи, компонента або сервісу, що використовує <SA11(02)\_\allowbreak{}ODP[01]\_\allowbreak{}інформацію>}


{\footnotesize\bfseries
No: 5\\
Name: sa\_\allowbreak{}11\_\allowbreak{}2\_\allowbreak{}b\_\allowbreak{}01\\
Type: integer\\
Default: 30
}

{\footnotesize Повинен розробник системи, компонента системи або системного сервісу виконувати моделювання загроз під час розробки системи, компонента або сервісу, що використовує засоби та методи}


{\footnotesize\bfseries
No: 6\\
Name: sa\_\allowbreak{}11\_\allowbreak{}2\_\allowbreak{}b\_\allowbreak{}02\\
Type: integer\\
Default: 30
}

{\footnotesize Повинен розробник системи, компонента системи або системного сервісу виконувати моделювання загроз під час подальшого тестування та оцінювання системи, компонента або сервісу, що використовує інструменти та методи}


{\footnotesize\bfseries
No: 7\\
Name: sa\_\allowbreak{}11\_\allowbreak{}2\_\allowbreak{}b\_\allowbreak{}03\\
Type: integer\\
Default: 30
}

{\footnotesize Повинен розробник системи, компонента системи або системної служби виконувати аналіз вразливостей під час розробки системи, компонента або служби, яка використовує інструменти та методи}


{\footnotesize\bfseries
No: 8\\
Name: sa\_\allowbreak{}11\_\allowbreak{}2\_\allowbreak{}b\_\allowbreak{}04\\
Type: integer\\
Default: 30
}

{\footnotesize Повинен розробник системи, системного компонента або системного сервісу виконувати аналіз вразливостей під час подальшого тестування та оцінювання системи, компонента або сервісу, що використовує інструменти та методи}


{\footnotesize\bfseries
No: 9\\
Name: sa\_\allowbreak{}11\_\allowbreak{}2\_\allowbreak{}c\_\allowbreak{}01\\
Type: integer\\
Default: 30
}

{\footnotesize Повинен розробник системи, системного компонента або системної служби виконувати моделювання загроз на широті та глибині під час розробки системи, компонента або сервісу}


{\footnotesize\bfseries
No: 10\\
Name: sa\_\allowbreak{}11\_\allowbreak{}2\_\allowbreak{}c\_\allowbreak{}02\\
Type: integer\\
Default: 30
}

{\footnotesize Зобов'язаний розробник системи, компонента системи або системної служби виконувати аналіз вразливостей під час подальшого тестування та оцінювання системи, компонента або сервісу, який проводить моделювання та аналіз на ширину та глибину}


{\footnotesize\bfseries
No: 11\\
Name: sa\_\allowbreak{}11\_\allowbreak{}2\_\allowbreak{}d\_\allowbreak{}01\\
Type: integer\\
Default: 30
}

{\footnotesize Повинен розробник системи, компонента системи або системної служби виконувати моделювання загроз під час розробки системи, компонента або сервісу, що дає змогу отримати докази, які відповідають критеріям прийнятності}


{\footnotesize\bfseries
No: 12\\
Name: sa\_\allowbreak{}11\_\allowbreak{}2\_\allowbreak{}d\_\allowbreak{}02\\
Type: integer\\
Default: 30
}

{\footnotesize Повинен розробник системи, компонента системи або системної служби виконувати моделювання загроз під час подальшого тестування та оцінювання системи, компонента або сервісу, що дає змогу отримати докази, які відповідають критеріям прийнятності }


{\footnotesize\bfseries
No: 13\\
Name: sa\_\allowbreak{}11\_\allowbreak{}2\_\allowbreak{}d\_\allowbreak{}03\\
Type: integer\\
Default: 30
}

{\footnotesize Повинен розробник системи, системного компонента або системної служби виконувати аналіз вразливостей під час розробки системи, компонента або служби, який надає докази, що відповідають прийнятності>}


{\footnotesize\bfseries
No: 14\\
Name: sa\_\allowbreak{}11\_\allowbreak{}2\_\allowbreak{}d\_\allowbreak{}04\\
Type: integer\\
Default: 30
}

{\footnotesize Критеріям повинен розробник системи, системного компонента або системної служби виконувати аналіз вразливостей під час подальшого тестування та оцінювання системи, компонента або сервісу, який надає докази, що відповідають <SA11(02)\_\allowbreak{}ODP[06] критеріям прийнятності}




\subsubsection{ATE\_\allowbreak{}IND.2 EAL2 Незалежне тестування EAL3 Незалежне тестування --- EAL4 EAL5 EAL6 зразок (SA-11(3))}


\vspace{10pt}
{\footnotesize\bfseries
No: 1\\
Name: sa\_\allowbreak{}11\_\allowbreak{}3\_\allowbreak{}a\_\allowbreak{}01\\
Type: integer\\
Default: 30
}

{\footnotesize Потрібен незалежний агент, який відповідатиме критеріям незалежності для перевірки правильності виконання плану оцінки безпеки розробника та доказів, отриманих під час тестування та оцінки}


{\footnotesize\bfseries
No: 2\\
Name: sa\_\allowbreak{}11\_\allowbreak{}3\_\allowbreak{}a\_\allowbreak{}02\\
Type: integer\\
Default: 30
}

{\footnotesize Потрібен незалежний агент, який відповідатиме критеріям незалежності для перевірки правильності виконання плану оцінювання конфіденційності розробника та доказів, отриманих під час тестування та оцінювання}


{\footnotesize\bfseries
No: 3\\
Name: sa\_\allowbreak{}11\_\allowbreak{}3\_\allowbreak{}b\\
Type: string\\
Default: nil
}

{\footnotesize Надано незалежному агенту достатньо інформації для завершення процесу перевірки, чи надано йому повноваження для отримання такої інформації}


{\footnotesize\bfseries
No: 4\\
Name: sa\_\allowbreak{}11\_\allowbreak{}3\_\allowbreak{}odp\\
Type: list\\
Default: []
}

{\footnotesize Визначені критерії незалежності, відповідати незалежний агент; яким повинен}




\subsubsection{Тестування та оцінювання розробника --- Тестування на проникнення (SA-11(5))}


\vspace{10pt}
{\footnotesize\bfseries
No: 1\\
Name: sa\_\allowbreak{}11\_\allowbreak{}5\_\allowbreak{}a\_\allowbreak{}01\\
Type: string\\
Default: nil
}

{\footnotesize Зобов'язаний розробник системи, системного компонента або системної служби виконувати тестування на проникнення на наступному рівні суворості: ширина}


{\footnotesize\bfseries
No: 2\\
Name: sa\_\allowbreak{}11\_\allowbreak{}5\_\allowbreak{}a\_\allowbreak{}02\\
Type: string\\
Default: nil
}

{\footnotesize Зобов'язаний розробник системи, системного компонента або системної служби виконувати тестування на проникнення на наступному рівні суворості: глибина}


{\footnotesize\bfseries
No: 3\\
Name: sa\_\allowbreak{}11\_\allowbreak{}5\_\allowbreak{}b\\
Type: string\\
Default: nil
}

{\footnotesize Зобов'язаний розробник системи, системного компонента або системної служби проводити тестування на проникнення в умовах <SA-11(05)\_\allowbreak{}ODP[03]\_\allowbreak{}обмежень>}




\subsubsection{Покриття Аналіз покриття (SA-11(7))}


\vspace{10pt}
{\footnotesize\bfseries
No: 1\\
Name: sa\_\allowbreak{}11\_\allowbreak{}7\_\allowbreak{}01\\
Type: string\\
Default: nil
}

{\footnotesize Повинен розробник системи, системного компонента або системної служби перевіряти, що обсяг тестування та оцінювання забезпечує повне покриття необхідних засобів контролю на ширину}


{\footnotesize\bfseries
No: 2\\
Name: sa\_\allowbreak{}11\_\allowbreak{}7\_\allowbreak{}02\\
Type: string\\
Default: nil
}

{\footnotesize Повинен розробник системи, системного компонента або системної служби перевіряти, що обсяг тестування та оцінювання забезпечує повне покриття необхідних засобів контролю на глибину}




\subsubsection{ДИНАМІЧНИЙ АНАЛІЗ КОДУ (SA-11(8))}


\vspace{10pt}
{\footnotesize\bfseries
No: 1\\
Name: sa\_\allowbreak{}11\_\allowbreak{}8\_\allowbreak{}01\\
Type: string\\
Default: nil
}

{\footnotesize Повинен розробник системи, системного компонента або системної служби використовувати інструменти динамічного аналізу коду для виявлення недоліків помилок}


{\footnotesize\bfseries
No: 2\\
Name: sa\_\allowbreak{}11\_\allowbreak{}8\_\allowbreak{}02\\
Type: string\\
Default: nil
}

{\footnotesize Зобов'язаний розробник системи, системного компонента або системної служби документувати результати аналізу}




\subsection{КЕРУВАННЯ РИЗИКАМИ ЛАНЦЮГА ПОСТАЧАННЯ (SA-12) [Вилучено]}
[Вилучено: включено до SR]

\vspace{10pt}
\textit{Немає параметрів для цього контролю.}
\vspace{10pt}


\subsection{ДОВІРЧІСТЬ (SA-13) [Вилучено]}
[Вилучено: включено до SA-8]

\vspace{10pt}
\textit{Немає параметрів для цього контролю.}
\vspace{10pt}


\subsection{АНАЛІЗ КРИТИЧНОСТІ (SA-14) [Вилучено]}
[Вилучено: включено до RA-9]

\vspace{10pt}
\textit{Немає параметрів для цього контролю.}
\vspace{10pt}


\subsection{ПРОЦЕСИ, СТАНДАРТИ ТА ІНСТРУМЕНТИ РОЗРОБКИ (SA-15)}
a.\\ 
b. Вимагати від розробника системи, системного компонента або системної служби слідувати документованому процесу розробки, який:\\ 
1. явно відповідає вимогам безпеки та приватності;\\ 
2. визначає стандарти й інструменти, які використовуються в процесі розробки;\\ 
3. документує конкретні параметри та конфігурації інструментарію, що використовуються в процесі розробки;\\ 
4. документує, управляє та забезпечує цілісність змін у процесі та/або інструментах, які використовуються в процесі розробки. Ознайомитися з процесом розробки, стандартами, інструментами, параметрами інструментарію і конфігураціями інструментів [Призначення: визначена організацією частота], щоб визначити, чи можуть вибрані й використовувані процеси, стандарти, інструменти, параметри та конфігурації інструментів задовольнити [Призначення: визначені організацією вимоги до безпеки та приватності].

\vspace{10pt}
{\footnotesize\bfseries
No: 1\\
Name: sa\_\allowbreak{}15\_\allowbreak{}odp\_\allowbreak{}01\\
Type: string\\
Default: \textquotedbl{}щорічно\textquotedbl{}
}

{\footnotesize Визначено періодичність перегляду процесу розробки, стандартів, інструментів, опцій інструментів та конфігурацій інструментів}


{\footnotesize\bfseries
No: 2\\
Name: sa\_\allowbreak{}15\_\allowbreak{}odp\_\allowbreak{}02\\
Type: string\\
Default: nil
}

{\footnotesize Визначені вимоги до безпеки, яким має відповідати процес, стандарти, інструменти, опції інструментів та конфігурації інструментів}


{\footnotesize\bfseries
No: 3\\
Name: sa\_\allowbreak{}15\_\allowbreak{}odp\_\allowbreak{}03\\
Type: integer\\
Default: 30
}

{\footnotesize Визначені вимоги до конфіденційності, яким має відповідати процес, стандарти, інструменти, опції інструментів та конфігурації інструментів; SA-15a.01[01] зобов'язаний розробник системи, системного компонента або системної служби дотримуватися задокументованого процесу розробки, який явно враховує вимоги безпеки; SA-15a.01[02] зобов'язаний розробник системи, системного компонента або системної служби дотримуватися задокументованого процесу розробки, який чітко враховує вимоги щодо конфіденційності; SA-15a.02[01] повинен розробник системи, системного компонента або системної служби дотримуватися задокументованого процесу розробки, який визначає стандарти, що використовуються в процесі розробки; SA-15a.02[02] повинен розробник системи, системного компонента або системного служби дотримуватися задокументованого процесу розробки, який визначає інструменти, що використовуються в процесі розробки; SA-15a.03[01] овинен розробник системи, системного компонента або системної служби дотримуватися задокументованого процесу розробки, який документує конкретний інструмент, що використовується в процесі розробки; SA-15a.03[02] повинен розробник системи, системного компонента або системної служби дотримуватися задокументованого процесу розробки, який документує конкретні конфігурації інструментів, що використовуються в процесі розробки; SA-15a.04 повинен розробник системи, системного компонента або системної служби дотримуватися задокументованого процесу розробки, який документує, управляє та забезпечує цілісність змін у процесі та/або інструментах, що використовуються під час розробки; SA-15b.[01] повинен розробник системи, системного компонента або системної служби дотримуватися задокументованого процесу розробки, в якому процес розробки, стандарти, інструменти, опції інструментів та конфігурації інструментів переглядаються з частотою, щоб визначити, що процес, стандарти, інструменти, опції інструментів та конфігурації інструментів, вибрані та застосовані, задовольняють вимоги безпеки; SA-15b.[02] повинен розробник системи, системного компонента або системної служби дотримуватися задокументованого процесу розробки, в якому процес розробки, стандарти, інструменти, опції інструментів та конфігурації інструментів переглядаються з частотою, щоб визначити, що процес, стандарти, інструменти, опції інструментів та конфігурації інструментів, обрані та застосовані, задовольняють вимоги щодо конфіденційності}




\subsubsection{ЗМЕНШЕННЯ ПОВЕРХНІ АТАКИ (SA-15(5))}


\vspace{10pt}
{\footnotesize\bfseries
No: 1\\
Name: sa\_\allowbreak{}15\_\allowbreak{}5\_\allowbreak{}01\\
Type: string\\
Default: nil
}

{\footnotesize Зобов'язаний розробник системи, системного компонента або системної служби зменшити поверхню атаки до межі}


{\footnotesize\bfseries
No: 2\\
Name: sa\_\allowbreak{}15\_\allowbreak{}5\_\allowbreak{}odp\\
Type: string\\
Default: nil
}

{\footnotesize Визначені порогові значення, до яких необхідно зменшити поверхню атаки}




\subsubsection{ПОСТІЙНЕ ВДОСКОНАЛЕННЯ (SA-15(6))}


\vspace{10pt}
{\footnotesize\bfseries
No: 1\\
Name: sa\_\allowbreak{}15\_\allowbreak{}6\_\allowbreak{}01\\
Type: string\\
Default: nil
}

{\footnotesize Зобов'язаний розробник системи, системного компонента або системної служби впроваджувати чіткий процес постійного вдосконалення процесу розробки}




\subsubsection{АНАЛІЗ ВРАЗЛИВОСТЕЙ АВТОМАТИЗОВАНИЙ (SA-15(7))}


\vspace{10pt}
{\footnotesize\bfseries
No: 1\\
Name: sa\_\allowbreak{}15\_\allowbreak{}7\_\allowbreak{}a\\
Type: string\\
Default: \textquotedbl{}щорічно\textquotedbl{}
}

{\footnotesize Зобов'язаний розробник системи, системного компонента або системного сервісу виконувати автоматизований аналіз вразливостей частота з використанням інструментарію}


{\footnotesize\bfseries
No: 2\\
Name: sa\_\allowbreak{}15\_\allowbreak{}7\_\allowbreak{}b\\
Type: string\\
Default: \textquotedbl{}щорічно\textquotedbl{}
}

{\footnotesize Зобов'язаний розробник системи, системного компоненту або системної служби визначати потенціал використання виявлених вразливостей частота}


{\footnotesize\bfseries
No: 3\\
Name: sa\_\allowbreak{}15\_\allowbreak{}7\_\allowbreak{}c\\
Type: string\\
Default: \textquotedbl{}щорічно\textquotedbl{}
}

{\footnotesize Повинен розробник системи, системного компоненту або системної служби визначати потенційні заходи зменшення ризику частота для наданих вразливостей}


{\footnotesize\bfseries
No: 4\\
Name: sa\_\allowbreak{}15\_\allowbreak{}7\_\allowbreak{}d\\
Type: list\\
Default: [\textquotedbl{}admin\textquotedbl{}, \textquotedbl{}security\_\allowbreak{}officer\textquotedbl{}]
}

{\footnotesize Повинен розробник системи, системного компонента або системної послуги надавати вихідні дані інструментів і результати аналізу частота персоналу або ролям}




\subsubsection{ПОВТОРНЕ ВИКОРИСТАННЯ ІНФОРМАЦІЇ ПРО ЗАГРОЗИ ТА ВРАЗЛИВОСТІ (SA-15(8))}


\vspace{10pt}
{\footnotesize\bfseries
No: 1\\
Name: sa\_\allowbreak{}15\_\allowbreak{}8\_\allowbreak{}01\\
Type: string\\
Default: nil
}

{\footnotesize Повинен розробниксистеми, системного компонента або системної служби використовувати моделювання загроз з подібних систем, компонентів або служб для інформування про поточний процес розробки}


{\footnotesize\bfseries
No: 2\\
Name: sa\_\allowbreak{}15\_\allowbreak{}8\_\allowbreak{}02\\
Type: string\\
Default: nil
}

{\footnotesize Повинен розробник системи, системного компонента або системних служб використовувати аналіз вразливостей аналогічних систем, компонентів або служб для інформування поточного процесу розробки}




\subsubsection{ПРОЦЕС, СТАНДАРТИ ТА ІНСТРУМЕНТИ ВИКОРИСТАННЯ РЕАЛЬНИХ ДАНИХ (SA-15(9)) [Вилучено]}
[Вилучено: включено до SA-3(2)]

\vspace{10pt}
\textit{Немає параметрів для цього контролю.}
\vspace{10pt}


\subsection{НАВЧАННЯ, ЩО НАДАЄТЬСЯ РОЗРОБНИКАМИ (SA-16)}
Вимагати від розробника системи, системного компонента або системної служби забезпечити наступне навчання щодо правильного використання та функціонування реалізованих функцій, заходів і механізмів безпеки та приватності [Призначення: визначене організацією навчання].

\vspace{10pt}
{\footnotesize\bfseries
No: 1\\
Name: sa\_\allowbreak{}16\_\allowbreak{}01\\
Type: string\\
Default: nil
}

{\footnotesize Повинен розробник системи, системного компонента або системної служби проводити навчання щодо правильного використання та експлуатації впроваджених функцій, засобів управління та/або механізмів безпеки та приватності}


{\footnotesize\bfseries
No: 2\\
Name: sa\_\allowbreak{}16\_\allowbreak{}odp\\
Type: string\\
Default: nil
}

{\footnotesize Визначено навчання щодо правильного використання та експлуатації впроваджених функцій безпеки та приватності, засобів контролю та/або механізмів, що надаються розробником системи, системного компонента або системної служби}




\subsection{ПРОЄКТ ТА АРХІТЕКТУРА РОЗРОБНИКА (SA-17)}
Вимагати від розробника системи, системного компонента або системної служби створення специфікації проєкту та архітектури безпеки та приватності, яка:\\ 
a. узгоджується з архітектурою безпеки та приватності організації яка є невід'ємною частиною корпоративної архітектури організації;\\ 
b. точно та повністю описує необхідні функції безпеки та приватності, а також розподіл заходів захисту між фізичними та логічними компонентами;\\ 
c. пояснює, як разом працюють окремі функції, механізми та служби безпеки для забезпечення необхідних можливостей безпеки та єдиного підходу до захисту.

\vspace{10pt}
{\footnotesize\bfseries
No: 1\\
Name: sa\_\allowbreak{}17\_\allowbreak{}a\_\allowbreak{}01\\
Type: integer\\
Default: 30
}

{\footnotesize Повинен розробник системи, системного компонента або системної служби створювати специфікації проєкту та архітектури безпеки, які відповідають архітектурі безпеки організації, що є невід'ємною частиною архітектури підприємства організації}


{\footnotesize\bfseries
No: 2\\
Name: sa\_\allowbreak{}17\_\allowbreak{}a\_\allowbreak{}02\\
Type: integer\\
Default: 30
}

{\footnotesize Повинен розробник системи, системного компонента або системної служби створювати проєкту та архітектури безпеки поиватності, які відповідають архітектурі приватності організації, що є невід'ємною частиною корпоративної архітектури організації}


{\footnotesize\bfseries
No: 3\\
Name: sa\_\allowbreak{}17\_\allowbreak{}b\_\allowbreak{}01\\
Type: string\\
Default: nil
}

{\footnotesize Зобов'язаний розробник системи, системного компонента або системної служби підготувати специфікацію проєкту та архітектуру безпеки, які точно і повно описують необхідну функціональність безпеки та розподіл засобів контролю між фізичними та логічними компонентами}


{\footnotesize\bfseries
No: 4\\
Name: sa\_\allowbreak{}17\_\allowbreak{}b\_\allowbreak{}02\\
Type: string\\
Default: nil
}

{\footnotesize Зобов'язаний розробник системи, системного компонента або системної служби підготувати специфікацію проєкту та архітектуру приватності, які точно і повно описують необхідну функціональність приватності та розподіл засобів контролю між фізичними та логічними компонентами}


{\footnotesize\bfseries
No: 5\\
Name: sa\_\allowbreak{}17\_\allowbreak{}c\_\allowbreak{}01\\
Type: string\\
Default: \textquotedbl{}автоматизований засіб моніторингу\textquotedbl{}
}

{\footnotesize Повинен розробник системи, системного компонента або системного сервісу створювати проектну специфікацію та архітектуру безпеки, які описують, як окремі функції, механізми та сервіси безпеки працюють разом для забезпечення необхідних можливостей безпеки та єдиного підходу до захисту}


{\footnotesize\bfseries
No: 6\\
Name: sa\_\allowbreak{}17\_\allowbreak{}c\_\allowbreak{}02\\
Type: string\\
Default: \textquotedbl{}автоматизований засіб моніторингу\textquotedbl{}
}

{\footnotesize Повинен розробник системи, системного компонента або системної служби створювати специфікацію проєкту та архітектуру приватності, які описують, як окремі функції, механізми та служби конфіденційності працюють разом для забезпечення необхідних можливостей приватності та єдиного підходу до захисту}




\subsubsection{Проєкт і архітектура безпеки розробника --- Формальна модель політики (SA-17(1))}


\vspace{10pt}
{\footnotesize\bfseries
No: 1\\
Name: sa\_\allowbreak{}17\_\allowbreak{}1\_\allowbreak{}a\_\allowbreak{}02\\
Type: list\\
Default: [\textquotedbl{}default\_\allowbreak{}deny\_\allowbreak{}rule\textquotedbl{}, \textquotedbl{}abac\_\allowbreak{}rule\_\allowbreak{}1\textquotedbl{}]
}

{\footnotesize Повинен розробник системи, системного компонента або системної служби, як невід'ємну частину процесу розробки, створювати формальну модель політики, що описує організаційну політику конфіденційності, яка підлягає виконанню}


{\footnotesize\bfseries
No: 2\\
Name: sa\_\allowbreak{}17\_\allowbreak{}1\_\allowbreak{}b\_\allowbreak{}01\\
Type: list\\
Default: [\textquotedbl{}default\_\allowbreak{}deny\_\allowbreak{}rule\textquotedbl{}, \textquotedbl{}abac\_\allowbreak{}rule\_\allowbreak{}1\textquotedbl{}]
}

{\footnotesize Повинен розробник системи, системного компонента або системної служби доводити, що формальна модель політики є внутрішньо узгодженою і достатньою для забезпечення дотримання визначених елементів політики безпеки організації при її впровадженні}


{\footnotesize\bfseries
No: 3\\
Name: sa\_\allowbreak{}17\_\allowbreak{}1\_\allowbreak{}b\_\allowbreak{}02\\
Type: list\\
Default: [\textquotedbl{}default\_\allowbreak{}deny\_\allowbreak{}rule\textquotedbl{}, \textquotedbl{}abac\_\allowbreak{}rule\_\allowbreak{}1\textquotedbl{}]
}

{\footnotesize Повинен розробник системи, системного компонента або системної служби доводити, що формальна модель політики є внутрішньо узгодженою і достатньою для забезпечення дотримання визначених елементів організаційної політики приватності при її впровадженні}




\subsubsection{Проєкт і архітектура безпеки розробника --- Компоненти, що необхідні для забезпечення безпеки (SA-17(2))}


\vspace{10pt}
{\footnotesize\bfseries
No: 1\\
Name: sa\_\allowbreak{}17\_\allowbreak{}2\_\allowbreak{}a\_\allowbreak{}01\\
Type: string\\
Default: nil
}

{\footnotesize Зобов'язаний розробник системи, системного компонента або системної служби визначати апаратне забезпечення, що має відношення до безпеки}


{\footnotesize\bfseries
No: 2\\
Name: sa\_\allowbreak{}17\_\allowbreak{}2\_\allowbreak{}a\_\allowbreak{}02\\
Type: string\\
Default: nil
}

{\footnotesize Зобов'язаний розробник системи, системного компонента або системної служби визначати програмне забезпечення, що має відношення до безпеки; SA-17(02)(a)[03] зобов'язаний розробник системи, системного компонента або системної служби визначати мікропрограмне забезпечення, що мають відношення до безпеки}


{\footnotesize\bfseries
No: 3\\
Name: sa\_\allowbreak{}17\_\allowbreak{}2\_\allowbreak{}b\\
Type: string\\
Default: nil
}

{\footnotesize Повинен розробник системи, системного компонента або системної служби надавати обґрунтування того, що визначення обладнання, програмного забезпечення та мікропрограмного забезпечення, що мають відношення до безпеки, є повним}




\subsubsection{Проєкт і архітектура безпеки розробника --- Формальна відповідність (SA-17(3))}


\vspace{10pt}
{\footnotesize\bfseries
No: 1\\
Name: sa\_\allowbreak{}17\_\allowbreak{}3\_\allowbreak{}a\_\allowbreak{}01\\
Type: string\\
Default: nil
}

{\footnotesize Зобов'язаний розробник системи, системного компонента або системної служби визначати апаратне забезпечення, що має відношення до безпеки}


{\footnotesize\bfseries
No: 2\\
Name: sa\_\allowbreak{}17\_\allowbreak{}3\_\allowbreak{}a\_\allowbreak{}02\\
Type: string\\
Default: nil
}

{\footnotesize Зобов'язаний розробник системи, системного компонента або системної служби визначати програмне забезпечення, що має відношення до безпеки}


{\footnotesize\bfseries
No: 3\\
Name: sa\_\allowbreak{}17\_\allowbreak{}3\_\allowbreak{}a\_\allowbreak{}03\\
Type: string\\
Default: nil
}

{\footnotesize Зобов'язаний розробник системи, системного компонента або системної служби визначати мікропрограми, що мають відношення до безпеки}


{\footnotesize\bfseries
No: 4\\
Name: sa\_\allowbreak{}17\_\allowbreak{}3\_\allowbreak{}b\\
Type: string\\
Default: nil
}

{\footnotesize Повинен розробник системи, системного компонента або системної служби надавати обґрунтування того, що визначення обладнання, програмного забезпечення та мікропрограмного забезпечення, що мають відношення до безпеки, є повним}


{\footnotesize\bfseries
No: 5\\
Name: sa\_\allowbreak{}17\_\allowbreak{}3\_\allowbreak{}c\\
Type: string\\
Default: nil
}

{\footnotesize Повинен розробник системи, системного компонента або системної служби демонструвати за допомогою неформальної демонстрації, що формальна специфікація верхнього рівня повністю охоплює інтерфейси до обладнання, програмного забезпечення та мікропрограмного забезпечення, що мають відношення до безпеки}


{\footnotesize\bfseries
No: 6\\
Name: sa\_\allowbreak{}17\_\allowbreak{}3\_\allowbreak{}d\\
Type: string\\
Default: nil
}

{\footnotesize Повинен розробник системи, системного компонента або системної служби доводити, що формальна специфікація верхнього рівня є точним описом впровадженого обладнання, програмного забезпечення та мікропрограмного забезпечення, що мають відношення до безпеки}


{\footnotesize\bfseries
No: 7\\
Name: sa\_\allowbreak{}17\_\allowbreak{}3\_\allowbreak{}e\\
Type: string\\
Default: \textquotedbl{}автоматизований засіб моніторингу\textquotedbl{}
}

{\footnotesize Повинен розробник системи, системного компонента або системної служби описувати релевантні для безпеки апаратні, програмні та мікропрограмні механізми, які не розглядаються у формальній специфікації верхнього рівня, але є суто внутрішніми для релевантних для безпеки апаратних, програмних та мікропрограмних засобів}




\subsubsection{Проєкт і архітектура безпеки розробника --- Неформальна відповідність (SA-17(4))}


\vspace{10pt}
{\footnotesize\bfseries
No: 1\\
Name: sa\_\allowbreak{}17\_\allowbreak{}4\_\allowbreak{}a\_\allowbreak{}01\\
Type: integer\\
Default: 30
}

{\footnotesize Повинен розробник системи, системного компонента або системної служби, як невід'ємну частину процесу розробки, створювати неформальну описову специфікацію верхнього рівня, яка визначає інтерфейси до релевантного для безпеки апаратного, програмного та мікропрограмного забезпечення з точки зору винятків}


{\footnotesize\bfseries
No: 2\\
Name: sa\_\allowbreak{}17\_\allowbreak{}4\_\allowbreak{}a\_\allowbreak{}02\\
Type: integer\\
Default: 30
}

{\footnotesize Повинен розробник системи, системного компонента або системної служби, як невід'ємну частину процесу розробки, створювати неформальну описову специфікацію верхнього рівня, яка визначає інтерфейси до релевантного для безпеки апаратного, програмного та мікропрограмного забезпечення в термінах повідомлень про помилки}


{\footnotesize\bfseries
No: 3\\
Name: sa\_\allowbreak{}17\_\allowbreak{}4\_\allowbreak{}a\_\allowbreak{}03\\
Type: integer\\
Default: 30
}

{\footnotesize Повинен розробник системи, системного компоненту або системної служби, як невід'ємну частину процесу розробки, створювати неформальну описову специфікацію верхнього рівня, яка визначає інтерфейси до релевантного для безпеки обладнання, програмного забезпечення та мікропрограмного забезпечення з точки зору наслідків}


{\footnotesize\bfseries
No: 4\\
Name: sa\_\allowbreak{}17\_\allowbreak{}4\_\allowbreak{}c\\
Type: string\\
Default: nil
}

{\footnotesize Повинен розробник системи, системного компонента або системної служби демонструвати за допомогою неформальної демонстрації, що описова специфікація верхнього рівня повністю охоплює інтерфейси до апаратного, програмного та мікропрограмного забезпечення, що мають відношення до безпеки}


{\footnotesize\bfseries
No: 5\\
Name: sa\_\allowbreak{}17\_\allowbreak{}4\_\allowbreak{}d\\
Type: string\\
Default: nil
}

{\footnotesize Повинен розробник системи, системного компонента або системної служби показувати, що описова специфікація верхнього рівня є точним описом інтерфейсів до релевантного для безпеки апаратного, програмного та мікропрограмного забезпечення}


{\footnotesize\bfseries
No: 6\\
Name: sa\_\allowbreak{}17\_\allowbreak{}4\_\allowbreak{}e\\
Type: string\\
Default: \textquotedbl{}автоматизований засіб моніторингу\textquotedbl{}
}

{\footnotesize Повинен розробник системи, системного компонента або системної служби описувати релевантні для безпеки апаратні, програмні та мікропрограмні механізми, які не розглядаються в описовій специфікації верхнього рівня, але є суто внутрішніми для релевантних для безпеки апаратних, програмних та мікропрограмних засобів}


{\footnotesize\bfseries
No: 7\\
Name: sa\_\allowbreak{}17\_\allowbreak{}4\_\allowbreak{}odp\\
Type: string\\
Default: nil
}

{\footnotesize Вибрано одне з наступних ЗНАЧЕНЬ ПАРАМЕТРА: \{неформальна описова специфікація, переконливий аргумент за допомогою формальних методів як можлива\}}




\subsection{ЗАХИСТ ТА ВИЯВЛЕННЯ ПІДРОБКИ (SA-18) [Вилучено]}
[Вилучено: включено до SR-9]

\vspace{10pt}
\textit{Немає параметрів для цього контролю.}
\vspace{10pt}


\subsection{СПРАВЖНІСТЬ КОМПОНЕНТА (SA-19) [Вилучено]}
[Вилучено: включено до SR-11]

\vspace{10pt}
\textit{Немає параметрів для цього контролю.}
\vspace{10pt}


\subsection{ІНДИВІДУАЛЬНА РОЗРОБКА КРИТИЧНИХ КОМПОНЕНТІВ (SA-20)}
Повторно реалізувати або налаштувати [Призначення: визначені організацією критичні компоненти системи].

\vspace{10pt}
{\footnotesize\bfseries
No: 1\\
Name: sa\_\allowbreak{}20\_\allowbreak{}01\\
Type: string\\
Default: nil
}

{\footnotesize Критична система повторно реалізувати або налаштувати на замовлення}


{\footnotesize\bfseries
No: 2\\
Name: sa\_\allowbreak{}20\_\allowbreak{}odp\\
Type: string\\
Default: nil
}

{\footnotesize Потрібно повторно реалізувати або налаштувати на замовлення критичні компоненти системи}




\subsection{ПЕРЕВІРКА РОЗРОБНИКА (SA-21)}
Вимагати, щоб розробник [Призначення: визначених організацією системи, компонент системи або послуги]:\\ 
a. мав відповідні дозволи доступу, як визначено призначеним [Призначення: визначеним організацією уповноваженим органом];\\ 
b. відповідає таким додатковим критеріям перевірки персоналу: [Призначення: визначені організацією додаткові критерії перевірки персоналу].

\vspace{10pt}
{\footnotesize\bfseries
No: 1\\
Name: sa\_\allowbreak{}21\_\allowbreak{}odp\_\allowbreak{}01\\
Type: string\\
Default: nil
}

{\footnotesize Визначена система, компонент системи або системна служба, до яких має доступ розробник}


{\footnotesize\bfseries
No: 2\\
Name: sa\_\allowbreak{}21\_\allowbreak{}odp\_\allowbreak{}02\\
Type: string\\
Default: nil
}

{\footnotesize Визначені офіційні обов'язки, покладені на розробника}


{\footnotesize\bfseries
No: 3\\
Name: sa\_\allowbreak{}21\_\allowbreak{}odp\_\allowbreak{}03\\
Type: list\\
Default: [\textquotedbl{}admin\textquotedbl{}, \textquotedbl{}security\_\allowbreak{}officer\textquotedbl{}]
}

{\footnotesize Визначені додаткові розробників; SA-21a. повинен розробник системи, системного компонента або системної служби мати відповідні повноваження доступу, як визначено призначеними уповноваженим органом організації; SA-21b. повинен розробник системи, системного компонента або системної служби відповідати додатковим критеріям перевірки персоналу. критерії перевірки персоналу}




\subsection{КОМПОНЕНТИ СИСТЕМИ, ЩО НЕ ПІДТРИМУЮТЬСЯ (SA-22)}
a. Замінювати компоненти системи, якщо підтримка компонентів більше не доступна розробнику, постачальнику або виробнику.\\ 
b. Надавати такі варіанти альтернативних джерел для подальшої підтримки непідтримуваних компонентів [Вибір (один або більше): внутрішня підтримка; [Призначення: підтримка, визначена організацією від зовнішніх постачальників]].

\vspace{10pt}
{\footnotesize\bfseries
No: 1\\
Name: sa\_\allowbreak{}22\_\allowbreak{}odp\_\allowbreak{}01\\
Type: string\\
Default: nil
}

{\footnotesize Вибрано одне або декілька з наступних ЗНАЧЕНЬ ПАРАМЕТРІВ: \{внутрішня підтримка; підтримка від зовнішніх постачальників\}}




\subsection{СПЕЦІАЛІЗАЦІЯ (SA-23)}
Покращення [Вибір (один або кілька): проектування; модифікація; збільшення; реконфігурація] на [Призначення: системи або системні компоненти, визначені організацією], які підтримують важливі служби або функції для підвищення надійності цих систем або компонентів.

\vspace{10pt}
{\footnotesize\bfseries
No: 1\\
Name: sa\_\allowbreak{}23\_\allowbreak{}odp\_\allowbreak{}01\\
Type: string\\
Default: nil
}

{\footnotesize Вибрано одне або декілька з наступних ЗНАЧЕНЬ ПАРАМЕТРІВ: \{модифікація дизайну; доповнення; реконфігурація\}}


{\footnotesize\bfseries
No: 2\\
Name: sa\_\allowbreak{}23\_\allowbreak{}odp\_\allowbreak{}02\\
Type: string\\
Default: nil
}

{\footnotesize Визначені системи або компоненти системи, підтримують важливі для місії послуги або функції}




\section{SC}
Клас заходів захисту SC --- ЗАХИСТ ІНФОРМАЦІЙНОЇ СИСТЕМИ ТА

Цей клас охоплює механізми захисту інформації під час її передачі мережами зв'язку, криптографічний захист та ізоляцію критичних компонентів.

\textbf{Перелік заходів захисту:} Політика та процедури захисту системи та комунікацій (SC-1); Розділення функцій (SC-2); Інтерфейси для непривілейованих користувачів (SC-2(1)); Відокремлення (SC-2(2)); Ізоляція функцій безпеки (SC-3); Ізоляція функцій забезпечення (SC-3(1)); Функції управління доступом та потоком (SC-3(2)); Мінімізація функціональності (SC-3(3)); З'єднання модулів зв'язність (SC-3(4)); Багаторівнева структура (SC-3(5)); Інформація в загальних системних ресурсах (SC-4); Рівні безпеки (SC-4(1)) [Вилучено]; Інформація в загальних системних багаторівнева або періодична обробка (SC-4(2)); Захист від атак «відмова в обслуговуванні» (SC-5); Обмеження внутрішніх користувачів (SC-5(1)); Захист від атак «відмова в обслуговуванні» продуктивність, пропускна здатність та надмірність (SC-5(2)); Виявлення та моніторинг (SC-5(3)); Доступність ресурсів (SC-6); Доступність ресурсів (SC-7); Фізично відділені підмережі (SC-7(1)) [Вилучено]; Публічний доступ (SC-7(2)) [Вилучено]; Точки доступу (SC-7(3)); Зовнішні комунікаційні служби (SC-7(4)); Відмова за замовчуванням - дозвіл за винятком (SC-7(5)); Відповідь на розпізнані помилки (SC-7(6)) [Вилучено]; Запобігання поділу тунелювання для віддалених пристроїв (SC-7(7)); Маршрутизація трафіку з автентифікованих проксі-серверів (SC-7(8)); Захист (SC-7(9)); Ізольовані від інших внутрішніх компонентів системи шляхом впровадження фізично відокремлених підмереж з керованими інтерфейсами до інших компонентів системи (SC-7(13)); Конфіденційність та цілісність передачі (SC-8); Конфіденційність та криптографічний захист (SC-8(1)); Попередня і постобробка (SC-8(2)); Конфіденційність та цілісність криптографічний захист повідомлень (SC-8(3)); Приховування або рандомізація комунікації (SC-8(4)); Конфіденційність та система розподілу (SC-8(5)); Конфіденційність передачі (SC-9) [Вилучено]; Відключення мережі (SC-10); Довірений канал зв'язку (SC-11); Встановлення ключами (SC-12); Криптографічний захист (SC-13); Захист громадського доступу (SC-14) [Вилучено]; Спільні обчислювальні пристрої та застосунки (SC-15); Передача атрибутів безпеки та приватності (SC-16); Сертифікати інфраструктури відкритих ключів (SC-17); Мобільний код (SC-18); Інтернет-протокол голосового зв'язку (SC-19) [Вилучено]; БЕЗПЕЧНА СЛУЖБА ІМЕН, АДРЕС (УПОВНОВАЖЕНЕ ДЖЕРЕЛО) (SC-20); БЕЗПЕЧНА СЛУЖБА ІМЕН, АДРЕС (УПОВНОВАЖЕНЕ ДЖЕРЕЛО) ДЖЕРЕЛО ДАНИХ ТА ЦІЛІСНІСТЬ (SC-21); Архітектура та забезпечення служби імен/адрес (SC-22); Автентифікація сесії (SC-23); Уведення у відомий стан (SC-24); Тонкі вузли (SC-25); ПРИМАНКА ДЛЯ ЗЛОВМИСНИКІВ (DECOYS) (SC-26); Незалежні від платформи застосунки (SC-27); Захист інформації в стані спокою (SC-28); ЗАХИСТ ІНФОРМАЦІЇ В СТАНІ СПОКОЮ | КРИПТОГРАФІЧНИЙ ЗАХИСТ; Гетерогенність (SC-29); Маскування та хибний напрям (SC-30); Аналіз прихованого каналу (SC-31); Поділ системи на частини (SC-32); Підготовка цілісності передачі (SC-33) [Вилучено]; Незмінювані виконавчі програми (SC-34); РОЗПІЗНАВАННЯ ПРИМАНОК ДЛЯ ЗЛОВМИСНИКІВ (HONEYCLIENT) (SC-35); Розподілена обробка та зберігання (SC-36); Позасмугові канали (SC-37); Безпека операцій (SC-38); Ізоляція процесу (SC-39); Захист бездротового з'єднання (SC-40); Доступ до портів та пристроїв введення, виведення (SC-41); Захист бездротового з'єднання доступ до портів і пристроїв введення/виведення можливості датчика та дані (SC-42); Обмеження використання (SC-43); Екрановані камери (SC-44); Синхронізація системи з часом (SC-45); Забезпечення виконання міждоменної політики (SC-46); Альтернативний шлях зв'язку (SC-47); Переміщення датчика (SC-48); Динамічно переміщуються до за (SC-48(1)); Примусове апаратне забезпечення виконання (SC-49); Примусове програмне забезпечення виконання (SC-50); Апаратний захист (SC-51).

\subsection{ПОЛІТИКА ТА ПРОЦЕДУРИ ЗАХИСТУ СИСТЕМИ ТА КОМУНІКАЦІЙ (SC-1)}
a. Розробити, задокументувати та поширити серед [Призначення: визначеного організацією персоналу або посадових осіб]:\\ 
1. 2. Політику захисту системи та комунікацій, яка:\\ 
(a) містить мету, сферу застосування, ролі, обов'язки, відповідальність керівництва, координацію між організаційними підрозділами та систему контролю відповідності (complaince);\\ 
(b) відповідає чинному законодавству, виконавчим наказам, директивам, нормам, політикам, стандартам та керівним принципам. Процедури для сприяння впровадженню політики в області захисту систем і комунікацій, а також пов'язаних з ними систем і засобів захисту зв'язку.\\ 
b. Призначити [Призначення: визначена організацією посадову особу] для управління політикою та процедурами захисту системи та комунікацій.\\ 
c. Переглядати та оновлювати:\\ 
1. поточну політику захисту системи та комунікацій [Призначення: визначена організацією частота];\\ 
2. поточні процедури захисту системи та комунікацій [Призначення: визначена організацією частота].

\vspace{10pt}
{\footnotesize\bfseries
No: 1\\
Name: sc\_\allowbreak{}1\_\allowbreak{}odp\_\allowbreak{}01\\
Type: list\\
Default: [\textquotedbl{}admin\textquotedbl{}, \textquotedbl{}security\_\allowbreak{}officer\textquotedbl{}]
}

{\footnotesize Визначено персонал або ролі, до яких має бути доведена політика захисту системи та комунікацій}


{\footnotesize\bfseries
No: 2\\
Name: sc\_\allowbreak{}1\_\allowbreak{}odp\_\allowbreak{}02\\
Type: list\\
Default: [\textquotedbl{}admin\textquotedbl{}, \textquotedbl{}security\_\allowbreak{}officer\textquotedbl{}]
}

{\footnotesize Визначено персонал або ролі, на які поширюються процедури захисту системи та комунікацій}


{\footnotesize\bfseries
No: 3\\
Name: sc\_\allowbreak{}1\_\allowbreak{}odp\_\allowbreak{}03\\
Type: string\\
Default: nil
}

{\footnotesize Вибрано жодне або декілька з наступних ЗНАЧЕНЬ ПАРАМЕТРІВ: \{рівень організації; рівень місії/бізнеспроцесу; рівень системи\}}


{\footnotesize\bfseries
No: 4\\
Name: sc\_\allowbreak{}1\_\allowbreak{}odp\_\allowbreak{}04\\
Type: list\\
Default: [\textquotedbl{}admin\textquotedbl{}, \textquotedbl{}security\_\allowbreak{}officer\textquotedbl{}]
}

{\footnotesize Визначено посадову особу, яка керуватиме політикою та процедурами захисту системи та комунікацій}


{\footnotesize\bfseries
No: 5\\
Name: sc\_\allowbreak{}1\_\allowbreak{}odp\_\allowbreak{}05\\
Type: list\\
Default: [\textquotedbl{}default\_\allowbreak{}deny\_\allowbreak{}rule\textquotedbl{}, \textquotedbl{}abac\_\allowbreak{}rule\_\allowbreak{}1\textquotedbl{}]
}

{\footnotesize Визначена періодичність перегляду та оновлення поточної політики захисту системи та комунікацій}


{\footnotesize\bfseries
No: 6\\
Name: sc\_\allowbreak{}1\_\allowbreak{}odp\_\allowbreak{}06\\
Type: list\\
Default: [\textquotedbl{}default\_\allowbreak{}deny\_\allowbreak{}rule\textquotedbl{}, \textquotedbl{}abac\_\allowbreak{}rule\_\allowbreak{}1\textquotedbl{}]
}

{\footnotesize Є події, які вимагають перегляду та оновлення поточної політики захисту системи та комунікацій}


{\footnotesize\bfseries
No: 7\\
Name: sc\_\allowbreak{}1\_\allowbreak{}odp\_\allowbreak{}07\\
Type: string\\
Default: \textquotedbl{}щорічно\textquotedbl{}
}

{\footnotesize Визначена періодичність перегляду та оновлення поточних процедур захисту системи та засобів зв'язку}




\subsection{РОЗДІЛЕННЯ ФУНКЦІЙ (SC-2)}
Розділяти функціональність користувача, включно зі службами, що призначені для користувача інтерфейсу, від функціональності системного управління.

\vspace{10pt}
{\footnotesize\bfseries
No: 1\\
Name: sc\_\allowbreak{}2\_\allowbreak{}01\\
Type: string\\
Default: nil
}

{\footnotesize Розділена функціональність користувача, включаючи сервіси користувацького інтерфейсу, від функціональності управління системою. ПОТЕНЦІЙНІ МЕТОДИ ТА ОБ'ЄКТИ ОЦІНКИ:}




\subsubsection{ІНТЕРФЕЙСИ ДЛЯ НЕПРИВІЛЕЙОВАНИХ КОРИСТУВАЧІВ (SC-2(1))}


\vspace{10pt}
{\footnotesize\bfseries
No: 1\\
Name: sc\_\allowbreak{}2\_\allowbreak{}1\_\allowbreak{}01\\
Type: string\\
Default: nil
}

{\footnotesize Запобігається представлення функціональності управління системою в інтерфейсі непривілейованим користувачам}




\subsubsection{ВІДОКРЕМЛЕННЯ (SC-2(2))}


\vspace{10pt}
{\footnotesize\bfseries
No: 1\\
Name: sc\_\allowbreak{}2\_\allowbreak{}2\_\allowbreak{}01\\
Type: string\\
Default: nil
}

{\footnotesize Зберігається інформація окремо від додатків та програмного забезпечення}




\subsection{ІЗОЛЯЦІЯ ФУНКЦІЙ БЕЗПЕКИ (SC-3)}
Ізолювати функції безпеки від інших функцій.

\vspace{10pt}
{\footnotesize\bfseries
No: 1\\
Name: sc\_\allowbreak{}3\_\allowbreak{}01\\
Type: string\\
Default: nil
}

{\footnotesize Ізольовані функції безпеки від інших функцій}




\subsubsection{ІЗОЛЯЦІЯ ФУНКЦІЙ ЗАБЕЗПЕЧЕННЯ (SC-3(1))}


\vspace{10pt}
{\footnotesize\bfseries
No: 1\\
Name: sc\_\allowbreak{}3\_\allowbreak{}1\_\allowbreak{}01\\
Type: string\\
Default: \textquotedbl{}автоматизований засіб моніторингу\textquotedbl{}
}

{\footnotesize Застосовуються механізми розділення апаратних засобів для реалізації ізоляції функцій безпеки}




\subsubsection{ФУНКЦІЇ УПРАВЛІННЯ ДОСТУПОМ ТА ПОТОКОМ (SC-3(2))}


\vspace{10pt}
{\footnotesize\bfseries
No: 1\\
Name: sc\_\allowbreak{}3\_\allowbreak{}2\_\allowbreak{}01\\
Type: string\\
Default: nil
}

{\footnotesize Ізольовані функції безпеки, що забезпечують управління доступом, які не пов'язані з безпекою}


{\footnotesize\bfseries
No: 2\\
Name: sc\_\allowbreak{}3\_\allowbreak{}2\_\allowbreak{}02\\
Type: string\\
Default: nil
}

{\footnotesize Ізольовані функції безпеки, що забезпечують контроль доступу, від інших функцій безпеки}


{\footnotesize\bfseries
No: 3\\
Name: sc\_\allowbreak{}3\_\allowbreak{}2\_\allowbreak{}03\\
Type: string\\
Default: nil
}

{\footnotesize Ізольовані функції безпеки, які не пов'язані з безпекою забезпечують контроль інформаційних потоків}


{\footnotesize\bfseries
No: 4\\
Name: sc\_\allowbreak{}3\_\allowbreak{}2\_\allowbreak{}04\\
Type: string\\
Default: nil
}

{\footnotesize Ізольовані функції безпеки, що забезпечують контроль інформаційних потоків, від інших функцій безпеки}




\subsubsection{МІНІМІЗАЦІЯ ФУНКЦІОНАЛЬНОСТІ (SC-3(3))}


\vspace{10pt}
{\footnotesize\bfseries
No: 1\\
Name: sc\_\allowbreak{}3\_\allowbreak{}3\_\allowbreak{}01\\
Type: integer\\
Default: 3
}

{\footnotesize Мінімізовано кількість функцій, не пов'язаних з безпекою, що входять до сфери ізоляції, яка містить функції безпеки}




\subsubsection{З'ЄДНАННЯ МОДУЛІВ ЗВ'ЯЗНІСТЬ (SC-3(4))}


\vspace{10pt}
{\footnotesize\bfseries
No: 1\\
Name: sc\_\allowbreak{}3\_\allowbreak{}4\_\allowbreak{}02\\
Type: string\\
Default: nil
}

{\footnotesize Реалізовані функції безпеки як значною мірою незалежні модулі, які мінімізують зв'язок між модулями}




\subsubsection{БАГАТОРІВНЕВА СТРУКТУРА (SC-3(5))}


\vspace{10pt}
{\footnotesize\bfseries
No: 1\\
Name: sc\_\allowbreak{}3\_\allowbreak{}5\_\allowbreak{}01\\
Type: string\\
Default: nil
}

{\footnotesize Реалізовані функції безпеки як багаторівнева структура, що мінімізує взаємодію між шарами дизайну та уникає будь-якої залежності нижчих шарів від функціональності або коректності вищих шарів}




\subsection{ІНФОРМАЦІЯ В ЗАГАЛЬНИХ СИСТЕМНИХ РЕСУРСАХ (SC-4)}
Запобігати несанкціонованій та ненавмисній передачі інформації через спільні системні ресурси.

\vspace{10pt}
{\footnotesize\bfseries
No: 1\\
Name: sc\_\allowbreak{}4\_\allowbreak{}01\\
Type: string\\
Default: nil
}

{\footnotesize Запобігається несанкціонована передача інформації через спільні системні ресурси}


{\footnotesize\bfseries
No: 2\\
Name: sc\_\allowbreak{}4\_\allowbreak{}02\\
Type: string\\
Default: nil
}

{\footnotesize Запобігається ненавмисна передача інформації через спільні системні ресурси. ПОТЕНЦІЙНІ МЕТОДИ ТА ОБ'ЄКТИ ОЦІНКИ:}




\subsubsection{РІВНІ БЕЗПЕКИ (SC-4(1)) [Вилучено]}
[Вилучено: включено до SC-4]

\vspace{10pt}
\textit{Немає параметрів для цього контролю.}
\vspace{10pt}


\subsubsection{ІНФОРМАЦІЯ В ЗАГАЛЬНИХ СИСТЕМНИХ БАГАТОРІВНЕВА АБО ПЕРІОДИЧНА ОБРОБКА (SC-4(2))}


\vspace{10pt}
{\footnotesize\bfseries
No: 1\\
Name: sc\_\allowbreak{}4\_\allowbreak{}2\_\allowbreak{}01\\
Type: string\\
Default: nil
}

{\footnotesize Запобігається несанкціонована передача інформації через спільні ресурси відповідно до процедур, коли системна обробка явно перемикається між різними рівнями класифікації інформації або категоріями безпеки}


{\footnotesize\bfseries
No: 2\\
Name: sc\_\allowbreak{}4\_\allowbreak{}2\_\allowbreak{}odp\\
Type: string\\
Default: nil
}

{\footnotesize Визначені процедури для запобігання несанкціонованій передачі інформації через спільні ресурси}




\subsection{ЗАХИСТ ВІД АТАК «ВІДМОВА В ОБСЛУГОВУВАННІ» (SC-5)}
a. [Призначення: захистити від; Обмежити] наслідки наступних типів подій відмови в обслуговуванні (DoS): [Призначення: визначені організацією типи подій відмови в обслуговуванні];\\ 
b. Застосувати наступні заходи захисту для досягнення мети відмови обслуговування [Призначення: заходи захисту визначені організацією, за типом події відмови в обслуговуванні].

\vspace{10pt}
{\footnotesize\bfseries
No: 1\\
Name: sc\_\allowbreak{}5\_\allowbreak{}odp\_\allowbreak{}01\\
Type: string\\
Default: nil
}

{\footnotesize Визначені типи подій відмов в обслуговуванні, від яких потрібно захищати або обмежувати; SC-05\_\allowbreak{}ODP[02] вибрано одне з наступних ЗНАЧЕНЬ ПАРАМЕТРІВ: \{захистити від; обмежити\}}




\subsubsection{ОБМЕЖЕННЯ ВНУТРІШНІХ КОРИСТУВАЧІВ (SC-5(1))}


\vspace{10pt}
{\footnotesize\bfseries
No: 1\\
Name: sc\_\allowbreak{}5\_\allowbreak{}1\_\allowbreak{}01\\
Type: string\\
Default: nil
}

{\footnotesize Обмежена можливість окремих осіб здійснювати атаки на відмову в обслуговуванні проти інших систем}


{\footnotesize\bfseries
No: 2\\
Name: sc\_\allowbreak{}5\_\allowbreak{}1\_\allowbreak{}odp\\
Type: list\\
Default: [\textquotedbl{}admin\textquotedbl{}, \textquotedbl{}security\_\allowbreak{}officer\textquotedbl{}]
}

{\footnotesize Визначені атаки на відмову в обслуговуванні, для яких необхідно обмежити можливість їх запуску окремими особами}




\subsubsection{ЗАХИСТ ВІД АТАК «ВІДМОВА В ОБСЛУГОВУВАННІ» ПРОДУКТИВНІСТЬ, ПРОПУСКНА ЗДАТНІСТЬ ТА НАДМІРНІСТЬ (SC-5(2))}


\vspace{10pt}
{\footnotesize\bfseries
No: 1\\
Name: sc\_\allowbreak{}5\_\allowbreak{}2\_\allowbreak{}01\\
Type: string\\
Default: nil
}

{\footnotesize Здійснюється управління ємністю, пропускною здатністю або іншими надлишковими ресурсами для обмеження наслідків інформаційних атак на відмову в обслуговуванні}




\subsubsection{ВИЯВЛЕННЯ ТА МОНІТОРИНГ (SC-5(3))}


\vspace{10pt}
{\footnotesize\bfseries
No: 1\\
Name: sc\_\allowbreak{}5\_\allowbreak{}3\_\allowbreak{}a\\
Type: string\\
Default: \textquotedbl{}автоматизований засіб моніторингу\textquotedbl{}
}

{\footnotesize Використовуються засоби моніторингу для виявлення ознак атак на відмову в обслуговуванні, спрямованих на систему або запущених з неї}


{\footnotesize\bfseries
No: 2\\
Name: sc\_\allowbreak{}5\_\allowbreak{}3\_\allowbreak{}b\\
Type: string\\
Default: nil
}

{\footnotesize Здійснюється моніторинг системних ресурсів для визначення наявності достатніх ресурсів для запобігання ефективним атакам на відмову в обслуговуванні}




\subsection{ДОСТУПНІСТЬ РЕСУРСІВ (SC-6)}
Забезпечити захист доступності ресурсів, виділивши [Призначення: визначені організацією ресурси], по [Вибір (один або кілька); пріоритет; квоти; [Призначення: визначені організацією заходи з безпеки]].

\vspace{10pt}
{\footnotesize\bfseries
No: 1\\
Name: sc\_\allowbreak{}6\_\allowbreak{}01\\
Type: string\\
Default: nil
}

{\footnotesize Захищено доступність ресурсів шляхом розподілу ресурсів за ВИБІРКОВИМ ЗНАЧЕННЯМ ПАРАМЕТРА(ів)}


{\footnotesize\bfseries
No: 2\\
Name: sc\_\allowbreak{}6\_\allowbreak{}odp\_\allowbreak{}01\\
Type: string\\
Default: nil
}

{\footnotesize Визначені ресурси, які необхідно виділити для захисту доступності ресурсів}


{\footnotesize\bfseries
No: 3\\
Name: sc\_\allowbreak{}6\_\allowbreak{}odp\_\allowbreak{}02\\
Type: string\\
Default: nil
}

{\footnotesize Вибрано одне або декілька з наступних ЗНАЧЕНЬ ПАРАМЕТРІВ: \{priority; quota; controls\}}


{\footnotesize\bfseries
No: 4\\
Name: sc\_\allowbreak{}6\_\allowbreak{}odp\_\allowbreak{}03\\
Type: string\\
Default: \textquotedbl{}автоматизований засіб моніторингу\textquotedbl{}
}

{\footnotesize Визначені засоби контролю для захисту доступності ресурсів (якщо вибрано)}




\subsection{ДОСТУПНІСТЬ РЕСУРСІВ (SC-7)}
a. Контролювати та управляти зв'язком на зовнішньому периметрі системи та на ключових внутрішніх периметрах всередині системи.\\ 
b. Реалізувати підмережі для загальнодоступних компонентів системи, які є [Вибір: фізично; логічно] відділені від внутрішніх мереж організації.\\ 
c. Підключатися до зовнішніх мереж або систем тільки через керовані інтерфейси, що складаються з пристроїв захисту периметру, і розташованих відповідно до архітектури безпеки та приватності організації.

\vspace{10pt}
\textit{Немає параметрів для цього контролю.}
\vspace{10pt}


\subsubsection{ФІЗИЧНО ВІДДІЛЕНІ ПІДМЕРЕЖІ (SC-7(1)) [Вилучено]}
[Вилучено: включено до SC-7]

\vspace{10pt}
\textit{Немає параметрів для цього контролю.}
\vspace{10pt}


\subsubsection{ПУБЛІЧНИЙ ДОСТУП (SC-7(2)) [Вилучено]}
[Вилучено: включено до SC-7]

\vspace{10pt}
\textit{Немає параметрів для цього контролю.}
\vspace{10pt}


\subsubsection{ТОЧКИ ДОСТУПУ (SC-7(3))}


\vspace{10pt}
{\footnotesize\bfseries
No: 1\\
Name: sc\_\allowbreak{}7\_\allowbreak{}3\_\allowbreak{}01\\
Type: integer\\
Default: 3
}

{\footnotesize Обмежена кількість зовнішніх мережевих підключень до системи}




\subsubsection{ЗОВНІШНІ КОМУНІКАЦІЙНІ СЛУЖБИ (SC-7(4))}


\vspace{10pt}
{\footnotesize\bfseries
No: 1\\
Name: sc\_\allowbreak{}7\_\allowbreak{}4\_\allowbreak{}a\\
Type: string\\
Default: nil
}

{\footnotesize Реалізовано керований інтерфейс для кожної зовнішньої телекомунікаційної послуги}


{\footnotesize\bfseries
No: 2\\
Name: sc\_\allowbreak{}7\_\allowbreak{}4\_\allowbreak{}b\\
Type: list\\
Default: [\textquotedbl{}default\_\allowbreak{}deny\_\allowbreak{}rule\textquotedbl{}, \textquotedbl{}abac\_\allowbreak{}rule\_\allowbreak{}1\textquotedbl{}]
}

{\footnotesize Встановлена політика потоку трафіку для кожного керованого інтерфейсу}


{\footnotesize\bfseries
No: 3\\
Name: sc\_\allowbreak{}7\_\allowbreak{}4\_\allowbreak{}c\_\allowbreak{}01\\
Type: string\\
Default: nil
}

{\footnotesize Захищена конфіденційність інформації, що передається через кожен інтерфейс}


{\footnotesize\bfseries
No: 4\\
Name: sc\_\allowbreak{}7\_\allowbreak{}4\_\allowbreak{}c\_\allowbreak{}02\\
Type: string\\
Default: nil
}

{\footnotesize Захищена цілісність інформації, що передається через кожен інтерфейс}


{\footnotesize\bfseries
No: 5\\
Name: sc\_\allowbreak{}7\_\allowbreak{}4\_\allowbreak{}d\\
Type: list\\
Default: [\textquotedbl{}default\_\allowbreak{}deny\_\allowbreak{}rule\textquotedbl{}, \textquotedbl{}abac\_\allowbreak{}rule\_\allowbreak{}1\textquotedbl{}]
}

{\footnotesize Задокументовано кожен виняток з політики управління трафіком з обґрунтуванням місії або бізнес-потреби, а також тривалості такої потреби}


{\footnotesize\bfseries
No: 6\\
Name: sc\_\allowbreak{}7\_\allowbreak{}4\_\allowbreak{}e\_\allowbreak{}01\\
Type: list\\
Default: [\textquotedbl{}default\_\allowbreak{}deny\_\allowbreak{}rule\textquotedbl{}, \textquotedbl{}abac\_\allowbreak{}rule\_\allowbreak{}1\textquotedbl{}]
}

{\footnotesize Переглядаються винятки з політики потоку трафіку частота}


{\footnotesize\bfseries
No: 7\\
Name: sc\_\allowbreak{}7\_\allowbreak{}4\_\allowbreak{}e\_\allowbreak{}02\\
Type: list\\
Default: [\textquotedbl{}default\_\allowbreak{}deny\_\allowbreak{}rule\textquotedbl{}, \textquotedbl{}abac\_\allowbreak{}rule\_\allowbreak{}1\textquotedbl{}]
}

{\footnotesize Потрібно видалити винятки з політики потоку трафіку, які більше не підтримуються чітко визначеною місією або бізнеспотребою}


{\footnotesize\bfseries
No: 8\\
Name: sc\_\allowbreak{}7\_\allowbreak{}4\_\allowbreak{}f\\
Type: string\\
Default: nil
}

{\footnotesize Запобігається несанкціонований обмін управління із зовнішніми мережами}


{\footnotesize\bfseries
No: 9\\
Name: sc\_\allowbreak{}7\_\allowbreak{}4\_\allowbreak{}g\\
Type: string\\
Default: nil
}

{\footnotesize Публікується інформація, яка дозволяє віддаленим мережам виявляти несанкціонований трафік площини керування з внутрішніх мереж}


{\footnotesize\bfseries
No: 10\\
Name: sc\_\allowbreak{}7\_\allowbreak{}4\_\allowbreak{}h\\
Type: string\\
Default: nil
}

{\footnotesize Фільтрується несанкціонований трафік з зовнішніх мереж. трафіком плану}


{\footnotesize\bfseries
No: 11\\
Name: sc\_\allowbreak{}7\_\allowbreak{}4\_\allowbreak{}odp\\
Type: list\\
Default: [\textquotedbl{}default\_\allowbreak{}deny\_\allowbreak{}rule\textquotedbl{}, \textquotedbl{}abac\_\allowbreak{}rule\_\allowbreak{}1\textquotedbl{}]
}

{\footnotesize Визначено періодичність перегляду винятків з політики управління інформаційними потоками}




\subsubsection{ВІДМОВА ЗА ЗАМОВЧУВАННЯМ - ДОЗВІЛ ЗА ВИНЯТКОМ (SC-7(5))}


\vspace{10pt}
\textit{Немає параметрів для цього контролю.}
\vspace{10pt}


\subsubsection{ВІДПОВІДЬ НА РОЗПІЗНАНІ ПОМИЛКИ (SC-7(6)) [Вилучено]}
[Вилучено: включено до SC-7(18)]

\vspace{10pt}
\textit{Немає параметрів для цього контролю.}
\vspace{10pt}


\subsubsection{ЗАПОБІГАННЯ ПОДІЛУ ТУНЕЛЮВАННЯ ДЛЯ ВІДДАЛЕНИХ ПРИСТРОЇВ (SC-7(7))}


\vspace{10pt}
{\footnotesize\bfseries
No: 1\\
Name: sc\_\allowbreak{}7\_\allowbreak{}7\_\allowbreak{}01\\
Type: string\\
Default: \textquotedbl{}автоматизований засіб моніторингу\textquotedbl{}
}

{\footnotesize Запобігається розділеному тунелюванню для віддалених пристроїв, що підключаються до систем організації, якщо розділене тунелюванню не захищено за допомогою засоби захисту}


{\footnotesize\bfseries
No: 2\\
Name: sc\_\allowbreak{}7\_\allowbreak{}7\_\allowbreak{}odp\\
Type: string\\
Default: nil
}

{\footnotesize Визначені гарантії безпечного прокладання розділеному тунелюванню}




\subsubsection{МАРШРУТИЗАЦІЯ ТРАФІКУ З АВТЕНТИФІКОВАНИХ ПРОКСІ-СЕРВЕРІВ (SC-7(8))}


\vspace{10pt}
{\footnotesize\bfseries
No: 1\\
Name: sc\_\allowbreak{}7\_\allowbreak{}8\_\allowbreak{}01\\
Type: string\\
Default: nil
}

{\footnotesize SC-07(08) \_\allowbreak{}ODP[01] внутрішній комунікаційний трафік> спрямовується до <SC-07(08) \_\allowbreak{}ODP[02] зовнішніх мереж> через автентифіковані проксі-сервери на керованих інтерфейсах}




\subsubsection{ЗАХИСТ (SC-7(9))}


\vspace{10pt}
{\footnotesize\bfseries
No: 1\\
Name: sc\_\allowbreak{}7\_\allowbreak{}9\_\allowbreak{}a\_\allowbreak{}01\\
Type: string\\
Default: nil
}

{\footnotesize Виявлено вихідний комунікаційний трафік, що становить загрозу для зовнішніх систем}


{\footnotesize\bfseries
No: 2\\
Name: sc\_\allowbreak{}7\_\allowbreak{}9\_\allowbreak{}a\_\allowbreak{}02\\
Type: string\\
Default: nil
}

{\footnotesize Заборонено вихідний комунікаційний трафік, що становить загрозу для зовнішніх систем}


{\footnotesize\bfseries
No: 3\\
Name: sc\_\allowbreak{}7\_\allowbreak{}9\_\allowbreak{}b\\
Type: string\\
Default: nil
}

{\footnotesize Перевіряється ідентичність внутрішніх пов'язаних з відмовою у зв'язку. користувачів,}




\subsubsection{ІЗОЛЬОВАНІ ВІД ІНШИХ ВНУТРІШНІХ КОМПОНЕНТІВ СИСТЕМИ ШЛЯХОМ ВПРОВАДЖЕННЯ ФІЗИЧНО ВІДОКРЕМЛЕНИХ ПІДМЕРЕЖ З КЕРОВАНИМИ ІНТЕРФЕЙСАМИ ДО ІНШИХ КОМПОНЕНТІВ СИСТЕМИ (SC-7(13))}


\vspace{10pt}
{\footnotesize\bfseries
No: 1\\
Name: sc\_\allowbreak{}7\_\allowbreak{}13\_\allowbreak{}01\\
Type: string\\
Default: \textquotedbl{}автоматизований засіб моніторингу\textquotedbl{}
}

{\footnotesize Ізольовані засоби, механізми та компоненти підтримки інформаційної безпеки від інших внутрішніх компонентів системи шляхом впровадження фізично відокремлених підмереж з керованими інтерфейсами до інших компонентів системи}


{\footnotesize\bfseries
No: 2\\
Name: sc\_\allowbreak{}7\_\allowbreak{}13\_\allowbreak{}odp\\
Type: string\\
Default: \textquotedbl{}автоматизований засіб моніторингу\textquotedbl{}
}

{\footnotesize Визначені інструменти, механізми та компоненти підтримки інформаційної безпеки, які мають бути ізольовані від інших внутрішніх компонентів системи}




\subsection{КОНФІДЕНЦІЙНІСТЬ ТА ЦІЛІСНІСТЬ ПЕРЕДАЧІ (SC-8)}
Забезпечити [Вибір (один або кілька): конфіденційність; цілісність] інформації, що передається.

\vspace{10pt}
\textit{Немає параметрів для цього контролю.}
\vspace{10pt}


\subsubsection{КОНФІДЕНЦІЙНІСТЬ ТА КРИПТОГРАФІЧНИЙ ЗАХИСТ (SC-8(1))}


\vspace{10pt}
{\footnotesize\bfseries
No: 1\\
Name: sc\_\allowbreak{}8\_\allowbreak{}1\_\allowbreak{}odp\\
Type: string\\
Default: nil
}

{\footnotesize Вибрано одне або декілька з наступних ЗНАЧЕНЬ ПАРАМЕТРІВ: \{запобігти несанкціонованому розголошенню інформації; виявити зміни в інформації\}}




\subsubsection{ПОПЕРЕДНЯ І ПОСТОБРОБКА (SC-8(2))}


\vspace{10pt}
{\footnotesize\bfseries
No: 1\\
Name: sc\_\allowbreak{}8\_\allowbreak{}2\_\allowbreak{}01\\
Type: string\\
Default: nil
}

{\footnotesize КОНФІДЕНЦІЙНІСТЬ ТА ЦІЛІСНІСТЬ ПЕРЕДАЧІ - ПОПЕРЕДНЯ І ПОСТОБРОБКА МЕТА ОЦІНКИ: Визначити, чи:}


{\footnotesize\bfseries
No: 2\\
Name: sc\_\allowbreak{}8\_\allowbreak{}2\_\allowbreak{}odp\\
Type: string\\
Default: nil
}

{\footnotesize Вибрано одне або декілька з наступних ЗНАЧЕНЬ ПАРАМЕТРІВ: \{конфіденційність; цілісність\}}




\subsubsection{КОНФІДЕНЦІЙНІСТЬ ТА ЦІЛІСНІСТЬ КРИПТОГРАФІЧНИЙ ЗАХИСТ ПОВІДОМЛЕНЬ (SC-8(3))}


\vspace{10pt}
{\footnotesize\bfseries
No: 1\\
Name: sc\_\allowbreak{}8\_\allowbreak{}3\_\allowbreak{}01\\
Type: string\\
Default: \textquotedbl{}AES-256-GCM\textquotedbl{}
}

{\footnotesize Впроваджено криптографічні механізми для захисту зовнішніх повідомлень, якщо інше не захищено альтернативні фізичні засоби контролю}


{\footnotesize\bfseries
No: 2\\
Name: sc\_\allowbreak{}8\_\allowbreak{}3\_\allowbreak{}odp\\
Type: string\\
Default: \textquotedbl{}автоматизований засіб моніторингу\textquotedbl{}
}

{\footnotesize Визначено альтернативні фізичні засоби контролю для захисту зовнішніх повідомлень}




\subsubsection{ПРИХОВУВАННЯ АБО РАНДОМІЗАЦІЯ КОМУНІКАЦІЇ (SC-8(4))}


\vspace{10pt}
{\footnotesize\bfseries
No: 1\\
Name: sc\_\allowbreak{}8\_\allowbreak{}4\_\allowbreak{}01\\
Type: string\\
Default: \textquotedbl{}AES-256-GCM\textquotedbl{}
}

{\footnotesize Застосовуються криптографічні механізми для приховування або рандомізації шаблонів комунікації, якщо інше не захищено альтернативні фізичні засоби контролю}


{\footnotesize\bfseries
No: 2\\
Name: sc\_\allowbreak{}8\_\allowbreak{}4\_\allowbreak{}odp\\
Type: string\\
Default: \textquotedbl{}автоматизований засіб моніторингу\textquotedbl{}
}

{\footnotesize Визначені альтернативні фізичні засоби контролю для захисту від несанкціонованого розкриття шаблонів комунікації}




\subsubsection{КОНФІДЕНЦІЙНІСТЬ ТА СИСТЕМА РОЗПОДІЛУ (SC-8(5))}


\vspace{10pt}
\textit{Немає параметрів для цього контролю.}
\vspace{10pt}


\subsection{КОНФІДЕНЦІЙНІСТЬ ПЕРЕДАЧІ (SC-9) [Вилучено]}
[Вилучено: включено до SC-8]

\vspace{10pt}
\textit{Немає параметрів для цього контролю.}
\vspace{10pt}


\subsection{ВІДКЛЮЧЕННЯ МЕРЕЖІ (SC-10)}
Завершити з'єднання з мережею, яке пов'язане із сеансом зв'язку в кінці сеансу або після [Призначення: визначений організацією період часу] бездіяльності.

\vspace{10pt}
{\footnotesize\bfseries
No: 1\\
Name: sc\_\allowbreak{}10\_\allowbreak{}odp\\
Type: integer\\
Default: 30
}

{\footnotesize Визначено період бездіяльності, після якого система розриває мережеве з'єднання, пов'язане з сеансом зв'язку; SC08(05)\_\allowbreak{}ODP[02] мережеве з'єднання, пов'язане з сеансом зв'язку, розірвано в кінці сеансу або після період часу бездіяльності}




\subsection{ДОВІРЕНИЙ КАНАЛ ЗВ'ЯЗКУ (SC-11)}
a. Надати [Вибір: фізично; логічно] ізольований надійний канал зв'язку для зв'язку між користувачем і довіреними компонентами системи.\\ 
b. Дозволити користувачам запросити довірений канал зв'язку для обміну даними між користувачем і наступними функціями безпеки системи, включно з, як мінімум, автентифікацією та повторною автентифікацією: [Призначення: визначені організацією функції безпеки].

\vspace{10pt}
{\footnotesize\bfseries
No: 1\\
Name: sc\_\allowbreak{}11\_\allowbreak{}odp\_\allowbreak{}01\\
Type: string\\
Default: nil
}

{\footnotesize Вибрано одне з наступних ЗНАЧЕНЬ ПАРАМЕТРІВ: \{фізично; логічно\}}




\subsection{ВСТАНОВЛЕННЯ КЛЮЧАМИ (SC-12)}
Встановити та управляти криптографічними ключами для криптографічних засобів, які використовуються в системі відповідно до [Призначення: визначені організацією вимоги до генерації, поширення, зберігання, доступу та знищення ключів].

\vspace{10pt}
{\footnotesize\bfseries
No: 1\\
Name: sc\_\allowbreak{}12\_\allowbreak{}01\\
Type: string\\
Default: \textquotedbl{}AES-256-GCM\textquotedbl{}
}

{\footnotesize Встановлюються криптографічні ключі, коли в системі використовується криптографія відповідно до < SC-12\_\allowbreak{}ODP вимог >}


{\footnotesize\bfseries
No: 2\\
Name: sc\_\allowbreak{}12\_\allowbreak{}02\\
Type: string\\
Default: \textquotedbl{}AES-256-GCM\textquotedbl{}
}

{\footnotesize Здійснюється управління криптографічними ключами, коли в системі використовується криптографія, відповідно до вимог }


{\footnotesize\bfseries
No: 3\\
Name: sc\_\allowbreak{}12\_\allowbreak{}odp\\
Type: string\\
Default: nil
}

{\footnotesize Визначені вимоги до генерації, розповсюдження, зберігання, доступу та знищення ключів}




\subsection{КРИПТОГРАФІЧНИЙ ЗАХИСТ (SC-13)}
a. Визначити [Призначення: використання криптографічних засобів, визначених організацією];\\ 
b. Впровадити [Завдання: визначені організацією види криптографії для кожного визначеного криптографічного використання].

\vspace{10pt}
{\footnotesize\bfseries
No: 1\\
Name: sc\_\allowbreak{}13\_\allowbreak{}01\\
Type: string\\
Default: \textquotedbl{}AES-256-GCM\textquotedbl{}
}

{\footnotesize КРИПТОГРАФІЧНИЙ ЗАХИСТ МЕТА ОЦІНКИ: Визначити, чи:}


{\footnotesize\bfseries
No: 2\\
Name: sc\_\allowbreak{}13\_\allowbreak{}odp\_\allowbreak{}01\\
Type: string\\
Default: \textquotedbl{}AES-256-GCM\textquotedbl{}
}

{\footnotesize Визначено використання криптографічних засобів}


{\footnotesize\bfseries
No: 3\\
Name: sc\_\allowbreak{}13\_\allowbreak{}odp\_\allowbreak{}02\\
Type: string\\
Default: \textquotedbl{}AES-256-GCM\textquotedbl{}
}

{\footnotesize Визначено типи криптографії для кожного вказаного криптографічного використання; SC-13a. ідентифіковано використання>; SC-13b. реалізовано типи криптографії для кожного вказаного криптографічного використання (визначеного в SC-13\_\allowbreak{}ODP[01]). криптографічне <SC-13\_\allowbreak{}ODP[01]}




\subsection{ЗАХИСТ ГРОМАДСЬКОГО ДОСТУПУ (SC-14) [Вилучено]}
[Вилучено: включено до AC-2, AC-3, AC-5, AC-6, SI-3, SI-4, SI-5, SI-7, SI-10]

\vspace{10pt}
\textit{Немає параметрів для цього контролю.}
\vspace{10pt}


\subsection{СПІЛЬНІ ОБЧИСЛЮВАЛЬНІ ПРИСТРОЇ ТА ЗАСТОСУНКИ (SC-15)}
a. Заборонити віддалену активацію спільних обчислювальних пристроїв (хмар) та застосунків з такими виключеннями: [Призначення: визначені організацією виключення, у яких дозволена віддалена активація].\\ 
b. Надати явну вказівку щодо використання користувачами фізично присутніми пристроями.

\vspace{10pt}
\textit{Немає параметрів для цього контролю.}
\vspace{10pt}


\subsection{ПЕРЕДАЧА АТРИБУТІВ БЕЗПЕКИ ТА ПРИВАТНОСТІ (SC-16)}
Пов'язувати [Призначення: визначені організацією атрибути безпеки та приватності] з інформацією, яка передається між системами та компонентами системи.

\vspace{10pt}
{\footnotesize\bfseries
No: 1\\
Name: sc\_\allowbreak{}16\_\allowbreak{}01\\
Type: list\\
Default: [\textquotedbl{}default\_\allowbreak{}deny\_\allowbreak{}rule\textquotedbl{}, \textquotedbl{}abac\_\allowbreak{}rule\_\allowbreak{}1\textquotedbl{}]
}

{\footnotesize Пов'язані атрибути інформацією, якою обмінюються системи; безпеки з}


{\footnotesize\bfseries
No: 2\\
Name: sc\_\allowbreak{}16\_\allowbreak{}02\\
Type: list\\
Default: [\textquotedbl{}default\_\allowbreak{}deny\_\allowbreak{}rule\textquotedbl{}, \textquotedbl{}abac\_\allowbreak{}rule\_\allowbreak{}1\textquotedbl{}]
}

{\footnotesize Пов'язані атрибути безпеки інформацією, якою обмінюються компоненти системи; з}


{\footnotesize\bfseries
No: 3\\
Name: sc\_\allowbreak{}16\_\allowbreak{}03\\
Type: list\\
Default: [\textquotedbl{}default\_\allowbreak{}deny\_\allowbreak{}rule\textquotedbl{}, \textquotedbl{}abac\_\allowbreak{}rule\_\allowbreak{}1\textquotedbl{}]
}

{\footnotesize Пов'язані атрибути приватності інформацією, якою обмінюються системи; з}


{\footnotesize\bfseries
No: 4\\
Name: sc\_\allowbreak{}16\_\allowbreak{}04\\
Type: list\\
Default: [\textquotedbl{}default\_\allowbreak{}deny\_\allowbreak{}rule\textquotedbl{}, \textquotedbl{}abac\_\allowbreak{}rule\_\allowbreak{}1\textquotedbl{}]
}

{\footnotesize Пов'язані атрибути приватності інформацією, якою обмінюються компоненти системи. з}


{\footnotesize\bfseries
No: 5\\
Name: sc\_\allowbreak{}16\_\allowbreak{}odp\_\allowbreak{}01\\
Type: list\\
Default: [\textquotedbl{}default\_\allowbreak{}deny\_\allowbreak{}rule\textquotedbl{}, \textquotedbl{}abac\_\allowbreak{}rule\_\allowbreak{}1\textquotedbl{}]
}

{\footnotesize Визначені атрибути безпеки, які будуть пов'язані з інформацією, що обмінюється}


{\footnotesize\bfseries
No: 6\\
Name: sc\_\allowbreak{}16\_\allowbreak{}odp\_\allowbreak{}02\\
Type: list\\
Default: [\textquotedbl{}default\_\allowbreak{}deny\_\allowbreak{}rule\textquotedbl{}, \textquotedbl{}abac\_\allowbreak{}rule\_\allowbreak{}1\textquotedbl{}]
}

{\footnotesize Визначені атрибути приватності, які будуть пов'язані з інформацією, що обмінюється}




\subsection{СЕРТИФІКАТИ ІНФРАСТРУКТУРИ ВІДКРИТИХ КЛЮЧІВ (SC-17)}
a. Випускати сертифікати відкритого ключа відповідно до [Призначення: визначеної організацією політики сертифікації];\\ 
b. Отримувати сертифікати відкритого ключа від затвердженого постачальника послуг.

\vspace{10pt}
\textit{Немає параметрів для цього контролю.}
\vspace{10pt}


\subsection{МОБІЛЬНИЙ КОД (SC-18)}
a. Визначати прийнятні та неприйнятні мобільні коди та технології мобільних кодів.\\ 
b. Проводити авторизацію, відстежувати та контролювати використання мобільного коду всередині системи.

\vspace{10pt}
\textit{Немає параметрів для цього контролю.}
\vspace{10pt}


\subsection{ІНТЕРНЕТ-ПРОТОКОЛ ГОЛОСОВОГО ЗВ'ЯЗКУ (SC-19) [Вилучено]}
[Вилучено: залежить від технології; розглядається як будь-яка інша технологія або протокол]

\vspace{10pt}
\textit{Немає параметрів для цього контролю.}
\vspace{10pt}


\subsection{БЕЗПЕЧНА СЛУЖБА ІМЕН, АДРЕС (УПОВНОВАЖЕНЕ ДЖЕРЕЛО) (SC-20)}
a. Надати додаткові дані автентифікації та перевірки цілісності джерела даних разом з офіційними даними розпізнавання імен, які система повертає у відповідь на запити дозволу імен/адрес.\\ 
b. Надати засоби для вказання статусу безпеки дочірніх зон і (якщо дочірня зона підтримує служби безпечного дозволу) забезпечити перевірку ланцюга довіри між батьківськими та дочірніми доменами при роботі в складі розподіленого ієрархічного простору імен.

\vspace{10pt}
\textit{Немає параметрів для цього контролю.}
\vspace{10pt}


\subsection{БЕЗПЕЧНА СЛУЖБА ІМЕН, АДРЕС (УПОВНОВАЖЕНЕ ДЖЕРЕЛО) ДЖЕРЕЛО ДАНИХ ТА ЦІЛІСНІСТЬ (SC-21)}
Зробити запит та виконати перевірку автентичності джерела даних і перевірку цілісності даних у відповідях на дозвіл імен/адрес, які система отримує від уповноважених джерел.

\vspace{10pt}
{\footnotesize\bfseries
No: 1\\
Name: sc\_\allowbreak{}21\_\allowbreak{}01\\
Type: string\\
Default: nil
}

{\footnotesize Реалізується запит перевірки автентичності джерела даних для відповідей на запит дозволу імен/адрес, які система отримує від авторитетних джерел}


{\footnotesize\bfseries
No: 2\\
Name: sc\_\allowbreak{}21\_\allowbreak{}02\\
Type: string\\
Default: nil
}

{\footnotesize Реалізується запит автентифікація походження даних на основі відповідей з дозволу імен/адрес, які система отримує від авторитетних джерел}


{\footnotesize\bfseries
No: 3\\
Name: sc\_\allowbreak{}21\_\allowbreak{}03\\
Type: string\\
Default: nil
}

{\footnotesize Реалізується запит перевірки цілісності даних для відповідей на запит дозволу імен/адрес, які система отримує від авторитетних джерел}


{\footnotesize\bfseries
No: 4\\
Name: sc\_\allowbreak{}21\_\allowbreak{}04\\
Type: string\\
Default: nil
}

{\footnotesize Виконується перевірка цілісності даних для відповідей на запит про дозвіл імен/адрес, які система отримує від авторитетних джерел}




\subsection{АРХІТЕКТУРА ТА ЗАБЕЗПЕЧЕННЯ СЛУЖБИ ІМЕН/АДРЕС (SC-22)}
Переконатися, що системи, які спільно надають послуги розпізнавання імен/адрес для організації, є відмовостійкими та забезпечують поділ внутрішніх і зовнішніх ролей.

\vspace{10pt}
{\footnotesize\bfseries
No: 1\\
Name: sc\_\allowbreak{}22\_\allowbreak{}01\\
Type: string\\
Default: nil
}

{\footnotesize Є системи, які спільно надають послуги з визначення імен/адрес для організації, відмовостійкими}


{\footnotesize\bfseries
No: 2\\
Name: sc\_\allowbreak{}22\_\allowbreak{}02\\
Type: string\\
Default: nil
}

{\footnotesize Реалізовано в системах, які спільно надають послуги з вирішення імен/адрес для організації, внутрішній розподіл ролей}


{\footnotesize\bfseries
No: 3\\
Name: sc\_\allowbreak{}22\_\allowbreak{}03\\
Type: string\\
Default: nil
}

{\footnotesize Реалізовано в системах, які спільно надають послуги з вирішення імен/адрес для організації, зовнішній розподіл ролей}




\subsection{АВТЕНТИФІКАЦІЯ СЕСІЇ (SC-23)}
Забезпечити автентифікацію сеансів зв'язку.

\vspace{10pt}
{\footnotesize\bfseries
No: 1\\
Name: sc\_\allowbreak{}23\_\allowbreak{}01\\
Type: string\\
Default: nil
}

{\footnotesize Захищено автентифікацію сеансів зв'язку}




\subsection{УВЕДЕННЯ У ВІДОМИЙ СТАН (SC-24)}
Увести систему в [Призначення: визначений організацією відомий стан системи] у разі [Призначення: визначені організацією типи збоїв системи] зі збереженням [Призначення: визначена організацією інформація про стан системи] при збої.

\vspace{10pt}
{\footnotesize\bfseries
No: 1\\
Name: sc\_\allowbreak{}24\_\allowbreak{}01\\
Type: string\\
Default: nil
}

{\footnotesize Типи системних збоїв на компонентах системи призводять до відомого стану системи, зберігаючи при цьому інформацію про стан системи у збої}


{\footnotesize\bfseries
No: 2\\
Name: sc\_\allowbreak{}24\_\allowbreak{}odp\_\allowbreak{}01\\
Type: string\\
Default: nil
}

{\footnotesize Визначені типи відмов системи, за яких компоненти системи переходять до відомого стану}


{\footnotesize\bfseries
No: 3\\
Name: sc\_\allowbreak{}24\_\allowbreak{}odp\_\allowbreak{}02\\
Type: string\\
Default: nil
}

{\footnotesize Відомий стан системи, до якого переходять компоненти системи у випадку її відмови}


{\footnotesize\bfseries
No: 4\\
Name: sc\_\allowbreak{}24\_\allowbreak{}odp\_\allowbreak{}03\\
Type: string\\
Default: nil
}

{\footnotesize Потрібно зберігати інформацію про стан системи у випадку її збою}




\subsection{ТОНКІ ВУЗЛИ (SC-25)}
Використовувати [Призначення: визначені організацією системні компоненти] з мінімальною функціональністю та зберіганням інформації.

\vspace{10pt}
{\footnotesize\bfseries
No: 1\\
Name: sc\_\allowbreak{}25\_\allowbreak{}01\\
Type: string\\
Default: nil
}

{\footnotesize Використовується мінімальна функціональність для компонентів системи}


{\footnotesize\bfseries
No: 2\\
Name: sc\_\allowbreak{}25\_\allowbreak{}02\\
Type: string\\
Default: nil
}

{\footnotesize Виділено мінімальне сховище інформації на компоненти системи}


{\footnotesize\bfseries
No: 3\\
Name: sc\_\allowbreak{}25\_\allowbreak{}odp\\
Type: string\\
Default: nil
}

{\footnotesize Потрібно використовувати компоненти системи з мінімальною функціональністю та обсягом зберігання інформації}




\subsection{ПРИМАНКА ДЛЯ ЗЛОВМИСНИКІВ (DECOYS) (SC-26)}
Вносити в систему компоненти, які спеціально призначені як об'єкти атак, з метою виявлення, відбиття й аналізу таких атак.

\vspace{10pt}
{\footnotesize\bfseries
No: 1\\
Name: sc\_\allowbreak{}26\_\allowbreak{}01\\
Type: string\\
Default: \textquotedbl{}автоматизований засіб моніторингу\textquotedbl{}
}

{\footnotesize Є в системах організації компоненти, спеціально розроблені для того, щоб стати мішенню зловмисних атак, і чи є в них засоби для виявлення таких атак}


{\footnotesize\bfseries
No: 2\\
Name: sc\_\allowbreak{}26\_\allowbreak{}02\\
Type: string\\
Default: \textquotedbl{}автоматизований засіб моніторингу\textquotedbl{}
}

{\footnotesize Є в організаційних компонентах системи, спеціально розроблені для того, щоб стати мішенню зловмисних атак, і чи є в них засоби для відбиття таких атак}


{\footnotesize\bfseries
No: 3\\
Name: sc\_\allowbreak{}26\_\allowbreak{}03\\
Type: string\\
Default: nil
}

{\footnotesize Включені в організаційні компоненти системи, спеціально розроблені для того, щоб бути мішенню зловмисних атак, для аналізу таких атак}




\subsection{НЕЗАЛЕЖНІ ВІД ПЛАТФОРМИ ЗАСТОСУНКИ (SC-27)}
Внести до системи: [Призначення: визначені організацією незалежні від платформи застосунки].

\vspace{10pt}
{\footnotesize\bfseries
No: 1\\
Name: sc\_\allowbreak{}27\_\allowbreak{}01\\
Type: string\\
Default: nil
}

{\footnotesize Включені незалежні від платформи додатки в системи організації}


{\footnotesize\bfseries
No: 2\\
Name: sc\_\allowbreak{}27\_\allowbreak{}odp\\
Type: string\\
Default: nil
}

{\footnotesize Визначені незалежні від платформи додатки, які мають бути включені в системи організації}




\subsection{ЗАХИСТ ІНФОРМАЦІЇ В СТАНІ СПОКОЮ (SC-28)}
Забезпечити [Вибір (один або кілька): конфіденційність; цілісність] [Призначення: визначена організацією інформація] в стані спокою.

\vspace{10pt}
{\footnotesize\bfseries
No: 1\\
Name: sc\_\allowbreak{}28\_\allowbreak{}odp\\
Type: string\\
Default: nil
}

{\footnotesize Вибрано одне або декілька з наступних ЗНАЧЕНЬ ПАРАМЕТРІВ: \{конфіденційність; цілісність\}}


{\footnotesize\bfseries
No: 2\\
Name: sc\_\allowbreak{}28\_\allowbreak{}odp\_\allowbreak{}02\\
Type: string\\
Default: nil
}

{\footnotesize Є інформація в стані спокою, яка потребує захисту}




\subsubsection{ЗАХИСТ ІНФОРМАЦІЇ В СТАНІ СПОКОЮ | КРИПТОГРАФІЧНИЙ ЗАХИСТ}


\vspace{10pt}
{\footnotesize\bfseries
No: 1\\
Name: sc\_\allowbreak{}28\_\allowbreak{}1\_\allowbreak{}01\\
Type: string\\
Default: \textquotedbl{}AES-256-GCM\textquotedbl{}
}

{\footnotesize Реалізовані криптографічні механізми для запобігання несанкціонованому розкриттю інформації, що знаходиться в стані спокою на системних компонентах або носіях}


{\footnotesize\bfseries
No: 2\\
Name: sc\_\allowbreak{}28\_\allowbreak{}1\_\allowbreak{}02\\
Type: string\\
Default: \textquotedbl{}AES-256-GCM\textquotedbl{}
}

{\footnotesize Реалізовані криптографічні механізми для запобігання несанкціонованій модифікації інформації, що знаходиться в стані спокою на системних компонентах або носіях}




\subsection{ГЕТЕРОГЕННІСТЬ (SC-29)}
Використовувати різноманітний набір інформаційних технологій для [Призначення: визначені організацією системні компоненти] при впровадженні системи.

\vspace{10pt}
{\footnotesize\bfseries
No: 1\\
Name: sc\_\allowbreak{}29\_\allowbreak{}01\\
Type: string\\
Default: nil
}

{\footnotesize Використовується різноманітний набір інформаційних технологій для компоненти системни при впровадженні системи}


{\footnotesize\bfseries
No: 2\\
Name: sc\_\allowbreak{}29\_\allowbreak{}odp\\
Type: string\\
Default: nil
}

{\footnotesize Визначені компоненти системи, які потребують різноманітного набору інформаційних технологій, що мають бути використані при впровадженні системи}




\subsection{МАСКУВАННЯ ТА ХИБНИЙ НАПРЯМ (SC-30)}
Використовувати [Призначення: визначені організацією методи маскування та хибного напряму] для [Призначення: визначені організацією системи] у [Призначення: визначений організацією період часу], щоб заплутати та ввести в оману зловмисників.

\vspace{10pt}
{\footnotesize\bfseries
No: 1\\
Name: sc\_\allowbreak{}30\_\allowbreak{}01\\
Type: integer\\
Default: 30
}

{\footnotesize Застосовуються методи маскування та хибного напряму для систем протягом періодів часу, щоб заплутати та ввести супротивника в оману. застосування методів}


{\footnotesize\bfseries
No: 2\\
Name: sc\_\allowbreak{}30\_\allowbreak{}odp\_\allowbreak{}01\\
Type: string\\
Default: nil
}

{\footnotesize Визначені методи маскування та хибного напряму, які будуть застосовані для того, щоб заплутати і ввести в оману супротивників, які потенційно можуть націлитися на системи}


{\footnotesize\bfseries
No: 3\\
Name: sc\_\allowbreak{}30\_\allowbreak{}odp\_\allowbreak{}02\\
Type: string\\
Default: nil
}

{\footnotesize Визначені системи, для яких повинні застосовуватися методи маскування та хибного напряму}


{\footnotesize\bfseries
No: 4\\
Name: sc\_\allowbreak{}30\_\allowbreak{}odp\_\allowbreak{}03\\
Type: integer\\
Default: 30
}

{\footnotesize Визначені часові періоди для маскування та хибного напряму}




\subsection{АНАЛІЗ ПРИХОВАНОГО КАНАЛУ (SC-31)}
a. Проводити аналіз прихованого каналу, щоб визначити ті аспекти комунікацій у системі, які володіють потенційними можливостями для реалізації прихованих каналів [Вибір (один або кілька): зберігання; синхронізації].\\ 
b. Оцінювати максимальну пропускну здатність цих каналів.

\vspace{10pt}
\textit{Немає параметрів для цього контролю.}
\vspace{10pt}


\subsection{ПОДІЛ СИСТЕМИ НА ЧАСТИНИ (SC-32)}
Розділити систему на [Призначення: визначені організацією системні компоненти], що розміщені в окремих фізичних доменах або середовищах на основі [Призначення: визначені організацією умови для фізичного поділу компонентів].

\vspace{10pt}
{\footnotesize\bfseries
No: 1\\
Name: sc\_\allowbreak{}32\_\allowbreak{}01\\
Type: string\\
Default: nil
}

{\footnotesize Розділена система на компоненти системи, що знаходяться в окремих ЗНАЧЕННЯ ВИБРАНОГО ПАРАМЕТРА доменах або середовищах на основі обставин для фізичного або логічного поділу компонентів. фізичного або логічного}


{\footnotesize\bfseries
No: 2\\
Name: sc\_\allowbreak{}32\_\allowbreak{}odp\_\allowbreak{}01\\
Type: string\\
Default: nil
}

{\footnotesize Повинні компоненти системи перебувати в окремих фізичних або логічних доменах або середовищах, виходячи з обставин фізичного або логічного поділу компонентів}


{\footnotesize\bfseries
No: 3\\
Name: sc\_\allowbreak{}32\_\allowbreak{}odp\_\allowbreak{}02\\
Type: string\\
Default: nil
}

{\footnotesize Вибрано одне з наступних ЗНАЧЕНЬ ПАРАМЕТРА: \{фізичний; логічний\}}


{\footnotesize\bfseries
No: 4\\
Name: sc\_\allowbreak{}32\_\allowbreak{}odp\_\allowbreak{}03\\
Type: string\\
Default: nil
}

{\footnotesize Визначені обставини для розділення компонентів}




\subsection{ПІДГОТОВКА ЦІЛІСНОСТІ ПЕРЕДАЧІ (SC-33) [Вилучено]}
[Вилучено: включено до SC-8]

\vspace{10pt}
\textit{Немає параметрів для цього контролю.}
\vspace{10pt}


\subsection{НЕЗМІНЮВАНІ ВИКОНАВЧІ ПРОГРАМИ (SC-34)}
У [Призначення: визначені організацією системні компоненти]:\\ 
a. Завантажити та виконати операційне середовище з апаратного носія, що працює в режимі лише для зчитування.\\ 
b. Завантажити та виконати [Призначення: визначені організацією застосунки] з апаратного носія, що працює в режимі лише для зчитування.

\vspace{10pt}
{\footnotesize\bfseries
No: 1\\
Name: sc\_\allowbreak{}34\_\allowbreak{}odp\_\allowbreak{}01\\
Type: string\\
Default: nil
}

{\footnotesize Визначені компоненти системи, для яких операційне середовище та додатки мають завантажуватися та виконуватися з апаратних носіїв, призначених лише для читання}


{\footnotesize\bfseries
No: 2\\
Name: sc\_\allowbreak{}34\_\allowbreak{}odp\_\allowbreak{}02\\
Type: string\\
Default: nil
}

{\footnotesize Визначено додатки, які мають завантажуватися та виконуватися з апаратних носіїв, призначених лише для читання; SC-34a. завантажується та виконується операційне середовище для системних компонентів з апаратного носія, доступного лише для читання; SC-34b. додатки для компонентів системи завантажуються та виконуються з апаратного носія, доступного лише для читання. ПОТЕНЦІЙНІ МЕТОДИ ТА ОБ'ЄКТИ ОЦІНКИ:}




\subsection{РОЗПІЗНАВАННЯ ПРИМАНОК ДЛЯ ЗЛОВМИСНИКІВ (HONEYCLIENT) (SC-35)}
Ввімкнути системні компоненти, які активно намагаються ідентифікувати мережевий шкідливий код та шкідливі вебсайти.

\vspace{10pt}
{\footnotesize\bfseries
No: 1\\
Name: sc\_\allowbreak{}35\_\allowbreak{}01\\
Type: string\\
Default: nil
}

{\footnotesize Ввімкнуті компоненти системи, які активно намагаються ідентифікувати мережевий шкідливий код та шкідливі вебсайти}




\subsection{РОЗПОДІЛЕНА ОБРОБКА ТА ЗБЕРІГАННЯ (SC-36)}
Розподіліть наведені нижче компоненти обробки та зберігання в кількох [Вибір: фізичні локації; логічні домени]: [Призначення: компоненти обробки та зберігання, визначені організацією].

\vspace{10pt}
{\footnotesize\bfseries
No: 1\\
Name: sc\_\allowbreak{}36\_\allowbreak{}odp\_\allowbreak{}01\\
Type: string\\
Default: nil
}

{\footnotesize Потрібно розподіляти компоненти обробки між кількома локаціями/доменами}


{\footnotesize\bfseries
No: 2\\
Name: sc\_\allowbreak{}36\_\allowbreak{}odp\_\allowbreak{}02\\
Type: string\\
Default: nil
}

{\footnotesize Вибрано одне з наступних ЗНАЧЕНЬ ПАРАМЕТРІВ: \{фізичні локації; логічні домени\}}


{\footnotesize\bfseries
No: 3\\
Name: sc\_\allowbreak{}36\_\allowbreak{}odp\_\allowbreak{}03\\
Type: string\\
Default: nil
}

{\footnotesize Потрібно розподіляти компоненти сховища між кількома локаціями/доменами}


{\footnotesize\bfseries
No: 4\\
Name: sc\_\allowbreak{}36\_\allowbreak{}odp\_\allowbreak{}04\\
Type: string\\
Default: nil
}

{\footnotesize Вибрано одне з наступних ЗНАЧЕНЬ ПАРАМЕТРІВ: \{фізичні локації; логічні домени\}}




\subsection{ПОЗАСМУГОВІ КАНАЛИ (SC-37)}
Використовувати [Призначення: визначені організацією позасмугові канали] для фізичного доставлення або електронної передачі [Призначення: визначена організацією інформація, системні компоненти або пристрої] до [Призначення: визначені організацією особи або системи].

\vspace{10pt}
{\footnotesize\bfseries
No: 1\\
Name: sc\_\allowbreak{}37\_\allowbreak{}01\\
Type: string\\
Default: nil
}

{\footnotesize Використовуються позасмугові канали для фізичної доставки або електронної передачі інформації, системних компонентів або пристроїв до осіб або систем}


{\footnotesize\bfseries
No: 2\\
Name: sc\_\allowbreak{}37\_\allowbreak{}odp\_\allowbreak{}02\\
Type: string\\
Default: nil
}

{\footnotesize Визначено інформацію, компоненти системи або пристрої для використання позасмугових каналів для фізичної доставки або електронної передачі}


{\footnotesize\bfseries
No: 3\\
Name: sc\_\allowbreak{}37\_\allowbreak{}odp\_\allowbreak{}03\\
Type: list\\
Default: [\textquotedbl{}admin\textquotedbl{}, \textquotedbl{}security\_\allowbreak{}officer\textquotedbl{}]
}

{\footnotesize Визначені особи або системи, до яких фізична доставка або електронна передача інформації, системних компонентів або пристроїв має бути досягнута за допомогою використання позасмугових каналів}




\subsection{БЕЗПЕКА ОПЕРАЦІЙ (SC-38)}
Впровадити [Призначення: визначені організацією заходи з безпеки операцій] для захисту ключової організаційної інформації протягом усього життєвого циклу розробки системи.

\vspace{10pt}
{\footnotesize\bfseries
No: 1\\
Name: sc\_\allowbreak{}38\_\allowbreak{}01\\
Type: string\\
Default: \textquotedbl{}автоматизований засіб моніторингу\textquotedbl{}
}

{\footnotesize Застосовуються засоби контролю заходів з безпеки операцій для захисту ключової інформації організації протягом життєвого циклу розробки системи}


{\footnotesize\bfseries
No: 2\\
Name: sc\_\allowbreak{}38\_\allowbreak{}odp\\
Type: string\\
Default: \textquotedbl{}автоматизований засіб моніторингу\textquotedbl{}
}

{\footnotesize Визначені засоби контролю заходів з безпеки операцій, які будуть застосовуватися для захисту ключової інформації організації протягом усього життєвого циклу розробки системи}




\subsection{ІЗОЛЯЦІЯ ПРОЦЕСУ (SC-39)}
Підтримувати окремий домен виконання для кожного процесу, що виконується в системі.

\vspace{10pt}
{\footnotesize\bfseries
No: 1\\
Name: sc\_\allowbreak{}39\_\allowbreak{}01\\
Type: string\\
Default: nil
}

{\footnotesize Підтримується окремий домен виконання для кожного процесу, що виконується в системі}




\subsection{ЗАХИСТ БЕЗДРОТОВОГО З'ЄДНАННЯ (SC-40)}
Забезпечити захист зовнішніх і внутрішніх [Призначення: визначені організацією бездротові з'єднання] від [Призначення: визначені організацією типи атак з параметрами сигналів або посилання на джерела для таких атак].

\vspace{10pt}
{\footnotesize\bfseries
No: 1\\
Name: sc\_\allowbreak{}40\_\allowbreak{}01\\
Type: string\\
Default: nil
}

{\footnotesize Захищені зовнішні бездротові з'єднання від типів атак на параметри сигналу або посилання на джерела таких атак. SC-40[02] захищені внутрішні бездротові з'єднання від типів атак на параметри сигналу або посилання на джерела для таких атак}


{\footnotesize\bfseries
No: 2\\
Name: sc\_\allowbreak{}40\_\allowbreak{}odp\_\allowbreak{}01\\
Type: string\\
Default: nil
}

{\footnotesize Потрібно захищати зовнішні бездротові з'єднання від певних типів атак на параметри сигналу}


{\footnotesize\bfseries
No: 3\\
Name: sc\_\allowbreak{}40\_\allowbreak{}odp\_\allowbreak{}02\\
Type: string\\
Default: nil
}

{\footnotesize Визначено типи атак на параметри сигналу або посилання на джерела таких атак, від яких потрібно захищати зовнішні бездротові з'єднання}


{\footnotesize\bfseries
No: 4\\
Name: sc\_\allowbreak{}40\_\allowbreak{}odp\_\allowbreak{}03\\
Type: string\\
Default: nil
}

{\footnotesize Потрібно захищати внутрішні бездротові з'єднання від певних типів атак на параметри сигналу}


{\footnotesize\bfseries
No: 5\\
Name: sc\_\allowbreak{}40\_\allowbreak{}odp\_\allowbreak{}04\\
Type: string\\
Default: nil
}

{\footnotesize Визначені типи атак на параметри сигналу або посилання на джерела таких атак, від яких потрібно захищати внутрішні бездротові з'єднання}




\subsection{ДОСТУП ДО ПОРТІВ ТА ПРИСТРОЇВ ВВЕДЕННЯ, ВИВЕДЕННЯ (SC-41)}
[Вибір: фізично або логічно] відключити або видалити [Призначення: визначені організацією, порти підключення або пристрої введення/виводу] у [Призначення: визначені організацією системи або системні компоненти].

\vspace{10pt}
{\footnotesize\bfseries
No: 1\\
Name: sc\_\allowbreak{}41\_\allowbreak{}odp\_\allowbreak{}01\\
Type: string\\
Default: nil
}

{\footnotesize Визначено порти підключення або пристрої вводу/виводу, які потрібно відключити або видалити}


{\footnotesize\bfseries
No: 2\\
Name: sc\_\allowbreak{}41\_\allowbreak{}odp\_\allowbreak{}02\\
Type: string\\
Default: nil
}

{\footnotesize Вибрано одне з наступних ЗНАЧЕНЬ ПАРАМЕТРІВ: \{фізично; логічно\}}


{\footnotesize\bfseries
No: 3\\
Name: sc\_\allowbreak{}41\_\allowbreak{}odp\_\allowbreak{}03\\
Type: string\\
Default: nil
}

{\footnotesize Визначені системи або компоненти системи з портами підключення або пристроями вводу/виводу, які потрібно відключити або видалити}




\subsection{Захист бездротового з'єднання Доступ до портів і пристроїв введення/виведення Можливості датчика та дані (SC-42)}
a. Заборонити дистанційну активацію можливостей зондування навколишнього середовища в системах організації або компонентах системи за такими виключеннями: [Призначення: визначені організацією виключення, в яких допускається дистанційна активація датчиків].\\ 
b. Забезпечити явну вказівку використання датчика для [Призначення: визначений організацією клас користувачів].

\vspace{10pt}
{\footnotesize\bfseries
No: 1\\
Name: sc\_\allowbreak{}42\_\allowbreak{}01\\
Type: string\\
Default: \textquotedbl{}автоматизований засіб моніторингу\textquotedbl{}
}

{\footnotesize Механізми, що перешкоджають МОЖЛИВОСТІ ДАТЧИКА ТА ДАНІ}


{\footnotesize\bfseries
No: 2\\
Name: sc\_\allowbreak{}42\_\allowbreak{}odp\_\allowbreak{}01\\
Type: string\\
Default: nil
}

{\footnotesize Вибрано одне або більше з наступних ЗНАЧЕНЬ ПАРАМЕТРІВ: \{використання пристроїв, що мають <SC- 42\_\allowbreak{}ODP[02] можливості зондування довкілля> на об'єктах, територіях або системах\}; дистанційна активація можливостей зондування довкілля на організаційних системах або системних компонентах з наступними винятками: <SC-42\_\allowbreak{}ODP[04] винятки, де дозволяється дистанційна активація датчиків\}}


{\footnotesize\bfseries
No: 3\\
Name: sc\_\allowbreak{}42\_\allowbreak{}odp\_\allowbreak{}02\\
Type: string\\
Default: nil
}

{\footnotesize Визначені можливості зондування навколишнього середовища в пристроях (якщо вибрано)}


{\footnotesize\bfseries
No: 4\\
Name: sc\_\allowbreak{}42\_\allowbreak{}odp\_\allowbreak{}03\\
Type: string\\
Default: nil
}

{\footnotesize Визначені об'єкти, зони або системи, на яких заборонено використання пристроїв, що мають можливості зондування навколишнього середовища (якщо вони були обрані)}


{\footnotesize\bfseries
No: 5\\
Name: sc\_\allowbreak{}42\_\allowbreak{}odp\_\allowbreak{}04\\
Type: string\\
Default: nil
}

{\footnotesize Визначено винятки, коли дозволено віддалену активацію датчиків (якщо вибрано)}




\subsection{ОБМЕЖЕННЯ ВИКОРИСТАННЯ (SC-43)}
a. Встановити обмеження на використання та рекомендації щодо впровадження для [Призначення: визначених організацією компонентів системи].\\ 
b. Проводити авторизацію, спостереження та контроль використання таких компонентів у системі.

\vspace{10pt}
{\footnotesize\bfseries
No: 1\\
Name: sc\_\allowbreak{}43\_\allowbreak{}odp\\
Type: string\\
Default: nil
}

{\footnotesize Визначені компоненти, для яких мають бути встановлені обмеження на використання та настанови щодо впровадження; SC-43a. встановлені обмеження на використання та настанови щодо впровадження для компонентів; SC-43b.[01] дозволено системі; SC-43b.[02] здійснюється моніторинг компонентів> в системі; SC-43b.[03] контролюється використання компонентів в системі. використання компонентів використання у <SC-43\_\allowbreak{}ODP}




\subsection{ЕКРАНОВАНІ КАМЕРИ (SC-44)}
Впровадити екрановані камери в [Призначення: визначену організацією систему, компонент системи або місце розташування].

\vspace{10pt}
{\footnotesize\bfseries
No: 1\\
Name: sc\_\allowbreak{}44\_\allowbreak{}01\\
Type: string\\
Default: nil
}

{\footnotesize Використовується в системі <SC-44\_\allowbreak{}ODP, компоненті або місці розташування> застосування екранованої камери системному можливість}


{\footnotesize\bfseries
No: 2\\
Name: sc\_\allowbreak{}44\_\allowbreak{}odp\\
Type: string\\
Default: nil
}

{\footnotesize Визначена система, компонент системи або місце, де має бути застосований потенціал екранованої камери}




\subsection{СИНХРОНІЗАЦІЯ СИСТЕМИ З ЧАСОМ (SC-45)}
Синхронізація системного годинника в системі та компонентах системи і між ними.

\vspace{10pt}
{\footnotesize\bfseries
No: 1\\
Name: sc\_\allowbreak{}45\_\allowbreak{}01\\
Type: integer\\
Default: 30
}

{\footnotesize Синхронізовані системні годинники всередині системи та між системами і системними компонентами}




\subsection{ЗАБЕЗПЕЧЕННЯ ВИКОНАННЯ МІЖДОМЕННОЇ ПОЛІТИКИ (SC-46)}
Впровадити механізм примусового виконання політики [Вибір: фізично; логічно] між фізичними та/або мережевими інтерфейсами для підключених доменів безпеки.

\vspace{10pt}
{\footnotesize\bfseries
No: 1\\
Name: sc\_\allowbreak{}46\_\allowbreak{}odp\\
Type: string\\
Default: nil
}

{\footnotesize Вибрано одне з наступних ЗНАЧЕНЬ ПАРАМЕТРІВ: \{фізично; логічно\}}




\subsection{АЛЬТЕРНАТИВНИЙ ШЛЯХ ЗВ'ЯЗКУ (SC-47)}
Встановіть [Призначення: альтернативні шляхи зв'язку, визначені організацією] для організаційного управління та контролю операцій системи.

\vspace{10pt}
{\footnotesize\bfseries
No: 1\\
Name: sc\_\allowbreak{}47\_\allowbreak{}01\\
Type: string\\
Default: nil
}

{\footnotesize Встановлені альтернативні шляхи зв'язку для системних операцій та контролю операцій системи}


{\footnotesize\bfseries
No: 2\\
Name: sc\_\allowbreak{}47\_\allowbreak{}odp\\
Type: string\\
Default: nil
}

{\footnotesize Визначені альтернативні шляхи зв'язку для системних операцій та контролю операцій системи}




\subsection{ПЕРЕМІЩЕННЯ ДАТЧИКА (SC-48)}
Перенесіть [Призначення: датчики та можливості моніторингу, визначені організацією] до [Призначення: місця, визначені організацією] за таких умов або обставин: [Призначення: умови або обставини, визначені організацією].

\vspace{10pt}
{\footnotesize\bfseries
No: 1\\
Name: sc\_\allowbreak{}48\_\allowbreak{}01\\
Type: string\\
Default: \textquotedbl{}автоматизований засіб моніторингу\textquotedbl{}
}

{\footnotesize Переміщуються датчики і засоби моніторингу до місць розташування за умов або обставин}


{\footnotesize\bfseries
No: 2\\
Name: sc\_\allowbreak{}48\_\allowbreak{}odp\_\allowbreak{}01\\
Type: string\\
Default: nil
}

{\footnotesize Визначені датчики та можливості необхідно перемістити; моніторингу, які}


{\footnotesize\bfseries
No: 3\\
Name: sc\_\allowbreak{}48\_\allowbreak{}odp\_\allowbreak{}02\\
Type: string\\
Default: \textquotedbl{}автоматизований засіб моніторингу\textquotedbl{}
}

{\footnotesize Визначені місця, куди будуть переміщені датчики та засоби моніторингу}


{\footnotesize\bfseries
No: 4\\
Name: sc\_\allowbreak{}48\_\allowbreak{}odp\_\allowbreak{}03\\
Type: list\\
Default: []
}

{\footnotesize Визначені умови або обставини для переміщення датчиків і можливостей моніторингу}




\subsubsection{ДИНАМІЧНО ПЕРЕМІЩУЮТЬСЯ ДО ЗА (SC-48(1))}


\vspace{10pt}
{\footnotesize\bfseries
No: 1\\
Name: sc\_\allowbreak{}48\_\allowbreak{}1\_\allowbreak{}01\\
Type: string\\
Default: \textquotedbl{}автоматизований засіб моніторингу\textquotedbl{}
}

{\footnotesize Датчики та засоби моніторингу динамічно переміщуються до місць розташування за умов або обставин. ПОТЕНЦІЙНІ МЕТОДИ ТА ОБ'ЄКТИ ОЦІНКИ:}




\subsection{ПРИМУСОВЕ АПАРАТНЕ ЗАБЕЗПЕЧЕННЯ ВИКОНАННЯ (SC-49)}
Впровадити механізми апаратного поділу та застосування політики між [Призначення: домени безпеки, визначені організацією].

\vspace{10pt}
{\footnotesize\bfseries
No: 1\\
Name: sc\_\allowbreak{}49\_\allowbreak{}01\\
Type: list\\
Default: [\textquotedbl{}default\_\allowbreak{}deny\_\allowbreak{}rule\textquotedbl{}, \textquotedbl{}abac\_\allowbreak{}rule\_\allowbreak{}1\textquotedbl{}]
}

{\footnotesize Впроваджено механізми апаратного розділення та застосування політик між доменами безпеки}


{\footnotesize\bfseries
No: 2\\
Name: sc\_\allowbreak{}49\_\allowbreak{}odp\\
Type: list\\
Default: [\textquotedbl{}default\_\allowbreak{}deny\_\allowbreak{}rule\textquotedbl{}, \textquotedbl{}abac\_\allowbreak{}rule\_\allowbreak{}1\textquotedbl{}]
}

{\footnotesize Визначені домени безпеки, які потребують апаратного розділення та механізмів забезпечення дотримання політики}




\subsection{ПРИМУСОВЕ ПРОГРАМНЕ ЗАБЕЗПЕЧЕННЯ ВИКОНАННЯ (SC-50)}
Впровадити програмне розділення та механізми застосування політики між [Призначення: домени безпеки, визначені організацією].

\vspace{10pt}
{\footnotesize\bfseries
No: 1\\
Name: sc\_\allowbreak{}50\_\allowbreak{}01\\
Type: list\\
Default: [\textquotedbl{}default\_\allowbreak{}deny\_\allowbreak{}rule\textquotedbl{}, \textquotedbl{}abac\_\allowbreak{}rule\_\allowbreak{}1\textquotedbl{}]
}

{\footnotesize Впроваджено програмне розділення та механізми застосування політик між доменами безпеки. ПОТЕНЦІЙНІ МЕТОДИ ТА ОБ'ЄКТИ ОЦІНКИ:}


{\footnotesize\bfseries
No: 2\\
Name: sc\_\allowbreak{}50\_\allowbreak{}odp\\
Type: list\\
Default: [\textquotedbl{}default\_\allowbreak{}deny\_\allowbreak{}rule\textquotedbl{}, \textquotedbl{}abac\_\allowbreak{}rule\_\allowbreak{}1\textquotedbl{}]
}

{\footnotesize Визначені домени безпеки, які потребують програмного розділення та механізмів забезпечення дотримання політик}




\subsection{АПАРАТНИЙ ЗАХИСТ (SC-51)}
a. Перевіряти правильність роботи [Призначення: визначені організацією функції безпеки та приватності].\\ 
b. Виконувати перевірку [Вибір (один або кілька): [Призначення: визначені організацією системні перехідні стани]; за командою користувача з відповідними повноваженнями; [Призначення: визначена організацією частота]].\\ 
c. Повідомляти [Призначення: визначені організацією персонал або посадові особи] про невдалі перевірки безпеки та приватності.\\ 
d. [Вибір (один або кілька): Вимкнути систему; Перезапустити систему; [Призначення: визначені організацією альтернативні дії]], коли виявляються аномалії.

\vspace{10pt}
{\footnotesize\bfseries
No: 1\\
Name: sc\_\allowbreak{}51\_\allowbreak{}odp\_\allowbreak{}01\\
Type: string\\
Default: nil
}

{\footnotesize Визначено компоненти системної прошивки, потребують апаратного захисту від запису; які}


{\footnotesize\bfseries
No: 2\\
Name: sc\_\allowbreak{}51\_\allowbreak{}odp\_\allowbreak{}02\\
Type: list\\
Default: [\textquotedbl{}admin\textquotedbl{}, \textquotedbl{}security\_\allowbreak{}officer\textquotedbl{}]
}

{\footnotesize Визначені уповноважені особи, які повинні виконувати процедури вимкнення та повторного ввімкнення апаратного захисту від запису; SC-51a. використовується апаратний захист від запису для компонентів мікропрограми системи; SC-51b.[01] впроваджено спеціальні процедури для уповноважених осіб для ручного вимкнення апаратного захисту від запису для модифікацій мікропрограми; SC-51b.[02] реалізовано спеціальні процедури для уповноважених осіб для повторного увімкнення захисту від запису перед поверненням до робочого режиму}




\section{SI}
Клас заходів захисту SI --- ЦІЛІСНІСТЬ СИСТЕМИ ТА ІНФОРМАЦІЇ

Цей клас спрямований на захист системи від шкідливого коду, виявлення вразливостей та запобігання несанкціонованим змінам інформації.

\textbf{Перелік заходів захисту:} Політика і процедури цілісності інформації (SI-1); Виправлення дефектів (SI-2); Централізоване управління (SI-2(1)) [Вилучено]; Автоматизоване виправлення дефектів (SI-2(2)); Час для усунення дефектів та орієнтири для коригувальних дій (SI-2(3)); Автоматичні засоби управління виправленнями (SI-2(4)); Виправлення дефектів автоматичне оновлення програмного забезпечення та вбудованого програмного забезпечення (SI-2(5)); Видалення попередніх версій програмного забезпечення та вбудованого програмного забезпечення (SI-2(6)); Захист від шкідливого коду (SI-3); Централізоване управління (SI-3(1)) [Вилучено]; Автоматичні оновлення (SI-3(2)) [Вилучено]; Захист від шкідливого користувачі (SI-3(3)) [Вилучено]; Захист від шкідливого коду привілейованими користувачами (SI-3(4)); Захист від шкідливого зберігання даних (SI-3(5)) [Вилучено]; Тестування та верифікація (SI-3(6)); Виявлення без підпису (SI-3(7)) [Вилучено]; Виявлення неавторизованих команд (SI-3(8)); Автентифікація віддалених команд (SI-3(9)) [Вилучено]; Аналіз шкідливого коду (SI-3(10)); Моніторинг системи (SI-4); МОНІТОРИНГ СИСТЕМИ ВИЯВЛЕННЯ ВТОРГНЕНЬ (IDS) (SI-4(1)); Автоматизовані механізми аналізу в реальному часі (SI-4(2)); Моніторинг системи механізми інтеграції (SI-4(3)); Моніторинг системи комунікацій (SI-4(4)); Системні сповіщення (SI-4(5)); Заборона для непривілейованих користувачів (SI-4(6)) [Вилучено]; Моніторинг системи підозрілі події (SI-4(7)); Захист інформації моніторингу (SI-4(8)) [Вилучено]; Тестування засобів і механізмів моніторингу (SI-4(9)); Моніторинг комунікацій (SI-4(10)); Моніторинг системи комунікацій (SI-4(11)); Моніторинг системи автоматизовані сповіщення (SI-4(12)); Аналіз трафіку та шаблонів подій (SI-4(13)); Моніторинг вторгнення (SI-4(14)); Перехід від бездротового зв'язку до провідних мереж (SI-4(15)); Моніторинг моніторингу (SI-4(16)); Моніторинг обізнаність (SI-4(17)); Аналіз трафіку та прихованої ексфільтрації (SI-4(18)); Особи, які представляють більший ризик (SI-4(19)); Привілейовані користувачі (SI-4(20)); Випробувальні терміни (SI-4(21)); Несанкціоновані послуги мережі (SI-4(22)); Пристрої на основі хоста (SI-4(23)); Індикатори компрометації (SI-4(24)); Аналіз мережевого трафіку (SI-4(25)); Попередження, рекомендації та директиви з безпеки (SI-5); Попередження, рекомендації та директиви з безпеки автоматичні попередження та рекомендації (SI-5(1)); Перевірка функцій безпеки та приватності (SI-6); Сповіщення про неуспішне проходження тестів з безпеки (SI-6(1)) [Вилучено]; Перевірка безпеки та функцій приватності автоматизована підтримка розподіленого тестування (SI-6(2)); Перевірка безпеки та функцій приватності повідомлення про результати перевірки (SI-6(3)); Цілісність програмного забезпечення, програмного забезпечення та інформації (SI-7); Перевірка цілісності (SI-7(1)); Автоматичні сповіщення про порушення цілісності (SI-7(2)); Інструменти цілісності з централізованим управлінням (SI-7(3)); Цілісність програмного забезпечення, вбудованого програмного забезпечення та інформації -- пакування з індикацією ознак її несанкціонованого розкриття (SI-7(4)) [Вилучено]; Автоматичні відповіді про порушення цілісності (SI-7(5)); Цілісність програмного забезпечення, вбудованого програмного забезпечення та інформації криптографічний захист (SI-7(6)); Інтеграція виявлення і реагування (SI-7(7)); Аудит важливих подій (SI-7(8)); Перевірка процесу завантаження (SI-7(9)); Захист завантажувального вбудованого програмного забезпечення (SI-7(10)); Обмежене середовище з обмеженими привілеями (SI-7(11)) [Вилучено]; Перевірка цілісності (SI-7(12)); Виконання коду в захищених середовищах (SI-7(13)) [Вилучено]; Двійковий або машинно-виконуваний код (SI-7(14)) [Вилучено]; Цілісність програмного забезпечення, вбудованого програмного забезпечення та інформації автентифікація коду (SI-7(15)); Термін виконання процесу без нагляду (SI-7(16)); Цілісність програмного забезпечення, вбудованого програмного забезпечення та інформації -- самозахист програм від самовільного виконання (SI-7(17)); Захист від спаму (SI-8); Централізоване управління (SI-8(1)) [Вилучено]; Автоматичні оновлення (SI-8(2)); Безперервне навчання (SI-8(3)); Обмеження на введення інформації (SI-9); Перевірка вводу інформації (SI-10); Перевірка вводу перевизначення (SI-10(1)); Перевірка вводу інформації помилок (SI-10(2)); Передбачувана поведінка (SI-10(3)); Часові взаємодії (SI-10(4)); Обмеження вхідних даних довіреними джерелами і затвердженими форматами (SI-10(5)); Профілактика вводу даних (SI-10(6)); Обробка помилок (SI-11); Управління та збереження інформації (SI-12); Управління та збереження інформації елементів персональних даних (SI-12(1)); Мінімізація використання персональних даних під час тестування, навчання та дослідженні (SI-12(2)); Управління інформації (SI-12(3)); Передбачуване запобігання збоїв (SI-13); Відповідальність за передачу функцій компонентів (SI-13(1)); Термін виконання процесу без нагляду (SI-13(2)) [Вилучено]; Запобігання передбачуваних функцій компонентів (SI-13(3)); Запобігання передбачуваних збоїв резервних компонентів та оповіщення (SI-13(4)); Запобігання передбачуваних аварійного перемикання (SI-13(5)); Нестійкість (SI-14); Оновлення з надійних джерел (SI-14(1)); Нестійка інформація (SI-14(2)); Нестійкі підключення (SI-14(3)); Фільтрація вихідних даних (SI-15); Захист пам'яті (SI-16); Відмовостійкі процедури (SI-17); Операції забезпечення якості даних (SI-18); Операції забезпечення підтримка (SI-18(1)); Тегування даних (SI-18(2)); Збирання (SI-18(3)); Операції запити (SI-18(4)); Повідомлення про виправлення чи видалення (SI-18(5)); Деідентифікація (SI-19); Збір (SI-19(1)); Архівація (SI-19(2)); Видалення (SI-19(3)); Видалення, маскування, шифрування, хешування або заміна прямих ідентифікаторів (SI-19(4)); Контроль статистичного розкриття (SI-19(5)); Диференційована конфіденційність (SI-19(6)); Перевірене програмне забезпечення (SI-19(7)); Мотивований порушник (SI-19(8)); Псування (SI-20); Оновлення інформації (SI-21); Різновиди інформації (SI-22); Фрагментація інформації (SI-23).

\subsection{ПОЛІТИКА І ПРОЦЕДУРИ ЦІЛІСНОСТІ ІНФОРМАЦІЇ (SI-1)}


\vspace{10pt}
{\footnotesize\bfseries
No: 1\\
Name: si\_\allowbreak{}1\_\allowbreak{}odp\_\allowbreak{}01\\
Type: list\\
Default: [\textquotedbl{}admin\textquotedbl{}, \textquotedbl{}security\_\allowbreak{}officer\textquotedbl{}]
}

{\footnotesize Визначено персонал або ролі, до яких має бути доведена політика цілісності системи та інформації}


{\footnotesize\bfseries
No: 2\\
Name: si\_\allowbreak{}1\_\allowbreak{}odp\_\allowbreak{}02\\
Type: list\\
Default: [\textquotedbl{}admin\textquotedbl{}, \textquotedbl{}security\_\allowbreak{}officer\textquotedbl{}]
}

{\footnotesize Визначено персонал або ролі, на які поширюються процедури цілісності системи та інформації }


{\footnotesize\bfseries
No: 3\\
Name: si\_\allowbreak{}1\_\allowbreak{}odp\_\allowbreak{}03\\
Type: string\\
Default: nil
}

{\footnotesize Вибрано одне або декілька з наступних ЗНАЧЕНЬ ПАРАМЕТРІВ: \{рівень організації; рівень місії/бізнеспроцесу; рівень системи\}}


{\footnotesize\bfseries
No: 4\\
Name: si\_\allowbreak{}1\_\allowbreak{}odp\_\allowbreak{}04\\
Type: list\\
Default: [\textquotedbl{}admin\textquotedbl{}, \textquotedbl{}security\_\allowbreak{}officer\textquotedbl{}]
}

{\footnotesize Визначено посадову особу, відповідальну за управління системою та політикою і процедурами цілісності інформації}


{\footnotesize\bfseries
No: 5\\
Name: si\_\allowbreak{}1\_\allowbreak{}odp\_\allowbreak{}05\\
Type: list\\
Default: [\textquotedbl{}default\_\allowbreak{}deny\_\allowbreak{}rule\textquotedbl{}, \textquotedbl{}abac\_\allowbreak{}rule\_\allowbreak{}1\textquotedbl{}]
}

{\footnotesize Визначено періодичність перегляду та оновлення поточної політики цілісності системи та інформації}


{\footnotesize\bfseries
No: 6\\
Name: si\_\allowbreak{}1\_\allowbreak{}odp\_\allowbreak{}06\\
Type: list\\
Default: [\textquotedbl{}default\_\allowbreak{}deny\_\allowbreak{}rule\textquotedbl{}, \textquotedbl{}abac\_\allowbreak{}rule\_\allowbreak{}1\textquotedbl{}]
}

{\footnotesize Є події, які вимагають перегляду та оновлення поточної політики цілісності системи та інформації}


{\footnotesize\bfseries
No: 7\\
Name: si\_\allowbreak{}1\_\allowbreak{}odp\_\allowbreak{}07\\
Type: integer\\
Default: 30
}

{\footnotesize Визначено частоту, з якою переглядаються та оновлюються поточні цілісності системи та інформації}




\subsection{ВИПРАВЛЕННЯ ДЕФЕКТІВ (SI-2)}


\vspace{10pt}
{\footnotesize\bfseries
No: 1\\
Name: si\_\allowbreak{}2\_\allowbreak{}odp\\
Type: integer\\
Default: 30
}

{\footnotesize Визначено період часу, протягом якого необхідно встановити оновлення програмного забезпечення, пов'язані з безпекою, після виходу оновлень; SI-02a.[01] виявлено недоліки системи; SI-02a.[02] повідомляється про недоліки системи; SI-02a.[03] виправлені недоліки системи; SI-02b.[01] перевіряються оновлення програмного забезпечення, пов'язані з усуненням недоліків, на ефективність перед встановленням; SI-02b.[02] перевіряються оновлення програмного забезпечення, пов'язані з виправленням дефектів, на наявність потенційних побічних ефектів перед встановленням; SI-02b.[03] перевіряються оновлення прошивки, пов'язані з усуненням недоліків, на ефективність перед встановленням; SI-02b.[04] перевіряються оновлення прошивки, пов'язані з усуненням недоліків, на наявність потенційних побічних ефектів перед встановленням; SI-02c.[01] встановлено оновлення програмного забезпечення, що стосуються безпеки, протягом часовий проміжок з моменту випуску оновлень; SI-02c.[02] встановлено оновлення мікропрограми, що стосуються безпеки, протягом часового періоду з моменту випуску оновлень; SI-02d. включено відновлення порушених прав у процес управління організаційною конфігурацією}




\subsubsection{ЦЕНТРАЛІЗОВАНЕ УПРАВЛІННЯ (SI-2(1)) [Вилучено]}
[Вилучено: перенесено до PL-09]

\vspace{10pt}
\textit{Немає параметрів для цього контролю.}
\vspace{10pt}


\subsubsection{АВТОМАТИЗОВАНЕ ВИПРАВЛЕННЯ ДЕФЕКТІВ (SI-2(2))}


\vspace{10pt}
{\footnotesize\bfseries
No: 1\\
Name: si\_\allowbreak{}2\_\allowbreak{}2\_\allowbreak{}01\\
Type: integer\\
Default: 30
}

{\footnotesize Встановлені на компонентах системи відповідні оновлення програмного забезпечення та мікропрограми, що стосуються безпеки, з частотою за допомогою автоматизованих механізмів}




\subsubsection{ЧАС ДЛЯ УСУНЕННЯ ДЕФЕКТІВ ТА ОРІЄНТИРИ ДЛЯ КОРИГУВАЛЬНИХ ДІЙ (SI-2(3))}


\vspace{10pt}
{\footnotesize\bfseries
No: 1\\
Name: si\_\allowbreak{}2\_\allowbreak{}3\_\allowbreak{}a\\
Type: integer\\
Default: 30
}

{\footnotesize Вимірюється час між виявленням дефекту та його усуненням; SI-02(03)(b) були встановлені орієнтири для вжиття коригувальних дій}


{\footnotesize\bfseries
No: 2\\
Name: si\_\allowbreak{}2\_\allowbreak{}3\_\allowbreak{}odp\\
Type: string\\
Default: nil
}

{\footnotesize Визначені контрольні коригувальних заходів; показники для вжиття}




\subsubsection{АВТОМАТИЧНІ ЗАСОБИ УПРАВЛІННЯ ВИПРАВЛЕННЯМИ (SI-2(4))}


\vspace{10pt}
{\footnotesize\bfseries
No: 1\\
Name: si\_\allowbreak{}2\_\allowbreak{}4\_\allowbreak{}odp\\
Type: string\\
Default: \textquotedbl{}автоматизований засіб моніторингу\textquotedbl{}
}

{\footnotesize Визначені компоненти системи, які потребують автоматизованих інструментів управління виправленнями для полегшення усунення дефектів; SI-02(04)] застосовуються автоматизовані засоби управління виправленнями для полегшення виправлення недоліків у компонентах}




\subsubsection{ВИПРАВЛЕННЯ ДЕФЕКТІВ АВТОМАТИЧНЕ ОНОВЛЕННЯ ПРОГРАМНОГО ЗАБЕЗПЕЧЕННЯ ТА ВБУДОВАНОГО ПРОГРАМНОГО ЗАБЕЗПЕЧЕННЯ (SI-2(5))}


\vspace{10pt}
{\footnotesize\bfseries
No: 1\\
Name: si\_\allowbreak{}2\_\allowbreak{}5\_\allowbreak{}01\\
Type: string\\
Default: nil
}

{\footnotesize SI-02(05) \_\allowbreak{}ODP[01] оновлення програмного забезпечення та мікропрограми, що стосуються безпеки>, встановлено автоматично до <SI-02(05) \_\allowbreak{}ODP[02] компонентів системи>}




\subsubsection{ВИДАЛЕННЯ ПОПЕРЕДНІХ ВЕРСІЙ ПРОГРАМНОГО ЗАБЕЗПЕЧЕННЯ ТА ВБУДОВАНОГО ПРОГРАМНОГО ЗАБЕЗПЕЧЕННЯ (SI-2(6))}


\vspace{10pt}
{\footnotesize\bfseries
No: 1\\
Name: si\_\allowbreak{}2\_\allowbreak{}6\_\allowbreak{}01\\
Type: string\\
Default: nil
}

{\footnotesize Видаляються попередні версії програмне забезпечення та компоненти мікропрограми після встановлення оновлених версій}


{\footnotesize\bfseries
No: 2\\
Name: si\_\allowbreak{}2\_\allowbreak{}6\_\allowbreak{}odp\\
Type: string\\
Default: nil
}

{\footnotesize Потрібно видаляти компоненти програмного забезпечення та мікропрограми після встановлення оновлених версій}




\subsection{ЗАХИСТ ВІД ШКІДЛИВОГО КОДУ (SI-3)}


\vspace{10pt}
{\footnotesize\bfseries
No: 1\\
Name: si\_\allowbreak{}3\_\allowbreak{}odp\_\allowbreak{}01\\
Type: string\\
Default: nil
}

{\footnotesize Вибрано одне або декілька з наступних ЗНАЧЕНЬ ПАРАМЕТРІВ: \{підписаний; непідписаний\}}


{\footnotesize\bfseries
No: 2\\
Name: si\_\allowbreak{}3\_\allowbreak{}odp\_\allowbreak{}02\\
Type: integer\\
Default: 30
}

{\footnotesize Визначено частоту, з якою механізми шкідливого коду виконують сканування}


{\footnotesize\bfseries
No: 3\\
Name: si\_\allowbreak{}3\_\allowbreak{}odp\_\allowbreak{}03\\
Type: string\\
Default: nil
}

{\footnotesize Вибрано одне або декілька з наступних ЗНАЧЕНЬ ПАРАМЕТРІВ: \{кінцева точка; точки входу та виходу з мережі\}}


{\footnotesize\bfseries
No: 4\\
Name: si\_\allowbreak{}3\_\allowbreak{}odp\_\allowbreak{}04\\
Type: string\\
Default: nil
}

{\footnotesize Вибрано одне або декілька з наступних ЗНАЧЕНЬ ПАРАМЕТРІВ: \{block malicious code; quarantitne malicious code; take action\}}


{\footnotesize\bfseries
No: 5\\
Name: si\_\allowbreak{}3\_\allowbreak{}odp\_\allowbreak{}05\\
Type: list\\
Default: [\textquotedbl{}login\textquotedbl{}, \textquotedbl{}logout\textquotedbl{}, \textquotedbl{}failed\_\allowbreak{}attempt\textquotedbl{}]
}

{\footnotesize Визначено дії, яких слід вжити у відповідь на виявлення шкідливого коду (якщо вибрано)}




\subsubsection{ЦЕНТРАЛІЗОВАНЕ УПРАВЛІННЯ (SI-3(1)) [Вилучено]}
[Вилучено: включено до PL-9]

\vspace{10pt}
\textit{Немає параметрів для цього контролю.}
\vspace{10pt}


\subsubsection{АВТОМАТИЧНІ ОНОВЛЕННЯ (SI-3(2)) [Вилучено]}
[Вилучено: включено до SI-03]

\vspace{10pt}
\textit{Немає параметрів для цього контролю.}
\vspace{10pt}


\subsubsection{ЗАХИСТ ВІД ШКІДЛИВОГО КОРИСТУВАЧІ (SI-3(3)) [Вилучено]}
[Вилучено: включено до AC-6(10)]

\vspace{10pt}
\textit{Немає параметрів для цього контролю.}
\vspace{10pt}


\subsubsection{ЗАХИСТ ВІД ШКІДЛИВОГО КОДУ ПРИВІЛЕЙОВАНИМИ КОРИСТУВАЧАМИ (SI-3(4))}


\vspace{10pt}
{\footnotesize\bfseries
No: 1\\
Name: si\_\allowbreak{}3\_\allowbreak{}4\_\allowbreak{}01\\
Type: string\\
Default: \textquotedbl{}автоматизований засіб моніторингу\textquotedbl{}
}

{\footnotesize Оновлюються механізми захисту від шкідливого коду лише за вказівкою привілейованого користувача}




\subsubsection{ЗАХИСТ ВІД ШКІДЛИВОГО ЗБЕРІГАННЯ ДАНИХ (SI-3(5)) [Вилучено]}
[Вилучено: включено до MP-7]

\vspace{10pt}
\textit{Немає параметрів для цього контролю.}
\vspace{10pt}


\subsubsection{ТЕСТУВАННЯ ТА ВЕРИФІКАЦІЯ (SI-3(6))}


\vspace{10pt}
{\footnotesize\bfseries
No: 1\\
Name: si\_\allowbreak{}3\_\allowbreak{}6\_\allowbreak{}a\\
Type: string\\
Default: \textquotedbl{}щорічно\textquotedbl{}
}

{\footnotesize Перевіряються механізми захисту від шкідливого коду частота шляхом введення в систему відомого доброякісного коду}


{\footnotesize\bfseries
No: 2\\
Name: si\_\allowbreak{}3\_\allowbreak{}6\_\allowbreak{}b\_\allowbreak{}01\\
Type: string\\
Default: nil
}

{\footnotesize Відбувається виявлення (доброякісний тест) коду}


{\footnotesize\bfseries
No: 3\\
Name: si\_\allowbreak{}3\_\allowbreak{}6\_\allowbreak{}b\_\allowbreak{}02\\
Type: string\\
Default: nil
}

{\footnotesize Відбувається відповідне звітування про інцидент}


{\footnotesize\bfseries
No: 4\\
Name: si\_\allowbreak{}3\_\allowbreak{}6\_\allowbreak{}odp\\
Type: string\\
Default: \textquotedbl{}щорічно\textquotedbl{}
}

{\footnotesize Визначено періодичність тестування механізмів захисту від зловмисного коду}




\subsubsection{ВИЯВЛЕННЯ БЕЗ ПІДПИСУ (SI-3(7)) [Вилучено]}
[Вилучено: включено до SI-3]

\vspace{10pt}
\textit{Немає параметрів для цього контролю.}
\vspace{10pt}


\subsubsection{ВИЯВЛЕННЯ НЕАВТОРИЗОВАНИХ КОМАНД (SI-3(8))}


\vspace{10pt}
{\footnotesize\bfseries
No: 1\\
Name: si\_\allowbreak{}3\_\allowbreak{}8\_\allowbreak{}a\\
Type: string\\
Default: nil
}

{\footnotesize Виявлено <SI-03(08) \_\allowbreak{}ODP[01] неавторизовані команди операційної системи> через інтерфейс прикладного програмування ядра на <SI-03(08) \_\allowbreak{}ODP[02] апаратних компонентах системи>}




\subsubsection{АВТЕНТИФІКАЦІЯ ВІДДАЛЕНИХ КОМАНД (SI-3(9)) [Вилучено]}
[Вилучено: включено до AC-17(10)]

\vspace{10pt}
\textit{Немає параметрів для цього контролю.}
\vspace{10pt}


\subsubsection{АНАЛІЗ ШКІДЛИВОГО КОДУ (SI-3(10))}


\vspace{10pt}
{\footnotesize\bfseries
No: 1\\
Name: si\_\allowbreak{}3\_\allowbreak{}10\_\allowbreak{}a\\
Type: string\\
Default: nil
}

{\footnotesize Використовуються інструменти та методи для аналізу характеристик та поведінки шкідливого коду}


{\footnotesize\bfseries
No: 2\\
Name: si\_\allowbreak{}3\_\allowbreak{}10\_\allowbreak{}b\_\allowbreak{}01\\
Type: string\\
Default: nil
}

{\footnotesize Включені результати аналізу шкідливого коду в організаційні процеси реагування на інциденти}


{\footnotesize\bfseries
No: 3\\
Name: si\_\allowbreak{}3\_\allowbreak{}10\_\allowbreak{}b\_\allowbreak{}02\\
Type: string\\
Default: nil
}

{\footnotesize Включені результати аналізу шкідливого коду в організаційні процеси виправлення недоліків}


{\footnotesize\bfseries
No: 4\\
Name: si\_\allowbreak{}3\_\allowbreak{}10\_\allowbreak{}odp\\
Type: string\\
Default: nil
}

{\footnotesize Визначені інструменти та методи, які будуть використовуватися для аналізу характеристик та поведінки шкідливого коду}




\subsection{МОНІТОРИНГ СИСТЕМИ (SI-4)}


\vspace{10pt}
{\footnotesize\bfseries
No: 1\\
Name: si\_\allowbreak{}4\_\allowbreak{}odp\_\allowbreak{}02\\
Type: list\\
Default: [\textquotedbl{}admin\textquotedbl{}, \textquotedbl{}security\_\allowbreak{}officer\textquotedbl{}]
}

{\footnotesize Визначені методи та способи, що використовуються для виявлення несанкціонованого використання системи}


{\footnotesize\bfseries
No: 2\\
Name: si\_\allowbreak{}4\_\allowbreak{}odp\_\allowbreak{}03\\
Type: list\\
Default: [\textquotedbl{}admin\textquotedbl{}, \textquotedbl{}security\_\allowbreak{}officer\textquotedbl{}]
}

{\footnotesize Визначена інформація про моніторинг системи, яка повинна надаватися персоналу або ролям}


{\footnotesize\bfseries
No: 3\\
Name: si\_\allowbreak{}4\_\allowbreak{}odp\_\allowbreak{}04\\
Type: list\\
Default: [\textquotedbl{}admin\textquotedbl{}, \textquotedbl{}security\_\allowbreak{}officer\textquotedbl{}]
}

{\footnotesize Визначено персонал або ролі, яким має надаватися інформація про моніторинг системи}


{\footnotesize\bfseries
No: 4\\
Name: si\_\allowbreak{}4\_\allowbreak{}odp\_\allowbreak{}05\\
Type: string\\
Default: \textquotedbl{}щорічно\textquotedbl{}
}

{\footnotesize Вибрано одне або декілька з наступних ЗНАЧЕНЬ ПАРАМЕТРІВ: \{за потребою; частота\}}




\subsubsection{МОНІТОРИНГ СИСТЕМИ ВИЯВЛЕННЯ ВТОРГНЕНЬ (IDS) (SI-4(1))}


\vspace{10pt}
{\footnotesize\bfseries
No: 1\\
Name: si\_\allowbreak{}4\_\allowbreak{}1\_\allowbreak{}01\\
Type: string\\
Default: \textquotedbl{}автоматизований засіб моніторингу\textquotedbl{}
}

{\footnotesize Підключені окремі засоби виявлення вторгнень загальносистемної системи виявлення вторгнень; до}


{\footnotesize\bfseries
No: 2\\
Name: si\_\allowbreak{}4\_\allowbreak{}1\_\allowbreak{}02\\
Type: string\\
Default: nil
}

{\footnotesize Об'єднані окремі інструменти виявлення вторгнень загальносистемну систему виявлення вторгнень. у}




\subsubsection{АВТОМАТИЗОВАНІ МЕХАНІЗМИ АНАЛІЗУ В РЕАЛЬНОМУ ЧАСІ (SI-4(2))}


\vspace{10pt}
\textit{Немає параметрів для цього контролю.}
\vspace{10pt}


\subsubsection{МОНІТОРИНГ СИСТЕМИ МЕХАНІЗМИ ІНТЕГРАЦІЇ (SI-4(3))}


\vspace{10pt}
{\footnotesize\bfseries
No: 1\\
Name: si\_\allowbreak{}4\_\allowbreak{}3\_\allowbreak{}01\\
Type: string\\
Default: \textquotedbl{}автоматизований засіб моніторингу\textquotedbl{}
}

{\footnotesize Використовуються автоматизовані інструменти та механізми для інтеграції інструментів та механізмів виявлення вторгнень у механізми контролю доступу}


{\footnotesize\bfseries
No: 2\\
Name: si\_\allowbreak{}4\_\allowbreak{}3\_\allowbreak{}02\\
Type: string\\
Default: \textquotedbl{}автоматизований засіб моніторингу\textquotedbl{}
}

{\footnotesize Використовуються автоматизовані інструменти та механізми для інтеграції інструментів та механізмів виявлення вторгнень у механізми контролю потоків}




\subsubsection{МОНІТОРИНГ СИСТЕМИ КОМУНІКАЦІЙ (SI-4(4))}


\vspace{10pt}
{\footnotesize\bfseries
No: 1\\
Name: si\_\allowbreak{}4\_\allowbreak{}4\_\allowbreak{}a\_\allowbreak{}01\\
Type: list\\
Default: []
}

{\footnotesize Визначені критерії незвичної або несанкціонованої діяльності або умови для вхідного трафіку зв'язку}


{\footnotesize\bfseries
No: 2\\
Name: si\_\allowbreak{}4\_\allowbreak{}4\_\allowbreak{}a\_\allowbreak{}02\\
Type: list\\
Default: []
}

{\footnotesize Визначені критерії незвичайної або несанкціонованої діяльності або умови для вихідного трафіку зв'язку}


{\footnotesize\bfseries
No: 3\\
Name: si\_\allowbreak{}4\_\allowbreak{}4\_\allowbreak{}b\_\allowbreak{}01\\
Type: string\\
Default: \textquotedbl{}щорічно\textquotedbl{}
}

{\footnotesize Здійснюється моніторинг вхідного комунікаційного трафіку частота на предмет незвичних або несанкціонованих дій або умов}


{\footnotesize\bfseries
No: 4\\
Name: si\_\allowbreak{}4\_\allowbreak{}4\_\allowbreak{}b\_\allowbreak{}02\\
Type: string\\
Default: \textquotedbl{}щорічно\textquotedbl{}
}

{\footnotesize Контролюється вихідний трафік зв'язку частота на предмет незвичних або несанкціонованих дій або умов. вихідного виявлення}




\subsubsection{СИСТЕМНІ СПОВІЩЕННЯ (SI-4(5))}


\vspace{10pt}
{\footnotesize\bfseries
No: 1\\
Name: si\_\allowbreak{}4\_\allowbreak{}5\_\allowbreak{}01\\
Type: list\\
Default: [\textquotedbl{}admin\textquotedbl{}, \textquotedbl{}security\_\allowbreak{}officer\textquotedbl{}]
}

{\footnotesize Відбувається оповіщення <SI-04(05) \_\allowbreak{}ODP[01] персоналу або ролей> при виникненні згенерованих системою <SI-04(05) \_\allowbreak{}ODP[02] індикаторів компрометації>}




\subsubsection{ЗАБОРОНА ДЛЯ НЕПРИВІЛЕЙОВАНИХ КОРИСТУВАЧІВ (SI-4(6)) [Вилучено]}
[Вилучено: включено до AC-6(10)]

\vspace{10pt}
\textit{Немає параметрів для цього контролю.}
\vspace{10pt}


\subsubsection{МОНІТОРИНГ СИСТЕМИ ПІДОЗРІЛІ ПОДІЇ (SI-4(7))}


\vspace{10pt}
{\footnotesize\bfseries
No: 1\\
Name: si\_\allowbreak{}4\_\allowbreak{}7\_\allowbreak{}a\\
Type: list\\
Default: [\textquotedbl{}admin\textquotedbl{}, \textquotedbl{}security\_\allowbreak{}officer\textquotedbl{}]
}

{\footnotesize Повідомляється персонал з реагування на інциденти про виявлені підозрілі події}


{\footnotesize\bfseries
No: 2\\
Name: si\_\allowbreak{}4\_\allowbreak{}7\_\allowbreak{}b\\
Type: list\\
Default: [\textquotedbl{}login\textquotedbl{}, \textquotedbl{}logout\textquotedbl{}, \textquotedbl{}failed\_\allowbreak{}attempt\textquotedbl{}]
}

{\footnotesize Вживаються найменш руйнівні дії при виявленні підозрілих подій}




\subsubsection{ЗАХИСТ ІНФОРМАЦІЇ МОНІТОРИНГУ (SI-4(8)) [Вилучено]}
[Вилучено: включено до SI-4]

\vspace{10pt}
\textit{Немає параметрів для цього контролю.}
\vspace{10pt}


\subsubsection{ТЕСТУВАННЯ ЗАСОБІВ І МЕХАНІЗМІВ МОНІТОРИНГУ (SI-4(9))}


\vspace{10pt}
{\footnotesize\bfseries
No: 1\\
Name: si\_\allowbreak{}4\_\allowbreak{}9\_\allowbreak{}01\\
Type: string\\
Default: \textquotedbl{}щорічно\textquotedbl{}
}

{\footnotesize Тестуються інструменти та механізми моніторингу вторгнень <SI-04(09)\_\allowbreak{} ODP періодичність>}


{\footnotesize\bfseries
No: 2\\
Name: si\_\allowbreak{}4\_\allowbreak{}9\_\allowbreak{}odp\\
Type: string\\
Default: \textquotedbl{}щорічно\textquotedbl{}
}

{\footnotesize Визначена періодичність тестування механізмів моніторингу вторгнень; інструментів і}




\subsubsection{МОНІТОРИНГ КОМУНІКАЦІЙ (SI-4(10))}


\vspace{10pt}
{\footnotesize\bfseries
No: 1\\
Name: si\_\allowbreak{}4\_\allowbreak{}10\_\allowbreak{}01\\
Type: string\\
Default: nil
}

{\footnotesize Передбачено, щоб <SI-04(10) \_\allowbreak{}ODP[01] зашифрований трафік зв'язку> був видимим для <SI-04(10) \_\allowbreak{}ODP[02] засобів та механізмів моніторингу системи>}




\subsubsection{МОНІТОРИНГ СИСТЕМИ КОМУНІКАЦІЙ (SI-4(11))}


\vspace{10pt}
{\footnotesize\bfseries
No: 1\\
Name: si\_\allowbreak{}4\_\allowbreak{}11\_\allowbreak{}01\\
Type: string\\
Default: nil
}

{\footnotesize Аналізується вихідний комунікаційний трафік на зовнішніх інтерфейсах системи для виявлення аномалій}


{\footnotesize\bfseries
No: 2\\
Name: si\_\allowbreak{}4\_\allowbreak{}11\_\allowbreak{}02\\
Type: string\\
Default: nil
}

{\footnotesize Аналізується вихідний трафік зв'язку в внутрішніх точках для виявлення аномалій}


{\footnotesize\bfseries
No: 3\\
Name: si\_\allowbreak{}4\_\allowbreak{}11\_\allowbreak{}odp\\
Type: string\\
Default: nil
}

{\footnotesize Визначені внутрішні точки в системі, в яких необхідно аналізувати комунікаційний трафік}




\subsubsection{МОНІТОРИНГ СИСТЕМИ АВТОМАТИЗОВАНІ СПОВІЩЕННЯ (SI-4(12))}


\vspace{10pt}
{\footnotesize\bfseries
No: 1\\
Name: si\_\allowbreak{}4\_\allowbreak{}12\_\allowbreak{}01\\
Type: list\\
Default: [\textquotedbl{}admin\textquotedbl{}, \textquotedbl{}security\_\allowbreak{}officer\textquotedbl{}]
}

{\footnotesize Оповіщається <SI-04(12) \_\allowbreak{}ODP[01] персонал або ролі> за допомогою <SI-04(12) \_\allowbreak{}ODP[02] автоматизованих механізмів>, коли <SI-04(12) \_\allowbreak{}ODP[03] дії, що викликають оповіщення> вказують на невідповідну або незвичну викликають оповіщення персоналу, або діяльність, що має вплив на безпеку або приватне життя}




\subsubsection{АНАЛІЗ ТРАФІКУ ТА ШАБЛОНІВ ПОДІЙ (SI-4(13))}


\vspace{10pt}
{\footnotesize\bfseries
No: 1\\
Name: si\_\allowbreak{}4\_\allowbreak{}13\_\allowbreak{}a\_\allowbreak{}01\\
Type: string\\
Default: nil
}

{\footnotesize Аналізується трафік для системи}


{\footnotesize\bfseries
No: 2\\
Name: si\_\allowbreak{}4\_\allowbreak{}13\_\allowbreak{}a\_\allowbreak{}02\\
Type: string\\
Default: nil
}

{\footnotesize Проаналізовано патерни подій для системи}


{\footnotesize\bfseries
No: 3\\
Name: si\_\allowbreak{}4\_\allowbreak{}13\_\allowbreak{}b\_\allowbreak{}01\\
Type: string\\
Default: nil
}

{\footnotesize Розроблені профілі, що представляють загальний трафік}


{\footnotesize\bfseries
No: 4\\
Name: si\_\allowbreak{}4\_\allowbreak{}13\_\allowbreak{}b\_\allowbreak{}02\\
Type: string\\
Default: nil
}

{\footnotesize Розроблені профілі, що представляють патерни подій}


{\footnotesize\bfseries
No: 5\\
Name: si\_\allowbreak{}4\_\allowbreak{}13\_\allowbreak{}c\_\allowbreak{}01\\
Type: string\\
Default: nil
}

{\footnotesize Використовуються профілі трафіку пристроїв системного моніторингу}


{\footnotesize\bfseries
No: 6\\
Name: si\_\allowbreak{}4\_\allowbreak{}13\_\allowbreak{}c\_\allowbreak{}02\\
Type: string\\
Default: nil
}

{\footnotesize Використовуються профілі подій при налаштуванні пристроїв системного моніторингу при налаштуванні}




\subsubsection{МОНІТОРИНГ ВТОРГНЕННЯ (SI-4(14))}


\vspace{10pt}
{\footnotesize\bfseries
No: 1\\
Name: si\_\allowbreak{}4\_\allowbreak{}14\_\allowbreak{}01\\
Type: string\\
Default: nil
}

{\footnotesize Використовується система виявлення бездротових вторгнень для виявлення несанкціонованих бездротових пристроїв}


{\footnotesize\bfseries
No: 2\\
Name: si\_\allowbreak{}4\_\allowbreak{}14\_\allowbreak{}02\\
Type: integer\\
Default: 3
}

{\footnotesize Використовується бездротова система виявлення вторгнень для виявлення спроб атак на систему}


{\footnotesize\bfseries
No: 3\\
Name: si\_\allowbreak{}4\_\allowbreak{}14\_\allowbreak{}03\\
Type: string\\
Default: nil
}

{\footnotesize Використовується бездротова система виявлення вторгнень для виявлення потенційних компрометації або порушень в системі}




\subsubsection{ПЕРЕХІД ВІД БЕЗДРОТОВОГО ЗВ'ЯЗКУ ДО ПРОВІДНИХ МЕРЕЖ (SI-4(15))}


\vspace{10pt}
\textit{Немає параметрів для цього контролю.}
\vspace{10pt}


\subsubsection{МОНІТОРИНГ МОНІТОРИНГУ (SI-4(16))}


\vspace{10pt}
{\footnotesize\bfseries
No: 1\\
Name: si\_\allowbreak{}4\_\allowbreak{}16\_\allowbreak{}01\\
Type: string\\
Default: nil
}

{\footnotesize Співвідноситься інформація з інструментів моніторингу та механізмів, що застосовуються в системі}




\subsubsection{МОНІТОРИНГ ОБІЗНАНІСТЬ (SI-4(17))}


\vspace{10pt}
{\footnotesize\bfseries
No: 1\\
Name: si\_\allowbreak{}4\_\allowbreak{}17\_\allowbreak{}01\\
Type: string\\
Default: nil
}

{\footnotesize Співвідноситься інформація, отримана в результаті моніторингу фізичної, кібернетичної діяльності та діяльності ланцюга поставок, з метою досягнення інтегрованої, загальноорганізаційної ситуаційної обізнаності}




\subsubsection{АНАЛІЗ ТРАФІКУ ТА ПРИХОВАНОЇ ЕКСФІЛЬТРАЦІЇ (SI-4(18))}


\vspace{10pt}
{\footnotesize\bfseries
No: 1\\
Name: si\_\allowbreak{}4\_\allowbreak{}18\_\allowbreak{}01\\
Type: string\\
Default: nil
}

{\footnotesize Аналізується вихідний комунікаційний трафік на зовнішніх по відношенню до системи інтерфейсах для виявлення прихованого витоку інформації}


{\footnotesize\bfseries
No: 2\\
Name: si\_\allowbreak{}4\_\allowbreak{}18\_\allowbreak{}02\\
Type: string\\
Default: nil
}

{\footnotesize Аналізується вихідний трафік зв'язку в внутрішніх точках для виявлення прихованого витоку інформації}


{\footnotesize\bfseries
No: 3\\
Name: si\_\allowbreak{}4\_\allowbreak{}18\_\allowbreak{}odp\\
Type: string\\
Default: nil
}

{\footnotesize Визначені внутрішні точки в системі, в яких необхідно аналізувати комунікаційний трафік}




\subsubsection{ОСОБИ, ЯКІ ПРЕДСТАВЛЯЮТЬ БІЛЬШИЙ РИЗИК (SI-4(19))}


\vspace{10pt}
{\footnotesize\bfseries
No: 1\\
Name: si\_\allowbreak{}4\_\allowbreak{}19\_\allowbreak{}01\\
Type: string\\
Default: nil
}

{\footnotesize Здійснюється <SI-04(19) \_\allowbreak{}ODP[01] додатковий моніторинг> щодо осіб, які були ідентифіковані джерелами <SI-04(19) \_\allowbreak{}ODP[02] як такі, що становлять підвищений рівень ризику>}




\subsubsection{ПРИВІЛЕЙОВАНІ КОРИСТУВАЧІ (SI-4(20))}


\vspace{10pt}
{\footnotesize\bfseries
No: 1\\
Name: si\_\allowbreak{}4\_\allowbreak{}20\_\allowbreak{}01\\
Type: string\\
Default: nil
}

{\footnotesize Реалізовано привілейованих користувачів. привілейованих додатковий моніторинг}


{\footnotesize\bfseries
No: 2\\
Name: si\_\allowbreak{}4\_\allowbreak{}20\_\allowbreak{}odp\\
Type: string\\
Default: nil
}

{\footnotesize Визначено додатковий користувачів; моніторинг}




\subsubsection{ВИПРОБУВАЛЬНІ ТЕРМІНИ (SI-4(21))}


\vspace{10pt}
{\footnotesize\bfseries
No: 1\\
Name: si\_\allowbreak{}4\_\allowbreak{}21\_\allowbreak{}01\\
Type: integer\\
Default: 30
}

{\footnotesize Здійснюється <SI-04(21) \_\allowbreak{}ODP[01] додатковий моніторинг> осіб під час <SI-04(21) \_\allowbreak{}ODP[02] випробувального терміну>}




\subsubsection{НЕСАНКЦІОНОВАНІ ПОСЛУГИ МЕРЕЖІ (SI-4(22))}


\vspace{10pt}
{\footnotesize\bfseries
No: 1\\
Name: si\_\allowbreak{}4\_\allowbreak{}22\_\allowbreak{}a\\
Type: string\\
Default: nil
}

{\footnotesize Виявлено послуги мережі, які не було авторизовано або схвалено відповідно до процесів авторизації або схвалення}




\subsubsection{ПРИСТРОЇ НА ОСНОВІ ХОСТА (SI-4(23))}


\vspace{10pt}
{\footnotesize\bfseries
No: 1\\
Name: si\_\allowbreak{}4\_\allowbreak{}23\_\allowbreak{}01\\
Type: string\\
Default: \textquotedbl{}автоматизований засіб моніторингу\textquotedbl{}
}

{\footnotesize Реалізовано <SI-04(23) \_\allowbreak{}ODP[01] механізми моніторингу на основі хостів> на <SI-04(23) \_\allowbreak{}ODP[02] компоненти системи>}




\subsubsection{ІНДИКАТОРИ КОМПРОМЕТАЦІЇ (SI-4(24))}


\vspace{10pt}
{\footnotesize\bfseries
No: 1\\
Name: si\_\allowbreak{}4\_\allowbreak{}24\_\allowbreak{}01\\
Type: string\\
Default: nil
}

{\footnotesize Виявлено індикатори компрометації, 04(24)\_\allowbreak{}ODP[01] джерелами>}


{\footnotesize\bfseries
No: 2\\
Name: si\_\allowbreak{}4\_\allowbreak{}24\_\allowbreak{}02\\
Type: string\\
Default: nil
}

{\footnotesize Збираються індикатори компрометації, що надаються джерелами}


{\footnotesize\bfseries
No: 3\\
Name: si\_\allowbreak{}4\_\allowbreak{}24\_\allowbreak{}03\\
Type: list\\
Default: [\textquotedbl{}admin\textquotedbl{}, \textquotedbl{}security\_\allowbreak{}officer\textquotedbl{}]
}

{\footnotesize Індикатори компрометації, надані <SI-04(24) джерелами>, поширюються на <SI-04(24) персонал або ролі>. надані <SI- \_\allowbreak{}ODP[01] \_\allowbreak{}ODP[02]}




\subsubsection{АНАЛІЗ МЕРЕЖЕВОГО ТРАФІКУ (SI-4(25))}


\vspace{10pt}
{\footnotesize\bfseries
No: 1\\
Name: si\_\allowbreak{}4\_\allowbreak{}25\_\allowbreak{}01\\
Type: string\\
Default: nil
}

{\footnotesize Забезпечується видимість мережевого трафіку на зовнішніх системних інтерфейсах для оптимізації ефективності пристроїв моніторингу}


{\footnotesize\bfseries
No: 2\\
Name: si\_\allowbreak{}4\_\allowbreak{}25\_\allowbreak{}02\\
Type: string\\
Default: nil
}

{\footnotesize Забезпечено видимість мережевого трафіку на ключових внутрішніх інтерфейсах системи для оптимізації ефективності пристроїв моніторингу}




\subsection{ПОПЕРЕДЖЕННЯ, РЕКОМЕНДАЦІЇ ТА ДИРЕКТИВИ З БЕЗПЕКИ (SI-5)}


\vspace{10pt}
{\footnotesize\bfseries
No: 1\\
Name: si\_\allowbreak{}5\_\allowbreak{}odp\_\allowbreak{}01\\
Type: string\\
Default: nil
}

{\footnotesize Визначені зовнішні організації, від яких необхідно постійно отримувати оповіщення, поради та директиви щодо безпеки системи}


{\footnotesize\bfseries
No: 2\\
Name: si\_\allowbreak{}5\_\allowbreak{}odp\_\allowbreak{}02\\
Type: list\\
Default: [\textquotedbl{}admin\textquotedbl{}, \textquotedbl{}security\_\allowbreak{}officer\textquotedbl{}]
}

{\footnotesize Вибрано одне або декілька з наступних ЗНАЧЕНЬ ПАРАМЕТРІВ: \{персонал або ролі; елементи; зовнішні організації\}}


{\footnotesize\bfseries
No: 3\\
Name: si\_\allowbreak{}5\_\allowbreak{}odp\_\allowbreak{}03\\
Type: list\\
Default: [\textquotedbl{}admin\textquotedbl{}, \textquotedbl{}security\_\allowbreak{}officer\textquotedbl{}]
}

{\footnotesize Визначено персонал або ролі, до яких мають бути доведені попередження, поради та директиви з безпеки (якщо визначено)}


{\footnotesize\bfseries
No: 4\\
Name: si\_\allowbreak{}5\_\allowbreak{}odp\_\allowbreak{}04\\
Type: string\\
Default: nil
}

{\footnotesize Визначені елементи в організації, до яких мають надсилатися оповіщення, поради та директиви з безпеки (якщо вони були обрані)}




\subsubsection{ПОПЕРЕДЖЕННЯ, РЕКОМЕНДАЦІЇ ТА ДИРЕКТИВИ З БЕЗПЕКИ АВТОМАТИЧНІ ПОПЕРЕДЖЕННЯ ТА РЕКОМЕНДАЦІЇ (SI-5(1))}


\vspace{10pt}
{\footnotesize\bfseries
No: 1\\
Name: si\_\allowbreak{}5\_\allowbreak{}1\_\allowbreak{}01\\
Type: string\\
Default: \textquotedbl{}автоматизований засіб моніторингу\textquotedbl{}
}

{\footnotesize Використовуються автоматизовані механізми для трансляції попередження та рекомендації інформації з питань безпеки по всій організації}


{\footnotesize\bfseries
No: 2\\
Name: si\_\allowbreak{}5\_\allowbreak{}1\_\allowbreak{}odp\\
Type: string\\
Default: \textquotedbl{}автоматизований засіб моніторингу\textquotedbl{}
}

{\footnotesize Визначені автоматизовані механізми, які використовуються для трансляції попередження та рекомендації інформації про безпеку в організації}




\subsection{ПЕРЕВІРКА ФУНКЦІЙ БЕЗПЕКИ ТА ПРИВАТНОСТІ (SI-6)}


\vspace{10pt}
{\footnotesize\bfseries
No: 1\\
Name: si\_\allowbreak{}6\_\allowbreak{}odp\_\allowbreak{}01\\
Type: string\\
Default: nil
}

{\footnotesize Визначені функції безпеки, які необхідно перевірити на коректність роботи}


{\footnotesize\bfseries
No: 2\\
Name: si\_\allowbreak{}6\_\allowbreak{}odp\_\allowbreak{}02\\
Type: string\\
Default: nil
}

{\footnotesize Визначені функції приватності, які потрібно перевіряти на коректність роботи}


{\footnotesize\bfseries
No: 3\\
Name: si\_\allowbreak{}6\_\allowbreak{}odp\_\allowbreak{}03\\
Type: string\\
Default: \textquotedbl{}щорічно\textquotedbl{}
}

{\footnotesize Вибрано одне або декілька з наступних ЗНАЧЕНЬ ПАРАМЕТРІВ: \{перехідні стани системи; за командою користувача з відповідним привілеєм; частота\}}


{\footnotesize\bfseries
No: 4\\
Name: si\_\allowbreak{}6\_\allowbreak{}odp\_\allowbreak{}04\\
Type: string\\
Default: nil
}

{\footnotesize Визначені перехідні стани системи, що вимагають перевірки функцій безпеки та конфіденційності; (якщо вибрано)}


{\footnotesize\bfseries
No: 5\\
Name: si\_\allowbreak{}6\_\allowbreak{}odp\_\allowbreak{}05\\
Type: string\\
Default: \textquotedbl{}щорічно\textquotedbl{}
}

{\footnotesize Визначена періодичність перевірки правильності роботи функцій безпеки та приватності; (якщо вибрано)}


{\footnotesize\bfseries
No: 6\\
Name: si\_\allowbreak{}6\_\allowbreak{}odp\_\allowbreak{}06\\
Type: list\\
Default: [\textquotedbl{}admin\textquotedbl{}, \textquotedbl{}security\_\allowbreak{}officer\textquotedbl{}]
}

{\footnotesize Визначено персонал або ролі, які мають бути сповіщені про невдалу перевірку безпеки та приватності}


{\footnotesize\bfseries
No: 7\\
Name: si\_\allowbreak{}6\_\allowbreak{}odp\_\allowbreak{}07\\
Type: list\\
Default: [\textquotedbl{}login\textquotedbl{}, \textquotedbl{}logout\textquotedbl{}, \textquotedbl{}failed\_\allowbreak{}attempt\textquotedbl{}]
}

{\footnotesize Вибрано одне або декілька з наступних ЗНАЧЕНЬ ПАРАМЕТРІВ: \{вимкнути систему; перезапустити систему; альтернативна дія (дії)\}}




\subsubsection{СПОВІЩЕННЯ ПРО НЕУСПІШНЕ ПРОХОДЖЕННЯ ТЕСТІВ З БЕЗПЕКИ (SI-6(1)) [Вилучено]}
[Вилучено: включено до SI-6]

\vspace{10pt}
\textit{Немає параметрів для цього контролю.}
\vspace{10pt}


\subsubsection{ПЕРЕВІРКА БЕЗПЕКИ ТА ФУНКЦІЙ ПРИВАТНОСТІ АВТОМАТИЗОВАНА ПІДТРИМКА РОЗПОДІЛЕНОГО ТЕСТУВАННЯ (SI-6(2))}


\vspace{10pt}
{\footnotesize\bfseries
No: 1\\
Name: si\_\allowbreak{}6\_\allowbreak{}2\_\allowbreak{}01\\
Type: string\\
Default: \textquotedbl{}автоматизований засіб моніторингу\textquotedbl{}
}

{\footnotesize Впроваджені автоматизовані механізми для розподіленого тестуванням функцій безпеки; підтримки}


{\footnotesize\bfseries
No: 2\\
Name: si\_\allowbreak{}6\_\allowbreak{}2\_\allowbreak{}02\\
Type: string\\
Default: \textquotedbl{}автоматизований засіб моніторингу\textquotedbl{}
}

{\footnotesize Проваджені автоматизовані механізми для розподіленого тестуванням функцій приватності. підтримки}




\subsubsection{ПЕРЕВІРКА БЕЗПЕКИ ТА ФУНКЦІЙ ПРИВАТНОСТІ ПОВІДОМЛЕННЯ ПРО РЕЗУЛЬТАТИ ПЕРЕВІРКИ (SI-6(3))}


\vspace{10pt}
{\footnotesize\bfseries
No: 1\\
Name: si\_\allowbreak{}6\_\allowbreak{}3\_\allowbreak{}01\\
Type: list\\
Default: [\textquotedbl{}admin\textquotedbl{}, \textquotedbl{}security\_\allowbreak{}officer\textquotedbl{}]
}

{\footnotesize Повідомляються результати перевірки функцій безпеки персоналу або ролям}


{\footnotesize\bfseries
No: 2\\
Name: si\_\allowbreak{}6\_\allowbreak{}3\_\allowbreak{}odp\\
Type: list\\
Default: [\textquotedbl{}admin\textquotedbl{}, \textquotedbl{}security\_\allowbreak{}officer\textquotedbl{}]
}

{\footnotesize Визначено персонал або ролі, призначені для отримання результатів перевірки функцій безпеки та приватності}




\subsection{ЦІЛІСНІСТЬ ПРОГРАМНОГО ЗАБЕЗПЕЧЕННЯ, ПРОГРАМНОГО ЗАБЕЗПЕЧЕННЯ ТА ІНФОРМАЦІЇ (SI-7)}
a. Впровадити механізми захисту від спаму в точках входу та виходу системи, щоб виявляти та протидіяти небажаним повідомленням.\\ 
b. Оновлювати механізми захисту від спаму, коли доступні нові механізми відповідно до організаційної політики та процедур управління конфігурацією.

\vspace{10pt}
{\footnotesize\bfseries
No: 1\\
Name: si\_\allowbreak{}7\_\allowbreak{}odp\_\allowbreak{}01\\
Type: string\\
Default: nil
}

{\footnotesize Визначено програмне забезпечення, яке потребує застосування засобів перевірки цілісності для виявлення несанкціонованих змін}


{\footnotesize\bfseries
No: 2\\
Name: si\_\allowbreak{}7\_\allowbreak{}odp\_\allowbreak{}02\\
Type: string\\
Default: nil
}

{\footnotesize Визначено прошивку, яка потребує застосування інструментів перевірки цілісності для виявлення несанкціонованих змін}


{\footnotesize\bfseries
No: 3\\
Name: si\_\allowbreak{}7\_\allowbreak{}odp\_\allowbreak{}03\\
Type: string\\
Default: nil
}

{\footnotesize Визначена інформація, яка потребує застосування засобів перевірки цілісності для виявлення несанкціонованих змін}


{\footnotesize\bfseries
No: 4\\
Name: si\_\allowbreak{}7\_\allowbreak{}odp\_\allowbreak{}04\\
Type: list\\
Default: [\textquotedbl{}login\textquotedbl{}, \textquotedbl{}logout\textquotedbl{}, \textquotedbl{}failed\_\allowbreak{}attempt\textquotedbl{}]
}

{\footnotesize Визначені дії, яких слід вжити при виявленні несанкціонованих змін у програмному забезпеченні}


{\footnotesize\bfseries
No: 5\\
Name: si\_\allowbreak{}7\_\allowbreak{}odp\_\allowbreak{}05\\
Type: list\\
Default: [\textquotedbl{}login\textquotedbl{}, \textquotedbl{}logout\textquotedbl{}, \textquotedbl{}failed\_\allowbreak{}attempt\textquotedbl{}]
}

{\footnotesize Визначені дії, яких слід вжити несанкціонованих змін у прошивці; при виявленні SI-07\_\allowbreak{}ODP[06] визначені дії, яких слід вжити несанкціонованих змін до інформації; при виявленні SI-07a.[01] використовуються засоби перевірки цілісності для виявлення несанкціонованих змін у програмному забезпеченні; SI-07a.[02] використовуються засоби перевірки цілісності для виявлення несанкціонованих змін у мікропрограмі; SI-07a.[03] використовуються засоби перевірки цілісності для виявлення несанкціонованих змін до інформації; SI-07b.[01] виконуються дії при виявленні несанкціонованих змін у програмному забезпеченні; SI-07b.[02] виконуються дії несанкціонованих змін у прошивці; при виявленні SI-07b.[03] виконуються дії несанкціонованих змін в інформації. при виявленні}




\subsubsection{ПЕРЕВІРКА ЦІЛІСНОСТІ (SI-7(1))}


\vspace{10pt}
\textit{Немає параметрів для цього контролю.}
\vspace{10pt}


\subsubsection{АВТОМАТИЧНІ СПОВІЩЕННЯ ПРО ПОРУШЕННЯ ЦІЛІСНОСТІ (SI-7(2))}


\vspace{10pt}
{\footnotesize\bfseries
No: 1\\
Name: si\_\allowbreak{}7\_\allowbreak{}2\_\allowbreak{}01\\
Type: list\\
Default: [\textquotedbl{}admin\textquotedbl{}, \textquotedbl{}security\_\allowbreak{}officer\textquotedbl{}]
}

{\footnotesize Застосовуються автоматизовані інструменти, які надають повідомлення персоналу або ролям при виявленні розбіжностей під час перевірки цілісності}


{\footnotesize\bfseries
No: 2\\
Name: si\_\allowbreak{}7\_\allowbreak{}2\_\allowbreak{}odp\\
Type: list\\
Default: [\textquotedbl{}admin\textquotedbl{}, \textquotedbl{}security\_\allowbreak{}officer\textquotedbl{}]
}

{\footnotesize Визначено персонал або ролі, яким необхідно повідомляти про виявлення розбіжностей під час перевірки цілісності}




\subsubsection{ІНСТРУМЕНТИ ЦІЛІСНОСТІ З ЦЕНТРАЛІЗОВАНИМ УПРАВЛІННЯМ (SI-7(3))}


\vspace{10pt}
{\footnotesize\bfseries
No: 1\\
Name: si\_\allowbreak{}7\_\allowbreak{}3\_\allowbreak{}01\\
Type: string\\
Default: nil
}

{\footnotesize Застосовуються інструменти цілісності з централізованим управлінням}




\subsubsection{ЦІЛІСНІСТЬ ПРОГРАМНОГО ЗАБЕЗПЕЧЕННЯ, ВБУДОВАНОГО ПРОГРАМНОГО ЗАБЕЗПЕЧЕННЯ ТА ІНФОРМАЦІЇ -- ПАКУВАННЯ З ІНДИКАЦІЄЮ ОЗНАК ЇЇ НЕСАНКЦІОНОВАНОГО РОЗКРИТТЯ (SI-7(4)) [Вилучено]}
[Вилучено: включено до SA-12]

\vspace{10pt}
\textit{Немає параметрів для цього контролю.}
\vspace{10pt}


\subsubsection{АВТОМАТИЧНІ ВІДПОВІДІ ПРО ПОРУШЕННЯ ЦІЛІСНОСТІ (SI-7(5))}


\vspace{10pt}
\textit{Немає параметрів для цього контролю.}
\vspace{10pt}


\subsubsection{ЦІЛІСНІСТЬ ПРОГРАМНОГО ЗАБЕЗПЕЧЕННЯ, ВБУДОВАНОГО ПРОГРАМНОГО ЗАБЕЗПЕЧЕННЯ ТА ІНФОРМАЦІЇ КРИПТОГРАФІЧНИЙ ЗАХИСТ (SI-7(6))}


\vspace{10pt}
{\footnotesize\bfseries
No: 1\\
Name: si\_\allowbreak{}7\_\allowbreak{}6\_\allowbreak{}01\\
Type: string\\
Default: \textquotedbl{}AES-256-GCM\textquotedbl{}
}

{\footnotesize Впроваджені криптографічні механізми для виявлення несанкціонованих змін у програмному забезпеченні}


{\footnotesize\bfseries
No: 2\\
Name: si\_\allowbreak{}7\_\allowbreak{}6\_\allowbreak{}02\\
Type: string\\
Default: \textquotedbl{}AES-256-GCM\textquotedbl{}
}

{\footnotesize Реалізовані криптографічні механізми несанкціонованих змін у прошивці; для виявлення}


{\footnotesize\bfseries
No: 3\\
Name: si\_\allowbreak{}7\_\allowbreak{}6\_\allowbreak{}03\\
Type: string\\
Default: \textquotedbl{}AES-256-GCM\textquotedbl{}
}

{\footnotesize Впроваджені криптографічні механізми несанкціонованих змін інформації. для виявлення}




\subsubsection{ІНТЕГРАЦІЯ ВИЯВЛЕННЯ І РЕАГУВАННЯ (SI-7(7))}


\vspace{10pt}
{\footnotesize\bfseries
No: 1\\
Name: si\_\allowbreak{}7\_\allowbreak{}7\_\allowbreak{}01\\
Type: string\\
Default: nil
}

{\footnotesize Виявлення змін включено до можливості організації реагування на інциденти}


{\footnotesize\bfseries
No: 2\\
Name: si\_\allowbreak{}7\_\allowbreak{}7\_\allowbreak{}odp\\
Type: string\\
Default: nil
}

{\footnotesize Визначені зміни в системі, що мають відношення до безпеки}




\subsubsection{АУДИТ ВАЖЛИВИХ ПОДІЙ (SI-7(8))}


\vspace{10pt}
{\footnotesize\bfseries
No: 1\\
Name: si\_\allowbreak{}7\_\allowbreak{}8\_\allowbreak{}01\\
Type: string\\
Default: nil
}

{\footnotesize Передбачена можливість аудиту потенційного порушення цілісності}




\subsubsection{ПЕРЕВІРКА ПРОЦЕСУ ЗАВАНТАЖЕННЯ (SI-7(9))}


\vspace{10pt}
{\footnotesize\bfseries
No: 1\\
Name: si\_\allowbreak{}7\_\allowbreak{}9\_\allowbreak{}01\\
Type: string\\
Default: nil
}

{\footnotesize Перевіряється цілісність процесу завантаження 07(09)\_\allowbreak{}ODP системних компонентів>. <SI-}


{\footnotesize\bfseries
No: 2\\
Name: si\_\allowbreak{}7\_\allowbreak{}9\_\allowbreak{}odp\\
Type: string\\
Default: nil
}

{\footnotesize Визначено компоненти системи, які потребують перевірки цілісності процесу завантаження}




\subsubsection{ЗАХИСТ ЗАВАНТАЖУВАЛЬНОГО ВБУДОВАНОГО ПРОГРАМНОГО ЗАБЕЗПЕЧЕННЯ (SI-7(10))}


\vspace{10pt}
{\footnotesize\bfseries
No: 1\\
Name: si\_\allowbreak{}7\_\allowbreak{}10\_\allowbreak{}01\\
Type: string\\
Default: \textquotedbl{}автоматизований засіб моніторингу\textquotedbl{}
}

{\footnotesize Реалізовано <SI-07(10) \_\allowbreak{}ODP[01] механізми> для захисту цілісності завантажувальної прошивки у <SI-07(10) \_\allowbreak{}ODP[02] системних компонентах>. компоненти системи, захисту цілісності які потребують завантажувальної}




\subsubsection{ОБМЕЖЕНЕ СЕРЕДОВИЩЕ З ОБМЕЖЕНИМИ ПРИВІЛЕЯМИ (SI-7(11)) [Вилучено]}
[Вилучено: включено до CM-7(6)]

\vspace{10pt}
\textit{Немає параметрів для цього контролю.}
\vspace{10pt}


\subsubsection{ПЕРЕВІРКА ЦІЛІСНОСТІ (SI-7(12))}


\vspace{10pt}
{\footnotesize\bfseries
No: 1\\
Name: si\_\allowbreak{}7\_\allowbreak{}12\_\allowbreak{}01\\
Type: string\\
Default: nil
}

{\footnotesize Перевіряється забезпечення, виконанням. цілісність програмне встановлене користувачем перед}


{\footnotesize\bfseries
No: 2\\
Name: si\_\allowbreak{}7\_\allowbreak{}12\_\allowbreak{}odp\\
Type: string\\
Default: nil
}

{\footnotesize Визначено програмне забезпечення, встановлене користувачем, яке потребує перевірки цілісності перед виконанням}




\subsubsection{ВИКОНАННЯ КОДУ В ЗАХИЩЕНИХ СЕРЕДОВИЩАХ (SI-7(13)) [Вилучено]}
[Вилучено: включено до CM-7(7)]

\vspace{10pt}
\textit{Немає параметрів для цього контролю.}
\vspace{10pt}


\subsubsection{ДВІЙКОВИЙ АБО МАШИННО-ВИКОНУВАНИЙ КОД (SI-7(14)) [Вилучено]}
[Вилучено: включено до CM-7(8)]

\vspace{10pt}
\textit{Немає параметрів для цього контролю.}
\vspace{10pt}


\subsubsection{ЦІЛІСНІСТЬ ПРОГРАМНОГО ЗАБЕЗПЕЧЕННЯ, ВБУДОВАНОГО ПРОГРАМНОГО ЗАБЕЗПЕЧЕННЯ ТА ІНФОРМАЦІЇ АВТЕНТИФІКАЦІЯ КОДУ (SI-7(15))}


\vspace{10pt}
{\footnotesize\bfseries
No: 1\\
Name: si\_\allowbreak{}7\_\allowbreak{}15\_\allowbreak{}01\\
Type: string\\
Default: \textquotedbl{}AES-256-GCM\textquotedbl{}
}

{\footnotesize Реалізовано криптографічні механізми для автентифікації програмного забезпечення або компонентів мікропрограми перед інсталяцією}


{\footnotesize\bfseries
No: 2\\
Name: si\_\allowbreak{}7\_\allowbreak{}15\_\allowbreak{}odp\\
Type: string\\
Default: \textquotedbl{}AES-256-GCM\textquotedbl{}
}

{\footnotesize Визначено компоненти програмного забезпечення або мікропрограми, які мають бути автентифіковані за допомогою криптографічних механізмів перед встановленням}




\subsubsection{ТЕРМІН ВИКОНАННЯ ПРОЦЕСУ БЕЗ НАГЛЯДУ (SI-7(16))}


\vspace{10pt}
{\footnotesize\bfseries
No: 1\\
Name: si\_\allowbreak{}7\_\allowbreak{}16\_\allowbreak{}01\\
Type: integer\\
Default: 30
}

{\footnotesize Заборонено процесам виконуватися без нагляду довше, ніж часовий період}


{\footnotesize\bfseries
No: 2\\
Name: si\_\allowbreak{}7\_\allowbreak{}16\_\allowbreak{}odp\\
Type: integer\\
Default: 30
}

{\footnotesize Визначено максимальний період часу, протягом якого процеси можуть виконуватися без нагляду}




\subsubsection{ЦІЛІСНІСТЬ ПРОГРАМНОГО ЗАБЕЗПЕЧЕННЯ, ВБУДОВАНОГО ПРОГРАМНОГО ЗАБЕЗПЕЧЕННЯ ТА ІНФОРМАЦІЇ -- САМОЗАХИСТ ПРОГРАМ ВІД САМОВІЛЬНОГО ВИКОНАННЯ (SI-7(17))}


\vspace{10pt}
{\footnotesize\bfseries
No: 1\\
Name: si\_\allowbreak{}7\_\allowbreak{}17\_\allowbreak{}odp\\
Type: integer\\
Default: 30
}

{\footnotesize Визначено елементи керування, які потрібно реалізувати для самозахисту програми під час виконання}




\subsection{ЗАХИСТ ВІД СПАМУ (SI-8)}


\vspace{10pt}
\textit{Немає параметрів для цього контролю.}
\vspace{10pt}


\subsubsection{ЦЕНТРАЛІЗОВАНЕ УПРАВЛІННЯ (SI-8(1)) [Вилучено]}
[Вилучено: включено до PL-9]

\vspace{10pt}
\textit{Немає параметрів для цього контролю.}
\vspace{10pt}


\subsubsection{АВТОМАТИЧНІ ОНОВЛЕННЯ (SI-8(2))}


\vspace{10pt}
{\footnotesize\bfseries
No: 1\\
Name: si\_\allowbreak{}8\_\allowbreak{}2\_\allowbreak{}01\\
Type: string\\
Default: \textquotedbl{}щорічно\textquotedbl{}
}

{\footnotesize Автоматично оновлюються механізми захисту від спаму частота}


{\footnotesize\bfseries
No: 2\\
Name: si\_\allowbreak{}8\_\allowbreak{}2\_\allowbreak{}odp\\
Type: string\\
Default: \textquotedbl{}щорічно\textquotedbl{}
}

{\footnotesize Визначено періодичність автоматичного механізмів захисту від спаму; оновлення}




\subsubsection{БЕЗПЕРЕРВНЕ НАВЧАННЯ (SI-8(3))}


\vspace{10pt}
{\footnotesize\bfseries
No: 1\\
Name: si\_\allowbreak{}8\_\allowbreak{}3\_\allowbreak{}01\\
Type: string\\
Default: \textquotedbl{}автоматизований засіб моніторингу\textquotedbl{}
}

{\footnotesize Впроваджені механізми захисту від спаму з можливістю навчання для більш ефективного визначення законого комунікаційного трафіку}




\subsection{ОБМЕЖЕННЯ НА ВВЕДЕННЯ ІНФОРМАЦІЇ (SI-9)}


\vspace{10pt}
\textit{Немає параметрів для цього контролю.}
\vspace{10pt}


\subsection{ПЕРЕВІРКА ВВОДУ ІНФОРМАЦІЇ (SI-10)}
Перевіряти дійсність [Призначення: визначена організацією введена інформація].

\vspace{10pt}
{\footnotesize\bfseries
No: 1\\
Name: si\_\allowbreak{}10\_\allowbreak{}01\\
Type: string\\
Default: nil
}

{\footnotesize Перевіряється дійсність синтаксису вхідної інформації}


{\footnotesize\bfseries
No: 2\\
Name: si\_\allowbreak{}10\_\allowbreak{}odp\\
Type: string\\
Default: nil
}

{\footnotesize Визначено вхідні дані до системи, які перевірки достовірності; потребують}




\subsubsection{ПЕРЕВІРКА ВВОДУ ПЕРЕВИЗНАЧЕННЯ (SI-10(1))}


\vspace{10pt}
{\footnotesize\bfseries
No: 1\\
Name: si\_\allowbreak{}10\_\allowbreak{}1\_\allowbreak{}a\\
Type: string\\
Default: nil
}

{\footnotesize Передбачена можливість ручного перевизначення валідації інформаційних входів}


{\footnotesize\bfseries
No: 2\\
Name: si\_\allowbreak{}10\_\allowbreak{}1\_\allowbreak{}b\\
Type: list\\
Default: [\textquotedbl{}admin\textquotedbl{}, \textquotedbl{}security\_\allowbreak{}officer\textquotedbl{}]
}

{\footnotesize Використання можливості ручного перевизначення обмежено лише уповноваженими особами}


{\footnotesize\bfseries
No: 3\\
Name: si\_\allowbreak{}10\_\allowbreak{}1\_\allowbreak{}c\\
Type: string\\
Default: nil
}

{\footnotesize Проводиться аудит перевизначення. використання можливості для ручного}


{\footnotesize\bfseries
No: 4\\
Name: si\_\allowbreak{}10\_\allowbreak{}1\_\allowbreak{}odp\\
Type: string\\
Default: nil
}

{\footnotesize Визначено авторизованих осіб, які можуть користуватися можливістю ручного перевизначення}




\subsubsection{ПЕРЕВІРКА ВВОДУ ІНФОРМАЦІЇ ПОМИЛОК (SI-10(2))}


\vspace{10pt}
{\footnotesize\bfseries
No: 1\\
Name: si\_\allowbreak{}10\_\allowbreak{}2\_\allowbreak{}01\\
Type: integer\\
Default: 30
}

{\footnotesize Переглядаються помилки валідації вхідних даних протягом часового періоду}


{\footnotesize\bfseries
No: 2\\
Name: si\_\allowbreak{}10\_\allowbreak{}2\_\allowbreak{}02\\
Type: integer\\
Default: 30
}

{\footnotesize Помилки валідації вводу вирішуються 10(02)\_\allowbreak{}ODP[02] часового проміжку>. протягом <SI-}




\subsubsection{ПЕРЕДБАЧУВАНА ПОВЕДІНКА (SI-10(3))}


\vspace{10pt}
{\footnotesize\bfseries
No: 1\\
Name: si\_\allowbreak{}10\_\allowbreak{}3\_\allowbreak{}02\\
Type: string\\
Default: nil
}

{\footnotesize Система поводиться задокументованим чином при отриманні недійсних вхідних даних}




\subsubsection{ЧАСОВІ ВЗАЄМОДІЇ (SI-10(4))}


\vspace{10pt}
{\footnotesize\bfseries
No: 1\\
Name: si\_\allowbreak{}10\_\allowbreak{}4\_\allowbreak{}01\\
Type: integer\\
Default: 30
}

{\footnotesize Враховується часова взаємодія між компонентами системи при визначенні відповідної реакції на невірні вхідні дані}




\subsubsection{ОБМЕЖЕННЯ ВХІДНИХ ДАНИХ ДОВІРЕНИМИ ДЖЕРЕЛАМИ І ЗАТВЕРДЖЕНИМИ ФОРМАТАМИ (SI-10(5))}


\vspace{10pt}
\textit{Немає параметрів для цього контролю.}
\vspace{10pt}


\subsubsection{ПРОФІЛАКТИКА ВВОДУ ДАНИХ (SI-10(6))}


\vspace{10pt}
{\footnotesize\bfseries
No: 1\\
Name: si\_\allowbreak{}10\_\allowbreak{}6\_\allowbreak{}01\\
Type: integer\\
Default: 30
}

{\footnotesize Визначено елементи керування, які потрібно реалізувати для самозахисту програми під час виконання}




\subsection{ОБРОБКА ПОМИЛОК (SI-11)}
a. Створити повідомлення про помилки, які надають інформацію, необхідну для реалізації виправних дій, без виявлення інформації, що може бути використана.\\ 
b. Показувати повідомлення про помилки лише [Призначення: визначений організацією персонал або посадові особи].

\vspace{10pt}
{\footnotesize\bfseries
No: 1\\
Name: si\_\allowbreak{}11\_\allowbreak{}odp\\
Type: list\\
Default: [\textquotedbl{}admin\textquotedbl{}, \textquotedbl{}security\_\allowbreak{}officer\textquotedbl{}]
}

{\footnotesize Визначено персонал або ролі, яким слід повідомляти про повідомлення про помилки; SI-11a. генеруються повідомлення про помилки, які надають інформацію, необхідну для коригувальних дій, без розкриття інформації, яка може бути використана; SI-11b. показувати повідомлення про помилки лише для персонал або ролі}




\subsection{УПРАВЛІННЯ ТА ЗБЕРЕЖЕННЯ ІНФОРМАЦІЇ (SI-12)}
Управляти та зберігати інформацію всередині системи та виводити інформацію із системи відповідно до чинного законодавства, виконавчих наказів, директив, правил, політик, стандартів, керівних принципів та експлуатаційних вимог.

\vspace{10pt}
{\footnotesize\bfseries
No: 1\\
Name: si\_\allowbreak{}12\_\allowbreak{}01\\
Type: list\\
Default: [\textquotedbl{}default\_\allowbreak{}deny\_\allowbreak{}rule\textquotedbl{}, \textquotedbl{}abac\_\allowbreak{}rule\_\allowbreak{}1\textquotedbl{}]
}

{\footnotesize Здійснюється управління інформацією в системі відповідно до чинних законів, наказів, директив, положень, політик, стандартів, інструкцій та експлуатаційних вимог}


{\footnotesize\bfseries
No: 2\\
Name: si\_\allowbreak{}12\_\allowbreak{}02\\
Type: list\\
Default: [\textquotedbl{}default\_\allowbreak{}deny\_\allowbreak{}rule\textquotedbl{}, \textquotedbl{}abac\_\allowbreak{}rule\_\allowbreak{}1\textquotedbl{}]
}

{\footnotesize Зберігається інформація в системі відповідно до чинних законів, указів Президента, директив, положень, політик, стандартів, інструкцій та експлуатаційних вимог}


{\footnotesize\bfseries
No: 3\\
Name: si\_\allowbreak{}12\_\allowbreak{}03\\
Type: list\\
Default: [\textquotedbl{}default\_\allowbreak{}deny\_\allowbreak{}rule\textquotedbl{}, \textquotedbl{}abac\_\allowbreak{}rule\_\allowbreak{}1\textquotedbl{}]
}

{\footnotesize Управління інформацією, що виводиться з системи, здійснюється відповідно до чинних законів, указів Президента, директив, положень, політик, стандартів, інструкцій та операційних вимог}


{\footnotesize\bfseries
No: 4\\
Name: si\_\allowbreak{}12\_\allowbreak{}04\\
Type: list\\
Default: [\textquotedbl{}default\_\allowbreak{}deny\_\allowbreak{}rule\textquotedbl{}, \textquotedbl{}abac\_\allowbreak{}rule\_\allowbreak{}1\textquotedbl{}]
}

{\footnotesize Зберігається інформація, що виводиться з системи, відповідно до чинних законів, наказів, директив, положень, політик, стандартів, інструкцій та експлуатаційних вимог}




\subsubsection{УПРАВЛІННЯ ТА ЗБЕРЕЖЕННЯ ІНФОРМАЦІЇ ЕЛЕМЕНТІВ ПЕРСОНАЛЬНИХ ДАНИХ (SI-12(1))}


\vspace{10pt}
{\footnotesize\bfseries
No: 1\\
Name: si\_\allowbreak{}12\_\allowbreak{}1\_\allowbreak{}01\\
Type: list\\
Default: [\textquotedbl{}admin\textquotedbl{}, \textquotedbl{}security\_\allowbreak{}officer\textquotedbl{}]
}

{\footnotesize Обмежується обробка персональних даних у життєвому циклі інформації в життєвому циклі інформації, елементами інформації, що ідентифікує особу}


{\footnotesize\bfseries
No: 2\\
Name: si\_\allowbreak{}12\_\allowbreak{}1\_\allowbreak{}odp\\
Type: list\\
Default: [\textquotedbl{}admin\textquotedbl{}, \textquotedbl{}security\_\allowbreak{}officer\textquotedbl{}]
}

{\footnotesize Визначені елементи персональних даних у життєвому циклі інформації}




\subsubsection{МІНІМІЗАЦІЯ ВИКОРИСТАННЯ ПЕРСОНАЛЬНИХ ДАНИХ ПІД ЧАС ТЕСТУВАННЯ, НАВЧАННЯ ТА ДОСЛІДЖЕННІ (SI-12(2))}


\vspace{10pt}
{\footnotesize\bfseries
No: 1\\
Name: si\_\allowbreak{}12\_\allowbreak{}2\_\allowbreak{}01\\
Type: list\\
Default: [\textquotedbl{}admin\textquotedbl{}, \textquotedbl{}security\_\allowbreak{}officer\textquotedbl{}]
}

{\footnotesize Використовуються методи для мінімізації використання персональних даних для досліджень; SI-12(02)[02] використовуються методи для мінімізації використання персональних даних для тестування}


{\footnotesize\bfseries
No: 2\\
Name: si\_\allowbreak{}12\_\allowbreak{}2\_\allowbreak{}03\\
Type: list\\
Default: [\textquotedbl{}admin\textquotedbl{}, \textquotedbl{}security\_\allowbreak{}officer\textquotedbl{}]
}

{\footnotesize Застосовуються методи для мінімізації використання персональних даних для навчання}




\subsubsection{УПРАВЛІННЯ ІНФОРМАЦІЇ (SI-12(3))}


\vspace{10pt}
{\footnotesize\bfseries
No: 1\\
Name: si\_\allowbreak{}12\_\allowbreak{}3\_\allowbreak{}01\\
Type: string\\
Default: nil
}

{\footnotesize Використовуються методи для знищення інформації після закінчення терміну зберігання}


{\footnotesize\bfseries
No: 2\\
Name: si\_\allowbreak{}12\_\allowbreak{}3\_\allowbreak{}02\\
Type: string\\
Default: nil
}

{\footnotesize Використовуються методи для знищення інформації після закінчення терміну зберігання}


{\footnotesize\bfseries
No: 3\\
Name: si\_\allowbreak{}12\_\allowbreak{}3\_\allowbreak{}03\\
Type: integer\\
Default: 30
}

{\footnotesize Використовуються методи для стирання інформації після закінчення періоду зберігання}




\subsection{ПЕРЕДБАЧУВАНЕ ЗАПОБІГАННЯ ЗБОЇВ (SI-13)}
a. Визначити середній час до збою (MTTF) для [Призначення: визначені організацією компоненти системи] в певних середовищах роботи.\\ 
b. Надати замінні компоненти системи та засоби для заміни активних компонентів резервними компонентами відповідно до [Призначення: визначені організацією критерії заміни].

\vspace{10pt}
{\footnotesize\bfseries
No: 1\\
Name: si\_\allowbreak{}13\_\allowbreak{}odp\_\allowbreak{}01\\
Type: integer\\
Default: 30
}

{\footnotesize Визначені компоненти системи, для яких необхідно визначити середній час до збою (MTTF)}


{\footnotesize\bfseries
No: 2\\
Name: si\_\allowbreak{}13\_\allowbreak{}odp\_\allowbreak{}02\\
Type: list\\
Default: []
}

{\footnotesize Визначені критерії заміни за середнім часом напрацювання до збою (MTTF), які будуть використовуватися для заміни активних і резервних компонентів; SI-13a. визначено середній час напрацювання до збою (MTTF) для системних компонентів у конкретних умовах експлуатації; SI-13b. передбачені замінні компоненти системи та засоби заміни активних і резервних компонентів відповідно до критеріїв заміни середнього часу напрацювання на відмову (MTTF) }




\subsubsection{ВІДПОВІДАЛЬНІСТЬ ЗА ПЕРЕДАЧУ ФУНКЦІЙ КОМПОНЕНТІВ (SI-13(1))}


\vspace{10pt}
{\footnotesize\bfseries
No: 1\\
Name: si\_\allowbreak{}13\_\allowbreak{}1\_\allowbreak{}01\\
Type: integer\\
Default: 30
}

{\footnotesize Виводяться компоненти системи з експлуатації шляхом передачі обов'язків компонентів на запасні компоненти не пізніше, ніж частка або відсоток середнього напрацювання на відмову}


{\footnotesize\bfseries
No: 2\\
Name: si\_\allowbreak{}13\_\allowbreak{}1\_\allowbreak{}odp\\
Type: integer\\
Default: 30
}

{\footnotesize Визначено частку або напрацювання до збою, відсоток середнього часу в межах якого обов'язки компонента системи можуть бути передані компоненту, що замінює його}




\subsubsection{ТЕРМІН ВИКОНАННЯ ПРОЦЕСУ БЕЗ НАГЛЯДУ (SI-13(2)) [Вилучено]}
[Вилучено: включено до SI-7 (16)]

\vspace{10pt}
\textit{Немає параметрів для цього контролю.}
\vspace{10pt}


\subsubsection{ЗАПОБІГАННЯ ПЕРЕДБАЧУВАНИХ ФУНКЦІЙ КОМПОНЕНТІВ (SI-13(3))}


\vspace{10pt}
{\footnotesize\bfseries
No: 1\\
Name: si\_\allowbreak{}13\_\allowbreak{}3\_\allowbreak{}01\\
Type: integer\\
Default: 30
}

{\footnotesize Ініціюються вручну передачі між активним та резервним компонентами системи, коли використання активного компонента досягає відсоток від середнього часу напрацювання до збою}


{\footnotesize\bfseries
No: 2\\
Name: si\_\allowbreak{}13\_\allowbreak{}3\_\allowbreak{}odp\\
Type: integer\\
Default: 30
}

{\footnotesize Визначено відсоток середнього часу напрацювання на відмову для передач, які потрібно ініціювати вручну}




\subsubsection{ЗАПОБІГАННЯ ПЕРЕДБАЧУВАНИХ ЗБОЇВ РЕЗЕРВНИХ КОМПОНЕНТІВ ТА ОПОВІЩЕННЯ (SI-13(4))}


\vspace{10pt}
{\footnotesize\bfseries
No: 1\\
Name: si\_\allowbreak{}13\_\allowbreak{}4\_\allowbreak{}a\\
Type: integer\\
Default: 30
}

{\footnotesize Успішно і прозоро встановлюються резервні компоненти протягом часу, якщо виявлено збої у роботі системних компонентів}




\subsubsection{ЗАПОБІГАННЯ ПЕРЕДБАЧУВАНИХ АВАРІЙНОГО ПЕРЕМИКАННЯ (SI-13(5))}


\vspace{10pt}
\textit{Немає параметрів для цього контролю.}
\vspace{10pt}


\subsection{НЕСТІЙКІСТЬ (SI-14)}
Реалізувати нестійкі [Призначення: визначені організацією компоненти системи та служби], які ініціюються у відомих станах і завершуються [Вибір (один або кілька): після закінчення сеансу використання; періодично з [Призначення: визначена організацією частота]].

\vspace{10pt}
{\footnotesize\bfseries
No: 1\\
Name: si\_\allowbreak{}14\_\allowbreak{}01\\
Type: string\\
Default: nil
}

{\footnotesize Реалізовано непостійні компоненти системи та сервіси, які ініціюються у відомому стані}


{\footnotesize\bfseries
No: 2\\
Name: si\_\allowbreak{}14\_\allowbreak{}odp\_\allowbreak{}01\\
Type: string\\
Default: nil
}

{\footnotesize Визначені непостійні компоненти системи та сервіси, які необхідно застосовувати}


{\footnotesize\bfseries
No: 3\\
Name: si\_\allowbreak{}14\_\allowbreak{}odp\_\allowbreak{}02\\
Type: string\\
Default: \textquotedbl{}щорічно\textquotedbl{}
}

{\footnotesize Вибрано одне або декілька з наступних ЗНАЧЕНЬ ПАРАМЕТРІВ: \{по закінченні сеансу використання; частота\}}


{\footnotesize\bfseries
No: 4\\
Name: si\_\allowbreak{}14\_\allowbreak{}odp\_\allowbreak{}03\\
Type: integer\\
Default: 30
}

{\footnotesize Визначено частоту завершення роботи непостійних компонентів і сервісів, які ініціюються у відомому стані (якщо вибрано)}




\subsubsection{ОНОВЛЕННЯ З НАДІЙНИХ ДЖЕРЕЛ (SI-14(1))}


\vspace{10pt}
{\footnotesize\bfseries
No: 1\\
Name: si\_\allowbreak{}14\_\allowbreak{}1\_\allowbreak{}01\\
Type: integer\\
Default: 30
}

{\footnotesize Програмне забезпечення та дані, що використовуються під час оновлення системних компонентів та служб, отримані з довірених джерел}


{\footnotesize\bfseries
No: 2\\
Name: si\_\allowbreak{}14\_\allowbreak{}1\_\allowbreak{}odp\\
Type: string\\
Default: nil
}

{\footnotesize Визначені надійні джерела для отримання програмного забезпечення та даних для оновлення системних компонентів і служб}




\subsubsection{НЕСТІЙКА ІНФОРМАЦІЯ (SI-14(2))}


\vspace{10pt}
{\footnotesize\bfseries
No: 1\\
Name: si\_\allowbreak{}14\_\allowbreak{}2\_\allowbreak{}b\\
Type: string\\
Default: nil
}

{\footnotesize Видаляється інформація, коли вона більше не потрібна}




\subsubsection{НЕСТІЙКІ ПІДКЛЮЧЕННЯ (SI-14(3))}


\vspace{10pt}
{\footnotesize\bfseries
No: 1\\
Name: si\_\allowbreak{}14\_\allowbreak{}3\_\allowbreak{}01\\
Type: string\\
Default: nil
}

{\footnotesize Встановлюються з'єднання з системою на вимогу}


{\footnotesize\bfseries
No: 2\\
Name: si\_\allowbreak{}14\_\allowbreak{}3\_\allowbreak{}odp\\
Type: integer\\
Default: 30
}

{\footnotesize Вибрано одне з наступних ЗНАЧЕНЬ ПАРАМЕТРА: \{завершення запиту; період невикористання\}}




\subsection{ФІЛЬТРАЦІЯ ВИХІДНИХ ДАНИХ (SI-15)}
Перевіряти інформацію, що виходить з [Призначення: визначені організацією програмні продукти та/або застосунки], щоб переконатися, що інформація відповідає очікуваному змісту.

\vspace{10pt}
{\footnotesize\bfseries
No: 1\\
Name: si\_\allowbreak{}15\_\allowbreak{}01\\
Type: string\\
Default: nil
}

{\footnotesize Перевіряється інформація, що виводиться з програмне забезпечення та/або додатки, щоб переконатися, що інформація відповідає очікуваному змісту}


{\footnotesize\bfseries
No: 2\\
Name: si\_\allowbreak{}15\_\allowbreak{}odp\\
Type: string\\
Default: nil
}

{\footnotesize Визначені програми та/або додатки, виведення інформації з яких потребує перевірки}




\subsection{ЗАХИСТ ПАМ'ЯТІ (SI-16)}
Виконати [Призначення: визначені організацією заходи безпеки] для захисту системної пам'яті від несанкціонованого коду, що виконується.

\vspace{10pt}
{\footnotesize\bfseries
No: 1\\
Name: si\_\allowbreak{}16\_\allowbreak{}01\\
Type: string\\
Default: nil
}

{\footnotesize Реалізовано контроль для захисту системної пам'яті від несанкціонованого виконання коду}


{\footnotesize\bfseries
No: 2\\
Name: si\_\allowbreak{}16\_\allowbreak{}odp\\
Type: string\\
Default: \textquotedbl{}автоматизований засіб моніторингу\textquotedbl{}
}

{\footnotesize Визначено засоби контролю для захисту системної пам'яті від несанкціонованого виконання коду}




\subsection{ВІДМОВОСТІЙКІ ПРОЦЕДУРИ (SI-17)}
Виконати [Призначення: визначені організацією відмовостійкі процедури], коли настають [Призначення: визначені організацією умови виявлення несправностей].

\vspace{10pt}
{\footnotesize\bfseries
No: 1\\
Name: si\_\allowbreak{}17\_\allowbreak{}01\\
Type: list\\
Default: []
}

{\footnotesize Реалізовано процедури захисту від збоїв при виникненні перелік умов збою. відмови, що вимагають}


{\footnotesize\bfseries
No: 2\\
Name: si\_\allowbreak{}17\_\allowbreak{}odp\_\allowbreak{}01\\
Type: string\\
Default: nil
}

{\footnotesize Визначені відмовостійкі процедури, пов'язані з умовами відмови}


{\footnotesize\bfseries
No: 3\\
Name: si\_\allowbreak{}17\_\allowbreak{}odp\_\allowbreak{}02\\
Type: list\\
Default: []
}

{\footnotesize Визначено перелік умов відмовостійких процедур}




\subsection{ОПЕРАЦІЇ ЗАБЕЗПЕЧЕННЯ ЯКОСТІ ДАНИХ (SI-18)}
a. Перевіряти точність, актуальність, своєчасність і повноту персональної інформації протягом її життєвого циклу [Завдання: частота, визначена організацією];\\ 
b. Виправляти або видаляти неточну або застарілу персональну інформацію.

\vspace{10pt}
{\footnotesize\bfseries
No: 1\\
Name: si\_\allowbreak{}18\_\allowbreak{}odp\_\allowbreak{}01\\
Type: list\\
Default: [\textquotedbl{}admin\textquotedbl{}, \textquotedbl{}security\_\allowbreak{}officer\textquotedbl{}]
}

{\footnotesize Визначено періодичність перевірки точності персональну інформацію протягом життєвого циклу інформації}


{\footnotesize\bfseries
No: 2\\
Name: si\_\allowbreak{}18\_\allowbreak{}odp\_\allowbreak{}02\\
Type: list\\
Default: [\textquotedbl{}admin\textquotedbl{}, \textquotedbl{}security\_\allowbreak{}officer\textquotedbl{}]
}

{\footnotesize Визначено періодичність перевірки актуальності персональної інформації протягом життєвого циклу інформації}


{\footnotesize\bfseries
No: 3\\
Name: si\_\allowbreak{}18\_\allowbreak{}odp\_\allowbreak{}03\\
Type: list\\
Default: [\textquotedbl{}admin\textquotedbl{}, \textquotedbl{}security\_\allowbreak{}officer\textquotedbl{}]
}

{\footnotesize Визначено періодичність перевірки актуальності персональної інформації протягом життєвого циклу інформації}


{\footnotesize\bfseries
No: 4\\
Name: si\_\allowbreak{}18\_\allowbreak{}odp\_\allowbreak{}04\\
Type: list\\
Default: [\textquotedbl{}admin\textquotedbl{}, \textquotedbl{}security\_\allowbreak{}officer\textquotedbl{}]
}

{\footnotesize Визначено періодичність перевірки повноти персональної інформації протягом життєвого циклу інформації; SI-18a.[01] перевіряється точність персональної інформації протягом життєвого циклу інформації частота; SI-18a.[02] перевіряється актуальність персональної інформації протягом життєвого циклу інформації < SI-18\_\allowbreak{}ODP[02] частота>; SI-18a.[03] перевіряється своєчасність персональної інформації протягом життєвого циклу інформації <частота SI-18\_\allowbreak{}ODP[03]>; SI-18a.[04] перевіряється повнота персональної інформації протягом життєвого циклу інформації частота; SI-18b. потрібно виправити або видалити неточну або застарілу персональну інформацію}




\subsubsection{ОПЕРАЦІЇ ЗАБЕЗПЕЧЕННЯ ПІДТРИМКА (SI-18(1))}


\vspace{10pt}
{\footnotesize\bfseries
No: 1\\
Name: si\_\allowbreak{}18\_\allowbreak{}1\_\allowbreak{}01\\
Type: list\\
Default: [\textquotedbl{}admin\textquotedbl{}, \textquotedbl{}security\_\allowbreak{}officer\textquotedbl{}]
}

{\footnotesize Використовуються автоматизовані механізми для виправлення або видалення персональної інформації яка є неточною, застарілою, неправильно визначеною щодо впливу або неправильно деідентифікованою}


{\footnotesize\bfseries
No: 2\\
Name: si\_\allowbreak{}18\_\allowbreak{}1\_\allowbreak{}odp\\
Type: list\\
Default: [\textquotedbl{}admin\textquotedbl{}, \textquotedbl{}security\_\allowbreak{}officer\textquotedbl{}]
}

{\footnotesize Визначені автоматизовані механізми, які використовуються для виправлення або видалення персональної інформації яка є неточною, застарілою, неправильно визначеною щодо впливу або неправильно деідентифікованою}




\subsubsection{ТЕГУВАННЯ ДАНИХ (SI-18(2))}


\vspace{10pt}
\textit{Немає параметрів для цього контролю.}
\vspace{10pt}


\subsubsection{ЗБИРАННЯ (SI-18(3))}


\vspace{10pt}
\textit{Немає параметрів для цього контролю.}
\vspace{10pt}


\subsubsection{ОПЕРАЦІЇ ЗАПИТИ (SI-18(4))}


\vspace{10pt}
{\footnotesize\bfseries
No: 1\\
Name: si\_\allowbreak{}18\_\allowbreak{}4\_\allowbreak{}01\\
Type: list\\
Default: [\textquotedbl{}admin\textquotedbl{}, \textquotedbl{}security\_\allowbreak{}officer\textquotedbl{}]
}

{\footnotesize Виправляється або видаляється персональну інформація на вимогу осіб або їхніх уповноважених представників. ПОТЕНЦІЙНІ МЕТОДИ ТА ОБ'ЄКТИ ОЦІНКИ:}




\subsubsection{ПОВІДОМЛЕННЯ ПРО ВИПРАВЛЕННЯ ЧИ ВИДАЛЕННЯ (SI-18(5))}


\vspace{10pt}
{\footnotesize\bfseries
No: 1\\
Name: si\_\allowbreak{}18\_\allowbreak{}5\_\allowbreak{}01\\
Type: list\\
Default: [\textquotedbl{}admin\textquotedbl{}, \textquotedbl{}security\_\allowbreak{}officer\textquotedbl{}]
}

{\footnotesize Одержувачі та фізичні особи повідомляються про виправлення або видалення персональної інформації}


{\footnotesize\bfseries
No: 2\\
Name: si\_\allowbreak{}18\_\allowbreak{}5\_\allowbreak{}odp\\
Type: list\\
Default: [\textquotedbl{}admin\textquotedbl{}, \textquotedbl{}security\_\allowbreak{}officer\textquotedbl{}]
}

{\footnotesize Визначені одержувачі персональних даних, які повинні бути повідомлені про виправлення або видалення персональних даних}




\subsection{ДЕІДЕНТИФІКАЦІЯ (SI-19)}
a. Видаліть такі елементи персональних даних з наборів даних: [Призначення: визначені організацією елементи персональних даних];\\ 
b. Оцініть [Призначення: деідентифікації. частота, визначена організацією] ефективність

\vspace{10pt}
{\footnotesize\bfseries
No: 1\\
Name: si\_\allowbreak{}19\_\allowbreak{}odp\_\allowbreak{}02\\
Type: string\\
Default: \textquotedbl{}щорічно\textquotedbl{}
}

{\footnotesize Визначено частоту, з якою слід оцінювати ефективність деідентифікації; SI-19a. вилучено елементи з наборів даних; SI-19b. оцінюється ефективність деідентифікації частота}




\subsubsection{ЗБІР (SI-19(1))}


\vspace{10pt}
{\footnotesize\bfseries
No: 1\\
Name: si\_\allowbreak{}19\_\allowbreak{}1\_\allowbreak{}01\\
Type: list\\
Default: [\textquotedbl{}admin\textquotedbl{}, \textquotedbl{}security\_\allowbreak{}officer\textquotedbl{}]
}

{\footnotesize Деідентифікується набір даних після збору шляхом відмови від збору персональної інформації}




\subsubsection{АРХІВАЦІЯ (SI-19(2))}


\vspace{10pt}
{\footnotesize\bfseries
No: 1\\
Name: si\_\allowbreak{}19\_\allowbreak{}2\_\allowbreak{}01\\
Type: list\\
Default: [\textquotedbl{}admin\textquotedbl{}, \textquotedbl{}security\_\allowbreak{}officer\textquotedbl{}]
}

{\footnotesize Заборонено архівування елементів персональної інформації, якщо ці елементи в наборі даних не будуть потрібні після того, як набір даних буде заархівовано}




\subsubsection{ВИДАЛЕННЯ (SI-19(3))}


\vspace{10pt}
{\footnotesize\bfseries
No: 1\\
Name: si\_\allowbreak{}19\_\allowbreak{}3\_\allowbreak{}01\\
Type: list\\
Default: [\textquotedbl{}admin\textquotedbl{}, \textquotedbl{}security\_\allowbreak{}officer\textquotedbl{}]
}

{\footnotesize Видаляються елементи персональної інформаціїз набору даних перед його оприлюдненням, якщо ці елементи в наборі даних не повинні бути частиною оприлюднення даних}




\subsubsection{ВИДАЛЕННЯ, МАСКУВАННЯ, ШИФРУВАННЯ, ХЕШУВАННЯ АБО ЗАМІНА ПРЯМИХ ІДЕНТИФІКАТОРІВ (SI-19(4))}


\vspace{10pt}
\textit{Немає параметрів для цього контролю.}
\vspace{10pt}


\subsubsection{КОНТРОЛЬ СТАТИСТИЧНОГО РОЗКРИТТЯ (SI-19(5))}


\vspace{10pt}
{\footnotesize\bfseries
No: 1\\
Name: si\_\allowbreak{}19\_\allowbreak{}5\_\allowbreak{}01\\
Type: list\\
Default: [\textquotedbl{}admin\textquotedbl{}, \textquotedbl{}security\_\allowbreak{}officer\textquotedbl{}]
}

{\footnotesize Не маніпулюють числовими даними так, щоб у результатах аналізу не можна було ідентифікувати жодну особу чи організацію}


{\footnotesize\bfseries
No: 2\\
Name: si\_\allowbreak{}19\_\allowbreak{}5\_\allowbreak{}02\\
Type: list\\
Default: [\textquotedbl{}admin\textquotedbl{}, \textquotedbl{}security\_\allowbreak{}officer\textquotedbl{}]
}

{\footnotesize Не маніпулюють таблицями непередбачених обставин таким чином, щоб у результатах аналізу не можна було ідентифікувати жодну особу чи організацію}


{\footnotesize\bfseries
No: 3\\
Name: si\_\allowbreak{}19\_\allowbreak{}5\_\allowbreak{}03\\
Type: list\\
Default: [\textquotedbl{}admin\textquotedbl{}, \textquotedbl{}security\_\allowbreak{}officer\textquotedbl{}]
}

{\footnotesize Не маніпулюють статистичними даними так, щоб за результатами аналізу не можна було ідентифікувати жодну особу чи організацію}




\subsubsection{ДИФЕРЕНЦІЙОВАНА КОНФІДЕНЦІЙНІСТЬ (SI-19(6))}


\vspace{10pt}
{\footnotesize\bfseries
No: 1\\
Name: si\_\allowbreak{}19\_\allowbreak{}6\_\allowbreak{}01\\
Type: list\\
Default: [\textquotedbl{}admin\textquotedbl{}, \textquotedbl{}security\_\allowbreak{}officer\textquotedbl{}]
}

{\footnotesize Запобігає розголошенню персональної інформації, додавання недетермінованого шуму до результатів математичних операцій до того, як результати будуть повідомлені}




\subsubsection{ПЕРЕВІРЕНЕ ПРОГРАМНЕ ЗАБЕЗПЕЧЕННЯ (SI-19(7))}


\vspace{10pt}
{\footnotesize\bfseries
No: 1\\
Name: si\_\allowbreak{}19\_\allowbreak{}7\_\allowbreak{}01\\
Type: string\\
Default: \textquotedbl{}AES-256-GCM\textquotedbl{}
}

{\footnotesize Виконується деідентифікація за допомогою перевірених алгоритмів}


{\footnotesize\bfseries
No: 2\\
Name: si\_\allowbreak{}19\_\allowbreak{}7\_\allowbreak{}02\\
Type: string\\
Default: \textquotedbl{}AES-256-GCM\textquotedbl{}
}

{\footnotesize Виконується деідентифікація за допомогою програмного забезпечення, яке пройшло валідацію для реалізації алгоритмів}




\subsubsection{МОТИВОВАНИЙ ПОРУШНИК (SI-19(8))}


\vspace{10pt}
{\footnotesize\bfseries
No: 1\\
Name: si\_\allowbreak{}19\_\allowbreak{}8\_\allowbreak{}01\\
Type: string\\
Default: nil
}

{\footnotesize Виконується тест мотивованого зловмисника для деідентифікованого набору даних, щоб визначити, чи залишаються ідентифіковані дані або чи можна повторно ідентифікувати деідентифіковані дані}




\subsection{ПСУВАННЯ (SI-20)}
Вбудуйте дані або можливості в такі системи або системні компоненти, щоб визначити, чи дані організації були викрадені або неналежним чином видалені з організації: [Призначення: визначені організацією системи або системні компоненти].

\vspace{10pt}
{\footnotesize\bfseries
No: 1\\
Name: si\_\allowbreak{}20\_\allowbreak{}01\\
Type: string\\
Default: nil
}

{\footnotesize Вбудовані дані або можливості в системи або компоненти системи, щоб визначити, чи були дані організації викрадені або неналежним чином видалені з організації}


{\footnotesize\bfseries
No: 2\\
Name: si\_\allowbreak{}20\_\allowbreak{}odp\\
Type: string\\
Default: nil
}

{\footnotesize Визначені системи або компоненти системи з даними або можливостями, що підлягають застосуванню}




\subsection{ОНОВЛЕННЯ ІНФОРМАЦІЇ (SI-21)}
Оновлюйте [Призначення: інформація, визначена організацією] з [Призначення: частота, визначена організацією] або згенеруйте інформацію за запитом і видаліть її, коли в ній більше не буде потреби.

\vspace{10pt}
{\footnotesize\bfseries
No: 1\\
Name: si\_\allowbreak{}21\_\allowbreak{}01\\
Type: string\\
Default: \textquotedbl{}щорічно\textquotedbl{}
}

{\footnotesize Інформація оновлюється частота або генерується на вимогу і видаляється, коли більше не потрібна. з якими потрібно оновлювати}


{\footnotesize\bfseries
No: 2\\
Name: si\_\allowbreak{}21\_\allowbreak{}odp\_\allowbreak{}01\\
Type: string\\
Default: nil
}

{\footnotesize Визначена інформація, яку потрібно оновити}


{\footnotesize\bfseries
No: 3\\
Name: si\_\allowbreak{}21\_\allowbreak{}odp\_\allowbreak{}02\\
Type: integer\\
Default: 30
}

{\footnotesize Визначені частоти, інформацію}




\subsection{РІЗНОВИДИ ІНФОРМАЦІЇ (SI-22)}
a. Визначити наступні альтернативні джерела інформації для [Завдання: основні функції та послуги, визначені організацією]: [Завдання: альтернативні, визначені організацією джерела інформації];\\ 
b. Використовуйте альтернативне джерело інформації для виконання основних функцій або послуг на [Призначення: визначені організацією системи або системні компоненти], коли основне джерело інформації пошкоджено або недоступне.

\vspace{10pt}
{\footnotesize\bfseries
No: 1\\
Name: si\_\allowbreak{}22\_\allowbreak{}odp\_\allowbreak{}01\\
Type: string\\
Default: nil
}

{\footnotesize Визначені альтернативні джерела інформації для основних функцій та послуг}


{\footnotesize\bfseries
No: 2\\
Name: si\_\allowbreak{}22\_\allowbreak{}odp\_\allowbreak{}02\\
Type: string\\
Default: nil
}

{\footnotesize Визначені основні функції та послуги, які потребують альтернативних джерел інформації}


{\footnotesize\bfseries
No: 3\\
Name: si\_\allowbreak{}22\_\allowbreak{}odp\_\allowbreak{}03\\
Type: string\\
Default: nil
}

{\footnotesize Визначені системи або компоненти системи, які потребують альтернативного джерела інформації для виконання основних функцій або послуг; SI-22a. визначені альтернативні джерела інформації для основних функцій та послуг; SI-22b. використовується альтернативне джерело інформації для виконання основних функцій або послуг у системах або компонентах системи, коли первинне джерело інформації пошкоджене або недоступне}




\subsection{ФРАГМЕНТАЦІЯ ІНФОРМАЦІЇ (SI-23)}


\vspace{10pt}
{\footnotesize\bfseries
No: 1\\
Name: si\_\allowbreak{}23\_\allowbreak{}odp\_\allowbreak{}01\\
Type: string\\
Default: nil
}

{\footnotesize Визначені обставини, інформації; які вимагають фрагментації}


{\footnotesize\bfseries
No: 2\\
Name: si\_\allowbreak{}23\_\allowbreak{}odp\_\allowbreak{}02\\
Type: string\\
Default: nil
}

{\footnotesize Визначена інформація, яка підлягає фрагментації}


{\footnotesize\bfseries
No: 3\\
Name: si\_\allowbreak{}23\_\allowbreak{}odp\_\allowbreak{}03\\
Type: string\\
Default: nil
}

{\footnotesize Визначені системи або компоненти системи, між якими має бути розподілена фрагментована інформація; SI-23a. за обставин, інформація є фрагментованою; SI-23b. за обставин фрагментована інформація розподіляється між системами або компонентами системи}




\section{SR}
Клас заходів захисту SR --- УПРАВЛІННЯ РИЗИКАМИ ЛАНЦЮГА

Цей клас встановлює вимоги для виявлення та мінімізації загроз, що виникають через зовнішніх постачальників продуктів та послуг.

\textbf{Перелік заходів захисту:} Політика та процедури управління ризиками ланцюга постачання (SR-1); План управління ризиками ланцюга постачання (SR-2); Створення команди постачання (SR-2(1)); Контроль ланцюга постачання і процесів (SR-3); Різні бази постачання (SR-3(1)); Обмеження шкоди (SR-3(2)); Перенесення заходів захисту управління ризиками ланцюга постачання до субпідрядників (SR-3(3)); Походження (SR-4); Ідентичність (SR-4(1)); Унікальна ідентифікація (SR-4(2)); Перевірка на справжність і відсутність внесення змін (SR-4(3)); Походження -- перевірка ланцюга цілісності (SR-4(4)); Стратегії придбання, інструменти і методи (SR-5); Належне постачання (SR-5(1)); Оцінка перед відбором, прийняття, модифікація чи оновлення (SR-5(2)); Оцінка постачальників (SR-6); Тестування та аналіз (SR-6(1)); Безпека операцій ланцюга постачання (SR-7); Повідомлення про порушення ланцюга постачання (SR-8); Захист від злому та виявлення (SR-9); Етапи чи системи розвитку життєвого циклу (SR-9(1)); Перевірка системи і компонентів системи (SR-10); Автентичність компоненту (SR-11); Автентичність компоненту (SR-11(1)); Утилізація компоненту (SR-12).

\subsection{ПОЛІТИКА ТА ПРОЦЕДУРИ УПРАВЛІННЯ РИЗИКАМИ ЛАНЦЮГА ПОСТАЧАННЯ (SR-1)}
a. Розробіть, задокументуйте та поширте [Призначення: персонал або ролі, визначені організацією]:\\ 
1. [Вибір (один або декілька): Рівень організації; Рівень місії/бізнес-процесу; рівень системи] політика управління ризиками ланцюга постачання, яка: a) Розглядає мету, сферу діяльності, ролі, відповідальність, зобов'язання керівництва, координацію між організаційними підрозділами та відповідність; b) Відповідає чинним законам, виконавчим наказам, директивам, положенням, політикам, стандартам і вказівкам;\\ 
2. Процедури для сприяння впровадженню політики управління ризиками ланцюга постачання та відповідних засобів контролю управління ризиками ланцюга постачання;\\ 
b. Призначити [Призначення: посадова особа, визначена організацією] для управління розробкою, документуванням і розповсюдженням політики та процедур управління ризиками ланцюга постачання;\\ 
c. Перегляньте та оновіть поточне управління ризиками ланцюга постачання:\\ 
1. Політика [Призначення: частота, визначена організацією] та наступне [Призначення: події, визначені організацією];\\ 
2. Процедури [Призначення: частота, визначена організацією] та наступні [Призначення: події, визначені організацією].

\vspace{10pt}
{\footnotesize\bfseries
No: 1\\
Name: sr\_\allowbreak{}1\_\allowbreak{}odp\_\allowbreak{}01\\
Type: list\\
Default: [\textquotedbl{}admin\textquotedbl{}, \textquotedbl{}security\_\allowbreak{}officer\textquotedbl{}]
}

{\footnotesize Визначено персонал або ролі, на які поширюється політика управління ризиками ланцюга постачання}


{\footnotesize\bfseries
No: 2\\
Name: sr\_\allowbreak{}1\_\allowbreak{}odp\_\allowbreak{}02\\
Type: list\\
Default: [\textquotedbl{}admin\textquotedbl{}, \textquotedbl{}security\_\allowbreak{}officer\textquotedbl{}]
}

{\footnotesize Визначено персонал або ролі, на які поширюються процедури управління ризиками ланцюга постачання}


{\footnotesize\bfseries
No: 3\\
Name: sr\_\allowbreak{}1\_\allowbreak{}odp\_\allowbreak{}03\\
Type: string\\
Default: nil
}

{\footnotesize Вибрано одне або декілька з наступних ЗНАЧЕНЬ ПАРАМЕТРІВ: \{рівень організації; рівень завдань/бізнеспроцесу; рівень системи\}}


{\footnotesize\bfseries
No: 4\\
Name: sr\_\allowbreak{}1\_\allowbreak{}odp\_\allowbreak{}04\\
Type: list\\
Default: [\textquotedbl{}admin\textquotedbl{}, \textquotedbl{}security\_\allowbreak{}officer\textquotedbl{}]
}

{\footnotesize Визначена посадова особа, відповідальна за розробку, документування та розповсюдження політики та процедур управління ризиками ланцюга постачання}


{\footnotesize\bfseries
No: 5\\
Name: sr\_\allowbreak{}1\_\allowbreak{}odp\_\allowbreak{}05\\
Type: list\\
Default: [\textquotedbl{}default\_\allowbreak{}deny\_\allowbreak{}rule\textquotedbl{}, \textquotedbl{}abac\_\allowbreak{}rule\_\allowbreak{}1\textquotedbl{}]
}

{\footnotesize Визначена періодичність перегляду та оновлення поточної політики управління ризиками ланцюга постачання}


{\footnotesize\bfseries
No: 6\\
Name: sr\_\allowbreak{}1\_\allowbreak{}odp\_\allowbreak{}06\\
Type: list\\
Default: [\textquotedbl{}default\_\allowbreak{}deny\_\allowbreak{}rule\textquotedbl{}, \textquotedbl{}abac\_\allowbreak{}rule\_\allowbreak{}1\textquotedbl{}]
}

{\footnotesize Є події, які вимагають перегляду та оновлення поточної політики управління ризиками ланцюга постачання}


{\footnotesize\bfseries
No: 7\\
Name: sr\_\allowbreak{}1\_\allowbreak{}odp\_\allowbreak{}07\\
Type: string\\
Default: \textquotedbl{}щорічно\textquotedbl{}
}

{\footnotesize Визначена періодичність перегляду та оновлення поточної процедури управління ризиками ланцюга постачання}




\subsection{ПЛАН УПРАВЛІННЯ РИЗИКАМИ ЛАНЦЮГА ПОСТАЧАННЯ (SR-2)}
a. Розробіть план управління ризиками ланцюга постачання, пов'язаними з дослідженнями та розробкою, проектуванням, виробництвом, придбанням, доставкою, інтеграцією, експлуатацією та обслуговуванням, а також утилізацією таких систем, компонентів системи або послуг для системи: [Призначення: системи, визначені організацією, системні компоненти або системні служби];\\ 
b. Перегляньте та оновіть план управління ризиками ланцюга постачання [Призначення: частота, визначена організацією] або за потреби для усунення загроз;\\ 
c. Захистіть план управління ризиками ланцюга постачання від несанкціонованого розголошення та модифікації.

\vspace{10pt}
{\footnotesize\bfseries
No: 1\\
Name: sr\_\allowbreak{}2\_\allowbreak{}odp\_\allowbreak{}01\\
Type: string\\
Default: nil
}

{\footnotesize Визначені системи, компоненти системи або системні послуги, для яких розробляється план управління ризиками ланцюга постачання}




\subsubsection{СТВОРЕННЯ КОМАНДИ ПОСТАЧАННЯ (SR-2(1))}


\vspace{10pt}
{\footnotesize\bfseries
No: 1\\
Name: sr\_\allowbreak{}2\_\allowbreak{}1\_\allowbreak{}01\\
Type: list\\
Default: [\textquotedbl{}admin\textquotedbl{}, \textquotedbl{}security\_\allowbreak{}officer\textquotedbl{}]
}

{\footnotesize Створена команда з управління ризиками ланцюга постачання, що складається з персонал, ролі та обов'язки для керівництва та підтримки діяльності з управління ризиками ланцюга постачання}




\subsection{КОНТРОЛЬ ЛАНЦЮГА ПОСТАЧАННЯ І ПРОЦЕСІВ (SR-3)}
a. Встановлення процесу або процесів для виявлення та усунення слабких місць або недоліків в елементах і процесах ланцюга постачання [Призначення: визначена організацією система або компонент системи] у координації з [Завдання: персонал ланцюга постачання, визначений організацією];\\ 
b. Використовуйте такі заходи захисту, щоб захистити систему, компонент системи або системну службу від ризиків ланцюга постачання та обмежити шкоду чи наслідки від подій, пов'язаних із ланцюгом постачання: [Призначення: заходи захисту ланцюга постачання, визначені організацією];\\ 
c. Задокументуйте обрані та впроваджені процеси та заходи захисту ланцюгом постачання у [Вибір: плани безпеки та приватності; план управління ризиками ланцюга постачання; [Призначення: документ, визначений організацією]].

\vspace{10pt}
{\footnotesize\bfseries
No: 1\\
Name: sr\_\allowbreak{}3\_\allowbreak{}odp\_\allowbreak{}01\\
Type: string\\
Default: nil
}

{\footnotesize Визначено систему або компонент системи, який потребує процесу або процесів для виявлення та усунення слабких місць або недоліків}


{\footnotesize\bfseries
No: 2\\
Name: sr\_\allowbreak{}3\_\allowbreak{}odp\_\allowbreak{}02\\
Type: list\\
Default: [\textquotedbl{}admin\textquotedbl{}, \textquotedbl{}security\_\allowbreak{}officer\textquotedbl{}]
}

{\footnotesize Визначено персонал ланцюга поставок, з яким необхідно координувати процес або процеси виявлення та усунення слабких місць або недоліків в елементах і процесах ланцюга постачання}


{\footnotesize\bfseries
No: 3\\
Name: sr\_\allowbreak{}3\_\allowbreak{}odp\_\allowbreak{}03\\
Type: string\\
Default: \textquotedbl{}автоматизований засіб моніторингу\textquotedbl{}
}

{\footnotesize Визначені засоби контролю ланцюга постачання, що застосовуються для захисту від ризиків ланцюга постачання для системи, системного компонента або системної послуги, а також для обмеження шкоди або наслідків від подій, пов'язаних з ланцюгом постачання}


{\footnotesize\bfseries
No: 4\\
Name: sr\_\allowbreak{}3\_\allowbreak{}odp\_\allowbreak{}04\\
Type: list\\
Default: [\textquotedbl{}admin\textquotedbl{}, \textquotedbl{}security\_\allowbreak{}officer\textquotedbl{}]
}

{\footnotesize Вибрано одне або декілька з наступних ЗНАЧЕНЬ ПАРАМЕТРІВ: \{плани безпеки та конфіденційності; план управління ризиками ланцюга постачання; документ\}; SR-03\_\allowbreak{}ODP[05] визначено документ, що ідентифікує обрані та впроваджені процеси та засоби контролю ланцюга постачання (якщо обрано); SR-03a.[01] запроваджено процес або процеси для виявлення та усунення слабких місць або недоліків в елементах та процесах ланцюга постачання; SR-03a.[02] процес або процеси виявлення та усунення слабких місць або недоліків в елементах та процесах ланцюга постачання системи або компонента системи координується/координуються з персоналом ланцюга постачання; SR-03b. застосовуються засоби контролю ланцюга постачання для захисту від ризиків ланцюга постачання для системи, системного компонента або системної послуги, а також для обмеження шкоди або наслідків від подій, пов'язаних з ланцюгом постачання; SR-03c. задокументовані обрані та впроваджені процеси та засоби контролю ланцюга постачання в ЗНАЧЕННЯ ВИБІРКОВОГО ПАРАМЕТРА(ів)}


{\footnotesize\bfseries
No: 5\\
Name: sr\_\allowbreak{}3\_\allowbreak{}odp\_\allowbreak{}05\\
Type: list\\
Default: [\textquotedbl{}admin\textquotedbl{}, \textquotedbl{}security\_\allowbreak{}officer\textquotedbl{}]
}

{\footnotesize Визначено документ, що ідентифікує обрані та впроваджені процеси та засоби контролю ланцюга постачання (якщо обрано); SR-03a.[01] запроваджено процес або процеси для виявлення та усунення слабких місць або недоліків в елементах та процесах ланцюга постачання; SR-03a.[02] процес або процеси виявлення та усунення слабких місць або недоліків в елементах та процесах ланцюга постачання системи або компонента системи координується/координуються з персоналом ланцюга постачання; SR-03b. застосовуються засоби контролю ланцюга постачання для захисту від ризиків ланцюга постачання для системи, системного компонента або системної послуги, а також для обмеження шкоди або наслідків від подій, пов'язаних з ланцюгом постачання; SR-03c. задокументовані обрані та впроваджені процеси та засоби контролю ланцюга постачання в ЗНАЧЕННЯ ВИБІРКОВОГО ПАРАМЕТРА(ів)}




\subsubsection{РІЗНІ БАЗИ ПОСТАЧАННЯ (SR-3(1))}


\vspace{10pt}
\textit{Немає параметрів для цього контролю.}
\vspace{10pt}


\subsubsection{ОБМЕЖЕННЯ ШКОДИ (SR-3(2))}


\vspace{10pt}
{\footnotesize\bfseries
No: 1\\
Name: sr\_\allowbreak{}3\_\allowbreak{}2\_\allowbreak{}odp\\
Type: string\\
Default: \textquotedbl{}автоматизований засіб моніторингу\textquotedbl{}
}

{\footnotesize Визначені засоби контролю для обмеження шкоди від потенційних супротивників ланцюга постачання; SR-03(02) застосовуються контроль для обмеження шкоди від потенційних супротивників, які ідентифікують та націлюються на ланцюг постачання організації}




\subsubsection{ПЕРЕНЕСЕННЯ ЗАХОДІВ ЗАХИСТУ УПРАВЛІННЯ РИЗИКАМИ ЛАНЦЮГА ПОСТАЧАННЯ ДО СУБПІДРЯДНИКІВ (SR-3(3))}


\vspace{10pt}
{\footnotesize\bfseries
No: 1\\
Name: sr\_\allowbreak{}3\_\allowbreak{}3\_\allowbreak{}01\\
Type: string\\
Default: \textquotedbl{}автоматизований засіб моніторингу\textquotedbl{}
}

{\footnotesize Включені засоби контролю, передбачені в контрактах з основними підрядниками, також і в контрактах з субпідрядниками}




\subsection{ПОХОДЖЕННЯ (SR-4)}
Документуйте, відстежуйте та підтримуйте справжнє походження таких систем, компонентів системи і пов'язаних даних: [Призначення: системи, визначені організацією, системні компоненти та пов'язані дані].

\vspace{10pt}
{\footnotesize\bfseries
No: 1\\
Name: sr\_\allowbreak{}4\_\allowbreak{}01\\
Type: string\\
Default: nil
}

{\footnotesize Задокументовано дійсне походження для систем, компонентів системи та пов'язаних з ними даних}


{\footnotesize\bfseries
No: 2\\
Name: sr\_\allowbreak{}4\_\allowbreak{}02\\
Type: string\\
Default: nil
}

{\footnotesize Відстежується дійсне походження для систем, компонентів системи та пов'язаних з ними даних}


{\footnotesize\bfseries
No: 3\\
Name: sr\_\allowbreak{}4\_\allowbreak{}03\\
Type: string\\
Default: nil
}

{\footnotesize Підтримується дійсне походження для систем, компонентів системи та пов'язаних з ними даних}


{\footnotesize\bfseries
No: 4\\
Name: sr\_\allowbreak{}4\_\allowbreak{}odp\\
Type: string\\
Default: nil
}

{\footnotesize Визначені системи, компоненти системи та пов'язані з ними дані, які потребують достовірного походження}




\subsubsection{ІДЕНТИЧНІСТЬ (SR-4(1))}


\vspace{10pt}
{\footnotesize\bfseries
No: 1\\
Name: sr\_\allowbreak{}4\_\allowbreak{}1\_\allowbreak{}01\\
Type: list\\
Default: [\textquotedbl{}admin\textquotedbl{}, \textquotedbl{}security\_\allowbreak{}officer\textquotedbl{}]
}

{\footnotesize Встановлена унікальна ідентифікація елементів ланцюга постачання, процесів та персоналу; SR-04(01)[02] підтримується унікальна ідентифікація елементів ланцюга постачання, процесів та персоналу}


{\footnotesize\bfseries
No: 2\\
Name: sr\_\allowbreak{}4\_\allowbreak{}1\_\allowbreak{}02\\
Type: string\\
Default: nil
}

{\footnotesize Зберігається унікальна ідентифікація систем та критично важливих системних компонентів для відстеження в ланцюгу постачання}


{\footnotesize\bfseries
No: 3\\
Name: sr\_\allowbreak{}4\_\allowbreak{}1\_\allowbreak{}odp\\
Type: list\\
Default: [\textquotedbl{}admin\textquotedbl{}, \textquotedbl{}security\_\allowbreak{}officer\textquotedbl{}]
}

{\footnotesize Визначені елементи ланцюга постачання, процеси та персонал, пов'язані з системами та критично важливими компонентами системи, які потребують унікальної ідентифікації}




\subsubsection{УНІКАЛЬНА ІДЕНТИФІКАЦІЯ (SR-4(2))}


\vspace{10pt}
{\footnotesize\bfseries
No: 1\\
Name: sr\_\allowbreak{}4\_\allowbreak{}2\_\allowbreak{}01\\
Type: string\\
Default: nil
}

{\footnotesize Встановлена унікальна ідентифікація систем та критичних системних компонентів для відстеження в ланцюгу постачання}


{\footnotesize\bfseries
No: 2\\
Name: sr\_\allowbreak{}4\_\allowbreak{}2\_\allowbreak{}odp\\
Type: string\\
Default: nil
}

{\footnotesize Визначені системи та критичні компоненти системи, які потребують унікальної ідентифікації для відстеження в ланцюгу постачання}




\subsubsection{ПЕРЕВІРКА НА СПРАВЖНІСТЬ І ВІДСУТНІСТЬ ВНЕСЕННЯ ЗМІН (SR-4(3))}


\vspace{10pt}
{\footnotesize\bfseries
No: 1\\
Name: sr\_\allowbreak{}4\_\allowbreak{}3\_\allowbreak{}01\\
Type: string\\
Default: \textquotedbl{}автоматизований засіб моніторингу\textquotedbl{}
}

{\footnotesize Застосовуються засоби контролю для перевірки того, що отримана система або компонент системи є справжніми}


{\footnotesize\bfseries
No: 2\\
Name: sr\_\allowbreak{}4\_\allowbreak{}3\_\allowbreak{}02\\
Type: string\\
Default: \textquotedbl{}автоматизований засіб моніторингу\textquotedbl{}
}

{\footnotesize Застосовуються засоби контролю для перевірки того, що отриману систему або компонент системи не було змінено}




\subsubsection{ПОХОДЖЕННЯ -- ПЕРЕВІРКА ЛАНЦЮГА ЦІЛІСНОСТІ (SR-4(4))}


\vspace{10pt}
{\footnotesize\bfseries
No: 1\\
Name: sr\_\allowbreak{}4\_\allowbreak{}4\_\allowbreak{}01\\
Type: string\\
Default: \textquotedbl{}автоматизований засіб моніторингу\textquotedbl{}
}

{\footnotesize Застосовуються засоби контролю для забезпечення цілісності системи та її компонентів}


{\footnotesize\bfseries
No: 2\\
Name: sr\_\allowbreak{}4\_\allowbreak{}4\_\allowbreak{}02\\
Type: string\\
Default: nil
}

{\footnotesize Проводиться метод аналізу для забезпечення цілісності системи та компонентів системи}




\subsection{СТРАТЕГІЇ ПРИДБАННЯ, ІНСТРУМЕНТИ І МЕТОДИ (SR-5)}
Використовуйте наступні стратегії придбання, контрактні інструменти та методи закупівель, щоб захистити від ризиків ланцюга постачання, визначити та пом'якшити їх: [Призначення: визначені організацією стратегії придбання, контрактні інструменти та методи закупівель].

\vspace{10pt}
{\footnotesize\bfseries
No: 1\\
Name: sr\_\allowbreak{}5\_\allowbreak{}01\\
Type: string\\
Default: nil
}

{\footnotesize Застосовуються стратегії, інструменти та методи для захисту від ризиків ланцюга постачання}


{\footnotesize\bfseries
No: 2\\
Name: sr\_\allowbreak{}5\_\allowbreak{}02\\
Type: string\\
Default: nil
}

{\footnotesize Застосовуються стратегії, інструменти та методи для виявлення ризиків ланцюга постачання}


{\footnotesize\bfseries
No: 3\\
Name: sr\_\allowbreak{}5\_\allowbreak{}03\\
Type: string\\
Default: nil
}

{\footnotesize Застосовуються стратегії, інструменти та методи для зменшення ризиків ланцюга постачання}


{\footnotesize\bfseries
No: 4\\
Name: sr\_\allowbreak{}5\_\allowbreak{}odp\\
Type: string\\
Default: nil
}

{\footnotesize Визначені стратегії закупівель, контрактні інструменти та методи закупівель для захисту, виявлення та пом'якшення ризиків ланцюга постачання}




\subsubsection{НАЛЕЖНЕ ПОСТАЧАННЯ (SR-5(1))}


\vspace{10pt}
{\footnotesize\bfseries
No: 1\\
Name: sr\_\allowbreak{}5\_\allowbreak{}1\_\allowbreak{}01\\
Type: string\\
Default: \textquotedbl{}автоматизований засіб моніторингу\textquotedbl{}
}

{\footnotesize Застосовуються засоби контролю для забезпечення адекватного постачання критично важливих компонентів системи}




\subsubsection{ОЦІНКА ПЕРЕД ВІДБОРОМ, ПРИЙНЯТТЯ, МОДИФІКАЦІЯ ЧИ ОНОВЛЕННЯ (SR-5(2))}


\vspace{10pt}
{\footnotesize\bfseries
No: 1\\
Name: sr\_\allowbreak{}5\_\allowbreak{}2\_\allowbreak{}01\\
Type: string\\
Default: nil
}

{\footnotesize Оцінюється система, компонент системи або послуги системи перед відбором}


{\footnotesize\bfseries
No: 2\\
Name: sr\_\allowbreak{}5\_\allowbreak{}2\_\allowbreak{}02\\
Type: string\\
Default: nil
}

{\footnotesize Оцінюється система, компонент системи або послуги системи перед прийняттям}


{\footnotesize\bfseries
No: 3\\
Name: sr\_\allowbreak{}5\_\allowbreak{}2\_\allowbreak{}03\\
Type: string\\
Default: nil
}

{\footnotesize Оцінюється система, компонент системи або послуги системи перед модифікацією}


{\footnotesize\bfseries
No: 4\\
Name: sr\_\allowbreak{}5\_\allowbreak{}2\_\allowbreak{}04\\
Type: string\\
Default: nil
}

{\footnotesize Оцінюється система, компонент системи або послуги системи перед оновленням}




\subsection{ОЦІНКА ПОСТАЧАЛЬНИКІВ (SR-6)}
Оцініть і перегляньте ризики ланцюга постачання, пов'язані з постачальниками або підрядниками, системою, системним компонентом або системною послугою, яку вони надають [Призначення: частота, визначена організацією].

\vspace{10pt}
{\footnotesize\bfseries
No: 1\\
Name: sr\_\allowbreak{}6\_\allowbreak{}01\\
Type: string\\
Default: \textquotedbl{}щорічно\textquotedbl{}
}

{\footnotesize Оцінюються та аналізуються ризики, пов'язані з ланцюгом постачання, які стосуються постачальників або підрядників та систем, компонентів системи або системних послуг, які вони надають < SR-06\_\allowbreak{}ODP частота >}


{\footnotesize\bfseries
No: 2\\
Name: sr\_\allowbreak{}6\_\allowbreak{}odp\\
Type: string\\
Default: \textquotedbl{}щорічно\textquotedbl{}
}

{\footnotesize Визначена періодичність оцінки та аналізу ризиків, пов'язаних з ланцюгом постачання, що стосуються постачальників або підрядників, а також систем, компонентів системи або системних послуг, які вони надають}




\subsubsection{ТЕСТУВАННЯ ТА АНАЛІЗ (SR-6(1))}


\vspace{10pt}
\textit{Немає параметрів для цього контролю.}
\vspace{10pt}


\subsection{БЕЗПЕКА ОПЕРАЦІЙ ЛАНЦЮГА ПОСТАЧАННЯ (SR-7)}
Використовуйте такі заходи захисту операційної безпеки (OPSEC), щоб захистити інформацію, пов'язану з ланцюгом постачання для системи, системного компонента чи системної служби: [Призначення: визначені організацією заходи захисту операційної безпеки (OPSEC)].

\vspace{10pt}
{\footnotesize\bfseries
No: 1\\
Name: sr\_\allowbreak{}7\_\allowbreak{}01\\
Type: string\\
Default: \textquotedbl{}автоматизований засіб моніторингу\textquotedbl{}
}

{\footnotesize Застосовуються засоби управління OPSEC для захисту інформації, пов'язаної з ланцюжком постачання для системи, системного компонента або системної служби}


{\footnotesize\bfseries
No: 2\\
Name: sr\_\allowbreak{}7\_\allowbreak{}odp\\
Type: string\\
Default: nil
}

{\footnotesize Визначені заходи захисту операційної безпеки (OPSEC) для захисту інформації, пов'язаної з ланцюжком поставок, для системи, системного компонента або системної служби}




\subsection{ПОВІДОМЛЕННЯ ПРО ПОРУШЕННЯ ЛАНЦЮГА ПОСТАЧАННЯ (SR-8)}
Затвердити угоди та процедури з суб'єктами, залученими до ланцюга постачання для системи, системного компонента або системної послуги для [Вибір (одного або кількох): повідомлення про порушення ланцюга постачання; результати оцінювання або аудитів; [Призначення: інформація, визначена організацією]].

\vspace{10pt}
{\footnotesize\bfseries
No: 1\\
Name: sr\_\allowbreak{}8\_\allowbreak{}01\\
Type: string\\
Default: nil
}

{\footnotesize Аналізу/тестування елементів, ПОВІДОМЛЕННЯ ПРО ПОРУШЕННЯ ЛАНЦЮГА ПОСТАЧАННЯ}


{\footnotesize\bfseries
No: 2\\
Name: sr\_\allowbreak{}8\_\allowbreak{}odp\_\allowbreak{}01\\
Type: string\\
Default: nil
}

{\footnotesize Вибрано одне або декілька з наступних ЗНАЧЕНЬ ПАРАМЕТРІВ: \{повідомлення про порушення ланцюга постачання; результати оцінок або аудитів\}}


{\footnotesize\bfseries
No: 3\\
Name: sr\_\allowbreak{}8\_\allowbreak{}odp\_\allowbreak{}02\\
Type: string\\
Default: nil
}

{\footnotesize Визначена інформація, для якої необхідно встановити угоди та процедури (якщо вибрано)}




\subsection{ЗАХИСТ ВІД ЗЛОМУ ТА ВИЯВЛЕННЯ (SR-9)}
Впровадити програму захисту від несанкціонованого доступу для системи, системного компонента або системної служби.

\vspace{10pt}
{\footnotesize\bfseries
No: 1\\
Name: sr\_\allowbreak{}9\_\allowbreak{}01\\
Type: string\\
Default: nil
}

{\footnotesize Реалізована програма захисту від несанкціонованого доступу для системи, компонента системи або системної служби}




\subsubsection{ЕТАПИ ЧИ СИСТЕМИ РОЗВИТКУ ЖИТТЄВОГО ЦИКЛУ (SR-9(1))}


\vspace{10pt}
{\footnotesize\bfseries
No: 1\\
Name: sr\_\allowbreak{}9\_\allowbreak{}1\_\allowbreak{}01\\
Type: string\\
Default: nil
}

{\footnotesize Застосовуються технології, інструменти та методи захисту від втручання протягом усього життєвого циклу розробки системи}




\subsection{ПЕРЕВІРКА СИСТЕМИ І КОМПОНЕНТІВ СИСТЕМИ (SR-10)}
Перевірте наступні системи або системні компоненти [Вибір (один або більше): випадковим чином обраних; на [Призначення: частота, визначена організацією], після [Призначення: визначені організацією ознаки необхідності перевірки]] для виявлення втручання: [Призначення: визначені організацією системи або компоненти системи].

\vspace{10pt}
{\footnotesize\bfseries
No: 1\\
Name: sr\_\allowbreak{}10\_\allowbreak{}odp\_\allowbreak{}01\\
Type: string\\
Default: nil
}

{\footnotesize Визначені системи або компоненти системи, які потребують перевірки}


{\footnotesize\bfseries
No: 2\\
Name: sr\_\allowbreak{}10\_\allowbreak{}odp\_\allowbreak{}02\\
Type: integer\\
Default: 30
}

{\footnotesize Вибрано одне або декілька з наступних ЗНАЧЕНЬ ПАРАМЕТРА: \{випадково; з частотою <SR-10\_\allowbreak{}ODP[03]; за наявності вказівок на необхідність перевірки\}}


{\footnotesize\bfseries
No: 3\\
Name: sr\_\allowbreak{}10\_\allowbreak{}odp\_\allowbreak{}03\\
Type: string\\
Default: \textquotedbl{}щорічно\textquotedbl{}
}

{\footnotesize Визначена періодичність проведення перевірок систем або компонентів системи (якщо вибрано)}


{\footnotesize\bfseries
No: 4\\
Name: sr\_\allowbreak{}10\_\allowbreak{}odp\_\allowbreak{}04\\
Type: string\\
Default: nil
}

{\footnotesize Визначені ознаки необхідності перевірки систем або компонентів системи (якщо вони були обрані)}




\subsection{АВТЕНТИЧНІСТЬ КОМПОНЕНТУ (SR-11)}
a. Розробити та впровадити політику та процедури боротьби з підробками, які включають засоби для виявлення та запобігання потраплянню підроблених компонентів у систему;\\ 
b. Повідомляти про підроблені системні компоненти [Вибір (один або кілька): джерело підробленого компонента; [Призначення: зовнішні звітні організації, визначені організацією]; [Призначення: персонал або ролі, визначені організацією]].

\vspace{10pt}
{\footnotesize\bfseries
No: 1\\
Name: sr\_\allowbreak{}11\_\allowbreak{}odp\_\allowbreak{}01\\
Type: list\\
Default: [\textquotedbl{}admin\textquotedbl{}, \textquotedbl{}security\_\allowbreak{}officer\textquotedbl{}]
}

{\footnotesize Вибрано одне або декілька з наступних ЗНАЧЕНЬ ПАРАМЕТРА: \{джерело підробленого компонента; зовнішні підзвітні організації; персонал або ролі\}}


{\footnotesize\bfseries
No: 2\\
Name: sr\_\allowbreak{}11\_\allowbreak{}odp\_\allowbreak{}02\\
Type: string\\
Default: nil
}

{\footnotesize Визначені зовнішні підзвітні організації, яким слід повідомляти про підроблені компоненти системи (якщо вони були обрані)}




\subsubsection{АВТЕНТИЧНІСТЬ КОМПОНЕНТУ (SR-11(1))}


\vspace{10pt}
{\footnotesize\bfseries
No: 1\\
Name: sr\_\allowbreak{}11\_\allowbreak{}1\_\allowbreak{}01\\
Type: list\\
Default: [\textquotedbl{}admin\textquotedbl{}, \textquotedbl{}security\_\allowbreak{}officer\textquotedbl{}]
}

{\footnotesize Підготовлений персонал або ролі до виявлення підроблених компонентів системи, включаючи апаратне, програмне та мікропрограмне забезпечення}


{\footnotesize\bfseries
No: 2\\
Name: sr\_\allowbreak{}11\_\allowbreak{}1\_\allowbreak{}odp\\
Type: list\\
Default: [\textquotedbl{}admin\textquotedbl{}, \textquotedbl{}security\_\allowbreak{}officer\textquotedbl{}]
}

{\footnotesize Визначено персонал або ролі, які потребують підготовки для виявлення підроблених компонентів системи (включаючи апаратне, програмне та мікропрограмне забезпечення)}




\subsection{УТИЛІЗАЦІЯ КОМПОНЕНТУ (SR-12)}
Утилізуйте [Призначення: визначені організацією дані, документація, інструменти або системні компоненти] за допомогою таких прийомів і методів: [Призначення: визначені організацією прийоми та методи].

\vspace{10pt}
{\footnotesize\bfseries
No: 1\\
Name: sr\_\allowbreak{}12\_\allowbreak{}01\\
Type: string\\
Default: nil
}

{\footnotesize Утилізуються дані, документація, інструменти або компоненти системи з використанням прийомів і методів}


{\footnotesize\bfseries
No: 2\\
Name: sr\_\allowbreak{}12\_\allowbreak{}odp\_\allowbreak{}01\\
Type: string\\
Default: nil
}

{\footnotesize Визначені дані, документація, інструменти або компоненти системи, які підлягають утилізації}


{\footnotesize\bfseries
No: 3\\
Name: sr\_\allowbreak{}12\_\allowbreak{}odp\_\allowbreak{}02\\
Type: list\\
Default: [\textquotedbl{}admin\textquotedbl{}, \textquotedbl{}security\_\allowbreak{}officer\textquotedbl{}]
}

{\footnotesize Визначені методи та способи утилізації даних, документації, інструментів або компонентів системи}




\section{PM}
Клас заходів захисту PM --- МЕНЕДЖМЕНТ ІНФОРМАЦІЙНОЇ БЕЗПЕКИ

Цей клас визначає загальноорганізаційні заходи для ефективного управління програмами інформаційної безпеки та приватності на рівні всього підприємства.

\textbf{Перелік заходів захисту:} Ролі програми інформаційної безпеки (PM-2); Ресурси забезпечення інформаційної безпеки та приватності (PM-3); Інвентаризація системи (PM-5); Архітектура підприємства (PM-7); Розвантаження (PM-7(1)); План захисту критичної інфраструктури (PM-8); Стратегія управління ризиками (PM-9); Процес авторизації (PM-10); Визначення завдань та процесів (PM-11); Програма інсайдерської загрози (PM-12); Безпека та приватність працівників (PM-13); Тестування, навчання та моніторинг (PM-14); Контакти з групами та асоціаціями з питань безпеки інформації та приватності (PM-15); Програма інформування про загрози (PM-16); Програма інформування про загрози (PM-16(1)); Захист публічної інформації у зовнішніх системах (PM-17); ПРОГРАМА (КОНЦЕПЦІЯ) ЗАБЕЗПЕЧЕННЯ ПРИВАТНОСТІ (PM-18); Керівні ролі програми приватності (PM-19); Система записів програми приватності (PM-20); Облік розкриття персональних даних (PM-21); Управління якістю персональних даних (PM-22); Орган управління персональними даними (PM-23); Орган з питань цілісності даних (PM-24); Мінімізація кількості персональних даних, що використовуються під час тестування, навчання та досліджень (PM-25); Управління скаргами (PM-26); Звітність з питань забезпечення приватності (PM-27); Оцінка ризиків (PM-28); Ролі керівників програми управління ризиками (PM-29); План управління ризиками ланцюга постачання (PM-30); План безперервного моніторингу (PM-31); Призначення (PM-32).

\subsection{Ролі програми інформаційної безпеки (PM-2)}
Призначити старшу посадову особу служби інформаційної безпеки, яка наділена відповідними завданнями та ресурсами для здійснення координації, розробки, впровадження та підтримки програми (концепції) інформаційної безпеки.

\vspace{10pt}
{\footnotesize\bfseries
No: 1\\
Name: pm\_\allowbreak{}2\_\allowbreak{}01\\
Type: list\\
Default: [\textquotedbl{}admin\textquotedbl{}, \textquotedbl{}security\_\allowbreak{}officer\textquotedbl{}]
}

{\footnotesize Призначено старшу посадову особу з питань інформаційної безпеки в установі}


{\footnotesize\bfseries
No: 2\\
Name: pm\_\allowbreak{}2\_\allowbreak{}02\\
Type: list\\
Default: [\textquotedbl{}admin\textquotedbl{}, \textquotedbl{}security\_\allowbreak{}officer\textquotedbl{}]
}

{\footnotesize Надано посадовій особі з інформаційної безпеки відомства повноваження та ресурси для координації загальноорганізаційної програми (концепції) з інформаційної безпеки}


{\footnotesize\bfseries
No: 3\\
Name: pm\_\allowbreak{}2\_\allowbreak{}03\\
Type: list\\
Default: [\textquotedbl{}admin\textquotedbl{}, \textquotedbl{}security\_\allowbreak{}officer\textquotedbl{}]
}

{\footnotesize Має старша посадова особа з питань інформаційної безпеки відомства місію та ресурси для розробки загальноорганізаційної програми інформаційної безпеки}


{\footnotesize\bfseries
No: 4\\
Name: pm\_\allowbreak{}2\_\allowbreak{}04\\
Type: list\\
Default: [\textquotedbl{}admin\textquotedbl{}, \textquotedbl{}security\_\allowbreak{}officer\textquotedbl{}]
}

{\footnotesize Забезпечено старшу посадову особу з інформаційної безпеки відомства необхідним та ресурсами для впровадження загальноорганізаційної програми інформаційної безпеки}


{\footnotesize\bfseries
No: 5\\
Name: pm\_\allowbreak{}2\_\allowbreak{}05\\
Type: list\\
Default: [\textquotedbl{}admin\textquotedbl{}, \textquotedbl{}security\_\allowbreak{}officer\textquotedbl{}]
}

{\footnotesize Забезпечено старшу посадову особу з інформаційної безпеки відомства необхідним та ресурсами для підтримки загальноорганізаційної програми (концепції) з інформаційної безпеки}




\subsection{Ресурси забезпечення інформаційної безпеки та приватності (PM-3)}
a.\\ 
b. Запровадити процес для забезпечення того, щоб плани дій та етапи програм безпеки та приватності, програм управління ризиками ланцюга постачання і пов'язаних систем організації:\\ 
1. розроблялися та підтримувалися;\\ 
2. задокументовані корегувальні заходи захисту адекватно реагували на ризики для операцій організацї і активів, фізичних осіб, інших організацій та держави;\\ 
3. оприлюднювалися відповідно до встановлених вимог до звітності. Переглядати плани дій та етапи для узгодженості з організаційною стратегією управління ризиками й організаційними пріоритетами щодо дій з реагування на ризики.

\vspace{10pt}
\textit{Немає параметрів для цього контролю.}
\vspace{10pt}


\subsection{Інвентаризація системи (PM-5)}
Розробити, відстежувати та звітувати про результати вимірювань показників продуктивності забезпечення безпеки інформації та приватності.

\vspace{10pt}
{\footnotesize\bfseries
No: 1\\
Name: pm\_\allowbreak{}5\_\allowbreak{}01\\
Type: list\\
Default: []
}

{\footnotesize Розроблено перелік систем організацї}


{\footnotesize\bfseries
No: 2\\
Name: pm\_\allowbreak{}5\_\allowbreak{}02\\
Type: list\\
Default: []
}

{\footnotesize Оновлюється frequency>. перелік оновлення систем переліку організації систем <PM-05\_\allowbreak{}ODP}


{\footnotesize\bfseries
No: 3\\
Name: pm\_\allowbreak{}5\_\allowbreak{}odp\\
Type: string\\
Default: \textquotedbl{}щорічно\textquotedbl{}
}

{\footnotesize Визначена періодичність організацї}




\subsection{Архітектура підприємства (PM-7)}
Визначити завдання інформаційної безпеки та приватності при розробці документуванні та оновленні плану захисту критичної інфраструктури та ключових ресурсів.

\vspace{10pt}
{\footnotesize\bfseries
No: 1\\
Name: pm\_\allowbreak{}7\_\allowbreak{}01\\
Type: string\\
Default: nil
}

{\footnotesize Розроблена архітектура інформаційної безпеки; підприємства з урахуванням}


{\footnotesize\bfseries
No: 2\\
Name: pm\_\allowbreak{}7\_\allowbreak{}02\\
Type: string\\
Default: nil
}

{\footnotesize Підтримується архітектура інформаційної безпеки; підприємства з урахуванням}


{\footnotesize\bfseries
No: 3\\
Name: pm\_\allowbreak{}7\_\allowbreak{}03\\
Type: string\\
Default: nil
}

{\footnotesize Розроблена архітектура конфіденційності; підприємства з урахуванням}


{\footnotesize\bfseries
No: 4\\
Name: pm\_\allowbreak{}7\_\allowbreak{}04\\
Type: string\\
Default: nil
}

{\footnotesize Підтримується архітектура конфіденційності; підприємства з урахуванням}


{\footnotesize\bfseries
No: 5\\
Name: pm\_\allowbreak{}7\_\allowbreak{}05\\
Type: string\\
Default: nil
}

{\footnotesize Розроблена архітектура підприємства з урахуванням ризиків для діяльності та активів організації, окремих осіб, інших організацій та держави в цілому}


{\footnotesize\bfseries
No: 6\\
Name: pm\_\allowbreak{}7\_\allowbreak{}06\\
Type: string\\
Default: nil
}

{\footnotesize Підтримується архітектура підприємства з урахуванням ризиків, що виникають в результаті цього для операцій та активів організації, окремих осіб, інших організацій та держави.,}




\subsubsection{РОЗВАНТАЖЕННЯ (PM-7(1))}


\vspace{10pt}
{\footnotesize\bfseries
No: 1\\
Name: pm\_\allowbreak{}7\_\allowbreak{}1\_\allowbreak{}01\\
Type: string\\
Default: nil
}

{\footnotesize Несуттєві функції або послуги вивантажуються на інші системи, компоненти системи або зовнішнього постачальника}


{\footnotesize\bfseries
No: 2\\
Name: pm\_\allowbreak{}7\_\allowbreak{}1\_\allowbreak{}odp\\
Type: string\\
Default: nil
}

{\footnotesize Визначені несуттєві функції або послуги, які потрібно розвантажити}




\subsection{ПЛАН ЗАХИСТУ КРИТИЧНОЇ ІНФРАСТРУКТУРИ (PM-8)}


\vspace{10pt}
{\footnotesize\bfseries
No: 1\\
Name: pm\_\allowbreak{}8\_\allowbreak{}01\\
Type: string\\
Default: nil
}

{\footnotesize Враховані питання інформаційної безпеки при розробці плану захисту критичної інфраструктури та ключових ресурсів; PM-08[02] розглядаються питання інформаційної безпеки в документації плану захисту критичної інфраструктури та ключових ресурсів}


{\footnotesize\bfseries
No: 2\\
Name: pm\_\allowbreak{}8\_\allowbreak{}03\\
Type: string\\
Default: nil
}

{\footnotesize Враховані питання інформаційної безпеки в оновленому плані захисту критичної інфраструктури та ключових ресурсів}


{\footnotesize\bfseries
No: 3\\
Name: pm\_\allowbreak{}8\_\allowbreak{}04\\
Type: string\\
Default: nil
}

{\footnotesize Враховані питання конфіденційності при розробці плану захисту критичної інфраструктури та ключових ресурсів}


{\footnotesize\bfseries
No: 4\\
Name: pm\_\allowbreak{}8\_\allowbreak{}05\\
Type: string\\
Default: nil
}

{\footnotesize Розглядаються питання конфіденційності в документації плану захисту критичної інфраструктури та ключових ресурсів}


{\footnotesize\bfseries
No: 5\\
Name: pm\_\allowbreak{}8\_\allowbreak{}06\\
Type: string\\
Default: nil
}

{\footnotesize Враховані питання конфіденційності в оновленому плані захисту критичної інфраструктури та ключових ресурсів}




\subsection{СТРАТЕГІЯ УПРАВЛІННЯ РИЗИКАМИ (PM-9)}
a. Розробити комплексну стратегію управління:\\ 
1. ризиками безпеки для операцій та активів організації, фізичних осіб, інших організацій і держави, пов'язаних з експлуатацією та використанням систем організації;\\ 
2. ризиками приватності для фізичних осіб, які можуть виникати внаслідок збирання, обміну, зберігання, передачі, використання та розпорядження персональними даними;\\ 
b. Реалізувати стратегію управління ризиками в масштабах організації.\\ 
c. Переглядати й оновлювати стратегію управління ризиками [Призначення: з визначеною організацією частотою] або, якщо потрібно, у разі змін в організації.

\vspace{10pt}
{\footnotesize\bfseries
No: 1\\
Name: pm\_\allowbreak{}9\_\allowbreak{}odp\\
Type: list\\
Default: [\textquotedbl{}admin\textquotedbl{}, \textquotedbl{}security\_\allowbreak{}officer\textquotedbl{}]
}

{\footnotesize Визначено періодичність перегляду та оновлення стратегії управління ризиками; PM-09a.01 розроблена комплексна стратегія управління ризиками безпеки для операцій та активів організації, окремих осіб, інших організацій та держави, пов'язаних з експлуатацією та використанням організаційних систем; PM-09a.02 розроблена комплексна стратегія управління ризиками для приватності осіб, що виникають внаслідок санкціонованої обробки інформації, що ідентифікує особу; PM-09b. стратегія управління ризиками послідовно впроваджується в організації; PM-09c. переглядається та оновлюється стратегія управління ризиками частота або в міру необхідності у зв'язку з організаційними змінами}




\subsection{ПРОЦЕС АВТОРИЗАЦІЇ (PM-10)}
a. Управляти станом безпеки та приватності інформаційних систем організації та середовищ, у яких ці інформаційні системи експлуатуються через процедури авторизації\\ 
b. Призначити окремих осіб для виконання певних ролей і обов'язків у рамках організаційного процесу управління ризиками.\\ 
c. Інтегрувати процеси авторизації в загальноорганізаційну програму управління ризиками.

\vspace{10pt}
\textit{Немає параметрів для цього контролю.}
\vspace{10pt}


\subsection{ВИЗНАЧЕННЯ ЗАВДАНЬ ТА ПРОЦЕСІВ (PM-11)}
Створити програму розвитку та вдосконалення спеціалістів з питань безпеки та приватності.

\vspace{10pt}
{\footnotesize\bfseries
No: 1\\
Name: pm\_\allowbreak{}11\_\allowbreak{}odp\\
Type: list\\
Default: [\textquotedbl{}admin\textquotedbl{}, \textquotedbl{}security\_\allowbreak{}officer\textquotedbl{}]
}

{\footnotesize Визначено періодичність перегляду завдань та бізнеспроцесів; PM-11a.[01] завдання та бізнес-процеси організації визначені з урахуванням інформаційної безпеки; PM-11a.[02] завдання та бізнес-процеси організації визначені з урахуванням права на приватність; PM-11a.[03] завдання та бізнес-процеси організації визначені з урахуванням ризиків для діяльності організації, її активів, окремих осіб, інших організацій та держави в цілому; PM-11b.[01] визначені потреби в захисті інформації, що випливають з визначених завдань та бізнес-процесів; PM-11b.[02] визначені потреби в обробці персональних випливають з визначеної місії та бізнес-процесів; PM-11c. переглядаються завдання та бізнес-процеси частота. даних, що}




\subsection{ПРОГРАМА ІНСАЙДЕРСЬКОЇ ЗАГРОЗИ (PM-12)}


\vspace{10pt}
{\footnotesize\bfseries
No: 1\\
Name: pm\_\allowbreak{}12\_\allowbreak{}01\\
Type: string\\
Default: nil
}

{\footnotesize Впроваджено програму інсайдерської (внутрішньої) загрози, яка передбачає наявність команди з обробки інцидентів, пов'язаних з внутрішньою дисципліною}




\subsection{БЕЗПЕКА ТА ПРИВАТНІСТЬ ПРАЦІВНИКІВ (PM-13)}


\vspace{10pt}
{\footnotesize\bfseries
No: 1\\
Name: pm\_\allowbreak{}13\_\allowbreak{}01\\
Type: string\\
Default: nil
}

{\footnotesize Існує програма розвитку та вдосконалення спеціалістів з питань безпеки}


{\footnotesize\bfseries
No: 2\\
Name: pm\_\allowbreak{}13\_\allowbreak{}02\\
Type: string\\
Default: nil
}

{\footnotesize Створена програма розвитку та вдосконалення спеціалістів з питань приватності}




\subsection{ТЕСТУВАННЯ, НАВЧАННЯ ТА МОНІТОРИНГ (PM-14)}
Запровадити програму інформування про загрози, яка містить можливості спільного обміну інформацією між організаціями для аналізу загроз.

\vspace{10pt}
\textit{Немає параметрів для цього контролю.}
\vspace{10pt}


\subsection{КОНТАКТИ З ГРУПАМИ ТА АСОЦІАЦІЯМИ З ПИТАНЬ БЕЗПЕКИ ІНФОРМАЦІЇ ТА ПРИВАТНОСТІ (PM-15)}


\vspace{10pt}
\textit{Немає параметрів для цього контролю.}
\vspace{10pt}


\subsection{ПРОГРАМА ІНФОРМУВАННЯ ПРО ЗАГРОЗИ (PM-16)}


\vspace{10pt}
{\footnotesize\bfseries
No: 1\\
Name: pm\_\allowbreak{}16\_\allowbreak{}01\\
Type: string\\
Default: nil
}

{\footnotesize Впроваджена програма інформування про загрози, яка передбачає можливість обміну інформацією між організаціями для розвідки загроз}




\subsubsection{ПРОГРАМА ІНФОРМУВАННЯ ПРО ЗАГРОЗИ (PM-16(1))}


\vspace{10pt}
{\footnotesize\bfseries
No: 1\\
Name: pm\_\allowbreak{}16\_\allowbreak{}1\_\allowbreak{}01\\
Type: string\\
Default: nil
}

{\footnotesize Automated mechanisms are employed to maximize the effectiveness of sharing threat intelligence information}




\subsection{ЗАХИСТ ПУБЛІЧНОЇ ІНФОРМАЦІЇ У ЗОВНІШНІХ СИСТЕМАХ (PM-17)}
a. Розробити політику та процедури для забезпечення того, щоб вимоги до захисту публічної (некласифікованої) інформації, яка обробляється, зберігається або передається у зовнішніх системах, здійснювалися відповідно до чинного законодавства.\\ 
b. Оновлювати політику та процедури [Призначення: з визначеною організацією частотою].

\vspace{10pt}
{\footnotesize\bfseries
No: 1\\
Name: pm\_\allowbreak{}17\_\allowbreak{}odp\_\allowbreak{}01\\
Type: list\\
Default: [\textquotedbl{}default\_\allowbreak{}deny\_\allowbreak{}rule\textquotedbl{}, \textquotedbl{}abac\_\allowbreak{}rule\_\allowbreak{}1\textquotedbl{}]
}

{\footnotesize Визначена періодичність перегляду та оновлення політики}


{\footnotesize\bfseries
No: 2\\
Name: pm\_\allowbreak{}17\_\allowbreak{}odp\_\allowbreak{}02\\
Type: list\\
Default: [\textquotedbl{}default\_\allowbreak{}deny\_\allowbreak{}rule\textquotedbl{}, \textquotedbl{}abac\_\allowbreak{}rule\_\allowbreak{}1\textquotedbl{}]
}

{\footnotesize Визначено періодичність перегляду та оновлення процедур; PM-17a.[01] розроблено політику, яка гарантує, що вимоги щодо захисту публічної (некласифікованої) інформації, яка обробляється, зберігається або передається в зовнішніх системах, виконуються відповідно до чинних законів, виконавчих наказів, директив, політик, нормативних актів та стандартів; PM-17a.[02] встановлені процедури для забезпечення виконання вимог щодо захисту публічної (некласифікованої) інформації, яка обробляється, зберігається або передається в зовнішніх системах, відповідно до чинних законів, наказів, директив, політик, нормативно-правових актів та стандартів; PM-17b.[01] переглядається та оновлюється політика частота; PM-17b.[02] переглядаються та оновлюються процедури частота}




\subsection{ПРОГРАМА (КОНЦЕПЦІЯ) ЗАБЕЗПЕЧЕННЯ ПРИВАТНОСТІ (PM-18)}
a.\\ 
b. Розробити та поширити загальноорганізаційну програму (концепцію) забезпечення приватності, яка:\\ 
1. містить опис структури програми забезпечення приватності та ресурсів, призначених для її реалізації;\\ 
2. містить огляд вимог до забезпечення приватності й опис засобів управління програмою забезпечення приватності та загальних заходів захисту, встановлених або запланованих для задоволення цих вимог;\\ 
3. визначає обов'язки посадової особи щодо приватності, а також визначає обов'язки інших посадових осіб і персоналу з питань забезпечення приватності;\\ 
4. описує зобов'язання керівництва, стратегічні цілі та завдання програми забезпечення приватності;\\ 
5. відображає координацію між організаційними структурами, відповідальними за різні аспекти приватності;\\ 
6. затверджена високопосадовцем, який є відповідальним (та підзвітним) за: управління ризиками приватності, що виникають при здійсненні операцій організації (включно із завданнями, функціями, іміджем і репутацією); організаційними активами, фізичними особами, іншими організаціями та країнами. Оновлювати програму [Призначення: за визначеною організацією частотою], а також в разі змін законодавства, змін в організації і виявлення проблем в ході реалізації програми або оцінювання заходів приватності.

\vspace{10pt}
{\footnotesize\bfseries
No: 1\\
Name: pm\_\allowbreak{}18\_\allowbreak{}odp\\
Type: list\\
Default: [\textquotedbl{}admin\textquotedbl{}, \textquotedbl{}security\_\allowbreak{}officer\textquotedbl{}]
}

{\footnotesize Визначена періодичність (концепції) приватності; оновлення плану програми PM-18a.[01] розроблено загальноорганізаційний план програми (концепції) приватності, який містить огляд програми приватності організації; PM-18a.01[01] план програми (концепції) приватності містить опис структури програми конфіденційності; PM-18a.01[02] план програми (концепції) приватності містить опис ресурсів, призначених для реалізації програми конфіденційності; PM-18a.02[01] план програми (концепції) приватності містить огляд вимог до програми конфіденційності; PM-18a.02[02] план програми (концепції) приватності містить опис наявних або запланованих засобів контролю для управління програмою (концепцією) приватності для виконання вимог програми; PM-18a.02[03] план програми (концепції) приватності містить опис загальних засобів контролю, що діють або заплановані для виконання вимог програми конфіденційності; PM-18a.03[01] в плані програми (концепції) приватності передбачена роль старшої посадової особи організації з питань приватності; PM-18a.03[02] план програми (концепції) приватності включає визначення та призначення ролей інших посадових осіб і співробітників, відповідальних за забезпечення приватності, та їхні обов'язки; PM-18a.04[01] план програми (концепції) приватності описує зобов'язання керівництва; PM-18a.04[02] в плані програми (концепції) приватності описано дотримання вимог; PM-18a.04[03] план програми (концепція) приватності описує стратегічні цілі та завдання програми приватності; PM-18a.05 план програми (концепція) приватності відображає координацію між підрозділами організації, відповідальними за різні аспекти приватності; PM-18a.06 затверджено план програми (концепцію) приватності вищою посадовою особою, яка несе відповідальність і підзвітність за ризики для приватності, яких зазнають операції організації (включно з місією, функціями, іміджем і репутацією), активи організації, окремі особи, інші організації та держава; PM-18a.[02] поширюється план програми (концепція) приватності; PM-18b.[01] оновлено план програми (концепцію) приватності частота; PM-18b.[02] оновлюється план програми (концепція) приватності відповідно до змін у державному законодавстві та політиці щодо приватності; PM-18b.[03] оновлюється план програми відповідно до змін в організації; PM-18b.[04] оновлюється план програми (концепція) приватності забезпечення приватності для вирішення проблем, виявлених під час реалізації плану або оцінок контролю за дотриманням приватності. (концепція) приватності}




\subsection{КЕРІВНІ РОЛІ ПРОГРАМИ ПРИВАТНОСТІ (PM-19)}
Призначити старшу посадову особу з питань забезпечення приватності з повноваженнями, завданням, підзвітністю і ресурсами для координації, розробки та реалізації відповідних вимог забезпечення приватності й управління ризиками приватності в рамках програми забезпечення приватності всієї організації.

\vspace{10pt}
{\footnotesize\bfseries
No: 1\\
Name: pm\_\allowbreak{}19\_\allowbreak{}01\\
Type: string\\
Default: \textquotedbl{}щорічно\textquotedbl{}
}

{\footnotesize Визначена періодичність приватності; оновлення плану програми}


{\footnotesize\bfseries
No: 2\\
Name: pm\_\allowbreak{}19\_\allowbreak{}02\\
Type: string\\
Default: nil
}

{\footnotesize Розроблено загальноорганізаційний план програми приватності, який містить огляд програми приватності для організації}


{\footnotesize\bfseries
No: 3\\
Name: pm\_\allowbreak{}19\_\allowbreak{}03\\
Type: string\\
Default: nil
}

{\footnotesize Містить план програми приватності опис структури програми приватності}


{\footnotesize\bfseries
No: 4\\
Name: pm\_\allowbreak{}19\_\allowbreak{}04\\
Type: string\\
Default: nil
}

{\footnotesize Містить план програми приватності опис ресурсів, призначених для реалізації програми приватності}


{\footnotesize\bfseries
No: 5\\
Name: pm\_\allowbreak{}19\_\allowbreak{}05\\
Type: string\\
Default: nil
}

{\footnotesize Містить план програми приватності огляд вимог до програми приватності}




\subsection{СИСТЕМА ЗАПИСІВ ПРОГРАМИ ПРИВАТНОСТІ (PM-20)}
підтримувати центральну вебсторінку ресурсу на головному загальнодоступному вебсайті організації, яка слугує центральним джерелом інформації про програму приватності організації та яка:\\ 
a. забезпечує доступ громадськості до інформації про діяльність щодо забезпечення приватності в організації та можливість комунікації з уповноваженою посадовою особою з питань забезпечення приватності;\\ 
b. оприлюднює організаційну політику забезпечення приватності на вебсайті організації або іншим чином;\\ 
c. використовує публічні адреси електронної пошти та/або телефонні лінії, щоб дати можливість громадськості надавати відгуки та/або направляти запитання щодо програми приватності в організації.

\vspace{10pt}
{\footnotesize\bfseries
No: 1\\
Name: pm\_\allowbreak{}20\_\allowbreak{}01\\
Type: string\\
Default: nil
}

{\footnotesize Ведеться центральна вебсторінка загальнодоступному вебсайті організації; на головному}


{\footnotesize\bfseries
No: 2\\
Name: pm\_\allowbreak{}20\_\allowbreak{}02\\
Type: list\\
Default: [\textquotedbl{}admin\textquotedbl{}, \textquotedbl{}security\_\allowbreak{}officer\textquotedbl{}]
}

{\footnotesize Слугує вебсторінка основним джерелом програму приватності організації; інформації про PM-20a.[01] забезпечує вебсторінка доступ громадськості до інформації про діяльність організації, пов'язану із захистом приватності; PM-20a.[02] забезпечує вебсторінка можливість громадськості спілкуватися з вищим посадовцем організації з питань приватності; PM-20b.[01] забезпечує вебсторінка публічний доступ до інформації організації щодо приватності; PM-20b.[02] забезпечує вебсторінка публічний доступ до звітів про приватність організації; PM-20c. є на веб-сторінці загальнодоступні адреси електронної пошти та/або номери телефонів, щоб громадськість могла надавати зворотній зв'язок та/або направляти запитання до відділів з питань приватності}




\subsection{ОБЛІК РОЗКРИТТЯ ПЕРСОНАЛЬНИХ ДАНИХ (PM-21)}
a. Забезпечити доступ громадськості до інформації із забезпечення приватності в організації та можливість комунікації з уповноваженою посадовою особою з питань забезпечення приватності щодо:\\ 
1. дати, характеру та мети кожного розкриття запису;\\ 
2. імені та адреси особи або організації, щодо яких було зроблено розкриття даних.\\ 
b. Обліковувати та зберігати випадки розкриття персональних даних протягом терміну дії запису або п'яти років після розкриття інформації.\\ 
c. Здійснювати облік випадків розкриття персональних даних, доступних особі, зазначеній у записі за запитом.

\vspace{10pt}
\textit{Немає параметрів для цього контролю.}
\vspace{10pt}


\subsection{УПРАВЛІННЯ ЯКІСТЮ ПЕРСОНАЛЬНИХ ДАНИХ (PM-22)}
Розробити та задокументувати загальноорганізаційну політику та процедури, які довзолять:\\ 
a. Проводити огляд точності, актуальності, своєчасності та повноти персональних даних протягом їх життєвого циклу;\\ 
b. Коригувати або видаляти неточну або застарілу інформацію;\\ 
c. Інформувати осіб або інші відповідні організації про внесення змін або видалення персональної інформації;\\ 
d. Оскаржувати відмови на запити щодо коригування чи видалення.

\vspace{10pt}
{\footnotesize\bfseries
No: 1\\
Name: pm\_\allowbreak{}22\_\allowbreak{}01\\
Type: list\\
Default: [\textquotedbl{}admin\textquotedbl{}, \textquotedbl{}security\_\allowbreak{}officer\textquotedbl{}]
}

{\footnotesize Розроблені та задокументовані загальноорганізаційні політики управління якістю персональних даних}


{\footnotesize\bfseries
No: 2\\
Name: pm\_\allowbreak{}22\_\allowbreak{}02\\
Type: list\\
Default: [\textquotedbl{}admin\textquotedbl{}, \textquotedbl{}security\_\allowbreak{}officer\textquotedbl{}]
}

{\footnotesize Розроблені та задокументовані загальноорганізаційні процедури управління якістю персональних даних; PM-22a.[01] передбачено в політиці перевірку точності інформації, що ідентифікує особу, протягом усього життєвого циклу інформації; PM-22a.[02] передбачено в політиці перегляд актуальності інформації, що ідентифікує особу, протягом життєвого циклу інформації; PM-22a.[03] передбачено в політиці перевірку своєчасності інформації, що ідентифікує особу, протягом усього життєвого циклу інформації; PM-22a.[04] передбачено в політиці перевірку повноти інформації, що ідентифікує особу, протягом життєвого циклу інформації; PM-22a.[05] передбачено в процедурах перевірку точності інформації, що ідентифікує особу, протягом усього життєвого циклу інформації; PM-22a.[06] передбачено в процедурах перегляд актуальності інформації, що ідентифікує особу, протягом усього життєвого циклу інформації; PM-22a.[07] процедури передбачають перевірку своєчасності інформації, що ідентифікує особу, протягом усього життєвого циклу інформації; PM-22a.[08] передбачено в процедурах перевірку повноти інформації, що ідентифікує особу, протягом життєвого циклу інформації; PM-22b.[01] передбачено в політиці виправлення або видалення неточної або застарілої персональної інформації, що ідентифікує особу; PM-22b.[02] передбачено в процедурах виправлення або видалення неточної або застарілої персональної інформації; PM-22c.[01] передбачено в політиці розсилання повідомлень про виправлену або видалену персональну інформацію фізичним особам або іншим відповідним суб'єктам; PM-22c.[02] передбачено в процедурах повідомлення про виправлення або видалення персональних даних фізичним особам або іншим відповідним суб'єктам; PM-22d.[01] передбачено в політиці оскарження негативних рішень щодо запитів на виправлення або видалення; PM-22d.[02] передбачені процедури оскарження негативних рішень щодо запитів на виправлення або видалення}




\subsection{ОРГАН УПРАВЛІННЯ ПЕРСОНАЛЬНИМИ ДАНИМИ (PM-23)}
створити орган управління персональними даними, на якого покладено [Призначення: визначені організацією функції] та виконання [Призначення: визначені організацією обов'язки].

\vspace{10pt}
{\footnotesize\bfseries
No: 1\\
Name: pm\_\allowbreak{}23\_\allowbreak{}01\\
Type: string\\
Default: nil
}

{\footnotesize Створено орган управління даними, що складається з ролей з обов'язками}


{\footnotesize\bfseries
No: 2\\
Name: pm\_\allowbreak{}23\_\allowbreak{}odp\_\allowbreak{}01\\
Type: list\\
Default: [\textquotedbl{}admin\textquotedbl{}, \textquotedbl{}security\_\allowbreak{}officer\textquotedbl{}]
}

{\footnotesize Визначені ролі органу управління персональними даними}


{\footnotesize\bfseries
No: 3\\
Name: pm\_\allowbreak{}23\_\allowbreak{}odp\_\allowbreak{}02\\
Type: list\\
Default: [\textquotedbl{}admin\textquotedbl{}, \textquotedbl{}security\_\allowbreak{}officer\textquotedbl{}]
}

{\footnotesize Визначені обов'язки органу управління персональними даними}




\subsection{ОРГАН З ПИТАНЬ ЦІЛІСНОСТІ ДАНИХ (PM-24)}
Створити орган з питань цілісності даних для здійснення:\\ 
a. Розгляду пропозицій щодо проведення відповідної програми або участі у ній.\\ 
b. Проведення огляду усіх поточних програм, в яких бере участь організація.

\vspace{10pt}
{\footnotesize\bfseries
No: 1\\
Name: pm\_\allowbreak{}24\_\allowbreak{}01\\
Type: integer\\
Default: 30
}

{\footnotesize Створено орган з питань цілісності даних; PM-24a. розглядає орган з питань цілісності даних пропозиції щодо проведення або участі у відповідній програмі; PM-24b. проводить орган з питань цілісності даних щорічну перевірку всіх програм співставлення, в яких агентство брало участь}




\subsection{МІНІМІЗАЦІЯ КІЛЬКОСТІ ПЕРСОНАЛЬНИХ ДАНИХ, ЩО ВИКОРИСТОВУЮТЬСЯ ПІД ЧАС ТЕСТУВАННЯ, НАВЧАННЯ ТА ДОСЛІДЖЕНЬ (PM-25)}
a. Розробити та впровадити політики та процедури, спрямовані на врегулювання питань використання персональних даних для внутрішнього тестування, навчання та досліджень.\\ 
b. Вжити заходи щодо обмеження або зведення до мінімуму кількості персональних даних, які використовуються для внутрішнього тестування, навчання та досліджень.\\ 
c. Надавати дозвіл на використання персональних даних, коли така інформація вимагається для внутрішнього тестування, навчання і досліджень.\\ 
d. Здійснювати огляд та оновлення політик та процедур, спрямованих на врегулювання питань використання персональних даних для внутрішнього тестування, навчання та досліджень [Призначення: з визначеною організацією частотою].

\vspace{10pt}
{\footnotesize\bfseries
No: 1\\
Name: pm\_\allowbreak{}25\_\allowbreak{}odp\_\allowbreak{}01\\
Type: list\\
Default: [\textquotedbl{}admin\textquotedbl{}, \textquotedbl{}security\_\allowbreak{}officer\textquotedbl{}]
}

{\footnotesize Визначено періодичність перегляду політик, які стосуються використання персональних даних для внутрішнього тестування, навчання та досліджень}


{\footnotesize\bfseries
No: 2\\
Name: pm\_\allowbreak{}25\_\allowbreak{}odp\_\allowbreak{}02\\
Type: list\\
Default: [\textquotedbl{}admin\textquotedbl{}, \textquotedbl{}security\_\allowbreak{}officer\textquotedbl{}]
}

{\footnotesize Визначено періодичність оновлення політик, які стосуються використання персональних даних для внутрішнього тестування, навчання та досліджень}


{\footnotesize\bfseries
No: 3\\
Name: pm\_\allowbreak{}25\_\allowbreak{}odp\_\allowbreak{}03\\
Type: list\\
Default: [\textquotedbl{}admin\textquotedbl{}, \textquotedbl{}security\_\allowbreak{}officer\textquotedbl{}]
}

{\footnotesize Визначено періодичність перегляду процедур, які стосуються використання персональних даних для внутрішнього тестування, навчання та досліджень}


{\footnotesize\bfseries
No: 4\\
Name: pm\_\allowbreak{}25\_\allowbreak{}odp\_\allowbreak{}04\\
Type: list\\
Default: [\textquotedbl{}admin\textquotedbl{}, \textquotedbl{}security\_\allowbreak{}officer\textquotedbl{}]
}

{\footnotesize Визначено періодичність оновлення процедур, які стосуються використання персональних даних для внутрішнього тестування, навчання та досліджень; PM-25a.[01] розроблені та задокументовані політики, які регулюють використання персональних даних для внутрішнього тестування; PM-25a.[02] розроблені та задокументовані політики, які стосуються використання персональних даних для внутрішнього навчання PM-25a.[03] розроблені та задокументовані політики, які регулюють використання персональних даних для внутрішніх досліджень; PM-25a.[04] розроблені та задокументовані процедури, які стосуються використання персональних даних для внутрішнього тестування; PM-25a.[05] розроблені та задокументовані процедури, які стосуються використання персональних даних для внутрішнього навчання; PM-25a.[06] розроблені та задокументовані процедури, які стосуються використання персональних даних для внутрішніх досліджень; PM-25a.[07] впроваджено політику, яка регулює використання персональних даних для внутрішнього тестування; PM-25a.[08] впроваджуються політики, які стосуються використання персональних даних для навчання; PM-25a.[09] впроваджуються політики, які стосуються використання персональної інформації для досліджень; PM-25a.[10] впроваджені процедури, які стосуються використання персональних даних для внутрішнього тестування; PM-25a.[11] впроваджені процедури, які стосуються використання персональної інформації для навчання; PM-25a.[12] впроваджені процедури, які стосуються використання особистої інформації для досліджень; PM-25b.[01] обмежено або зведено до мінімуму кількість персональних даних, що використовуються для цілей внутрішнього тестування; PM-25b.[02] обмежено або зведено до мінімуму обсяг інформації, що ідентифікує особу, яка використовується для внутрішніх навчальних цілей; PM-25b.[03] обмежено або зведено до мінімуму обсяг персональних даних, що використовуються для внутрішніх досліджень; PM-25c.[01] дозволено використання внутрішнього тестування; персональних даних для PM-25c.[02] дозволено використання внутрішнього навчання; персональних даних для PM-25c.[03] дозволено необхідне використання персональних даних для внутрішніх досліджень; PM-25d.[01] переглядаються політики частота; PM-25d.[02] оновлюються політики частота; PM-25d.[03] переглядаються процедури частота; PM-25d.[04] оновлюються процедури частота}




\subsection{УПРАВЛІННЯ СКАРГАМИ (PM-26)}
Впровадити процес отримання та реагування на скарги, проблеми чи запитання від фізичних осіб щодо організаційної практики забезпечення приватності, який охоплює:\\ 
a. механізми, які легко використовувати та які є легкодоступними для громадськості;\\ 
b. усю інформацію, необхідну для успішного подання скарг;\\ 
c. механізми відстеження, що забезпечують отримання всіх скарг та їх вчасний і належний розгляд протягом [Призначення: визначений організацією період часу];\\ 
d. підтвердження отримання скарг, заявлених проблем чи запитань від фізичних осіб протягом [Призначення: визначений організацією період часу];\\ 
e. надання відповідей на отримані скарги, заявлені проблеми чи запитання від фізичних осіб протягом [Призначення: визначений організацією період часу].

\vspace{10pt}
{\footnotesize\bfseries
No: 1\\
Name: pm\_\allowbreak{}26\_\allowbreak{}01\\
Type: string\\
Default: nil
}

{\footnotesize Впроваджено процес отримання скарг, занепокоєнь або запитань від фізичних осіб про безпеку та конфіденційність в організації}


{\footnotesize\bfseries
No: 2\\
Name: pm\_\allowbreak{}26\_\allowbreak{}02\\
Type: integer\\
Default: 30
}

{\footnotesize Впроваджено процес реагування на скарги, занепокоєння або запитання від фізичних осіб про безпеку та конфіденційність в організації; PM-26a.[01] включає процес управління скаргами механізми, які є простими у використанні для громадськості; PM-26c.[02] включає процес управління скаргами механізми, які є легкодоступними для громадськості; PM-26b. містить процес управління скаргами всю інформацію, необхідну для успішного подання скарг; PM-26c.[01] включає процес управління скаргами механізми відстеження, які гарантують, що всі скарги будуть розглянуті протягом періоду часу; період для підтвердження PM-26c.[02] включає процес управління скаргами механізми відстеження, щоб гарантувати, що всі скарги розглядаються протягом часового періоду; PM-26d. передбачає процес управління скаргами підтвердження отримання скарг, занепокоєнь або запитань від фізичних осіб протягом часового періоду; PM-26e. включає процес управління скаргами реагування на скарги, занепокоєння або питання від фізичних осіб протягом часового періоду}


{\footnotesize\bfseries
No: 3\\
Name: pm\_\allowbreak{}26\_\allowbreak{}odp\_\allowbreak{}01\\
Type: integer\\
Default: 30
}

{\footnotesize Визначено період часу, протягом якого мають бути розглянуті скарги (в тому числі звернення або питання) від фізичних осіб}


{\footnotesize\bfseries
No: 4\\
Name: pm\_\allowbreak{}26\_\allowbreak{}odp\_\allowbreak{}02\\
Type: integer\\
Default: 30
}

{\footnotesize Визначено період часу, протягом якого мають бути оброблені скарги (в тому числі звернення або питання) від фізичних осіб}


{\footnotesize\bfseries
No: 5\\
Name: pm\_\allowbreak{}26\_\allowbreak{}odp\_\allowbreak{}03\\
Type: integer\\
Default: 30
}

{\footnotesize Визначено часовий отримання скарг}


{\footnotesize\bfseries
No: 6\\
Name: pm\_\allowbreak{}26\_\allowbreak{}odp\_\allowbreak{}04\\
Type: string\\
Default: nil
}

{\footnotesize Визначено термін для відповіді на скарги}




\subsection{ЗВІТНІСТЬ З ПИТАНЬ ЗАБЕЗПЕЧЕННЯ ПРИВАТНОСТІ (PM-27)}
a. Визначити та задокументувати:\\ 
1. припущення, що впливають на оцінку ризиків, реагування на ризики та моніторинг ризиків;\\ 
2. обмеження, що впливають на оцінку ризиків, реагування на ризики та моніторинг ризиків;\\ 
3. пріоритети та компроміси, які розглядаються організацією для здійснення управління ризиками;\\ 
4. стійкість організації до ризиків.\\ 
b. Поінформувати [Призначення: визначений організацією персонал] про результати визначення ризиків.\\ 
c. Переглядати та оновлювати підходи щодо визначення ризиків [Призначення: з визначеною організацією частотою].

\vspace{10pt}
{\footnotesize\bfseries
No: 1\\
Name: pm\_\allowbreak{}27\_\allowbreak{}odp\_\allowbreak{}01\\
Type: string\\
Default: nil
}

{\footnotesize Визначені звіти з питань забезпечення приватності}


{\footnotesize\bfseries
No: 2\\
Name: pm\_\allowbreak{}27\_\allowbreak{}odp\_\allowbreak{}02\\
Type: string\\
Default: nil
}

{\footnotesize Визначені органи нагляду за дотриманням приватності}


{\footnotesize\bfseries
No: 3\\
Name: pm\_\allowbreak{}27\_\allowbreak{}odp\_\allowbreak{}03\\
Type: list\\
Default: [\textquotedbl{}admin\textquotedbl{}, \textquotedbl{}security\_\allowbreak{}officer\textquotedbl{}]
}

{\footnotesize Визначені посадові особи, відповідальні за контроль і дотриманням програми приватності}


{\footnotesize\bfseries
No: 4\\
Name: pm\_\allowbreak{}27\_\allowbreak{}odp\_\allowbreak{}04\\
Type: list\\
Default: [\textquotedbl{}admin\textquotedbl{}, \textquotedbl{}security\_\allowbreak{}officer\textquotedbl{}]
}

{\footnotesize Визначена періодичність перегляду та оновлення звітів про приватність; PM-27a. розроблено приватності; PM-27a.01 передаються звіти з питань забезпечення приватності до наглядових органів, щоб продемонструвати підзвітність законодавчим, регуляторним та політичним мандатам щодо приватності; звіти з питань PM-27a.02[01] поширюються звіти про конфіденційність серед посадових осіб; PM-27a.02[02] поширюються звіти з питань забезпечення приватності серед іншого персоналу, відповідального за контроль за дотриманням програми конфіденційності; PM-27b. переглядаються та оновлюються звіти з питань забезпечення приватності частота}




\subsection{ОЦІНКА РИЗИКІВ (PM-28)}


\vspace{10pt}
{\footnotesize\bfseries
No: 1\\
Name: pm\_\allowbreak{}28\_\allowbreak{}odp\_\allowbreak{}01\\
Type: list\\
Default: [\textquotedbl{}admin\textquotedbl{}, \textquotedbl{}security\_\allowbreak{}officer\textquotedbl{}]
}

{\footnotesize Визначено персонал, який отримуватиме результати визначення ризиків}


{\footnotesize\bfseries
No: 2\\
Name: pm\_\allowbreak{}28\_\allowbreak{}odp\_\allowbreak{}02\\
Type: list\\
Default: [\textquotedbl{}admin\textquotedbl{}, \textquotedbl{}security\_\allowbreak{}officer\textquotedbl{}]
}

{\footnotesize Визначено періодичність перегляду та міркувань щодо структуризації ризиків; PM-28a.01[01] визначені та задокументовані припущення, що впливають на оцінку ризиків; PM-28a.01[02] визначені та задокументовані припущення, що впливають на реагування ризиків; PM-28a.01[03] визначені та задокументовані припущення, що впливають на моніторинг ризиків; PM-28a.02[01] визначені та задокументовані обмеження, що впливають на оцінку ризиків; PM-28a.02[02] визначені та задокументовані обмеження, що впливають на реагування на ризики; PM-28a.02[03] визначені та задокументовані обмеження, що впливають на моніторинг ризиків; оновлення PM-28a.03[01] визначені та задокументовані пріоритети, розглядаються організацією для управління ризиками; які PM-28a.03[02] визначені та задокументовані компроміси, розглядаються організацією для управління ризиками; які PM-28a.04 визначена та задокументована організаційна толерантність до ризиків; PM-28b. поширюються результати діяльності з фреймворкінгу ризиків серед персоналу ; PM-28c. переглядаються та оновлюються міркування фреймінгу ризиків частота. щодо}




\subsection{РОЛІ КЕРІВНИКІВ ПРОГРАМИ УПРАВЛІННЯ РИЗИКАМИ (PM-29)}
a. Розробити план управління ризиками ланцюга постачання, пов'язаного з розробкою, придбанням, обслуговуванням та утилізацією систем, компонентів системи та послуг для системи.\\ 
b. Реалізувати план управління ризиками ланцюга постачання послідовно та наскрізно по всій організації.\\ 
c. Переглядати й оновлювати план управління ризиками ланцюга постачання [Призначення: з визначеною організацією частотою] або, якщо потрібно, у разі змін в організації.

\vspace{10pt}
\textit{Немає параметрів для цього контролю.}
\vspace{10pt}


\subsection{ПЛАН УПРАВЛІННЯ РИЗИКАМИ ЛАНЦЮГА ПОСТАЧАННЯ (PM-30)}


\vspace{10pt}
{\footnotesize\bfseries
No: 1\\
Name: pm\_\allowbreak{}30\_\allowbreak{}odp\\
Type: list\\
Default: [\textquotedbl{}admin\textquotedbl{}, \textquotedbl{}security\_\allowbreak{}officer\textquotedbl{}]
}

{\footnotesize Призначено старшу посадову особу, відповідальну за управління ризиками ланцюга постачання; PM-30a.[01] узгоджує старша посадова особа, відповідальна за управління ризиками ланцюга постачання, процеси управління інформаційною безпекою та конфіденційністю з процесами стратегічного, операційного та бюджетного планування; PM-30a.[02] створена посада (функція) ризик-менеджер; PM-30a.[03] розглядає та аналізує керівник з управління ризиками (функція) ризики з точки зору всієї організації; PM-30a.[04] забезпечує керівник (функція) з управління ризиками узгоджене управління ризиками в межах всієї організації. PM-30a.[05] стратегія управління ризиками ланцюга постачання враховує ризики, пов'язані з придбанням систем; PM-30a.[06] стратегія управління ризиками ланцюга поставок враховує ризики, пов'язані з придбанням компонентів системи; PM-30a.[07] стратегія управління ризиками ланцюга поставок враховує ризики, пов'язані з придбанням системних послуг; PM-30a.[08] тратегія управління ризиками ланцюга поставок враховує ризики, пов'язані з обслуговуванням систем; PM-30a.[09] стратегія управління ризиками ланцюга поставок враховує ризики, пов'язані з обслуговуванням компонентів системи; PM-30a.[10] стратегія управління ризиками ланцюга поставок враховує ризики, пов'язані з обслуговуванням системних послуг; PM-30a.[11] враховує стратегія управління ризиками ланцюга постачання ризики, пов'язані з утилізацією систем; PM-30a.[12] враховує стратегія управління ризиками ланцюга постачання ризики, пов'язані з утилізацією компонентів системи; PM-30a.[13] стратегія управління ризиками ланцюга поставок враховує ризики, пов'язані з утилізацією системних послуг; PM-30b. стратегія управління ризиками ланцюга послідовно впроваджується в організації; PM-30c. переглядається та оновлюється стратегія управління ризиками ланцюга постачання частота або в міру необхідності у зв'язку з організаційними змінами поставок}




\subsection{ПЛАН БЕЗПЕРЕРВНОГО МОНІТОРИНГУ (PM-31)}
Розробити план безперервного моніторингу в масштабах всієї організації та впровадити програми безперервного моніторингу, які включають:\\ 
a. Встановити відповідні показників для моніторингу в масштабах всієї організації [Призначення: визначені організацією показники];\\ 
b. Встановити [Призначення: частота, визначеної організацією] для здійснення моніторингу та [Призначення: періодичність, визначена організацією] проведення оцінки ефективності контролю;\\ 
c. Постійний моніторинг визначених організацією показників відповідно до стратегії безперервного моніторингу;\\ 
d. Співставлення та аналіз інформації, отриманої в результаті здійснення моніторингу, та контрольних оцінок;\\ 
e. Заходи реагування на результати аналізу оцінок контролю та моніторингових даних;\\ 
f. Звітування про стан безпеки та приватності систем організації перед [Призначення: визначеним організацією персоналом чи посадовою особою] [Призначення: з визначеною організацією періодичністю].

\vspace{10pt}
{\footnotesize\bfseries
No: 1\\
Name: pm\_\allowbreak{}31\_\allowbreak{}01\\
Type: list\\
Default: [\textquotedbl{}admin\textquotedbl{}, \textquotedbl{}security\_\allowbreak{}officer\textquotedbl{}]
}

{\footnotesize Розроблена загальноорганізаційна стратегія безперервного моніторингу; PM-31a. впроваджуються програми безперервного моніторингу, які включають встановлення параметрів, що підлягають моніторингу; PM-31b.[01] впроваджено програми безперервного моніторингу, які встановлюють частоту для моніторингу; PM-31b.[02] впроваджуються програми безперервного моніторингу, які встановлюють частоту для оцінки ефективності контролю; про стан PM-31c. впроваджуються програми безперервного моніторингу, які включають моніторинг параметрів на постійній основі відповідно до стратегії безперервного моніторингу; PM-31d.[01] впроваджуються програми безперервного моніторингу, які включають співставлення інформації, отриманої в результаті контрольних оцінок та моніторингу; PM-31d.[02] впроваджуються програми постійного моніторингу, які включають аналіз інформації, отриманої в результаті контрольних оцінок та моніторингу; PM-31e.[01] впроваджуються програми безперервного моніторингу, які передбачають заходи реагування на аналіз інформації, отриманої в результаті оцінки результатів контролю; PM-31e.[02] впроваджуються програми безперервного моніторингу, які передбачають заходи реагування на результати аналізу інформації, отриманої під час моніторингу; PM-31f.[01] впроваджено програми безперервного моніторингу, які передбачають звітування про стан безпеки систем організації перед персонал або ролі частота; PM-31f.[02] впроваджені програми безперервного моніторингу, які передбачають звітування про стан конфіденційності організаційних систем перед персонал або ролі частота}


{\footnotesize\bfseries
No: 2\\
Name: pm\_\allowbreak{}31\_\allowbreak{}odp\_\allowbreak{}01\\
Type: string\\
Default: nil
}

{\footnotesize Визначені параметри для безперервного моніторингу в масштабах всієї організації}


{\footnotesize\bfseries
No: 3\\
Name: pm\_\allowbreak{}31\_\allowbreak{}odp\_\allowbreak{}02\\
Type: string\\
Default: \textquotedbl{}щорічно\textquotedbl{}
}

{\footnotesize Визначено періодичність моніторингу}


{\footnotesize\bfseries
No: 4\\
Name: pm\_\allowbreak{}31\_\allowbreak{}odp\_\allowbreak{}03\\
Type: string\\
Default: \textquotedbl{}щорічно\textquotedbl{}
}

{\footnotesize Визначена періодичність оцінки ефективності контролю}


{\footnotesize\bfseries
No: 5\\
Name: pm\_\allowbreak{}31\_\allowbreak{}odp\_\allowbreak{}04\\
Type: list\\
Default: [\textquotedbl{}admin\textquotedbl{}, \textquotedbl{}security\_\allowbreak{}officer\textquotedbl{}]
}

{\footnotesize Визначено персонал або ролі для звітування про стан безпеки систем організації}


{\footnotesize\bfseries
No: 6\\
Name: pm\_\allowbreak{}31\_\allowbreak{}odp\_\allowbreak{}05\\
Type: list\\
Default: [\textquotedbl{}admin\textquotedbl{}, \textquotedbl{}security\_\allowbreak{}officer\textquotedbl{}]
}

{\footnotesize Визначено персонал або ролі для звітування про стан конфіденційності систем організації}


{\footnotesize\bfseries
No: 7\\
Name: pm\_\allowbreak{}31\_\allowbreak{}odp\_\allowbreak{}06\\
Type: string\\
Default: \textquotedbl{}щорічно\textquotedbl{}
}

{\footnotesize Визначено періодичність звітування про стан безпеки систем організації}


{\footnotesize\bfseries
No: 8\\
Name: pm\_\allowbreak{}31\_\allowbreak{}odp\_\allowbreak{}07\\
Type: string\\
Default: \textquotedbl{}щорічно\textquotedbl{}
}

{\footnotesize Визначено періодичність звітування конфіденційності систем організації}




\subsection{ПРИЗНАЧЕННЯ (PM-32)}
Включіть ролі й обов'язки з безпеки та приватності в опис посади в організації.

\vspace{10pt}
{\footnotesize\bfseries
No: 1\\
Name: pm\_\allowbreak{}32\_\allowbreak{}01\\
Type: string\\
Default: nil
}

{\footnotesize Аналізуються допоміжні послуги або функції, необхідні для виконання місії, для забезпечення того, щоб інформаційні ресурси використовувалися відповідно до їхнього призначення}




\section{PT}
Клас заходів захисту PT --- ПОВНОВАЖЕННЯ

Цей клас фокусується на дотриманні законодавства щодо захисту персональних даних, отриманні згоди та забезпеченні прав суб'єктів даних.

\textbf{Перелік заходів захисту:} Pt (PT-1); Повноваження на обробку персональних даних (PT-2); Повноваження на тегування даних (PT-2(1)); Повноваження автоматизація (PT-2(2)); Цілі обробки персональних даних (PT-3); Тегування даних (PT-3(1)); Автоматизація (PT-3(2)); Згода на обробку персональних даних (PT-4); Індивідуальна згода на обробку персональних даних (PT-4(1)); Своєчасна згода на обробку персональних даних (PT-4(2)); Відкликання (PT-4(3)); Повідомлення про конфіденційність (PT-5); Повідомлення про конфіденційність повідомлення про конфіденційність (PT-5(1)); Повідомлення про конфіденційність (PT-5(2)); Система записів повідомлень про конфіденційність (PT-6); Система записів повідомлень про конфіденційність звичайне використання (PT-6(1)); (PT-6(2)); Спеціальні категорії персональних даних (PT-7); Спеціальні категорії персональних соціального страхування (PT-7(1)); Інформація про першу поправку (PT-7(2)); Вимоги до відповідності (PT-8).

\subsection{PT (PT-1)}
a. Розробіть, задокументуйте та поширте [Призначення: персонал або ролі, визначені організацією]:\\ 
1. [Вибір (один або декілька): Рівень організації; Рівень місії/бізнес-процесу; рівень системи], обробки персональних даних та політики прозорості, який: a) розглядає мету, сферу діяльності, ролі, відповідальність, зобов'язання керівництва, координацію між організаційними підрозділами та відповідність; b) відповідає чинним законам, розпорядженням, директивам, положенням, політикам, стандартам і рекомендаціям.\\ 
2. Процедури для реалізації політики обробки та прозорості персональних даних, а також пов'язані засоби контролю;\\ 
b. Призначте [Призначення: посадову особу, визначену організацією] для керування розробкою, документуванням і розповсюдженням політики й процедур щодо обробки персональних даних та прозорості;\\ 
c. Перегляньте та оновіть поточні процедури обробки та прозорість персональних даних:\\ 
1. Політика [Призначення: частота, визначена організацією] і наступні [Призначення: події, визначені організацією];\\ 
2. Процедури [Призначення: частота, визначена організацією] та наступні [Призначення: подія, визначена організацією].

\vspace{10pt}
{\footnotesize\bfseries
No: 1\\
Name: pt\_\allowbreak{}1\_\allowbreak{}odp\_\allowbreak{}01\\
Type: list\\
Default: [\textquotedbl{}admin\textquotedbl{}, \textquotedbl{}security\_\allowbreak{}officer\textquotedbl{}]
}

{\footnotesize Визначено персонал або ролі, на які поширюється політика обробки персональних даних та забезпечення прозорості}


{\footnotesize\bfseries
No: 2\\
Name: pt\_\allowbreak{}1\_\allowbreak{}odp\_\allowbreak{}02\\
Type: list\\
Default: [\textquotedbl{}admin\textquotedbl{}, \textquotedbl{}security\_\allowbreak{}officer\textquotedbl{}]
}

{\footnotesize Визначено персонал або ролі, на які поширюються процедури обробки персональних даних та політики прозорості}


{\footnotesize\bfseries
No: 3\\
Name: pt\_\allowbreak{}1\_\allowbreak{}odp\_\allowbreak{}03\\
Type: string\\
Default: nil
}

{\footnotesize Вибрано одне або декілька з наступних ЗНАЧЕНЬ ПАРАМЕТРІВ: \{рівень організації; рівень місії/бізнеспроцесу; рівень системи\}}


{\footnotesize\bfseries
No: 4\\
Name: pt\_\allowbreak{}1\_\allowbreak{}odp\_\allowbreak{}04\\
Type: list\\
Default: [\textquotedbl{}admin\textquotedbl{}, \textquotedbl{}security\_\allowbreak{}officer\textquotedbl{}]
}

{\footnotesize Визначено посадову особу, яка керуватиме політикою та процедурами обробки персональних даних, а також політикою та процедурами прозорості}


{\footnotesize\bfseries
No: 5\\
Name: pt\_\allowbreak{}1\_\allowbreak{}odp\_\allowbreak{}05\\
Type: list\\
Default: [\textquotedbl{}admin\textquotedbl{}, \textquotedbl{}security\_\allowbreak{}officer\textquotedbl{}]
}

{\footnotesize Визначена періодичність перегляду та оновлення поточної політики обробки та прозорості інформації, що ідентифікує особу}


{\footnotesize\bfseries
No: 6\\
Name: pt\_\allowbreak{}1\_\allowbreak{}odp\_\allowbreak{}06\\
Type: list\\
Default: [\textquotedbl{}admin\textquotedbl{}, \textquotedbl{}security\_\allowbreak{}officer\textquotedbl{}]
}

{\footnotesize Є події, які вимагають перегляду та оновлення поточної політики обробки персональних даних та прозорості}


{\footnotesize\bfseries
No: 7\\
Name: pt\_\allowbreak{}1\_\allowbreak{}odp\_\allowbreak{}07\\
Type: list\\
Default: [\textquotedbl{}admin\textquotedbl{}, \textquotedbl{}security\_\allowbreak{}officer\textquotedbl{}]
}

{\footnotesize Визначена частота, з якою переглядаються та оновлюються поточні процедури обробки персональних даних та забезпечення прозорості}




\subsection{ПОВНОВАЖЕННЯ НА ОБРОБКУ ПЕРСОНАЛЬНИХ ДАНИХ (PT-2)}
a. визначити та задокументувати [Призначення: повноваження, визначені організацією], які дозволяють [Призначення: обробку, визначену організацією] персональної інформації;\\ 
b. обмежити [Призначення: обробку, визначену організацією] персональної інформації лише таким чином, яким дозволено (тільки до того, що дозволено)

\vspace{10pt}
{\footnotesize\bfseries
No: 1\\
Name: pt\_\allowbreak{}2\_\allowbreak{}odp\_\allowbreak{}01\\
Type: list\\
Default: [\textquotedbl{}admin\textquotedbl{}, \textquotedbl{}security\_\allowbreak{}officer\textquotedbl{}]
}

{\footnotesize Визначені повноваження щодо надання дозволу на обробку (визначені в PT-02\_\allowbreak{}ODP[02]) персональних даних}


{\footnotesize\bfseries
No: 2\\
Name: pt\_\allowbreak{}2\_\allowbreak{}odp\_\allowbreak{}02\\
Type: list\\
Default: [\textquotedbl{}admin\textquotedbl{}, \textquotedbl{}security\_\allowbreak{}officer\textquotedbl{}]
}

{\footnotesize Визначено тип обробки персональних даних}


{\footnotesize\bfseries
No: 3\\
Name: pt\_\allowbreak{}2\_\allowbreak{}odp\_\allowbreak{}03\\
Type: list\\
Default: [\textquotedbl{}admin\textquotedbl{}, \textquotedbl{}security\_\allowbreak{}officer\textquotedbl{}]
}

{\footnotesize Визначено тип обробки підлягають обмеженню; PT-02a. визначено та задокументовано орган, який дозволяє обробку персональних даних; PT-02b. обробка персональних обмежується лише таким чином, яким дозволено. персональних даних, що даних,}




\subsubsection{ПОВНОВАЖЕННЯ НА ТЕГУВАННЯ ДАНИХ (PT-2(1))}


\vspace{10pt}
{\footnotesize\bfseries
No: 1\\
Name: pt\_\allowbreak{}2\_\allowbreak{}1\_\allowbreak{}01\\
Type: list\\
Default: [\textquotedbl{}admin\textquotedbl{}, \textquotedbl{}security\_\allowbreak{}officer\textquotedbl{}]
}

{\footnotesize Теги даних, що містять <PT-02(01) \_\allowbreak{}ODP[01] санкціонована обробка>, прикріплені до <PT-02(01) \_\allowbreak{}ODP[02] елементів інформації, що іденти0фікує особу>}




\subsubsection{ПОВНОВАЖЕННЯ АВТОМАТИЗАЦІЯ (PT-2(2))}


\vspace{10pt}
{\footnotesize\bfseries
No: 1\\
Name: pt\_\allowbreak{}2\_\allowbreak{}2\_\allowbreak{}01\\
Type: list\\
Default: [\textquotedbl{}admin\textquotedbl{}, \textquotedbl{}security\_\allowbreak{}officer\textquotedbl{}]
}

{\footnotesize Управління дотриманням санкціонованої обробки персональних даних здійснюється за допомогою <PТ02(02)\_\allowbreak{}ODP автоматизовані механізми обробки персональних даних>}


{\footnotesize\bfseries
No: 2\\
Name: pt\_\allowbreak{}2\_\allowbreak{}2\_\allowbreak{}odp\\
Type: list\\
Default: [\textquotedbl{}admin\textquotedbl{}, \textquotedbl{}security\_\allowbreak{}officer\textquotedbl{}]
}

{\footnotesize Визначені автоматизовані механізми, які використовуються для управління дотриманням санкціонованої обробки інформації, що ідентифікує особу}




\subsection{ЦІЛІ ОБРОБКИ ПЕРСОНАЛЬНИХ ДАНИХ (PT-3)}
a. визначити та задокументувати [Призначення: цілі, визначені організацією] для обробки персональних даних;\\ 
b. описати мету (цілі) у публічних повідомленнях про конфіденційність і політиках організації;\\ 
c. обмежити [Призначення: обробку, визначену організацією] персональних даних лише тією, яка сумісна з визначеною ціллю(ями);\\ 
d. відстежувати зміни в обробці персональних даних та впроваджувати [Завдання: визначені організацією механізми], щоб гарантувати, що будь-які зміни вносяться відповідно до [Завдання: визначені організацією вимоги].

\vspace{10pt}
{\footnotesize\bfseries
No: 1\\
Name: pt\_\allowbreak{}3\_\allowbreak{}odp\_\allowbreak{}01\\
Type: list\\
Default: [\textquotedbl{}admin\textquotedbl{}, \textquotedbl{}security\_\allowbreak{}officer\textquotedbl{}]
}

{\footnotesize Визначено мету (цілі) обробки персональних даних}


{\footnotesize\bfseries
No: 2\\
Name: pt\_\allowbreak{}3\_\allowbreak{}odp\_\allowbreak{}02\\
Type: list\\
Default: [\textquotedbl{}admin\textquotedbl{}, \textquotedbl{}security\_\allowbreak{}officer\textquotedbl{}]
}

{\footnotesize Визначена обробка персональних даних, яка підлягає обмеженню}


{\footnotesize\bfseries
No: 3\\
Name: pt\_\allowbreak{}3\_\allowbreak{}odp\_\allowbreak{}03\\
Type: list\\
Default: [\textquotedbl{}admin\textquotedbl{}, \textquotedbl{}security\_\allowbreak{}officer\textquotedbl{}]
}

{\footnotesize Визначені механізми, які мають бути впроваджені для забезпечення того, щоб будь-які зміни в персональних данних, вносилися відповідно до вимог}


{\footnotesize\bfseries
No: 4\\
Name: pt\_\allowbreak{}3\_\allowbreak{}odp\_\allowbreak{}04\\
Type: list\\
Default: [\textquotedbl{}admin\textquotedbl{}, \textquotedbl{}security\_\allowbreak{}officer\textquotedbl{}]
}

{\footnotesize Визначені вимоги до зміни обробкиперсональних даних; PT-03a. визначено та задокументовано мету (цілі) обробки персональних даних; PT-03b.[01] описана мета (цілі) в публічних конфіденційність організації; PT-03b.[02] описана мета (цілі) в політиці організації; PT-03c обробка персональних даних, обмежується лише тим, що є сумісним з визначеною метою (цілями); PT-03d.[01] здійснюється моніторинг змін в обробці персональних даних; PT-03d.[02] впроваджено механізми для забезпечення того, щоб будь-які зміни вносилися відповідно до вимог. повідомленнях про}




\subsubsection{ТЕГУВАННЯ ДАНИХ (PT-3(1))}


\vspace{10pt}
{\footnotesize\bfseries
No: 1\\
Name: pt\_\allowbreak{}3\_\allowbreak{}1\_\allowbreak{}01\\
Type: list\\
Default: [\textquotedbl{}admin\textquotedbl{}, \textquotedbl{}security\_\allowbreak{}officer\textquotedbl{}]
}

{\footnotesize Теги даних, що містять <PT-03(01) \_\allowbreak{}ODP[01] цілі обробки>, приєднані до <PT-03(01) \_\allowbreak{}ODP[02] елементів інформації, що ідентифікує особу>}




\subsubsection{АВТОМАТИЗАЦІЯ (PT-3(2))}


\vspace{10pt}
{\footnotesize\bfseries
No: 1\\
Name: pt\_\allowbreak{}3\_\allowbreak{}2\_\allowbreak{}01\\
Type: list\\
Default: [\textquotedbl{}admin\textquotedbl{}, \textquotedbl{}security\_\allowbreak{}officer\textquotedbl{}]
}

{\footnotesize Відстежуються цілі обробки персональних даних за допомогою автоматизованих механізмів}


{\footnotesize\bfseries
No: 2\\
Name: pt\_\allowbreak{}3\_\allowbreak{}2\_\allowbreak{}odp\\
Type: list\\
Default: [\textquotedbl{}admin\textquotedbl{}, \textquotedbl{}security\_\allowbreak{}officer\textquotedbl{}]
}

{\footnotesize Визначені автоматизовані механізми відстеження цілей обробки персональних даних}




\subsection{Згода на обробку персональних даних (PT-4)}
Впроваджувати [Призначення: інструменти або механізми, визначені організацією], щоб окремі особи давали згоду на обробку їх персональних даних до її збору, що полегшить прийняття обґрунтованих рішень особами.

\vspace{10pt}
{\footnotesize\bfseries
No: 1\\
Name: pt\_\allowbreak{}4\_\allowbreak{}01\\
Type: list\\
Default: [\textquotedbl{}admin\textquotedbl{}, \textquotedbl{}security\_\allowbreak{}officer\textquotedbl{}]
}

{\footnotesize Впроваджено інструменти або механізми для надання фізичними особами згоди на обробку їхніх персональних даних до її збору, які сприяють прийняттю фізичними особами поінформованих рішень}


{\footnotesize\bfseries
No: 2\\
Name: pt\_\allowbreak{}4\_\allowbreak{}odp\\
Type: list\\
Default: [\textquotedbl{}admin\textquotedbl{}, \textquotedbl{}security\_\allowbreak{}officer\textquotedbl{}]
}

{\footnotesize Визначені інструменти або механізми, які мають бути застосовані для надання особами згоди на обробку їхніх персональних даних}




\subsubsection{ІНДИВІДУАЛЬНА ЗГОДА НА ОБРОБКУ ПЕРСОНАЛЬНИХ ДАНИХ (PT-4(1))}


\vspace{10pt}
{\footnotesize\bfseries
No: 1\\
Name: pt\_\allowbreak{}4\_\allowbreak{}1\_\allowbreak{}odp\\
Type: list\\
Default: [\textquotedbl{}admin\textquotedbl{}, \textquotedbl{}security\_\allowbreak{}officer\textquotedbl{}]
}

{\footnotesize Визначені механізми адаптації для обробки окремих елементів дозволів персональних даних; PT-04(01) передбачені механізми, які дозволяють особам пристосовувати дозволи на обробку до вибраних елементів персональних даних}




\subsubsection{СВОЄЧАСНА ЗГОДА НА ОБРОБКУ ПЕРСОНАЛЬНИХ ДАНИХ (PT-4(2))}


\vspace{10pt}
{\footnotesize\bfseries
No: 1\\
Name: pt\_\allowbreak{}4\_\allowbreak{}2\_\allowbreak{}01\\
Type: list\\
Default: [\textquotedbl{}admin\textquotedbl{}, \textquotedbl{}security\_\allowbreak{}officer\textquotedbl{}]
}

{\footnotesize Надаються механізми згоди особам частота та в поєднанні з обробкою персональних даних}




\subsubsection{ВІДКЛИКАННЯ (PT-4(3))}


\vspace{10pt}
{\footnotesize\bfseries
No: 1\\
Name: pt\_\allowbreak{}4\_\allowbreak{}3\_\allowbreak{}01\\
Type: list\\
Default: [\textquotedbl{}admin\textquotedbl{}, \textquotedbl{}security\_\allowbreak{}officer\textquotedbl{}]
}

{\footnotesize Впроваджено інструменти або механізми, які дозволяють фізичним особам відкликати згоду на обробку їхніх персональних даних}


{\footnotesize\bfseries
No: 2\\
Name: pt\_\allowbreak{}4\_\allowbreak{}3\_\allowbreak{}odp\\
Type: list\\
Default: [\textquotedbl{}admin\textquotedbl{}, \textquotedbl{}security\_\allowbreak{}officer\textquotedbl{}]
}

{\footnotesize Визначені інструменти або механізми для відкликання згоди на обробку персональних даних}




\subsection{ПОВІДОМЛЕННЯ ПРО КОНФІДЕНЦІЙНІСТЬ (PT-5)}
Впровадити повідомлення про конфіденційність особам, чиї персональні дані обробляються в системі, які:\\ 
a. доступні окремим особам під час першої взаємодії з організацією, та згодом [Призначення: частота, визначена організацією];\\ 
b. виражені простою мовою;\\ 
c. визначають орган, який надає дозвіл на обробку персональних даних;\\ 
d. визначають цілі, для яких мають оброблятися персональні дані;\\ 
e. включають [Призначення: інформація, визначена організацією].

\vspace{10pt}
{\footnotesize\bfseries
No: 1\\
Name: pt\_\allowbreak{}5\_\allowbreak{}odp\_\allowbreak{}01\\
Type: list\\
Default: [\textquotedbl{}admin\textquotedbl{}, \textquotedbl{}security\_\allowbreak{}officer\textquotedbl{}]
}

{\footnotesize Визначена частота, з якою повідомлення надається особам на рівні первинної взаємодії з організацією}


{\footnotesize\bfseries
No: 2\\
Name: pt\_\allowbreak{}5\_\allowbreak{}odp\_\allowbreak{}02\\
Type: list\\
Default: [\textquotedbl{}admin\textquotedbl{}, \textquotedbl{}security\_\allowbreak{}officer\textquotedbl{}]
}

{\footnotesize Визначена інформація, яка повинна бути включена до повідомлення про обробку персональних даних; PT-05a.[01] направляється особам повідомлення про обробку персональних даних таким чином, щоб вони могли ознайомитися з ним при першій взаємодії з організацією; PT-05a.[02] направляється повідомлення фізичним особам про обробку пермональних даних, таким чином, щоб це повідомлення було згодом доступне фізичним особам частота; PT-05b. направляється фізичним особам повідомлення про обробку персональних даних, яке є чітким, легким для розуміння та містить інформацію про обробку персональних даних простою мовою; PT-05c. направляється фізичним особам повідомлення про обробку персональних даних, яке визначає орган, що надає дозвіл на обробку персональних даних; PT-05d. направляється повідомлення фізичним особам про обробку персональних даних, в якому вказується мета, з якою буде оброблятися персональна інформація; PT-05e. направляється повідомлення фізичним особам про обробку персональних даних, які включають інформацію}




\subsubsection{ПОВІДОМЛЕННЯ ПРО КОНФІДЕНЦІЙНІСТЬ ПОВІДОМЛЕННЯ ПРО КОНФІДЕНЦІЙНІСТЬ (PT-5(1))}


\vspace{10pt}
{\footnotesize\bfseries
No: 1\\
Name: pt\_\allowbreak{}5\_\allowbreak{}1\_\allowbreak{}01\\
Type: list\\
Default: [\textquotedbl{}admin\textquotedbl{}, \textquotedbl{}security\_\allowbreak{}officer\textquotedbl{}]
}

{\footnotesize Надається повідомлення про обробку персональних даних особам у той час і в тому місці, де особа надає персональні дані, у зв'язку з якою здійснюється дія з даними, або періодичність обробки даних частота}


{\footnotesize\bfseries
No: 2\\
Name: pt\_\allowbreak{}5\_\allowbreak{}1\_\allowbreak{}odp\\
Type: list\\
Default: [\textquotedbl{}admin\textquotedbl{}, \textquotedbl{}security\_\allowbreak{}officer\textquotedbl{}]
}

{\footnotesize Визначена періодичність подання обробку персональних даних; повідомлення про}




\subsubsection{ПОВІДОМЛЕННЯ ПРО КОНФІДЕНЦІЙНІСТЬ (PT-5(2))}


\vspace{10pt}
{\footnotesize\bfseries
No: 1\\
Name: pt\_\allowbreak{}5\_\allowbreak{}2\_\allowbreak{}01\\
Type: string\\
Default: nil
}

{\footnotesize Включаються повідомлення про конфіденційність у форми, які збирають інформацію, що буде зберігатися в системі записів Закону про конфіденційність, або ж заяви про конфіденційність надаються на окремих формах, які можуть зберігатися у приватних осіб}




\subsection{СИСТЕМА ЗАПИСІВ ПОВІДОМЛЕНЬ ПРО КОНФІДЕНЦІЙНІСТЬ (PT-6)}
для систем, які обробляють інформацію, яка зберігатиметься в системі записів Закону про конфіденційність:\\ 
a. розробити проект системи повідомлень про записи відповідно до вказівок OMB і подати нову та суттєво змінену систему повідомлень про записи до OMB та відповідних комітетів Конгресу для попереднього розгляду;\\ 
b. опублікувати систему записів повідомлень у Державному реєстрі;\\ 
c. зберігайте повідомлення системи записів точними, оновленими та в обсязі відповідно до впровадженої політики.

\vspace{10pt}
\textit{Немає параметрів для цього контролю.}
\vspace{10pt}


\subsubsection{СИСТЕМА ЗАПИСІВ ПОВІДОМЛЕНЬ ПРО КОНФІДЕНЦІЙНІСТЬ ЗВИЧАЙНЕ ВИКОРИСТАННЯ (PT-6(1))}


\vspace{10pt}
{\footnotesize\bfseries
No: 1\\
Name: pt\_\allowbreak{}6\_\allowbreak{}1\_\allowbreak{}01\\
Type: string\\
Default: \textquotedbl{}щорічно\textquotedbl{}
}

{\footnotesize Переглядаються всі звичайні види використання, опубліковані в повідомленні системи записів періодичність, для забезпечення постійної точності, а також для забезпечення того, щоб звичайні види використання і надалі були сумісними з метою, для якої була зібрана інформація}


{\footnotesize\bfseries
No: 2\\
Name: pt\_\allowbreak{}6\_\allowbreak{}1\_\allowbreak{}odp\\
Type: string\\
Default: \textquotedbl{}щорічно\textquotedbl{}
}

{\footnotesize Визначено періодичність перегляду всіх звичайних видів використання, опублікованих у системі обліку повідомлень}




\subsubsection{(PT-6(2))}


\vspace{10pt}
{\footnotesize\bfseries
No: 1\\
Name: pt\_\allowbreak{}6\_\allowbreak{}2\_\allowbreak{}01\\
Type: string\\
Default: \textquotedbl{}щорічно\textquotedbl{}
}

{\footnotesize Всі винятки із Закону про конфіденційність, заявлені для системи записів, переглядаються частота, щоб переконатися, що вони залишаються доречними та необхідними відповідно до закону}


{\footnotesize\bfseries
No: 2\\
Name: pt\_\allowbreak{}6\_\allowbreak{}2\_\allowbreak{}02\\
Type: string\\
Default: \textquotedbl{}щорічно\textquotedbl{}
}

{\footnotesize Всі винятки із Закону про конфіденційність, заявлені для системи записів, переглядаються частота, щоб переконатися, що вони були оприлюднені як нормативні акти}


{\footnotesize\bfseries
No: 3\\
Name: pt\_\allowbreak{}6\_\allowbreak{}2\_\allowbreak{}03\\
Type: string\\
Default: \textquotedbl{}щорічно\textquotedbl{}
}

{\footnotesize Всі винятки із Закону про конфіденційність, заявлені для системи записів, переглядаються частота, щоб переконатися, що вони точно описані в повідомленні про систему записів}


{\footnotesize\bfseries
No: 4\\
Name: pt\_\allowbreak{}6\_\allowbreak{}2\_\allowbreak{}odp\\
Type: string\\
Default: \textquotedbl{}щорічно\textquotedbl{}
}

{\footnotesize Визначено періодичність перегляду всіх винятків із Закону про конфіденційність, заявлених для системи записів}




\subsection{СПЕЦІАЛЬНІ КАТЕГОРІЇ ПЕРСОНАЛЬНИХ ДАНИХ (PT-7)}
Застосувати [Призначення: умови обробки, визначені організацією] для певних категорій персональних даних.

\vspace{10pt}
{\footnotesize\bfseries
No: 1\\
Name: pt\_\allowbreak{}7\_\allowbreak{}01\\
Type: list\\
Default: [\textquotedbl{}admin\textquotedbl{}, \textquotedbl{}security\_\allowbreak{}officer\textquotedbl{}]
}

{\footnotesize Застосовуються умови обробки до певних категорій персональних даних}


{\footnotesize\bfseries
No: 2\\
Name: pt\_\allowbreak{}7\_\allowbreak{}odp\\
Type: list\\
Default: [\textquotedbl{}admin\textquotedbl{}, \textquotedbl{}security\_\allowbreak{}officer\textquotedbl{}]
}

{\footnotesize Визначені умови обробки, що застосовуються до певних категорій персональних даних}




\subsubsection{СПЕЦІАЛЬНІ КАТЕГОРІЇ ПЕРСОНАЛЬНИХ СОЦІАЛЬНОГО СТРАХУВАННЯ (PT-7(1))}


\vspace{10pt}
{\footnotesize\bfseries
No: 1\\
Name: pt\_\allowbreak{}7\_\allowbreak{}1\_\allowbreak{}a\_\allowbreak{}01\\
Type: string\\
Default: nil
}

{\footnotesize При обробці системою номерів соціального страхування усувається непотрібний збір, зберігання та використання номерів соціального страхування}


{\footnotesize\bfseries
No: 2\\
Name: pt\_\allowbreak{}7\_\allowbreak{}1\_\allowbreak{}a\_\allowbreak{}02\\
Type: list\\
Default: [\textquotedbl{}admin\textquotedbl{}, \textquotedbl{}security\_\allowbreak{}officer\textquotedbl{}]
}

{\footnotesize Вивчаються альтернативи використанню номерів соціального страхування в якості персонального ідентифікатора, коли система обробляє їх}


{\footnotesize\bfseries
No: 3\\
Name: pt\_\allowbreak{}7\_\allowbreak{}1\_\allowbreak{}b\\
Type: list\\
Default: [\textquotedbl{}admin\textquotedbl{}, \textquotedbl{}security\_\allowbreak{}officer\textquotedbl{}]
}

{\footnotesize Не відмовляється система при обробці номерів соціального страхування в індивідуальних правах, пільгах або привілеях, передбачених законом, через відмову особи розкрити свій номер соціального страхування}


{\footnotesize\bfseries
No: 4\\
Name: pt\_\allowbreak{}7\_\allowbreak{}1\_\allowbreak{}c\_\allowbreak{}01\\
Type: list\\
Default: [\textquotedbl{}admin\textquotedbl{}, \textquotedbl{}security\_\allowbreak{}officer\textquotedbl{}]
}

{\footnotesize При обробці системою номерів соціального страхування кожну особу, яку просять розкрити свій номер соціального страхування, інформують про те, чи є таке розкриття обов'язковим чи добровільним, яким законодавчим чи іншим органом запитується такий номер, і як він буде використовуватися}


{\footnotesize\bfseries
No: 5\\
Name: pt\_\allowbreak{}7\_\allowbreak{}1\_\allowbreak{}c\_\allowbreak{}02\\
Type: list\\
Default: [\textquotedbl{}admin\textquotedbl{}, \textquotedbl{}security\_\allowbreak{}officer\textquotedbl{}]
}

{\footnotesize При обробці системою номерів соціального страхування кожну особу, яку просять розкрити свій номер соціального страхування, інформують про те, яким законодавчим чи іншим органом запитується цей номер}


{\footnotesize\bfseries
No: 6\\
Name: pt\_\allowbreak{}7\_\allowbreak{}1\_\allowbreak{}c\_\allowbreak{}03\\
Type: list\\
Default: [\textquotedbl{}admin\textquotedbl{}, \textquotedbl{}security\_\allowbreak{}officer\textquotedbl{}]
}

{\footnotesize При обробці системою номерів соціального страхування кожну особу, яку просять розкрити свій номер соціального страхування, інформують про те, як він буде використовуватися}




\subsubsection{ІНФОРМАЦІЯ ПРО ПЕРШУ ПОПРАВКУ (PT-7(2))}


\vspace{10pt}
{\footnotesize\bfseries
No: 1\\
Name: pt\_\allowbreak{}7\_\allowbreak{}2\_\allowbreak{}01\\
Type: list\\
Default: [\textquotedbl{}admin\textquotedbl{}, \textquotedbl{}security\_\allowbreak{}officer\textquotedbl{}]
}

{\footnotesize Заборонена обробка інформації, що описує, як будь-яка особа реалізує права, гарантовані Першою поправкою, за винятком випадків, коли це прямо дозволено законом або особою, або якщо вона не стосується та входить до сфери санкціонованої діяльності правоохоронних органів}




\subsection{ВИМОГИ ДО ВІДПОВІДНОСТІ (PT-8)}
Коли система чи організація обробляє інформацію з метою проведення програми відповідності необхідно:\\ 
a. отримати схвалення Ради з цілісності даних для проведення програми відповідності;\\ 
b. розробити та укласти договір комп'ютерної відповідності;\\ 
c. незалежним чином перевіряти інформацію, надану програмою відповідності, перш ніж вживати негативних заходів проти особи;\\ 
d. повідомляти осіб і надати їм можливість оскаржити висновки, перш ніж вживати проти них негативних заходів.

\vspace{10pt}
\textit{Немає параметрів для цього контролю.}
\vspace{10pt}


\end{document}
